% 昆虫记
% 昆虫记.tex

\documentclass[12pt,UTF8]{ctexbook}

% 设置纸张信息。
\usepackage[a4paper,twoside]{geometry}
\geometry{
	left=25mm,
	right=25mm,
	bottom=25.4mm,
	bindingoffset=10mm
}

% 设置字体，并解决显示难检字问题。
\xeCJKsetup{AutoFallBack=true}
\setCJKmainfont{SimSun}[BoldFont=SimHei, ItalicFont=KaiTi, FallBack=SimSun-ExtB]

% 目录 chapter 级别加点（.）。
\usepackage{titletoc}
\titlecontents{chapter}[0pt]{\vspace{3mm}\bf\addvspace{2pt}\filright}{\contentspush{\thecontentslabel\hspace{0.8em}}}{}{\titlerule*[8pt]{.}\contentspage}

% 设置 part 和 chapter 标题格式。
\ctexset{
	chapter/name={第,章},
	chapter/number={\arabic{chapter}}
}

% 设置古文原文格式。
\newenvironment{yuanwen}{\bfseries\zihao{4}}

% 设置署名格式。
\newenvironment{shuming}{\hfill\bfseries\zihao{4}}

% 注脚每页重新编号，避免编号过大。
\usepackage[perpage]{footmisc}

\title{\heiti\zihao{0} 昆虫记}
\author{J·H·法布尔[法]}
\date{}

\begin{document}

\maketitle
\tableofcontents

\frontmatter

\chapter{前言}



\mainmatter

\chapter{论祖传}

人人都有自己的才能和自己的性格。有的时候这种性格看起来好像是从我们的祖先
那里遗传下来的，然而要想再追究这些性格是来源于何处，却又是一件非常非常困难的
事情。
例如，有一天看到一个牧童，他正低声地数着一颗颗小石子，计算这些小石子的总
数，把这当做一种消遣，于是他长大后竟然成了十分著名的教授，最后，他也许可以成
为数学家。另外又有一个孩子，他的年龄比起别的小孩子们也大不了多少，别的孩子们
只注意玩闹的事情，然而他却不和别的小孩子们在一起玩儿，而是整日幻想一种乐器的
声音，于是当他独自一人的时候，竟听到一种神秘的合奏曲子了。可见这个小孩是很有
音乐天才的。第三个小孩，长得又小又瘦，年龄也很小，也许他吃面包和果酱时，还会
不小心涂到脸上，但他竟然有他独自的爱好——喜欢雕塑粘土，制成各种各样的小模型，
这些小模型被他雕塑得各具形态。如果这个小孩子运气好的话，他将来总有一天会成为
一名著名的雕刻家的。
我知道，在背后议论别人的私事，是十分让人讨厌的一种行为，但是我想也许大家
能允许我来讲一番，并借这个机会来介绍我自己和我的研究。
在我很小很小的时候，我已经有一种与自然界的事物接近的感觉。如果你认为我的
这种喜欢观察植物和昆虫的性格是从我的祖先那里遗传下来的，那简直是一个天大的笑
话，因为，我的祖先们都是没有受过教育的乡下佬，对其他的东西都一无所知。他们唯
一知道和关心的，就是他们自己养的牛和羊。在我的祖父辈之中，只有一个人翻过书本
儿，甚至就连他对于字母的拼法在我看来也是十分不可信的。至于如果要说到我曾经受
过什么专门的训练，那就更谈不上了，从小就没有老师教过我，更没有指导者，而且也
常常没有什么书可看。不过，我只是朝着我眼前的一个目标不停地走，这个目标就是有
朝一日在昆虫的历史上，多少加上几页我对昆虫的见解。
回忆过去，在很多年以前，那时候我还是一个不懂事的小孩子，那时我才刚刚学会
认字母，然而，我对于当时我那种初次学习的勇气和决心，至今都感到非常骄傲。
我记得很清楚的一次经历是我第一次去寻找鸟巢和第一次去采集野菌的情景，当时
那种高兴的心情真令我直到今天还难以忘怀。
记得有一天，我去攀登离我家很近的一座山。在这座山顶上，有一片很早就引起我
浓厚兴趣的树林，从我家的小窗子里看出去，可以看见这些树木朝天立着，在风中摇摆，
在雪里弯腰，我很早就想能有机会跑到这些树林那儿去看一看了。这一次的爬山，爬了
好长的时间，而我的腿又很短，所以爬的速度十分缓慢，草坡十分陡峭，就跟屋顶一样。
忽然，在我的脚下，我发现了一只十分可爱的小鸟。我猜想这只小鸟一定是从它藏
身的大石头上飞下来的。不到一会儿工夫，我就发现了这只小鸟的巢。这个鸟巢是用干
草和羽毛做成的，而且里面还排列着六个蛋。这些蛋具有美丽的纯蓝色，而且十分光亮，
这是我第一次找到鸟巢，是小鸟们带给我许多的快乐中的第一次。我简直高兴极了，于
是我伏在草地上，十分认真地观察它。
这时候，母鸟十分焦急的在石上飞来飞去，而且还“塔克！塔克！”地叫着，表现
出一种十分不安的样子。我当时年龄还太小，甚至还不能懂得它为什么那么痛苦，当时
我心里想出了一个计划，我首先带回去一只蓝色的蛋，作为纪念品。然后，过两星期后
再来，趁着这些小鸟还不能飞的时候，将它们拿走。我还算幸运，当我把蓝鸟蛋放在青
苔上，小心翼翼地走回家时，恰巧遇见了一位牧师。
他说“呵！一个萨克锡柯拉的蛋！你是从哪里捡到这只蛋的？”
我告诉他前前后后捡蛋的经历，并且说：“我打算再回去拿走其余的蛋，不过要等
到当新生出的小鸟们刚长出羽毛的时候。”
“哎，不许你那样做！”牧师叫了起来；“你不可以那么残忍，去抢那可怜母鸟的
孩子。现在你要做一个好孩子，答应我从此以后再也不要碰那个鸟巢。”
从这一番谈话当中，我懂得了两件事。第一件，偷鸟蛋是件残忍的事。第二件，鸟
兽同人类一样，它们各自都有各自的名字的。
于是我自己问自己道：“在树林里的，在草原上的，我的许多朋友，它们是叫什么
名字呢？萨克锡柯拉的意思是什么呢？”
几年以后，我才晓得萨克锡柯拉的意思是岩石中的居住者，那种下蓝色蛋的鸟是一
种被称为石鸟的鸟儿。
有一条小河沿着我们的村子旁边悄悄地流过，在河的对岸，有一座树林，全是光滑
笔直的树木，就像高高耸立的柱子一般，而且地上铺满了青苔。
在这座树林里，我第一次采集到了野菌。这野菌的形状，猛一眼看上去，就好像是
母鸡生在青苔上的蛋一样。还有许多别的种类的野菌形状不一，颜色也各不相同。有的
形状长得像小铃儿，有的形状长得像灯泡，有的形状像茶杯，还有些是破的，它们会流
出像牛奶一样的泪，有些当我踩到它们的时候，变成蓝蓝的颜色了。其中，有一种最稀
奇的，长得像梨一样，它们顶上有一个圆孔，大概是一种烟筒吧。我用指头在下面一戳，
会有一簇烟从烟筒里面喷出来，我把它们装满了好大一袋子，等到心情好的时候，我就
把它们弄得冒烟，直到后来它们缩成一种像火绒一样的东西为止。
在这以后，我又好几次回到这片有趣的树林。我在乌鸦队里，研究真菌学的初步功
课，通过这种采集所得到的一切，是呆在房子里不可能获得的。
在这种一边观察自然与一边做试验的方法相结合的情况之下，我的所有功课，除两
门课，差不多都学过了。我从别人那里，只学过两种科学性质的功课，而且在我的一生
中，也只有这两种：一种是解剖学，一种是化学。
第一种是我得力于造诣很深的自然科学家摩根·斯东，他教我如何在盛水的盆中看
蜗牛的内部结构。这门功课的时间很短，但是能学到很多东西。
我初次学习化学时，运气就比较差了。在一次实验中，玻璃瓶爆炸，使多数同学受
了伤，有一个人眼睛险些儿瞎了，老师的衣服也被烧成了碎片，教室的墙上沾污了许多
斑点。后来，我重新回到这间教室时，已经不是学生而是教师了，墙上的斑点却还留在
那里。这一次，我至少学到了一件事，就是以后我每做一种试验，总是让我的学生们离
开远一点。
我有一个最大的愿望，就是想在野外建立一个试验室。当时我还处于在为每天的面
包问题而发愁的生活状况下，这真是一件不容易办到的事情！我几乎四十年来都有这种
梦想，想拥有一块小小的土地，把土地的四面围起来，让它成为我私人所有的土地；寂
寞、荒凉、太阳曝晒、长满荆草，这些都是为黄蜂和蜜蜂所喜好的环境条件。在这里，
没有烦扰，我可以与我的朋友们，如猎蜂手，用一种难解的语言相互问答，这当中就包
含了不少观察与试验呢。
在这里，也没有长的旅行和远足，以至于白白浪费了时间与精力，这样我就可以时
时留心我的昆虫们了！
最后，我实现了我的愿望。在一个小村落的幽静之处，我得到了一小块土地。这是
一块哈麻司，这个名字是给我们洽布罗温司的一块不能耕种，而且有许多石子的地方起
的。那里除了一些百里香，很少有植物能够生长起来。如果花费功夫耕耘，是可以长出
东西的，可是实在又不值得。不过到了春天会有些羊群从那里走过，如果碰巧当时下点
雨，也是可以生长一些小草的。
然而，我自己专有的哈麻司，却有一些掺着石子的红土，并且曾经被人粗粗地耕种
过了。有人告诉我说，在这块地上生长过葡萄树，于是我心里真有几分懊恼，因为原来
的植物已经被人用二脚叉弄掉了，现在已经没有百里香了。百里香对于我也许有用，因
为可以用来做黄蜂和蜜蜂的猎场，所以我不得已又把它们重新种植起来。
这里长满了偃卧草、刺桐花、以及西班牙的牡莉植物——那是长满了橙黄色的花，
并且有硬爪般的花序的植物。在这些上面，盖着一层伊利里亚的棉蓟，它那耸然直立的
树枝干，有时长到六尺高，而且末梢还长着大大的粉红球，还带有小刺，真是武装齐备，
使得采集植物的人不知应从哪里下手摘取才好。在它们当中，有穗形的矢车菊，长了好
长一排钩子，悬钩子的嫩芽爬到了地上。假使你不穿上高筒皮鞋，就来到有这么多刺的
树林里，你就要因为你的粗心而受到惩罚了。
这就是我四十年来拼命奋斗得来的属于我的乐园啊！
在我的这个稀奇而又冷清的王国里，是无数蜜蜂和黄蜂的快乐的猎场，我从来没有
在单独的一块地方，看见过这么多的昆虫。各种生意都以这块地为中心，来了猎取各种
野味的猎人、泥土匠、纺织工人、切叶者、纸板制造者，同时也有石膏工人在拌和泥灰，
木匠在钻木头，矿工在掘地下隧道，以及牛的大肠膜(用来隔开金箔)工人，各种各样的
人都有。
快看啊！这里有一种会缝纫的蜜蜂。它剥下开有黄花底的刺桐的网状线，采集了一
团填充的东西，很骄傲地用它的腮(即颚)带走了。它准备到地下，用采来的这团东西储
藏蜜和卵。那里是一群切叶蜂，在它们的身躯下面，带着黑色的，白色的，或者血红色
的，切割用的毛刷，它们打算到邻近的小树林中，把树叶子割成圆形的小片用来包裹它
们的收获品。这里又是一群穿着黑丝绒衣的泥水匠蜂，它们是做水泥与沙石工作的。在
我的哈麻司里我们很容易在石头上发现它们工作用的工具。另外，这有一种野蜂，它把
窝巢藏在空蜗牛壳的盘梯里。还有一种，把它的蛴螬安置在干燥的悬钩子的秆子的木髓
里。第三种，利用干芦苇的沟道做它的家。至于第四种，住在泥水匠蜂的空隧道中，而
且连租金都用不着付。还有的蜜蜂生着角，有些蜜蜂后腿头上长着刷子，这些都是用来
收割的。
我的哈麻司的墙壁建筑好了，到处可以看到成堆成堆的石子和细沙，这些全是建筑
工人们堆弃下来的，并且不久就被各种住户给霸占了。泥水匠蜂选了个石头的缝隙，用
来做它们睡眠的地方。若是有凶悍的蜥蜴，一不小心压到它们的时候，它们就会去攻击
人和狗。它们挑选了一个洞穴，伏在那里等待路过的蜣螂。黑耳毛的鸫鸟，穿着白黑相
间的衣裳，看上去好像是黑衣僧，坐在石头顶上唱简单的歌曲。那些藏有天蓝色的小蛋
的鸟巢，会在石堆的什么地方才能找到呢？当石头被人搬动的时候，在石头里面生活的
那些小黑衣僧自然也一块儿被移动了。我对这些小黑衣僧感到十分惋惜，因为它们是很
可爱的小邻居。至于那个蜥蜴，我可不觉得它可爱，所以对于它的离开，我心里没有丝
毫的惋惜之情。
在沙土堆里，还隐藏了掘地蜂和猎蜂的群落，令我感到遗憾的是，这些可怜的掘地
蜂和猎蜂们后来无情地被建筑工人给无辜地驱逐走了。但是仍然还有一些猎户们留着，
它们成天忙忙碌碌，寻找小毛虫。还有一种长得很大的黄蜂，竟然胆大包天地敢去捕捉
毒蜘蛛，在哈麻司的泥土里，有许多这种相当利害的蜘蛛居住着。而且你可以看到，还
有强悍勇猛的蚂蚁，它们派遣出一个兵营的力量，排着长长的队伍，向战场出发，去猎
取它们强大的俘虏。
此外，在屋子附近的树林里面，住满了各种鸟雀。它们之中有的是唱歌鸟，有的是
绿莺，有的是麻雀，还有猫头鹰。在这片树林里有一个小池塘，池中住满了青蛙，五月
份到来的时候，它们就组成振耳欲聋的乐队。在居民之中，最最勇敢的要数黄蜂了，它
竟不经允许地霸占了我的屋子。在我的屋子门口，还居住着白腰蜂。每次当我要走进屋
子里的时候，我必须十分小心，不然就会踩到它们，破坏了它们开矿的工作。在关闭的
窗户里，泥水匠蜂在软沙石的墙上建筑土巢。我在窗户的木框上一不小心留下的小孔，
被它们利用来做门户。在百叶窗的边线上，少数几只迷了路的泥水匠蜂建筑起了蜂巢。
午饭时候一到，这些黄蜂就翩然来访，它们的目的，当然是想看看我的葡萄成熟了没有。
这些昆虫全都是我的伙伴，我的亲爱的小动物们，我从前和现在所熟识的朋友们，
它们全都住在这里，它们每天打猎，建筑窝巢，以及养活它们的家族。而且，假如我打
算移动一下住处，大山离我很近，到处都是野草莓树、岩蔷薇和石楠植物，黄蜂与蜜蜂
都是喜欢聚集在那里的。我有很多理由，使我为了乡村而逃避都市，来到西内南，做些
除杂草和灌溉莴苣的事情。



神秘的池塘



当我面对池塘，凝视着它的时候，我可从来都不觉得厌倦。在这个绿色的小小世界
里，不知道会有多少忙碌的小生命生生不息。在充满泥泞的池边，随处可见一堆堆黑色
的小蝌蚪在暖和的池水中嬉戏着，追逐着；有着红色肚皮的蝾螈也把它的宽尾巴像舵一
样地摇摆着，并缓缓地前进；在那芦苇草丛中，我们还可以找到一群群石蚕的幼虫，它
们各自将身体隐匿在一个枯枝做的小鞘中——这个小鞘是用来作防御天敌和各种各样意
想不到的灾难用的。
在池塘的深处，水甲虫在活泼地跳跃着，它的前翅的尖端带着一个气泡，这个气泡
是帮助它呼吸用的。它的胸下有一片胸翼，在阳光下闪闪发光，像佩带在一个威武的大
将军胸前的一块闪着银光的胸甲。在水面上，我们可以看到一堆闪着亮光的“蚌蛛”在
打着转，欢快地扭动着，不对，那不是“蚌蛛”，其实那是豉虫们在开舞会呢！离这儿
不远的地方，有一队池鳐正在向这边游来，它们那傍击式的泳姿，就像裁缝手中的缝针
那样迅速而有力。
在这个地方你还会见到水蝎，只见它交叉着两肢，在水面上悠闲地做出一副仰泳的
姿势，那神态，仿佛它是天底下最伟大的游泳好手。还有那蜻蜓的幼虫，穿着沾满泥巴
的外套，身体的后部有一个漏斗，每当它以极高的速度把漏斗里的水挤压出来的时候，
借着水的反作用力，它的身体就会以同样的高速冲向前方。
在池塘的底下，躺着许多沉静又稳重的贝壳动物。有时候，小小的田螺们会沿着池
底轻轻地、缓缓地爬到岸边，小心翼翼地慢慢张开它们沉沉的盖子，眨巴着眼睛，好奇
地展望这个美丽的水中乐园，同时又尽情地呼吸一些陆上空气；水蛭们伏在它们的征服
物上，不停地扭动着它们的身躯，一副得意洋洋的样子；成千上万的孑孓在水中有节奏
地一扭一曲，不久的将来它们会变成蚊子，成为人人喊打的坏蛋。
乍一看，这是一个停滞不动的池塘，虽然它的直径不超过几尺，可是在阳光的孕育
下，它却犹如一个辽阔神秘而又丰富多彩的世界。它多能打动和引发一个孩子的好奇心
啊！让我来告诉你，在我的记忆中的第一个池塘怎样深深地吸引了我，激发起我的好奇
心。
我小的时候，家里很穷。除了我妈妈继承的一所房子和一块小小的荒芜的园子之外，
几乎什么也没有了。“我们将怎么生活下去呢？”这个严重的问题，常常会挂在我爸爸
妈妈的嘴边。
你听说过“大拇指”的故事吗？那个“大拇指”藏在他父亲的矮凳子下，偷听他父
亲和母亲所说的一些关于生活窘迫的对话。我就很像那个“大拇指”。但是我没有像他
那样，可以藏在凳子底下，我是伏在桌子上一面假装睡着了，一面偷听他们的谈话。幸
运的是，我所听到的，并不像“大拇指”的父亲所说的那种使人心寒的话，相反地，那
是一个美妙的计划。我听了以后，心中涌起一阵难以形容的快乐和欣慰。
“如果我们来养一群小鸭”，妈妈说，“将来一定可以换得不少钱。我们可以买些
油脂回来，让亨利天天照料它们，把它们喂得肥肥的。”
“太好了！”父亲高兴地说道，“让我们来试试吧。”
那天晚上，我做了一个美妙的梦。我和一群可爱的小鸭子们一起漫步到池畔，它们
都穿着鲜黄色的衣裳，活泼地在水中打闹、洗澡。我在旁边微笑地看着它们洗澡，耐心
地等它们洗痛快，然后带着它们慢悠悠地走回家。半路上，我发现其中有一只小鸭累了，
就小心翼翼地把它捧起来放在篮子里面，让它甜甜地睡觉。
没想到就在两个月之后，我的美梦就实现了：我们家里养了二十四只毛茸茸的小鸭
子。鸭子自己不会孵蛋，常常由母鸡来孵。可怜的老母鸡分不出孵的是自己的亲骨肉还
是别家的“野孩子”，只要是那圆溜溜、和鸡蛋差不多样子的蛋，它都很乐意去孵，并
把孵出来的小生物当作自己的亲生孩子来对待。负责孵育我们家的小鸭的是两只黑母鸡，
其中一只是我们自己家的，而另一只是向邻居借来的。
我们家的那只黑母鸡，每天陪着小鸭们玩，不厌其烦地和它们做游戏玩耍，让它们
快乐健康地长大。我往一只木桶里盛了些水，大约有两寸高，这个木桶就成了小鸭们的
游泳池。只要是晴朗的日子，小鸭们总是一边沐浴着温暖的阳光，一边在木桶里洗澡嬉
戏，显得无比的美满、和谐和舒适，令旁边的黑母鸡羡慕不已。
两星期以后，这只小小的木桶渐渐地不能满足小鸭们的要求了。它们需要大量的水，
这样它们才能在里面自由自在地翻身跳跃，它们还需要许多小虾米、小螃蟹、小虫子之
类作为它们的食物。而这些食物通常大量地蕴藏在互相缠绕的水草中，等候着它们自己
去猎取。现在我感觉到取水是个大问题。因为我们家住在山上，而从山脚下带大量的水
上来是非常困难的。尤其是在夏天，我们自己都不能痛快地喝水，哪里还顾得了那些小
鸭呢？
虽然在我们家附近也有一口井，可那是一口半枯的井，每天要供四五家邻居轮流使
用，还有学校里的校长先生养的那头驴子，它总是贪得无厌地对着那井水大口大口地喝
水，那口井很快就被喝干了。直到整整一昼夜之后，才看见有井水渐渐地升起来，恢复
到原来的样子。在这么艰难的水荒中，我们可怜的小鸭子自然就没有自由嬉水的份了。
不过，在那山脚下，有一条潺潺的小溪。那倒是小鸭们的天然乐园。可是从我们家
到那小溪，必须穿过一条村里的小路，可是我们不能走那条小路，因为在那条路上我们
很可能会碰到几只凶恶的猫和狗，它们会毫不犹豫地冲散小鸭们的队伍，使我没法把它
们重新聚拢在一起。于是，我只得另谋出路。我想起在离山不远的地方，有一块很大的
草地和一个很不小的池塘。那是一个很荒凉很偏僻的地方，没有什么猫狗的打扰，的确
可以成为小鸭们的乐园。
我第一天做牧童，心中又快活又自在。不过有件事很令我难受，那就是赤裸裸的双
脚，渐渐地起泡了。因为跑了太多的路，我又不能把箱子里那双鞋子拿出来穿。那双宝
贵的鞋子只有在过节的时候才能穿。我赤裸的脚不停地在乱石杂草中奔跑，伤口越来越
大了。
小鸭们的脚似乎也受不了这么折腾，因为它们的蹼还没有完全长成，还远不够坚硬。
当它们走在这么崎呕的山路上时，不时地发出“呷呷——”的叫声，似乎是在请求我允
许它们休息一下。每当这个时候，我也只得满足它们的要求，招呼它们在树荫下歇歇脚，
否则恐怕它们再也没有力气走完剩下的路了。
我们终于到达了目的地。那池水浅浅的，温温的，水中露出的土丘就好像是一个个
小小的岛屿。小鸭们一到那儿就飞奔过去忙碌地在岸上寻找食物。吃饱喝足后，它们会
下到水里去洗澡。洗澡的时候，它们常常会把身体倒竖起来，前身埋在水里，尾巴指向
空中，仿佛在跳水中芭蕾。我美滋滋地欣赏着小鸭们优美的动作，看累了，就看看水中
别的景物。
那是什么？在泥土上，我看到有几段互相缠绕着的绳子又粗又松，黑沉沉的，像熏
满了的烟灰。如果你看到它，可能会以为它是从什么袜子上拆下来的绒线。于是我想：
可能是哪位牧羊女在水边编一只黑色的绒线袜子，突然发现某些地方漏了几针，不能往
下编了，埋怨了一阵子后，就决心全部拆掉，重新开始，而在她拆得不耐烦的时候，就
索性把这编坏的部分全丢在水里。这个推测看起来合情合理。
我走过去，想拾一段放到手掌里仔细观察，没想到这玩意儿又粘又滑，一下子就从
我的手指缝里滑走了。我花费了好大的劲，就是捉不住它，并且有几段绳子的结突然散
了，从里面跑出一颗颗小珠子，只有针尖般大小，后面拖着一条扁平的尾巴，我一下子
认出它们了，那是我所熟悉的一种动物的幼虫。它就是青蛙的幼虫——蝌蚪。
在这里我还看到了许多别的生物。其中有一种不停地在水面上打旋，它的黑色的背
部在阳光下发着亮光。每当我伸手去捉它们的时候，它们似乎早就预料到危险来临似的，
不等我碰它们，就逃得无影无踪了。我本想捉几个放到碗里面仔细研究，可惜就是捉不
到它们。
看啦！在那池水深处，有一团绿绿的、浓浓的水草，我轻轻拨开一束水草，看到立
刻有许多水珠争先恐后地浮到水面聚成一个大大的水泡。我想在这厚厚的水草底下一定
藏着什么奇怪的生物。我继续往下探索，看到了许多贝壳像豆子一样扁平，周围冒着几
个涡圈；有一种小虫看上去像戴了羽毛；还有一种小生物舞动着柔软的鳍片，像穿着华
丽的裙子在跳舞。我也不知道它们为什么这样不停地游来游去，也不晓得它们叫什么，
我只能出神地对着这个神秘、玄妙的水池，浮想联翩。
池水通过小小的渠道缓缓地流入附近的田地，那儿长着几棵赤杨，我又在那儿发现
了美丽的生物，那是一只甲虫，像核桃那么大，身上带着一些蓝色。那蓝色是如此的赏
心悦目，使我联想起了那天堂里美丽的天使，她的衣服一定也是这种美丽的蓝色。我怀
着虔诚的心情轻轻地捉起它，把它放进了一个空的蜗牛壳，用叶子把它塞好。我要把它
带回家中，细细欣赏一番。
接着我的注意力又被别的东西吸引住了。清澈又凉爽的泉水源源不断地从岩石上流
下来，不停地滋润着这个池塘。泉水先流到一个小小的潭里，然后汇成一条小溪。我看
着看着就突发奇想，觉得这样让溪水默默地流过就太可惜了。可以把它看作一个小小的
瀑布，去推动一个磨。于是，我就开始着手做一个小磨。我用稻草做成轴，用两个小石
块支着它，不一会儿就完工了。这个磨子做得很成功，只可惜当时没有小伙伴和我一起
玩，只有几只小鸭来欣赏我的杰作。
这个小小的成功大大地激发了我的创造欲望，一发不可收拾。我又计划筑一个小水
坝，那里有许多乱石可以利用，我耐心地挑选着可以用来筑坝的石块，挑着挑着，我忽
然发现了一个奇迹，它使我再也无心去继续建造水坝的事了。
当我打开一个大石头时，有一个小拳头那么大的窟窿，从窟窿里面发出一簇簇光环，
好像是一簇簇钻石的小面在阳光照耀下闪着耀眼的光，又好像是教堂里彩灯上垂下来的
一串串晶莹剔透的珠子。
多么灿烂而美丽的东西啊。它使我想起孩子们躺在打禾场的干草上所讲的神龙传奇
的故事。神龙是地下宝库的守护者，它们守护着不计其数的奇珍异宝。现在在我眼前闪
光的这些东西，会不会就是神话中所说的皇冠和首饰呢？难道它们就是蕴藏在这些砖石
中吗？在这些破碎的砖石中，我可以搜集到许多发光的碎石，这些都是神龙赐给我的珍
宝啊！我仿佛觉得神龙在召唤我，要给我数不清的金子。在潺潺的泉水下，我看见许多
金色的颗粒，它们都粘在一片细砂上。我俯下身子仔细观察，发现这些金粒在阳光下随
着泉水打着转，这真是金子吗？真是那可以用来制造二十法郎金币的金子吗？对一个贫
穷的家庭来说，这金币是多么的宝贵啊！
我轻轻地拣起一些细砂，放在手掌中。这发光的金粒数量很多，但是颗粒却很小，
得把麦杆用唾沫浸湿了，才能用来沾住它们。我不得不放弃这项麻烦的工作。我想一定
有一大块一大块的金子深藏在山石中，可以等到以后我来把山炸毁了再说，这些小金粒
太微不足道了，我才不去拣它们呢！
我继续把砖石打碎，看看里面还有什么，可是这下我看到的不是珠宝，我看到有一
条小虫从碎片里爬出来，它的身体是螺旋形的，带着一节一节的疤痕，像一条蜗牛在雨
天的古墙里蜿蜒着爬到墙外，那有节疤的地方显得格外沧桑和强壮。我不知道它们是怎
样钻进这些砖石内部的，也不知道它们钻进去干嘛。
为了纪念我发现的“宝藏”，再加上好奇心的驱使，我把砖石装在口袋里，塞得满
满的。这时候，天快黑了，小鸭们也吃饱了，于是我对它们说：“来，跟着我，我们得
回家了。”
我的脑海里装满了幻想，脚跟的疼痛早已忘记了。
在回去的路上，我尽情地想着我的蓝衣甲虫，像蜗牛一样的甲虫，还有那些神龙所
赐的宝物。可是一踏进家门，我就回过神来，父母的反应令我一下子很失望。他们看见
了我那膨胀的衣袋里面尽是一些没有用处的砖石，我的衣服也快被砖石撑破了。
“小鬼，我叫你看鸭子，你却自顾自地去玩耍，你捡那么多砖石回来，是不是还嫌
我们家周围的石头不够多啊？赶紧把这些东西扔出去！”父亲冲着我吼道。
我只好遵照父亲的命令，把我的那些珍宝、金粒、羊角的化石和天蓝色的甲虫统统
抛在门外的废石堆里，母亲看着我，无奈地叹了口气。
“孩子，你真让我为难。如果你带些青菜回来，我倒也不会责备你，那些东西至少
可以喂喂兔子，可这种碎石，只会把你的衣服撑破，这种毒虫只会把你的手刺伤，它们
究竟能给你什么好处呢？蠢货！准是什么东西把你迷住了！”
可怜的母亲，她说得不错，的确有一种东西把我迷住了——那是大自然的魔力。几
年后，我知道了那个池塘边的“钻石”其实是岩石的晶体；所谓的“金粒”，原来也不
过是云母而己，它们并不是什么神龙赐给我的宝物。尽管如此，对于我，那个池塘始终
保持着它的诱惑力，因为它充满了神秘，那些东西在我看来，其魅力远胜于钻石和黄金。

玻璃池塘

你有一处建在房子里面的小池塘吗？在那个小池塘里，你可以随时观察水中生物生
活的每一个片断。它没有像户外的池塘那么大，也没有太多的生物，可这些恰恰又为观
察提供了有利条件。除此之外，还不会有行人来打扰你专注的观察。其实这并不是什么
天方夜谈，这是很容易实现的。
我的户内池塘是在铁匠和木匠的合作下造成的：先用铁条做好池架，把它装在木头
做的基座上面。池上面盖着一块可以活动的木板，下面的池底是铁做的，底上有一个排
水的小洞。池的四周镶着玻璃。这是一个设计得相当不错的玻璃池，就放在我的窗口，
它的体积大约有十到十二加仑。
我先往池里放进一些滑腻腻的硬块。那是一种分量很重的东西，表面长着许多小孔，
看上去很像珊瑚礁。硬块上面盖着许多绿绿的绒毛般的苔藓，这苔藓能够使水保持清洁，
为什么呢？让我们来看一看吧。
动物在水池里和我们在空气中一样，要吸入新鲜的气味，同时，排出废气(二氧化碳)。
这些废气是不适宜人呼吸的。而植物刚好相反，它们吸入二氧化碳。所以池中的水草就
吸收这种不可以呼吸的废气，经过一番工作后，释放出可以供动物呼吸的氧气。
如果你在充满阳光的池边站一会儿，你就能观察到这种变化，在有水草的珊瑚礁上，
那一点点发亮的闪烁的星光，好像是绿苗遍地的草坪上点缀着的零零碎碎的珍珠。这些
珍珠不断地消逝，又接连不断地出现，它们会倏然在水面上飞散开来，好像水底下发生
了小小的爆炸，冒出一串串的气泡。
水草分解了水中的二氧化碳，得到炭元素，炭可以用来制造淀粉。淀粉是生物细胞
所不可缺少的东西。营养物水草所吐出来的废气是新鲜的氧气。这些氧气一部分溶解在
水中，供给水中的生物呼吸，一部分离开水面跑到空气中。你在外面看到的像珍珠一样
的气泡就是氧气！
我注视着池水中的气泡，作了一番遐想：在许多许多年以前，陆地刚刚脱离了海洋，
那时草是第一棵植物，它吐出第一口氧气，供给生物呼吸。于是各种各样的动物相继出
现了，而且一代一代繁衍、变化下去，一直形成今天的生物世界。我的玻璃池塘似乎在
告诉我一个行星航行在没有氧气的空间里的故事。
------------------








石蚕



我往我的玻璃池塘里放进一些小小的水生动物，它们叫石蚕。确切地说，它们是石
蚕蛾的幼虫，平时很巧妙地隐藏在一个个枯枝做的小鞘中。
石蚕原本是生长在泥潭沼泽中的芦苇丛里的。在许多时候，它依附在芦苇的断枝上，
随芦苇在水中漂泊。那小鞘就是它的活动房子，也可以说是它旅行时随身带的简易房子。
这活动房子其实可以算得上是一个很精巧的编织艺术品，它的材料是由那种被水浸
透后剥蚀、脱落下来的植物的根皮组成的。在筑巢的时候，石蚕用牙齿把这种根皮撕成
粗细适宜的纤维，然后把这些纤维巧妙地编成一个大小适中的小鞘，使它的身体能够恰
好藏在里面。有时候它也会利用极小的贝壳七拼八凑地拼成一个小鞘，就好像一件小小
的百纳衣；有时候，它也用米粒堆积起来。布置成一个象牙塔似的窝，这算是它最华丽
的住宅了。

暴徒的袭击

石蚕的小鞘不但是它的寓所，同时还是它的防御工具。我曾在我的玻璃池塘里看到
一幕有趣的战争，鲜明地证实了那个其貌不扬的小鞘的作用。
玻璃池塘的水中原本潜伏着一打水甲虫，它们游泳的姿态激起了我极大的兴趣。有
一天，我无意中撒下两把石蚕，正好被潜在石块旁的水甲虫看见了，它们立刻游到水面
上，迅速地抓住了石蚕的小鞘，里面的石蚕感觉到此次攻击来势凶猛，不易抵抗，就想
出了金蝉脱壳的妙计，不慌不忙地从小鞘里溜出来，一眨眼间就逃得无影无踪了。
野蛮的水甲虫还在继续凶狠地撕扯着小鞘，直到知道早已失去了想要的食物，受了
石蚕的骗，这才显出懊恼沮丧的神情，无限留恋又无可奈何地把空鞘丢下，去别处觅食
了。
可怜的水甲虫啊！它们永远也不会知道聪明的石蚕早已逃到石底下，重新建造它的
新鞘，准备着你们的下一次袭击了。

潜水艇——石蚕

石蚕靠着它们的小鞘在水中任意遨游，它们好像是一队潜水艇，一会儿上升，一会
儿下降，一会儿又神奇地停留在水中央。它们还能靠着那舵的摆动随意控制航行的方向。
我不由想到了木筏，石蚕的小鞘是不是有木筏那样的结构，或是有类似于浮囊作用
的装备，使它们不致于下沉呢？
我将石蚕的小鞘剥去，把它们分别放在水上。结果小鞘和石蚕都往下沉。这是为什
么呢？
原来，当石蚕在水底休息时，它把整个身子都塞在小鞘里。当它想浮到水面上时，
它先拖带着小鞘爬上芦梗，然后把前身伸出鞘外。这时的小鞘的后部就留出一段空隙，
石蚕靠着这一段空隙便可以顺利往上浮。就好像装了一个活塞，向外拉时就跟针筒里空
气柱的道理一样。这一段装着空气的鞘就像轮船上的救生圈一样，靠着里面的浮力，使
石蚕不致于下沉。所以石蚕不必牢牢地粘附在芦苇枝或水草上，它尽可以浮到水面上接
触阳光，也可以在水底尽情遨游。
不过，石蚕并不是十分擅长游泳的水手，它转身或拐弯的动作看上去很笨拙。这是
因为它只靠着那伸在鞘外的一段身体作为舵桨，再也没有别的辅助工具了，当它享受了
足够的阳光后，它就缩回前身，排出空气，渐渐向下沉落了。
我们人类有潜水艇，石蚕也有这样一个小小的潜水艇。它们能自由地升降，或者停
留在水中央——那就是当它们在慢慢地排出鞘内的空气的时候。虽然它们不懂人类博大
精深的物理学，可这只小小的鞘造得这样的完美，这样的精巧，完全是靠它们的本能。
大自然所支配的一切，永远是那么巧妙和谐。
------------------








蜣螂



一、圆球
蜣螂第一次被人们谈到，是在过去的六七千年以前。古代埃及的农民，在春天灌溉
农田的时候，常常看见一种肥肥的黑色的昆虫从他们身边经过，忙碌地向后推着一个圆
球似的东西。他们当然很惊讶地注意到了这个奇形怪状的旋转物体，像今日布罗温司的
农民那样。
从前埃及人想象这个圆球是地球的模型，蜣螂的动作与天上星球的运转相合。他们
以为这种甲虫具有这样多的天文学知识，因而是很神圣的，所以他们叫它“神圣的甲虫”。
同时他们又认为，甲虫抛在地上滚的球体，里面装的是卵子，小甲虫是从那里出来的。
但是事实上，这仅是它的食物储藏室而已。里面并没有卵子。
这圆球并不是什么可口的食品。因为甲虫的工作，是从土面上收集污物，这个球就
是它把路上与野外的垃圾，很仔细地搓卷起来形成的。
做成这个球的方法是这样的：在它扁平的头的前边，长着六只牙齿，它们排列成半
圆形，像一种弯形的钉把，用来掘割东西。甲虫用它们抛开它所不要的东西，收集起它
所选拣好的食物。它的弓形的前腿也是很有用的工具，因为它们非常的坚固，而且在外
端也长有五颗锯齿。所以，如果需要很大的力量去搬动一些障碍物，甲虫就利用它的臂。
它左右转动它有齿的臂，用一种有力的扫除法，扫出一块小小的面积。于是，在那堆集
起了它所耙集来的材料。然后，再放到四支后爪之间去推。这些腿是长而细的，特别是
最后的一对，形状略弯曲，前端还有尖的爪子。甲虫再用这后腿将材料压在身体下，搓
动、旋转，使它成为一个圆球形。一会儿，一粒小丸就增到胡桃那么大，不久又大到像
苹果一样。我曾见到有些贪吃的家伙，把圆球做到拳头那么大。
食物的圆球做成后，必须搬到适当的地方去。于是甲虫就开始旅行了。它用后腿抓
紧这个球，再用前腿行走，头向下俯着，臀部举起，向后退着走。把在后面堆着的物件，
轮流向左右推动。谁都以为它要拣一条平坦或不很倾斜的路走。但事实并非如此！它总
是走险峻的斜坡，攀登那些简直不可能上去的地方。这固执的家伙，偏要走这条路。这
个球，非常的重，一步一步艰苦的推上，万分留心，到了相当的高度，而且它常还是退
着走的。只要有一些不慎重的动作，劳力就全白费了：球滚落下去，连甲虫也被拖下来
了。再爬上去，结果再掉下来。它这样一回又一回地向上爬，一点儿小故障，就会前功
尽弃，一根草根能把它绊倒，一块滑石会使它失足。球和甲虫都跌下来，混在一起，有
时经过一、二十次的继续努力，才得到最后的成功。有时直到它的努力成为绝望，才会
跑回去另找平坦的路。
有的时候，蜣螂好像是一个善于合作的动物，而这种事情是常常发生的。当一个甲
虫的球已经做成，它离开它的同类，把收获品向后推动。一个将要开始工作的邻居，看
到这种情况，会忽然抛下工作，跑到这个滚动的球边上来，帮球主人一臂之力。它的帮
助当然是值得欢迎的。但它并不是真正的伙伴，而是一个强盗。要知道自己做成圆球是
需要苦工和忍耐力的！而偷一个已经做成的，或者到邻居家去吃顿饭，那就容易多了。
有的贼甲虫，用很狡猾的手段，有的简直施用武力呢！
有时候，一个盗贼从上面飞下来，猛地将球主人击倒。然后它自己蹲在球上，前腿
靠近胸口，静待抢夺的事情发生，预备互相争斗。如果球主人起来抢球，这个强盗就给
它一拳，从后面打下去。于是主人又爬起来，推摇这个球，球滚动了。强盗也许因此滚
落。那末，接着就是一场角力比赛。两个甲虫互相扯扭着，腿与腿相绞，关节与关节相
缠，它们角质的甲壳互相冲撞，摩擦，发出金属互相摩擦的声音，胜利的甲虫爬到球顶
上，贼甲虫失败几回被驱逐后，只有跑开去重新做自己的小弹丸。有几回，我看见第三
个甲虫出现，像强盗一样抢劫这个球。
但也有时候，贼竟会牺牲一些时间，利用狡猾的手段来行骗。它假装帮助这个被驱
者搬动食物，经过生满百里香的沙地，经过有深车轮印和险峻的地方，但实际上它用的
力却很少，它做的大多只是坐在球顶上观光，到了适宜于收藏的地点，主人就开始用它
边缘锐利的头，有齿的腿向下开掘，把沙土抛向后方，而这贼却抱住那球假装死了。土
穴越掘越深，工作的甲虫看不见了。即使有时它到地面上来看一看，球旁睡着的甲虫一
动不动，觉得很安心。但是主人离开的时间久了，那贼就乘这个机会，很快的将球推走，
同小偷怕被人捉住一样快。假使主人追上了它——这种偷盗行为被发现了——它就赶快
变更位置，看起来好像它是无辜的，因为球向斜坡滚下去了，它仅是想止住它啊！于是
两个“伙伴”又将球搬回，好像什么事情都没有发生一样。
假使那贼安然逃走了，主人艰苦做起来的东西，只有自认倒霉。它揩揩颊部，吸点
空气，飞走，重新另起炉灶。我颇羡慕而且嫉妒它这种百折不挠的品质。
最后，它的食品才平安的储藏好了。储藏室是在软土或沙土上掘成的土穴。做的如
拳头般大小，有短道通往地面，宽度恰好可以容纳圆球。食物推进去，它就坐在里面，
进出口用一些废物塞起来，圆球刚好塞满一屋子，肴馔从地面上一直堆到天花板。在食
物与墙壁之间留下一个很窄的小道，设筵人就坐在这里，至多两个，通常只是自己一个。
神圣甲虫昼夜宴饮，差不多一个礼拜或两个礼拜，没有一刻停止过。

二、梨
我已经说过，古代埃及人以为神圣甲虫的卵，是在我刚才叙述的圆球当中的。这个
已经我证明不是如此。关于甲虫被放卵的真实情形，有一天碰巧被我发现了。
我认识一个牧羊的小孩子，他在空闲的时候，常来帮助我。有一次，在六月的一个
礼拜日，他到我这里来，手里拿着一个奇怪的东西，看起来好像一只小梨，但已经失掉
新鲜的颜色，因腐朽而变成褐色。但摸上去很坚固，样子很好看，虽然原料似乎并没有
经过精细的筛选。他告诉我，这里面一定有一个卵，因为有一个同样的梨，掘地时被偶
然弄碎，里面藏有一粒像麦子一样大小的白色的卵。
第二天早晨，天色才刚刚亮的时候，我就同这位牧童出去考察这个事实。
一个神圣甲虫的地穴不久就被找到了，或者你也知道，它的土穴上面，总会有一堆
新鲜的泥土积在上面。我的同伴用我的小刀铲向地下拼命的掘，我则伏在地上，因为这
样容易看见有什么东西被掘出来。一个洞穴掘开，在潮湿的泥土里，我发现了一个精制
的梨。我真是不会忘记这是我第一次看见，一个母甲虫的奇异的工作呢！当挖掘古代埃
及遗物的时候，如果我发现这神圣甲虫是用翡翠雕刻的，我的兴奋却也不见得更大呢。
我们继续搜寻，于是发现了第二个土穴。这次母甲虫在梨的旁边，而且紧紧抱着这
只梨。这当然是在它未离开以前，完工毕事的举动，用不着怀疑，这个梨就是蜣螂的卵
子了。在这一个夏季，我至少发现了一百个这样的卵子。
像球一样的梨，是用人们丢弃在原野上的废物做成的，但是原料要比较精细些，为
的是给蛴螬预备好食物。当它从卵里跑出来的时候，还不能自己寻找食物，所以母亲将
它包在最适宜的食物里，它可以立刻大吃起来，不至于挨饿。
卵是被放在梨的比较狭窄的一端的。每个有生命的种子，无论植物或动物，都是需
要空气的，就是鸟蛋的壳上也分布着无数个小孔。假如蜣螂的卵是在梨的最后部分，它
就闷死了，因为这里的材料粘得很紧，还包有硬壳。所以母甲虫预备下一间精制透气的
小空间，薄薄的墙壁，给它的小蛴螬居住，在它生命最初的时候，甚至在梨的中央，也
有少许空气，当这些已经不够供给柔弱的小蛴螬消耗，它要到中央去吃食，已经很强壮，
能够自己支配一些空气了。
当然，梨子大的一头，包上硬壳子，也是有很好的理由的。蜣螂的地穴是极热的，
有时候温度竟达到沸点。这种食物，经过三四个礼拜之后，就会干燥，不能吃了。如果
第一餐不是柔软的食物，而是石子一般硬得可怕的东西，这可怜的幼虫就会因为没有东
西吃，而饿死了。在八月的时候，我就找到了许多这样的牺牲者，这苦东西烤在一个封
闭的炉内，要减少这种危险，母甲虫就拼命用它强健而肥胖的前臂，压那梨子的外层，
把它压成保护的硬皮，如同栗子的硬壳，用以抵抗外面的热度。在酷热的暑天，管家婆
会把面包摆在闭紧的锅里，保持它的新鲜。而昆虫也有自己的方法，实现同样的目的：
用压力打成锅子的样子来保藏家族的面包。
我曾经观察过甲虫在巢里工作，所以知道它是怎样做梨子的。
它收集建筑用的材料，把自己关闭在地下，可以专心从事当前的任务，这材料大概
是由两种方法得来的。照常例，在天然环境下，甲虫用常法搓成一个球推向适应的地点。
当推行的时候，表面已稍微有些坚硬，并且粘上了一些泥土和细沙，这在后来是很多见
的，不只在离收集材料很近的地方，可以寻找到用来储藏的场所，在这种情况下，它的
工作不过是捆扎材料，运进洞而已。后来的工作，却尤其显得稀奇。有一天，我见它把
一块不成形的材料隐藏到地穴中去了。第二天，我到达它的工作场地时，发现这位艺术
家正在工作，那块不成形的材料已成功的变成了一个梨，外形已经完全具备，而且是很
精致的做好了。
梨紧贴着地板的部分，已经敷上了细沙。其余的部分，也已磨光得像玻璃一样，这
表明它还没有把梨子细细的滚过，不过是塑成形状罢了。
它塑造这只梨时，用大足轻轻敲击，如同先前在日光下塑造圆球一样。
在我自己的工作室里，用大口玻璃瓶装满泥土，为母甲虫做成人工的地穴，并留下
一个小孔以便观察它的动作，因此它工作的各项程序我都可以看得见。
甲虫开始是做一个完整的球，然后环绕着梨做成一道圆环，加上压力，直至圆环成
为一条深沟，做成一个瓶颈似的样子。这样，球的一端就做出了一个凸起。在凸起的中
央，再加压力，做成一个火山口，即凹穴，边缘是很厚的，凹穴渐深，边缘也渐薄，最
后形成一个袋。它把袋的内部磨光，把卵产在当中，包袋的口上，即梨的尾端，再用一
束纤维塞住。
用这样粗糙的塞子封口是有理由的，别的部分甲虫都用腿重重的拍过，只有这里不
拍。因为卵的层端朝着封口，假如塞子重压深入，蛴螬就会感到痛苦。所以甲虫把口塞
住，却不把塞子撞下去。

三、甲虫的生长
甲虫在梨里面产卵约一个星期或十天之后，卵就孵化成蛴螬了，它毫不迟疑地开始
吃四周的墙壁，它聪明异常，因为它总是朝厚的方向去吃，不致把梨弄出小孔，使自己
从空隙里掉出来。不久它就变得很肥胖了，不过样子实在很难看，背上隆起，皮肤透明，
假如你拿它来朝着光亮看，能看见它的内部器官。如果是古代埃及人有机会看见这肥白
的蛴螬，在这种发育的状态之下，他们是不会猜想到将来甲虫会具有的那些庄严和美观
了。
当第一次脱皮时，这个小昆虫还未长成完全的甲虫，虽然全部甲虫的形状，已经能
辨别出来了。很少有昆虫能比这个小动物更美丽，翼盘在中央，像折叠的宽阔领带，前
臂位于头部之下。半透明的黄色如蜜的色彩，看来真如琥珀雕成的一般。它差不多有四
个星期保持这个状态，到后来，重新再脱掉一层皮。
这时候它的颜色是红白色，在变成檀木的黑色之前，它是要换好几回衣服的，颜色
渐黑，硬度渐强，直到披上角质的甲胄，才是完全长成的甲虫。
这些时候，它是在地底下梨形的巢穴里居住着的。它很渴望冲开硬壳的甲巢，跑到
日光里来。但它能否成功，是要依靠环境而定的。
它准备出来的时期，通常是在八月份。八月的天气，照例是一年之中最干燥而且最
炎热的。所以，如果没有雨水来软一软泥土，要想冲开硬壳，打破墙壁，仅凭这只昆虫
的力量，是办不到的，它是没有法子打破这坚固的墙壁的。因为最柔软的材料，也会变
成一种不能通过的坚壁，烧在夏天的火炉里，早已成为硬砖头了。
当然，我也曾做过这种试验，将干硬壳放在一个盒子里，保持其干燥，或早或迟，
听见盒子里有一种尖锐的摩擦声，这是囚徒用它们头上和前足的耙在那里刮墙壁，过了
两三天，似乎并没有什么进展。
于是我加入一些助力给它们中的一对，用小刀戳开一个墙眼，但这两个小动物也并
没有比其余的更有进步。
不到两星期，所有的壳内都沉寂了。这些用尽力量的囚徒，已经死了。
于是我又拿了一些同从前一样硬的壳，用湿布裹起来，放在瓶里，用木塞塞好，等
湿气浸透，才将里面的潮布拿开，重新放到瓶子里。这次试验完全成功，壳被潮湿浸软
后，遂被囚徒冲破。它勇敢地用腿支持身体，把背部当作一条杠杆，认准一点顶和撞，
最后，墙壁破裂成碎片。在每次试验中，甲虫都能从中解放出来。
在天然环境下，这些壳在地下的时候，情形也是一样的。当土壤被八月的太阳烤干，
硬得像砖头一样，这些昆虫要逃出牢狱，就不可能了。但偶尔下过一阵雨，硬壳回复从
前的松软，它们再用腿挣扎，用背推撞，这样就能得到自由。
刚出来的时候，它并不关心食物。这时它所最需要的，是享受日光。跑到太阳里，
一动不动地取暖。
一会儿，它就要吃了。没有人教它，它也会做，像它的前辈一样，去做一个食物的
球，也去掘一个储藏所，储藏食物，一点不用学习，它就完全会从事它的工作。
------------------








蝉



一、蝉和蚁
我们大多数人对于蝉的歌声，总是不大熟悉的，因为它是住在生有洋橄榄树的地方，
但是凡读过拉封敦的寓言的人，大概都记得蝉曾受过蚂蚁的嘲笑吧。虽然拉封敦并不是
谈到这个故事的第一人。
故事上说：整个夏天，蝉不做一点事情，只是终日唱歌，而蚂蚁则忙于储藏食物。
冬天来了，蝉为饥饿所驱，只有跑到它的邻居那里借一些粮食。结果他遭到了难堪的待
遇。
骄傲的蚂蚁问道：“你夏天为什么不收集一点儿食物呢？”蝉回答道：“夏天我歌
唱太忙了。”
“你唱歌吗？”蚂蚁不客气地回答：“好啊，那么你现在可以跳舞了”，然后它就
转身不理它了。
但在这个寓言中的昆虫，并不一定就是蝉，拉封敦所想的恐怕是螽斯，而英国常常
把螽斯译为蝉。
就是在我们村庄里，也没有一个农夫，会如此没常识地想象冬天会有蝉的存在。差
不多每个耕地的人，都熟悉这种昆虫的蛴螬，天气渐冷的时候，他们堆起洋橄榄树根的
泥土，随时可以掘出这些蛴螬。至少有十次以上，他见过这种蛴螬从土穴中爬出，紧紧
握住树枝，背上裂开，脱去它的皮，变成一只蝉。
这个寓言是造谣，蝉并不是乞丐，虽然它需要邻居们很多的照应。每到夏天，它成
阵地来到我的门外唱歌，在两棵高大筱悬木的绿荫中，从日出到日落，那粗鲁的乐声吵
得我头脑昏昏。这种振耳欲聋的合奏，这种无休无止的鼓噪，使人任何思想都想不出来
了。
有的时候，蝉与蚁也确实打一些交道，但是它们与前面寓言中所说的刚刚相反。蝉
并不靠别人生活。它从不到蚂蚁门前去求食，相反的倒是蚂蚁为饥饿所驱乞求哀恳这位
歌唱家。我不是说哀恳吗？这句话，还不确切，它是厚着脸皮去抢劫的。
七月时节，当我们这里的昆虫，为口渴所苦，失望地在已经枯萎的花上，跑来跑去
寻找饮料时，蝉则依然很舒服，不觉得痛苦。用它突出的嘴——一个精巧的吸管，尖利
如锥子，收藏在胸部——刺穿饮之不竭的圆桶。它坐在树的枝头，不停的唱歌，只要钻
通柔滑的树皮，里面有的是汁液，吸管插进桶孔，它就可饮个饱了。
如果稍许等一下，我们也许就可以看到它遭受到的意外的烦扰。因为邻近很多口渴
的昆虫，立刻发现了蝉的井里流出的浆汁，跑去舔食。这些昆虫大都是黄蜂、苍蝇、蛆
蜕、玫瑰虫等，而最多的却是蚂蚁。
身材小的想要到达这个井边，就偷偷从蝉的身底爬过，而主人却很大方地抬起身子，
让它们过去。大的昆虫，抢到一口，就赶紧跑开，走到邻近的枝头，当它再转回头来时，
胆子比从前变大来了，它忽然就成了强盗，想把蝉从井边赶走。
最坏的罪犯，要算蚂蚁了。我曾见过它们咬紧蝉的腿尖，拖住它的翅膀，爬上它的
后背，甚至有一次一个凶悍的强徒，竟当着我的面，抓住蝉的吸管，想把它拉掉。
最后，麻烦越来越多，无可奈何，这位歌唱家不得已抛开自己所做的井，悄然逃走
了。于是蚂蚁的目的达到，占有了这个井。不过这个井也干得很快，浆汁立刻被吃光了。
于是它再找机会去抢劫别的井，以图第二次的痛饮。
你看，真正的事实，不是与那个寓言相反吗？蚂蚁是顽强的乞丐，而勤苦的生产者
却是蝉呢！

二、蝉的地穴
我有很好的环境可以研究蝉的习惯，因为我是与它同住的。七月初，它就占据了靠
我屋子门前的那棵树。我是屋里的主人，门外就它是最高的统治者，不过它的统治无论
怎样总是不会让人觉得舒服。
蝉初次被发现是在夏至。在行人很多，有太阳光照着的道路上，有好些圆孔，与地
面相平，大小约如人的手指。在这些圆孔中，蝉的蛴螬从地底爬出来，在地面上变成完
全的蝉。它们喜欢特别干燥而阳光充沛的地方。因为蛴螬有一种有力的工具，能够刺透
焙过的泥土与沙石。
当我考察它们的储藏室时，我是用手斧来开掘的。
最使人注意的，就是这个约一寸口径的圆孔，四边一点尘埃都没有，也没有泥土堆
积在外面。大多数的掘地昆虫，例如金蜣，在它的窝巢外面总有一座土堆。蝉则不同，
是由于它们工作方法的不同。金蜣的工作是在洞口开始，所以把掘出来的废料堆积在地
面；但蝉蛴螬是从地底上来的。最后的工作，才是开辟门口的生路，因为当初并没有门，
所以它不是在门口堆积尘土的。
蝉的隧道大都是深达十五至十六寸，一直通行无阻，下面的部分较宽，但是在底端
却完全关闭起来。在做隧道时，泥土搬移到哪里去了呢？为什么墙壁不会崩裂下来呢？
谁都以为蝉是用了有爪的腿爬上爬下的，而这样却会将泥土弄塌了，把自己房子塞住。
其实，它的举措简直像矿工或是铁路工程师一样。矿工用支柱支持隧道，铁路工程
师利用砖墙使地道坚固。蝉的聪明同他们一样，它在隧道的墙上涂上水泥。这种粘液是
藏在它身子里的，用它来做灰泥，地穴常常建筑在含有汁液的植物须上的，它可以从这
些根须取得汁液。
能够很容易的在穴道内爬上爬下，对于它是很重要的，因为当它爬出去到日光下的
时候，它必须知道外面的气候如何。所以它要工作好几个星期，甚至一个月，才做成一
道坚固的墙壁，适宜于它上下爬行。在隧道的顶端，它留着手指厚的一层土，用以保护
并抵御外面空气的变化，直到最后的一霎那。只要有一些好天气的消息，它就爬上来，
利用顶上的薄盖，以便测知气候的状况。
假使它估计到外面有雨或风暴——当纤弱的蛴螬脱皮的时候，这是一件最重要的事
情——它就小心谨慎地溜到隧道底下。但是如果气候看来很温暖，它就用爪击碎天花板，
爬到地面上来了。
在他肿大的身体里面，有一种液汁，可以利用它避免穴里面的尘土。当它掘土的时
候，将液汁倒在泥土上，使它成为泥浆。于是墙壁就更加柔软了。蛴螬再用它肥重的身
体压上去，便把烂泥挤进干土的缝隙里。因此，当它在顶端出口处被发现时，身上常有
许多湿点。
蝉的蛴螬，初次出现在地面上时，常常在附近徘徊，寻找适当的地点脱掉身上的皮
——一棵小矮树，一丛百里香，一片野草叶，或者一枝灌木枝——找到后，它就爬上去，
用前足的爪紧紧地握住，丝毫不动。
于是它外层的皮开始由背上裂开，里面露出淡绿色的蝉。当时头先出来，接着是吸
管和前腿，最后是后腿与翅膀。此时，除掉身体的最后尖端，身体己完全蜕出了。
然后，它会表演一种奇怪的体操，身体腾起在空中，只有一点固着在旧皮上，翻转
身体，使头向下，花纹满布的翼，向外伸直，竭力张开。于是用一种差不多看不清的动
作，又尽力将身体翻上来，并且前爪钩住它的空皮，用这种运动，把身体的尖端从鞘中
脱出，全部的过程大约需要半个小时。
在短时期内，这个刚被释放的蝉，还不十分强壮。它那柔软的身体，在还没具有足
够的力气和漂亮的颜色以前，必须在日光和空气中好好地沐浴。它只用前爪挂在已脱下
的壳上，摇摆于微风中，依然很脆弱，依然是绿色的。直到棕色的色彩出现，才同平常
的蝉一样。假定它在早晨九点钟取得树枝，大概在十二点半，弃下它的皮飞去。那壳有
时挂在枝上有一两月之久。

三、蝉的音乐
蝉是非常喜欢唱歌的。它翼后的空腔里带有一种像钹一样的乐器。它还不满足，还
要在胸部安置一种响板，以增加声音的强度。的确，有种蝉，为了满足音乐的嗜好，牺
牲了很多。因为有这种巨大的响板，使得生命器官都无处安置，只得把它们压紧到身体
最小的角落里。当然了，要热心委身于音乐，那么只有缩小内部的器官，来安置乐器了。
但是不幸得很，它这样喜欢的音乐，对于别人，却完全不能引起兴趣。就是我也还
没有发现它唱歌的目的。通常的猜想以为它是在叫喊同伴，然而事实明显，这个意见是
错误的。
蝉与我比邻相守，到现在已有十五年了，每个夏天差不多有两个月之久，它们总不
离我的视线，而歌声也不离我的耳畔。我通常都看见它们在筱悬木的柔枝上，排成一列，
歌唱者和它的伴侣比肩而坐。吸管插到树皮里，动也不动地狂饮，夕阳西下，它们就沿
着树枝用慢而且稳的脚步，寻找温暖的地方。无论在饮水或行动时，它们从未停止过歌
唱。
所以这样看起来，它们并不是叫喊同伴，你想想看，如果你的同伴在你面前，你大
概不会费掉整月的功夫叫喊他们吧！
其实，照我想，便是蝉自己也听不见所唱的歌曲。不过是想用这种强硬的方法，强
迫他人去听而已。
它有非常清晰的视觉。它的五只眼睛，会告诉它左右以及上方有什么事情发生，只
要看到有谁跑来，它会立刻停止歌唱，悄然飞去。然而喧哗却不足以惊扰它。你尽管站
在它的背后讲话，吹哨子、拍手、撞石子。就是比这种声音更轻微，要是一只雀儿，虽
然没有看见你，应当早已惊慌得飞走了。这镇静的蝉却仍然继续发声，好像没事儿人一
样。
有一回，我借来两枝乡下人办喜事用的土铳，里面装满火药，就是最重要的喜庆事
也只要用这么多。我将它放在门外的筱悬木树下。我们很小心的把窗打开，以防玻璃被
震破。在头顶树枝上的蝉，看不见下面在于什么。
我们六个人等在下面，热心倾听头顶上的乐队会受到什么影响。“碰！”枪放出去，
声如霹雷。
一点没有受到影响，它仍然继续歌唱。它既没有表现出一点儿惊慌扰乱之状，声音
的质与量也没有一点轻微的改变。第二枪和第一枪一样，也没有发生影响。
我想，经过这次试验，我们可以确定，蝉是听不见的，好像一个极聋的聋子，它对
自己所发的声音是一点也感觉不到的！

四、蝉的卵
普通的蝉喜欢把卵产在干的细枝上，它选择最小的枝，粗细大都在枯草与铅笔之间。
这些小枝干，垂下的很少，常常向上翘起，并且差不多已经枯死了。
蝉找到适当的细树枝，即用胸部尖利的工具，把它刺上一排小孔——这样的孔好像
是用针斜刺下去的，把纤维撕裂，使其微微挑起。如果它不被打扰与损害，在一根枯枝
上，常常被刺成三十或四十个孔。
它的卵就产在这些小孔里，这些小穴是一种狭窄的小径，一个个的斜下去。每个小
穴内，普通的约有十个卵，所以总数约有三百或四百个。
这是一个蝉的很好的家族。然而它所以产这么多卵，其理由是为防御一种特别的危
险，必须要生产出大量的蛴螬，预备将会被毁坏掉一部分。经过多次的观察，我才知道
这种危险是什么。就是一种极小的蚋，拿它们的大小相比较，蝉简直是庞然大物呢！
蚋和蝉一样，也有穿刺工具，位于身体下面靠近中部的地方，伸出来时和身体成直
角。蝉卵刚产出，蚋立刻就会把它毁坏。这真是蝉的家族中的灾祸！大怪物只须一踏，
就可轧扁它们，然而它们竟镇静异常，毫无顾忌，置身于大怪物之前，真令人惊讶之至。
我曾见过三个蚋顺序地排列着，同时预备掠夺一个倒楣的蝉。
蝉刚装满一个小穴的卵，移到稍高处，另外做穴时，蚋立刻就会到那里去，虽然蝉
的爪可以够得着它，然而它却镇静而无恐，像在自己家里一样，它们在蝉卵之上，加刺
一个孔，将自己的卵产进去。蝉飞回去时，它的孔穴内，多数已加进了别人的卵，这些
冒充的家伙能把蝉的卵毁坏掉。这种成熟得很快的蛴螬——每个小穴内一个——即以蝉
卵为食，代替了蝉的家族。
几世纪的经验，这可怜的蝉的母亲仍一无所知。它的大而锐利的眼睛，并非看不见
这些可怕的恶人，鼓翼其旁。它当然知道有其它昆虫跟在后面，然而它仍然不为所动，
宁肯让自己做牺牲。它要轧碎这些坏种子是非常容易的，不过它竟不改变原来的本能，
解救它的家族，以免遭破坏。
从放大镜里，我曾见过蝉卵的孵化过程。开始很像极小的鱼，眼睛大而黑，身体下
面有一种鳍状物。由两个前腿连在一起组成。这种鳍有些运动力，可以帮助蛴螬冲出壳
外。并且帮它走出有纤维的树枝，而这恰恰是比较困难的事情。
鱼形蛴螬到穴外后，立刻把皮脱去。但脱下的皮会形成一种线，蛴螬依靠它附着在
树枝上。它在未落地以前，就在这里进行日光浴，用腿踢着，试试它的精力，有时则又
懒洋洋地在绳端摇摆。
等到触须自由了，可以左右挥动，腿可以伸缩，在前面的能够张合其爪，身体悬挂
着，只要有一点微风，就摇摆不定，在空气中翻跟斗。我所看到的昆虫中再没有比这个
更为奇观的了。
不久，它就落到地面上来。这个像跳蚤一般大小的小动物，在它的绳索上摇荡，以
防在硬地面上摔伤。身体渐渐地在空气中变硬。现在它开始该投入到严肃的实际生活中
去了。
此时，它当前仍有着千重危险。只要有一点儿风，就能把它吹到硬的岩石上，或车
辙的污水中，或不毛的黄沙上，或粘土上，硬得它不能钻下去。
这个弱小的动物，有如此迫切的需要藏身，所以必须立刻钻到地底下寻觅藏身之所。
天气是冷起来了，迟缓一些就有死亡的危险。它不得不四处寻找软土，毫无疑问，它们
之中有许多在没有找到合适的地方之前就死去了。
最后，它寻找到适当的地点，用前足的钩爬挖掘地面。从放大镜中，我看见它挥动
斧头向下掘，并将士抛出地面。几分钟后，土穴完成，这个小生物钻下去，埋藏了自己，
此后就再也看不见了。
未长成的蝉的地下生活，至今还是未发现的秘密，我们所知道的，只是它未成长爬
到地面上来以前，地下生活经过了许多时间而已，它的地下生活大概是四年。此后，日
光中的歌唱不到五个星期。
四年黑暗的苦工，一月日光中的享乐，这就是蝉的生活，我们不应厌恶它歌声中的
烦吵浮夸。因为它掘土四年，现在忽然穿起漂亮的衣服，长起与飞鸟可以匹敌的翅膀，
在温暖的日光中沐浴着。那种钹的声音能高到足以歌颂它的快乐，如此难得，而又如此
短暂。
------------------








泥水匠蜂



一、选择造屋的地点
有很多种昆虫都非常喜欢在我们的屋子旁边建筑它们的巢穴，在这些昆虫中最能够
引起人们兴趣的，要首推那种叫舍腰蜂的动物了。为什么呢？主要原因在于，舍腰蜂有
着十分美丽而动人的身材，非常聪明的头脑，还有一点应该注意的就是它那种非常奇怪
的窠巢。但是，知道舍腰蜂这种小昆虫的人却是很少的。甚至有的时候，它们住在某一
家人的火炉的旁边，但是，这户人家都对这个小邻居竟然一无所知。为什么呢？主要是
由于它那种天生下来就具备了的，安静，而且平和的本性。的确，这个小东西居住得十
分隐避，很难引起人们的注意。因此，连它自己的主人都不知道它就住在自己的家里，
算得上是自家成员之一。然而，讨厌吵闹，而且特别怕麻烦的人类，和这些隐避性很强
的小动物相比，要想使它出名，倒是件很容易就能达到的事情。现在，就让我来把这个
谦逊的、默默无闻的小动物，从不知名中提拔出来吧！
舍腰蜂是一种非常怕冷的动物。它搭建起自己的帐蓬，在那帮助橄榄树茁壮成长，
鼓励着蝉儿纵悦高歌的太阳光下建筑自己的安乐之居。甚至有的时候，为了它们整个家
族的需要，为了让大家都觉得比在阳光下更加温暖舒适一些，它们常常找到我们人类的
门上，要求和我们一起作伴。不用敲开人们的大门，询问一下主人是否同意它们和大家
同住在一个屋檐下，便自作主张，举家迁移进来，并且定居下来享受生活。舍腰蜂平常
的居所，主要是一些农夫们的单独的茅舍。在那茅屋的门外，大部分都生长着一些高大
挺拔的无花果树。这些果树的树荫遮盖着一口小小的水井。舍腰蜂在具体确定它的住所
的时候，主要会选择一个能够暴露在夏日里的炎热之下的地点，并且，如果有可能的话，
最好能够有一只大一点儿的火炉，还要有一些能够燃烧使用的柴火，这些条件对于舍腰
蜂而言都是必要的，不可缺少的。这是由它的天性所决定的。到了寒冷的冬天的夜晚，
火炉中喷射出来的温暖无比的火焰，对于它的选择，有着十分重要的影响力。因此，每
当看到从烟筒里面出来的黑炭，舍腰蜂就会欣喜若狂，因为它们知道那里便是一个可以
考虑选择的地方。因为，那里将会提供给它所必需的温暖与安逸。但是，相反的，要是
烟筒里面并没有什么黑炭的话，那么它是绝对不会信任这种地方的，也绝对不会选择这
样的地方来建筑自己的家。因为舍腰蜂会利用它的头脑作出判断，这问屋子里的主人们
一定是在里面忍受着饥寒交迫的悲惨境域。
在七、八月里的大暑天中，这位小客人，忽然出现了。它在找寻着适合它做窠的地
点。舍腰蜂一点儿也不为这间屋子里面的一切喧闹行为所惊动和扰乱。而住在屋子里的
人们也一点儿都注意不到它。他们互相都没有注意到，因此也就互无干扰了。舍腰蜂只
不过在有的时候，利用它那尖锐的目光，有的时候，又利用它那灵敏十足的触须，视察
一下已经变得乌黑的天花板、木缝、烟筒等。但是，特别受到它关注的是火炉的旁边。
这是它从不轻易放过的地方。甚至，它连烟筒内部都要仔仔细细地视察一遍。它可是一
种细致入微的小动物，一旦视察工作完毕，并且已经决定了建巢的地点以后，它们便立
即飞走了。然后，不久就会带着少量的泥土又飞回来，开始建筑它的房子的底层了。于
是，筑造家园的工作便正式破土动工了。
舍腰蜂所选择的地点各不相同，也是非常奇怪的一个特点。炉子内部的温度最适合
那些小蜂了，因此，舍腰蜂所中意的地点，至少得是烟筒内部的两侧，其高度大约是二
十寸或者差不多的地方。不过，尽管这个地点可以说是一个非常舒服的藏身之妙处，但
是，世上没有十分完善的东西，它也有相当不少的缺点。由于巢是建在烟筒的内部的，
那么自然便会有烟在里面。如果烟要是喷到蜂巢上面，那么，巢中的舍腰蜂就会被“污
染”了，会被弄成棕色的或者是黑色的，就好像烟筒里被熏过的砖石一样。假使火炉里
的火焰烧不到蜂巢，那还不是一件最要紧的事。最重要的事是小黄蜂有可能会被闷死在
粘土罐子里。不过，不用替它们担心，它们的母亲似乎早就已经知道这些事情了，因为
这位母亲总是把它自己的家族安排在烟筒的适当位置上。它们选定的位置非常宽大。在
那个地方，除了烟灰以外，其它的东西都是很难到达的。
虽然舍腰蜂样样都当心，时刻都仔细、谨慎。但是“智者千虑，必有一失”。它如
此地认真，但还是有一件很危险的事情在等待着它们。这件事有的时候会发生，那就是
当舍腰蜂正在建造它的房屋的时候，如果在这个关健时刻，有一阵蒸汽或者是烟幕的侵
扰，那么，它刚刚造成的一半的房子，便不得不半途而废。于是，它们要么停工一些时
候，要么就全日停工不干。特别是在这家的主人在煮、洗衣服的日子里，这种事情发生
的可能性最大，危险性也最大。一天从早到晚，大盆子里不停地滚沸着，炉灶里的烟灰、
大盆和木桶里面的大量蒸汽，一起混合成为浓厚的云雾。这给蜂巢造成了严重的威胁。
这个时候舍腰蜂就会面临着家毁人亡的危险。
我以前曾经听别人说过，河鸟在回巢的时候，总是要飞过水坝下的大瀑布。这一点
听起来会让人觉得河鸟已经算得上是一种相当有勇气、有胆量的小动物了。但是，与之
相比的舍腰蜂也毫不示弱，甚至，它的勇敢已经超过河鸟。它在回巢的时候，牙齿间总
是要含着一块用于建造它的巢穴的泥土。要想到达它的施工工地，它要从浓厚的烟灰的
云雾中穿越过去。但是，那层烟幕简直太厚重了，舍腰蜂冲进去以后，就完全都看不见
它那小小的身影了。虽然看不见它那小小的躯体，但是能够听见一阵不太规则的呜呜的
声音。这是什么声音呢？这不是别的什么声音，这是它在一边工作，一边低唱的歌声。
因此，我们可以断定，舍腰蜂肯定还呆在里面，而且它很快乐，高高兴兴地从事着它的
本职工作，不知劳苦地建筑着它自己的住所。看得出来，它对自己的劳动很满意，也很
乐意从事这项工作，在这层厚厚的云雾里，它很神秘地进行着它自己的工作。忽然，低
低的劳动之歌停止了。不一会儿它飞出来了，从那层充满神秘色彩的浓雾里飞出来了，
它安然无恙，什么伤也没有得。毕竟这是它的本能嘛！差不多每天它都要经过很多次这
种十分危险的事情，直到它把巢最终建好，把食物都储藏好，最后把自家的大门关上为
止。然后，它才休息一下。这个小东西为了自己的家园也真够不辞辛苦的了！
每一次，只有我一个人能够看到舍腰蜂在我的炉灶里不停地忙碌着，建造住所，储
备食物。这大概是因为我比较细心。记得我第一次看到它们的时候，是有一天我在煮、
洗衣服的时候。本来，那个时候，我是在爱维浓(Avignon)学院里教书的。那天，时间已
经将近两点钟了，几分钟之内，外面就会敲鼓催促我去给羊毛工人们做演讲了。就在这
个时候，忽然，我看见了一个非常奇怪而且轻灵的小昆虫。它从由木桶里升腾起来的蒸
汽中穿飞出来。这只小动物的身体很有意思，当中的部分非常的瘦小，但是后部却是非
常肥大的，而这两个部分之间，竟然是由一根长线连接起来的。多么奇妙的小东西啊！
这个小东西就是舍腰蜂，这是我第一次没有用观察的眼光来看它。于是，便有了第一印
象。
在初次相识之后，我对家里的这个小客人一直报有非常浓厚的兴趣。我非常热心地
希望能和这个小不点儿客人互相熟识，作一些交流。于是，我便嘱咐我的家人，在我不
在家的时候，不要去主动打扰它们，破坏它们的正常生活。瞧，我多么注意保护这个没
有受到邀请的不速之客呀！事情发展的良好态势已然胜过了我所希望的那样。当我回到
家里的时候，发现它一点儿也没有受到什么打扰，而且一个个都安然无恙。它仍然呆在
蒸汽的后面，努力地进行着它自己的工作，为自己的家而辛苦。由于我想要观察一下舍
腰蜂的建筑以及它的建筑才能，还有它的食物的性质，以及幼小的黄蜂的进化及其生长
过程等等，因此，我把炉灶中的火焰给弄灭了。我这样做的目的主要是减少烟灰的量。
差不多将近两小时，我非常仔细地注视着它。
但是，从这以后，不知道是什么原因，差不多将近四十年来，我的屋子里，再也没
有这样小的客人光临了，一点儿也见不到它们的踪影了。有关舍腰蜂的进一步的知识，
我还是从我的邻居家的炉灶旁边的蜂巢里得出来的。
通过细心观察我发现，在这个小小的动物身上，有一种十分孤癖的流浪的习性。这
一点使得它和其它大多数黄蜂，以及蜜蜂是不尽相同的。一般情况下，它总是选择好一
个地点，自己筑起一个显得特别孤独的巢穴。同时，在舍腰蜂自己养活自己的地方，是
很少能见到它自己家族的成员及亲属的。在距离我们城南不远的地方，经常可以看到这
种小动物。但是，这个小东西，宁愿挑选农民那充满烟灰的屋子里的炉灶来筑造自己的
小家，也不喜欢那些城镇居民的雪白的别墅里的炉灶。我所到过的任何地方所看到的舍
腰蜂，都没有像我们村里这么多的。与此同时，我们村里的屋子都很有特点。我们村上
的茅屋都有一定的倾斜性，而且茅屋都被日光晒成了黄色，这使得它们看上去都很有特
色。
事实是很明显的，泥水匠蜂选择烟筒做为自己的住所。这一点是不用置疑的了。但
是，它之所以为自己选择这样一个地方，倒并不是意味着它贪图安逸与享乐。因为，很
显然，选择这样的地方可不是什么特舒服的地方。这种地方更需要这种小动物加倍地努
力，能够具备更好的才能。而且，在这种地方工作，是有很大危险性的。因为时常有险
情发生，需要冒一定的危险，甚至是生命的危险。从这一点来看，说它选择烟筒建巢是
为了自己的安逸，那可真的要大大的冤枉了我们这位小客人了。它选择这样的地点来筑
巢建穴，主要意图还完全是为了它的整个家族来考虑的，而并非出于私利。它不希望只
是自己舒服就可以了，应该是大家共同享福，共同舒适，那才是它们真正要达到的目的。
因而可以说，舍腰蜂还是一种比较热爱家庭的动物，家庭责任感很强。当然了，舍腰蜂
选择烟筒还有一个很重要的原因，那就是舍腰蜂及它的家族成员对温度的要求比较高，
这主要是由于本能的原因，它们的住所必须建在很温暖的地方，这一点和其它的黄蜂、
蜜蜂是很不相同的。
我记得有一次去一家丝厂，在那里我见到过一个舍腰蜂的巢。它把自己的巢建在机
房里，并且为自己选择了刚好是在大锅炉的上面的天花板上的一个地方。看来，它真是
慧眼独具啊！它为自己选择的这个地点，整个一年，无论寒暑，也无论春夏秋冬的变迁，
温度计所显示出的温度，总是不变的一百二十度，只是要除去晚上的时间，还有那些放
假的日子。很显然，在这些日子里，锅炉里并没有加热，所以，温度当然会随之有所变
化的。这个事实很明显地告诉我们，这个小小的动物对温度真是要求很高啊！而且，它
也是个非常会为自己挑选地点的家伙。
还有，在乡下的那些蒸酒的屋子里，我也曾经不止一次地看到过许多舍腰蜂的巢穴。
而且，凡是那些可以选择的、方便它们安居与行动的地方，都已经被它们占满了。甚至，
连那些帐薄堆积的地方，都被它们占据下来了。蒸酒房里的温度，和刚才提到的丝厂里
的温度相差得不是太多。大约有一百一十三度左右。这些温度字再次告诉我们，这种泥
水匠蜂甚至足可以在那种使油棕树生长的热度下生存。
这样看来，锅，还有炉灶，当然也就很自然地成了舍腰蜂最理想的家和首选对象了。
但是，除了这些首屈一指的地方以外，舍腰蜂也不厌弃一些其它可以选择的地点。它非
常希望居住在任何可以让它觉得舒适、安逸的角落里面。比如说，在养花房里，在厨房
的天花板上，可关闭窗户的凹进去的地方，还有就是茅舍中卧室的墙上等等。至于建造
自己窠巢的地基，这一点，它并不放在心上。为什么呢？因为，在平常它的多孔的巢穴，
一般都是建筑在石壁或者是木头上的。这些地方相对而言，还是比较坚实的。因而，它
们似乎并不是很关心房屋的基础。不过，也有的时候我曾经看到过它把自己的巢筑在葫
芦的内部，或者在皮帽子里，砖的缝隙之中，或者是装麦子用的空袋子里，还有的时候，
它建巢在铅管里面。
记得有一次，我在接近学院的一个农夫的家里所看到的事情，更让人觉得特别的新
奇。在这个农夫的家里，有一个特别宽大的炉灶。在装有这个炉灶的大房间里，在炉灶
上的一排锅里，正煮着农工们要喝的汤，还有一些供牲畜们食用的东西。过了一会儿，
工人们都从田地里收工回家了。累了一天，他们的肚子肯定饿坏了。回来后，他们便迫
不及待地，不声不响地，在一边非常迅速地吞食着他们的食品及汤。他们为了要享受这
休工用饭的大约半小时的舒适，干脆摘下了戴在头上妨碍吃饭的帽子，随后也脱去了他
们的上衣，随手把它们挂在一个木钉的上面。这吃饭的时间，对于农工们而言，虽然是
短暂的，但是，要是让泥水匠蜂去占据工人们刚刚脱下的衣物，却又是绰绰有余的了。
在这些衣物中、草帽里边，被它们视为最合适的地方，它们抢先去占领它。那些上衣的
褶缝，则被视为最佳的地点。与此同时，泥水匠蜂的建筑工作也就马上破土动工。这时，
一个工人已经吃完了他的饭，从饭桌旁边站了起来，抖了抖他自己的衣服。另外一个人
也站起来，走了过来，摘下自己的草帽，也抖了一下。这样几下抖动便去掉了舍腰蜂刚
刚初具规模的窠巢，就是在这个时候，在这么短暂的时间里，它的蜂巢居然已经有一个
橡树果子那样大了，真让人始料不及。它们可真是一些让人惊奇的小动物。
那个农夫家，有一位专门烹调食物的女人。她对于泥水匠蜂这种动物可是一点儿好
感也没有。她报怨说这些可恶的小东西常常跑出来，弄脏了许多的东西。天花板、墙壁，
还有烟筒上，经常被涂满了泥，非常烦人，打扫起来很费力气的。但是，在衣服和窗幔
上，情况就大不相同了。这个女人每天都会用一根竹子，使劲地敲打窗幔，以保持它的
清洁。所以，在这些地方情况会稍好一些，略微干净一些。但是，驱逐这些扰人的小动
物是多么地不容易啊！赶走了一次，第二天早晨它又会一样地跑回来来做巢。它可真是
个执着的小家伙，总是不厌其烦地从事着它的本能工作。

二、它的建筑物
事实上，我也非常同情这个农家厨役，很能理解她的烦恼。但是，我同时感到遗憾
的是，我不能代替她的位置。对此，我无能为力。如果，我能够以某种力量，使得这种
小动物安安静静地固定地在某一稳定的地点建屋居住，那该有多好啊，我肯定会特别高
兴的。这样一来既便它把家俱弄满了泥土，那也是不碍事的！我更希望能够知道它的那
种巢的命运。如果这个巢是做在不太稳固的东西上，比如，在衣服上，或是在窗幔上，
那么它们该怎么办呢？
泥水匠蜂的窠巢是利用硬的灰泥制做而成的。一般它的巢都围绕在树枝的四周。由
于是灰泥组成的，所以它就能够非常坚固地附着在上面。但是，泥水匠蜂的窠巢，只是
用泥土做成的，没有加水泥，或者是其它什么更能让它坚固的基础。那么，它怎么解决
这些问题呢？
建筑上的材料，并没有什么特殊的。只是潮湿的泥土，从那种湿地上取来的。因此，
河边的粘土是最合适的选择。但是，在我们这样一个多沙石的村庄里面，河道非常的少。
然而在我自己的小园子里，我在种植蔬菜的区域里，挖掘了一些小沟渠，以便更好地种
植。因此，有的时候，有一点儿水，便会整天在沟里流。因而，这里便经常会有舍腰蜂
的身影出没。它们在这里选择适宜的泥土，于是在无事可做的时候，我就可以观察这些
建筑家了。这里倒是一个很好的观察地点。
临近沟渠的时候，它当然就会注意到这件可喜的事情，于是就匆匆忙忙地跑过来取
水边这一点点十分宝贵的泥土。它们不肯轻意放过这没有湿气的时节极为珍稀的发现。
那么它们是怎样掘取这里的泥土的呢？它们用下颚刮取沟渠旁边那层表面光滑的泥，足
直立起来，双翼还振动着，把它那黑色的身体抬举得相当的高。我的管家妇在这泥土的
旁边做工。她把她的裙子非常小心谨慎地提起来，以免弄脏了。但是事实上，却很少能
够不沾上污渍。可是这样一群不停地搬取着泥土的黄蜂，原本应该是很脏的，但是事实
上它们的身上竟然连一点儿泥迹都没有。之所以会这样，它们自然有它们自己聪明的办
法。它们会把身子提起来，这样就能使它们全身上下一点儿泥污也沾染不上。除去它们
的足尖以及用于工作的下颚之外，其它的地方都看不到什么泥迹之类的脏东西。
这样，用不了多长时间，一个泥球就制做成功了。差不多能有豌豆那么大。然后，
泥水匠蜂会用牙齿把它衔住，飞回去，在它自己的建筑物上再增加上一层。这项工作完
成以后，它歇也不歇一下，便继续投入新的工作之中。接着飞回来，再做第二个泥球。
在一天中，天气最为炎热的时候，只要那片泥土未干，仍然是潮湿的，那么，泥水匠蜂
的工作就会不停地坚持下去。
除了我这园中的小小的沟渠边这片潮湿的泥土以外，在村子里，最好的地点，就算
是村里的人牵着驴子去饮水的那片泉水旁边了。在这个地方，无论什么时候都有潮湿的
黑色的烂泥。哪怕是那种最热的太阳，最强烈的风，都不可能把这片泥土吹干。这种泥
泞不堪的地方，对于那些走路的人来说，是非常不方便的，也是极不受欢迎的。然而，
舍腰蜂却不一样的。它非常喜欢到这个地方来，因为这里的泥土质量非常好，它也很喜
欢在驴子的蹄旁做小泥丸。每次它都会有丰富的收获。
和泥水匠蜂这位粘土建筑家不一样，黄蜂并不把泥土先做成水泥，它就这样把现成
的泥土拿走，直接应用于建筑。所以，黄蜂的巢建造得很不结实，更不稳定，完全不能
抵挡气候的千变万化。只要有一点儿水滴落上去，蜂巢就会变软，变成了和原来一样的
泥土。要是有一阵狂风大雨的话，它的巢穴就会被打成泥浆。这主要是因为，这种蜂巢
实际上只不过是由干了的烂泥做成的，一旦浸了水以后，就会马上变成和原来一样的软
泥，自然巢穴也就不复存在了，还须再次辛苦地重建家园。
事实是很显然的，即便是幼小的舍腰蜂一点儿也不惧怕寒冷，不怕雨水把蜂巢打得
粉碎，那蜂巢也必须建在避雨的地方。这就是为什么这种小动物喜欢选择人类居住的屋
子，特别是选择温暖的烟筒里面来建筑自己的住所的缘故。看来，安全是很重要的。
在最后一项装饰工作——那遮盖起它辛苦制造的建筑的各层——还没有完全成功之
前，舍腰蜂的窠巢确实具有一种非常自然的美感。它有一些小窠穴，有的时候它们互相
并列成一排，那种形状有一点儿像口琴。不过，那些小巢穴，还是以那种互相堆叠起来
成层的居多。有的时候，数一下有十五个小巢穴；有的时候，有十个；有时，又减少至
三、四个，甚至仅有一个。
舍腰蜂的巢穴的形状和一个圆筒子差不多。它的口稍微有点儿大，底部又稍小一些。
大的有一寸多长，半寸多宽，蜂巢有一个非常别致的表面，它是经过了非常仔细的粉饰
而形成的。在这个表面上，有一列线状的凸起围绕在它的四周，就好像金线带子上的线
一样。每一条线，就是建筑物上的一层。这些线的形状，是由于用泥土盖起每一层已经
造好的巢穴而显露出来的。数一数它们，就可以知道，在黄蜂建筑它的时候，来回旅行
了一共有多少次。它们通常是十五到二十层之间。每一个巢穴，这位辛辛苦苦的不辞辛
劳的建筑家在建筑它时，大概须用二十次来往返复搬运材料。可见，它们有多少勤劳！
蜂窠的口当然是朝着上面的。如果一个罐子的口是朝下的，那么，它还能盛下什么
东西呢？当然什么也盛不下了。道理也便就在这里。黄蜂的窠穴，也并不是什么特殊的
东西，不过就像一个罐子而已，其中预备盛储的食物便是：一堆小蜘蛛。
这些巢穴，一一建造好了以后，黄蜂便往里面塞满了蜘蛛。等它们自己产下卵以后，
便把它们全部封闭好。但是，这时候，它依然保存着美观的外表。这种外表一直要保持
到黄蜂认为巢穴的数量已经足够多了的时候为止。于是，黄蜂会把整个巢穴的四周，再
堆上一层泥土，以便使它能够更加坚固一些，从而可以起到保护的作用。这一次黄蜂在
工作时，也不进行什么周密的计算了。因此，它做得特别不精巧，更不像从前做巢那样，
铺加以相当的修饰之物。黄蜂能带回多少泥土，就往上面堆积多少泥土。只要能够堆积
得上去就可以了，再没有更多的修补、装璜的动作了。泥土一旦取了回来，便堆放到原
来的巢穴上去。然后，就那么非常漫不经心地轻轻地敲几下，使这些泥土可以铺开。这
一层包裹物质，一下子把建筑物的美观统统都掩盖住了。这最后一道工序完成以后，蜂
巢的最后形状就形成了。此时此刻的蜂巢就好像是一堆泥，一堆人们抛掷到墙壁上的泥。

三、它的食物
现在，我们都已经很清楚这个装食物的罐子是怎样形成的了。接下来，我们必须知
道的是，在这个罐子里边，究竟都隐藏了一些什么东西。
幼小的舍腰蜂，是以各种各样的蜘蛛作为食物的。甚至，在同一窠巢中，其食品的
形状各个都不相同，因为，各种各样的蜘蛛，都可以充当食品，只是个头不要过大。否
则就装不到罐子里去了。在幼蜂的各种食品中，那种后背上有三个交叉着的白点的十字
蜘蛛，是最为常见的美味佳肴。这其中的理由，我觉得应该是很简单的。因为，黄蜂不
是那种跑到离家很远的地方去千里迢迢地捕猎食物的动物。它只不过经常在住所的附近
地区游猎而已。而在它的住宅的近区内，这种有交叉纹的蜘蛛是最容易寻找得到的。
对于幼蜂而言，那种生长着毒爪的蜘蛛，要算是最最危险的野味儿了。假使蜘蛛的
身体特别的大，就需要黄蜂拥有更大的勇气和更多的技艺，才能够征服它。这可不是一
件容易的事情！而且，蜂巢的地方太小，也盛不下这么大的一个东西。所以，黄蜂只得
放弃猎取大个儿的蜘蛛，不去干这种费时、费力、又不讨好的傻事。还是更实际些吧。
于是，它只得选择去猎取那些较小一些的蜘蛛为食。如果，它偶然会碰上一群可以猎食
的蜘蛛，那么它总是很聪明，从来也不贪多，只选择其中最小的那一个。但是，虽然个
儿头都是较小，但它的俘虏的身材还是差别比较大。因此，大小的不同，就会影响到数
目的不同。在这个巢穴里面，盛有一打蜘蛛，而在另外一个巢穴里面，只藏着五个或者
六个蜘蛛。
黄蜂专选那些个儿小的蜘蛛，还有一个理由，那就是，在它还没有把猎物装入它的
巢穴里之前，它先得把那个蜘蛛杀死。它所要采取的行动，有以下几步：它先是突然一
下子落到蜘蛛的身上，以快取胜，差不多连翅膀都还没来得及停下来，就要把这个小蜘
蛛带走。其它的昆虫所采用的什么麻醉的方法等等，这个小动物可是一点儿也不知道。
这个小小的食物，一旦被储藏起来，就很容易变坏。幸好这个蜘蛛的个子小，一顿就可
以把它全部吃掉。要是换了一只大一些的蜘蛛，一顿是不可能吃完的，只能分成几次吃。
这样的话，这个蜘蛛是一定要腐烂的。这样烂了的食品就会毒害窠巢里其它的幼虫，这
对整个家族是不利的。
我经常能够看到，黄蜂的卵并不是放在蜂案的上面。而是在蜂案里面储藏着的第一
个蜘蛛的身上。差不多都是这样的，完全没有什么例外。黄蜂都是把第一个被捉到的蜘
蛛放在最下层，然后把卵放到它的上面的，再把别的蜘蛛放在顶上。用了这种聪明的办
法以后，小幼虫就能先吃掉那些比较陈旧的死蜘蛛，然后再吃那些比较新鲜的。这样一
来，蜂案里面储藏的食物也就没有什么时间足以变坏了。这不失为一种很安全的办法。
蜂的卵总是放在蜘蛛的身上的某一部分的。蜂卵的包含头的一端，放在靠近蜘蛛最
肥的地方。这对于幼虫是很好的。因为，一经孵化以后，幼虫就可以直接吃到最柔软、
最可口和最有营养的食物了。因此，这是一个很聪明的主意。应该说，大自然赋予了黄
蜂一种相当巧妙的天性。这样的一个有经济头脑的动物，一口食物也不浪费掉。等到它
完全吃光这个蜘蛛的时候，一堆蜘蛛什么也剩不下来了。这种大嚼的生活要经过八天到
十天之久。
在一顿美餐之后，蛴螬就开始做它的茧了。那是一种纯洁的白丝袋，异常而又精致。
还有一些东西，能够使这个幼虫的丝袋更加坚实。这些东西，可以用作保护之用。于是，
蛴螬就又从它身体里生出一种像漆一样的流质。这种流质慢慢地浸入丝的网眼里，然后
会渐渐地变硬，成为一种很光亮的保护漆。此时，幼虫又会在它正在做的茧下面增加一
个硬的填充物，使得一切都十分妥当。
这一项工作完成以后，这个茧呈现出琥珀的黄颜色，很容易让人联想到那种洋葱头
的外皮。因为，它和洋葱头有着同样细致的组织，同样的颜色，同样的透明感，而且，
它和洋葱头一样，如果用指头摸一摸，便会立即发出沙沙的响声，完整的昆虫就从这个
黄茧里孵化出来。早一点或是迟一点，这要随气候的变化而变化，各有不同。
当黄蜂在蜂巢中把东西储藏好以后，如果我们打算和它开一次玩笑的话，就立刻会
显露出黄蜂的本能是如何的机械了。
在它辛辛苦苦地把它自己的巢穴做好以后，便带回了它的第一个蜘蛛。黄蜂会马上
把它拖进巢里，然后收藏起来，立刻，又在它的身体的最肥大的部位产下一个卵。做好
了这一切以后，它便又飞了出去，继续它的第二次野外旅行和捕猎。当它不在家里的时
候，我从它的巢穴里，把那只死蜘蛛连同那个卵一起都取走了。就算和这只黄蜂开个小
小的玩笑吧。不知道它会有什么样的反应。
我们很自然地就会想到，如果这个小动物稍微有一点儿头脑的话，那么，这个蜘蛛
和卵的失踪，它是一定能够发觉得到的，而且应该会感到奇怪的！蜂卵虽然是小的，但
是，它是被放在那个大的蜘蛛的身体上的。那么，当我们的这个小东西回来以后，发现
巢穴里面是空的，它会怎么做呢？将有什么举动呢？它将很有理智地行动，再产下一个
卵，以补偿它所失去的那一个吗？事实上这些都不是，它的举动是非常不合情理的。
现在，这个小东西所做的事情，只不过是又带回了一只蜘蛛，非常坦然地再次把它
放到那巢穴里边去。对于其它的事情一律不理睬，就好像并没有发生过什么意外一样。
似乎它根本就没有看到自己的孩子已经丢失了。那只刚刚捕获的蜘蛛也已经丢了。它没
有发现这一切的不幸，也并没有表现出什么吃惊、诧异、着急、不知所措之类的失意的
表情。这以后，居然若无其事地一只又一只地盲目地往窠里继续传带着蜘蛛。每当它把
巢里的猎物和卵都安排妥当了以后，便又飞了出去，继续盲目地执着地奋斗着。每次在
它飞出去的时候，我都会把这些蜘蛛和蜂卵悄悄地拿出来。因此，它每一次游猎回来，
储藏室里实际上总是空着的。就这样，它十分固执而徒劳地忙碌了整整两天时间。它一
心打算要使劲努力，无论如何也要争取装满这个不知为什么永远也装不满的食物瓶子。
我呢，也和它一样，不屈不挠地坚持了有两天的工夫。一次又一次耐心地把巢穴里的蜘
蛛和卵取出来。想要看看这个执着的小傻瓜究竟要等何时才能终结它这种看起来毫无意
义的工作。当这个傻乎乎的小动物完成了它的第二十次任务的时候，也就是到了第二十
次的收获物送来的时候，这位辛苦多时的猎人大概以为这罐子已经装够了——或许也是
因为这么多次的旅行，疲倦了——于是，它便自认为非常小心而且谨慎地把自己的巢穴
封锁了起来，然而，实际上，里面却完全是空的！什么东西都没有。它忙碌了这么久，
事实上它根本意识不到这一点，真是让人可怜啊！
在任何情况下，昆虫的智慧都是非常有限的。这一点是毫无疑问的，无论是哪一种
临时的困难，昆虫，这种小小的的动物，都是无力加以很好而且迅速的解决的。无论是
哪一个种类的昆虫，都同样的不能对抗。这一点，我可以列举出一大堆的例子来，证明
昆虫是一种完全没有理解能力的动物。当然，同时，它也是一种不具有意识的动物。虽
然它们的工作是那么异常的完备。经过长时间的经验和观察，使我不能不断定它们的劳
动，既不是自动的，也不是有意识的。它们的建筑、纺织、打猎、杀害，以及麻醉它们
的捕获物，都和消化食物或是分泌毒汁一样，其方法和目的完全都是不自知的。所以，
我相信这样一点，即这些动物对于它们所具有的特殊的才能，完全是莫名其妙的，既不
知也不觉。
动物的本能是不能改变的。经验不能指导它们；时间也不能使它们的无意识有一丝
一毫的觉醒。如果它们只有单纯的本能，那么，它们便没有能力去应付大千世界，应付
大自然环境的变化。环境是要经常有所变化的，意外的事情有很多，也时常会发生，正
因为如此，昆虫需要具备一种特殊的能力，来教导它，从而让它们自己能够清楚什么是
应该接受的，什么又是应该拒绝的。它需要某种指导。这种指导，它当然是具备的，不
过，智慧这样一个名词，似乎太精细了一点，在这里是不适用的。于是，我预备叫它为
辨别力。
那么，昆虫，能够意识到它自己的行动吗？能，但同时也不能。如果它的行动是由
于它所拥有的本能而引起的，那么它就不能知道自己的行动。如果它的行动是由于辨别
力而产生的结果，那么，它就能意识到。
比如，舍腰蜂利用软土来建造巢穴，这一点就是它的本能。它常常是如此建造巢穴
的，从一生下来就会。既不是时间，也不是生活的奋斗与激励，能够使得它模仿泥水匠
蜂，用那种细沙的水泥去建造它自己的巢，这并不是它的本能。
黄蜂的这个泥巢，一定要建在一种隐避之处，以便抵抗自然风雨的侵袭。在最初的
时候，大概那种石头下面可以隐匿的地方就能够被认为是相当合适了。但是，当它发现
还有更好的其它的地方可以选择时，它便会立刻去占据下来，然后搬到人家的屋子里边
去住。那么，这一种就属于辨别力了。
黄蜂利用蜘蛛作为它的子女的食物，这就是它本能的一种。表现没有其它的任何方
法，能够让这只黄蜂明白，小蟋蟀也是一样的好，和蜘蛛一样可以当作食物。不过假设
那种长有交叉白点的蜘蛛缺少了，那它也不会让它的宝宝挨饿的。它会选择其它类型的
蜘蛛，将其捕捉回来，给它的子女吃。那么，这种就是辨别力。
在这种辨别力的性质之下，隐伏了昆虫将来进步的可能性。

四、它的来源
舍腰蜂又给我们带来了另外一个问题。它寻找着我们房子里的火炉的热量。这是因
为它的巢穴是用软土建筑起来的，潮湿会把它给弄成泥浆而无法居住。所以，基于上述
原因，它必须要有一个干燥的隐避场所。因此，热量，也是黄蜂生活中所必要的。
那么，它是不是一个侨民呢？或许它是从海边被卷过来的？它是从有枣椰树的陆地
来到生长洋橄榄的陆地的？如果真的是这样的话，那么它们也就自然地会觉得我们这个
地方的太阳不够温暖，也就必须要寻找一些地方，比如说火炉，作为人工取暖的地方了。
这样就可以解释它的习性了，为什么它能和别的种类的黄蜂有如此大的差别。而且这种
蜂都是避人的。
在它还没有到我们这里来做客以前，它的生活是什么样子的呢？在没有房屋以前，
它住在什么样的地方呢？没有烟筒的时候，它把蛴螬隐藏在哪里呢？
也许，当古代山上的居民用燧石做武器，剥掉羊皮做衣服，用树枝和泥土造屋子的
时候，这些屋子也已经早就有舍腰蜂的足迹了。也许，它们的巢就建筑在一个破盆里面，
那是我们的祖先用手指取粘土制做成的。或者它就在狼皮及熊皮做的衣服的褶缝里边筑
巢。我感到非常奇怪，当它们在用树枝和粘土造成的粗糙的壁上做巢的时候，它们是否
选择那些靠近烟筒的地点呢？这些烟筒，虽然和我们现在所使用的烟筒不同，但是，在
不得已的时候，那些烟筒也是可以利用的。
如果说，舍腰蜂在古代的时候，的确和那些最古老的人们共同在这个地方居住过，
那么它所经历和见到的进步，就真的是不少了，而且，它所得到的文明的利益也真正是
不少了。它已经把人类不断增进的幸福转变成为自己的了。当人类社会，发明出在房屋
的屋顶上铺上天花板的法子，想出在烟筒上加上管子的主意以后，我们便可以想象得到，
这个怕冷的动物就会悄悄地对自己自言自语道：
“这是如何的舒适啊！让我们在这里撑开帐篷吧！”
但是，我们还应该追究得更远一些。在小屋没有出现以前，在壁龛也不常见之前，
甚至是在人类还没有出现之前，舍腰蜂又是在哪里造房子的呢？这个问题，当然不是一
个孤立的问题。我们还可以提出这样的问题：燕子和麻雀在没有窗子、烟筒等东西以前，
它们又是在哪里筑巢的呢？
燕子、麻雀、舍腰蜂是在人类出现以前就已经存在了的动物。显然，它们的工作是
不能依靠人类劳作的。当这里还没有人类的时候，它们各自必定已经具有了高超的建筑
技艺了。
三四十年来，我都常常问我自己，在那个时候舍腰蜂住在哪里的问题。
在我们的屋子外面，我们找不到它们的窠巢的痕迹。在房子外边，在空旷的广场、
在荒丘的草地里，我们都没有找到舍腰蜂的住处。
但是，最后，我长时间的研究结果表明，一个帮助我的机会出现了。
在西往南的采石场上，有许多碎石头子和很多的废弃物，堆积在这里有很长时间了。
据说已经有几百年的时间了。在这个乱石堆上，沉积了几个世纪的污泥，几百年的风风
雨雨，将这些乱石堆摆在人们面前。田鼠也在那里生活着。在我寻找这些宝藏的时候，
我有三次发现了在乱石堆中的舍腰蜂的巢。
这三个巢与在我们屋子里发现的完完全全一个样，材料当然也是泥土，而用以保护
的外壳，也是相同的泥土。
这个地方的危险性，并没有促使这位建筑家丝毫的一点儿进步。我们有时——不过
很少——看到舍腰蜂的巢筑在石堆里和不靠着地的平滑的石头下面。
在它们还没有侵入我们的屋子以前，它们的窠巢一定是建筑在这类地方的。
然而，这三个巢的形状，是很凄惨的，湿气已经把它们给侵蚀坏了，茧子也被弄得
粉碎了。周围也没有厚厚的土保护着它，它们的幼虫也已经牺牲了——已经被田鼠或别
的动物吃光了。
这个荒废的景象，使我惊疑的来到我邻居的屋外。是否能够真为舍腰蜂建巢的地点，
挑选一个适当的位置呢？事实很显然，母蜂不愿意这么做，并且也不至于被驱逐到这么
绝望的地步。同时，如果气候使它不能从事它祖先的生活方式，那么，我想，我可以断
言，它就是一个侨民。它很可能是从遥远的异国他乡，侨居到这个地方来的侨民。也可
能是另一种移民，那种背井离乡的移民。也可能是难民，是为了生计，不得不远走他乡，
被其他地方收养的难民。
事实的确如此，它是从炎热的、干燥的、缺水的、沙漠式的地方来的，在它们那里，
雨水不多，雪简直是没有的。
我相信，舍腰蜂是从非洲来的。
很久以前，它经过了西班牙，又经过了意大利，来到了我们这里，它可以说的上是
千里迢迢，也可以说它是不远万里、不辞辛苦地到我们这里来。
它不会越过长着洋橄榄树的地带，再往北去。它的祖籍是非洲，而现在它又归入了
我们布曼温司。
在非洲，据说它常把巢穴建筑在石头的下面，而在马来群岛，听说也有它们的同族、
同宗，它们是住在屋子里的。
从世界的这一边，来到世界的那一边，从世界的南边来到世界的北边，从地球的南
边——非洲，来到地球的北边——欧洲！最后又来到马来群岛。它的嗜好都是一个样的：
蜘蛛、泥巢，还有人类的屋顶。
假如我是在马来群岛，我一定要翻开乱石堆，翻找它居住的巢穴。这时，我会很高
兴地在一块平滑的石头下面，发现它的巢穴，发现它的住所——原来它的位置，就在这
些石头的下面。
------------------








螳螂



一、打猎
在南方有一种昆虫，与蝉一样，很能引起人的兴趣，但不怎么出名，因为它不能唱
歌。如果它也有一种钹，它的声誉，应比有名的音乐家要大得多，因为它在形状上与习
惯上都十分的不平常。它将是一名出色的乐手。
多年以前，在古希腊时期，这种昆虫叫做螳螂，或先知者。农夫们看见它半身直起，
立在太阳灼烧的青草上，态度很庄严，宽阔的、轻纱般的薄翼，如面膜似的拖曳着，前
腿形状如臂，伸向半空，好像是在祈祷，在无知识的农夫看来，它好像是一个女尼，所
以后来，就有人称呼它为祈祷的螳螂了。
这个错误再大没有了！那种貌似真诚的态度是骗人的，高举着的似乎是在祈祷的手
臂，其实是最可怕的利刃，无论什么东西经过它的身边，它便立刻原形毕露，用它的凶
器加以捕杀。它真是凶猛如饿虎，残忍如妖魔，它是专食活的动物的。看来，在它温柔
的面纱下，隐藏着十分吓人的杀气。
如果单从外表上看来，它并不令人生畏，相反，看上去它相当美丽，它有纤细而优
雅的姿态，淡绿的体色，轻薄如纱的长翼。颈部是柔软的，头可以朝任何方向自由转动。
只有这种昆虫能向各个方向凝视，真可谓是眼观六路。它甚至还有一个面孔。这一切都
构成了这样一个小动物的温柔。
螳螂天生就有着一副娴美而且优雅的身材。不仅如此，它还拥有另外一种独特的东
西，那便是生长在它的前足上的那对极具杀伤力，并且极富进攻性的冲杀、防御的武器。
而它的这种身材和它这对武器之间的差异，简直是太大了，太明显了，真让人难以相信，
它是一种温存与残忍并存的小动物。
见过螳螂的人，都会十分清楚地发现，它的纤细的腰部非常的长。不光是很长，还
特别的有力呢。与它的长腰相比，螳螂的大腿要更长一些。而且，它的大腿下面还生长
着两排十分锋利的像锯齿一样的东西。在这两排尖利的锯齿的后面，还生长着一些大齿，
一共有三个。总之，螳螂的大腿简直就是两排刀口的锯齿。当螳螂想要把腿折叠起来的
时候，它就可以把两条腿分别收放在这两排锯齿的中间，这样是很安全的，不至于自己
伤到自己。
如果说螳螂的大腿像是两排刀口的锯齿的话，那么它的小腿可以说是两排刀口的锯
子。生长在小腿上的锯齿要比长在大腿上的多很多。而且，小腿上的锯齿和大腿上的有
一些不太相同的地方。小腿锯齿的末端还生长着尖而锐的很硬的钩子，这些小钩子就像
金针一样。除此以外，锯齿上还长着一把有着双面刃的刀，就好像那种成弯曲状的修理
各种花枝用的剪刀一样。
对于这些小硬钩，我有着许多不堪回首的记忆。每次想到它们，都有一种难受的感
觉。记得从前曾经有过许多次这样的经历。在我到野外去捕捉螳螂的时候，经常遭到这
个小动物的强有力的自我保护与还击，总是捉它不成，反过来倒中了这个小东西的十分
厉害的“暗器”，被它抓住了手。而且，它总是抓得很牢，不轻易松开，让我自己无法
从中解脱出来，只有想其他的方法，请求别的人前来相助，帮我摆脱它的纠缠。所以，
在我们这种地方，或许再也没有什么其它的昆虫比这种小小的螳螂更难以对付，更难以
捕捉的了。螳螂身上的武器、暗器很多，因此，它在遇到危险的时候，可以选择多种方
法来自我保护。比如，它有如针的硬钩，可以用镰钩去钩你的手指；它长有锯齿般的尖
刺，可以用它来扎、刺你的手；它还有一对锋利无比、而且十分健壮的大钳子。这对大
钳子对你的手有相当的威力，当它挟住你的手时，那滋味儿可不太好受啊！综上所述，
这种种有杀伤力的方法，让你很难对付它。要想活捉这个小动物，还真得动一番脑筋，
费一番周折呢！否则，捉住它将是不可能的。这个小东西不知要比人类小多少倍，但却
能威胁住人类。
平时，在它休息、不活动的时候，这个异常勇猛的捕捉其它昆虫的机器，只是将身
体蜷缩在胸坎处，看上去，似乎特别的平和，不至于有那么大的攻击性，甚至会让你觉
得，这个小动物简直是一只热爱祈祷的温和的小昆虫。但是，它可不总是这样的，否则
的话，它身上具备的那些进攻、防卫的武器也就派不上什么用场了。只要是有其它的昆
虫从它们的身边经过，无论是什么样的昆虫，也无论它们是无意路过，还是有意地侵袭，
螳螂的那副祈祷和平的相貌便会一下子烟消云散了。这个刚才还是蜷缩着休息的小动物，
立刻便伸展开它身体的三节，于是，那个可怜的路过者，还没有完全反应过来，便己糊
里糊涂地成了螳螂利钩之下的俘虏了。它被重压在螳螂的两排锯齿之间，移动不得。然
后，螳螂很有力地把钳子夹紧，一切战斗就都结束了。无论是蝗虫，还是蚱蜢，或者甚
至是其它更加强壮的昆虫，都无法逃脱这四排锋利的锯齿的宰割。于是，一旦被捉，只
好束手就擒了。它可真是个了不得的杀虫机器。
假如你想到原野里面去详尽地研究，观察螳螂的习性，那几乎是不可能的。因此，
也就不得不把螳螂拿到室内来进行观察、分析和研究。如果把螳螂放在一个用铜丝盖住
的盆里面，再往盆里加上一些沙子，那么，这只螳螂将会生活得十分快乐和满意。我所
要做的，只是提供给它充足而又新鲜的食物就可以了。有了它必须的食品，它会生活得
更满意。因为我想要做一些试验，测量一下螳螂的筋力究竟能够有多大，所以，我不仅
仅是提供一些活的蝗虫或者是活的蚱蜢给螳螂吃，同时，还必须供给它一些最大个儿的
蜘蛛，以使它的身体更加强壮。至于我的观察、研究，以下便是在我做了上述工作以后，
所观察到的情形。
有这样一只不知危险、无所畏惧的灰颜色的蝗虫，朝着那只螳螂迎面跳了过去。后
者，也就是那只螳螂，立刻表现出异常愤怒的态度，接着，反应十分迅速地做出了一种
让人感到特别诧异的姿势，使得那只本来什么也不怕的小蝗虫，此时此刻也充满了恐惧
感。螳螂表现出来的这种奇怪的面像，我敢肯定，你从来也没有见到过。螳螂把它的翅
膀极度地张开，它的翅竖了起来，并且直立得就好像船帆一样。翅膀竖在它的后背上，
螳螂将身体的上端弯曲起来，样子很像一根弯曲着手柄的拐杖，并且不时地上下起落着。
不光是动作奇特，与此同时，它还会发出一种声音。那声音特别像毒蛇喷吐气息时发出
的声响。螳螂把自己的整个身体全都放置在后足的上面。显然，它已经摆出了一副时刻
迎接挑战的姿态。因为，螳螂已经把身体的前半部完全都竖起来了，那对随时准备东挡
西杀的前臂也早已张了开来，露出了那种黑白相间的斑点。这样一种姿势，谁能说不是
随时备战的姿势呢？
螳螂在做出这种令谁都惊奇的姿势之后，一动不动，眼睛瞄准它的敌人，死死盯住
它的俘虏，准备随时上阵，迎接激烈的战斗。哪怕那只蝗虫轻轻地、稍微移动一点位置，
螳螂都会马上转动一下它的头，目光始终不离开蝗虫。螳螂这种死死的盯人战术，其目
的是很明显的，主要就是利用对方的惧怕心理，再继续把更大的惊恐纳入这个不久以后
就将成为牺牲者的对手心灵深处，造成“火上加油”的效果，给对手施加更重的压力。
螳螂希望在战斗未打响之前，就能让面前的敌人因恐惧心理而陷于不利地位，达到使其
不战自败的目的。因此，螳螂现在需要虚张声势一番，假装什么凶猛的怪物的架势，利
用心理战术，和面前的敌人进行周旋。螳螂真是个心理专家啊！
看起来，螳螂的这个精心安排设计的作战计划是完全成功的。那个开始天不怕、地
不怕的小蝗虫果然中了螳螂的妙计，真的是把它当成什么凶猛的怪物了。当蝗虫看到螳
螂的这副奇怪的样子以后，当时就有些吓呆了，紧紧地注视着面前的这个怪里怪气的家
伙，一动也不动，在没有弄清来者是谁之前，它是不敢轻易地向对方发起什么攻势的。
这样一来，一向擅于蹦来跳去的蝗虫，现在，竟然一下子不知所措了，甚至连马上跳起
来逃跑也想不起来了。已经慌了神儿的蝗虫，完全把“三十六计，走为上策”这一招儿
忘到脑后去了。可怜的小蝗虫害怕极了，怯生生地伏在原地，不敢发出半点声响。生怕
稍不留神，便会命丧黄泉，在它最害怕的时候，它甚至莫明其妙地向前移动，靠近了螳
螂。它居然如此地恐慌，到了自己要去送死的地步。看来螳螂的心理战术是完全成功了。
当那个可怜的蝗虫移动到螳螂刚好可以碰到它的时候，螳螂就毫不客气，一点儿也
不留情地立刻动用它的武器，用它那有力的“掌”重重的地击打那个可怜虫，再用那两
条锯子用力地把它压紧。于是，那个小俘虏无论怎样顽强抵抗，也无济于事了。接下来，
这个残暴的魔鬼胜利者便开始咀嚼它的战利品了。它肯定是会感到十分得意的。就这样，
像秋风扫落叶一样地对待敌人，是螳螂永不改变的信条。
在蜘蛛捕捉食物、降服敌人的时候，它通常采取的办法是：首先，一上来便先发制
人，猛烈地刺击敌人的颈部，让它中毒。这样做的好处是对手中了毒，自然也就没有了
力气，也就不能继续抵抗防卫了。先下手为强嘛！与此相同的，螳螂在攻击蝗虫的时候，
也是首先重重地、不留情面地击打对方的颈部。受了一顿狂轰乱炸的痛捶之后，再加上
先前万分的恐惧，蝗虫的运转能力逐渐下降，动作慢慢地迟缓下来。也许是已经被打蒙
了的原因吧。这种办法既有效又非常的实用。螳螂就是利用这种办法，屡屡取得战斗的
胜利。无论是杀伤并食用和它一样大小的动物，还是对付比自己还要大一些的昆虫，这
种办法都是十分有效的。不过，最让人感到奇怪的，就是这么一只小个儿的昆虫，竟然
是一种十分贪吃的动物，能吃掉这么多的食物。
那些爱掘地的黄蜂们，算得上是螳螂的美餐之一了，因此常常受到螳螂的光顾。螳
螂经常出没于黄蜂的地穴附近。因此，在黄蜂的窠巢近区看到螳螂的身影屡屡出现，便
不足为奇了。螳螂总是埋伏在蜂窠的周围，等待时机，特别是那种能获得双重报酬的好
机会。为什么说是双重报酬呢？原来，有的时候，螳螂等待的不仅仅是黄蜂本身，因为
黄蜂自己的身上常常也会携带一些属于它自己的俘虏。这样一来，对于螳螂而言，不就
是双份的俘虏，双重报酬了吗？不过，螳螂并不总是这么走运的，也有不太幸运的时候。
有时，它也会常常什么都等不到，竟无功而返。主要原因是，黄蜂已经有所疑虑，从而
有所戒备了，方让螳螂失望而归。但是，也有个别掉以轻心者虽已发觉但仍不当心的，
被螳螂看准时机，一举将其抓获。这些命运悲惨的黄蜂为什么会遭到螳螂的毒手呢？因
为，有一些刚从外面回家的黄蜂，它们振翅飞来，有一些粗心大意，对早已埋伏起来的
敌人毫无戒备。当突然发觉大敌当前时，会被猛地吓了一跳，心里会稍稍迟疑一下，飞
行速度忽然减慢下来。但是，就在这千钧一发的关键时刻，螳螂的行动简直是迅雷不及
掩耳。于是，黄蜂一瞬间便坠入那个两排锯齿的捕捉器中——即螳螂的前臂和上臂的锯
齿之中了。螳螂就是这样出其不备，以快致胜的。接下来，那个不幸的牺牲者就会被胜
利者一口一口地蚕食掉。又成了螳螂的一顿美餐。
记得有一次，我曾看见过这样有趣的一幕。有一只黄蜂，刚刚俘获了一只蜜蜂，并
把它带回到自己的储藏室里，正在享用这只蜜蜂体内的蜜汁。不料，正在它吃得高兴的
时候，遭到了一只凶悍的螳螂的突然袭击。它无力还击，便束手就擒了。这只黄蜂正在
吃蜜蜂的嗉袋里储藏的蜜，但是螳螂的双锯，在不经意中，竟然有力地夹在了它的身上。
可是，就是在这种被俘虏的关键时刻，无论怎样的惊吓、恐怖和痛苦，竟然不能让这只
贪吃的小动物停止继续吸食蜜蜂体内的蜜汁。它依然在甜食着那芬香诱人的蜜汁。这真
是太奇异了，真是人为财死，鸟为食亡啊！
螳螂，这样一种凶狠恶毒、有如魔鬼一般的小动物，它的食物的范围并不仅仅局限
于其它种类的所有昆虫。螳螂的气概虽然特别神圣，但是，或许你想不到，因为这实在
是让人不可思议。事实上，螳螂还是一种自食其同类的动物呢。也就是说，螳螂是会吃
螳螂的，吃掉自己的兄弟姐妹。而且，在它吃的时候，面不改色，心不跳，十分泰然自
若，那副样子，简直和它吃蝗虫，吃蚱蜢的时候一模一样，仿佛这是天经地义的事情。
并且，与此同时，围绕在食同类的螳螂旁边围观的观众们，也没有任何反应，没有任何
抵抗的行动。不仅如此，这些观众还纷纷跃跃欲试，时刻准备着，一旦有了机会，它们
也会做同样的事情，也同样地毫不在乎，仿佛顺理成章似的。然而在事实上，螳螂甚至
还具有食用它丈夫的习性。这可真让人吃惊！在吃它的丈夫的时候，雌性的螳螂会咬住
它丈夫的头颈，然后一口一口地吃下去。最后，剩余下来的只是它丈夫的两片薄薄的翅
膀而已。这真人难以置信。
螳螂真的是比狼还要狠毒十倍啊！听说，即便是狼，也不吃它们的同类。那么，螳
螂真的是很可怕的动物了！

二、它的巢
虽然我们的螳螂是如此的凶猛而又可怕，它身上有那么多的杀伤性很强的武器，还
有那么凶恶的捕食方法，甚至它居然要以自己的同类为食。尽管如此，螳螂也和人类是
一样的，不光有缺点和不足之处，还拥有很多自己的优点。比如，螳螂能够建造十分精
美的巢穴，这便是螳螂众多优点中很突出的一个。
螳螂建造的窠巢，在有太阳光照耀的地方随处都可以找得到。比如，石头堆里，木
头块下，树枝上，枯草丛里，一块砖头底下，一条破布下，或者是旧皮鞋的破皮子上面
等等。总之一句话，在任何东西上，只要那个东西上有凸凹不平的表面，都可以作为非
常坚固的地基。螳螂就是利用这样的地基建巢的。
螳螂的巢，大小约有一两寸长，不足一寸宽。巢的颜色是金黄色的，样子很像一粒
麦子。这种巢是由一种泡沫很多的物质做成的。但是，不久以后，这种多沫的物质就逐
渐变成固体了，而且慢慢地变硬了。如果燃烧一下这种物质，便会产生出一种像燃烧丝
质品一样的气味儿。螳螂巢的形状各不相同。这主要是因为巢所附着的地点不同，因而
巢随着地形的变化而变化，会有不同形状的巢存在。但是，不管巢的形状多么干变万化，
它的表面总是凸起的。这一点是不变的。
整个的螳螂巢，大概可以分成三个部分。其中的一部分是由一种小片做成的，并且
排列成双行，前后相互覆盖着，就好像屋顶上的瓦片一样。这种小片的边沿，有两行缺
口，是用来做门路的。在小螳螂孵化的时候，就是从这个地方跑出来的。至于其它部分
的墙壁，全都是不能穿过的。
螳螂的卵在巢穴里面堆积成好几层。其中每一层，卵的头都是向着门口的。前面我
已经提到过了，那道门有两行，分成左、右两边。所以，在这些幼虫中，有一半是从左
边的门出来的，其余的则从右边的门出来。
有这样一个事实是值得注意的，那就是母螳螂在建造这个十分精致的巢穴的时候，
也正是它产卵的时候。在这个时候，从母螳螂的身体里，会排泻出一种非常有粘性的物
质。这种物质和毛虫排泻出来的丝液很相象。这种物质在排泻出来以后，将与空气互相
混合在一起，于是就会变成泡沫。然后，母螳螂会用身体末端的小杓，把它打起泡沫来。
这种动作，特别像我们用叉子搅打鸡蛋蛋白一样。打起来的泡沫是灰白色的，与肥皂沫
十分相似。开始的时候，泡沫是有粘性的。但是过了几分钟以后，粘性的泡沫就变成了
固体。
母螳螂就是在这种泡沫的海洋中产卵，繁衍后代的，每当它产下一层卵以后，它就
会往卵上覆盖上一层这样的泡沫。于是，很快地，这层泡沫就将变成固体了。
在新建的巢穴的门外面，有一层材料，把这个巢穴封了起来。看上去，这层材料和
其它的材料并不一样——那是一层多孔、纯洁无光的粉白状的材料。这与螳螂巢内部，
其它部分的灰白颜色是完全不一样的。这就好像面包师们把蛋白、糖和面粉搅合在一起，
用来作饼干外衣的混合物一样。这样一种雪白色的外壳，是很容易破碎的，也很容易脱
落下来。当这层外壳脱落下来的时候，螳螂巢的门口，完全裸露在外，可以看出来，门
的中间装着两行板片。不久以后，风吹雨打会把它侵蚀，将它剥成小片。于是这些小片
会逐渐地脱落下去。所以，在旧巢上，就看不见它的痕迹了。
至于这两种材料，虽然它们从外表上看来，一点儿也不一样，但是实际上，它们的
质地是完全一样的。它们只不过是同样原质的东西的两种不同的表现形式罢了。螳螂用
它身上的杓打扫着泡沫的表面，然后，撇掉表面上的浮皮，使其形成一条带子，覆盖在
巢穴的背面。这看起来，就像那种冰霜的带一样。因此，这种物质实际上仅仅是粘性物
质的最薄、最轻的那一部分。它看上去之所以会比较白一些，主要是因为它的泡沫比较
细巧，光的反射力比较强罢了。
这可真是一个非常奇异的操作方式。它相当有自己的一套方法，可以很迅速、很自
然地做成一种角质的物质。于是，螳螂的第一批卵就生产在这种物质上面了。
螳螂真是一种很能干的动物，也是一种很有建筑才能的动物。产卵时，它排泻出用
于保护的泡沫，制造出柔软得像糖一样的包被物，同时，它还能制做出一种遮盖用的薄
片，以及通行用的小道。而在进行这一切工作的时候，螳螂都只是在巢的根脚处站立着，
一动也不动，用不着移动身体，就在它背后建筑起一座了不起的建筑物，而它自己对这
个建筑物连看都不看一眼。它那粗壮而有力的大腿，在这件事的整个过程中，竟然没有
用武之地，发挥不了什么作用，一点儿都没有做什么。这所有一切的繁杂工作，完全都
是这部小机器自己完成的。
母亲的工作成功完成以后，就洗手不干了，放开一切，跑走了。我总是对它抱着一
线希望，盼望着它有朝一日能够回来看一下，以便表示一些它对整个家族的生产地的爱
护和关切之情。但是，我的这个希望总也得不到实现。看起来很显然，这一点对于它竟
然没有一点儿兴趣了。它真的是一去不回头了。
所以，根据这一事实，我便得出了这个结论：螳螂都是些没有心肝的东西，尽干一
些残忍、恶毒达到极点的事情。比如，它以自己的丈夫作为美餐，而且，它居然还会抛
弃它自己的子女，弃家出走且永不归还。
螳螂卵的孵化，通常都是在有太阳光的地方进行的，而且，大约是在六月中旬，上
午十点钟的时候。
在前面的文章里，我已经告诉过大家了，在这个螳螂巢里，只有一小部分可以为螳
螂幼虫当作出路。这一部分指的就是窠巢里面那一带鳞片的地方。再仔细地观察一下，
你就会发现在每一个鳞片的下面，都可以看见一个物体，而且稍微有一点儿透明的小块
儿。在这个小块儿的后面，紧接着的就是两个大大的黑点。那不是别的什么东西，就是
那个可爱的小动物的一对小眼睛了。幼小的螳螂蛴螬，静静地伏卧在那个薄薄的片下面。
如果仔细地看一下，就会发现它现在差不多已经有将近一半的身体解放了出来。下面，
再看看这个小东西的身体是什么样的吧。它身体的颜色主要是黄色，又带有一些红的颜
色。它长了一个十分肥胖而且很大的脑袋。从这个幼虫的外面的皮肤来看，非常容易地
就能分辨出它的那对特别大的眼睛来。幼虫的小嘴是贴在它的胸部的，腿又是和它的腹
部紧紧相贴的。这只小幼虫，从它的外形上看，除了它那些和腹部紧贴着的腿以外，其
它部分都能够让人联想到另一种动物的状态，那就是那种刚刚才离开巢穴的蝉的最初期
的状态。这种状态和此时的螳螂幼虫像极了。
和蝉一样，为了图方便，但是更重要的是为了安全起见，幼小的螳螂刚一降临到这
个世界上来，的确有穿上一层结实的外套的必要。要是幼虫打算从巢穴中非常狭小而又
弯曲的那条小道里爬出来，如果它想要完全地把自己的小腿伸展开来，那将是不太可能
的一件事了。这主要是因为，如果它完全伸展开身体，高高翘起那尚还缺乏力量的用来
杀戮敌人的长矛，然后再竖立起它那十分灵敏的触须，那么这样一来，它自己就完全把
自己的道路给阻挡住了，以至它无法前行，而根本不可能从通道中爬行出来。正因为如
此，这个小动物，在它刚刚降临到这个世界上的时候，它是被团团包裹在一个襁褓之中
的，那种形状就好像一只小船一般。
在小幼虫刚刚降生，出现在巢中的薄片下面不久以后，它的头便逐渐地变大，一直
膨胀到形状像一粒水泡一样为止。这个有力气的小生命，在出生后不久，就开始靠自己
的力量努力生存。它一刻也不停地一推一缩地努力地解放着自己的躯体。就这样，每每
做一次动作的时候，它的脑袋就要稍稍变大一些。最终的结果是，它的胸部的外皮终于
破裂了。于是，它便更加努力，“乘胜追击”。它摆动得更加剧烈了，也更加快了。它
挣扎着，用尽浑身解数，不停歇地弯曲扭动着它那副小小的躯干。看来，它是义无反顾
地下定决心要挣脱掉这一件外衣的束缚，想马上看到外面的大千世界究竟是个什么样子。
渐渐地，首先得到解放的是它的腿和触须。然后，它继续不懈地努力。又进行了几次摆
动与挣扎以后，终于，它的目的和企图就完全实现了。
有几百只小螳螂，它们同时团团地拥挤在不太宽敞的巢穴之中，这场景，倒真的算
得上是一种不可多得的奇观呢！当巢中的螳螂幼虫还没有集体打破外衣，还没有集体冲
出襁褓，变成螳螂的形态之前，首先暴露出它的那双小眼睛。我们很少会见到哪一个单
独的小动物独自行动。情况恰恰相反，就好像存在什么统一行动的信号一样，每当这信
号传达出来的时候，速度非常之快，几乎所有的卵差不多在同一时刻孵化出来，一起打
破它们的外衣，从硬壳中抽出身体来。因此，也就是在一刹那之间，螳螂巢穴的中部顿
时如同召开大会一样，无数个幼虫一下子集合起来，挤满了这个不太大的地方。它们近
乎狂热一般地爬动着，似乎很兴奋、很急切地要马上脱掉这件困扰它们生活的讨厌的外
衣。在这之后，它们或者是不小心跌落，或者是使劲地爬行到巢穴附近的其它的枝叶上
面去。再过几天以后，就会在巢穴中又发现一群幼虫，它们同样要进行与前辈们相同的
工作，直到它们全都孵化出来。于是，繁衍就这样不停地继续下去。
然而，有一点非常不幸！这些可怜的小幼虫竟然孵化到了一个布满了危险与恐怖的
世界上来，也许它们自己还并不清楚明白这一点。曾经有过好多次，我在门外边的围墙
内，或者是在树林中的那些幽静的地方，看到螳螂的卵在孵化，一个个小幼虫破壳而出。
我总有一种美好的愿望，希望能够尽自己微薄的力量，好好地保护这些可爱的小生命，
让它们能够平平安安而且快快乐乐地生活在这个世界上。但是，很不幸，这种愿望总是
会成为泡影。已经至少有二十次了，(实际上比这要多得多)我总是看到那种非常残暴的
景象，总是亲眼目睹那令人恐惧的一幕。这些还不知道什么叫危险的小幼虫，在它们乳
臭未干的时候，便惨遭杀戮，还没来得及体验一下生活，体会一下生命的宝贵，就已经
结束了年幼的生命，真是可怜啊！螳螂虽然产下了许多卵，但是事实上，它并没有生产
极大数目的卵。至少，它所产下的那些卵，还不足以抵御那些早已在巢穴门口埋伏多时，
一旦幼虫出现，便会不失时机地加以杀戮的强大的敌人。
对于螳螂幼虫而言，它们的最具杀伤力的天敌，要算是蚂蚁了。几乎每一天，我都
会有意无意地看到，一只只蚂蚁不厌其烦地光临到螳螂巢穴的旁边，非常耐心，而且信
心十足地等待时机的成熟，以便立即采取先下手为强的行动。我一看到它们，就千方百
计地帮着螳螂驱赶它们。可是，无济于事。我的能力经常都驱逐不了它们。因为，它们
常常是先人一步，率先占据有利的位置。看来，蚂蚁的时间观念还是很强的，但是，虽
然它们早早就静候在大门之外，可它们却很难深入到巢穴的内部去。这主要是因为，螳
螂巢穴的四围有一层硬硬的厚壁，这便形成了十分坚固的壁垒，因而，蚂蚁对此束手无
策。它们的智慧还不足以使它们想出冲破这一层屏障的办法。不过，它们总是埋伏在巢
穴的门口，静候着它们的俘虏。
因此，螳螂幼虫的处境实在是非常危险的。只要它一不小心跨出自家大门一步，那
么，马上就会坠入深渊，葬送了自己的生命。因为守候在巢边的蚂蚁是不会轻易放过任
何一顿美餐的。一旦有猎物探出头来，便立刻将其擒住，然后再扯掉幼虫身上的外衣，
将其毫不客气地切成碎片。在这场战斗中，你可以看到，那些只能利用随意的乱摆来进
行自我保护的小动物，和那些前来俘虏食品的非常凶猛、残忍的大队的强盗们展开激烈
的拼杀。小动物们尽管非常弱小，但是仍然坚持着、挣扎着，不放弃对求生的渴望。但
是，这种挣扎与那些凶恶之众相比，显得多么可怜呢！用不了多长时间，也就是一小会
儿的工夫，这场充满血腥的大屠杀便宣告终结了。残杀过后，剩余下来的，只不过是碰
巧有幸能够逃脱敌人的恶爪的少数几个幸存者而已。其它的小生命，都已经变成了蚂蚁
口中之食了。就这样，一个原本的人丁兴旺的家族就衰败了。
这是非常奇怪的事情。前面我们曾提到过，螳螂是一种十分凶残的动物。它不仅以
锋利的杀伤性武器去攻击其它的动物，以猎取食物，而且还居然以自己的同类为食，并
且在食用自己的同胞骨肉时，竞然还那样心安理得坦坦荡荡。然而，就是这种可以被视
为昆虫中的灾害的螳螂，如今，在它刚刚拥有生命的初期，本身也要牺牲在昆虫中个儿
头最小的蚂蚁的魔爪下，这难道不奇妙吗！大自然造物真是让人不可思议啊！这种小小
的恶魔，眼睁睁地看着它自己的家族被这样毁掉，它自己的兄弟姐妹被这么一群小小的
侏儒所欺凌，所吞食，却对此束手无策。只能傻傻地目送亲人们远离这个充满危险的世
界。
不过，这样的情形并不会持续相当长的时期。因为，遭到不测的只是那些刚刚问世，
刚刚从卵中孵化出来的幼虫而己。但是，当这些幼虫开始和空气相接触以后，用不了多
长时间，便会马上变得非常的强壮。这样一来，渐渐地，它自己就具备了能够自我保护
的能力了，再也不是那些任人宰割的可怜虫了！
在它长大一些以后，情况就大不相同了。它从蚂蚁的群里，快速地走过去，它所经
过的地方，原来任意行凶的敌人们都纷纷跌倒下来，再也不敢去攻击和欺负这个已经长
大了的“弱者”了。螳螂在行进的时候，把它的前臂放置在胸前，作出一副自卫的警戒
状态。它那种骄傲的态度和不可小视的神气，早已经把这群小小的蚂蚁吓倒了。它们再
也不敢轻举妄动了，有些甚至已经望风而逃了。
但是，事实上，螳螂的敌人，不只是这些小个子的蚂蚁，还有许多其它的敌人。这
些天敌可不是那么容易就能吓倒的。比如说，那种居住在墙壁上面的小形的、灰色的蜥
蜴，就很难对付。对于小小螳螂的自卫和恐吓的姿势，它是全然不在意的。小蜥蜴进攻
螳螂的方法主要是用它的舌尖，一个一个地舐起那些刚刚幸运地逃出蚂蚁虎口的小昆虫。
虽然一个小螳螂还不能填满蜥蜴的嘴，但是，从它的面目表情便可以很清楚地看出来，
那味道却是非常之好。看来，它相当满意。每吃掉一个，蜥蜴的眼皮总是要微微一闭，
这的确是一种极端满足的表现。然而，对于那些年轻的、仍不走运的少年螳螂而言，它
们真可谓“才出龙潭，又入虎穴”啊！
不仅仅是在卵孵化出来以后是如此的危险，甚至就是在卵还没有发育出来以前，它
们就已经处于万分危险之中了。有这样一只小个儿的野蜂(challis)，它随身携带着一种
刺针，其尖利的程度，足可以刺透螳螂的由泡沫硬化以后而形成的巢穴，这样一来，螳
螂的血统，就如同蝉的子孙后代一样，遭受到相同的命运。这样一位外来的客人，并没
有受到谁的邀请，就在螳螂的巢穴中擅自决定产下自己的卵。它的卵的孵化也要比这巢
穴的主人的卵提前一步。于是，螳螂的卵就会顺其自然地受到侵略者的骚扰，被侵略者
吞食掉。比如说螳螂产下一千枚卵，那么，最后剩留下来的，没有遭受噩运而被残酷地
毁灭了的，大概也就只有一对而已了。
这样一来，便形成了下面这条生物链。螳螂以蝗虫为食，蚂蚁又会吃掉螳螂，而蚂
蚁又是鸡的食品。但是，等到了秋天的时候，鸡长大了，长肥了，我又会把鸡做成佳看
吃掉，这可真有趣！
或许螳螂、昨晚、蚂蚁，甚至是其它个儿头更小一些的动物，食用之后都可以增加
人类的脑力。它们采用一种非常奇妙但又见不到的方法，提供给我们的大脑某种有益的
物质。然后，作为我们人类思想之灯的油料。它们的精力慢慢地发达起来，然后贮蓄起
来，并且一点一点地传送到我们身体的各个部位，流进我们的血脉里。它们滋养着我们
身上的不足之处。我们就是生存在它们的死亡之上的。世界本来就是一个永无穷尽的循
环着的圆环。各种物质完结以后，在此基础上，各种物质又纷纷重新开始一切；从某种
意义上讲，各种物质的死，就是各种物质的生。这是一个十分深刻的哲学道理。
很多年前，人们总是习惯性地把螳螂的巢看做是一种充满迷信的东西。在布罗温司
这个地方，螳螂的巢，被人们视为医治冻疮的一种丹灵妙药。大多数人拿一个螳螂的巢，
然后把它劈开成两半，挤出里面的浆汁来，涂抹在痛楚的部位。农村里的人常说，螳螂
巢的功效，就仿佛有什么神奇的魔力一样。然而，我自己从没有感到它有什么功效。
与此同时，也有一些人盛传，认为螳螂巢医治牙痛非常有效。假如你有了它，也就
用不着再怕什么牙痛了。一般情况下，妇女们常常在月夜里到野外去收集它，然后，很
小心地收藏在杯碗橱子的角落里，或者是把它们缝在一个袋子里面，仔细珍藏起来。如
果附近的邻居们，谁要是患了牙痛病的话，就会跑过来，借用它。妇女们管它叫做“铁
格奴(Tigno)”。
如果是脸肿了的病人，他们会说：“请你借给我一些铁格奴，好吗？我现在痛得厉
害呢！”另外的一个就会赶快放下手里的针线活儿，拿出这个宝贝东西来。
而她会很慎重地对朋友说：“你随便做什么都可以，但是不要摘掉它。我只有这么
一个了，而且，现在又是没有月亮的时候！”
没有想到，农民们的这种心理上的简单而幼稚的反应，竟然被十九世纪的一位英国
医生兼科学家所超越。他曾经告诉过我们如此荒唐可笑的事情。他说在那个时候，如果
一个小孩子在树林里迷了路，他可以询问螳螂，让它指点道路。并且，他还说道：“螳
螂会伸出它的一足，指引给他正确的道路，而且很少或是从不会出错的。”
------------------








蜜蜂、猫和红蚂蚁



我希望能够了解更多的关于蜜蜂的故事。我曾听人说起过蜜蜂有辨认方向的能力，
无论它被抛弃到哪里，它总是可以自己回到原处。于是我想亲自试一试。
有一天，我在屋檐下的蜂窝里捉了四十只蜜蜂，叫我的小女儿爱格兰等在屋檐下，
然后我把蜜蜂放在纸袋里，带着它们走了二里半路，接着打开纸袋，把它们抛弃在那里，
看有没有蜜蜂飞回来。
为了区分飞到我家屋檐下的蜜蜂是否是被我扔到远处的那群，我在那群被抛弃的蜜
蜂的背上做了白色的记号。在这过程中，我的手不可避免地被刺了好几口，但我一直坚
持着，有时候竟然忘记了自己的痛，只是紧紧地按住那蜜蜂，把工作做完，结果有二十
多只损伤了，当我打开纸袋时，那些被闷了好久的蜜蜂一拥而出地向四面飞散，好像在
区分该从哪个方向回家一样。
放走蜜蜂的时候，空中吹起了微风。蜜蜂们飞得很低，几乎要触到地面，大概这样
可以减少风的阻力，可是我想，它们飞得这样低，怎么可以眺望到它们遥远的家园呢？
在回家的路上，我想到它们面临的恶劣环境，心里推测它们一定都我不到回家的方
向了。可是没等我跨进家门，爱格兰就冲过来，她的脸红红的，看上去很激动。她冲着
我喊道：
“有两只蜜蜂回来了！在两点四十分的时候到达巢里，还带来了满身的花粉。”
我放蜜蜂的时间是两点整。也就是说，在三刻钟左右的时间里，那两只小蜜蜂飞了
二里半路，这还不包括采花粉的时间。
那天天快黑的时候，我们还没见到其它蜜蜂回来。可是第二天当我检查蜂巢时，又
看见了十五只背上有白色记号的蜜蜂回到巢里了。这样，二十只中有十七只蜜蜂没有迷
失方向，它们准确无误地回到了家，尽管空中吹着逆向的风，尽管沿途尽是一些陌生的
景物。但它们确确实实地回来了。也许是因为它们怀念着巢中的小宝贝和丰富的蜂蜜。
凭借这种强烈的本能，它们回来了。是的，这不是一种超常的记忆力，而是一种不可解
释的本能，而这种本能正是我们人类所缺少的。

猫
我一直没有相信过这样一种说法，即猫也和蜜蜂一样，能够认识自己的归途。直到
有一天我家的猫的确这样做了，我才不得不相信这一事实。
有一天，我在花园里看见一只并不漂亮的小猫，薄薄的毛皮下显露着一节一节的脊
背，瘦骨嶙峋的。那时我的孩子们还都很小，他们很怜惜这只小猫，常塞给它一些面包，
一片一片还都涂上了牛乳。小猫很高兴地吃了好几片，然后就走了。尽管我们一直在它
后面温和地叫着它，“咪咪，咪咪——”，它还是无怨无悔地走了。可是隔了一会儿，
小猫又饿了。它从墙头上爬下来，又美美地吃了几片。孩子们怜惜地爱抚着它瘦弱的身
躯，眼里充满了同情。
我和孩子们作了一次谈话，我们达成一致，决定驯养它。后来，它果然不负众望，
长成一只小小的“美洲虎”——红红的毛，黑色的斑纹，虎头虎脑的，还有锋利的爪子。
它的小名叫做“阿虎”。后来阿虎有了伴侣，她也是从别处流浪而来的。他们俩后来生
了一大堆小阿虎。不管我家有什么变迁，我一直收养着它们，大约有二十多年了。
第一次搬家时，我们很为它们担忧，假如遗弃这些我们所宠爱的猫，它们将再度遭
受流浪的生活。可是如果把它们带上的话，雌猫和小猫们还能稳住气，保持安静，可两
只大雄猫——一只老阿虎，一只小阿虎在旅途上是一定不会安静的。最后我们决定这样：
把老阿虎带走，把小阿虎留在此地，替它另外找一个家。
我的朋友劳乐博士愿意收留小阿虎。于是某天晚上，我们把这只猫装在篮子里，送
到他家去。我们回来后在晚餐席上谈论起这只猫，说它运气真不坏，找到了一户人家。
正说着，突然一个东西从窗口跳进来。我们都吓了一跳，仔细一看，这团狼狈不堪的东
西快活而亲切地用身体在我们的腿上蹭着，这正是那只被送掉的小阿虎。
第二天，我们听到了关于它的故事：它刚到劳乐博士家里，就被锁在一间卧室里。
当它发现自己已在一个陌生的地方做了囚犯时，它就发狂一般地乱跳。一会儿跳到家俱
上，一会儿跳到壁炉架上，撞着玻璃窗，似乎要把每一样东西都撞坏。劳乐夫人被这个
小疯子吓坏了，赶紧打开窗子，于是它就从窗口里跳了出来。几分钟之后，它就回到了
原来的家。这可不是容易的事啊，它几乎是从村庄的一端奔到另一端，它必须经过许多
错综复杂的街道，其间可能遭遇到几千次的危险，或是碰到顽皮的孩子，或是碰到凶恶
的狗，还有好几座桥，我们的猫不愿意绕着圈子去过桥，它决定拣取一条最短的路径，
于是它就勇敢地跳入水中——他那湿透了的毛告诉了我们一切。
我很可怜这只小猫，它对它的家是如此的忠心。我们都同意带它一起走，正当我们
担心它在路上会不安分的时候，这个难题竟自动解决了。几天之后，我们发现它已经僵
硬地躺在花园里的矮树下。有人已经替我把它毒死了。是谁干的呢？这种举动可不会是
出自好意！
还有那只老阿虎。当我们离开老屋的时候，却怎么也找不到它了。于是我们另外给
车夫两块钱，请他负责找那老阿虎，无论什么时候找到它，都要把它带到新家这边来。
当车夫带着最后一车家俱来的时候，他把老阿虎带来了。他把它藏在自己的座位底下。
当我打开这活动囚箱，看到这前两天就被关进去的囚徒的时候，我真不能相信它就是我
的老阿虎了。
它跑出来的时候，活像一只可怕的野兽，它的脚爪不停在张舞着，口里挂着口水，
嘴唇上沾满了白沫，眼睛充满了血，毛已经倒竖起来，完全已经没有了原来的阿虎的神
态和风采。难道它发疯了吗？我仔细把它察看了一番。我终于明白了，它没有疯，只是
被吓着了。可能是车夫捉他的时候把它吓坏了，也可能是长途的旅行把它折磨得精疲力
尽。我不能确定到底是什么原因，但显而易见的是，它的性格大变，它不再口中常常念
念有词，不再用身体擦我们的腿了，只有一副粗暴的表情和深沉的忧郁。慈爱的抚慰也
不能消除它的苦痛了。终于有一天，我们发现它死了，躺在火炉前的一堆灰上，忧郁和
衰老结束了它的生命。如果它精力还足够的话，它会不会回到我们的老房子去呢？我不
敢断定。但是，这样一个小生灵，因为衰老的体力不允许它回到老家，终于得了思乡病，
忧郁而死，这总是一件令人感慨的事吧！
当我们第二次搬家的时候，阿虎的家族己完全换了一批了：老的死了，新的生出来
了。其中有一只成年的小阿虎，长得酷像它的先辈。也只有它会在搬家的时候增加我们
的麻烦。至于那些小猫咪和母亲们，是很容易制服的。只要把它们放在一只篮子里就行
了，小阿虎却得被单独放在另一只篮子里，以免它把大家都闹得不太平。这样一路上总
算相安无事。到了新居后，我们先把母猫们抱出篮子。它们一出篮子，就开始审视和检
阅新屋，一间一间地看过去，靠着它们粉红色的鼻子，它们嗅出了那些熟悉的家俱的气
味。它们找到了自己的桌子、椅子和铺位，可是周围的环境确实变了。它们惊奇地发出
微微的“喵喵”声，眼睛里时时闪着怀疑的目光。我们疼爱地抚摸着它们，给它们一盆
盆牛奶，让它们尽情享用。第二天它们就像在自己家里一样习惯了。
可是轮到我们的小阿虎，情形却完全不同了。我们把它放到阁楼上，让它渐渐习惯
新环境，那儿有好多空屋可以让它自由地游玩。我们轮流陪着它，给它加倍的食物，并
时时刻刻把其余的猫也捉上去和它作伴。我们想让它知道，它并不是独自一个在这新屋
里。我们想尽了一切办法，让它忘掉原来的家。果然，它似乎真的忘记了。每当我们抚
摸它的时候，它显得非常温和驯良，一叫它，它就会“咪咪”地叫着过来，还把背弓起
来。这样关了一个星期，我们觉得应该恢复它的自由了，于是把它从阁楼上放了出来。
它走进了厨房，和别的猫一同站在桌子边。后来它又走进了花园。我的女儿爱格兰紧紧
地盯着它看它有什么异样的举动，只见它做出一副非常天真的样子，东张张，西望望，
最后仍回到屋里。太好了，小阿虎再也不会出逃了。
第二天，当我们唤它的时候，任凭我们叫了多少声“咪咪咪咪——”，就是没有它
的影子！我们到处找，呼唤它，丝毫没有结果。骗子！骗子！我们上了它的当！它还是
走了，我说它是回到老家去了。可是家里其他人都不相信。
我的两个女儿为此特意回了一次老家。正如我说的那样，她们在那里找到了小阿虎。
她们把它装在篮子里又带了回来，虽然天气很干燥，也没有泥浆，可它的爪子上和腹部
都沾满了沙泥，无疑它一定是渡过河回老家去的，当它穿过田野的时候，泥土就粘在了
它湿漉漉的毛上，而我们的新屋，距离原来的老家，足足有四里半呢！
我们把这个逃犯关在阁楼上，整整两个星期之后，再放它出来。可是还不到一天工
夫，它又跑回去了，对于它的前途，我们只能听天由命了。后来有一位老屋的邻居来看
我们，说起小阿虎，说他有一次看到我们的猫口里叼着一只野兔，躲在篱笆下。是啊，
再也没有人喂给它食物了，它得用自己的力量去寻找食物。后来我就再也没有听到过它
的消息了。它的结局一定是挺悲惨的。它变成了强盗，当然要遭到强盗的命运。
这些真实的故事证明了猫和泥匠蜂一样，有着辨别方向的本领。鸽子也是这样，当
它们被送到几百里以外的时候；它们还能回来找到自己的老巢。还有燕子，还有许多别
的鸟也是这样。让我们再回到昆虫的问题上吧。蚂蚁和蜜蜂是最相似的一对昆虫，我很
想知道它们是不是像蜜蜂一样有着辨别方向的本领。

红蚂蚁
在一块废墟上，有一处地方是红蚂蚁的山寨。红蚂蚁是一种既不会抚育儿女也不会
出去寻找食物的蚂蚁，它们为了生存，只好用不道德的办法去掠夺黑蚂蚁的儿女，把它
们养在自己家里，将来这些被它们占为已有的蚂蚁就永远沦为了奴隶。
夏天的下午，我时常看见红蚂蚁出征的队伍，这队伍大约有五、六码长。当它们看
见有黑蚂蚁的巢穴时，前面的队伍出现一阵忙乱。几只间谍似的蚂蚁先离开了队伍往前
走。一队的蚂蚁仍旧列着队伍蜿蜒不停地前进，有时候有条不紊地穿过小径，有时在荒
草的枯叶中若隐若现。
最后，它们终于找到了黑蚂蚁的巢穴，就长驱直入地进入到小蚂蚁的卧室里，把它
们抱出了巢。在巢内，红蚂蚁和黑蚂蚁有过一番激烈的厮杀，最终黑蚂蚁败下阵来，无
可奈何地让强盗们把自己的孩子抢走。
我再讲一下它们一路上怎样回去的情形吧。
有一天我看见一队出征的蚂蚁沿着池边前进，那时天刮着大风，许多蚂蚁被吹落了，
白白地做了鱼的美餐。这一次鱼又多吃了一批意外的食物——黑蚂蚁的婴儿。显然蚂蚁
不会像蜜蜂那样，会选择另一条路回家，它们只会沿着原路回家。
我不能把整个下午都消耗在蚂蚁身上，所以我叫小孙女拉茜帮我监视它们。她喜欢
听蚂蚁的故事，也曾亲眼看到红蚂蚁的战争，她很高兴接受我的嘱托。凡是天气不错的
日子里，小拉茜总是蹲在园子里，瞪着小眼睛往地上张望。
有一天，我在书房里听到拉茜的声音，“快来快来！红蚂蚁已经走到黑蚂蚁的家里
去了！”
“你知道它们走的是哪条路吗？”
“是的，我已经做了记号。”
“什么记号？你怎么做的？”
“我沿路撒了小石子。”
我急忙跑到园子里，拉茜说得没错。红蚂蚁们正沿着那一条白色的石子路凯旋呢！
我取了一张叶子，截走几只蚂蚁，放到别处。这几只就这样迷了路，其它的，凭着它们
的记忆力顺着原路回去了。这证明它们并不是像蜂那样，直接辨认回家的方向，而是凭
着对沿途景物的记忆找到回家的路的。所以即使它们出征的路程很长，需要几天几夜，
但只要沿途不发生变化，它们也照旧回得来。
------------------








开隧道的矿蜂



矿蜂是细长形的蜜蜂，它们的身材大小不同，大的比黄蜂还大，小的比苍蝇还小。
但是它们有一个共同的特征，那就是它的腹部的底端有一条明显的沟，沟里藏着一根刺，
遇到敌人来侵犯时，这根刺可以沿着沟来回地移动，以保护自己。我这里要讲的是关于
矿蜂中的一种有红色斑纹的蜂。雌蜂的斑纹是很美丽夺目的，细长的腹部被黑色和褐色
的条纹环绕着。至于它的身材，大约和黄蜂差不多。
它的巢往往建在结实的泥土里面，因为那里没有崩溃的危险。比如，我们家院子里
那条平坦的小道就是它们最理想的屋基。每到春天，它们就成群结队地来到这个地方安
营扎寨。每群数量不一，最大的差不多有上百只黄蜂。这地方简直成了它们的大都市。
每只蜜蜂都有自己单独的一个房间。这个房间除了它自己以外，谁也不可以进去。
如果有哪只不识趣的蜜蜂想闯进别人的房间，那么主人就会毫不客气地给它一剑。因此，
大家都各自守着自己的家，谁也不冒犯谁，这个小小的社会充满了和平的气氛。
一到四月，它们的工作就不知不觉地开始了。唯一可以显而易见地证明它们在工作
的，是那一堆堆新鲜的小土山。至于那些劳动者，我们外人是很少有机会看到的。它们
通常是在坑的底下忙碌着，有时在这边，有时在那边。我们在外面可以看到，那小土堆
渐渐地有了动静，先是顶部开始动，接着有东西从顶上沿着斜坡滚下来，一个劳动者捧
着满怀的废物，把它们从土堆顶端的开口处抛到外面来，而它们自己却用不着出来。
五月到了，太阳和鲜花带来了欢乐。四月的矿工们，这时已经演变成勤劳的采蜜者
了。我们常常看到它们满身披着黄色的尘土停在土堆上，而那些土堆现在已变得像一只
倒扣着的碗了，那碗底上的洞就是它们的入口了。
它们的地下建筑离地面最近的部分是一根几乎垂直的轴，大约有一支铅笔那么粗，
在地面下约有六寸到十二寸深，这个部分就算是走廊了。
在走廊的下面，就是一个小小的巢。每个小巢大概有四分之三寸长，呈椭圆形。那
些小巢有一个公共的走廊通到地上。
每一个小巢内部都修聋得很光滑，很精致。我们可以看出一个个淡淡的六角形的印
子，这就是它们作最后一次工程时留下的痕迹。它们用什么工具来完成这么精细的工作
呢？是它们的舌头。
我曾经试图往巢里面灌水，看看会有什么后果，可是水一点儿也流不到巢里去。这
是因为斑纹蜂在巢上涂了一层唾液，这层唾液像油纸一样包住了巢，在下雨的日子里，
巢里的小蜜蜂就再也不会被弄湿了。
斑纹蜂一般在三、四月里筑巢。那时候天气不大好，地面上也缺少花草。它们在地
下工作，用它的嘴和四肢代替铁锹和把子。当它们把一堆堆的泥粒带到地面上后，巢就
渐渐地做成了。最后用它的铲子——舌头，涂上一层唾液。当快乐的五月到来时，地下
的工作已经完毕，那和照的阳光和灿烂的鲜花也已经开始向它们招手了。
田野里到处可以看到蒲公英、野蔷薇、雏菊花等，在花丛里尽是些忙忙碌碌的蜜蜂。
它们带上花蜜和花粉后，就兴高采烈地回去了。它们一回到自己的城市里，就会立即改
变飞行方式，它们很低地盘旋着，好像对这么多外观酷似的地穴产生了迟疑，不知道哪
个才是自己真正的家。但是没过一会儿，它们就各自认清了自己的记号，很快地，准确
无误地钻了进去。
斑纹蜂也像其它蜜蜂一样，每次采蜜回来，先把尾部塞入小巢，刷下花粉，然后一
转身，再把头部钻入小巢，把花蜜洒在花粉上，这样就把劳动成果储藏起来了。虽然每
一次采的花蜜和花粉都微乎其微，但经过多次的采运，积少成多，小巢内已经变得很满
了。接着斑纹蜂就开始动手制造一个个“小面包”，“小面包”是我给那些精巧的食物
起的名字。
斑纹蜂，开始为它未来的子女们预备食品了。她把花粉和花蜜搓成一粒粒豌豆大小
的“小面包。”这种“小面包”和我们吃的小面包大不一样：它的外面是甜甜的蜜质，
里面充满了干的花粉，这些花粉不甜，没有味道。这外面的花蜜是小蜜蜂早期的食物，
里面的花粉则是小蜜蜂后期的食物。
斑纹蜂做完了食物，就开始产卵。它不像别的蜜蜂产了卵后就把小巢封起来，它还
要继续去采蜜，并且看护它的小宝宝。
小蜜蜂在母亲的精心养护和照看之下渐渐长大了。当它们作茧化蛹的时候，斑纹蜂
就用泥把所有的小巢都封好。在它完成这项工作以后，也到了该休息的时候了。
如果没有什么意外发生的话，在短短的两个月之后，小蜜蜂就能像它们的妈妈一样
去花丛中玩耍了。

温厚长者和小强盗
可是斑纹蜂的家并不像想象中那样安逸，在它们周围埋伏着有许多凶恶的强盗。其
中有一种蚊子，虽然小得微不足道，却是矿蜂的劲敌。
这种蚊子是什么样的呢？它的身体不到五分之一寸长，眼睛是红黑色的；脸是白色
的；胸甲是黑银灰色的，上面有五排微小的黑点儿，长着许多刚毛；腹部是灰色的；腿
是黑色的，像一个又凶恶又奸诈的杀手。
在我所观察到的这一群蜂的活动范围内，就有许多这样的蚊子。这些蚊子在太阳底
下时能找到一个隐蔽的地方潜伏起来。等到斑纹蜂携带着许多花粉过来时，蚊子就紧紧
地跟在它后面，跟着打转、飞舞。忽然，斑纹蜂俯身一冲，冲进自己的屋子。立刻，蚊
子也跟着在洞口停下，头向着洞口，就这样等了几秒钟，蚊子纹丝不动。
它们常常这样面对着面，彼此只隔一个手指那么宽的距离僵持着。但彼此都显得十
分镇定。斑纹蜂这温厚的长者，只要它愿意，它完全有能力把门口那个破坏它家庭的小
强盗打倒，它可以用嘴把它咬阵，可以用刺把它刺得遍体鳞伤，可它并没有这么做。它
任凭那小强盗安然地埋伏在那里。至于那小强盗呢？虽然有强大的对手在它眼前虎视耽
耽，而那可恶的小蚊子尽管知道斑纹蜂只要举手之劳就可以把它撕碎，可它丝毫没有恐
惧的样子。
不久斑纹蜂就飞走了。蚊子便开始行动了。它飞快地进入了巢中，像回到自己的家
里那样不客气。现在它可以在这储藏着许多粮食的小巢里胡作非为了。因为这些巢都还
没有封好。它从从容容地选好一个巢，把自己的卵产在那个巢里。在主人回来之前，它
是安全的，谁也不会来打扰它，而在主人回来之时，它早已完成任务，拍拍屁股逃之夭
夭了。它会再在附近找一处藏身之处，等候着第二次盗窃的机会。
几个星期后，让我们再来看看斑纹蜂藏在巢里的花粉团吧，我们将发现这些花粉团
已被吃得狼藉一片。在藏着花粉的小巢里，我们会看到几条尖嘴的小虫在蠕动着——它
们就是蚊子的小宝宝，在它们中间，我们有时候也会发现几条斑纹蜂的幼虫——它们本
该是这房子的真正的主人，却已经饿得很瘦很瘦了。那帮贪吃的入侵者剥夺了原该属于
它们的一切。这可怜的小东西渐渐地衰弱，渐渐地萎缩，最后竟完全消失了。那凶恶的
蚊子的幼虫就一口一口把这尸体也吞下去了。
小蜜蜂的母亲虽然常常来探望自己的孩子，可是它似乎并没有意识到巢里已经发生
了天翻地覆的变化。它从不会把这陌生的幼虫杀掉，也不会毫不犹豫地把它们抛出门外，
它只知道巢里躺着它亲爱的小宝贝。它认真地小心地把巢封好，好像自己的孩子正在里
面睡觉一样。其实，那时巢里已经什么都没有了，连那蚊子的宝宝也早已趁机飞走了。
多么可怜的母亲啊！

老门警
斑纹蜂的家里如果没有碰到意外，也就是说没有像刚才我说的那样被蚊子所偷袭，
那么它们大约应有十个姐妹。为了节约时间和劳动力，它们不再另外挖隧道，只要把它
们的母亲遗留下来的老屋拿过来继续用就是了。大家都客客气气地从同一个门口进出，
各自做着自己的工作，互不打扰。不过在走廊的尽头，它们有各自的家，每一个家包括
一群小屋，那就是它们自己挖的，不过那走廊是公用的。
让我们来看看它们是怎样来来去去地忙碌的吧。当一只采完花蜜的蜜蜂从田里回来
的时候，它的腿上都沾满了花粉。如果那时门正好开着，它就会立刻一头钻进去。因为
它忙得很，根本没有空闲时间在门口徘徊。有时候会有几只蜜蜂同时到达门口的情况，
可那隧道的宽度又不允许两只蜂并肩而行，尤其是在大家都满载花粉的时候，只要轻轻
一触就会把花粉都掉到地上，半天的辛勤劳动就都白废了。于是它们定了一个规矩：靠
近洞口的一个赶紧先进去，其余的依次在旁边排着队等候。第一个进去后，第二个很快
地跟上，接着是第三个，第四个，第五个……大家都排着队很有秩序地进去。
有时候也会碰到这样的情况：一只蜂刚要出来，而另一只正要进去。在这种情况下，
那只要进去的蜂会很客气地让到一边，让里面的那只蜜蜂先出来，每只蜜蜂在自己的同
类面前，都表现得非常有风度，有礼貌。有一次我看到一只蜂已经从走廊到达洞口，马
上要出来了，忽然，它又退了回去，把走廊让给刚从外面回来的蜂。多有趣啊！这种互
助的精神实在令人佩服，有了这样一种精神，它们的工作才能很快地进行。
让我们把眼睛睁大些仔细地观察，还有比这更有趣的事呢！当一只蜜蜂从花田里采
了花粉回到洞口的时候，我们可以看到一块堵住洞口的活门忽然落下，开出一条通路来。
当外来的蜂进去以后，这活门又升上来把洞口堵住。同样，当里面的蜜蜂要出来的时候，
这活门也是先降下，等里面的蜜蜂飞出去后，又升上来关好。
这个像针筒的活塞一般忽上忽下的东西究竟是什么呢？这是一只蜂，是这所房子的
门卫。它用它的大头顶住了洞口。当这所房子的居民要进进出出的时候，它就把“门栓”
一拔，也就是说，它立刻退到一边，那儿的隧道特别宽大，可以容得下两只蜂。当别的
蜜蜂都通过了，这“门警”又上来用头顶住洞口。它一动不动地守着门，那样的尽心尽
责，除非它不得不去驱除一些不知好歹的不速之客，否则它是不会擅自离开岗位的。
当这位门警偶尔走出洞口的时候，让我们趁机仔仔细细地看看它吧。我们发现它和
其它蜂一样，不过它的头长得很扁，它的衣服是深黑色的，并且有着一条条的纹路。身
上的绒毛已经看不出来了，它本该有的那种美丽的红棕色的花纹也没有了。这一套破碎
的衣服似乎告诉了我们一切。
这一只用自己的身躯顶住门口充当老门警的蜜蜂看起来比谁都显得沧桑和年老。事
实上它正是这所屋子的建筑者，现在的工蜂的母亲，现在的幼虫的祖母。就在三个月之
前，它还挺年轻的，那时候它正在独自辛辛苦苦地建筑这座房子。现在它算是告老退休
了——不，这不是退休，它还要发挥它的余热，用它的全力来保护着这个家呢。
你还记得那多疑的小山羊的故事吗。它从门缝里往外张望一下，然后对门外的狼说，
“你是我们的妈妈吗？请你把白腿伸给我看，如果你的腿是黑色的，我们就不开门。”
我们这位老祖母的警惕绝不亚于那小山羊。它对每一位来客说：
“把你的蜜蜂黑脚伸给我看，否则我就不让你进来。”
只有当它认出这是它家的一员时，它才会开门，否则它是决不会让任何外客进入到
它家里去的。
你看！在洞旁走过一只蚂蚁，它是一个大胆的冒险家。它很想知道这个散发着一阵
阵蜂蜜香味的地方究竟是怎样的。
“滚开！”老蜜蜂摇了摇头说道。
蚂蚁被它吓了一跳，悻悻地走开了。也幸亏它走开了。如果它仍逗留在蜂房旁的话，
老蜜蜂就要离开它的岗位，飞过去不客气地追击它了。
也有一种不擅长挖隧道的蜜蜂，也就是樵叶蜂，它要寻找人家从前挖掘好的隧道。
斑纹蜂的隧道对它再适合不过了。那些以前受蚊子偷袭，被蚊子占据的斑纹蜂的巢一直
是空着的。因为蚊子让它们家绝了后，整个家都已经败落了。于是樵叶蜂就可以顺理成
章地占据这个空巢，来个废物利用了。为了找到这样的空巢，以便于让它们放那些用枯
叶做成的蜜罐，这帮樵叶蜂常到我的这种斑纹蜂的领地里来巡视。有时候它似乎找到目
标了，可还没等它的脚站稳，它的嗡嗡声已引起了门警的注意。门警立刻冲出洞来，在
门口作了几下手势，告诉它这洞早就有主人了。樵叶蜂明白了它的意思，立即飞到别处
去找房子了。
有时候没等门警出来，樵叶蜂已经迫不及待地把头伸了进去。于是做门警的老祖母
立刻把头顶上来塞住通路，并且发出一个并不十分严厉的信号，以示警告，樵叶蜂立即
明白了这屋子的所有权，很快就离开了。
有一种“小贼”，它是樵叶蜂的寄生虫，有时候会受到斑纹蜂的教训。有次我亲眼
看到它受了一顿重罚。这鲁莽的东西一闻进隧道便为非作歹，以为自己进了樵叶蜂的家
了。可是不一会儿，它立刻发现自己犯了一个大错误，它闯进的是斑纹蜂的家。它碰到
了守门的老祖母，受了一顿严厉的惩罚。于是它急急忙忙地往外逃。同样，其它野心勃
勃又没有头脑的傻瓜，如果想闯进斑纹蜂的家，毫无疑问它将受到同样的待遇。
有时候守门的蜜蜂也会和另外一位老祖母发生争执。七月中旬，是蜜蜂们最忙的时
候。这时候我们会看到两种迥然不同的蜂群：那就是老蜜蜂和年轻的母蜂。年轻的母蜂
又漂亮又灵敏，忙忙碌碌地从花间飞到巢里，又从巢里飞向花间。而那些老蜂，失去了
青春，失去了活力。只是从一个洞口踱到另一个洞口。看上去就像迷失了方向找不到自
己的家。这些流浪者究竟是谁呢？它们就是那些受了可恶的小强盗的蒙骗而失去家庭的
老蜜蜂。当初夏来临的时候，老蜜蜂终于发现从自己的巢里钻出来的是可恶的蚊子，这
才恍然大悟、痛心疾首。可是这已经太晚了，它已经变成了无家可归的孤老。它只好委
屈地离开自己的老家，到别处去另谋生路了。看看哪一家需要一个管家或是需要一个门
警。可是那些幸福的家庭早已有了自己的祖母来打点一切了。而这些老祖母往往对外来
找工作，抢自己饭碗的老蜜蜂心存敌意，往往会给它一个不客气的答复。的确，一个家
只需要一个门警就早已足够了。来了两个的话，反而把那原本就不宽敞的走廊给堵住了。
有些时候，两个老祖母之间真的会发生一阵恶斗。当流浪的老蜜蜂停在别家门口的
时候，这家的看门老祖母一方面紧紧守着门，一方面张牙舞爪地向外来的老蜂挑战，而
胜的那一方，往往是那身心疲惫，悲伤孱弱的老孤蜂。
这些无家可归的老蜜蜂后来怎样了呢？它们一天一天地衰老下去，渐渐数目也少了
起来，最后全部绝迹了。它们有的是被那些灰色的小蜥蜴吃掉了，有的是饿死的，有的
是老死了，还有的是万念俱灰，心力衰竭而死。
至于那守门的老祖母，它似乎从来不休息，在清晨天气还很凉快的时候，它已经到
达它的岗位，到了中午，正是工蜂们采蜜工作最忙的时候，许许多多蜜蜂从洞口飞进飞
出，它仍旧守护在那里；到了下午，外边很热，工蜂都不去采蜜，留在家里建造新的巢，
这时候，老祖母仍旧在上面守着门。在这种闷热的时候，它连瞌睡都不打一下，它不能
打瞌睡，这个家的安全都靠它了。
到了晚上，甚至是深夜，别的蜜蜂都休息了，它还像白天一样忙，防备着夜里的盗
贼。
在它小心的守护下，整个蜂巢的安全可以一直持续到五月以后。如果那蚊子来抢巢，
让它来吧，老祖母会立即冲出去和它拼个你死我活。但它们不会来。因为在明年冬天到
来之前，它们还是躲在茧子里的蛹。
虽然没有蚊子，其它的寄生虫类也不少。它们也很可能来侵犯蜂巢。但是，奇怪的
是，我天天认真观察那个蜂巢，却从没有在它的附近发现什么蜂类的敌人。整个夏天它
都那么安静而平和。可见那些暴徒己深知老祖母的厉害。同时也可见老祖母是如何的警
觉了。
------------------








萤



一、它的外科器具
在昆虫的各种类型中，很少有能够发光的。但其中有一种是以发光而出名的。这个
稀奇的小动物的尾巴上像挂了一盏灯似的，用来表达它对快乐生活的美好祝愿。即便是
我们不曾与它相识，不曾见过它在黑夜中从草丛上飞过，不曾见过它从圆月上落下来。
就像一点点火星儿一样，那么，至少从它的名字上，我们可以多少对它有一些了解。古
代的时候，希腊人曾经把它叫做亮尾巴，是很形象的一个名字。现代，科学家们则给它
起了一个新的名字，叫做萤火虫(Iamqyris)。
事实上，萤无论如何也不是蠕虫，即便是从它的外表上来看，它也不能算作是蠕虫。
它生长着六只短短的腿，而且，它能够知道如何去利用这些短足。从某种意义上讲，它
可真算得上一位真正的闲游家。雄性的萤到了它发育完全的时候，会生长出翅盖，像真
的甲虫一样。不过，事实上，它也就是甲虫的一种。然而，雌性的萤却不怎么引人注意，
它对于飞行的快乐，却是一无所知的。真的有些可怜，它不能完全懂得世界上还有自由
飞行这种快乐可以享受。它终身都处于幼虫的状态，也就是说处于一种不完全的形状，
似乎永远也长不大。可是，即便在这种状态下，使用“蠕虫”这个名词也是很不贴切，
不得当的。我们法国人，经常用“像蠕虫一样的粉光”一语来表示这些没有一点儿保护
和遮掩的动物，不过现在，萤却是有衣服的。可以说，外皮就是它的衣服，它用自己的
外皮来保护自己。而且，它的外皮还具有很丰富的颜色呢。它全身是黑棕色的，但是胸
部有一些微红。在它的身体的每一节的边沿部位，还装饰着一些粉红色的斑点。像这样
色彩丰富的衣服，蠕虫是不会穿的！
尽管如此，我们还是继续把它叫做发光的蠕虫。因为，它的这个名字是全世界的人
早就已熟知的了。
萤，有两个最有意思的特点：第一、就是它获取食物的方法；第二、就是它的尾巴
上有灯。
有一位著名的法国研究食物的科学家曾经说过：“告诉我，你吃的是什么东西，那
么我就会告诉你，你究竟是什么东西。”
同样的问题，都应该对任何昆虫提出，我们想要研究的东西就是这些昆虫们的生活
习性——因为，有关昆虫的食品供给方面的知识，是动物生活中最主要的问题，就是我
们人类常说的“民以食为天”，它也就成了我们应该重点研究的问题之一。虽然，从萤
的外表来看，它似乎是一个纯洁善良而可爱的小动物，但是，事实上，它却是一个凶猛
无比的食肉动物。它是一个善于猎取山珍野味的猎人，而且，它的捕猎方法是十分凶恶
的。看来，它的外表也像其他一些昆虫一样具有一定的欺骗性。通常，它的俘虏对象主
要是一些蜗牛。这个事实，早就被人们认识到了。而人们所不知道的，鲜为人知的，只
是萤的那种有些稀奇古怪的捕捉食品的方法。这个方法，我在其他的地方还没有看到过
相同的例子，可见其独特性并非一般。
萤的这种捕食方法，具体情况如下所述：在它开始捉食它的俘虏以前，它总是先要
给它打一针麻醉药，使这个小猎物失去知觉，从而也就失去了防卫抵抗的能力，以便它
捕捉并食用。这就好比我们人类在动手术之前，在病床上，先接受麻醉，从而渐渐失去
知觉而不感到疼痛一样。在一般情况下，萤所猎取的食物，都是一些很小很小的蜗牛。
很少能捕捉到比樱桃大的蜗牛。在气候非常炎热的时候，就会在路旁边的枯草或者是麦
根上，聚集着大群的蜗牛，像集体纳凉一般。也许是酷热难耐的原因，它们一动也不动
地群伏在那些地方。生怕动一动，就会觉得热气逼人。它们就是这样静止着，懒洋洋地
经过炎热的夏天的。在这些地方，我常常会看到，一些萤在咀嚼它们那已经失去知觉的
俘虏。萤就是在这些摇摆不定的物体上把它们麻醉了的。
除了上述路边的枯草、麦根等地方以外，萤也常常选择一些其他可以获得食物的地
方出没或停留。比如说，它常到一些又凉又潮湿的阴暗的沟渠附近去蹓跶。因为在这些
地方，经常会有一些杂草丛生。在那里经常可以找到大量的蜗牛。这可是难得的美餐啊！
饱饱地吃上几顿是没问题的。通常在这些地方，萤把它们的俘虏在地上杀死，就像人们
说的就地处决一样，干净利落地结束战斗。然后，获得丰厚的战利品。在我自家的屋子
里，我也可以经常制造出这样的地方，来吸引萤到这里来捕食。因此，待在家里，我便
能制造出一个战场，并且，更便了我非常仔细地观察它们的一举一动。
那么，下面我就来叙述一下这种奇怪的情形吧。我拿了一个大玻璃瓶，然后再找一
点儿小草，把草放到大玻璃瓶子里面，再往里边放进几只萤，还有一些蜗牛。我取的蜗
牛，大小比较适中，不算特别的大，也不算特别的小。这一切准备工作就绪以后，我们
所需要继续进行的工作，就是等待，而且，必须要耐心地等待。不过，最为重要的一点
是必须十分留心，时刻注视着玻璃瓶中发生的一切动静，哪怕是最微小的动作，也不能
轻易放过。因为，整个事情的发生，是在非常不经意的时候，而且，时间持续也不长，
几乎就是一会儿的时间。所以，必须目不转睛地紧紧盯住瓶中的这些生灵。
不一会儿，玻璃瓶中就有事情要发生了。萤已经开始注意到它的牺牲品了。看起来，
蜗牛对于萤而言，有极强的，难以抗拒的吸引力。按照通常情况下，蜗牛的习性，除去
外套膜的边缘的地方，它的身体会微微露出一点儿以外，其余的躯体全部都隐藏在它的
家中——即它背上的壳子里面，可能它觉得这样会更安全一些。于是，这位猎人跃跃欲
试，准备发起总攻了。它先做的事情，就是把自己身上随身携带着的兵器迅速地抽出来。
这件兵器是何等地细小啊，要是没有放大镜的帮助，简直是一点儿也看不出来。萤的身
上长有两片颚，它们分别弯曲起来，再合拢到一起，就形成了一把钩子，一把尖利、细
小，像一根毛发一样的钩子。如果把它放到显微镜下面观察，就可以发现，在这把钩子
上有一条沟槽。如此而已，这件武器并没有什么其他更特别的地方。然而，这可是一件
有用的兵器，是可以致对手于死地的夺命宝刀。
这个小小的昆虫，正是利用这样一件兵器，在蜗牛的外膜上，不停地、反复地刺击。
但是，萤所表现出来的态度很平和，神情也很温和，并不像恶狠般凶猛，乍一看起来，
好像并不是猎人在捕猎食物，在咬它的俘虏，倒好像是两个动物在亲呢接吻一般。当小
孩子在一起互相戏弄对方的时候，他们常常用两个手指头，互相握住对方的皮肤，轻轻
地揉搓。这种动作，一般情况下，我们常用“扭”这个字眼儿来表示。因为，事实上，
这种动作近乎相互搔痒，而并不是那样重重地打。现在，就让我们来使用“扭”这个字
吧。一提到动物，除去那些最简单平实的字以外，通常在言语中使用的字，可以说，好
多都没有用处。那么，我们就可以说，萤是在“扭”动蜗牛，这大概更贴切一些。
萤在扭动蜗牛时，颇有它自己的方法。你会看到它一点儿也不着急，不慌不忙，很
有章法。它每扭动一下对方，总是要停下来一小会儿。仿佛是要审看一下，这一次扭动
产生了何等效果一般。萤扭动的次数并不是很多，顶多有五六次足以。就这么几下，就
能让蜗牛动弹不得，失去了一切知觉而不醒人世。再后来，也就是在萤开始吃战利品的
时候，还要再扭上几下。看起来，这几下扭动更至关重要。但是，至于究竟萤为什么要
如此这般行事，我就不能确定其真正的原因了。确实在最初的时候，只要轻轻地几下，
就足以使蜗牛慢慢地不能感觉一切了。那么，它为什么在食用前，还要来上几下呢？我
不得而知，至今仍是个迷。由于萤的动作非常迅速而敏捷，如同闪电一般快，就已经将
毒汁从沟槽中传送到蜗牛的身上了。这只是一个瞬间的动作，要非常仔细地观察才能觉
察到。
当然，有一点是不用怀疑的，那就是，在萤对蜗牛进行刺击时，蜗牛一点也不会感
觉到痛苦。关于这一点，我曾经做过一次小小的试验。在一只萤进攻一只蜗牛的时候，
当萤只扭了四五次以后，我马上迅速地把那只受了毒汁破害的蜗牛拿开。然后，用一根
很小很小的针去刺激这可怜虫的皮肤。那一点儿被刺伤了的肉，竟然一点儿也没有收缩
的迹象。这就已经很清楚地表明，此时此刻，这只蜗牛已经一点儿活气也没有了。它是
不会感觉到痛苦的。因为它已经沉浸到另一个世界里去了。还有一次，我非常偶然地看
到一只可怜的蜗牛遭受到萤的攻击。当时，这只蜗牛正在向前自由自在地爬行着。它的
足慢慢地蠕动着，触角也伸得很长。忽然，由于一下子的刺激和兴奋，这只蜗牛自己乱
动了几下。但是，马上这一切就静止了下来，它的足也不再向前慢爬了。整个身体的前
部也全然失去了它刚才的那种温文尔雅的曲线。它的长长的触角也变得软了，不再向上
伸展着了，而是拖垂到下边来，就像一根已经坏了的手杖一样，再也感受不到什么东西
了。从这种种表面的现象来看，这只蜗牛已经死了，已经真的到另一个世界里去了。
然而，实际上，这只蜗牛并没有真正悲惨地死去。我完全有办法，能让它重新活过
来，我可以给它第二次生命的机会。就在这个可怜的、假死的蜗牛既不生，又不死的两
三天内，我每天都坚持给它洗浴，清洁身体，特别是伤口。就在几天以后，奇迹出现了。
这只被萤无情地伤害得很惨重的、几乎一命呜呼了的家伙，恢复到了以前的状态中，它
已经能够自由地爬来爬去了。而且，它的知觉也已经恢复正常了。因为当我用小针刺击
它的肉时，它立刻就会有反应，小小的躯体马上就会缩到背壳里去藏了起来，这充分说
明它已经恢复知觉了，就像什么都不曾发生过一样。它是完全可以爬行了，那对长长的
触角重新又伸展开来，好像并没有发生过什么特别意外的事情一样。而且，它还精神倍
增。在它失去知觉的日日夜夜里，它仿佛进入了一种什么都不知晓的沉醉的状态，一切
都惊动不了它，而现在则大不一样了。它醒了，而且完全苏醒了过来，从死亡中复活了，
奇迹般地逃离魔爪，获得了第二次生命。
在人类社会的科学中，人们已经发明创造了在外科手术时不会让人感觉到疼痛的方
法，并且这种方法在医学实践上，已经被证明是非常成功的了。然而，在人类还没有找
到这种方法之前，萤以及很多其他种类的动物，就已经通过实践，实行了好几个世纪之
久了。所不同的是，外科医生在手术前，让我们闻以太或者是其他麻醉剂。而那些昆虫
采用的方法是，利用它们天生就长着的毒牙，向别的动物注射极小量的特别的毒药，以
达到让别的动物失去知觉的目的。
当我们偶然想到蜗牛具有那样温柔、平和而无害的天性，可萤却要采取向它注射毒
汁以麻醉它的特别方法来制服它，并且以它为食物，似乎总有些奇怪的感觉。但是，我
想我可以知道萤利用这种方法猎取蜗牛的种种鲜为人知的理由。
假如蜗牛只是在地上爬行，甚至是蜷缩在自己的壳子里，那么，对于它的种种攻击，
是一点儿也不困难的。原因是，蜗牛背上背的壳儿上并没有任何遮盖的东西加以保护。
而且，蜗牛身体的前部是完全暴露在外面的，也是毫无遮挡的。但是，实际上，事情并
不这么简单。蜗牛不仅仅单单地在上爬行，它经常是置身在较高的而且不稳定的地方。
比如，它喜欢趴在草杆的顶上，或者是待在很光滑的石面上。它贴身在这些地方，就无
需什么其他的保护了。因为，这些地点本身就为它提供了再好不过的天然保护。其原因
是，当蜗牛把自己的身体紧紧地依附在这些东西上时，这些东西就起到了盖子的作用。
于是，便能起到很好的保护作用了。不过，只要稍微一不留神，有一点儿没有遮盖严密，
一旦被萤发现，它的钩子可一点儿也不讲情面。只要有机可乘，它是一定要钻空子的。
总之，萤的钩子总会有办法可以触及到蜗牛的身体。然后，一下子钩住，释放出毒液，
蜗牛便会失去知觉。萤就可以安安稳稳地找个地方坐下来，享用它的美餐了。
不过，蜗牛深居高处，也是有一定危险的。当它爬在草杆上时，很容易掉下来。哪
怕稍微有一点儿扭动，或者是挣扎，蜗牛就会移动它的身体。一旦蜗牛落到地面上，那
么，萤就不得不去选择一个更好的猎物。所以，在萤捕捉蜗牛时，必须要使它没有丝毫
的痛苦感，失去知觉，让它动晃不得，从而也就不能轻易地逃跑了。因此，萤在进攻蜗
牛时，一定要采取触得很轻微的方法，以免惊动了这只蜗牛，让它从摇摆不定的草杆上
摇落下来。否则的话，可就前功尽弃，白费了一番心思，让到手的食物给溜走了。因此，
我想，萤之所以具备这样的外科器具，理由不过如此而已吧！除此以外，我再也想不到
其他的理由了。

二、蔷薇花饰物
萤不仅仅是在草木的枝干上结束战斗，使它的俘虏逐渐失去知觉，而且，也在这种
存在一定危险性的地方，就地把它解决处理掉，也就是要把它给完全吃掉。所以，萤的
食品的获得可并不是件简单容易的事呢！
那么，萤在吃蜗牛时，又是采用怎样奇特方法呢？它真的是在食用它吗？是不是要
先把蜗牛分割成一片一片的，或者是割成一些小碎片或碎粒什么的，然后再去慢慢地、
细细地咀嚼品味它呢？我猜想，它并不是以这样普通的方式食用它的。因为，我从来也
没有在这些动物体内，找到过任何这种小粒的食物。这就证明萤的吃，并不是通常的狭
义上的吃的意思，它只不过是以另一种方式，来解决问题罢了。具体方法是这样的，它
要将蜗牛先制造成非常稀薄的肉粥，然后才开始饮用。就像蝇吃小的幼虫一样，它能够
在还没有吃之前，先把它弄成流质，然后再痛快地享用。
更为具体的情形和做法是这样的。萤先使蜗牛失去知觉，无论蜗牛的身体大小如何。
在开始的时候，总是常常只有一只的四分之一大小。客人们也二三两两地跑过来了。它
们和主人毫无争吵，全部聚集到一起，准备和主人一起分享食物。过了两三天以后，如
果把蜗牛的身体翻转过来，把它的面孔朝下面放置，那样，它体内盛的东西，就会像锅
里的羹一样流出来。这个时候，萤的膳食已经结束了。它所饮下的只不过是一些其他动
物已经吃剩下的东西。因而，一只蜗牛被众虫同时分享了。
事实是很显然的。和前面我们已经看到过的“扭”的动作相似，它们经过几次轻轻
的咬，蜗牛的肉就已经变成了肉粥。然后，许多客人一起跑过来共同享用。很随意地，
每一位客人都一口一口地把它吃掉。而且，每一位客人都利用自己的一种消化素把它做
成汤。能够应用这样一种方法，说明萤的嘴是非常柔软的。萤在用毒牙给蜗牛注射毒药
的同时，也会注入——引进其他的无物质到蜗牛的体内，以便蜗牛身上固体的肉能够变
成流质。这样一来，这种流质很适合萤那柔软的嘴，使它吃得更加方便自如。
蜗牛被关闭在我的玻璃瓶里，虽然有的时候，它所处的地位不是特别稳固，但是，
它还是非常仔细小心的。有的时候，蜗牛爬到了瓶子的顶部，而那顶口是用玻璃片盖住
的。于是，它为了能在那里停留得更加稳固、踏实一些，它就利用那自己随身携带着的
粘性液体，粘在那个玻璃片上。这样一来，的确是非常稳定安全的。不过一定要多用一
些粘液，不然的话，哪怕稍微少用了一点儿粘液，都将是十分危险的。即便是微微地动
一点儿，也足以使它的壳脱离那个玻璃片，掉到瓶子底下去。
萤常常要利用一种爬行器——为了弥补它自己腿部，以及足部力量的不足——爬到
瓶子的顶部去，先仔细地观察一下蜗牛的动静，然后，作一下判断、和选择，寻找可以
下来的地方。然后，就那么迅速地轻轻一咬，就足以使对手失去知觉了。这一切都发生
在一瞬间。于是，一点儿也不拖延，萤开始抓紧时间来制造它的美味佳肴——肉粥，以
准备作为数日内的食品。
当萤一阵风卷残云以后，便吃得很饱了。剩下的蜗牛壳也就完全空了。但是，这个
空壳依然是粘在玻璃片上的，并没有脱落到瓶底上来。而且，壳的位置也一点儿都没有
改变，这都是粘液作用的结果。那个牺牲了的隐居者一点儿也不加以反抗，就这样静悄
悄地，不知不觉地任人宰割，最终，变成了别人嘴里营养丰富、美不胜收的大餐。就在
它那受到攻击的原地，逐渐流干了身体的全部，成了一个空空如也的壳儿。这种详细的
情形，向我们表明了这样一个事实，即萤的这种麻醉式的咬伤，是何等的有效。因此，
可以说，萤处理蜗牛的方法是十分的巧妙的。
萤要想顺利地完成自己的任务，实现自己的目的，比如，爬到悬在半空中的玻璃片
上去，或者是爬行到草杆上去，必须要具备一种特别的爬行足或其他什么有利的器官，
以便使它自己不至于在还未触及到猎物时，就先从高空跌落下来，从而半途而废。显然
它现有的笨拙的足是不够用的，这就决定它需要辅助的东西。
把一只萤放到放大镜下面进行仔细的观察研究，我们就可以很容易地发现，在萤的
身上，的确生长着这种特别的器官。大自然在创造它的时候很公平，非常细心，并没有
忘掉赐给它必要的工具。在萤的身体下面，接近它尾巴的地方，有一块白点，通过放大
镜可以清楚地看到。这主要是由一些一打以上的短小的细管，或者是指头组合而成的。
有的时候，这些东西合拢在一起形成为一团，而有的时候，它们则张开，成蔷薇花的形
状。就是这精细的结构，这些隆起来的指头，帮助了萤，使得它能够牢牢地吸伏在非常
光滑的表面上，与此同时，还可以帮助它向前爬行。如果萤想使自己紧紧地吸到玻璃片
上，或者是草杆上，那么，它就会放开那些指头，让蔷薇花绽放开来。在支撑物上，这
些指头放开得很大。萤就利用它自己自然的粘力而牢固地附着在那些它想停留的支撑物
体上。而且，当萤想在它所待的地方爬行时，它便让那些指头相互交错地一张一缩。这
样一来，萤就可以在看起来很危险的地方自由地爬行了。
那些长在萤身上的，构成蔷薇花形的指头，是不长节的，但是，它们每一个都可以
向各个不同的方向随意地转动。事实上，与其说它们像是指头，倒不如说它们更像一根
根细细的管子。因为，这一个比喻要更加合适，贴切一些。要是说它们像指头的话，它
们却并不能拿起什么东西。它们只能是利用其粘附力而附着在其他东西上。它们的作用
很大，除掉粘附，以及在危险处爬行这两大功能外，它还具有第三种功能，那就是它们
能当海绵以及刷子使用。在萤饱餐一顿以后，当它休息的时候，它便会利用这种自动的
小刷子，在头上、身上到处进行扫刷和清洁工作，这样既方便，又卫生，它之所以能够
如此自如地利用身体的这一器官，主要是因为那刺有着很好的柔韧性，使用起来相当便
利。在它饱餐之后，舒舒服服地休息一下，再用刷子一点一点，从身体的这一端刷到另
外一端，而且非常仔细、认真，几乎哪个部位都不会被遗漏掉。可以说，它是一种非常
爱清洁，注意文明修身的小动物。从它那副神采奕奕，得意洋洋，而舒服的表情来判断，
这个小动物对这个清理个人卫生的事情还是非常重视的，也非常的有兴趣去做的。刚一
开始的时候，我们当然会产生某种疑问：为什么这个小东西在拂拭自己的时候，是如此
的专心致志，而且如此地当心呢？答案是显而易见的。把一只蜗牛做成一顿肉粥，而且
花费了很多心思，用很多天的工夫去饮食它，肯定会把自己的身体弄得出奇的肮脏，那
么，认认真真地在饱餐之后，把自己的身体好好清洗一番，让自己焕然一新，是很有必
要的。

三、它的灯
如果，我们的萤，除了利用那种类似于接吻一样的动作——轻轻地扭动几下，来施
行麻醉术以外，就再也不具备什么其他的才能了，那么它的名声就不会有如此之大了，
以至于所有的人都知道它的大名。因此，它必定还具有一些其他的特殊本领，与其他动
物区别开来，比如特异功能什么的。那究竟它还有什么样的奇特本领呢？
众所周知，它的身上还带有一盏灯。它会在自己的身上点燃这盏灯。在黑夜中为自
己留一盏灯，照耀着自己行进的路程。这就是它成名的最重要的原因之一了。
雌性萤那个发光的器官，生长在它身体最后三节的地方。在前两节中的每一节下面
发出光来，形成了宽宽的节形。而位于第三节的发光部位比前两节要小得多，只是有两
个小小的点，发出的光亮可以从背面透射出来，因而在这个小昆虫的上下面都可以看得
见光。从这些宽带和小点上，发出的光是微微带蓝色的、很明亮的光。
而雄性的萤则不一样，它与雌萤相比，只有雌黄那些灯中的小灯，也就是说，只有
尾部最后一节处的两个小点。而这两个小点，在萤类的全族之中，差不多全都具备，从
萤还处于幼小的蛴螬的时代开始，就已经具备这两个用于发光的小点。此后，随着萤的
成长，它们也随着身体的生长不断地长大。在萤的一生中都不改变。这两个小点，经常
是无论在身体的上面，还是下面，都可以看见。但是雌萤所特有的那两条宽带子则不同，
它只能在下面发光的。这就是雄、雌的主要区别之一。
我曾经在我的显微镜下，观察过这两条发光的带子。在萤的皮上，有一种白颜色的
涂料，形成了很细很细的粒形物质。于是，光就是发源于这个地方。在这些物质的附近，
更是分布着一种非常奇特的器官，它们都有短干，上面还生长着很多细枝。这种枝干散
布在发光物体上面，有时还深入其中。
我很清楚地知道，光亮是产生于萤的呼吸器官的。世界上有一些物质，当它和空气
相混合以后，立即便会发出亮光，有的时候，甚至还会燃烧，产生火焰。此等物质，被
人们称为“可燃物”。而那种和空气相混合便能发光或者产生火焰的作用，则通常被人
们称之为“氧化作用”。萤能够发光，便是这种氧化作用的一个很好的例证与说明。萤
的灯就是氧化的结果。那种形如白色涂料的物质，就是经过氧化作用以后，剩下的余物。
氧化作用所需要的空气，是由连接着萤的呼吸器官的细细的小管提供的。至于那种发光
的物质的性质，至今尚无人知晓其答案。
但是，另外有一个问题，我们是知道得比较详细的。我们清楚地知道，萤完全有能
力调节它随身携带的亮光。也就是说，它可以随意地将自己身上的光放大一些，或者是
调暗一些，或者是干脆熄灭它。
那么，这个聪明的小动物，究竞是怎样行动才达到它调节自身光亮的目的呢？经过
观察我了解到，如果萤身上的细管里面流入的空气量增加了，那么它发出来的光亮度就
会变得更强一些；要是哪天萤不高兴了，把气管里面的空气的输送停止下来，那么，光
的亮度自然就会变得很微弱，甚至是熄灭了。
一些外界的刺激，将会对气管产生影响。这盏精致的小灯——萤身后最后一节上的
两个小点，哪怕只有一点点的侵扰，立刻就会熄灭。这一点我深有体会，每次当我想要
捕捉那些十分幼稚可爱的小动物的时候，它们总是爱和我玩捉迷藏的游戏。我明明就在
刚才，清清楚楚地看见它在草丛里发光，并且飞旋着，但是，只要我的脚步稍微有一点
儿不经意，发出一点儿声响，或者是我不知不觉地触动了旁边的一些枝条，那个光亮立
刻就会消失掉，这个昆虫自然也就不见了。我也就失去了捕捉对象，又浪费了一次机会。
然而，雌萤的光带，即便是受到了极大的惊吓与扰动，都不会产生多么大的影响。
比如说，把一个雌萤放在一个铁丝做的笼子里面，空气是完全可以流通的。然后，我们
再在铁笼子旁边放上一枪。就是这样爆烈的声音，竟然也毫无结果。萤似乎什么也没有
听到，或是听到了，也置之不理。它的光亮依然如故，丝毫变化都没有。于是，我又换
了一种方法试探。我取了一个树枝，而且还把冷水洒到它们的身上去，但是，这种种方
法都失败了。各种刺激居然都不奏效。没有一盏灯会熄灭，顶多是把光亮稍微停一下。
但是，这种情况是很少发生的。然后，我又拿了我的一个烟斗，往铁笼子里吹进一阵烟
去。这一吹，那光亮停止的时间长久了一些。还有一些竟然停熄掉了。但是，即刻之间
便又点着了。等到烟雾全部散去以后，那光亮便又像刚才一样明亮了。假如把它们拿在
手掌上，然后轻轻地一捏。只要你捏得不是特别的重，那么，它们的光亮并不会减少得
很多。总之，到目前为止，我们根本就没有什么办法，能让它们全体熄灭光亮。
从各个方面来看，毫无疑问，萤千真万确地能够控制并且调节它自己的发光器官，
随意地使它更明亮，或更微弱，或熄灭。不过，在某一种我们还不知道的环境之下，它
也会失去它这种自我调节的能力。如果我们从它发光的地方，割下一片皮来，把它放在
玻璃瓶或管子里面，虽然并没有像在活着的萤体上那么明亮耀眼，但是，它也还是能够
从容地发出亮光的。因为，对于发光的物质而言，是并不需要什么生命来支持的。原因
在于，能够发光的外皮，直接和空气相接触而起作用。因此，气管中氧气的流通也就不
必要了。就是在那种含有空气的水中，这层外皮发出的光也和在空气中发出的光同样明
亮。如果是在那种已经煮沸过的水里，由于空气已经被“驱逐”出来了，于是，发出的
光就会渐渐地熄灭了。再没有更好的证据来证明萤的光是氧化作用的结果了。
萤发出来的光，是白色而且平静的。另外，它的光对于人的眼睛一点儿也不刺激很
柔和。这种光看过以后，便会很自然地让人联想到，它们简直就像那种从月亮里面掉落
下来的一朵朵可爱的洁白的小花朵，充满诗情画意的温馨。虽然这种光亮十分灿烂，但
是同时它又是很微弱的。假使在黑暗之中，我们捉住一只细小的萤，然后把萤的光向一
行油印的字上照过去，于是我们便会很容易地辨别出一个一个的字母，甚至也可分辨出
不是很长的词来。不过，超过了这份光亮所涉及到的比较狭小的范围以外，那就什么都
看不清楚了。不过，这样的灯，这样吝啬的光亮，不久就会令读书人厌倦的。
但是，这些能够发出光亮的小动物，这些本该是心中一片光明的小昆虫，在事实上
却是一群心理很黑暗的家伙。它们对于整个家族的感情是完全不存在的。家庭对于它们
而言，是无足轻重的。柔情对于它们也是没有丝毫实际意义的。它们能够随处地产卵。
有的时候，产在地面上；有的时候，产在草叶上。无论何时何地，它都可以随意散播自
己的子孙后代。真可谓四处闯荡，四海为家，随遇而安。而且，在它们产下卵以后，就
再也不去注意它们了，随它们自生自灭，自然生长去了。
从生到死，萤总是放着亮光的。甚至连它的卵也是要发光的。幼虫也是如此。当寒
冷的气候马上就要降临的时候，幼虫就会立刻钻到地面下边去，但是并不钻得很深。假
如我从地面下，把它轻轻地掘起来，它的小灯仍然还是亮着的。就是在土壤的下面，它
的小灯还是点着的，永远为自己留一盏希望的灯！
------------------








被管虫



一、衣冠齐整的毛虫
当春天来临的时候，只要长着一对眼睛，就可以看清楚世界上任何东西的人，在破
旧的墙壁和尘土飞扬的大路上，或者是在那些空旷的土地上，都能够发现一种比较奇怪
的小东西。
那是一个小小的柴束，不知因为什么，它能自己自由自在地行动，一跳一跳地向前
走动。没有生命的东西变成了有生命的东西，不会活动的居然能够跳动了。这究竟是怎
么一回事呢？
这一点的确非常稀奇，而且很令人感到奇怪。不过如果我们靠近些仔细地看一看，
很快就能解开这个谜了。
在那些会动的柴束中，有一条特别漂亮、特别好看的毛虫。在它的身上装饰着白色
和黑色的条纹。大概它是正在寻找自己的食物。也许它是正在寻找一个可以让它安全地
化成蛾的适当的地点。对于它这些让人猜测不适的动作，通过以后它自己的所做所为就
能了解清楚了。
它很懦怯的朝前方急切地行走着，它总是穿着树枝做成的奇异的服装，完全把自己
的身体遮挡住了。只有除掉头和长有六只短足的前部暴露在外面。
它只要受到一点小小的惊动，就会本能地隐藏到这层壳里去，而且一动也不动了。
生怕一不小心被其他的东西侵害了，这显然是它的一种自我保护的本能。
这就是一束柴枝会走动的秘密的答案，它就是柴把毛虫，它是属于被管虫一类的。
为了防御气候的变化，这个既非常害怕寒冷而且又全身裸体的被管虫，建筑起了一
个属于它自己的很轻便的，又很舒服的隐避的场所，一个能够移动的安全的茅草屋。
在它还没有变成蛾的时候，一刻也不敢贸然离开这间茅屋。这确实要比那种装有轮
盘的草屋要好一些，它完全像是由一种特殊的材料制做而成的隐士们穿的保护外衣。
邓内白，山谷里的农夫，穿着一种用兰草带子紧紧地扎住的外衣，而且是羊皮的，
它是皮板朝里，羊皮朝外的。特别是居住在深山里的农夫，尤其是黄土高坡上的农夫，
这种穿着打扮更是常见，他们的头上还要系一条白色的羊肚毛巾。相比较而言，被管虫
的外衣，比这种打扮还要草率一些，因为它们只是拿一个简简单单的柴枝做成一件再朴
素不过的外衣而已，没有任何过多的装饰物品。可见，它们是多么不拘于小节啊！四月
里，在我们家的作坊上面有很多昆虫，在这些墙上能够发现很多的被管虫，它们都向我
提供了十分详尽的常识，如果它是在蛰伏的状态下，这就表示了它们不久就要变成蛾子
了。这是一个最好的机会，它使我能够直接地仔细地观察一下它的柴草的外衣。
这些外衣形状都是一个样子的，真的很像一个纺锤，大约有一寸半那么长。那位于
前端的细枝是固定的，而末端则是分散开的，它们就是这样排列着的，要是没有什么其
他更好的可以当作保护的地方，那么这里就是可以抵挡日光与雨水侵袭的避难场所了。
在并不认识它以前，乍一看上去，它真像一捆普通的草束。不过只是用草束这两个
字还并不能来正确地形容它的样子，因为麦茎实在是很少见得到的。
它的这件外衣的主要的材料是那些光滑的、柔韧的、富有木髓的小枝和小叶，其次
则是那些草叶和柏树的鳞片枝等，最后如果材料不够用了，就采用那些干叶的碎片和碎
枝。
总之，小毛虫遇到什么就使用什么，只要它是轻巧的、柔韧的、光滑的、干燥的、
大小适当的就可以了。所以，它的要求还不算苛刻。
它所利用的材料完全都是依照其原有的形状，一点儿都不加以改变。也就是说既保
持原有材料的性质，又保持原有材料的形状。
一些过长的材料，它也不修整一下，并使其成为适合的、适当的长度。造屋顶的板
条也直接被它拉过来使用。它的工作只不过是把前面固定了就行了。这在它是很简单易
行的。
因为要是想让旅行中的毛虫可以自由地行动，特别是在它装上新枝的时候，仍然能
够使它的头和足可以自如地活动，这个匣子的前部必须用一种特别的方法装置而成。仅
仅是用树枝装饰成的匣子对它而言是不适用的，理由很简单，因为它的枝特别长而且还
很硬实，这就大大防碍了这位勤劳的工人的工作，使它不能正常地尽职尽责。
它所需要的，是必须拥有一个柔软的前部，使得它可以向任何方向自由地转来转去，
从而可以很高兴地完成本职工作。
所以那些硬树枝，在离开毛虫前部相当远的地方，就中止了，取而代之的是一种领
圈，那里的丝带只是用一种碎木屑来衬托，这样一来，也就增加了材料的强度和韧性，
从而不妨碍毛虫的弯曲性。这样一个能够使毛虫自由行动和弯曲的领圈是非常重要，而
且绝对不可缺少的。以至于无论它的做法有怎样的不同，而所有的被管虫都要用到它。
在柴束前部，那张装得下可以自由转动的头部，触摸起来让人觉得很柔软，它的内
部是用纯丝织成的网，外面包裹着绒状的木屑。这种木屑，是毛虫在割碎那些干草的时
候得到的。
我把草匣的外层轻轻地剥掉，将它撕碎，就会发现里面有很多极细的枝干，我曾经
仔细地数过，大概有八十多个呢。在这里面，从靠近毛虫的这一端到那一端，我又发现
了同样的内衣，在把它的外衣打开以前，只有中部与前端是可以看见的，而现在则可以
看到全部了。这种内衣全都是由坚韧的丝做成的，这种丝的韧性是很强的，人用手拉都
不能把它拉断。这是一种光滑的组织，其内部是美丽的白颜色，外部是褐色的而且是有
皱纹的，还有细碎的木屑分散地装饰在上面。
于是，我要看看毛虫是如何做成这件精巧的外衣的了。这件外衣内外共有三层，它
们互相按一定次序叠加在一起。第一层是极细的绫子，它可以和毛虫的皮肤直接相接触；
第二层是粉碎的木屑，用来保护衣服上的丝，并使之坚韧；最后一层是小树枝做成的外
壳。
虽然各种被管虫全都穿上了这种三层的衣服，不过各个种族的外壳却有所不同。比
如，有一种，六月底我在靠近屋子旁边的尘土飞扬的大路上遇见的，它的壳无论从形式
还是从做法上来看，比前边提到过的那一种都要更加高明一些。它外面的厚披是用很多
片材料制做而成的，比如那种空心树杆的断片，细麦杆的小片，还有那些青草的碎叶等。
在壳的前部，简直找不到一点儿枯叶的痕迹。我先前所说的那一种，是常常有的，但那
足以妨碍其美观。在它的背部，也没有什么长的突出物，长出外皮之后，除去颈部的领
圈之外，这个毛虫的全身都武装在那个用细杆做成的壳里面。总体上的差别并不是很大，
不过最显著的一点差异就是它有比较完整的美观。
还有一种身材比较小，衣服穿得比较简易一些的被管虫，在冬天快要结束的时候，
在墙上或树上，在树皮多皱的老树上，比如洋橄榄树，或榆树上，常常可以发现它的踪
迹。当然在其他的地方也会见到。它的壳非常小，常常还不到一寸的五分之二长。它随
意地拾起一些干草，然后把它们平行地粘合起来，除去丝质的内壳以外，这就是它的全
身衣服的材料。
衣服要穿的更经济，更便宜，而且看上去更漂亮、更美丽，那是难度更大的事情了。

二、良母
如果我们在四月的时候捉几条幼小的被管虫，把它们放在铁丝罩子里面，关于它们
的一些事实，我们可以看得更多一点，也可以观察的更清楚一点了。
这时它们中的多数还是处在蛹的时代，等待着有朝一日变成蛾子。但是它们并不都
是那么安份守己地静静地待着，有的比较活跃好动一些，它们会很自豪地慢慢地爬到铁
丝格子上去。在那里，它们会用一种丝质的小垫子，把自身的身体固定好，无论是对它
们而言还是对我而言，都要耐心地等待几个星期，然后，才会有一些事情要发生。
到了六月底的时候，雄性的幼虫从它的壳里跑出来了，它已经不再是什么毛虫了，
而已经变成蛾子了。
这个壳，即一束细杆，你应当记得，它有两个出口，一个在前面，另一个在后面。
前面的一个，是这个毛虫很谨慎、很当心地制做的，是永远封闭着的，因为毛虫要利用
这一端钉在支持物上，以便使蛹得以固定在上面。因此，孵化的蛾必须从后面的口钻出
来，在毛虫还没有变化成蛾子之前，要在壳内先转一个身。然后，才会慢慢地出来。
虽然雄蛾只穿着一件十分简单的黄灰色的衣服，只有苍蝇差不多大小的翼翅，然而，
它却是异常漂亮的。它们长有羽毛状的触须，翼边还挂着细须头。
至于雌蛾，则很少能够在一些比较显眼暴露的地方捕捉到，而且，它们当中的大多
数是很难见得到的。
比别的昆虫迟几天以后，它才会从壳里姗姗来迟地钻出来，其形状简直是难看到了
极点，这个怪物也就是雌蛾。当你刚刚见到它的时候，甚至会惊吓地叫起来。也许它的
样子会吓你一跳的。没有一个人能够马上就看惯眼前这个凄惨的情形。它的难看程度并
不比那些毛虫差些。它没有长翅膀，什么都没有，在它背的中央，连毛也没有，光秃秃、
圆溜溜的。人们简直难得睁开眼睛看它一眼。在它圆圆的有装饰的体端，戴有一顶灰白
色的小帽子，第一节上，在背部的中央，长着一个大大的、长方形的黑斑点——这便是
它身体惟一的装饰物，母被管虫放弃了蛾类所有的一切美丽。这就是雌蛾，这个怪物般
的雌蛾形像。
当它离开它的蛹壳的时候，就在里面生卵。于是，母亲的茅屋(即她的大衣)就留传
给它的后代子孙了。她的卵产的很多，所以这产卵的时间也很长，要经过三十个小时以
上。
产完卵后，它将门关闭起来，使其免受外来的一些侵扰，从而获得一种安全感。为
了达到这个目的，某种填塞物是必要的。于是这位溺爱的母亲，在她一贫如洗、穷困潦
倒的情况下就只能利用它仅有的衣服了。也就是说要利用戴在它体端的那顶丝绒帽子，
塞住门口，以保母子平安，安然无恙地生活。
最后，它所做的还不限于此，它还要拿自己的身体来做屏障。经过一次激烈的震动
以后，它死在这个新屋的门前，留在那里慢慢地干掉，即在死后，它还依然留守在阵地，
为了下一代，死了也甘心。别看它外表上看起来丑陋不堪，但实际上它的内心、它的精
神是很伟大的。
假如破开外面的壳，我们可以看到那里面储存有蛹的外衣，除去前面蛾子钻出来的
地方留下的孔以外，一点也没有受到损坏，雄蛾要从这个狭小的隧道中出来的时候，会
感觉到它的翼和羽毛是很笨重的负担，而且对它形成了一定的阻力。
因此，当毛虫还处在蛹的时代时，就拼命地朝门口奔跑出将近一半的旅程来。最后
终于成功地撞出琥珀色的外衣来，在它的前面，出现了一块开阔的场所，可以允许它自
由地飞行了。
但是，母蛹不长翼，也不生羽毛，用不着经过这种艰难的步骤。
她的圆筒形的身体是裸露出来的，和毛虫没有多少区别。所以可以容许她在狭小的
隧道中爬出爬进，一点儿困难也没有。因此它把外衣抛弃在后面——抛在壳里面，作为
盖着茅草的屋顶。
同时，还有一种非常深谋远虑的举动，足以表现出她对于卵的命运有极其深切的关
心，事实上它们已经好像是被装在桶里面了，在它脱下的羊皮纸状的袋子里，母蛾已经
非常有办法地把卵产在里面了，直到把它装满为止，但是仅仅把它的房子与丝绒帽子遗
传给子孙，这并不能让它感到满足，最后的举动，还要把自己的皮也奉献出来留给子孙
后代，在它身上，“可怜天下父母心”这句话得到了最好的体现。
我想方便地观察这件事情的过程，于是我曾有一次从柴草的外壳里捡来一只装满卵
的蛹袋，并把它放在玻璃管中。在七月的第一个星期里，我忽然发现我竟然拥有了一个
被管虫的兴旺的大家族。它们孵化的速度是如此之快，差不多有四十多只以上的新生的
毛虫，竟在我没有看见的时候，在我还没来得及注意的时候，统统都穿上衣服了。
它们穿的衣服特别像波斯人戴的头巾，由光亮的白绒制做而成，说的普通一点，通
俗一些，就像一种白棉的礼帽，只是没有帽缨子。
不过说起来很奇怪，它们的这项帽子不是戴在头顶上的，而是从尾部一直披到前面
来的，它们在这玻璃管里非常得意地跑来跑去，因为这是属于它们自己的广大的屋子啊！
因此，我就想要看一看这顶帽子，究竟是由哪种材料做成的，织造的初步手续又是什么
样的。
幸运得很，蛹袋是不大会变空的。在里面，我又找到了它们第二个大家族，其数目
和先前跑出去的差不多。大概总有五打或六打的卵在里面。
我把那些已经穿好衣服的毛虫拿走，只留下这些裸露着身体的新客房在玻璃管里面，
它们有鲜红的头部，身体的其余部分全都是灰白色的，全身还不足一寸的二十五分之一
长。
我等待的时间并不长久，从第二天开始，这些小动物，慢慢地，成群结队地，开始
离开它们的蛹袋，用不着把这些摇篮弄破，只从它们母亲在当中弄破的口中出来就行了。
虽然它们都有洋葱头般的、漂亮的琥珀色，但是，没有一个把它拿来用做衣服的材
料，也没有一个利用那些柔软摇床的毛绒，谁都可能会以为这种材料可以做成这些伯冷
的动物的毛毯，但是事实上没有一个小动物去利用它一下。
它们一起冲到柴枝壳粗糙的外面，那是我故意为它们留下来的，而且直接靠近那个
装有卵的蛹袋，于是这些小动物们开始感觉到它们面临的情况有些不对头。于是便产生
了一种迫切感。
在你还未进入世界去打猎的时候，首先要做的是必须穿好自己的衣服，这一点对于
这些小动物们同样是适用的。它们也一样地焦急，恨不得马上攻破这个令人厌倦的陈旧
的老壳，赶紧穿上准备好的安全的外衣。
它们之中有的注意到了已经咬裂开的细枝，撕下那柔软的洁白的内层，有的很大胆，
深入到空茎的隧道，在黑暗中努力收集一些材料，它们的勇敢当然会有所报酬的，它们
得到了极其优等的材料，用这些织成雪白的衣服，还有一些毛虫加入了一些它们所选择
的东西，制做成了杂色的衣服，于是雪白的颜色给黑的微粒玷污了。
小毛虫制做衣服的工具就是它们的大头，其形状很像一把剪刀，并且它还长有五个
坚硬的利齿，这把剪刀的刀口靠得很紧凑。虽然它实际上很小，但它却很锋利，刀很快，
能夹住也能剪断各种纤维。
把它放在显微镜下可以清楚地观察到，小毛虫的这把剪刀竟然是有机械的、正确的，
而且是强有力的奇异标本。
如果羊也具备这样的工具的话，并且与它的身体成一定的比例，那么羊也就可以不
光吃草而也能吃树干了。由此可见，小毛虫的头可不能等闲视之啊！
观察这些被管虫的幼虫，制造棉花一样的灰白色的礼帽，这一点很能够启发人们的
智慧。无论是它们工作的行程，或是它们所应用的方法，都有很多的事情值得人们注意
一下。它们太微小了，它们也太纤弱了。当我用放大镜观察它时，必须非常小心，非常
仔细，既不敢使劲呼吸、喘粗气，也不敢大声说话，哪怕稍有一点不小心，就会惊扰了
它们，也许会把它们移动了位置，或者也可能是把它给一口气吹跑了。
别看这个小东西是如此微小，但是，它可是一位有着高超的制造毛毯技术的专家，
这个刚刚生下来一小会儿的小孤儿，竟然天生的知道怎样从它母亲留给它的旧衣服上裁
剪下自己的衣服来。它所采用的方法，我现在可以告诉人们，不过在此之前，我必须先
交代一点关于它的死去的母亲的事情。
我已经说过铺在蛹袋里的毛绒被，它很像一只鸭绒的床铺，软软乎乎，舒舒服服的，
小毛虫钻出卵以后，就睡在这张床上面休息一会儿，从而取得适当的温暖，并为到外面
的世界中去工作做好准备。
野鸭会脱下身上的绒毛，用它为子孙后代做成一张华丽舒适的床。母兔则会剪下身
上那些最柔软的毛，为它的新出生的儿女做成一张温暖的垫褥。母的被管虫也做着同样
的事情。看来，天下的母亲总还是有一定的共性的，这种共性也是它们的本能所决定的，
那就是无私地疼爱自己的儿女。
母亲会用一块柔软的充塞物，给小毛虫做成温暖的外衣，这材料非常的精细而且美
观。从显微镜下仔细地观察，可以看到上面有一点一点的鳞状片体，这就是它为小儿女
们制做衣服的最好的呢绒材料。小幼虫不久就会在壳里出现，因此要给它们准备好一个
温暖的屋子，让它们可以在里面自由地游戏玩耍。在它们还没有进入到广大的世界里去
之前，可以在里面修养，积蓄力量。所以母蛾像母兔、母鸭一样从身上取下毛来，为儿
女不辞辛劳地建造一片美好的天地。
这大概是以一种非常机械的方式进行的，好像是连续不断地磨擦墙壁而且并不是有
意识的有心的举动一样，然而的确没有理由向我们证实确是如此。甚至连最蠢笨的母亲
也有它自己的先见之明。这位看上去似乎有毛病的蛾子翻来覆去地打着滚，在狭窄的通
道中跑来跑去，想方设法地把自己身上的毛弄下来，给它的家族制做舒适的床铺。
有些书上说，小被管虫自从有了生命以后，就会吃掉它们的母亲。事实上，我却始
终也没有看到过这种情形发生。而且也不知道这个说法是怎样传说起来的。事实上，它
已经为它的家族奉献、牺牲了那么多，最后自己只留下干干的、薄薄的一个条，还不够
许多小子孙们的一口食物。实际上我的小被管虫们，它们是不吃母亲的。我看到的是它
们自从穿上衣服以后，一直到自己开始吃食的时候，没有一个曾经咬到自己的已死的母
亲的身上。

三、聪明的裁缝
现在我要详细地讲一讲这些小幼虫的衣服了。
卵的孵化是在七月初开始的，小幼虫的头部和身体的上部呈现出鲜明的黑色，下面
的两节，是带棕色的，其他部分都是灰灰的琥珀色。它们是一些十分精锐的小生物，跑
来跑去的脚步是很短小的，而且也是很快的。
它们从孵化地点的袋里钻出来以后，有一段时间，它们仍然需要待在从它们的母亲
身上取来的绒毛堆里。这里要比它们钻出来的那个袋子里更加空旷舒适一些。它们待在
绒毛堆里，有些在休息，有些十分忙乱，有的比较心急一些的已经开始练习行走了。它
们全体在离开外壳以前，都在修身养性，增强体质，以迎接未知世界风雨的洗礼。
在这个看上去比较奢华的地点，它们却并不留恋。等到它们的精力逐渐充沛起来，
就纷纷爬出来散布在壳上面。随后积极的工作就开始了，逐渐将自己穿着打扮起来。食
物问题以后才会想起来解决，目前却只有穿衣服是最要紧的事情，看来这些小家伙把脸
面上的事看得很重。
蒙坦穿上他父亲从前曾经穿过的衣服时，常常说：“我穿起我父亲的衣服了。”如
今，幼被管虫同样地穿起自己母亲的衣服(这同样必须记清，不是它身上的皮，而是它的
衣服)。它们从树枝的外壳，也就是我有时称做屋子，有时称做衣服的那种东西，剥取下
一些适当的材料，然后开始利用这些材料，给自己做衣服。它所用的材料都是小枝中的
木髓，特别是裂开的几枝，主要是因为它的髓是更容易取到的原故。
它们制做衣服的方法倒是非常值得注意的。这个小动物所采用的方法，真是出乎我
们人类的想象力，它是那样的灵巧，那样的细致，那样的精心，这种填塞物都被弄成极
其微小的圆球。那么这些小圆球是怎样连接在一起的呢？这位小裁缝需要一种支持物，
作为一个基础。而这个支持物又不能是从毛虫自己的身体上得来的。这个困难，并不能
难倒这些聪明的小家伙的；它们把小圆球聚集起来弄成一堆，然后依次用丝将它们一个
个绑起来。于是，困难就这样被克服了。你已经知道了，毛虫是能从自己身上吐出丝来
的，就像蜘蛛能吐丝织网一样。采用这种方法，把圆球或微粒连接在同一根丝上，做成
一种十分好看的花环，等到足够长了以后，这个花环就围绕在这个小动物的腰间，留出
六只脚，以便行动自由，末梢再用丝捆住，于是就形成了一根圈带，围绕在这个小幼虫
的身上。
这个圈带就是所有工作的起点和幼虫所需的支持物，完成第一道工序以后，小幼虫
再用大腮从壳上取下树心，固定上去，使它增长增大，于是就形成了一件完全的外衣。
这些碎树心或圆球，有时被放置在顶上，有时又被放在底下或旁边，不过通常都是放在
前边的时候居多。没有其他的设计，要比这个花环的做法更好了。外衣刚一做出来的时
候，是平的，后来把它扣住以后就像带子，圈在小毛虫的身体上。
最初工作的起点已经完成了，然后它会继续纺织下去。于是，那个最初的圈带逐渐
成为披肩、背心和短衫，后来成为长袍，几个小时以后，就完全变成一件雪白的崭新的
大衣了。
还要感谢它的母亲的关心，小幼虫得以免去光着身子到处跑来跑去的危险。假如它
不放弃那个旧的壳，那么，它们要想获得新的衣服将有很大的困难呀，因为草束和有心
髓的枝杆不是随处都可以找到的呢！然而，除非它们曝露而死，看来迟早它们总会找到
它们穿的衣服的，因为它们能利用随便什么材料，只要能找得到，什么都行。在玻璃管
中，我对于这些新生的小幼虫也曾做过好几回这样的试验。
从一种蒲公英的茎里，它毫不犹豫的挖出雪白的心髓，然后将它做成洁净的长袍子，
比它的母亲遗留给它的旧衣服所做成的要精致得多。有时还有更好的衣服，是用一种特
殊植物的心髓织造而成的。这一回的衣服上面饰有细点，像一粒粒的结晶块，或白糖的
颗粒。这可真正算是我们裁缝制作家的杰出作品了。
第二种材料，是我提供给它们的。那是一张吸墨纸，同样的，我的小幼虫也毫不犹
豫地割碎其表面，用它做成一件纸衣服，它们对这种新奇的材料非常高兴，也非常感兴
趣。当我再给它们提供那种原来的柴壳当做服装的材料时，它们竟然不予理睬，弃而不
顾，选取这种吸墨纸来继续做它们的衣服。
对于别的小幼虫，我什么东西也没有提供给它们，然而它们并没有就此失败。它们
非常聪明，采用了另一种方法，急急地去割碎那个瓶塞，使其成为小碎块，然后将这些
小碎块割成极其微小的颗粒，好像它们和它们的祖先也曾经利用过这种材料一样，因为
看上去这些小幼虫对这些材料并不陌生。这种稀奇的材料，也许毛虫们从来没有利用过，
然而它们把这些材料拿来做成衣服，竟然与其他材料做成的毫无差别。这些小幼虫的所
做所为真是让人感到惊奇！
从而我已经知道了它们能够接受干而轻的植物材料了，于是我决定换一种方法做试
验。用动物与矿物的材料来试试，我割下一片大孔雀蛾的翅膀，把两个裸体的小毛虫放
在上面。它们两个先是迟疑了好长时间。然后其中的一个就决心要利用这块奇怪的地毯，
一天的工夫都不到，它就穿起了它亲手用大孔雀蛾的鳞片做成的灰色的绒衣了。
第二回，我又拿来一些软的石块，其柔软的程度，只要轻轻一碰，就能破碎到如同
蝴蝶翼上的粉粒。在这种材料上，我放了四个需要衣服的毛虫。有一个很快就决定把自
己打扮起来，开始为自己缝制衣服。它的金属的衣服，有彩虹一样发出各种颜色的亮光，
闪烁在小毛虫的外壳上。这当然是很贵重，而且非常华丽的，只不过有点太笨重了。在
这样一个金属物的重压之下，小毛虫的行走变得非常辛苦，非常缓慢。不过，东罗马的
皇帝在国家有重大仪式的时候，也得如此呢！
为了满足本能上的迫切需要，幼小的毛虫也不顾忌这种蠢笨的行动了。穿衣服的需
要太迫切了，与其光着身子还不如纺织一些矿物好一些。爱美之心虫也有之，它也愿意
把自己打扮得漂漂亮亮的。吃的东西对于它并没有像穿的东西那样重要，只顾穿衣打扮，
外表好看，是这些小毛虫的共性与天性。假如先将它关起来两天，然后再换去它的衣服，
将它放在它喜欢吃的食物面前，比如一片山柳菊的叶子，它一定先做一件衣服，这是必
然的，因为一件衣服穿在身上后，它才会放心地去满足它的饮食需要。
它们对于衣服如此需要，并不是因为有特别寒冷的感觉，而是因为这种毛虫的先见。
别的毛虫在冬天，都是把自己隐藏在厚厚的树叶里，有的藏在地下的案穴里避寒，有的
在树枝的裂缝里，这是怕寒的毛虫。但是，我们所说的被管虫却安然地暴露在空气当中。
它不怕寒，也不怕冷，它从有生之日起，就学会了怎样预防冬季的寒冷。
受到秋天细雨的威胁以后，它又开始做外层的柴壳，开始时做得很草率、很不用心，
参差不齐的草茎和一片片的枯叶，混杂在一起，没有次序地缀在颈部后面的衬衣上，头
部必须仍然是柔软的，可让毛虫向任何方向自由转动。这些不整齐的第一批材料，并不
妨害建筑物后来的整齐。当这件长袍在前面增长起来的时候，那些材料便被甩到后边去
了。
经过一段时候以后，碎叶渐渐的加长，并且小毛虫也更细心地选择材料。各种材料
都被它直排的铺下去。它铺置草茎时的敏捷与精巧，真令人大吃一惊。人们不仅惊异地
发现小毛虫的动作如此之快，如此之轻巧，而且做的还很认真实在，铺垫的如此舒适，
这是一些大的昆虫都无法比拟的。真的不能小看它呀！
它将这些东西放在它的腮和脚之间，不停地搓卷，然后用下腮很紧的把它们含住，
在末端削去少许，立即贴在长袍的尾端。它的这种做法或许是要使丝线能粘的更坚固、
更结实些，和铅管工匠在铅管接合的尾梢锉去一点的意思是一个样的。
于是，在还没有放到背上以前，小毛虫用腮的力气，将草管竖起来，并且在空中舞
动它，吐丝口就立即开始工作，将它粘在适当的地方。于是，毛虫也不再摸索行动，也
不再移动，一切手续都已完成了。等到寒冷的气候来临的时候，保护自己的、温暖的外
壳已经作好了，所以，它可以安心地过自己的日子了。
不过这衣服内部的丝毡，并不很厚实，但能使它感到很舒服安逸。等到春天来临以
后，它可以利用闲暇的时间，加以改良，使它又厚又密，而且变得很柔软。就是我们拿
去它的外壳，它也不再重新制造了，它只管在衬衣上加上新层，甚至到不能再加为止。
这件长袍非常柔软，宽松而且多皱，又舒适、又美观。它既没有保护，也没有隐避之所，
然而它以为这并不要紧。做木工的时候已经一去不复返了，该是装饰室内的时候了，它
只一心一意地装饰它的室内，填充房子——即衬它的长袍，而房子已经没有了。它将要
凄惨地死去，被蚂蚁咬得粉碎，成为蚂蚁的一顿美餐。这就是本能过分顽固的结果呢！
------------------








樵叶蜂



如果你在园子里漫步，会发现丁香花或玫瑰花的叶子上，有一些精致的小洞，有圆
形的，也有椭圆形的，好像是被谁用巧妙的手法剪过了一般。有些叶子的洞实在太多了，
最后只剩下了叶脉。这是谁干的呢？它们又为什么要这样做呢？是因为好吃，还是好玩
呢？这些事都是樵叶蜂干的，它们用嘴巴作剪刀，靠眼睛和身体的转动，剪下了小叶片。
它们这么做，既不是觉得好吃，也不是为了好玩，而是这些剪下来的小叶片在它们的生
活中实在太重要了。它们把这许多小叶片凑成一个个针箍形的小袋，袋里可以储藏蜂蜜
和卵。每一个樵叶蜂的巢都有一打针箍形的袋，那些袋一个个地重叠在一起。
我们常看到的那种樵叶蜂是白色的，带着条纹。它常常寄居在蚯蚓的地道里，如果
你走到泥滩边，蹲下身子细细搜索，不难发现这样的地道。樵叶蜂并不利用地道的全部
作自己的居所，因为地道的深处往往又阴暗又潮湿，而且不适合排泄废物，有时还会遭
受昆虫的暗袭。所以它用靠近地面七八寸长的那段作自己的居所。
樵叶蜂一生中会碰到许多天敌，那地道毕竟不是一个安全坚固的防御工事，它用什
么办法来加强对自己家园的护卫呢？你瞧，那些剪下来的碎叶又派上大用场了。它用剪
下的许多零零碎碎的小叶片，把深处给堵塞了。这些用来堵塞的小叶片，都是樵叶蜂漫
不经心地从叶子上剪下来的。所以看上去非常零碎，不像筑巢用的那些小叶片一样整齐。
在樵叶蜂的防御工事之上有一叠小巢，大约有五六个。这些小巢由樵叶蜂所剪的小
叶片筑成，这些筑巢用的小叶片比那些做防卫工事的碎片，要求要高得多，它们必须是
大小相当、形状整齐的碎叶，圆形叶片用来作巢盖，椭圆形叶片用来做底和边缘。
樵叶蜂的小叶片，都是用它那把小刀——嘴角剪成的。为了适应巢的各部分的要求，
它往往用这把剪刀剪出大小不同的小叶片。对于巢的底部，它往往精心去设计，一点儿
也不含糊。如果一张较大的叶片不能完全吻合地道的截面的话，它会用两三张较小的椭
圆的叶片凑成一个巢底，一直到紧密地与地道截面吻合为止，决不留一点空隙。
做巢盖子的是一张正圆形的叶片。它好像是用圆规精确地规划过的，可以完美元缝
地盖在小巢上。
一连串的小巢做成后，樵叶蜂就着手剪许多大小不等的叶片，搓成一个栓塞把地道
塞好。
最值得我们思量的是，樵叶蜂没有任何可以用来当模子用的工具。它怎么能够剪下
这么多精确的叶子呢？它有可以依照的模型吗？还是它有什么特殊仪器可以测量呢？有
人推测，樵叶蜂的身体可以当作圆规来使用，一端固定住，即尾部固定在叶片某一点上；
另一端，也就是它们的头部，像圆规的脚一样在叶片上转动。这样就可以剪下一个标准
的圆。就像我们的手臂那样，固定肩肿挥动起来就是一个圆形。但是我们的手臂不会像
樵叶蜂那样巧妙又精确地画出大小一样的圆圈，这些用来做巢的盖子的圆叶片，恰好天
衣无缝地盖在巢上，非常完美。而小巢在地道的下面，它们不知道随时测量小巢的大小，
它们只靠摸索得到的感觉，来决定这只小巢所需要的叶盖大小。
圆形的叶片，不能剪得太大或太小。太大了盖不下，太小了会跌落在小巢内，使卵
活活闷死。你不用担心樵叶蜂的技术，它能很熟练地从叶子上剪下符合要求的叶片，虽
然没有什么模子，却是那么精确。樵叶蜂为什么有这么深厚的几何学基础呢？
冬天的一个晚上，我们都围着炉子坐着。我想起樵叶蜂剪叶片的事情，于是我设计
了一个小实验。
“明天是赶集的日子，你们中有人得出去采办这个星期所要用的东西。我们厨房里
一只天天要用的罐子的盖子被猫打破了。我要求他买一只盖子回来，不大不小，恰好能
盖上我们那个罐子。在去买之前，我们允许他仔细地把那罐子的大小估计一下，但不可
以用任何东西来测量，然后明天到集市上，凭记忆力选择合适的盖子。”大家听了都面
露难色，谁也不敢立即站出来承担这项任务。
的确，这似乎是一件很难的事。可是樵叶蜂的工作比我们这事更难以估测，它没有
看到自己的巢盖，根本没有这样一个印象；它也不能像我们选择盖子似的在摊贩的一大
堆罐中，靠着互相比较来选择一个最合适的盖子。对樵叶蜂来说，它必须在离家很远的
地方，毫不犹豫地剪下一片大圆叶，使它恰好能做巢的盖子。我们觉得很难的事，对它
来说像小孩游戏一样稀松平常。我们如果不用测量工具的话，比如绳子之类，或一个模
型或是一个图样，我们就很难选择一个大小适宜的盖子。可樵叶蜂什么都不需要。对于
如何治家，它们的确比我们聪明得多。
在实用几何学问题上，椎叶蜂的确胜过我们。当我看到樵叶蜂的巢和盖子，再观察
了其他昆虫在“科技”方面创造的奇迹——那些都不是我们的结构学所能解释的，我不
得不承认我们的科学还远不及它们。
------------------








采棉蜂和采脂蜂



我们知道，有许多蜜蜂像樵叶蜂一样自己不会筑巢，只会借居别的动物遗留或抛弃
的巢作自己的栖身之所。有的蜜蜂会借居泥匠蜂的故居，有的会借居于蚯蚓的地道中或
蜗牛的空壳里，有的会占据矿蜂曾经盘踞过的树枝，还有的会搬进掘地蜂曾经居住过的
砂坑。在这些借居它屋的蜜蜂中有一种采棉蜂，它的借居方式尤其奇特。它在芦枝上做
一个棉袋，这个棉袋便成了它的绝佳的睡袋；还有一种叫采脂蜂，它在蜗牛的空壳里塞
上树胶和树脂，经过一番装修，就可以当房间用了。
泥匠蜂很匆忙地用泥土筑成了“水泥巢”，就算大功告成了；木匠蜂在枯木上钻了
一个九英寸深的孔也开始心满意足地过日子了。尽管它们的家很粗糙，它们还是以采蜜
产卵为第一重要的大事，没有时间去精心装修它们的居室，屋子只要能够遮风挡雨就行
了；而另几类蜜蜂可算得上是装饰艺术大师，像樵叶蜂在蚯蚓的地道中做一串盖着叶片
的小巢，像采棉蜂在芦枝中做一个小小的精致的棉袋，使原来的地道和芦枝中别有一番
风情，令人不由得拍案叫绝。
看到那一个个洁白细致的小棉袋，我们可以知道采棉蜂是不适宜做掘土的工作的，
它们只能做这种装修工作。棉袋做得很长也很白，尤其是在没有灌入蜜糖的时候，看起
来像一件轻盈精致的艺术品。我想没有一个鸟巢可以像采棉蜂做的棉袋那样清洁、精巧
的。它是怎样把一个个棉花小球集中起来，拼成一个针箍形的袋子的呢？它也没有其他
特殊的工具，只有和泥匠蜂、樵叶蜂一样的灵巧的嘴，但它们的工作无论从方式上还是
从成果上看，都截然不同。
我们很难看清楚采棉蜂在芦枝内工作的情形，它们通常在毛蕊花、蓟花、鸢尾草上
采棉花，那些棉花早已没有水分了，所以将来不会出现难看的水痕。
它是这样工作的：它先停在植物的干枝上，用嘴巴撕去外表的皮，采到足够的棉花
后，用后足把棉花压到胸部，成为一个小球，等到小球有一粒豌豆那么大的时候，它再
把小球放到嘴里，衔着它飞走了。如果我们有耐心等待的话，将会看到它一次次地回到
同一棵植物上采棉，直到它的棉袋做完。
采棉蜂会把采到的棉花分成不同的等级，以适应袋中各个部分不同的需要。有一点
它们很像鸟类。鸟类为使自己的巢结实一些，会用硬硬的树枝卷成架子；又为了要使巢
温暖舒适些，而且宜于孵育小鸟，会用不同的羽毛填满巢的底部。采棉蜂也是这样做它
的巢，它用最细的棉絮衬在巢的内部，入口处用坚硬的树枝或叶片做“门”和“窗。”
我看不到采棉蜂在树枝上做巢的情形，但我却看到了它怎样做“塞子”，这个“塞
子”其实就是它的巢的“屋顶”。它用后足把棉花撕开并铺开，同时用嘴巴把棉花内的
硬块撕松，然后一层一层地叠起来，并用它的额头把它压结实。这是一种很粗的工作。
推想起来，它做别的部分的精细工作时，大概也是用这种办法。
有几只采棉蜂在做好屋顶后，怕不可靠，还要把树枝间的空隙填起来。它们利用了
所有能够得到的材料：小粒的沙土、一撮泥、几片木屑、一小块水泥，或是各种植物的
断枝碎屑。这巢的确是一个坚固的防御工事，任何敌人都无法攻进去。
采棉蜂藏在它巢内的蜂蜜是一种淡黄色的胶状颗粒，所以它们不会从棉袋里渗出来。
它的卵就产在这蜜上。不久，幼虫孵出来了。它们刚睁开眼睛，就发现食物早巳准备好
了，就把头钻进花蜜里，大口大口地吃着，吃得很香，也渐渐变得很肥。现在我们已经
可以不去照看它了。因为我们知道，不久它就会织起一个茧子，然后变成一只像它们母
亲那样的采棉蜂。
另外有一种蜜蜂，它们也是利用人家现成的房子，稍作改造变为自己的居住之处，
那就是采脂蜂。在矿石附近的石堆上，常常可以看到坐着吃各种硬壳果的蜗牛。它们吃
完后就跑了，石堆上留下一堆空壳。在这中间我们很可能找到几只塞着树脂的空壳，那
就是采脂蜂的巢了。竹蜂也利用蜗牛壳做巢，不过它们是用泥土做填充物的。
关于采脂蜂巢内的情形我们很难知道。因为它的巢总是做在蜗牛壳的螺旋的末端，
离壳口有很长的距离，从外面根本看不到里面的构造。我拿起一只壳照了照，看上去挺
透明的，也就是说这是只空壳，以后很可能被某个采脂蜂看中，在此安家落户，于是我
把它放回原处，让它作为将来的采脂蜂的巢。我又换一只照照，结果发现第二节是不透
明的，看来这里面一定有些东西。是什么呢？是下雨时冲进去的泥土？还是死了的蜗牛？
我不能确定。于是我在末端的壳上弄一个小洞，我看见了一层发亮的树脂，上面还嵌着
沙粒，一切都真相大白了：我得到的正是采脂蜂的巢。
采脂蜂往往在蜗牛壳中选择大小适宜的一节作它的巢。在大的壳中，它的巢就在壳
的末端。在小的壳中，它的巢就筑在靠进壳口的地方。它常常用细砂嵌在树胶上做成有
图案的薄膜。起初我也不知道这就是树胶。这是一种黄色半透明的东西，很脆，能溶解
在酒精中，燃烧的时候有烟，并且有一股强烈的树脂气味。你可以根据这些特点，判断
出来脂蜂用的是树干里流出来的树脂。
在用树脂和砂粒做成的盖子下，还有第二道防线，用砂粒、细枝等做的壁垒，这些
东西把壳的空隙都填得严严实实的。采棉蜂也有着类似的防御工程。不过，采脂蜂这种
工程只有在大的壳中才有，因为大的壳中空隙较多。在小的壳中，如果它的巢离入口处
不远，那它就用不着筑第二道防线了。
在第二道防线后面就是小房间了。在采脂蜂所选定的一节壳的末尾，共有两间小屋，
前屋较大，有一只雄蜂，后室较小，有一只雌蜂——采脂蜂的雄蜂比雌蜂要大。有一件
事科学家们至今仍无法解释，那就是母蜂怎能预先知道它所产的卵将是雌的还是雄的呢？
也就是它们怎么保证产在前屋的卵将来是只雄蜂，而产在后屋的卵一定会变成雌蜂呢？
有时候，采脂蜂筑巢的时候，一个小小的疏忽会造成下一代的一个大悲剧。让我们
来看看这只倒霉的采脂蜂吧！它选择了一只大的壳，把巢筑在壳的末端，但是从入口处
到巢的一段空间它忘记用废料来填充。前面我们提到过有一种竹蜂也是把巢筑在蜗牛壳
里的，它往往不知道这壳的底部已有了主人，一看到这个壳里还有一段空隙，就把巢筑
在这段空间里，并且用厚厚的泥土层把入口处封好。七月来了，悲剧就开始了。后面采
脂蜂巢里的蜂已经长大，它们咬破了胶膜，冲破了防线，想解放自己。可是，它们的通
路早已被一个陌生的家庭堵住了。它们试图通知那些邻居，让它们暂时让一让，可是无
论它们怎么闹，外屋那邻居始终没有动静。是不是它们故意装作听不见呢？不是的，竹
蜂的幼虫此时还正在孕育中，至少要到明年春天才能长成呢！难怪它们一直无动于衷。
采脂蜂无法冲破泥土的防线，一切都完了。它们只能让自己活活地饿死在洞里。这只能
怪那粗心的母亲，如果它们早能料到这一点，那么这悲剧也就不会发生了。如果那粗心
的母亲得知是自己活活杀死了孩子们，不知道该有多恨自己！不幸的遭遇并不能使采脂
蜂的后代学乖，事实上，时时有采脂蜂犯这样的错误，这与科学家所说的“动物不断地
从自己的错误和经验中学习和改进”的理论不符合。不过也难怪，你想，那些被关在壳
里的小蜂们永远地埋在了里面，没有一个能生还，这件事也随着小蜂们的死去而永远地
埋在了泥土里，成了无人能知的千古奇冤，更不用说让采脂蜂的后代吸取教训了。
------------------








西班牙犀头的自制



我希望你还记得神圣甲虫，它消耗掉它的时间，做成即可以当食物，又可以当梨形
窝巢的基础的圆球。
我已经指出，这种形状对于小甲虫的利处和害处，因为圆形是顶好的形状，可以保
存好食物使其不干也不硬。
经过长时间的观察这种甲虫的工作，我开始怀疑我极力赞扬它的本能，或许是我估
计错误了。它们是否真的关心它们的小幼虫，并且替它们预备下最柔软最合适的食物呢？
甲虫做球是它们自己的职业啊！它要继续在地底做球不是很奇怪吗？一个动物生着长而
弯的腿，用它把球在地上滚来滚去是很便利的。无论在那里，自然要从事自己所喜欢的
职业。自己想干的工作，就一定要干好，只有这样才能在自然界中求生存，才能在大自
然中繁衍后代，一代一代地生存下去。
它并不顾及它自己的幼虫，或许它做成梨形的外壳这件事仅仅是碰巧了而已。
为了要圆满地解决这个疑难问题，我还观察过一种清道的甲虫，在它的日常工作中，
它非常不熟悉做球这种工作。可是，到了产卵期，它突然改变了以往的习惯，将自己储
存的所有食物都统统做成圆圆的一个团。这一点表明这不仅仅是习惯而已，而是真的关
心它的幼虫，因而选择圆形的球做为它的窠巢。
如今，在我的住所附近，就有这样一种甲虫。它是甲虫中最漂亮的，个子最大的。
虽然不如神圣甲虫那么魁伟，它的名字就是——西班牙犀头。
它最显著、最特别的地方，就是它胸部的陡坡和头上长的角。
这种甲虫是圆的，而且很短，当然也就不适合做神圣甲虫所做的那些运动。它的腿
不足以供做球使用。稍有一点点惊扰，它的腿就本能地卷缩在自己身体的下面，它不像
一个勇敢者，也不像神圣甲虫那样，有一个勇敢者的气魄。
它们一点也不像搓滚弹丸的工具，它们那种发育不全的形象，表明它们缺乏挖掘性，
这足以使我们清楚它是不能带着一个滚动的圆球走路的。
的确，犀头的性格很不活泼。有一次，在夜里，或在黄昏的月光下，它寻找到食物，
就在原来的地点挖开一个洞穴。它的这种挖掘草率的很，其最大的也只能藏下一个苹果。
在这里，它逐渐堆下刚刚才找来的食品和食料，至少一直要堆积到洞穴的门口。
它的大量的食物要堆积为不成形的一大堆，这就足以证明这个犀头的贪食、贪吃和
馋嘴了。食物能够吃多长的时间，它自身也就在这地底下待多长时间，一直待到吃完所
存的食物为止。
等它把所有存储的食物全都吃完以后，它的食品仓库空了，它这才又重新跑出来，
再去寻找新鲜的食物，然后再另挖掘一个洞穴，重复它那种存了吃，吃了再出来找的周
期性运动。
实实在在地说，它只不过是一个清道夫，是一个肥料的收集者而已。总之，它没有
什么特别的本事，是一个平庸之虫。
对于搓捏圆球的技术，它明显表现出特别的外行。而且，它的短而笨的腿，也极其
的不适合干这种技术性的工作。
在五六月之间，产卵的时候到了，这个昆虫则变成了非常擅长于选择最柔软的材料，
选择最舒适的环境，为它顺利产卵打下一个良好环境的能手了。
它开始为它的家族制做食物，只要在一个地方找到，如果它认为是最好的，它立刻
就把它们埋在地下，它从不旅行，从不搬运，从不做任何添加配制工作，也从不进行再
加工。
然而，我看到这个洞穴，比它自己吃食的临时的洞穴，挖掘得更宽大一些，而且建
筑得也比较精细。
我觉得在这种野外的环境里，要想仔细观察犀头的一些生活习惯，以及它的生长过
程，是非常不容易的，所以后来我就将它放到我的昆虫屋里面，这样，我可以更加认真、
更加仔细地观察。这为我自己提供了许多的方便。
起初，这个可怜的昆虫，因为被我俘虏了，所以有一些胆怯，它可能认为大难即将
来临。当它做好了洞穴以后，自己出入洞穴时，也还是提心吊胆，唯恐自己被再次伤害
了一样。然而从这以后，它也就逐渐的胆壮起来，在一夜之间，将我提供给它的食物全
部储存起来了。
在一个星期快要过去的时候，我掘起昆虫屋中的泥土。我发现，我见过的它储存食
物的洞穴显现出来了，这是一个很大的厅堂，一个很大的仓库。它的屋顶并不很整齐，
四壁也是很普通的，地板差不多是平平坦坦的。
在一个角上，找个圆孔，从这里一直通往倾斜的走廊，这个走廊一直通到土面上。
这个房子——这个昆虫的别墅——用新鲜的泥土掘成的一个大洞。它的墙壁，曾经被很
仔细的压过，很认真地装饰过。这也就足以抵抗我在做试验时所引起的地震了。并且很
容易就能看到这个昆虫以及它所有的技能，它不遗余力，用尽所有的掘地力量，来做一
个永久的家。可是它的餐室却仅仅是一个土穴，墙壁做的也不那么坚固。
当它从事这个大型建筑的建设的时候，我想，它的丈夫，或者是它的伴侣一定会来
帮助它的，至少我常常看见它和它的丈夫一同待在一个洞穴里。我也相信这个帮助会使
它的妻子更加勤快，丈夫和妻子可以一起收集并储存食物。因为夫妻二人同做一件事情，
同干一件工作，自然要快得多，至少比一个人干事要快得多，但是等到屋子里储备满了，
足够它生活以后，它的丈夫也就隐退了。这位丈夫就跑回到土面上来，到别的地方去安
身了。它对这个家庭应做的工作，应尽的职责也就全部结束了，尽到了一个丈夫应尽的
职责，就此结束了对这个家庭的义务。
那么，在许多食物放下去的土屋中，我所看到的是什么样的呢？是一大堆小土块，
互相堆叠在一起吗？但是，一点儿都不对。实际上不是想象中的样子。我只看到单独的
一个很大的土块，除掉一条小路以外，储存食物的那一个屋子，全都被塞满了。
这种圆堆块没有一定的形状，有的大小像吐绶鸡的蛋，有的像普通的洋葱头。有的
是差不多是完整的圆形。这使我想起了荷兰的那种圆形硬酪。有的是圆形而上部微微有
点突起。然而，无论是哪一种，其表面都是很光滑的，呈现出精致的曲线。
这位母亲，不辞辛苦地一次一次地带去很多很多的材料，收集在一起并搓成一个大
团。它的做法是，捣碎这许多的小堆，将它们合在一起，并把它们揉合起来，同时也踩
踏它们。有好几回我都曾经见到它在这个巨大的球顶上。当然，这个球要比神圣的甲虫
做的那个大得多，两个互相比较一下，后者只不过是个小小的弹丸而已。它也有时在约
四寸直径的凸面上徘徊，它敲它、拍它、打它、揉它、含它，使它变得坚固而且平坦。
我只有一次见过如此新奇的景观，而且只有一次。这是多么难得的机会啊！但是当它一
见到我的时候，立刻就滚到弯曲的斜坡下不见了。它发现，它的所做所为已被人注意到
了，完全暴露身份和目标，所以它就逃之大吉了。
我得力于一排墨纸盖住的玻璃瓶，在这里我发现了许许多多有趣的事情。
第一我发现了这个大球的雕饰过程——常常是很整齐的，无论其倾斜程度的差异如
何——这并不是由于搓滚的方法而形成的。
事实上我已经知道，这么大的体积决不能滚进这个差不多已经被塞满了的洞里去。
而且这个昆虫的力量也不足以移动这么大堆的东西。
我每次到瓶边观察时，所得到的证据都是一个样子的，我常常看到母虫爬到球顶上，
看看这里、又看看那里，看看这边、又看看那边，它轻轻地敲，轻轻地拍，尽量使之光
滑，似乎没有见过它有想移动这个球的意思。
事实明确地证实，制球是并不采用搓滚的方法的。
最后已经准备好了。这就像面包工人将面粉团分成许多许多的小块，每一块将来都
将成为面包。这犀头甲虫也是一样的做法。它用头部锋利的边缘，及前爪的利齿，划开
圆形的裂口，从大块上随意割下小小的一块来。在做这次工作的时候，一点犹豫都没有，
也不重复改做一下。它从不在这里加上一点，或者在那里去掉一点。直接了当，只要一
次切割，它就得到适当的一块了。
其次，就是如何使球有一定的形状。它竭力将球抱在那双短臂之间，叫人看起来它
很不适于做这项工作的，只用压力把它做成圆块。它很庄严、很正重地，在不成形的一
块食物上爬上爬下，向左爬，向右爬，向前爬，向后爬，不停地爬，耐心地一再触摸，
最后经过二十四小时以上的工作，终于有棱有角的东西变圆了，像成熟的梅子一样大小。
在它狭小的技术操作室里，简直就没有什么余地可以自由地转动一下了。这位又矮
又胖的艺术家完成此项工作，竞然会没有动摇它的底面。但是经过相当的时间与耐力以
后，它竟然做成了确实适当的圆球。从如此笨拙的工具与有限的地位而论，这看来似乎
是不可能的。
它亲切有味地用足摩擦圆球的表面，再经过很长的时间，最后它终于满意了。然后，
它爬到圆顶上面，慢慢地压，压出一个浅浅的穴来，就在这个盆样的孔穴里它产下一个
卵。
于是，它非常当心，非常精细地把这个盒子的边缘合拢起来，以遮盖它产下的那个
卵，再把边缘挤向顶上，使之略略尖细而突出。最后，这个球就做成椭圆形的了。
这个昆虫于是又开始从事第二个小块的工作，制造的方法完全相同。余下的，又重
新做第三个乃至第四个，你当然记得，神圣甲虫用很熟悉的方法只做一个梨形的巢。
它的洞穴中隐藏着三四个蛋形的球，一个紧靠着一个，而且组合、排列都很有规则，
细小的一端全都朝着上面。
它经过长期的工作以后，谁都要以为它也像神圣的甲虫一样，跑出来寻找自己的食
物去了。然而事实却不是这样的，它没有那样做。它没有跑出去，也没有去寻找食物，
而是在那里一动不动地守着，并且自打它钻入地下以后，它一点食物也没有吃过，它像
宇宙间任何母亲一样，一样的无私，一样的奉献。对自己的子女只有爱护、关怀与牺牲。
它没有，而且也不肯去碰一碰那为自己的子女预备下的食物。它宁愿自己挨饿，宁
愿自己受痛苦，也不愿意自己的小幼虫将来感受到一点儿痛苦，这是多么了不起的奉献
精神啊：在昆虫的世界里也充分体现了母爱是最伟大的。
它不出去的目的，当然是为了看守这几个为子女建筑下的摇篮。因为这是这个家族
生存的基本条件之一。这是它们的房子，是它们的小别墅，是它们生活在世上的惟一栖
身的地方。因而要仔细地看护它。
神圣的甲虫的梨正是因为母亲的离开，而遭到损坏的，当母亲离开不久，梨就已破
裂开了。经过一个相当长的时期以后，就不成形状了，就这样，一个家被毁掉了。
但是这个甲虫的蛋，可以保存完好，并长时间地保存，因为它有母亲的关心爱护，
母亲的一份责任感，才使它们的蛋完好地保存下来。
它从这一个跑到那一个上，再从那一个跑到另一个上，看看它们，听听它们，唯恐
它们有什么闪失，受到了什么外来的侵害。就像人类母亲对自己怀里的婴儿一样，关怀
得无微不至。这小甲虫真是一个好母亲。
它修补这一处，然后又修补那一处，生怕它的小幼虫受到什么干扰，受到外来的欺
辱。我们的眼睛看不出什么不足的地方，它虽然很笨拙而且有角，有足，但是在黑暗中
竟然比我们的视觉在日光中还要灵敏，还要看的清楚，这一点我们可以感觉的到。只要
有细微的破裂，它立刻就会跑过去，赶紧地修补一下，惟恐空气会透进去，干掉它的卵。
它在摇篮当中狭窄的过道里跑出跑进，为的是保护它的卵，它仔细观察，认真巡视，
假如我们打扰它，破坏它正常的生活，它就立刻用体尖抵住翼尖壳的边缘，做出柔软的
沙沙之声，如同和平的鸣声，又像发出强烈的抗议一般。
它就是这样，辛辛苦苦地关心着它的摇篮。有时候它实在困了，也会在旁边睡上一
小会儿，但时间不会太长的，只是打一会儿盹而已，决不会高枕无忧的睡上一大觉。这
位母亲就是这样在看守它的卵，为它的后代做出无私的奉献，为儿女操碎了一颗心。
犀头在地下室中，有着一个昆虫所稀有的特点，那就是照顾自己家庭的快乐。这是
多么伟大的母爱呀！这是一个奉献者的自豪。
它在自己弄下的缺口处。听见它的幼虫在壳内爬动，争取自由。当这个小囚犯，伸
直了腿，弯曲了腰，想推开压在自己头上的天花板时，它的母亲会意识到，小幼虫一天
天长大了，要独立生活了，该自己去世界上闯荡一番了。这位小幼虫自己出来，感受自
由与生命的美好。
即然有建造修理的本领，为什么不能打碎它呢？然而我不能做出肯定的回答，因为
我没有见到过这种事情发生。或许可以说这个母虫，被关在无法逃脱的玻璃瓶子里，所
以它一直守在巢中，因为它没有任何行动的自由。不过，假使如此，它对摩擦工作与长
时间的观察难道不感到焦急吗？这个工作显然对于它很自然，形成了它已经习惯了的一
部分生活了。
假如它急切的想恢复自由，它当然要在瓶中爬上爬下，毫无休止地忙碌。但是，我
只看见它常常是很平静的，也很安心地待在它的圆球旁。
为了要得到确切的第一手资料，为了得到确切的真象，所以我随时去察看玻璃瓶中
的现象如何。
如果它要休息，它可以任意的钻入沙土中，到处都可以隐藏它的身体，如果需要饮
食，也可以出来取得新鲜食物，然而既不是休息，也不是日光与饮食，可以便它离开它
自己的家族片刻。它只坐镇在那里，直到最后一个圆球破裂开，我常见它总是坐在摇篮
旁边的，那分安静，那分重担在肩的责任感很让我感动。
大概有四个月的时间，它不吃任何食物，它已不像最初为了照顾家族时那么贪嘴了，
而这时它竟然对于长时间的坐守，有非常惊人的自制力了。
母鸡伏在它的蛋上，忘记饮食数星期以后，自己的蛋才变成小鸡，然而犀头却要忘
记饮食达到一年的三分之一那么久。
夏天过去了。人类和牲畜都很希望下几场雨，终于下来了，地上积了很深的水。
于是，在我们布罗温司酷热干燥、生命不安的夏季过后，我们有凉爽的气候来使它
复活了。
石南开放了它的红色钟形的花，海葱绽放穗状的花朵，草莓树的珊瑚色果子也已经
开始变软了，神圣甲虫和犀头也裂开外层的包壳，跑到地面上来，享受一下一年来最后
这几天的好天气了。
刚刚解放出来的犀头家族，与它们的母亲一起，逐渐地来到地面。大概有三四个，
最多的是五个。
公的犀头生有比较长的角，很容易就能分辨出来。
母的犀头与母亲则很难分别。因此它们之间，很容易混淆。
不久，又有一种突然的改变发生了。从前牺牲一切的母亲，现在对于家族的利益，
已不再那么关心了。
自此它们各自开始管理自己的家和自己的利益了。它们彼此之间也就不相互照应了。
目前虽然母甲虫对家族漠不关心，但我们都不能因此而忘记它四个月来辛辛苦苦的
看护，除掉蜜蜂、黄蜂、蚂蚁等外来的干涉和侵犯。自己能养儿育女，关心它们的健康，
直到长成之后，据我所知，再没有别的昆虫能够做到这些了。
它独自一个毫无外来帮助，为每个孩子预备摇篮似的食物，并且尽心修补，以防止
其破裂，使摇篮十分安全。这是一个母亲无私的奉献。
它的情感如此的浓厚与执着，使它失掉了一切的欲望和饮食的需要。
在洞穴的黑暗里看护它的骨肉达到四个月之久。细心地看护着它的卵。
它在的子女们未得到解放出来之前，它决不恢复户外的快乐生活。
我们竟从田野中愚蠢的清道夫身上，看到最深切的关于母性本能的例子，不禁对这
种小昆虫产生了无限的敬意。
------------------








两种稀奇的蚱蜢



一、恩布沙
海是生物最初出现的地方，至今还存在许多种奇形怪状的动物，让人们无法统计出
它们的具体数目，也分不清它们的具体种类。这些动物界原始的模型，保存在海洋的深
处。这就是我们常说的，海洋是人类无价的宝库，它是人类生存的重要条件之一。
但是，在陆地上，从前的奇形动物，差不多都已经灭绝了，只有少数的还遗留下来，
能留到现在的大多都是一些昆虫类的动物。其中之一就是那种祈祷的螳螂，关于它特有
的形状和习性，我已经在前文对你们说过了。别一种则是恩布沙。
这种昆虫，在它的幼虫时代，大概要算布罗温司省内最怪的动物了。它是一种细长，
摇摆不定的奇形的昆虫。它的形状和任何昆虫都不一样，没有看惯的人，决不敢用手指
去碰触它。我的近邻的小孩，看了这个奇怪的昆虫以后，看到它这个奇异的模样，留下
了很深的印象，他们叫它为“小鬼”。他们想象它和妖法魔鬼等等多少有些关系。从春
季到五月，或是到秋天，有时在有阳光和温暖的冬天，可以遇见它们，虽然从不集成大
群。
荒地上坚韧的草丛，可以受到日光照耀，并且有石头可以遮风的矮丛树，都是畏寒
的恩布沙最喜欢的住宅。
我要尽我一切的可能告诉你们，它看起来像什么样子。它身体的尾部常常向背上卷
起，曲向背上，形成一个钩的形状，身体的下面，即钩的上面，铺垫着许多叶状的鳞片，
并排列成三行。
这个钩架在四只长而细的，形如高跷的腿上；每只足的大腿和小腿连接之处，有一
个弯的、突出的刀片，这个刀片与屠夫切肉常用的那种刀片相仿。
在高跷或四足蹬上的钩的前面，有很长而且很直的胸部突起。形状圆而且很细，像
一根草一样，草干的末梢，有猎狩的工具，是完全类似螳螂的那种猎具。
这里有比较尖利的鱼叉，还有一个残酷的老虎钳，生长着如锯子似的牙齿。上臂做
成的钳口中间有一道沟，两边各有五只长长的钉，当中也有小锯齿。臂做成的钳口也有
同样的沟，但是锯齿比较细巧，比较密一些，而且很整齐。
在它休息的时候，前臂的锯齿嵌在上臂的沟里。它的整体就像一架可以加工的机器，
有锯齿、有老虎钳、有沟、有道，如果这部机器再稍微大一点，那它就成了一部令人可
畏的刑具了。
它的头部也和这种机器相辅相承。这是一个多么怪异的头啊！尖形的面孔，卷曲而
长的胡须，巨大而且突出的眼睛，在它们中间还有短剑的锋口；在前额，有一种从未见
过的东西——一种高的僧帽一样的东西，一种向前突出的精美的头饰，向左向右分开，
形成尖起的翅膀。
为什么这个“小鬼”要这样像古代占卜家一样戴着奇形怪状的尖帽子呢？它的用途
在不久以后我们就会知道的。
在这个时候，这动物的颜色是普通的，大抵为灰色，待发育以后，就会变为装饰着
灰绿、白、与粉红的条纹。
如果你在丛林中遇见这个奇怪的东西，它在四只长足上动荡，头部向着你不停地摇
摆，转动它的僧帽，凝视着你的眉头。
在它的尖脸上，你似乎可以看到要遭受危险的形象。但是，如果你想要捉到它，这
种恐吓姿势，马上就会不见了。
它高举的胸部就会低下去，竭力用大步逃之夭夭，并且它的武器会帮助它握着小树
枝。假如你有比较熟练的眼光，它就很容易被捉住，关在铁丝笼子里。
起初，我不知道应该如何喂养它们。我的“小鬼”又很小，最多只有一两个月大。
我捉大小适宜的蝗虫给它们吃，我选取了其中最小的一些喂给它吃。
“小鬼”不但不要它们，而且还惧怕它们，无论那个无思想的蝗虫怎样很温和地靠
近它，都会受到很坏的待遇。
尖帽子低下来，愤怒的一捅，使蝗虫滚跌开去。
因此可知，这个魔术家的帽子实际上是自卫的武器。雄羊用它的前额来冲撞，和它
的对手进行搏斗，同样的，思布沙也在用它的僧帽来和它的对手进行抵抗。
第二次，我喂给它一个活的苍蝇，这个恩布沙立即就接受了它，把它当成一次酒席
上的佳肴。当苍蝇走近它的时候，早己守候着的恩布沙掉转了它的头，弯曲了胸部，给
苍蝇猛然一叉，把它夹在两条锯子之间。就连老猫扑捉老鼠也没有这样的迅速。
我惊奇地发现，一只苍蝇不仅可供给它一餐，而且足够整日食用，甚至可以连着吃
上几天。这种相貌凶恶的昆虫，竞然是极其节食的动物。
我开始以为它们是一个个的魔鬼，但是，后来发现它们的食量像病人一样少。经过
一个时期以后就连小蝇也不能引诱它们了。在冬天的几个月里，它完全是断食的。到了
春天，才又准备吃一些小量的米蝶和蝗虫。它们总在颈部攻击俘虏，如螳螂一般。
幼小的恩布沙，被关在笼子里时，有一种非常特殊的习性。
在铁丝笼里，它的态度从最初一直到最后，都是一样的，而且是一种顶奇怪的态度。
它用它那四只后足的爪，紧握着铁丝倒悬着，纹丝不动，活像一只倒挂在横杠上的小金
丝猴一样，它的背部向下，整个的身体就挂在那四个点上。如果它想移动一下，前面的
鱼叉就会张开，向外伸展开去，然后，紧握住另一根铁丝，朝怀里拉过来。
用这种方法将这个昆虫在铁丝上拽动，仍然是背朝下的，于是鱼叉两口合拢，缩回
来放在胸前。
这种倒悬的位置，对于我们而言一定会很难受的，也是很不容易做到的，要是人很
可能就会得病的，要么是高血压，要么是脑出血。但是，恩布沙保持这样的姿势的时间
并不短，它在铁丝笼里，可以持续十个月以上，竟然毫无改变。
苍蝇在天花板上，确实也是这种姿势的，但是它有休息的时间，它累了就要休息一
会，养足了精神以后，再做这种动作。它在空中飞动，用平常的习惯走路，沐浴在阳光
中。
恩布沙则完全相反，它保持这种奇怪的姿势，达到十个月以上，绝不休息。它悬挂
在铁丝网上，背部朝下，猎取、吃食、消化、睡眠、经过昆虫生活所有的经历，直至最
后死亡。它爬上去时年纪还很轻，而落下来的时候，已经是年老的尸首了。
它这个习惯的动作，应该注意的是只有处在俘囚期的时候才会如此，并不是这种昆
虫天生的、固有的习惯。因为在户外，除去很少的时候，它站在草上时是背脊向上的，
并不是倒悬着的。
和这种行为相似的，我还知道另外一个稀奇的例子，比起这个还要特别一些。这就
是一种黄蜂和蜜蜂，在夜晚休息时的姿态。有一种特别的黄蜂——生有红色的前脚的
“泥蜂”八月底的时候在我的花园里非常之多，它们很喜欢在薄荷草上睡眠。在傍晚薄
暮时，特别是在窒闷的日子里，暴风雨正在酝酿，大风大雨即将来临的时候，可是，我
们却能见到一个奇怪的睡眠者——仍然在那里安祥地熟睡着。
大概在晚上休息时，它的睡眠姿态没有比这个更奇怪的了。当你见到它以后就会觉
得特别的稀奇古怪了。它用颚咬入薄荷草的茎内，方的茎比圆的茎更能握的牢固一些，
它只用嘴咬住它，身体却笔直地横在空中，腿折叠着，它和树干成直角，这昆虫把全身
的重量，完完全全的放置在它的大腮上。
泥蜂利用它强有力的颚这样睡觉，身体伸展在空中。如果按动物的这种情形来推测，
我们从前对于休息的固有观念就要被推翻了。
任凭风暴狂欢，树枝摇摆，这位睡眠者并不被这摇晃的吊床所烦扰，至多是在某个
时候用前足抵住这摇动的枝干罢了。也许黄蜂的颚像鸟类的足趾一般，具有极强的把握
力，比风的力量还要强大许多。
尽管如此，有好几种黄蜂和蜜蜂都是采用这种奇怪的姿势来睡眠的——用大腮咬住
枝干，身体伸直，腿缩着。
大约在五月中旬，那时候恩布沙已经发育完整了。它的体态和服饰比螳螂更引入注
目。它还保留着一点幼稚时代的怪相——垂直的胸部，膝上的武器和它身体下面的三行
鳞片。但是它现在已经不能卷成钩子，它现在看起来也文雅多了：大型灰绿色的翅膀，
粉红色的肩头，矫捷的飞翔，下面的身体装饰着白色和绿色的条纹。
雄的恩布沙，是一个花花公子，和有些蛾类相似，更是夸张地用羽毛状的触须修饰
着自己。
在春天，农人们遇见恩布沙的时候，他们总以为是看到了螳螂——这个秋天的女儿
了。
它们外表很相像，以致人们都怀疑它们的习性也是一样的。因为外观一样，又都是
昆虫类的动物，所以人们没有认真仔细观察，也没有考察过它们的行动坐卧，所以就猜
测它们的生活习惯是一样的。
但是，事实上因为它的那种异常的甲胄，会使人们想到恩布沙的生活方式甚至比螳
螂要凶狠得多。但是，这种想法却错了，这个误解对恩布沙是不公平的，无调查研究的
结论是靠不住的。
尽管它们都具有一种作战的姿态，但是，恩布沙却是一个比较和平友好的动物呢！
它不是一个好斗好战的恶劣的凶手。
把它们关在铁丝罩里，无论是半打(一打是十二只，半打是六只)或者只有一对，它
们没有一刻忘掉柔和的态度。它们之间都是和平友好，互利相处的。
甚至到发育完成的时候，它们几个也是互相体谅，互相谦让，互不侵犯的。它们吃
的东西比较少，每天的食物只有两三只苍蝇就足够了。
食量大的小动物，当然是好争斗的。吃得饱的动物，把争斗当做一种消化食物的手
段，同时也是一种健身的方式。争强好胜，事事不让人，从来不吃亏，这是典型的弱肉
强食者的特点，它从来就是见便宜就占，见利益就争，见好事就抢。螳螂一见到蝗虫立
刻就会兴奋起来，于是战争就不可避免地开始了。螳螂立刻就扑向蝗虫，但是蝗虫也不
示弱，两者你争我斗，蝗虫用利齿欲扑向螳螂，但螳螂用它尖利的双夹给蝗虫以有力的
反扑。你争我斗的场面，十分精彩。
但是，节食的恩布沙，是个和平的使者，它从不和邻居们争斗，也从不用做鬼的形
状，去恐吓外来者。它也从不像螳螂那样，和邻居们争夺地盘。它从不突然张开翅膀，
也不像毒蛇那样作喷气、吐舌状。它从来也不吃掉自己的兄弟姐妹。更不像螳螂那样，
吞食自己的丈夫。这种惨无人道的事情，它是从来不做的。
这两种昆虫的器官，是完全一样的。所以这种性格上的不同，与身体的形状无关，
与其外表也无关。或许可以说是由于食物的差异而造成的。
无论是人还是动物，淳朴的生活总可以使性格变得温和一些，随和一些。这些都可
以营造一个和平共处的好环境。但是，自奉太厚了，就要开始残忍起来。贪食者吃肉又
饮酒——这是野性勃发的普遍原因——从不能像自制的隐士一样温和平静。它是吃些面
包，在牛奶里浸浸，这样简单的生活。它是一个普普通通的昆虫，它是平和、温柔、和
善的。而螳螂则是十足的贪食者。
虽然我的解释已经很清楚明白了，但是还有人可能会提出更深一层的问题。
这两种昆虫有完全相同的形状，想来一定也有同样的生活需要，而为什么，一种如
此的贪食，而另一种又如此有节制呢？它们在态度方面，如同别的昆虫已经告诉我们的
一样，嗜好和习性并不完全取决于自身的形状，以及身体结构，而是在决定物质的定律
方面，还有决定本能的定律存在。

二、白面孔螽斯
在我所居住的区域里的螽斯是白面孔的。无论在其善于歌唱，还是在其庄严的丰彩
上，它都可以算得上是蚱蜢类中的首领。它生有灰色的身体，一对强有力的大腮，以及
宽阔的象牙色的面孔。
如果要想捕捉它，这并不是什么难做到的事，也并不烦人。在夏天最炎热的时候，
我们常可以见到它在长长的草上来回跳跃。特别是在岩石下面，那里有松树生长着。
希腊字Dectikog(即白面孔螽斯、Decticns的语源)的意义是咬，喜欢咬。白面孔螽
斯因此取了这个名字。
它确实是善于咬的昆虫。假如有一种强壮的蚱蜢抓住了你的指头，你可是要当心一
点儿，它会把你的指头咬出血来，咬得你生疼，甚至有时疼痛难忍。它那强有力的颚仿
佛是凶猛的武器。当我要捕捉它时，我必须非常小心提防它，否则随时都有被它咬伤的
危险和被它咬破的可能。它那两颊突出的大型肌肉，显然是用来切碎它捕捉的、硬皮的
捕获物时用的。
把白面孔螽斯关在我的笼子里，我发现蝗虫蚱蜢等任何新鲜的肉食，都符合它们的
需要。特别是那种长着兰色翅膀的蝗虫，尤其适合它的嗜好。
当把食物放进笼子里时，常常会引起一阵骚动。特别是在它们饿极了的时候，它们
一步一步的很笨重地向前突进。因为受长颈的阻碍，它不能很敏捷地行动。有些蝗虫立
刻就被捉住，有的乱飞，乱蹦，乱跳，有的急了跳到笼子的顶上，逃出这螽斯所能捉捕
到的范围之外。因为它的身体很笨重，不能爬得那么高。不过蝗虫也只能是延长它们自
己的生命而已，最终也无法逃脱被白面螽斯蚕食的恶运。它们或因疲倦、或因被下面的
绿色食物所引诱，纷纷从上面跑下来，于是立刻就会被螽斯所捕获，成为其口中之美食。
这种螽斯，虽然智力很低下，然而却会用一种科学的杀戳方法。如同我们在别的地
方见到的一样，它常常先刺捕猎物的颈部，然后再咬住主宰它运动的神经，使它立刻失
去抵抗的能力。和其他肉食动物一样如哺乳动物虎、猎豹等等，它们都是先将所捕捉的
猎物的喉头管咬住，使其停止呼吸，丧失反抗力后，再一点点地享用它的肉体。
这是一种很聪明的方法，因为蝗虫是很难杀死的。有时虽然蝗虫的头已经掉下来了，
但它的躯体依然还能够跳动不已。我曾经见过几只蝗虫，已经被吃掉一半了，还不断地
乱跳，居然被它逃走了。
因它嗜好蝗虫，以及有些对于未成熟的谷类有害的种族，所以这类螽斯多一些，对
于农业也许有相当的益处。
不过现在它对于土地上保存果实的帮助，是非常薄弱的。它带给我们的主要的兴趣，
事实上是那些远古遗留下来的纪念物。它留给我们一些现今已经不用了的习性。
我应该谢谢白面孔螽斯，使我再次知道了关于幼小螽斯的一两件事情。
它产下的卵，并不和蝗虫，螳螂一样，把它们装在硬沫做成的桶里，它也不像蝉那
样，将它们产在树枝的洞穴里。
这种螽斯将卵像植物种子一般，种植在土壤里。母的白面螽斯身体的尾部有一种器
官，可以帮助它在土面上掘下一个小小的洞穴。在这个洞穴内，产下若干个卵，将洞穴
四周的土弄松一些，用这种器具，将土推入洞中，就像我们用手杖将土填入洞穴一样。
用这样一种方法，它将这个小土井盖好，再将上面的土弄平整。
然后，它到附近的地方散一会儿步，以作消遣和放松。用不了多长时间，它就会回
到先前产卵的那个地方，靠近原来的地点——这是它记得很清楚的——又重新开始工作。
如果我们注意观察它一个小时，就可以看到这种全部的动作，不下五次以上，连附
近的散步也包括在内。它产卵的地点，常是靠得很近的。
各种工作都已经完成以后，我察看这种小穴。只有卵放在那里，没有小室或壳来保
护它们。通常约有六十个，颜色大部分是紫灰色的，形状如同棱一样。
我开始观察螽斯的工作，就想看看它的卵子孵化的情形，于是在八月底的时候，我
取来很多的卵，放在一个里面铺有一层沙土的玻璃瓶子中。它们在里面度过八个月的时
间，感受不到气候变化带来的痛苦：没有风暴，没有大雨，没有大雪，也没有它们在户
外必须经受到的、过度炎热的光照和日晒。
六月来临时，瓶中的卵，还没有表现出开始孵化的征兆。和九个月以前，我刚把它
们取来的时候一样，即不发皱，也不变色，反而表现出极其健康的外观。在六月里，小
螽斯在原野里经常可以遇到了，有的，甚至已发育的很大了。因此我很怀疑，究竟是什
么理由使它迟延下来的。
于是，就产生了一种意见，这种螽斯的卵，如同植物一样，被种在土地里，是毫无
保护地暴露在雨雪之中的。
在我瓶子里的卵，在比较干燥的状况下，度过了一年的三分之二的时间。因为它们
本来是像植物种子一样散播着的。它的孵化大概也需要潮，需要适合它的一切孵化条件，
如同种子发芽时需要潮湿一样。这时，我开始决定要试一试。
我将从前取来的卵，分出一部分，放在我的玻璃管里，在它们上面，薄薄的加上一
层细细的潮湿的沙子。然后把玻璃管用湿棉花塞好，以保持里面的湿度。无论谁看见我
的试验，都会以为我是那种在试验种子的植物学家。
我的希望可以实现了。在温暖的，潮湿的环境之下，卵不久就表示出要孵化的迹象，
它们渐渐的，一点点地涨大，壳显然就要分裂开了。我花费了两个星期的功夫，每个小
时我都很认真仔细、不知疲倦地守候着它，想看看小螽斯跑出卵来的情形，以解决遗留
在我心中很长时间的疑问。
那个疑问是这样的。这种螽斯，按照惯例，是埋在土下边约一寸深的地方，现在这
个新生的小螽斯，夏初时在草地上跳跃，发育得完全一样，长有一对很长的触须，细得
如同发丝一般；并且身后生有两条十分异常的腿——像两条跳跃用的支撑杆，对于走路
是很不方便的障碍。
我很想知道，这个柔弱的小动物，携带着这样笨重的行李，当它到地面上来时，其
间所有的工作，是怎样进行的呢？它用什么东西从土中开出一条小道路来呢？它有遇到
一粒小沙就会折断的触角，少许的力量就会断脱的长腿，这个小动物是显然不可能从土
坑中解放出来的。
我已经告诉过你们：蝉和螳螂，一个从它的枝头、一个从它的巢出来时，穿有一种
保护物，就像一件大衣一样。
我想，这个小螽斯，从沙土里钻出来的时候，一定也有比出生以后，在草间跳跃时
所穿的还要简单而且又紧又窄的衣服，作为一种保护。
我的估计并没有错。这时候，白面孔螽斯，和别的昆虫一样，的确穿有一件保护外
衣。这个细小的，肉白色的，小动物，已经长在一个鞘里了，六个足平置胸前，向后伸
直。
为了让出来时比较容易一些，它的大腿绑在身旁；另一半不太方便的器官——触须
——一动也不动地压在包袋里面。
它的颈弯向胸部。大的黑点——是它的眼睛，那毫无生气而且十分肿大的面孔，使
人以为那是盔帽。颈部则因头弯曲的关系，十分开阔。它的筋脉同时微微地跳动着，时
张时合，因为有了这种突出的、可以跳动的筋脉，新生的螽斯的头部才能自由转动。依
赖颈部推动潮湿的沙土，挖掘出一个小洞穴。于是筋脉张开，成为球状，紧塞在洞里，
在它的幼虫移动它的背，并推土时，可以有足够的力量。
如此，进一步的步骤已经成功了，球泡的每一次涨起，对于小螽斯在洞中的爬动，
都是很有帮助的。
看到这个柔软的小动物，身上还是没有什么颜色，移动着它那膨胀的颈部，攒掘土
壁，真是可怜。
它的肌肉还没有达到强健的时候，这真无益于与硬石的斗争啊！不过经过不懈的奋
斗它却居然获得了最终的成功。
一天早晨，这块地方，已经做成了小小的孔道，不是直的，约有一寸深，宽阔得像
一根柴草。一般用这样的方法，这个疲倦的昆虫终于可以达到地面上了。
在还没有完全脱离土壤以前，这位奋斗者也要休息一会儿，以恢复它这次旅行后的
精力。再作一次最后的拼搏，竭力膨胀头后面突出的筋脉，以突破那个保护它已经很久
的鞘。这个动物就这样将外衣抛弃了。
于是，这是一个幼小的螽斯了，它还是灰色的，但是，第二天就渐渐变黑了，同发
育完全的螽斯比较起来简直是成了一个黑奴了。不过它成熟时的象牙面孔是天生的，在
大腿之下，有一条窄窄的白斑纹。
在我面前发育的螽斯啊！在你面前展开的生命是太凶险了。
你的许多亲属们，在尚没有得到自由之前，就因疲倦而死去了。在我的玻璃管中，
我看到了好多螽斯因受到沙粒的阻碍而放弃了尚未成功的奋斗。
它的身上长有一种绒毛，欲将它的尸体包裹起来。如果我不去帮助它，到地面上来
的旅行会更加危险，因为屋子外面的泥土更加粗糙，已经被太阳晒硬了。
这个有白条纹的黑鬼，在我给它的莴苣菜叶上咬啮，在我给它居住的笼子里跳跃着，
我可以很容易地豢养它。
不过它已不能再提供给我更多的知识了，所以，我就恢复了它的自由，以报答它教
给我的那些知识，我送给它这个房子——玻璃管，还有花园里的那些蝗虫。
因为它教给我蚱蜢在离开产卵的地点时，穿着一件临时的保护衣服，将那些最笨、
最重的部份，如它的长腿和它的触角等等，全都包在鞘里。它又告诉我这种略微伸缩、
干尸状的动物，为了它旅行的方便，它的头颈上生有一种瘤，或者说是颤动的泡口——
是一种原来就生成的机器，在我最初观察螽斯的时候，我并没有看见它用它作为走路的
帮助。
------------------








黄蜂



一、它们的聪明和愚笨
在九月里的一天，我和我的小儿子保罗跑出去，想去瞧一瞧黄蜂的巢。
小保罗的眼力非常好，再加上特别集中的注意力，这些都有助于我们的观察很好地
进行。我们两个饶有兴趣地欣赏着小径两旁的风景。
忽然，小保罗指着不远的地方，冲着我喊了起来：“看！一个黄蜂的巢。就在那边，
一个黄蜂的巢，比什么都要更清楚呢！”果然，在大约二十码以外的地方，小保罗看见
一种运动得非常快的东西，一个一个地从地面上飞跃起来，立即迅速地飞去，好像那些
草丛里面隐避着小小的即将爆发的火山口，马上要将它们一个个喷出来一般。
我们小心谨慎地慢慢地跑近那个地点，生怕一不小心，惊动了这些凶猛的动物，引
起它们对我们的注意和攻击，那样的话，后果可是不堪设想的。
在这些小动物们的住所的门边，有一个圆圆的裂口。口的大小大约可容下人的大拇
指。同居一室者来来去去，进进出出，摩肩接踵地向相反的方向飞去飞回，不停地忙碌
着。
突然，“噗”的一声，我不觉吃了一惊，但是马上又醒悟过来了。我忽然想起现在
我们正处于一个很不安全的时刻。要是我们太靠近去观察它们的行踪，就会引起不良的
后果。因为，这样的不速之客会让它们感到不安，会激怒这些容易发脾气的战士来袭击
我们。因此，我们不敢再多观察了。再观察下去就意味着要“牺牲”更多的东西了。
我和小保罗记住了那个地点，以便日落后再来观察。到了夕阳西下的时候，这个巢
里的居住者，全体都应该从野外回家了。那样，我们就可以更好地观察了。
当一个人决定要征服黄蜂的巢时，如果他的这一举动，没有经过谨慎而细致的思考
的话，那么这种行动简直就是一种冒险的事情。半品脱的石油，九寸长的空芦管，一块
有相当坚实度的粘土，这些构成了我的全部武器装备。还有一点必须提到的是，以前的
几次小小的观察研究，稍稍积累了一点儿成功的经验。这所有的一切物品与经验对我而
言，是最简单，同时也是再好不过了。
有一种方法对我是至关重要的，那就是窒息的方法。除非我打算用我所不能够忍受
的牺牲的方法。否则，我必须掌握窒息的方法。当瑞木特要把一个活的黄蜂的巢放在玻
璃匣子内，观察里面的同居者的习性的时候，他不是亲自行动，而是选择了另一种方法，
雇佣了一个帮手，协助他进行实验。这个帮手经常从事这种痛苦不堪的工作。为了获得
优厚的报酬，情愿牺牲自己的皮肤，为科学家们提供有偿性服务。但是，我可是打算牺
牲自己的皮肤的。
在还没有挖出我所要的蜂巢之前，我仔细思考了两次。然后，才开始我的计划。我
首先将蜂巢里的居民窒闷住，死了的黄蜂就不能刺人了。这是一个残忍的方法，但也是
一个十分安全的方法，可以让我不至于身处危险之中。我采用的是石油，因为它的刺激
作用不至于过于猛烈。
因为我要做一次观察，所以我希望能留下一部分不死的黄蜂，否则的话老是观察死
了的对象，就前功尽弃了。现在的问题只是在于如何把石油倒进有蜂巢的穴里去。蜂巢
穴的出入孔道大约有九寸长，而且差不多和地面是平行的，一直通到地底下的窠巢。假
如把石油直接倒入隧道的口上，这便是一个大大的错误，而且将会带来极其严重的不良
后果。为什么呢？主要原因是，这样少量的石油，会被泥土吸收进去，而无法到达地下
的窠巢。这样一来，到了第二天，当我们凭着想象，以为这时挖掘、凿开窠巢一定是很
安全的时候，就会遇到很大的危险。我们就会碰到一群火上浇油般的黄蜂，在我们的铁
铲下回旋，从而对我们造成一定的威胁。
早已准备好的九寸长的空芦管可以阻止这一不幸事件的发生。把这根空芦管插进差
不多九寸长的隧道里面的时候，就形成了一根自动引水管。于是石油可以顺着导管流人
土穴中，一点儿也不会漏掉，而且，速度还很快。然后，我们再用一块事先已经捏好的
泥土，像瓶塞子一样，塞住出入的孔道口，断绝这些黄蜂的后路。我们所要做的工作就
到此为止了，剩下的就只是等待了。
当我们准备做这项工作的时候，是昏暗的月夜，正是九点钟，小保罗和我一起出去。
我们只带了一盏灯，还有一篮子需要用到的工具。当时，远远的还可以听见农家的狗还
在互相吠叫着，猫头鹰在橄榄树的高枝上叫着，蟋蟀在浓密的草丛中不停地奏着动听的
音乐。小保罗和我则在谈论着昆虫。他热切而好奇地向我提出很多问题。为了不让他失
望，我将我所知道的一切告诉了他，帮助他学习，以丰富他的知识，满足的他的兴趣。
这样一个快乐的猎取黄蜂的夜晚，让我们忘记了睡眠和被黄蜂攻击时的痛苦。
将芦管插入土穴中是一件非常精巧的工作，需要一些技巧。因为孔道的方向是无从
知晓的，需要颇费一番猜疑和试探。而且有时候，黄蜂保卫室里的门卫会突然警觉地飞
出来，毫不客气地攻击正在进行这项工作而且没有防备的人的手掌。为了防止这种措手
不及的不幸事情发生，我和小保罗中的一个人，在一旁守卫，时刻警惕着，并用手帕不
停地驱赶着进攻的敌人。这样一来，即使最后有一个人的手上不幸被命中，隆起了一块，
就是很疼痛，也算是一个理想的，不很大的代价，尚可以忍受。
在石油流入土穴中以后，我们便听到地下传来的众蜂惊人的喧哗声。然后，很快地，
我们用湿泥将孔道封闭起来，一次一次地用脚踏实，使封口坚不可摧，从而使它们无路
可逃。现在，没有什么其他的事可以做了。于是，我和小保罗就跑回去睡觉休息了。
第二天清晨，我们带了一把锄头和一把铁铲，重新又回到了老地方。早一点儿去比
较好些，因为可能有很多黄蜂夜里是在外面游荡的，它们有可能在我们挖土的时候飞回
来，这就糟糕了，因为这对我们又将是一种威胁。另外，清晨的冷气，可以多少削弱一
些它们的凶恶和威风。
在孔道之前，芦管依然还插在那边，我和小保罗挖了一条壕沟，宽度刚好能容下我
们俩，行动很方便。于是，我们从沟道的两边开始挖。很小心地一片一片铲去。后来，
挖了差不多有二十寸深，蜂巢便暴露出来了。它吊在土穴的屋脊当中，一点儿也没有被
损坏，完好地吊在那里，这真让我们感到高兴。
这真是一个壮观美丽的建筑啊！它大得简直像一个大南瓜。除去顶上的一部分以外，
各方面全都是悬空的，顶上生长有很多的根，其中多数是茅草根，穿透了很深的“墙壁”
进入墙内，和蜂巢结在一起，非常坚实。如果那地方的土地是软的，它的形状就成圆形，
各部分都会同样的坚固。如果那地方的土地是沙砾的，那黄蜂掘凿时就会遇到一定的阻
碍，蜂巢的形状就会随之有所变化，至少会不那么整齐。
在低巢和地下室的旁边，常常留有手掌宽的一块空隙，这块面积是宽阔的街道。这
些建筑者，在这里可以行动自由，继续不停地进行它们各自的工作，用它们自己的双手，
使它们的窠巢更大更坚固。通向外面的那条孔道，也通向这里。在蜂巢的下面，还有一
块更大一些的空隙，其形状是圆的，就如同一个大圆盆，在蜂巢扩建新房时，可以增大
其体积。这个空穴，还有另外一个用途，那就是盛废弃物品的垃圾箱。看来这里的基本
建设还是较为齐全的。
这个地穴是黄蜂们用自己的“双手”亲自挖掘出来的。关于这一点，是用不着怀疑
的。因为如此之大、如此整齐的洞穴，在自然界是没有现成的。当初，第一个开辟这个
巢的黄蜂，也许是利用了鼹鼠所做的洞穴，加以借用，以图开始创建的便利。可是，筑
巢的绝大部分工作都是黄蜂亲自操作的。然而，事实上，并没有一些挖出的泥土堆积在
蜂巢的大门之外。那么，黄蜂们挖出的泥土被搬运到哪里去了呢？答案是：它们已经被
弃散在不引人注意的广阔的野外去了。有成千上万的黄蜂参与挖掘这个壮丽的建筑物，
必要的时候，还要将它扩大。这千百万只黄蜂，飞到外面来的时候，每一个身上都附带
着一粒土屑，抛散在离开窠巢很远的各处土地上。因此，挖出的泥土的痕迹一点儿也看
不到了。所以，蜂巢看上去像一片净土一样。
黄蜂的巢是用一种薄而柔韧的材料做成的。这种材料是木头的碎粒，很像一种棕色
的纸。它的上面有一条条的带，其颜色视所用木头的不同而不同。如果蜂巢是用整张
“纸”做的，就可以稍稍抵御寒冷，起到保暖的作用。但是，黄蜂就像做气球的人一样，
它们懂得温度可以利用各层外壳中所含有的空气来保持。所以，黄蜂把它们的低巢做成
宽的鳞片的形状，一片一片松松的铺起来，显出很多的层次来，整个蜂巢形成一种粗粗
的毛毯状，厚厚的，而且多孔，其内部含有大量的空气。这样一来，外壳里的温度，在
天气很热时，一定是很高的。
大黄蜂——黄蜂们的领导，在一样的原则下，建筑它自己的巢。在杨柳的树孔中，
或者是在空的壳层里，它用木头的碎片，做成脆弱的黄色的纸板。它就利用这种材料来
包裹它自己的窠。一层一层相互地重叠起来，就像个凸起的大鳞片一样，可以想象这有
多么保暖！这个大鳞片的中间有充分的空隙，空气停留在里边也不流动。
黄蜂们的动作常常与物理学和几何学的定理相吻合。它们可以利用空气——这个不
良导体来保持它们家里的温度。它们早在人类还未曾想到做毛毯之前就已经做出来了，
而且技艺还很高，它们在建筑窠巢的外墙时，只要极小的外围，就足以造出很多的房间，
它们的小房间也同样如此，其面积与材料都非常经济。
然而，尽管这些建筑家们有如此之聪明智慧，但是，令我们感到奇怪的是，当它们
遇到最小的困难时，居然会束手无策，愚笨无比。一方面，它们得益于大自然的本能指
导它们像科学家一样地行动；而另一方面，很显然它们完全不具备反省的能力，其智力
是相当低下的。关于这一个事实，我已用各种各样的试验加以证明了。
黄蜂碰巧将自己的房子安置在我家花园的路旁边，于是，我便可以利用一个玻璃罩
来做试验了。在原野里的时候，我无法利用这种器具，因为乡下的小孩子们很快就会把
它打破，而破坏了我精心准备的试验。有一天晚上，天已经黑了，黄蜂也已经回家了。
我弄平了泥土，放上一个玻璃罩罩住黄蜂的洞口。第二天早晨，黄蜂们习惯性地开始工
作。当它们发觉自己的飞行受到阻止时，它们是否能够在玻璃罩的边下挖掘出另外一条
道路呢？是不是这些能够掘出广阔洞穴的刚强的动物知道只要创造一条很短的地道，便
可以让它们重获自由呢？这便是我们的问题关键之所在了。那么，结果如何呢？
第二天早晨，我看到温暖耀眼的阳光已经落在玻璃罩上了。这些工作者们已经成群
地由地下上来，急于要出去寻觅它们的食物。但是，它们一次又一次地撞在透明的“墙
壁”上跌落下来，重新又上来。就这样，成群地团团飞转不停地尝试，丝毫不想放弃。
其中有一些，舞跳得疲倦了，脾气暴躁地乱走一阵，然后重新又回到住宅里去了。有一
些，当太阳更加炽热的时候，代替前者来乱撞。就这样轮换着倒班。但是，最终没有一
只黄蜂大智大勇，能够伸出手足，到玻璃罩四周的边沿下边抓、挖泥土，开辟新的谋生
之路。这就说明它们是不能设法逃脱的。它们的智慧是多么有限啊。
这个时候，有少数在外面过夜的黄蜂，从原野归来了。它们围绕着玻璃罩盘旋飞舞，
一直迟疑徘徊，不知如何是好。有一个带头决定往玻璃罩的下边去挖，其他的黄蜂也随
着学它的样子。于是，大家齐动手，很快地，一条新的通路很容易地开辟出来了。它们
也就跑了进去，终于到家了。于是，我用土将这条新辟之路堵住。假设从里面能够看出
这条狭窄的通路，当然可以帮助罩内的黄蜂轻而易举地逃走。我很愿意让这些囚徒通过
自己的观察和努力争得自由的光荣，享受阳光沐浴的欢乐，领会大自然的优美。
无论黄蜂的理解能力是如何的差劲。我想它们的逃脱，现在应该是可能实现的了。
那些刚刚进去的黄蜂当然会指引一下路径，它们会指教其他的黄蜂向玻璃罩下边挖，以
便尽快地逃离牢笼。
然而，事实却并不那么乐观。我非常失望，可爱的黄蜂们居然没有一点儿要从经验
和实例上仿效学习的企图。在那个玻璃罩里，一点儿没有要继续挖掘出逃之路——地道
的迹象。这些小昆虫们只是依旧团团乱飞，毫无计划，毫无目的，它们只是盲目地乱碰
乱撞，挤作一团不知究竞发生了什么意外。每天都有很多可怜的黄蜂死于饥饿和炎热之
下。一个星期以后，很遗憾，没有一个活的黄蜂能够侥幸存活下来，全军覆没了。一堆
死尸铺在地面上，其状况尤其惨烈。
从原野里返回的黄蜂们可以另辟新路，毫不费力地回到自己的家中。其原因是，从
泥土外面可以嗅知它们的家，并去寻找它。这是黄蜂自然本能想方设法投入家的怀抱的
一种表现，或者说是它们的一种防御方法。这是不需要任何思想和解释的。自从小小的
黄蜂初次降临到这个世界上的时候开始，地面上的一切阻碍，对于每一个黄蜂而言，都
已经很熟悉了。
但是，对于那些不幸被罩在玻璃罩里的黄蜂，就没有这种本能来帮助它们逃离险境
了。它们的目的是明确的，单一的。它们想到阳光里面去，到野外去觅食。它们被罩在
玻璃罩里，在这个透明的牢狱中，能够看到日光，它们便被蒙骗了，以为自己的目的已
经完全达到。虽然它们几经努力，一往无前，继续不断地和玻璃罩相抗衡、相碰撞，心
中抱有无限希望，想朝着日光，飞得再远一点儿。以便能觅到急需的食物。可是事实上
那是无用的。在它们以往的经历中，没有任何经验和实践指导它们遇到这种情况时，应
该如何行事。于是它们走投无路，别无选择，只能盲目地固守着它们生来就惯有的老习
性，从而生的希望越来越小，而逐渐将自己推向无奈的死亡。

二、它们的几种习性
假如我们掀开蜂巢的厚包，便可以看到里面隐藏有许多的蜂房，那好几层的小房间，
上下排列着，中间用稳固坚实的柱子紧密连在一起，层数是不一定的。在一定季节的后
期，大概是十层，或者是更多一些。各个小房间的口都是向下的。在这个种类看起来很
奇怪的小世界里，幼蜂无论是睡眠还是饮食，都是脑袋朝着下边生长的，即倒挂着的。
这一层一层的楼即蜂房层，有广大的空间把它们分隔开。在外壳与蜂房之间，有一
条门路与各个部分是相通的。经常有许多的守护者进进出出，负责照顾蜂巢中的幼虫。
在外壳的一边，矗立着这个丰富多彩的都市的大门，一个没有经过什么过多装饰的裂口，
隐藏在被包着的薄鳞片中。直面对着这个大门的，就是那从地穴深处直通到外面的大千
世界的隧道的进出口。
在黄蜂的社会中，生活着数量众多的黄蜂。它们的全部生命完全都投入到不辞劳苦
的工作之中。它们的主要职责就是，当人口不断增加的时候，就不停地扩建蜂巢，以便
新的公民居住。尽管它们并没有自己的幼虫，可它们呵护巢内的幼虫，却是极小心勤勉，
无微不至的。
为了能观察到它们的工作状况，以及快到冬天的时候它们之中会有什么事情发生，
我在十月里，把少许巢的小片放在盖子下面，里面居住着很多的卵和幼虫，并且还有一
百个以上的工蜂在细心地看护着它们。
为了便于进行观察，我将蜂房分隔开来，让小房间的口朝着上面，然后并排放着。
这样颠倒的排列，看起来似乎并没有使我的这些囚徒们烦恼，它们很快地就从被打扰的
情形下适应过来，恢复了原来的空间状态，重新开始忙碌而辛勤地工作，似乎从来没有
什么事情发生过一样。
事实上，它们当然需要再建筑一点东西。所以，我便选择了一块软木头送给它们，
并且用蜂蜜来喂养它们，满足它们的需要。用一个拿铁丝盖着的大泥锅来代替隐藏蜂巢
的土穴。再盖上一个可以移动的纸板做的圆顶形的东西，使得内部相当黑暗。当然，当
我需要亮一些时就把它移开。
黄蜂继续进行它们自己的日常工作，就好像从来没有受到过任何的扰乱一样。工蜂
们一面照料着蜂巢中的蜂宝宝，与此同时，又要照顾好它们自己的房子。它们一起努力
加油，开始慢慢地筑起一道新的铜墙铁壁。这墙壁围绕着它们最封闭的蜂房。看起来，
它们似乎是打算重新再建筑一个新的外壳，作为新的外壳，来代替那个被我用铁铲毁坏
了的旧外壳。但是，这些工蜂们并不是简单地修修补补，它们是从被我破坏了的那个地
方开始它们的工作。它们很快就筑成了一个弧形的纸鳞片似的房顶，然后，用它遮盖住
大约三分之一的蜂房。如果这个小蜂巢不曾遭到我的破坏的话，那么这些工蜂们搭建起
的这个屋顶足可以连接到外壳呢。它们亲手做成的一个房顶，还不够大，只能遮盖住整
个小房间的一部分而已。
至于我事先为它们精心准备好的那块软木头，它们根本不予理睬，甚至连碰都不曾
碰上一下，仿佛它根本就不存在一样。或许这种“新型”的材料，对于黄蜂而言，用起
来很不方便。它们宁愿放弃它，而继续选用那已经废弃不用了的旧巢，这样更加方便，
而且更加得心应手一些。因为在这些旧的小巢内，不必辛辛苦苦地重新制做纤维，因为
它们是已经做好了的，方便实用。只需拿来主义即可。而且，它们也不用浪费很多的唾
液，只需相当少的唾液，再用它们的大腮仔细咀嚼几下，然后便形成了上等质地的浆糊，
这是相当好的建筑材料。
下一步，它们一起把不居住的小房间统统毁得粉碎。然后，利用这些碎物，做成一
种似天篷一样的东西。如果有必要的话，它们也会再次利用同样的方法，筑造出新的小
房间，以便居住，活动之用。
与它们齐心协力筑造屋顶的工作相比较，更加有趣味的要算是喂养蛴螬幼虫了。刚
才还是一个个粗暴刚强，卖力气的战士，这会儿就摇身一变，成了温柔、体贴的小保姆。
看到这些，谁也不会感到厌倦和反感的。一下子，充满了战斗气息的军营一样的窠巢，
立刻变成了温馨的育婴室了。真是妙趣横生啊！
喂养好可爱的、柔弱的小宝宝，可是需要相当的耐心与细致的。假如我们只将注意
力集中到一个正在忙碌工作的黄蜂身上，我们就可以清楚地观察到，在它的嗉囊里，充
满了蜜汁。它停在一个小房间的前面，它的样子特别有意思，它把它小小的头慢慢地伸
到洞口里面去，然后再用它的触须的尖儿去轻轻地碰一碰里面的一个小幼虫。那个小宝
宝慢慢地清醒了过来，似乎看到了那个黄蜂递送进来的触须，于是向它微微地张开小嘴。
它的样子，特别像一只刚刚初生不久，羽毛尚未丰满起来，乳臭未干的小鸟，正在向着
刚刚辛辛苦苦为它觅食而归的妈妈伸出小嘴，急切地索要食品一般，不觉得让人感到一
阵温馨。
不一会儿，这个刚刚从梦中苏醒过来的小宝宝，将它的小脑袋摇来摆去的，渴望着
能够马上探索到它急切需要得到的食物，这可以算是它的本能天性了。然而它又是盲目
地探寻着，一次次试探着外面的黄蜂为它们提供的食物。可以想象小宝宝的急切心情网，
终于两张小嘴接触到了。一滴浆汁从“小保姆”的嘴里流出来，流了进那个被看护者的
小嘴里。仅仅这一点点就足够一个小宝宝享用了。现在，该轮到第二个黄蜂婴儿进食了。
于是，这个小保姆又赶快马不停蹄地跑到别处去，继续履行它神圣的职责。
小宝宝们通过口对口地交接食物后，享受到大部分的蜜汁。但是，进食并没有完全
告终，它们还没有享用完呢。因为，在喂食的时候，幼虫的胸部会暂时膨胀起来，其作
用就如同一块围嘴或餐巾纸一样，从嘴里流出来的东西全都滴落在它上面。这样等保姆
走后，小宝宝们就会在它们自己的颈根上舐来舐去，吮吸着滴在胸部的蜜汁，尽情地享
受着美味的食物，不浪费一点儿食物。大部分的蜜汁咽下之后，幼虫胸部的鼓胀便会自
然而然地消失掉。然后，幼虫会稍微往蜂巢里缩进去一点，继续回到它的甜蜜的梦乡里。
当黄蜂在我的笼子里喂养小宝宝时，小幼虫们的头是朝上的，从它们的小嘴里遗漏
出来的东西，自然会滴落在它们的围嘴上面。至于在蜂巢里喂养它们的时候，它们的小
脑袋则是朝下的，但是，我并不心存怀疑，那便是在这样头向下的位置上，小幼虫的餐
巾纸仍能发挥作用，而且，功效是一样的。这是因为幼虫在蜂巢中时，它的头不是直的，
而是略微有一点弯度的。因此，它们嘴里溢出来的蜜汁很可能是堆积在那块小小的围嘴
上。而且，溢出的蜜汁是非常粘稠的，很快就会粘在围嘴上。与此同时，细心的小保姆
就是再放下一部分食品在这个地方，也是有可能的。所以，无论小幼虫的头是朝下的还
是朝上的。无论那块围嘴是在嘴的上边，还是在嘴的下边，这都不会阻碍围嘴充分发挥
它的作用。其主要原因是这种食品非常有粘性，可以牢牢地附着在围嘴上。因而可以说，
这块小小的围嘴简直就是一个又方便又及时的小碟子，它可以减少喂食工作的困难，避
免许多不必要的麻烦。为我们的小保姆们提供方便，使得它们又省力又省时，而且，它
还可以使得小幼虫们能够舒适、宁静地享用它们的美味佳看，直至一饱口福，满足为止。
还有一个好处，那就是不致让小宝宝们吃得太饱，撑坏了小肚皮而夭折。
如果是在野外，置身于大自然中，每当一年快要结束时，也是果品数量非常少的时
候，有些青黄不接。在这种情况下，大多数的小保姆挑选其它的食物来继续喂养小幼虫。
它们大多选择苍蝇，先将它们一一切碎，然后再喂给小幼虫们食用。但是，在我为它们
制作的笼子中，一概不选择其他的东西作为幼虫的食品，我只为它们提供充足的有营养
的蜜汁。
吃了这些蜜汁以后，所有的看护者和被看护者似乎都变得精力旺盛起来。而且，一
旦有什么不速之客突然闯进蜂房里，进行袭击侵略，那么它们将很不幸地立刻被处以死
刑。显然，黄蜂分明是一种不好客的生灵，从不厚待宾客，更不允许其他动物随意侵扰
自己的家园。即使是那种叫拖足蜂的蜂，形状和颜色与黄蜂是极其相像的，如果它们稍
一走近黄蜂，来分享它们的蜜汁，那么它们的这种企图很快就会破灭，马上被觉察出来，
于是，黄蜂们群起而攻之，直到致其于死地为止。拖足蜂的外貌并不能欺骗黄蜂那敏感
的目光，如果拖足蜂反应不迅速，没有及时退避，那么就会大难临头引来杀身之祸，被
黄蜂残酷地处死。所以，擅自闯入黄蜂的巢，实在不是一件明智的事情。即使来访的客
人的外表与它们极为相似，工作也与它们大同小异，几乎可以说是团体中的一份子，都
是绝对不行的。黄蜂是不会轻易地放过任何不请自来，不知趣的所谓客人的。因此，其
它动物还是退避三舍，回避为佳。
我已经一而再、再而三地看到黄蜂对于客人们的野蛮的待遇，假如闯入境内的不速
之客，是个相当有杀伤力，而且凶猛无比的家伙，当它受到群攻而牺牲后，其尸首便会
马上被众蜂拖到蜂巢以外，抛弃在下面的垃圾堆里边。但是，黄蜂似乎不会轻易地动用
它那有毒的短剑来攻击其他的动物，还是比较手下留情的。如果我把一个锯蝇的幼虫抛
到黄蜂群里面，对于这条绿黑色的小龙一样的侵入者，黄蜂们表示出很大的兴趣，它们
一定是感到很奇怪。接下来，它们便向它发起进攻，把它弄伤，但是并不利用它们带毒
的针去刺伤它。然后众蜂齐用力，要把它拖出巢去。与此同时，这条“小龙”也不服输，
不断进行抵抗，用它的钩子钩住蜂房。有时利用它的前足，也有时利用它的后足。然而，
最终这条可怜的“小龙”还是因为伤势太重，而且它还很软弱，最终被有力的黄蜂拉了
出来。这条“小龙”很惨，小小的身体上充满了血迹，被一直拖到垃圾堆上去。黄蜂们
驱赶这样一条并无什么力气的可怜虫可并不轻松，耗费了足有两个小时的时间呢！
如果，与此相反，我放的不是一个弱小的幼虫，而放一个住在樱桃树孔里的一种相
对比较魁伟的幼虫在蜂巢里面，结果就不同了。立刻会有五六只黄蜂拥上来，纷纷用有
毒的针去刺它的身体。不大一会儿，差不多几分钟以后，这只较强壮有力的幼虫终于也
难逃恶运，一命呜呼了。但是，接下来便产生出了一个问题。这具笨重的尸体，很难把
它搬运到巢外去。所以，黄蜂们觉得动它不成，便选择了其它的方法，比如吃掉它，或
者，至少要设法使它的体重有所减少。因此，它们便一直吃它，直到吃剩下的那部分可
以被拖动为止。然后，还是要把它拖到外面去，抛弃掉。

三、它们悲惨的结果
有了如此凶猛而又残酷的方法来抵御外来闯入者的入侵，还有如此巧妙而又温柔的
喂食方法，我制作的笼子里的小幼虫们一天一天茁壮成长着，黄蜂的家族日益兴旺起来。
不过，当然也存在例外的现象。黄蜂的窠巢里，也有一些非常柔弱、不走运的小幼虫，
它们长大成人，还未经历世间的风雨，沐浴阳光的温暖，便早早地夭折而去了。
我通过观察，发现了那些柔弱的病者，亲眼目睹它们不能继续享用蜜汁，不能进食，
渐渐地，一点点憔悴下去，衰弱下去。那些小保姆们早已比我更清楚地知晓了这一切。
它们十分无奈地把头轻轻弯下来，朝着那些可怜的患病者，用触须很小心地去试听一下，
最后得出结论，证明这些病者的确是不可医治，无法挽救了。于是，慢慢地，这个弱小
的生命逐步走向生命的尽头，快到死的程度了。最终，被毫不怜惜地从小房间里拖到蜂
巢的外面去。在充满野蛮气息的黄蜂的社会里，久病者不过仅仅是一块没有用处的垃圾
而已，越是赶紧拖出去越好，否则的话，就有蔓延传染的可能。对于黄蜂而言，那将是
很可怕的事情。但是这还不是最坏的可能。因为，随着冬天渐渐来临，黄蜂们大都已经
预感到它们将来的命运。它们深知，末日就在眼前了。
十一月里，非常寒冷的夜，使得蜂巢的内部起了变化。大搞基础建设的热情逐渐衰
退了。到储蜜的地方，从事储蜜工作的黄蜂不再频繁地到那里去了。整个家庭，所有黄
蜂全都逐渐地放任自流了。幼虫由于饥饿而大张着它们的小嘴，然而，等到的只不过是
非常迟缓的救济品，或者干脆没有小保姆愿意光临到这里来给它们喂食。深深的惆怅牢
牢占据了那些小保姆的心灵，它们从前的那份工作热情也不见了，最终竟转化为厌恶。
它们知道，再过不久的时间，一切就将变成不可能了。那么，小保姆还有什么存在价值
吗？还会有什么好处吗？当然，在看护蜂的心中，答案是否定的。于是，饥饿的时候来
临了。噩运降临到小幼虫的头上，它们悲惨而孤独地死去。因此，从前的那些温柔体贴
的小保姆一变而成为不可思议的凶残的刽子手了。
那些小保姆会对自己说：“我们没有必要留下许许多多的孤儿。不久以后，等我们
都离开了这里以后，还能有谁来照顾这些可怜的后代呢？没有。既然是这样的结果，那
还不如让我们亲手来把这些卵和小幼虫统统杀死。这样一个十分残暴的结果，总比那种
慢慢被饥饿煎熬而死要强得多，长痛不如短痛嘛！”
接下来的一幕闹剧，便是一场凶残的大屠杀行动。黄蜂们残忍地咬住了小幼虫颈项
的后面，然后粗暴地把它们一个个从小房间里拖出来，拉到蜂巢的外面去，抛到外面土
穴底下的垃圾堆里，其情景简直是惨不忍睹！
那些小保姆，也就是工蜂，在把幼虫从小房间中强行拖拉出来时，那种情形之残酷，
就好像这些幼虫都是一些从外面来的生客一般，或者是一群已经死掉了的尸体。它们野
蛮地拖着小幼虫的尸体，并且还要将它们的尸体扯碎。至于那些小卵，则会被工蜂们撕
扯开来，最后把它们吃掉。
在此之后，这些小保姆，即刽子手，毫无生气地留着它们自己的生命。一天一天地，
我带着无比的惊奇，注视着我的这些昆虫的最终的结局。非常出乎我的意料，这些工蜂
忽然间都死掉了。它们跑到上面，跌倒下来，仰卧着，从此再也没有爬起来，就如同触
了电一般。它们也有它们自己的生命周期。它们被时间这个无情无义的毒品毒死了。就
算是一只钟表内的机器，当它的发条被放开到最后一圈时，也是会如此的。
工蜂是老了！然而，母蜂是蜂巢中最迟生出来的，它们既年轻，又强壮。所以，当
严冬降临，威胁到它们时，它们还仍有能力来抵挡一阵。至于那些末日已经临近的，很
容易地就能从它们的外表的病态上分辨出来。在它们的背上，是有尘土沾附着的。在它
们尚健壮，还年青的时候，它们一旦发现有尘土附着在身上，就会不停地拂拭，把它们
黑色、黄色的外衣清洁得十分光亮。然而，当它们有病时，也就无心注意卫生清洁了。
因为已经无暇顾及了。它们或是停留在阳光底下，一动也不动，或者很迟缓地踱来踱去。
它们已经不再拂拭它们的衣裳了，因为这已不再重要了，也没有任何意义了。
这种对装束的不在意，就是一种不祥的征兆。过了两三天以后，这个身上带有尘土
的动物，便最后一次离开它自己的巢穴。它跑出来，主要是打算再最后享受一点日光的
温暖。忽然，它跌倒在地上，一动也不动，再也不能够重新爬起来了。它尽量避免自己
死在它所热爱和生存的巢里。这是因为，在黄蜂中，有一种不成文的“法律”规定，那
就是巢里是要绝对保持干净整洁的。这个生命即将结束的黄蜂，要自行解决它自己的葬
礼。它把自已跌落在土穴下面的坑里。由于要保持清洁卫生，这些苦行主义者，不愿意
自己死在蜂房里。至于那些剩余下来的，还没有死去的黄蜂，它们仍然要保留这种习惯，
直到它们最终的结局为止。这形成了一种不曾被摒弃的法律条文。无论黄蜂世界中的人
口如何增加，或是减少，这一传统总是要保持遵守的。
我的笼子里，一天天地空起来了。虽然这个屋子仍然是暖和的，而且里面还储备有
很多的蜜汁，供剩下来的那些健康者食用。但是，到了圣诞节的时候，仅仅剩下了约一
打的雌蜂。到一月六日，连最后剩余下来的黄蜂也全都死掉了。
那么，这种死亡是从哪里来的呢？让我的黄蜂统统都倒毙了。它们并没有受过饿，
也没有挨过冻，更没有经历过离家的痛苦。那么，它们究意是为了什么而死的呢？
我们不应该归罪于囚禁，即便是在野外，也会发生同样的事情。在十二月末的时候，
我曾到野外去观察过很多的蜂巢，都曾经发生过同样的情况。大量数目的黄蜂，必须要
死亡，这并不是因为碰到了什么意外情况，也并不是因为疾病的担扰，或是因为某种气
候的摧残影响，而是由于一种不可逃脱的命运，这种命运摧残着它们，这和鼓舞着它们
生活下去的力量是一样有力的。不过，它们这样的生命，对于我们人类倒是很有好处的。
一只母黄蜂可以创造出一个拥有三万居民的城市。假如全体黄蜂都存活下来，那么，可
想而知，这将是一场多么大的灾难啊！若是那样的话，黄蜂就可以在野外构造自己的王
国，并且称王施虐了。
到了后期，蜂巢自己会毁灭的。一种将来会变成形状平庸的蛾子的毛虫，一种赤色
的小甲虫，还有一种身着鳞状的金丝绒外衣的小幼虫，它们都是有可能攻击，毁灭蜂巢
的小动物。它们会利用锋利的牙齿，咬碎一层层小巢的地板，使得整个蜂巢内的所有住
房全部崩塌毁坏。最后，剩下来的只有几把尘土和几片棕色的纸片。到了第二年春天到
来的时候，黄蜂们便又可以废物利用，白手起家，发挥大自然在建筑房屋方面赋予它们
的高度的灵性和悟性，建造起属于它们自己的新家园。新的结构精巧而且十分坚固的城
池，其中居住着约有三万居民——一个庞大的家族。它们将一切从零开始。它们将继续
繁衍后代，喂养小宝宝，继续抵御外来的侵略，与大自然抗争，为自己的安全而战斗，
为蜂巢内部生活的快乐而贡献自己的一份力量。生命不息奋斗不止！
------------------








蛴螬的冒险



一、蜂螨
围绕着卡本托拉斯(Carpenras)乡下沙土地的高堤一带，是黄蜂和蜜蜂最喜欢光临的
地方了。它们为什么会如此喜欢这个地方呢？究其原因，主要是因为这一地区的阳光非
常充足，而且这一带还非常容易开凿，很适合黄蜂和蜜蜂在这里安居乐业。在五月份这
样的天气，主要有两种蜜蜂特别的多。它们都是泥水匠蜂，是地下的一个个小屋的建造
者。其中的一种蜜蜂，它们在自己的住宅门口，建筑起一道自认为固若金汤的防御用的
壁垒——一个土筒。它的里面留有空白，而且整个筒是呈弧形的。筒的长和宽就像人的
一个手指头一样。有时候，会有很多蜜蜂飞到这一带来定居，当它们发现了这一个个斜
形的土手指的装饰以后，谁都会感到奇怪，不知道这是一些什么东西。
还有另外一种蜜蜂，也就是我们大家经常能够见到的，它们的名字叫做掘地蜂。它
们的走廊的外口没有什么手指形的防御壁垒，而是直接暴露在外面的。旧墙的石头之间
的缝隙之中，废弃的房舍，或者是沙石上头显露的表面，这些地方都非常适合于掘地蜂
的工作。但是，最最理想，最最适宜的地方，它们常常成群结伙地奔赴的地方，要算是
那些地面上突起的，朝着南方的直路。因为我经常可以看到它们开凿低路的处所。
这里的面积有好多码宽呢，而且墙上常常穿有很多很多的小孔，以致于这块地看起
来，呈海绵的形状。这些小小的洞孔，大概是用锥子戳出来的，因为它们是那么的整齐，
每一个孔穴都与盘曲的走廊相通相连，差不多都有四五寸深。蜂巢是在这底下的。如果
我们打算观察一下这种蜜蜂的工作情况，那么我们一定要在五月的下半月到它们的工作
场上来，但是千万要注意，必须保持一定的距离，要离开相当的远。这主要是出于安全
考虑。于是，我们会发现它们一群一群地汇合在一起，喧哗着，并且，众蜂齐努力，以
一种让人惊讶的毅力，从事着关于食物和蜂巢的各项工作。
但是，我来到这个被掘地蜂占领了的地方，次数最多的，要算是在八、九月间了，
这个时候正好是快乐、自由的夏天休假的时期。在这样的季节里，在靠近掘地蜂窠巢的
地方，都显得非常宁静。一切的工作都早已进行完毕了。在缝隙之中，有很多的蜘蛛拥
挤地呆在那里面，或者有丝管子伸入到蜜蜂的走廊里。从前住满了蜂，到处都熙熙攘攘、
热热闹闹的，现在仿佛变成了凄惨、悲凉的废墟一般。这其中的理由，我们谁也无从知
晓。距离大地表面约有数寸深的下面，有成千的幼虫被封闭在它们的土室之中。它们全
都静静地等候着春天的来临。当然，这些如此柔弱而又不能自己起来保护自己的幼虫，
是这样的肥胖，这样的有吸引力。一定足以引诱某种寄生者，或者招来某种正处在饥饿
之中到处寻觅食品的外来昆虫。这件事是非常值得研究注意一下的。
有两个事实，引起了我的留心。有一些非常丑陋的苍蝇，它们身上的颜色是半黑半
白的，这些苍蝇慢慢地从一个洞穴飞到另外一个洞穴里。它们这样飞来飞去的目的是要
表明它们在那些地方产卵。其中，有一些卵是挂在网上的，都早已干枯而死了。而在其
它的地方，比如，在堤上的蜘蛛网上，也挂了许多某种甲虫——蜂螨的尸体。在这些尸
体当中，有雌的也有雄的。不过，仍然还有少数是有生命的。雌性的甲虫，一定是伸入
到了蜂的住宅里面，而且，毫无疑问，它们一定是在蜂的窠巢中产下自己的卵。
如果我们慢慢地，稍稍掘开堤的表面，我们就会惊奇地发现更多有趣的东西。在八
月之初的时候，我们看到的是：顶上有一层的小房间，它们的样子和底下的蜂巢相比，
大不一样，相差甚远。之所以有这种区别，主要是因为这是由两种不一样的蜂建造而成
的。其中有一种是已经在前面提到过的掘地蜂，另外一种，有一个很动听的名字，叫竹
蜂。
掘地蜂组成了一只先锋队。挖掘地道的工作完全由它们负责承包下来。它们懂得，
必须选择适宜的地方来建造它们自己的住所。然后，如果它们无论是因为什么样的事情
发生，都会离开它们辛苦建筑起来的外部的小房间，那么，竹蜂就会紧随它们之后跑进
来，占据这一方难得的宝地。竹蜂就将利用很粗糙的土壁，把走廊分割成大小并不完全
相等的、毫无艺术特色的许多的小房间。这便是它们所能设计出的惟一的建筑构思了。
由此可见，它们是多么投机取巧，而且还很缺少艺术的灵感。
掘地蜂建造的窠巢，却做得非常整洁，而且还进行了非常精心别致的粉饰和装修。
所以，我们可以认为它们从事的工作是颇具艺术性的，它们自身具有高超的艺术创造才
能。它们很会利用适当的土壤，把窠巢构造得连任何一个普通的敌害都无法轻易地入侵。
也正是因为这个原因，这种蜜蜂的幼虫是不会做茧的。它们只是“赤身裸体”地躺在温
暖的小房间中享福，那里面光滑得如同粉饰过一样。
然而，竹蜂的小房间里却不一样了。那里是需要一定的东西来加以保护的。原因就
在于，竹蜂的窠巢是建筑在土壤的表面上的，做得非常的草率肤浅，而且只有相当薄的
墙壁做保垒。因此，和掘地蜂的幼虫不同，竹蜂的幼虫是包在非常坚固的厚厚的虫茧里
的。这样一来，有两点好处。一方面，厚厚的茧可以保护幼虫不致于和草率而建的巢里
的墙壁相碰撞而受到不必要的伤害。另一方面，也可以使得小幼虫能躲过闯进来的仇敌
的爪牙，不致于还在襁褓之中，就遭到不测而夭折离世。
在这样的堤上，居住着两种不同的蜜蜂。我们很容易就可以分辨出哪一种蜂巢属于
哪一种蜜蜂。很显然，在掘地蜂的窠巢里，隐藏着“一丝不挂”的赤裸小幼虫；而在竹
蜂的窠巢中，则有用坚实的茧包裹着的小幼虫。
同时，这两种不同的蜜蜂，都各自有它们特殊的寄生者，或是不速之客。竹蜂的寄
生者，是那种身上黑白相间的蝇。总是能够在蜂巢隧道的门口发现这种蝇。它们闯进到
窠巢中，然后产下一些自己的卵。掘地蜂的寄生者是蜂螨。我们经常可以在堤面上发现
很多这种甲虫的尸首。
如果我把竹蜂的小房间拿开，我便可以观察到掘地蜂的家了。在一些小房间中居住
着正在成长之中的昆虫。还有一部分小房间中，住满了掘地蜂的幼虫。也有一些小房间
中，大多数是藏着一个蛋形的壳。这种壳分成了好几节，上面还有突出来的呼吸孔。这
种壳特别的薄，而且还很脆，非常易碎。它的颜色是琥珀色的，非常透明。因此，从外
边看，可以很清楚地看到，里面有一个已经发育完全的蜂螨在挣扎着，好像极其渴望自
由，希望能早日从里面解放出来。
那么，这个很奇特的壳到底是个什么东西呢？看起来，它并不太像某一种甲虫的壳。
这个寄生者，是怎样来到这个蜂巢里面的呢？
从它的地理位置上来看，简直是不能侵入的。而且，若是使用放大镜进行仔细观察，
也看不出什么受过伤害的痕迹。经过了三年之久的周密而细致的观察，我终于寻找到了
这些问题的答案。于是，在我记录的昆虫的生活史上，便又增加了最为奇怪有趣的一页。
下面就是我所研究出来的一些结果。
蜂螨，即便是在它发育完整的时期内，也只不过有一两天的寿命而已，它的全部生
命，是在掘地蜂的门口度过的。而这短暂的生命，除去要繁殖子孙后代以外；其余的什
么也没有了。
蜂螨也具备其它的动物所有的消化器官，但是，它究竟要不要吃食物呢？我表示怀
疑。对于雌甲虫而言，它惟一的愿望，便是要产下它的小宝宝。等这件大事做完以后，
它便寿终正寝，放心地离开这个世界了。那么，雄性又怎么样呢？它们要在这种土穴上
伏上一两天之后，也同样命归九泉了。这也就为我们提供了一个问题的答案，为什么在
蜂的住宅旁边的那片蜘蛛网上，悬挂着那么多莫名其妙的尸首。现在可知道它们的来源
了。
人们乍一看，肯定都会以为这种甲虫在它产卵的时候，一定是要一个小房间一个小
房间地全都跑遍，在每一个蜜蜂的幼虫身上，都要产下一个卵。可是，事实并非如此，
在我观察的过程中，我曾经在蜜蜂的隧道里面仔仔细细地搜寻过，最后发现，蜂螨只将
所有的卵产在蜂巢的门口里边，积累成一堆，距离门口差不多有一到两寸远的地方。这
些卵全部都是白颜色的，其形状呈蛋形。它们的体积都很小，互相之间轻轻地粘连在一
起。至于它们到底有多少数目，暂时算它们一共有两千多个吧，我觉得这个数目还不能
算是过高的估计。
这一事实和任何人对于它必然要想象到的相反，它们并不把卵产在蜂巢的里面，而
仅仅将它们产在蜜蜂住宅的门口之内，而且还堆成一小堆。不只如此，它们的母亲也不
储备布置一些起保护作用的东西，也不考虑为它们防御冬天的寒冷，也不替它们关上这
扇进出孔道的大门，以便抵御前来侵略、攻击、伤害它们的成千上万的敌人。总之它们
的母亲在把它们产下之后，便抛开它们，让它们独自在这个世界上闯荡。这其中的原因
是：在冬日的严寒还不曾到来之前，这条开着口的隧道，便遭受着蜘蛛及其它更为凶悍
的侵略者们的践踏，被它们所侵占，因而那些可怜的卵也就成了侵犯者的可口美餐了。
为了能看得更清楚仔细一些，我把若干的卵放在一个盒子里面。大约到了九月，它
们还没有孵化出来的时候，我想象着，它们会立刻就跑开去，到处寻找掘地蜂的小房间。
然而，事实告诉我，我完全估计错了。这一群幼小的蛴螬——小小的黑色动物，还不到
一寸的二十五分之一长——虽然它们拥有强有力的健壮的腿，但竟然利用不上。它们并
不跑散开，而是非常混乱地相处在一起，和脱下来的卵壳混杂在一起生活。于是，我在
它们面前故意悄悄地放了一块带有蜂巢的土块，想看看它们会采取怎样的行动，可是结
果却无济于事，一点儿也不能诱惑这些小动物移动一丝一毫。要是我采取行动，强行把
其中的几个挪开一些，它们便会立即又跑回去，继续躲在其它的同伴里面，和它们混居
生活在一起。
最后，在冬天的时候，我跑到了卡本托拉斯的野外，到那里去观察掘地蜂居住的地
方。我想通过这样的观察来确定是否在自然的状态之下，蜂螨的幼虫们仍然也同样是如
此，即在孵化之后，不分散开居住，而是混杂地在一起生活。到底是不是一样的呢？是
的，完全一样。在野外的情况与我的小盒子里的情况真的一模一样。我看到那些在野外
的蜂螨的幼虫也同样是累积成一堆，并且也是和它们的卵壳混住在一起的。
到现在为止，我还不能回答这样的疑问：蜂螨究竟是怎么进到蜜蜂的小房间里面来
的呢？还有它们又是怎样走进另一种并不属于自己的壳里去的呢？

二、第一次的冒险
观察过幼小的蜂螨的外表以后，我便立刻就能感觉到，它们的生活习性一定是非常
特殊的，也一定挺有意思的。
经过仔细的观察，我发现很难使得蜂螨在很一般的平面上边轻轻移动一下。在蜂螨
的幼虫所居住的地方，很显然，它们要冒着一定的跌落下去的危险。怎样才能防备这种
危险的事情发生呢？这个问题对于蜂螨幼虫而言，是轻而易举地就能解决掉的。因为它
们天生就长着一对非常强有力的大腮，弯曲而且尖利；它们还生有强壮的腿，以及成为
能够活动的爪；还生长有很多的硬毛和尖尖的针；并且，它们生来就有一对坚硬的长钉，
有着锋利而且很坚硬的尖子，其形状和样子都很像一种犁头，它可以牢牢地刺入任何光
滑的土里。还不止这些呢。除了上面提到的这些器官以外，它还可以吐出一种粘性很强
的液汁，即便是没有其它的任何东西存在，单单是这种液汁，也是可以把它紧紧地粘住，
不致于滑下去的。可见，它的自我保护能力是很强的。
我曾经一再绞尽脑汁，具思苦想着一个问题，究竟存在什么样的因素，使得这些幼
小的蛴螬决定要居住在这里呢？可是，我怎么也想不出答案来。于是，我便只有非常急
切地等待着大自然的气候能够尽快转暖一些，以便能很快地找到答案。
到了四月底的时候，被我禁闭在牢笼中的蛴螬，以前本来一直是躺着不动的，躲避
在像海绵一样的卵壳堆里边睡觉。现在则不同了，它们忽然活动起来了。开始时，它们
在度过严冬的盒子里，到处爬走着。它们急急匆匆的动作，以及它们那不知疲倦的精气
神儿，都表明它们似乎正在寻觅一些东西一些它们急切需要的东西。看起来，这些东西
自然是它们的食物了。因为，这些蜂螨的幼虫们是在九月底进行孵化的，一直到现在，
四月底。虽然它们总是处在麻木不仁的状态之中，但是差不多足足有七个月的时间，没
有获取一点儿有营养的东西来强壮身体。从孵化的时候开始，虽然这些小动物是具有生
命的“精灵”，但是它们就像是注定被判了七个月的徒刑一样，什么事情也不能做，只
能保持着一种姿势。
与此同时，当我看到它们一个个如此兴奋充满激情的时候，我便自然而然地猜想到，
驱使这些有生命的小动物如此忙忙碌碌地工作的原因，一定是饥饿，也只有饥饿才能让
它们本能地不辞辛苦。
这些匆忙寻找食物的小动物，它们真正需要的食物不过是蜂巢中的储藏品罢了。有
什么理由可以这样说呢？因为到了后期的时候，我们是在这些蜂巢中找到那些蜂螨的。
现在这些储藏品不仅限于蜜蜂的幼虫食用，也供蜂螨们分享了。
我所提供给它们的，是里面藏着蜜蜂幼虫的蜂巢。我甚至把蜂螨直接地放到蜂巢里
边去。总之，我利用各种东西，采用各种方法，希望能引起它们的食欲。但是，事实上，
我的努力仍然是一点儿结果也没有。于是，我故意用了一种方法，利用蜂蜜进行试探。
为了能够找到储藏着蜜汁的蜂巢，我花去了五月份的大部分时间。
在找到了我所希望的蜂巢以后，我把其中蜜蜂的幼虫拿了出去，然后再把蜂螨的幼
虫放到蜂巢中储备的蜂蜜里。然而事实再一次打击了我。简直没有任何其它的试验比这
一次失败得更为严重了。幼虫们根本就不去饮食那些蜜汁，更糟糕的是，它们反而被这
种粘性的东西粘住了，以至于在其中被闷死了，这太出乎我的意料了！
于是，我非常失望地说道：“我提供给你们现成的蜂窠，幼虫，还有蜜汁，难道这
些还不足够吗！你们还不满足吗？那么，你们这些丑恶的小东西，到底需要些什么东西
呢？”
不过，最终还是让我发现了它们真正需要的东西了。原来，它们并不需要什么特殊
的东西。它们是要掘地蜂亲自把它们带到蜂巢里边去。
我在前面早就已经提到过，当四月来临的时候，在蜂巢的门内居住的一堆幼虫，已
经开始表现出一点儿活动的迹象了，它们蠢蠢欲动。仅仅几天以后，它们便已经不在那
个地方停留了。真是非常怪异的小动物。它们牢牢地，死不放手地攀附在蜜蜂的毛上，
于是，便被带到了野外去，甚至已经被带到很遥远的地方了。
当掘地蜂经过蜂巢的门口的时候，无论它是要出远门，还是刚从远游中归来，睡在
门口，已经等待许久的蜂螨的幼虫，便会立刻爬到蜜蜂的身上去。它们爬进掘地蜂的绒
毛里面，抓得十分紧，无论这只掘地蜂要飞到多么遥远的地方去，它们一点儿也不担心
自己有跌落到地上去的危险。因为它们抓得太紧了。之所以要采用这样的方法，它们惟
一的目的就是想借助蜜蜂强壮的身体，将它们带到那些储有丰富的蜜的巢里去。
当一个人第一次发现这种情形的时候，他一定会以为这种喜欢冒风险的小幼虫，可
能要在蜜蜂的身上先寻觅到一些食物。但是，事实并非如此，蜂螨的幼虫伏卧在蜂的细
细的绒毛里面，和蜜蜂的身体是呈直角状的。它的头朝向里面，尾巴却朝向外面，呆在
靠近蜜蜂肩头的地方。它们选择好了地点以后，便不再随便地移来移去。如果，它们真
的打算在蜜蜂的身上吃些什么东西的话，那么它们应该是这里那里地到处跑动，寻找一
下哪一部分是最鲜嫩的。然而，事实并不是这样。它们总是要固着在蜜蜂身上，并且是
最硬的那部分，一般是在靠近蜜蜂翅膀下面的部位，有时也附着在蜜蜂的头上。它们攀
住一根毛以后，纹丝不动。因此，在我看来，事实总归是事实，是不以人的意志为转移
的。这些小甲虫之所以如此附着在蜜蜂的身体上，它们的目的仅仅是打算让蜜蜂把它们
带到即将要建造起来的蜂巢里去。
不过，在飞行的时候，这位未来的寄生者，必须要紧紧抓牢它的主人的毛才行。无
论蜜蜂是正在花叶中穿梭飞行时，如何地急速，还是它在向窠巢里飞的时候，如何地摩
擦，甚至无论它是正在用足清洁它的身体的时候，小幼虫都必须抓得很紧才行，这样才
能确保安全。
不久以前，我们曾经怀疑过，究竟是什么东西，可以使蜂螨的幼虫依附在蜜蜂的身
上呢？现在已经知道答案了，那便是生长在蜜蜂身上的绒毛。
现在，我们可以知道长在蜂螨身上的那两根大钉有什么用途了。它们合拢起来，便
可以紧紧握住蜜蜂身上的毛，比起那些最精密的人工的钳子来，还要更加精密得多。
同时，我们也可以知道那些粘液的价值了。它能帮助这个小动物更加牢固地伏在蜜
蜂的身上。而且，我们也可以了解幼虫足上长着的尖针和硬毛的作用了。它们都是用来
插入蜜蜂的软毛里，使得它本身的地位更加稳固的。
我们越是想这些当幼虫爬在平面上时似乎毫无用武之地的“组成设备”，越是会感
觉到对于这些“机器”的惊异。当这个柔弱的小动物，在它冒着危险去周游大干世界的
时候，竟然能够利用如此多的器具，防止它从蜜蜂的身上跌落下来，是多么的奇妙啊！

三、第二次的冒险
在五月二十一日这一天，我到卡本托拉斯去，想看一看蜂螨进入蜂巢时的门路。
这件工作很不容易做，需要用尽全力去完成。在野外广大的地面上，有一群蜜蜂，
像是受了日光的刺激似的，正在那边疯狂乱舞着。就在我正用缭乱的眼光观察它们的动
作时，忽然在狂乱的蜂群中间响起了一种单调而可伯的喧哗声。然后，就像闪电一样迅
猛，掘地蜂飞身而起，到处去寻找食物。与此同时，另外一群成千的蜜蜂正飞回家来。
它们身上或是携带着已经采好了的蜜汁，或者带回了用来建造它们的蜂巢用的泥土。
在那个时候，我已经具备了一些关于这类昆虫的知识，了解了一些有关它们的习性。
我以为，无论是谁有意无意地闯入它们的群里，或者只是轻轻地碰一碰它们的住宅，那
么它马上就会遭到成千锥子的狂刺而身亡。有一次，我去观察大黄蜂的蜂房，由于距离
太近了，立刻就起了一阵恐惧的颤抖，那种感觉真的是一辈子也忘不了！
然而，不管如何困难，要想知道我所渴望知道的事情，我就必须进入到这种可怕的
蜂群里，而且，必须在那里站上几个钟头，必要的时候，甚至是一整天的时间。我必须
盯着它们的工作，目不转睛，把放大镜拿在手上，站在它们当中一动也不动，观察着蜂
巢里将会发生什么样的事情。与此同时，面套、手套等等，以及其它各类遮盖保护的东
西全都不能使用，其原因是，我的手指和眼睛一定要完全不受任何防碍。其余的一切都
不管，即便是我离开蜂巢时，脸上被刺得都快让人认不出来了，也不能带各种遮盖的东
西。
那一天，我决定要解决那个困扰了我很长时间的问题了。
我用我的网子捉住了几只掘地蜂，这竟然让我感到十分的满意。因为这几只蜂的身
上都栖息着蜂螨的幼虫，这也正如我一直所希望的那样。
我先把衣服扣紧，然后，突入这群蜜蜂的中心。我拿了锄头锄了几下，然后取下一
块泥。令我感到非常诧异的，就是我居然一点儿也没有受到攻击和伤害。
在我第二次开始的时候，花的时间比第一次还要更长一些，但是，仍然是同样的结
果。我并没有受一点儿伤，也并没有一个蜜蜂利用它的尖针来刺击我一下。这以后，我
也就没什么可担心害怕的了。于是，我就大胆地长时间地停留在蜂巢前面，揭起土块，
拿出里面的蜂蜜，赶走其中的蜜蜂。在这一过程中，始终没有引起比那种喧哗更为可怕
的事情发生。为什么呢？这主要是因为掘地蜂是一种比较爱好和平的动物。每当它们的
巢内部被扰乱的时候，它们便会马上离开老家，转移地点，躲避到其它的地方去。即便
是有的时候受了一点儿伤，它们也不会使用它们的尖针，只有当它们被人捉住的时候，
它们才会用一下。
不过我得向这个泥水匠蜂的缺乏勇气表示谢意。虽然我并没有进行一点儿防御，但
是，居然能够在这些喧闹的蜂群中，在一块石头上，安安静静地坐着，并且还能随意地
观察它们的巢达到几个小时之久，却并没有被刺击过一针。这时，有一些乡下人经过这
个地方，看见我居然很安静地坐在蜂群之中。于是，他们便问我是不是对它们施加了什
么魔法。
就是这样，我观察了很多的蜜蜂蜂巢。其中有些蜂巢还是敞开着的，里面多少储备
着一些蜜汁。还有一些蜂巢已经用土掩盖了起来。而里面的东西，是大不一样的。有的
时候，我看到的是蜜蜂的幼虫；有的时候，我又看到其它种类的稍为肥大的幼虫；另外
一些时候，还能看见一个卵飘浮在蜜汁的表面上。这个卵呈非常美丽的白颜色，它是圆
柱形的，而且稍微有一点儿弯曲，差不多有一寸的五分之二或者六分之一长，这就是掘
地蜂的卵。
在少数的小房间中，我看到这种虫卵浮在蜂蜜的表面上，在其它的许多小房间中，
我看到的更多的就是幼小的蜂螨的蛴螬，它们伏卧在蜜蜂的卵上，就好像是伏在一种木
筏上一样。它的形状和大小都和刚刚孵化出来的时候是一样的，在这个蜜蜂巢里，敌人
已经卧在家门口了。
它是在什么时候并且是用什么方法进去的呢？在很多小房间中，经过仔细观察，我
简直找不出一点儿它们可以闯进去的缝隙。因为这些小房间全都封闭得很严密；所以我
推测这位寄生者一定是在储藏蜂蜜的仓库还没有关上大门之前就进去了。可是另一方面，
我看到门没有关上的小房间里面藏满了蜂蜜，但是没有发现有卵浮在那上面，也从没有
蜂螨的幼虫在里面留宿。所以，这些幼虫一定是在蜜蜂产卵的时候，或者是后来蜜蜂封
门的时候进去的。我凭借我的一点儿经验断定，幼虫在进入小房间的时候，一定是在蜜
蜂产卵在蜜上时的一瞬间。
如果我拿了一个里面装满了蜂蜜，表面上还浮着一个卵的小房间来；然后再拿上几
只蜂螨的幼虫，把它们一起放到玻璃罩里面进行观察。然而它们却很少会跑到蜂巢里边
去，它们也不能够安然地跑到“木筏”上边去！围绕着这个“木筏”的蜂蜜看来对它们
而言是太危险了。即使有那么一两只幼虫碰巧跑近了这个蜜湖，那么它们一看到这粘性
很大的东西，或者稍一涉足其中，便马上会千方百计地设法逃离这个危险的地方。可是，
经常有一些不太走运的幼虫，一不小心，就跌落到蜂的窠巢里面，很不幸运地被闷死。
所以，我们根据这一点便可以断定，蜂螨的幼虫是绝对不会离开蜜蜂的毛的，特别是在
蜜蜂呆在小房间里或靠近小房间的时候，更要牢牢地依附于蜜蜂的身体。这是因为，只
要和蜂蜜的表面稍有一点儿接触，就可以使得柔弱的小幼虫窒息而死亡。
我们必须记清楚这一点，那就是幼小的蜂螨是在封闭的小房间中发现的，而且它们
一定是呆在蜜蜂的卵上面的。这个小小的卵，不仅可以给这个小小的动物当做一个安全
的木筏，以便它能漂浮在这个可怕的、不可信任的蜜湖中，而且，这个小小的卵也就成
为了幼虫的第一顿美味佳看。
然而，要想到达这只漂浮在蜜湖中心的，而且又将成为它的食品的木筏，这只蜂螨
的幼小的蛴螬必须要避免与蜂蜜接触。否则的话后果会不堪设想的。要想达到目的，做
成这件事情，只有一个可以选择的方法。这个聪明的小幼虫，趁着蜜蜂还在产卵的空当
儿，从它的身上迅速地一下子滑落到了那个卵上。这样一来，目的达到了。幼虫便和卵
一起作伴，共同浮在蜜上了。因为这只由蜜蜂产下的卵太小了，不能同时乘载超过一个
以上的幼虫。因此，我们在一个蜂室里面，只能看到一个蜂螨的幼虫。
蜂螨幼虫的这种动作，在我们人类看来，好像是异常有灵性似的。但是，如果我们
再继续研究昆虫的话，那么，它们还将为我们提供更多的这样有灵感的例子。
可以说，当蜜蜂产下卵以后，把卵放在蜜汁上时，同时也就把它们的小天敌——蜂
螨的幼虫一起放到了小房间里面，然后，蜜蜂要非常仔细地用土把小房间的门给密封起
来。于是，一切需要它做的工作都做完了。然后，第二个小房间是做在第一个小房间的
旁边的。大概也要经历和前面相同的过程。就照这样，不停地继续下去。一直要继续到
隐蔽在蜜蜂的绒毛中的寄生者统统安居下来，这才算结束。
现在，让我们抛开这个实际上有些苦恼的母亲，不管它在做着怎样无结果的工作，
把我们的注意力转移一下，转到这些用聪明的方法最终得到膳宿的蜂螨的幼虫身上，看
看它对我们的试验会做出什么样的反应。
让我们想一想看，如果我们一旦将一只有蜂螨的幼虫在小房间上面的盖子拿下来，
那么会有什么样的事情将要发生呢？
卵还是十分完好的，一点儿也没有受到破坏。但是，好景不长。不久以后，蜂螨幼
虫的破坏工作便开始了。我们可以观察到，幼虫朝着一个长有小黑点的白卵上跑去。最
后，它忽然停下来了，由于它有六只脚，所以身体可以停得很稳。然后，它利用长在自
己身上的大腮的尖钩咬住了那个卵身上的薄皮儿，用尽自己浑身的力气，猛烈地拉扯着，
直到那个卵被它拉破为止。于是，卵里面的东西便都流了出来。那只得胜的幼虫，见了
这种东西非常满意，立刻高兴地把它享用光。这个小小的寄生虫，一生第一次使用它的
大腿的尖钩，原来是在拉破蜂卵的时候。
蜂螨的幼虫真是天生聪明啊！想出这样奇妙的方法来。利用了这种巧妙的方法，幼
虫便可以在它寄生的蜂巢的小房间中为所欲为，毫无顾及了。它可以任意地享用蜜汁。
这是因为，在蜜蜂的幼虫在孵化过程中，也是需要蜜汁来增加营养的但是，在孵化中，
蜂卵吸收的一点点东西，是不能在日后提供给两者一起享用的。因此只要蜂螨的幼虫在
拉扯卵皮的时候，越快越用力就越好。这样一来，“僧多粥少”的这个困难就不存在了。
蜂螨的幼虫之所以要破坏蜜蜂的卵，还有另外一个重要的原因。那就是因为蜂卵有
一种特殊的滋味。这种很吸引蜂螨幼虫的滋味，驱使着小幼虫在它的第一餐就要享用这
个香甜可口的小卵。这个小幼虫，在把卵撕破的初期，饮食到的是从卵里流出来的诱人
的浆汁。接连过了好几天以后，幼虫继续加油努力，把卵的裂口撕扯得更大一些这样一
来，幼虫就可以得寸进尺了。它将继续享受卵内部的流质，直到满足为止。
在幼虫吸食蜜蜂卵的过程中，储备在蜂卵周围的甜美的蜜汁，却一点儿也诱惑不了
贪吃的蜂螨幼虫，它理都不理睬一下，也不去碰它们一下。因此，可以这样讲，蜜蜂的
卵对于蜂螨幼虫而言，是绝对重要的，它是幼虫的必需食品。因而小小的蜂卵，不仅仅
可以当作蜂螨幼虫的一叶小舟，使得它在蜜湖中安全地行进，更重要的是，它还是幼虫
相当有营养的食品，为幼虫的茁壮成长提供条件。
等到整整一个星期以后，这个可怜的小蜂卵已不再剩下什么东西了，只留下了一个
空空无物的干壳。一个生命就这样悄悄地结束了。而到了这个时候，蜂螨幼虫的第一顿
大餐也已经享用完毕。幼虫也茁壮成长，差不多有原来的两倍大了。它的形状也发生了
一些不大不小的变化。它的背部裂开了，形成了自己的第二种形状，长成了一只简单的
甲虫。小幼虫从那个裂缝中解脱出来，然后落到蜂蜜上。从它身上脱下来的那个壳，还
依然停留在原来的那个小“木筏”上面。但是，在不久以后，它们都被掩没在蜜浪之中
了。
此时此刻，蜂螨幼虫的历史便画上了一个圆满的句号！
------------------








蟋蟀



一、家政
居住在草地里的蟋蟀，差不多和蝉是一样有名气的。它们在有数的几种模范式的昆
虫中，表现是相当不错的。它之所以如此名声在外，主要是因为它的住所，还有它出色
的歌唱才华。只占有这其中的一项，是不足以让它们成就如此大的名气的。一位动物故
事学家拉封丹，对于它只谈了简单的几句，仿佛并没有注意到这种小动物的天才与名气。
另外，还有一位法国寓言作家曾经写过一篇关于蟋蟀的寓言故事，但是很可惜，太
缺乏真实性和含蓄一些的幽默感。而且，这位寓言作家在这个蟋蟀的故事中写到：蟋蟀
并不满意，在叹息它自己的命运！事实可以证明，这是一个多么错误的观点。因为，无
论是什么样的人，只要曾经亲自研究过蟋蟀，观察过它们的生活情况，哪怕仅仅是一点
表面上的观察与研究，都会感觉到蟋蟀对于自己的住所，以及它们天生的歌唱才能，是
非常满意而又愉快的。是的，这两点所给它们带来的名气真的足以让它们感到庆幸了。
在这个故事的结尾处，他承认了蟋蟀的这种满足感。他写到：

“我的舒适的小家庭，是个快乐的地方，
如果你想要快乐的生活，就隐居在这里面吧！”

在我的一位朋友所做的一首诗中，给了我另一种感觉。我觉得这首诗所要表达的更
具有真实性，更加有力地表现出蟋蟀对于生活的热爱。
下面就是我的朋友写的这首诗：

曾经有个故事是讲述动物的，
一只可怜的蟋蟀跑出来，
到它的门边，
在金黄色的阳光下取暖，
看见了一只趾高气扬的蝴蝶儿。

她飞舞着，
后面拖着那骄傲的尾巴，
半月形的兰色花纹，
轻轻快快地排成长列，
深黄的星点与黑色的长带，
骄傲的飞行者轻轻地拂过。

隐士说道：飞走吧，
整天到你们的花里去徘徊吧，
不论菊花白，
玫瑰红，
都不足与我低凹的家庭相比。

突然，
来了一阵风暴，
雨水擒住了飞行者，
她的破碎的丝绒衣服上染上了污点儿，
她的翅膀被涂满了烂泥。

蟋蟀藏匿着，
淋不到雨，
用冷静的眼睛看着，
发出歌声。
风暴的威严对于它毫不相关，
狂风暴雨从它的身边无碍地过去。

远离这世界吧！
不要过分享受它的快乐与繁华，
一个低凹的家庭，
安逸而宁静，
至少可以给你以不须忧虑的时光。

从这首诗里，我们可以认识一下可爱的蟋蟀了。
我经常可以在蟋蟀住宅的门口看到它们正在卷动着它们的触须，以便使它们的身体
的前面能够凉快一些，后面能更加暖和一些。它们一点儿也不妒嫉那些在空中翩翩起舞
的各种各样的花蝴蝶。相反的，蟋蟀反倒有些怜惜它们了。它们的那种怜悯的态度，就
好像我们常看到的一样，那种有家庭的人，能体会到有家的欢乐的人，每当讲到那些无
家可归，孤苦伶仃的人时，都会流露出一样的怜悯之情。蟋蟀也从来不诉苦、不悲观，
它一向是很乐观的、很积极向上的，它对于自己拥有的房屋，以及它的那把简单的小提
琴，都相当的满意和欣慰。从某种意义上，可以这样说，蟋蟀是个地道正宗的哲学家。
它似乎清楚地懂得世间万事的虚无漂渺，并且还能够感觉到那种躲避开盲目地、疯狂地
追求快乐的人的扰乱的好处。
对了，这样来描写我们的蟋蟀，无论如何，总应该是正确的。不过，仍然需要用几
行文字，以便把蟋蟀的优点公之于众。自从那个动物故事学家拉封丹，忽略了它们以后，
蟋蟀已经等待了很长的时间了，等待着人们对它加以描述，加以介绍，加以重视。它们
的朋友——人类忽略了它们。
对于我，一个自然学者而言，前面提到的两篇寓言中，最为重要的一点，乃是蟋蟀
的窠穴，教训便建筑这在上面。
寓言作家在诗中谈到了蟋蟀的舒适的隐居地点；而拉封丹，也赞美了它的在他看来
是低下的家庭。所以，从这一点讲，最能引起人们注意的，毫无疑问，就是蟋蟀的住宅。
它的住宅，甚至吸引了诗人的目光来观察它们，尽管他们常常很少能做注意到真正存在
的事物。
确实，在建造窠穴以及家庭方面，蟋蟀可以算是超群出众的了。在各种各样的昆虫
之中，只有蟋蟀在长大之后，拥有固定的家庭，这也算是它辛苦工作的一种报酬吧！在
一年之中最坏的时节，大多数其它种类的昆虫，都只是在一个临时的隐避所里暂且躲避
身形，躲避自然界的风风雨雨。因此，它们的隐避场所得来的方便，在放弃它的时候，
也并不会觉得可惜。
这些昆虫在很多时候，也会制造出一些让人感到惊奇的东西，以便安置它们自己的
家。比如，棉花袋子，用各种树叶制做而成的篮子，还有那种水泥制成的塔等等。有很
多的昆虫，它们长期在埋伏地点伏着，等待着时机，以捕获自己等待已久的猎物。例如，
虎甲虫。它常常挖掘出一个垂直的洞，然后，利用它自己平坦的、青铜颜色的小脑袋，
塞住它的洞口。如果一旦有其它种类的昆虫涉足到这个具有迷惑性的、诱捕它们的大门
上时，那么，虎甲虫就会立刻行动，毫不留情地掀起门的一面来捕捉它。于是，这位很
不走运的过客，就这样落入虎甲虫精心伪装起来的陷阱里，不见踪影了。
另外一个例子，是蚁狮。它会在沙子上面，做成一个倾斜的隧道。这里的牺牲者是
蚂蚁。蚂蚁一旦误入歧途，便会从这个斜坡上不由自主地滑下去，然后，马上就会被一
阵乱石击死。这条隧道中守候猎物的猎者，把颈部做成了一种石弩。
但是，上面提到的例子统统都只是一种临时性的避难所或是陷阱而已，实在不是什
么长久之计。
经过辛辛苦苦的劳作构造出的家，昆虫住在里面，无论是朝气蓬勃、生机盎然的春
天，或者是在寒风刺骨、漫天雪飘的冬令时节，都让昆虫无比地依赖，不想迁移到其它
的任何地方去居住。这样一个真正的居住之所，是为了安全以及舒适而建筑的，是从长
远的角度考虑的，而并不是像前面所提到的那样，家是为了狩猎而建的。或是所谓的
“育儿院”之类的延期行为。那么，只有蟋蟀的家是为了安全和温馨而建造的了。在一
些有阳光的草坡上，蟋蟀就是这个隐逸者的场院的所有者。正当其它的或许正在过着孤
独流浪的生活，或许是卧在露天地里，或许是埋伏在枯树叶、石头和老树的树皮底下的
昆虫正为没有一个稳定的家庭而烦恼时，蟋蟀却成了大自然中的一个拥有固定居所的优
越的居民。由此可见，它是有远见意识的。
要想做成一个稳固的住宅，并不那么简单不过，现在对于蟋蟀、兔子，最后是人类，
已经不再是什么大问题了。在与我的住地相距不太远的地方，有狐狸和獾猪的洞穴，它
们绝大部分只是由不太整齐的岩石构建而成的，而且一看就知道这些洞穴都很少被修整
过。对于这类动物而言，只要能有个洞，暂且偷生，“寒窑虽破能避风雨”也就可以了。
相比之下，兔子要比它们更聪明一些。如果，有些地方没有任何天然的洞穴可以供兔子
们居住，以便躲避外界所有的侵袭与烦扰，那么，它们就会到处寻找自己喜欢的地点进
行挖掘。
然而，蟋蟀则要比它们中的任何一位更聪明得多。在选择住所时，它常常轻视那些
偶然碰到的以天然的隐避场所为家的种类。它总是非常慎重地为自己选择一个最佳的家
庭住址。它们很愿意挑选那些排水条件优良，并且有充足而温暖的阳光照射的地方。凡
是这样的地方。都被视为佳地，要优先考虑选取。蟋蟀宁可放弃那种现成的天然而成的
洞穴，因为，这些洞都不合适，而且它们都建造得十分草率，没有安全保障。有时，其
它条件也很差。总之这种洞不是首选对象。蟋蟀要求自己的别墅每一点都必须是自己亲
手挖掘而成的，从它的大厅一直到卧室，无一例外。
除去人类以外，至今我还没有发现哪种动物的建筑技术要比蟋蟀更加高超。即便是
人类，在混合沙石与灰泥使之凝固，以及用粘土涂抹墙壁的方法尚未发明之前，也不过
是以岩洞为隐避场所，和野兽进行战斗和大自然进行博击。那么，为什么这样一种非常
特殊的本能，大自然单单把它赋予了这种动物呢？最为低下的动物，但却可以居住得非
常完美和舒适。它拥有自己的一个家，有很多被文明的人类所不知晓的优点：它拥有安
全可靠的躺避隐藏的场所；它有享受不尽的舒适感，同时，在属于它自己的家的附近地
区，谁都不可能居住下来，成为它们的邻居。除了我们人类以外，没有谁可以与蟋蟀相
比。
令人感到不解和迷惑的是，这样一种小动物，它怎么会拥有这样的才能呢？难道说，
大自然偏向它们，赐予了它们某种特别的工具吗？当然，答案是否定的。蟋蟀，它可不
是什么掘凿技术方面的一流专家。实际上，人们也仅仅是因为看到蟋蟀工作时的工具非
常柔弱，所以才对蟋蟀有这样的工作结果，建造出这样的住宅感到十分惊奇的。
那么，是不是因为蟋蟀的皮肤过于柔嫩，经不起风雨的考验，才需要这样一个稳固
的住宅呢？答案仍然是否定的。因为，在它的同类兄弟姐妹中，也有和它一样，有柔美
的、感觉十分灵敏的皮肤，但是，它们并不害怕在露天底下呆着，并不怕暴露于大自然
之中。
那么，它建筑它那平安舒适的住所的高超才能，是不是由于它的身体结构上的原因
呢？它到底有没有进行这项工作的特殊器官呢？答案又是否定的。在我住所的附近地区，
分别生活着三种不同的蟋蟀。这三种蟋蟀，无论是外表、颜色，还是身体的构造，和一
般田野里的蟋蟀是非常相像的。在开始时，刚一看到它们，经常就把它们当成田野中的
蟋蟀。然而，就是这些由一个模子刻出来的同类，竟然没有一个晓得究竟怎样才能为自
己挖掘一个安全的住所。其中，有一只身上长有斑点的蟋蟀，它只是把家安置在潮湿地
方的草堆里边；还有一只十分孤独的蟋蟀，它自个儿在园丁们翻土时弄起的土块上，寂
寞地跳来跳去，像一个流浪汉一样；而更有甚者，如波尔多蟋蟀，甚至毫无顾忌，毫不
恐惧地闯到了我们的屋子里来，真是不请自来的客人，不顾主人的意愿。从八月份到九
月份，它独自呆在那些既昏暗又特别寒冷的地方，小心翼翼地唱着歌。
如果再继续前面已经提到过的那些问题，将是毫无意义的。因为那些问题的答案统
统都是否定的。蟋蟀自然形成的本能，从来也不为我们提供有关答案的原因所在。如果
寄希望于从蟋蟀的体态，身体结构，或是工作时所利用的工具上来寻找答案，来解释那
些答案，同样是不可能的。长在昆虫身上的所有的东西，没有什么能够提供给我们一些
满意的解释与答案，或者是能够让我们知晓一些原因，给不了我们任何有力的帮助。
在这四种相互类似的蟋蟀中，只有一种能够挖掘洞穴。于是，我们可以得知，蟋蟀
本能的由来，我们尚不可得知。
难道会有谁不晓得蟋蟀的家吗？哪一个人在他还是小孩子的时候，没有到过这位隐
士的房屋之前去观察过呢？无论你是怎样的小心，脚步是如何地轻巧，这个小小的动物
总能发觉，总能感觉到你的来访。然后，它立刻警觉起来，并且有所反映，马上躲到更
加隐避的地方去。而当你好不容易才接近这些动物的定居地时，此时此刻，这座住宅的
门前已经是空空如也了，很让人失望。
我想，凡是有过如此经历的人，谁都会知道，如何把这些隐匿者从躲藏处诱惑出来
的方法。你可以拿起一根草，把它放到蟋蟀的洞穴里去，轻轻地转动几下。这样一来，
小蟋蟀肯定会认为地面上发生了什么事情。于是，这只已经被搔痒了，而且已经有些恼
怒了的蟋蟀，将从后面的房间跑上来。然后，停留在过道中，迟疑着，同时，鼓动着它
的细细的触须认真而警觉地打探着外面的一切动静。然后，它才渐渐地跑到有亮光的地
方来，只要这个小东西一跑到外面来，便是自投罗网，很容易就会被人捉到。因为，前
面发生的一系列事情，已经把我们这只可怜的小动物的简单的小小头脑给弄迷糊了，毕
竞它的智力水平是何等低下啊！假如这一次，小蟋蟀逃脱掉了，那么，它将会很疑虑，
很机警，时刻提高它的警惕性，不肯再轻易地冒险，从躲避的地方跑出来。在这种情况
下，就不得不选择其它的应付手段了。比如，你可以利用一杯水，把蟋蟀从洞穴中冲出
来。
想起我们的孩童时代，那个时候真的是值得人怀念与羡慕。我们跑到草地里去，到
处捉拿蟋蟀这种昆虫。捉到以后，就把它们带回家里。就把它们放在笼子里供养。采来
一些新鲜的莴苣叶子来养活它们。这真是一种莫大的童趣啊！
现在，回过头来谈谈我这里的情况吧。为了能够更好地研究它们，我到处搜寻着它
们的窠穴。孩童时代发生的事情，就仿佛昨天刚刚发生过一般。当我的另一个小同伴—
—小保罗——一个在利用草须方面，可以称为专家级的孩子，在很长时间地实施他的战
略战术之后，忽然，他十分激动而兴奋地叫起来：“我捉住它了！我捉住它了！一只可
爱的小蟋蟀！”
“动作快一点儿，”我对小保罗说道，“我这里有一个袋子。我的小战俘，你快快
跳进去吧，你可以在袋子里面安心居住。里面有充足的饮食。不过，有个条件，那就是，
你可一定不要让我们失望啊！你一定要赶快告诉我们一些事情，一些我们渴望知道而且
正在苦苦寻觅的答案。而这些事情中，需要你做的头一件便是：把你的家给我看一看。”

二、它的住屋
在那些青青的草丛之中，不注意的话，就会不为人知地隐藏着一个有一定倾斜度的
隧道。在这里，即便是下了一场滂沱的暴雨，也会立刻就干了的。这个隐蔽的隧道，最
多不过有九寸深的样子，宽度也就像人的一个手指头那样。隧道按照地形的情况和性质，
或是弯曲，或是垂直。差不多如同定律一样，总是要有一叶草把这间住屋半遮掩起来，
其作用是很明显的，如同一所罩壁一样，把进出洞穴的孔道遮避在黑暗之中。蟋蟀在出
来吃周围的青草的时候，决不会去碰一下这一片草。那微斜的门口，仔细用扫帚打扫干
净，收拾得很宽敝。这里就是它们的一座平台，每当四周的事物都很宁静的时候，蟋蟀
就会悠闲自在地聚集在这里，开始弹奏它的四弦提琴了。多么温馨的促夏消暑音乐啊！
屋子的内部并不奢华，有暴露的，但是并不粗糙的墙。房子的住户有很多空闲的时
间去修整太粗糙的地方。隧道的底部就是卧室，这里比别的地方修饰的略微精细些，并
且宽敞些。大体上说，这是个很简单的住所，非常清洁，也不潮湿，一切都符合卫生标
准。从另一方面来说，假如我们考虑到蟋蟀用来掘土的工具十分简单，那么可以说这真
是一个伟大的工程了。如果想要知道它是怎样做的，它是什么时候开始这么大的工程的，
我们一定要回溯到蟋蟀刚刚下卵的时候。
蟋蟀像黑螽斯一样，只把卵产在土里，深约四分之三寸，它把它们排列成群，总数
大约有五百到六百个。这卵真是一种惊人的机器。孵化以后，它看起来很像一只灰白色
的长瓶子，瓶顶上有一个因而整齐的孔。孔边上有一顶小帽子，像一个盖子一样。去掉
盖子的原因，并不是蛴螬在里面不停地冲撞，把盖子弄破了，而是因为有一种环绕着的
线——一种抵抗力很弱的线，它自己会自动裂开。
卵产下两个星期以后，前端出现两个大的蛴螬，是一个呆在襁褓中的蛴螬，穿着紧
紧的衣服，还不能完全辨别出来。你应当记得，螽斯也以同样的方法孵化，当它来到地
面上时，也一样穿着一件保护身体的紧紧的外衣。蟋蟀和螽斯是同类动物，虽然事实上
并不需要，但它也穿着一件同样的制服。螽斯的卵留在地下有八个月之久，它要想从地
底下出来必须同已经变硬了的土壤搏斗一番，因此需要一件长衣保护它的长腿。但是蟋
蟀整体上比较短粗，而且卵在地下也不过几天，它出来时无非只要穿过粉状的泥土就可
以了。用不着和土地相抗挣。因为这些理由，它不需要外衣，于是它就把这件外衣抛弃
在后面的壳里了。
当它脱去襁褓时，蟋蟀的身体差不多完全是灰白色的，它开始和眼前的泥土战斗了。
它用它的大腮将一些毫无抵抗力的泥土咬出来，然后把它们打扫在一旁或干脆踢到后面
去，它很快就可以在土面上享受着阳光，并冒着和它的同类相冲突的危险开始生活，它
是这样弱小的一个可怜虫，还没有跳蚤大呢！
二十四小时以后，它变成了一个小黑虫，这时它的黑檀色足以和发育完全的蟋蟀相
媲美，它全部的灰白色到最后只留下来一条围绕着胸部的白肩带，它身上生有两个黑色
的点。在这两点中上面的一点，就在长瓶的头上，你可以看见一条环绕着的，薄薄的、
突起的线。壳子将来就在这条线上裂开。因为卵是透明的，我们可以看见这个小动物身
上长着的节。现在是应该注意的时候了，特别是在早上的时候。
好运气是关爱带来的，如果我们不断地到卵旁边去看，我们会得到报酬的。在突起
的线的四周，壳的抵抗力会渐渐消失，卵的一端逐渐分裂开，被里面的小动物的头部锥
动，它升起来，落在一旁，像小香水瓶的盖子一样，战俘就从瓶子里跳了出来。
当它出去以后，卵壳还是长形的，光滑、完整、洁白，馆子似的盖子挂在口上的一
端。鸡卵破裂，就是小鸡用嘴尖上的小硬瘤撞破的；蟋蟀的卵做得更加巧妙，和象牙盒
子相似，能把盖子打开。它的头顶，已经足可以做这件工作了。
我们上面说过，盖子去掉以后，一个幼小的蟋蟀跳出来，这句话还不十分精确。它
是非常灵敏和活泼的，不时用长的而且经常颤动的触须打探四周发生的情况，并且很性
急地跑来跳去。当有一天，它长胖了，不能如此放肆了，那才真有些滑稽呢！
现在我们要看一看母蟋蟀为什么要产下这么多的卵。这是因为多数的小动物是要被
处以死刑的。它们常遭到别的动物大量的残忍的大屠杀，特别是小形的灰螅蜴和蚂蚁的
杀害。蚂蚁这种讨厌的流寇，常常不留一只蟋蟀在我们的花园里。它一口就能咬住这可
怜的小动物，然后狼吞虎咽地将它们吞咽下去。
唉，这个可恨的恶人，请想想看，我们还将蚂蚁放在比较高级的昆虫当中，还为它
写了很多的书，更对它大加赞美。称赞之声，不绝于耳。自然学者对它们很推祟，而且
其名誉日益增加。这样看来，动物和人一样，引起人们注意的最绝妙的方法，就是损害
别人。
那些从事十分有益处的清洁工作的甲虫，并不能引来人们的注意与称赞，甚至无人
去理睬它们；而吃人血的蚊虫，却是每个人都知道的；同时人们也知道那些带着毒剑，
暴躁而又虚夸的黄蜂，以及专做坏事的蚂蚁。后者在我们南方的村庄中，常常会跑到人
们的家里面弄坏椽子，而且它们在做这些坏事时，还像品尝无花果一样高兴。
我花园里的蟋蟀，已经完全被蚂蚁残杀殆尽，这就使得我不得不跑到外面的地方去
寻找它们。八月里在落叶下，那里的草还没有完全被太阳晒枯干，我看到幼小的蟋蟀，
已经长得比较大了，全身已经都是黑色了，白肩带的痕迹一点也没有存留下来，在这个
时期，它的生活是流浪式的，一片枯叶，一块扁石头，已经足够它去应付大千世界中的
一些事情了。
许多从蚂蚁口中逃脱而残生的蟋蟀，现在又作了黄蜂的牺牲品。它们猎取这些旅行
者，然后把它们埋在地下。其实只要蟋蟀提前几个星期做好防护工作，它们就没有这种
危险了。但是它们从来也没想到过这点，总是死守着旧习惯，仿佛视死如归的样子。
一直要到十月末，寒气开始袭人时，蟋蟀才开始动手建造自己的巢穴。如果以我们
对养在笼子里的蟋蟀的观察来判断，这项工作是很简单的。挖穴并不在裸露的地面上进
行，而是常常在莴苣叶——残留下来的食物——掩盖的地点。或者是其它的能代替草叶
的东西，似乎为了使它的住宅秘密起见，这些掩盖物是不可缺少的。
这位矿工用它的前足扒着土地，并用大腮的钳子，咬去较大的石块。我看到它用强
有力的后足蹬踏着土地，后腿上长有两排锯齿式的东西。同时，我也看到它清扫尘土将
其推到后面，把它倾斜地铺开。这样，就可以知道蟋蟀挖掘巢穴的全部方法了。
工作开始做的很快。在我笼子里的土中，它钻在下面一呆就是两个小时，而且隔一
小会儿，它就会到进出口的地方来。但是它常常是向着后面的，不停地打扫着尘土。如
果它感到劳累了，它可以在还没完成的家门口休息一会儿，头朝着外面，触须特别无力
地摆动，一副倦怠的样子。不久它又钻进去，用钳子和耙继续劳作。后来，休息的时间
渐渐加长，这使我感到有些不耐烦了。
这项工作最重要的部分已经完成了。洞口已经有两寸多深了，足够满足一时之需。
余下的事情，可以慢慢地做，今天做一点，明天再做一点，这个洞可以随天气的变冷和
蟋蟀身体的长大而加大加深。如果冬天的天气比较暖和，太阳照射到住宅的门口，仍然
还可以看见蟋蟀从洞穴里面抛散出泥土来。在春天尽情享乐的天气里，这住宅的修理工
作仍然继续不已。改良和装饰的工作，总是经常地不停歇地在做着，直到主人死去。
四月的月底，蟋蟀开始唱歌，最初是一种生疏而又羞涩的独唱，不久，就合成在一
起形成美妙的奏乐，每块泥土都夸赞它是非常善于演奏动听的音乐的乐者。我乐意将它
置于春天的歌唱者之魁首。在我们的荒废了的土地上，在百里香和欧薄荷繁盛的开花时
节，百灵鸟如火箭般飞起来，打开喉咙纵情歌唱，将优美的歌声，从天空散布到地上。
而呆在下面的蟋蟀，它们也禁不住吸引，放声高歌一曲，以求与相知者相应和。它们的
歌声单调而又无艺术感，但它的这种艺术感和它生命复苏的单调喜悦相协调，这是一种
警醒的歌颂，为萌芽的种子和初生的叶片所了解、所体味。对于这种二人合奏的乐曲，
我们应该判定蟋蟀是优秀中的胜者。它的数目和不间断的音节足以使它当之无愧。百灵
鸟的歌声停止以后，在这些田野上，生长着青灰色的欧薄荷，这些在日光下摇摆着芳香
的批评家，仍然能够享受到这样朴实的歌唱家的一曲赞美之歌，从而伴它们度过每一刻
寂寞的时光。多么有益的伴侣啊！它给大自然以美好的回报。

三、它的乐器
为了科学的研究，我们可以很坦率地对蟋蟀说道：“把你的乐器给我们看看。”像
各种有价值的东西一样，它是非常简单的。它和螽斯的乐器很相像，根据同样的原理，
它不过是一只弓，弓上有一只钩子，以及一种振动膜。右翼鞘遮盖着左翼鞘，差不多完
全遮盖着，只除去后面和转折包在体侧的一部分，这种样式和我们原先看到的蚱蜢、螽
斯，及其同类相反。蟋蟀是右边的盖着左边的，而蚱蜢等，是左边的盖着右边的。
两个翼鞘的构造是完全一样的。知道一个也就知道另一个了。它们分别平铺在蟋蟀
的身上。在旁边，突然斜下成直角，紧裹在身上，上面还长有细脉。
如果你把两个翼鞘揭开，然后朝着亮光仔细地留意，你可以看到它是极其淡的淡红
色，除去两个连接着的地方以外，前面是一个大的三角形，后面是一个小的椭圆，上面
生长有模糊的皱纹，这两个地方就是它的发声器官，这里的皮是透明的，比其它的地方
要更加紧密些，只是略带一些烟灰色。
在前一部分的后端边隙的空隙中有五条或是六条黑色的条纹，看来好像梯子的台阶。
它们能互相磨擦，从而增加与下面弓的接触点的数目，以增强其振动。
在下面，围绕着空隙的两条脉线中的一条，呈肋状。切成钩的样子的就是弓，它长
着约一百五十个三角形的齿，整齐得几乎符合几何学的规律。
这的确可以说是一件非常精致的乐器。弓上的一百五十个齿，嵌在对面翼鞘的梯级
里面，使四个发声器同时振动，下面的一对直接摩擦，上面的一对是摆动摩擦的器具，
它只用其中的四只发音器就能将音乐传到数百码以外的地方，可以想象这声音是如何的
急促啊！
它的声音可以与蝉的清澈的鸣叫相抗衡，并且没有后者粗糙的声音。比较来说，蟋
蟀的叫声要更好一些，这是因为它知道怎样调节它的曲调。蟋蟀的翼鞘向着两个不同的
方向伸出，所以非常开阔。这就形成了制音器，如果把它放低一点，那么就能改变其发
出声音的强度。根据它们与蟋蟀柔软的身体接触程度的不同，可以让它一会儿能发出柔
和的低声的吟唱，一会儿又发出极高亢的声调。
蟋蟀身上两个翼盘完全相似，这一点是非常值得注意的。我可以清楚地看到上面弓
的作用，和四个发音地方的动作。但下面的那一个，即左翼的弓又有什么样的用处呢？
它并不被放置在任何东西上，没有东西接触着同样装饰着齿的钩子。它是完全没有用处
的，除非能将两部分器具调换一下位置，那下面的可以放到上面去。如果这件事可以办
到的话，那么它的器具的功用还是和以前相同，只不过这一次是利用它现在没有用到的
那只弓演奏了。下面的胡琴弓变成上面的，但是所演奏出来的调子还是一样的。
最初我以为蟋蟀是两只弓都是有用的，至少它们中有些是用左面那一只的。但是观
察的结果恰恰与我的想象相反。我所观察过的蟋蟀(数目很多)都是右翼鞘盖在左翼鞘上
的，没有一只例外。
我甚至用人为的方法来做这件事情。我非常轻巧地，用我的钳子，使蟋蟀的左翼鞘
放在右翼鞘上，决不碰破一点儿皮。只要有一点技巧和耐心，这件事情是容易做到的。
事情的各方面都得很好，肩上没有脱落，翼膜也没有皱褶。
我很希望蟋蟀在这种状态下仍然可以尽情歌唱，但不久我就失望了。它开始回复到
原来的状态。我一而再而三地摆弄了好几回，但是蟋蟀的顽固终于还是战胜了我的摆布。
后来我想这种试验应该在翼鞘还是新的、软的时候进行，即在蛴螬刚刚蜕去皮的时
候。我得到刚刚脱化的一只幼虫，在这个时候，它未来的翼和翼鞘形状就像四个极小的
薄片，它短小的形状和向着不同方向平铺的样子，使我想到面包师穿的那种短马甲，这
蛴螬不久就在我的面前，脱去了这层衣服。
小蟋蟀的翼鞘一点一点长大，渐渐变大，这时还看不出哪一扇翼鞘盖在上面。后来
两边接近了，再过几分钟，右边的马上就要盖到左边的上面去了。于是这时是我加以干
涉的时候了。
我用一根草轻轻地调整其鞘的位置，使左边的翼鞘盖到右边的上面。蟋蟀虽然有些
反抗，但是最终我还是成功了。左边的翼鞘稍稍推向前方，虽然只有一点点。于是我放
下它，翼鞘逐渐在变换位置的情况下长大。蟋蟀逐渐向左边发展了。我很希望它使用它
的家族从未用过的左琴弓来演奏出一曲同样美妙动人的乐曲。
第三天，它就开始了。先听到几声摩擦的声音，好像机器的齿轮还没有切合好，正
在把它调整一样。然后调子开始了，还是它那种固有的音调。
唉，我过于信任我破坏自然规律的行为了。我以为已造就了一位新式的奏乐师，然
而我一无所获。蟋蟀仍然拉它右面的琴弓，而且常常如此拉。它因拼命努力，想把我颠
倒放置的翼鞘放在原来的位置，导致肩膀脱臼，现在它已经经过自己的几番努力与挣扎，
把本来应该在上面的翼鞘又放回了原来的位置上，应该放在下面的仍放在下面。我想把
它做成左手的演奏者的方法是缺乏科学性的。它以它的行动来嘲笑我的做法，最终，它
的一生还是以右手琴师的身份度过的。
乐器已讲得够多了，让我们来欣赏一下它的音乐吧！蟋蟀是在它自家的门口唱歌的，
在温暖的阳光下面，从不躲在屋里自的欣赏。翼鞘发出“克利克利”柔和的振动声。音
调圆满，非常响亮、明朗而精美，而且延长之处仿佛无休止一样。整个春天寂寞的闲暇
就这样消遣过去了。这位隐士最初的歌唱是为了让自己过得更快乐些。它在歌颂照在它
身上的阳光，供给它食物的青草，给它居住的平安隐避之所。它的弓的第一目的，是歌
颂它生存的快乐，表达它对大自然恩赐的谢意。
到了后来，它不再以自我为中心了，它逐渐为它的伴侣而弹奏。但是据实说来，它
的这种关心并没收到感谢的回报，因为到后来它和它的伴侣争斗的很凶，除非它逃走，
否则它的伴侣会把它弄成残废，甚至将吃掉它一部分的肢体。不过无论如何，它不久总
要死的，就是它逃脱了好争斗的伴侣，在六月里它也是要死亡的。听说喜欢听音乐的希
腊人常将它养在笼子里，好听它们的歌唱。然而我不信这回事，至少是表示怀疑。第一，
它发出的略带烦嚣的声音，如果靠近听久了，耳朵是受不了的，希腊人的听觉恐怕不见
得爱听这种粗糙的，来自田野间的音乐吧！
第二，蝉是不能养在笼子里面的，除非我们连洋橄榄或榛系木一齐都罩在里面。但
是只要关一天，就会使这喜欢高飞的昆虫厌倦而死。
将蟋蟀错误地作为蝉，好像将蝉错误地当作蚱蜢一样，并不是不可能的。如果如此
形容蟋蟀，那么是有一定道理的。它被关起来是很快乐的，并不烦恼。它长住在家里的
生活使它能够被饲养，它是很容易满足的。只要它每天有莴苣叶子吃，就是关在不及拳
头大的笼子里，它也能生活的很快乐，不住地叫。雅典小孩子挂在窗口笼子里养的，不
就是它吗？
布罗温司的小孩子，以及南方各处的小孩子们，都有同样的嗜好。至于在城里，蟋
蟀更成为孩子们的珍贵财产了。这种昆虫在主人那里受到各种恩宠，享受到各种美味佳
肴。同时，它们也以自己特有的方式来回报好心的主人，为他们不时地唱起乡下的快乐
之歌。因此它的死能使全家人都感到悲哀，足可以说明它与人类的关系是多么亲密了。
我们附近的其它三种蟋蟀，都有同样的乐器，不过细微处稍有一些不同。它们的歌
唱在各方面都很像，不过它们身体的大小各有不同。波尔多蟋蟀，有时候到我家厨房的
黑暗处来，是蟋蟀一族中最小的，它的歌声也很细微，必须要侧耳静听才能听得见。
田野里的蟋蟀，在春天有太阳的时候歌唱，在夏天的晚上，我们则听到意大利蟋蟀
的声音了。它是个瘦弱的昆虫，颜色十分浅淡，差不多呈白色，似乎和它夜间行动的习
惯相吻合。如果你将它放在手指中，你就会怕把它捏扁。它喜欢呆在高高的空气中，在
各种灌木里，或者是比较高的草上，很少爬下地面来。在七月到十月这些炎热的夜晚，
它甜蜜的歌声，从太阳落山起，继续至半夜也不停止。
布罗温司的人都熟悉它的歌声，最小的灌木叶下也有它的乐队。很柔和很慢的“格
里里，格里里”的声音，加以轻微的颤音，格外有意思。如果没有什么事打扰它，这种
声将会一直持续并不改变，但是只要有一点儿声响，它就变成迷人的歌者了。你本来听
见它在你面前很靠近的地方，但是忽然你听起来，它已在十五码以外的地方了。但是如
果你向着这个声音走过去，它却并不在那里，声音还是从原来的地方传过来的。其实，
也并不是这样的。这声音是从左面，还是从后面传来的呢？一个人完全被搞糊涂了，简
直辨别不出歌声发出的地点了。
这种距离不定的幻声，是由两种方法造成的。声音的高低与抑扬，根据下翼鞘被弓
压迫的部位而不同，同时，它们也受翼鞘位置的影响。如果要发较高的声音，翼鞘就会
抬举得很高；如果要发较低的声音，翼鞘就低下来一点。淡色的蟋蟀会迷惑来捕捉它的
人，用它颤动板的边缘压住柔软的身体，以此将来者搞昏。
在我所知道的昆虫中，没有什么其它的歌声比它更动人、更清晰的了。在八月夜深
人静的晚上，可以听到它。我常常俯卧在我哈麻司里迷迭香旁边的草地上，静静地欣赏
这种悦耳的音乐。那种感觉真是十分的惬意。
意大利蟋蟀聚集在我的小花园中，在每一株开着红花的野玫瑰上，都有它的歌颂者，
欧薄荷上也有很多。野草莓树、小松树，也都变成了音乐场所。并且它的声音十分清澈，
富有美感，特别动人。所以在这个世界中，从每棵小树到每根树枝上，都飘出颂扬生存
的快乐之歌。简直就是一曲动物之中的“欢乐颂”！
高高的在我头顶上，天鹅飞翔于银河之间，而在地面上，围绕着我的，有昆虫快乐
的音乐，时起时息。微小的生命，诉说它的快乐，使我忘记了星辰的美景，我已然完全
陶醉于动听的音乐世界之中了。那些天眼，向下看着我，静静的，冷冷的，但一点也不
能打动我内在的心弦。为什么呢？因为它们缺少一个大的秘密——生命。确实，我们的
理智告诉我们：那些被太阳晒热的地方，同我们的一样，不过终究说来，这种信念也等
于一种猜想，这不是一件确实无疑的事。
在你的同伴里，相反的啊，我的蟋蟀，我感到生命的活力，这是我们土地的灵魂，
这就是为什么我不看天上的星辰，而将注意力集中于你们的夜歌的原因了。一个活着的
微点——最小最小的生命的一粒，它的快乐和痛苦，比无限大的物质，更能引起我的无
限兴趣，更让我无比地热爱你们！
------------------








娇小的赤条蜂



细细的腰，玲珑的身材，腹部分成两节，下面大，上面小，中间好像是用一根细线
连起来，黑色的肚皮上面围着一丝红色的腰带：这就是赤条蜂。
赤条蜂的巢穴是建筑在疏松的极容易钻通的泥土里的。小路的两旁，太阳照耀着的
泥滩上，那些地方的草长得很稀疏，都是黄蜂最理想的住所。在春季，四月初的时候，
我们总可以在这样的地方找到它们。
赤条蜂通常在泥土里筑一个垂直的洞，好像一口井，口径只有鹅毛管那么粗，约有
二寸深，洞底是一个孤立的小房间，专为产卵用的。黄蜂建巢的时候，总是静静地、慢
慢地工作着，丝毫没有什么热烈或兴奋的样子。像别的蜂一样，它用前足作耙，用嘴巴
作挖掘的工具。有时候我们就可以听到，从洞底发出一声尖利、刺耳的摩擦声，这是因
为它遇到了一颗极不容易搬去的沙粒引起翅膀和全身剧烈振动的缘故。每隔短短的十几
分钟，我们就可以看到赤条蜂在洞口出现，嘴里衔着一些垃圾或是一颗沙粒。它总把这
种垃圾丢到几寸以外的地方，这样可以保持自己的居所和周围环境的整洁干净。
有几颗沙粒会被区别对待。赤条蜂们会对它们进行特殊的处理，使它们免遭被远远
抛出去的命运。这些沙粒被赤条蜂们堆在洞的附近，将来会另有重大的用途。当赤条蜂
把洞完全挖好了，它就在这小沙滩上察看有没有适合它需要的沙粒。如果没有，它就到
附近去找，直到找到为止。它需要的是一粒扁平的，比它的洞口稍大一些的沙粒，它可
以把这个沙粒盖在洞口，做成一扇门。第二天它从外面猎取一条毛毛虫回来，就不慌不
忙地把门打开，把猎物拖进去。这门看起来和其它沙粒完全一样，谁也不会想到它底下
会藏着食物，藏着一只赤条蜂的家，只有它自己才能辨别出它的家。它打开门，不紧不
但地把猎物放到洞底后，就开始在上面产卵，然后再用它以前藏在附近的沙粒把洞口堵
住。这听起来似乎有点像《阿里巴巴与四十大盗》中“芝麻开门”的故事。
赤条蜂所猎取的食物是一种灰蛾的幼虫。这种虫大部分都是生活在地底下的，赤条
蜂又是怎么把它捉到的呢？让我们来看看吧。有一天当我散步回来的时候，正好看到一
只赤条蜂在一丛百里香底下忙碌着。我立刻在它附近的地上躺下。我的出现并没有把它
吓走。它先飞到我的衣袖上停了一会儿，断定我不会伤害它之后，又飞回到百里香丛中
去了。从过去的经验我知道这意味着什么：它忙得很，没有时间来考虑我这个不速之客。
赤条蜂把百里香根部的泥土挖去，又把周围的小草拔掉，然后把头钻进它挖松的土
块里。它匆匆忙忙地从这里飞到那里，向每一条裂缝里张望。而不是在为自己筑巢，而
是在寻找地底下的食物，活像一只猎狗在寻找洞里的野兔一般。
灰蛾的幼虫觉察到了上面的动静，决定离开自己的巢，爬到地面上来看看到底发生
了什么事。这一念之差就决定了它的命运。那赤条蜂是早已准备就绪，就等着它的出现
了。果然，灰蛾的幼虫一露出地面，赤条蜂就冲出去一把将它抓住了，然后伏在它的背
上，像一个熟练的外科医生一样，不慌不忙地用刺把毛虫的每一节都刺一下。从前到后
一节一节地往下刺，一点儿也不遗漏。它那熟练的动作，让人想起游刃有余的屠夫。
赤条蜂的技巧令科学家们都自愧不如。它可以靠观察去推断人类所从不知道的事情。
它很熟悉它的俘虏的神经系统，它知道往哪些神经中枢上扎刺，可以使它的俘虏神经麻
木而不致于死亡。它从哪里学到的这种知识呢？我们人类有学校，有老师，有各种书籍，
我们通过这些逐渐地积累知识懂得大自然的许多奥妙。可是赤条蜂是怎么知道这些复杂
的知识的呢？而且不用练习就掌握了那么熟练的技术，难道是在它们出生前，冥冥之中
就有神灵赋予它们这种本领了吗？大自然多神奇啊，当我们孜孜不倦地探索它的秘密时，
它早就有条不紊地安排好了一切！
我还要告诉你另一幕赤条蜂和毛毛虫大战的故事，这也是我亲眼看见的。那是在五
月里，我看着赤条蜂在一条光洁的路旁为它的巢作最后一步的清除工作。它的几码外的
地方已经麻醉好了一条毛毛虫，当它清除好那条街道并且把洞开得足够大后，它就出去
搬毛毛虫了。
它很容易地找到了那条毛毛虫，那条毛毛虫就躺在地上。可是糟糕的是，蚂蚁也正
在猎取那条毛毛虫。赤条蜂不愿意和蚂蚁分享这条毛毛虫，可是要把蚂蚁赶掉，也不是
一件容易的事，再三考虑之后，它认为自己的能力实在有限，还是不要作无谓的牺牲吧。
于是它决定放弃这条毛毛虫，再到别处去寻找食物。
它在离巢大约十尺以内的地方，一步一步慢慢地走着，察看着泥土，不时地用它那
弯着的触须，在地面上挥动，像一名执着的士兵用探雷针寻找着地雷。在烈日的下面，
我观察了它整整三个钟头！要找到一条毛毛虫是多么困难啊，尤其是在急需的时候。
即使是对人，这也是一件困难的工作。我一心要帮助它，替它找到一条毛毛虫，因
为我想看它怎样麻痹毛毛虫。
于是我就想起我的老朋友法维，他是我的园丁，正在那里照料花园，于是我把他招
呼过来。
“快来快来！我想要几条灰色毛毛虫！”我把事情飞快地向他解释了一下。他明白
了，马上去找虫子。他挖掘着莴苣的根，耙着草莓里的泥，察看着鸢尾草丛的边缘。我
非常信任他，他的眼力和聪明，因为那么多年以来，大家都认为他是一个出色的园丁。
过了好久，也没见法维拿毛毛虫过来。
“喂，法维，毛毛虫呢？”
“先生，我一条也找不着。”
“怎么会？！把你们所有的人都叫过来！克兰亚、爱格兰，你们都来！到莴苣田里
来！帮我找毛毛虫！”
于是全家都出动了，每个人都很努力，可是毫无结果。三个钟头过去了。我们中间
谁也没找着毛毛虫。
赤条蜂也没有找着毛毛虫，它已经很疲倦了，我看到它很果断地在地面上有裂缝的
地方寻找。它尽着它最大的努力寻找，甚至把杏核般大的泥块搬开。可是不久它又离开
了这些地方。于是我开始怀疑，赤条蜂捕获不到猎物，不是因为找不到毛毛虫，而是因
为虽然它知道毛毛虫在哪儿，却没办法捉到它们，可能因为毛毛虫早有防备，把巢挖得
很深，而赤条蜂没有能力把虫子从地底下挖出来。我真愚昧，为什么早没有想到这一点？
难道这样一个经验丰富的猎取家会盲目地浪费精力吗？当然不会。
此刻赤条蜂又在挖另一个地方了，可是不久它又放弃了，正像它所尝试过的许多地
方一样。我决心要帮它的忙，于是我就继续它的工作，用小刀朝那儿挖下去。可是什么
也没有，于是我也放弃了那块地方。
可是不久赤条蜂又回来了，在我挖过的地方继续往下挖。我明白它的意思了，我为
它创造了条件，重新激起它对这个地方的信心。
“滚开，你这个笨头笨脑的家伙！”赤条蜂似乎在说“让我来告诉你这里到底有没
有毛毛虫！”
于是我按照赤条蜂指引的方向挖下去，果然挖出了一条毛毛虫。太好了！聪明的赤
条蜂！你没有辜负我对你的信任！
照这种办法，我挖到了第二条毛毛虫，不久，第三条、第四条都被我挖到了。我发
现赤条蜂所挖掘的地方，都是光秃秃的几个月前曾经翻松过的泥土。除此之外，再也没
有什么记号可以显示出毛毛虫的所在地了。好了，法维、克兰亚、爱格兰和你们其余所
有的人，你们还有什么话说？你们找了三个钟头都没找到一条毛毛虫，而这只聪明的赤
条蜂，却提供给了我足够的毛毛虫，同时，我也为自己对赤条蜂的信任和了解而沾沾自
喜。是啊，我能够懂得它的心思，能够和它密切配合，互补长短，那一堆丰盛的“战果”
就是我们之间天衣无缝的完美合作的最好证明。

袭击毛毛虫
我把第五条毛毛虫留给赤条蜂，当时我正躺在地上，和这位屠夫靠得很近，所以没
有一个细小的动作能逃过我的眼睛。现在我要把我眼前所发生的情景一段一段地记下来。
1．赤条蜂用它的嘴巴夹住了毛毛虫的颈部，毛毛虫剧烈地挣扎着，扭动着身体。赤
条蜂却不慌不忙，自己让到一边，以避免剧烈的冲撞。它的刺扎在毛毛虫的头和第一节
之间的关节上，那是毛毛虫的皮最嫩的地方。这是最性命攸关的一下，这一下可以使毛
毛虫完全受赤条蜂的控制。
2．赤条蜂突然离开毛毛虫，躺倒在地上，剧烈地扭动着，不停地打着滚，抖动着足，
拍打着翅膀，像是在垂死挣扎。我以为它也被毛毛虫扎了一下，受了致命的伤。看着它
的生命就要这样结束，我对它充满了无限的同情。可是它突然又恢复正常了，扇扇翅膀，
理理须发，又活灵活现地回到猎物旁。刚才看到的那一幕，其实正是它庆祝胜利的表现，
而不是像我想的那样受了伤。
3．赤条蜂抓住了毛毛虫的背部，抓的部位比第一次稍微低些，然后开始用刺扎它身
体的第二个体节，仍旧是刺在下方。它一节一节地往下刺。头三节上有脚，接着两节没
有脚，再以后四节又有脚，不过那不是真正的脚，充其量也只能算是沈状突起物。一共
有九节。但早在赤条蜂第一针刺下后，毛毛虫已经没有多大的抵抗力了。
4．最后赤条蜂把钳子般的嘴巴张到最大的宽度，钳住了毛毛虫的头，有节奏地轻轻
压榨它，但尽量不使它受伤。每压一次，赤条蜂就要停一下，看看毛毛虫有什么反应。
这样一停、一等、一压，循环往复地进行着。这种控制大脑的手术不能做得太猛烈，否
则毛毛虫很可能会死掉。说来很奇怪，赤条蜂为什么并不想让它死掉呢？
现在外科医生的手术已经结束，毛毛虫瘫趴在地上。它不会动了，几乎没有生命，
只有一息尚存。它任凭被赤条蜂拖到洞里，不做也不能做丝毫的反抗。当赤条蜂把卵产
在它身上后，它也没有能力伤害在它身上成长的赤条蜂的幼虫。这就是赤条蜂所做的麻
醉工作的目的：它是在为未来的婴儿预备食物。它把毛毛虫拖到洞里以后，就在它身上
产一个卵。等到幼虫从卵里孵化出来，就可以把毛毛虫当作食物。想想看，如果毛毛虫
还能动弹会有什么后果？只要它轻轻一转身，就能轻而易举地把赤条蜂的卵压破！当然，
毛毛虫是不会动了。可是它又不能完全死掉，因为如果它死了，尸体很快就会腐烂，不
适宜做赤条蜂幼虫的食物了。所以赤条蜂用它的毒刺刺进毛毛虫的每一节神经中枢，使
它失去运动的能力，半死不活地苟且残喘下去，自动地为幼虫将来的食物“保鲜”。赤
条蜂想得多周到啊！不过，等你看到它把猎物拉回家的过程，你会发现它对事物考虑的
周到程度还远不止这些。它想到毛毛虫的头部还没有受伤，嘴巴还能动。当它被赤条蜂
拖着走的时候，它能够咬住地上的草，从而阻碍赤条蜂继续把它往前拖，所以赤条蜂还
得想办法把毛毛虫的头部也麻痹掉。这次它不再用它的毒刺，因为那会致毛毛虫于死地。
它连续不断地压榨和摩擦毛毛虫的头部，这种方法实在是十分恰到好处，毛毛虫很快便
失去了知觉，它被折腾晕了。
虽然我们很佩服赤条蜂的技巧，但同时也不禁为毛毛虫们捏把汗。毕竟，被赤条蜂
弄得求生不得，求死不能是一件很残忍的事，不过如果我们是农夫，我们就不会对毛毛
虫深表同情了。它是农作物和花草最可怕的敌人！它们白天弯曲着身子在洞里睡大觉，
晚上爬出来为非作歹，咬植物的根和茎。无论是观赏用的花草还是食用的蔬菜，都是它
们的美餐。有时候，一棵幼苗无缘无故地枯萎了，如果你把它轻轻地拔起来，就会发现
它的根部受了伤。原来，晚上的时候，恶贯满盈的毛毛虫到过这儿，它那剪刀一样的嘴
巴把这棵幼苗咬伤了。它和另一种白色的毛毛虫(金虫的幼虫)一样的坏。只要它到了甜
菜园里，这菜园就要蒙受巨大的损失。它的的确确是一种害虫。赤条蜂把它杀死，是在
为民除害。我们就完全没有必要对毛毛虫产生同情了。
------------------








西西斯



我希望你们听过了关于清道的甲虫做球的奇怪的事情，还不至于厌倦。我已经告诉
过你们神圣甲虫和西班牙的犀头，现在我想再讲一些这种动物的其它种类。在昆虫的世
界里，我们遇到过许多模范母亲，现在只是为了好玩，来注意一回好的父亲吧！
除非在高等动物中，好的父亲是很少见的。在这方面，鸟类是优秀的，而人类最能
尽这种义务。低级动物当中，父亲对家族中的事情是漠不关心的。很少有昆虫是这种定
律的例外。这种无情，在高级动物的世界中是要被厌恶的，而且它们幼小的动物不需要
长时间的看护。而对于昆虫的父亲来说这是可以原谅的。因为只要有个适当的地点，新
生昆虫就可以十分健康地成长，很可能无需帮助而得到食物。例如粉蝶为了种族的安全，
只要把卵产在菜叶上，父亲的责任心又有什么用呢？母亲有利用植物的本能，是不需要
帮助的。在产卵的时候，也是不需要父亲在一边保护的。
许多昆虫都采用一种简单的养育法。即它们先要找一个餐室，当作幼虫卵化后的家，
或者先找一个地方，使幼虫自己能觅到适当的食物，在出生后食用。在这种情况下，它
们是不需要父亲的，所以父亲通常到死都没有给它的后代成长工作以丝毫的帮助。
然而事情也不是常常以这种原始的方式进行的。有些种族为它们的家庭预备下妆奁，
作它们将来的食宿。蜜蜂和黄蜂特别善于营造小巢，例如口袋、小瓶等，并在里面装满
蜜，它们还十分善于造筑土穴，储藏着野味，给蛴螬做食物。
这种伟大的建筑巢穴和收集食物的工作，要花去它全部的生命，而这工作却是母亲
单独一人做的。这工作消磨它的时间，耗去它的生命。父亲则沉醉于日光下，懒惰地站
在工作场之外，只是看着它的勤劳的伴侣在从事艰苦的工作。
为什么它不帮助一下呢？事实上它从没有帮助过。为什么它不学学燕子夫妻，它们
都带一些草和一些泥土到巢里，还带一些小虫给小鸟。而雄性昆虫一点也没做那种事。
也许它借口比较软弱，无以作辩解。这是个无聊的议论。因为在叶子上割下一块，从植
物上摘下一些棉花，从泥土中收集一点水泥，完全是它的力量所能做到的。它很可以像
工人一样地帮助雌虫，它很适于收集一些材料，再由更智慧的雌虫建筑起来。它不做的
真正原因，只是因为它不愿做而已。
多数从事劳作的昆虫，竟然都不知道做父亲的责任，这是很令男人们感到奇怪的。
谁都为了幼虫发展最高才能的需要而努力，但这些父亲们仍然愚钝如蝴蝶，对于家族是
很少费力的。我们每一次都不能回答下面的问题：为什么这种昆虫，有这个特别的本能，
而别的昆虫就没有呢？
当我们看到清道夫甲虫有这种高贵的品质而收蜜的昆虫却没有时，我们非常的惊奇
和难以理解，好多种清道夫甲虫善于负起家政的重任，并知道两人共同工作的价值。例
如蜣螂夫妻，它们共同预备蛴螬的食物，父亲帮助它的伴侣在制造腊肠般的食物时，助
以强有力的轧榨工作。
它们就是形成家族共同劳作的习惯的最好的榜样，在普遍的自私的情形中，是最稀
罕的一个例外。
关于这件事，经过我长期的研究，在这个例子之外，我又可以增加三个另外的例子，
全都是清道夫甲虫合作的事实。
这三个中的一个是锡赛弗斯，它是搓丸药者中最小最勤劳的一个。它在它们当中最
活泼、最灵敏，并且毫不介意在危险的道路上倾倒和翻跟斗，在那里它固执地爬起来，
但又重新倒下去。正是因为那些狂乱的体操，所以人们给它起了一个名字，叫西西弗。
我想你们总该知道，一个可怜的人变得很著名，一定要经过很多很艰难的奋斗。它
被迫把一块大石头滚上高山，每次好容易到山顶时，那石头又滑落，滚到山脚下。我很
喜欢这个神话，这是我们当中许多人的历史，就拿我自己来说，刻苦的爬峻峭的山坡已
经五十多年了，把我的精力浪费在为了安全地得到每天的面包的挣扎里。面包一经滑去，
就滚下去，落到深渊里，很难摆稳。
我现在所谈及的西西弗，就不知道有这种困难，它在陡峭的山坡上毫无挂念地滚着
粮食，有时供给它自己，有时供给它的子女。在我们这些地方，它是很少见的，如果没
有我们从前几次提起过的助手，我也没有地方得到这么多观察物来研究。
我的小儿子保罗，年龄才七岁。他是我猎取昆虫的热心的同伴，而且比任何同龄的
小孩，更清楚地知道蝉、蝗虫、蟋蟀的秘密，尤其是清道夫甲虫。他锐利的眼光能在二
十步以外，辨别出地上隆起的土堆，哪一个是甲虫的巢穴，哪一个不是。他的灵敏的耳
朵可以听到螽斯细微的歌声，这是我完全办不到的。他帮助我看和听，我则把意见给他
用以作交换，他是很注意接受我的意见的。
小保罗有他自己养虫子的笼子，神圣的甲虫在里面做巢。他自己的花园，和手帕差
不多大小，能在里面种些豆子，但他常常将它们挖起来，看看小根长了一点没有。他的
林地上，有四株小槲树，只有手掌那么高，一旁还连着槲树子，在供给它养料。这是研
究昆虫之余极好的休息，对于昆虫研究的进步是毫无妨碍的。
五月将近的时候，有一天保罗和我起得很早，因为太早了，出去时连早饭都没有吃，
我们在山脚下的草场上，在羊群曾经走过的地方寻找。在这里，我们找着了西西弗，保
罗非常热心地搜索，不久我们就得到了足足好几对，收获真是不少。
使它们安居下来所需要的是一个铁丝的罩子，沙土的床，以及食物的供给——为了
这个我们也变成清道夫了。这些动物是很小的，还不及樱桃核大。形状也很奇怪！一个
短而肥的身体，后部是尖的，足很长，伸开来和蜘蛛的脚很像。后足更长，呈弯曲形，
挖土和搓小球时最有用。
不久，建设家族的时候到了。父亲和母亲同样热心地从事着搓卷、搬动和贮藏食物
的工作，都是为了它们的子女。它们利用前足的刀子，随意的从食物上割下小块来。夫
妻俩一同工作，一次次地抚拍和挤压，做成了一粒豌豆大的球。
和在神圣甲虫的工作场里一样，它们把圆球做成正确的圆形，是用不着机械的力量
来滚这球的。材料在没有移动之前，甚至在没有拾起之前，就已经做成圆形了。现在我
们又有了一个圆形学家，善于制作制造和保藏食物的最好的样式——圆形。
球不久就制造成功了。现在须用力地滚动，使它具有一层硬壳，保护里面柔软的物
质，使它不致变得太干燥。我们可以从大一些的身段上辨别出在前面全幅武装的母亲。
它将长长的后足放在地上，前足放在球上，将球向自己的身边拉，向后退着走。父亲处
在相反的方位，头向着下面，在后面推，这与神圣甲虫两个在一起工作时的方法相同，
不过目的是两样的，西西弗夫妻是为蛴螬搬运食物，而大的滚梨者(即神圣甲虫)则是准
备食物为自己在地下大嚼用的。
这一对在地面上走过，它们没有固定的目标，只是一直走下去，不管横在路中央的
障碍物。这样倒退着走，障碍固然是免不了的，但是即使看到了，它们也不会绕过障碍
走。它甚至做顽固的尝试，想爬过我的铁丝笼子。这是一种费时而且不可能完成的工作，
母亲的后足抓住铁丝网将球向它拉过来，然后用前足抱住它，把球抱在空中。父亲觉得
无物可推就抱住了球，伏在上面，把它身体的重量，加在球上，不再费什么力气了。这
种努力维持下去，未免太难了。于是球和骑在上面的昆虫，滚成一团，掉落到地上。母
亲从上面惊异地看着下面，不久就下来了，扶好这个球，重新做这个不可能成功的尝试。
一再的跌落之后，才放弃攀爬这个铁丝网。
就是在平地运输的时候，也不是完全没有困难的。差不多每分钟都会碰到隆起的石
头堆，货物就会翻倒。正在奋力推的昆虫也翻倒了，仰卧着把脚乱踢。不过这只是小事
情，很小很小的事情。西西弗是常常翻倒的，它并不在意。甚至有人也许以为它是喜欢
这样的。然而无论如何，球是变硬了，而且相当的坚固。跌倒、颠簸等都是程序单上的
一部分。这种疯狂的跳跃往往要继续几个小时！
最后母亲以为工作已经完毕，跑到附近找个适当的地点贮存球。父亲留守，蹲在实
物的上面。如果它的伴侣离开太久，它就用它高举的后足灵活地搓球，用以解闷儿，它
处置它珍贵的小球时，如同演戏者处置他的球一样。它用它变形的腿实验那个球是否完
整。那种高举的样子，无论谁看了，都不会怀疑它生活得很满足——父亲会保障它子女
将来的幸福的满足。
它好像是在说：“我搓成的这块圆球，是我给我的儿子们做面包的！”
它高高举起那个球，给每个人看看，这个是它工作的成果。这时候，母亲已经找到
了埋藏的地方，开始的一小部分工作已经做了，已经开始做下一个浅穴了。可以将球推
进浅穴。守卫的父亲一刻也不离开，母亲则用足和头挖土。不久，地穴已足够容纳那个
球了。它始终坚决地把球靠近自己，它一定感到在洞穴做成以前，一定要前后左右地把
它摇动摇动，以免寄生物的侵害。如果把它放在洞穴边上，一直到这个家完成，它害怕
会有什么不幸的事发生。因为有很多蚊蝇和别的动物，会出其不意地来攫取，因此不能
不格外当心。
于是圆球已经有一半放在还没有完成的土穴里了。母亲在下面，用足把球抱住往下
拉，父亲在上面，轻轻地往下放，而且还要注意落下去的泥土会不会把穴堵住，一切都
进行得很顺利。掘凿进行着，继续往下放球，非常的小心，一个虫往下拉，一个管理着
落下去的速度，并清除那些阻碍工作的事物。经过再进一步的努力，球和矿工都到地下
去了。以后所要做的事，是把从前做好的事再做一回，而且我们必须再等半天或几个小
时。
如果我们仔细等待，我们可以见到父亲又单独到地面上来了，它蹲在靠近土穴的沙
土。母亲为了尽它的伴侣不能帮助它的责任，常常要到第二天才出现。最后它也出来了，
父亲才离开它打瞌睡的地点，同它一道走。这对重新联合在一起的夫妇，又回到它们从
前找到食物的地方，休息一会，又收集起材料来。于是它们俩又重新工作，又一起塑模
型，运输和储藏球。
我对于这种恒心很是佩服。然而我不敢公然宣布，这是甲虫确定的习惯。无疑，有
许多甲虫是轻浮的，没有恒心的。但不要紧，我所看见的这点，关于西西弗爱护家庭的
习惯，已经使我看重了它们。
该是我们查看土巢的时候了。它并不很深，我看到墙边有一个小空隙，宽度足以让
母亲在球旁转动。寝室很小，这告诉我们父亲是不能在那里留很久的。当工作室准备好
了的时候，它一定要跑出去，请女雕刻家来继续工作。
地窖中只储藏着一只球，一件艺术的杰作。和神圣甲虫的梨形状相同，不过小得多。
因为小，球表面的光滑和圆形的标准，更加令人吃惊，最宽的地方，直径也只有一寸的
二分之一到四分之三。
另外还有一次对西西弗的观察。在我的铁丝笼下有六对，它们做了五十七个梨，每
个当中都有一颗卵——每一对平均有九个以上的蛴螬。神圣甲虫远不及这个数目。什么
原因使它产下这么多的后代呢？我看只有一个理由，就是父亲和母亲共同工作，一个家
族的负担，一人的精力不足以应付，两人分担起来就不觉得太重了。
------------------








捕蝇蜂



你已经知道了赤条蜂和黄蜂怎样麻痹毛毛虫或蟋蟀来喂自己的孩子，然后怎样封闭
洞口，离开巢飞到别处去。不过并不是每一种蜂都是这样生活的，现在你将要听到另一
种蜂，它们每天用新鲜的食物喂它的孩子，这就是捕蝇蜂。
这种蜜蜂喜欢在明亮的阳光下和蔚蓝的天空中选择最轻最松的泥土做它的巢。我有
时候会在一片没有树荫的广场上观察它们。天气很热，要避免烈日的煎熬只有躺在小沙
堆后面，把头钻进兔子洞，或是为自己预备一把大伞。我就采取了后一种办法，如果大
家愿意在七月快要结束的时候来和我一同坐在这样的大伞下，那么他(或她)也可以和我
一起饱饱眼福：
一只捕蝇蜂突然飞来，毫不犹豫地停在某个地方，这地方在我看来和别处没什么不
同。它前足上长着一排排的硬毛，会使你想起一把扫帚，一个刷子或一个钉锚。它用前
脚工作，用四只后脚支持着自己的身体。它先把沙耙起，然后向后拂去，它的动作非常
快，使这些连续不断的沙子看上去像不住地流水一样流到七八寸以外的地方。这种沙粒
的飞射要维持五分钟或十分钟左右。
和这些沙粒堆在一起的，还有木屑、腐烂的叶片和其它废料的碎屑。捕蝇蜂把这些
垃圾一一用嘴搬掉，这就是它工作的目的。它要使它家门前的沙都是又轻又细的“高级
沙粒”，没有任何粗重的杂质。这样，当它为孩子们捕了蝇回来的时候，它就可以很容
易地打开一条通路，把猎物带到洞里去。这种清洁工作，它总是在空闲的时候做。譬如，
在猎物已经储藏了许多，足够它的孩子们吃一段时间的时候，它没有必要再出去觅食，
那么它就会抽出时间来清除垃圾，像一个出色的家庭主妇那样。我们可以看出，在它认
真勤勉地工作的时候，它显得非常快乐和满足。也许这正是一个做母亲的看到孩子们在
自己盖的屋子里成长起来所涌上心头的喜悦吧！
如果我们用一把小刀，在母蜂所刮的沙地上挖下去，我们首先会发现一条隧道，有
一个手指那么粗，或许有八寸到十二寸那么长，接着就是一个小屋。小屋的大小足以容
纳三个胡桃，可是到目前为止，这里面只有一个蝇和一只白色的小卵，那就是捕蝇蜂的
卵。大约二十四个小时之后，这个卵就能孵化成一条小虫，小虫出来后靠吃母亲为它准
备好的死蝇长大。
大约两三天之后，捕蝇蜂的幼虫就快把那死蝇吃完了。这时的母蜂离家并不很远，
你可以看到它有时从花蕊里吸几口蜜汁充饥，有时快乐地坐在火热的沙地上——它是在
看守自己的家。它会常常在家门口耙去一些沙，然后，又飞走了，过了一段时间再来。
可是不管它在外面呆多久，它总不会忘记估算一下它家里的小屋中的食物还能维持多久，
作为一个母亲，它的本能会告诉它什么时候它孩子的食物快吃完了，于是它就回到自己
的巢里。至于这巢，上面已经提到过，在外面看来和其它沙地是一样的，没有明显的洞
口或什么标记，可是它自己却清清楚楚地知道它的巢在哪一点。它每次回来探望孩子，
总不忘带些丰盛的礼物回来，这次它带回一只大蝇，它把蝇送进地下的家后，自己又出
来，直到需要它再送第三只蝇的时候再下去，这中间隔的时间是很短的，因为幼虫的胃
口始终很好，如果母蜂稍有懈怠的话，它的孩子就要挨饿了。
这样维持了两个星期，幼虫不停地在成长。食物的需求量也越来越大，母蜂也不断
地送食物进来。在第二个星期末，幼虫已经长得很肥胖了。母亲加倍努力地寻找食物以
供养这老也吃不饱的孩子，直到它完全长大，不再需要别人给它准备食物为止。有一次
我算了一下，一条幼虫长大过程中所需要吃的蝇加起来达八十二只之多。
我有时候怀疑，这种蜂为什么不像其它蜂那样预先储藏好食物，把洞封好，自己也
就可以离开，何必老耐着性子守在洞口呢？可能是因为它捕回的死蝇不能藏得太久的缘
故吧。可是它又为什么不像黄翅蜂一样把蝇麻痹，而是把它杀死呢？我推测可能是因为
蝇毕竟和毛毛虫、蟋蟀不大一样，它是那样的轻，那样的软，放不了多久就会缩得没有
了。所以这东西必须吃新鲜的，否则就没有什么价值了。另一个原因可能是蝇非常灵敏，
必须擒得快，不像那呆头呆脑的毛毛虫和庞然大物似的蟋蟀，目标明显，动作又不灵便，
让母蜂有充分的时间去麻痹它们。捕蝇蜂在必要的时候，须随时应用它的爪子、嘴巴或
刺，这样捕捉来的蝇当然不能随心所欲地让它半死不活了，要么让蝇逃掉了，要么捉个
死蝇。母蜂当然选择后者。
要观察到捕蝇蜂袭击苍蝇可不是件容易的事。因为它总是在离巢很远的地方捕捉的。
可是“踏破铁蹄无觅处，’得来全不费功夫”，有一次，我就无意之中看到了这精彩的
一幕，饱了眼福。那天我张着伞坐在烈日下。享受着伞的阴影的不只我一个，还有各种
马蝇也躲在我的伞下休息。它们平静地歇在张着的伞顶上。我在伞下没有事情做，就欣
赏着它们大大的金色眼睛来消磨时间。那些眼睛在我的伞下闪闪发光，好像一颗颗宝石。
有时候伞的某一部分被晒得太热了，它们就不得不转移阵地，移到没有被太阳光晒到的
那部分。我很喜欢看它们这种严肃的动作。
有一天，我正在伞下打瞌睡，突然，“梆！”的一声，张着的伞忽然像皮鼓似地被
击了一下。
“怎么回事？”我一下子清醒了许多。
大概是一颗榆树的果子掉到伞上了吧！我想。
可是“梆——梆——梆！”一声接着一声地传来。是哪个爱搞恶作剧的家伙把种子
或石子往我伞上扔？我离开了我的伞荫，四处巡视了一下，什么也没有。那声音又响起
来了，我抬头往伞顶一望，原来如此！我终于明白是怎么回事了。附近的捕蝇蜂发现我
这里有这么多肥美的食物，都飞过来捕取猎物，一切都像我希望的那样进行着，我只需
静静地坐着观看便是了。
每隔十五分钟左右，就有一只捕蝇蜂飞进来，直向伞顶冲去，发出一声重击。于是
战争就在伞顶上展开了。那是多么精彩和紧张啊，大家打得难分难解，不分上下，使你
辨不清谁是袭击者谁是自卫者。不过这种争执并没有维持很久。不一会儿，捕蝇蜂就用
双腿夹着它的俘虏飞走了。奇怪的是，这愚蠢的蝇群还不肯离开这危险的场所——的确，
外面实在太热了，与其被晒死，还不如在里面“今朝有酒今朝醉”，先尽情享受再说。
现在，让我来观察这只带着战利品回去的蜂吧。当它接近自己家的时候，突然发出
一种尖锐的嗡嗡声，听来颇有点凄凉的意味，好像十分不安。这声音一直继续着，直到
它降落到地面上为止。它先在地面上方盘旋了一会，然后又小心翼翼地降落，如果它那
敏锐的眼睛发现了一些什么不正常的情形，它就要降低下落的速度，在上面盘旋几秒钟，
飞上去又飞下来，然后像一支箭一般地飞开去了。稍过一会儿，我们就知道它为什么犹
豫不决了。不到一会儿，它又回来了。这次它先在高处巡视一遍，然后慢慢降落到地上
某一点——这一点在我看来实在没什么特别之处。
我想它大概是随便降落在这一点上的，降落之后，它还得慢慢地寻找自己巢的入口。
可是后来我发现自己又低估了捕蝇蜂。它恰好不偏不倚地降落在自己的巢上。它把前面
的沙扒开一些，再用头一顶，便顺利地拖着它的猎物进巢了。它进去后，旁边的沙粒立
刻又堆上洞口把它堵住。这和我从前所看到的无数次捕蝇蜂回巢的情形一样。我常常惊
异于蜂类为什么毫不犹豫地找到它的巢的入口，虽然那入口处和旁边的地方完全一样，
没有任何可以辨别的记号。
捕蝇蜂回巢的时候，并不是每次要在空中盘旋许多时候，它之所以要这样不停盘旋
是因为它看到自己的巢被一种巨大的危险所笼罩。它那种凄凉的嗡嗡声是表示它内心的
忧愁和恐惧。在没有危险的时候，它绝对不会发出这种声音。那么它的敌人是谁呢？原
来是一只小小的蝇，外表看上去十分的软弱无能。而这捕蝇蜂，它虽是蝇类的天敌，大
马蝇的刽子手，但当它发现自己被这种小蝇所监视的时候，它竟然会吓得不敢进洞去。
而事实上那只小蝇小得像一个不够它的幼虫吃一顿的侏儒。
这情形似乎像猫怕老鼠一样让人费解。为什么捕蝇蜂不冲下去把这个无耻的小蝇赶
走呢？这个我无法解释。也许这其貌不扬的小蝇自有它的厉害之处，在茫茫宇宙中占有
相当的地位，就像许多凶猛的动物一样。大自然的规则常常是我们人类不能了解的。
我以后还会讲到这种蝇把卵产在捕蝇蜂放在巢内的猎物上。它的幼虫孵出来以后就
掠夺了捕蝇蜂幼虫的养料。如果食物不够的话，它们还会毫不手软地把捕蝇蜂的幼虫当
作美食吃掉。所以，它决不是一种微不足道的小蝇，而是一个无情的杀手。捕蝇蜂那么
害怕它，不是没有道理的。那么这种小蝇是怎么把卵产在捕蝇蜂的卵上的呢？这非常值
得研究。
它从不走近捕蝇蜂的巢，只是耐着性子等着捕蝇蜂拖着丰美的猎物回来。当捕蝇蜂
半个身体钻进洞穴的时候，它就冲下去附在那只死的马蝇身上，当捕蝇蜂艰难地拖着马
蝇进洞的时候，这种小蝇就以惊人的速度在马蝇身上产下一个卵，有时候会连续地产下
两三个。捕蝇蜂从前半身钻进洞到完全把猎物拖进洞，前后只不过是一瞬间的功夫，可
就是在这一瞬间的时间内，小蝇已经完成了它的任务。现在它可以在洞旁的阳光里蹲下，
为它的第二次偷袭作准备了。
平时总有三四个这样的小蝇同时出现在一个巢的附近。对于进巢的入口它们往往知
道得比较清楚。它们那暗红的肤色、大而红的眼睛、以及它们惊人的耐心，常常使我想
到绑票的情形。那些歹徒穿着黑色的衣服，头上包着红布，静静地躲在阴暗的角落里，
等候机会来拦住过路的客人。
那可怜的捕蝇蜂是因为看到了这种歹徒在家门口才犯踌躇的。它知道那帮歹徒一定
会干坏事的。但是最后它还是飞进了自己的家。于是这些小歹徒似的蝇们便飞起来紧紧
跟着它，它向前，它们也向前；它后退，它们也后退；它无法使它们离开。最后它终于
撑不住了，不得不歇歇脚，那些小歹徒也跟着歇下来，但仍然虎视耽耽地跟在背后。于
是捕蝇蜂又飞起来，带着一声愤怒的呜咽声。这些无耻的小歹徒仍然厚着脸皮紧迫不舍。
捕蝇蜂只好另外想了一个办法，另找了一条路，以比较高的速度飞行，希望它的敌人跟
不上它而最终迷失方向。没想到那伙小歹徒早已料到这一招，折回到洞口等它回来。果
然，不一会儿，以为已经摆脱了危险的捕蝇蜂们回来了，这帮小歹徒赶紧起身直追。母
蜂的耐心已经没有了，最后终于被它们找着了产卵的机会。
好在我们刚才所讲的那只捕蝇蜂没有遭到这种不幸，所以让我们来结束这一章吧。
蜂的幼虫吃着母亲留给它们的粮食，慢慢地长大。过了两个星期，它就开始做茧了。可
是在它身体内没有足够的丝，所以它必须掺入沙粒以增加它的硬度。它把残余的食物堆
积到小屋的一角，先把地面扫清，然后在墙和墙之间搭起白色的美丽的丝来。它先把丝
攀成一个网，然后开始第二步工作。
它在网的中央做一个吊床。这吊床好像一个袋子，一端封闭，一端留有小孔。捕蝇
蜂的幼虫半个身体伸在床外，用嘴巴一粒一粒地挑选沙粒，太大的沙粒它看不上眼，会
一下子把它丢开。选好后，它再把沙一粒一粒地衔进去，很均匀地铺在吊床袋的四周，
就像泥水匠把石子嵌入灰泥一样。
到现在为止，茧子的一端还是开着口的，它必须把它封上。它用丝织成一顶帽子，
大小恰巧能盖住茧子的开口处，在这上面，也嵌进一粒一粒的沙。现在茧子可以说做完
了。不过捕蝇蜂在茧里还要做一番修葺工作。它在墙上涂一层浆液，为的是避免让自己
柔嫩的皮肤被沙粒擦伤或者蹭破。在这之后它就可以安心地睡大觉了。不久它将变成一
只成年的捕蝇蜂，就像它的母亲一样。
------------------








寄生虫



在八九月里，我们应该到光秃秃的、被太阳灼得发烫的山峡边去看看，让我们找一
个正对太阳的斜坡，那儿往往热得烫手，因为太阳已经把它快烤焦了。恰恰是这种温度
像火炉一般的地方，正是我们观察的目标。因为就是在这种地方，我们可以得到很大的
收获。这一带热土，往往是黄蜂和蜜蜂的乐土。它们往往在地下的土堆里忙着料理食物
——这里堆上一堆象鼻虫、蝗虫或蜘蛛，那里一组组分列着蝇类和毛毛虫类，还有的正
在把蜜贮藏在皮袋里、土罐里、棉袋里或是树叶编的瓮里。
在这些默默地埋头苦干的蜜蜂和黄蜂中间，还夹杂着一些别的虫，那些我们称之为
寄生虫。它们匆匆忙忙地从这个家赶到那个家，耐心地躲在门口守候着，你别以为它们
是在拜访好友，它们这些鬼鬼祟祟的行为决不是出于好意，它们是要找一个机会去牺牲
别人，以便安置自己的家。
这有点类似于我们人类世界的争斗。劳苦的人们，刚刚辛辛苦苦地为儿女积蓄了一
笔财产，却碰到一些不劳而获的家伙来争夺这笔财产。有时还会发生谋杀、抢劫、绑票
之类的恶性事件，充满了罪恶和贪婪。至于劳动者的家庭，劳动者们曾为它付出了多少
心血，贮藏了多少他们自己舍不得吃的食物，最终也被那伙强盗活活吞灭了。世界上几
乎每天都有这样的事情发生，可以说，哪里有人类，哪里就有罪恶。昆虫世界也是这样，
只要存在着懒惰和无能的虫类，就会有把别人的财产占为已有的罪恶。蜜蜂的幼虫们都
被母亲安置在四周紧闭的小屋里，或呆在丝织的茧子里，为的是可以静静地睡一个长觉，
直到它们变为成虫。可是这些宏伟的蓝图往往不能实现，敌人自有办法攻进这四面不通
的堡垒。每个敌人都有它特殊的战略——那些绝妙又狠毒的技巧，你根本连想都想不到。
你看，一只奇异的虫，靠着一根针，把它自己的卵放到一条蛰伏着的幼虫旁边——这幼
虫本是这里真正的主人；或是一条极小的虫，边爬边滑地溜进了人家的巢，于是，蛰伏
着的主人格永远长睡不醒了，因为这条小虫立刻要把它吃掉了。那些手段毒辣的强盗，
毫无愧意地把人家的巢和茧子作为自己的巢和茧子，到了来年，善良的女主人已经被谋
杀，抢了巢杀了主人的恶棍倒出世了。
看看这一个，身上长着红白黑相间的条纹，形状像一只难看而多毛的蚂蚁，它一步
一步地仔细地考察着一个斜坡，巡查着每一个角落，还用它的触须在地面上试探着。你
如果看到它，一定会以为它是一只粗大强壮的蚂蚁，只不过它的服装要比普通的蚂蚁漂
亮。这是一种没有翅膀的黄蜂，它是其它许多蜂类的幼虫的天敌。它虽然没有翅膀，可
是它有一把短剑，或者说是一根利刺。只见它踯躅了一会儿，在某个地方停下来，开始
挖和扒，最后居然挖出了一个地下巢穴，就跟经验丰富的盗墓贼似的。这巢在地面上并
没有痕迹，但这家伙能看到我们人类所看不到的东西。它钻到洞里停留了一会儿，最后
又重新在洞口出现。这一去一来之间，它已经干下了无耻的勾当：它潜进了别人的茧子，
把卵产在那睡得正酣的幼虫的旁边，等它的卵孵化成幼虫，就会把茧子的主人当作丰美
的食物。
这里是另外一种虫，满身闪耀着金色的、绿色的、蓝色的和紫色的光芒。它们是昆
虫世界里的蜂雀，被称作金蜂，你看到它的模样，决不会相信它是盗贼或是搞谋杀的凶
手。可它们的确是用别的蜂的幼虫作食物的昆虫，是个罪大恶极的坏蛋。
这十恶不赦的金蜂并不懂得挖人家墙角的方法，所以只得等到母蜂回家的时候溜进
去。你看，一只半绿半粉红的金蜂大摇大摆地走进一个捕蝇蜂的巢。那时，正值母亲带
着一些新鲜的食物来看孩子们。于是，这个“侏儒”就堂而皇之地进了“巨人”的家。
它一直大摇大摆在走到洞的底端，对捕蝇蜂锐利的刺和强有力的嘴巴似乎丝毫没有惧意。
至于那母蜂，不知道是不是不了解金蜂的丑恶行径和名声，还是给吓呆了，竟任它自由
进去。来年，如果我们挖开捕蝇蜂的巢看看，就可以看到几个赤褐色的针箍形的茧子，
开口处有一个扁平的盖。在这个丝织的摇篮里，躺着的是金蜂的幼虫。至于那个一手造
就这坚固摇篮的捕蝇蜂的幼虫呢？它已完全消失了。只剩下了一些破碎的皮屑了。它是
怎么消失的？当然是被金蜂的幼虫吃掉了！
看看这个外貌漂亮而内心奸恶的金蜂，它身上穿着金青色的外衣，腹部缠着“青铜”
和“黄金”织成的袍子，尾部系着一条蓝色的丝带。当一只泥匠蜂筑好了一座弯形的巢，
把入口封闭，等里面的幼虫渐渐成长，把食物吃完后，吐着丝装饰着它的屋子的时候，
金蜂就在巢外等候机会了。一条细细的裂缝，或是水泥中的一个小孔，都足以让金蜂把
它的卵塞进泥匠蜂的巢里去。总之，到了五月底，泥匠蜂的巢里又有了一个针箍形的茧
子，从这个茧子里出来的，又是一个口边沾满无辜者的鲜血的金蜂，而泥匠蜂的幼虫，
早被金蜂当作美食吃掉了。
正如我们所知道的那样，蝇类总是扮演强盗或小偷或歹徒的角色。虽然它们看上去
很弱小，有时候甚至你用手指轻轻一撞，就可以把它们全部压死。可它们的确祸害不小。
有一种小蝇，身上长满了柔软的绒毛，娇软无比，只要你轻轻一摸就会把它压得粉身碎
骨，它们脆弱得像一丝雪片，可是当它们飞起来时有着惊人的速度。乍一看，只是一个
迅速移动的小点儿。它在空中徘徊着，翅膀震动得飞快，使你看不出它在运动，倒觉得
是静止的。好像是被一根看不见的线吊在空中。如果你稍微动一下，它就突然不见了。
你会以为它飞到别处去了，怎么找都没有。它到哪儿去了呢？其实，它哪儿都没去，它
就在你身边。当你以为它真的不见了的时候，它早就回到原来的地方了。它飞行的速度
是如此之快，使你根本看不清它运动的轨迹，那么它又在空中干什么呢？它正在打坏主
意，在等待机会把自己的卵放在别人预备好的食物上。我现在还不能断定它的幼虫所需
要的是哪一种食物：蜜、猎物，还是其它昆虫的幼虫？
有一种灰白色的小蝇，我对它比较了解，它蜷伏在日光下的沙地上，等待着抢劫的
机会。当各种蜂类猎食回来，有的衔着一只马蝇，有的衔着一只蜜蜂，有的衔着一只甲
虫，还有的衔着一只蝗虫。大家都满载而归的时候，灰蝇就上来了，一会儿向前，一会
儿向后，一会儿又打着转，总是紧跟着蜂，不让它从自己的眼皮底下溜走。当母蜂把猎
物夹在腿间拖到洞里去的时候，它们也准备行动了。就在猎物将要全部进洞的那一刻，
它们飞快地飞上去停在猎物的末端，产下了卵。就在那一眨眼的工夫里，它们以迅雷不
及掩耳之势完成了任务。母蜂还没有把猎物拖进洞的时候，猎物已带着新来的不速之客
的种子了，这些“坏种子”变成虫子后，将要把这猎物当作成长所需的食物，而让洞的
主人的孩子们活活饿死。
不过，退一步想，对于这种专门掠夺人家的食物吃人家的孩子来养活自己的蝇类，
我们也不必对它们过于指责。一个懒汉吃别人的东西，那是可耻的，我们会称他为“寄
生虫”，因为它牺牲了同类来养活自己，可昆虫从来不做这样的事情。它从来不掠取其
同类的食物，昆虫中的寄生虫掠夺的都是其它种类昆虫的食物，所以跟我们所说的“懒
汉”还是有区别的。你还记得泥匠蜂吗？没有一只泥匠蜂会去沾染一下邻居所隐藏的蜜，
除非邻居已经死了，或者已经搬到别处去很久了。其它的蜜蜂和黄蜂也一样。所以，昆
虫中的“寄生虫”要比人类中的“寄生虫”要高尚得多。
我们所说的昆虫的寄生，其实是一种“行猎”行为。例如那没有翅膀，长得跟蚂蚁
似的那种蜂，它用别的蜂的幼虫喂自己的孩子，就像别的蜂用毛毛虫、甲虫喂自己的孩
子一样。一切东西都可以成为猎手或盗贼，就看你从怎样的角度去看待它。其实，我们
人类是最大的猎手和最大的盗贼。他们偷吃了小牛的牛奶，偷吃了蜜蜂的蜂蜜，就像灰
蝇掠夺蜂类幼虫的食物一样。人类这样做是为了抚育自己的孩子。自古以来人类不也总
是想方设法地把自己的孩子拉扯大，而且往往不择手段吗——这不是很像灰蝇吗？
------------------








新陈代谢的工作者



有许多昆虫，它们在这世界上做着极有价值的工作，尽管它们从来没有因此而得到
相应的报酬和相称的头衔。当你走近一只死鼹鼠，看见蚂蚁、甲虫和蝇类聚集在它身上
的时候，你可能会全身起鸡皮疙瘩，拔腿就跑。你一定会以为它们都是可怕而肮脏的昆
虫，令人恶心。事实并不是这样的，它们正在忙碌着为这个世界做清除工作。让我们来
观察一下其中的几只蝇吧，我们就可以知道它们的所作所为是多么的有益于人类，有益
于整个自然界了。
你一定看见过碧蝇吧？也就是我们通常所说的“绿头苍蝇”。它们有着漂亮的金绿
色的外套，发着金属般的光彩，它们还有一对红色的大眼睛。
当它们嗅出在很远的地方有死动物的时候，会立即赶过去在那里产卵。几天以后，
你会惊讶地发现那动物的尸体变成了液体，里面有几千条头尖尖的小虫子，你一定会觉
得这种方法实在有点令人反胃，可是除此之外，还有什么别的更好更容易的方法消灭腐
烂发臭的动物的尸体，让它们分解成元素被泥土吸收而再为别的生物提供养料呢？是谁
能够使死动物的尸体奇迹般地消失，变成一滩液体的呢？正是碧蝇的幼虫。
如果这尸体没有经过碧蝇幼虫的处理，它也会渐渐地风干，这样的话，要经过很长
一段时间才会消失。碧蝇和其它蝇类的幼虫一样，有一种惊人的本事，那就是能使固体
物质变成液体物质。有一次我做了一个试验，把一块煮得很老的蛋白扔给碧蝇作食物，
它马上就把这块蛋白变成一滩像清水一样的液体。而这种使它能够把固体变成液体的东
西，是它嘴里吐出来的一种酵母素，就好像我们胃里的胃液能把食物消化一样。碧蝇的
幼虫就靠着这种自己亲手制作的肉汤来维持自己的生命。
其实，能做这种工作的，除了碧蝇之外，还有灰肉蝇和另一种大的肉蝇。你常常可
以看到这种蝇在玻璃窗上嗡嗡飞着。千万不要让它停在你要吃的东西上面，要不然的话，
它会使你的食物也变得充满细菌了。不过你可不必像对待蚊子一样，毫不客气地去拍死
它们，只要把它们赶出去就行了。因为在房间外面，它们可是大自然的功臣。它们以最
快的速度，用曾经活过的动物的尸体产生新的生命，它们使尸体变成一种无机物质被土
壤吸收，使我们的土壤变得肥沃，从而形成新一轮的良性循环。
------------------








松毛虫



在我那个园子里，种着几棵松树。每年毛毛虫都会到这松树上来做巢，松叶都快被
它们吃光了。为了保护我们的松树，每年冬天我不得不用长叉把它们的巢毁掉，搞得我
疲惫不堪。
你这贪吃的小毛虫，不是我不客气，是你太放肆了。如果我不赶走你，你就要喧宾
夺主了。我将再也听不到满载着针叶的松树在风中低声谈话了。不过我突然对你产生了
兴趣，所以，我要和你订一个合同，我要你把你一生的传奇故事告诉我，一年、两年，
或者更多年，直到我知道你全部的故事为止。而我呢，在这期间不来打扰你，任凭你来
占据我的松树。
订合同的结果是，不久我们就在离门不远的地方，拥有了三十几只松毛虫的巢。天
天看着这一堆毛毛虫在眼前爬来爬去，使我不禁对松毛虫的故事更有了一种急切了解的
欲望。这种松毛虫也叫作“列队虫”，因为它们总是一只跟着一只，排着队出去。
下面我开始讲它的故事：
第一，先要讲到它的卵。在八月份的前半个月，如果我们去观察松树的枝端，一定
可以看到在暗绿的松叶中，到处点缀着一个个白色的小圆柱。每一个小圆柱，就是一个
母亲所生的一簇卵。这种小圆柱好像小小的手电筒，大的约有一寸长，五分之一或六分
之一寸宽，裹在一对对松针的根部。这小筒的外貌，有点像丝织品，白里略透一点红，
小筒的上面叠着一层层鳞片，就跟屋顶上的瓦片似的。
这鳞片软得像天鹅绒，很细致地一层一层盖在筒上，做成一个屋顶，保护着筒里的
卵。没有一滴露水能透过这层屋顶渗进去。这种柔软的绒毛是哪里来的呢？是松毛虫妈
妈一点一点地铺上去的。它为了孩子牺牲了自己身上的一部分毛。它用自己的毛给它的
卵做了一件温暖的外套。
如果你用钳子把鳞片似的绒毛刮掉，那么你就可以看到盖在下面的卵了，好像一颗
颗白色珐琅质的小珠。每一个圆柱里大约有三百颗卵，都属于同一个母亲。这可真是一
个大家庭啊！它们排列得很好看，好像一颗玉蜀黍的穗。无论是谁，年老的或年幼的，
有学问的还是没文化的，看到松蛾这美丽精巧的“穗”，都会禁不住喊道：“真好看啊！”
多么光荣而伟大母亲啊！
最让我们感兴趣的东西，不是那美丽的珐琅质的小珠本身，而是那种有规则的几何
图形的排列方法。一只小小的蛾知道这精妙的几何知识，这难道不是一件令人惊讶的事
吗？但是我们愈和大自然接触，便愈会相信大自然里的一切都是按照一定的规则安排的。
比如，为什么一种花瓣的曲线有一定的规则？为什么甲虫的翅鞘上有着那么精美的花纹？
从庞然大物到微乎其微的小生命，一切都安排得这样完美，这是不是偶然的呢？似乎不
大可能吧？是谁在主宰这个世界呢？我想冥冥之中一定有一位“美”的主宰者在有条不
紊地安排着这个缤纷的世界。我只能这样解释了。
松蛾的卵在九月里孵化。在那时候，如果你把那小筒的鳞片稍稍掀起一些，就可以
看到里面有许多黑色的小头。它们在咬着，推着它们的盖子，慢慢地爬到小筒上面，它
们的身体是淡黄色的，黑色的脑袋有身体的两倍那么大。它们爬出来后，第一件事情就
是吃支持着自己的巢的那些针叶，把针叶啃完后，它们就落到附近的针叶上。常常可能
会有三四个小虫恰巧落在一起，那么，它们会自然地排成一个小队。这便是未来大军的
松毛虫雏形。如果你去逗它们玩，它们会摇摆起头部和前半身，高兴地和你打招呼。
第二步工作就是在巢的附近做一个帐篷。这帐篷其实是一个用薄绸做成的小球，由
几片叶子支持着。在一天最热的时候，它们便躲在帐里休息，到下午凉快的时候才出来
觅食。
你看松毛虫从卵里孵化出来还不到一个小时，却已经会做许多工作了：吃针叶、排
队和搭帐篷，仿佛没出娘胎就已经学会了似的。
二十四小时后，帐篷已经像一个榛仁那么大。两星期后，就有一个苹果那么大了。
不过这毕竟是一个暂时的夏令营。冬天快到的时候，它们就要造一个更大更结实的帐篷。
它们边造边吃着帐篷范围以内的针叶。也就是说，它们的帐篷同时解决了它们的吃住问
题。这的确是一个一举两得的好办法。这样它们就可以不必特意到帐外去觅食。因为它
们还很小，如果贸然跑到帐外，是很容易碰到危险的。
当它们把支持帐篷的树叶都吃完了以后，帐篷就要塌了。于是，像那些择水草而居
的阿拉伯人一样，全家会搬到一个新的地方去安居乐业。在松树的高处，它们又筑起了
一个新的帐篷。它们就这样辗转迁徒着，有时候竟能到达松树的顶端。
也就是这时候，松毛虫改变了它们的服装。它们的背上面上了六个红色的小圆斑，
小圆斑周围环绕着红色和绯红色的刚毛。红斑的中间又分布着金色的小斑。而身体两边
和腹部的毛都是白色的。
到了十一月，它们开始在松树的高处，木枝的顶端筑起冬季帐篷来。它们用丝织的
网把附近的松叶都网起来。树叶和丝合成的建筑材料能增加建筑物的坚固性。全部完工
的时候，这帐篷的大小相当于半加仑的容积，它的形状像一个蛋。巢的中央是一根乳白
色的极粗的丝带，中间还夹杂着绿色的松叶。顶上有许多圆孔，是巢的门，毛毛虫们就
从这里爬进爬出。在矗立在帐外的松叶的顶端有一个用丝线结成的网，下面是一个阳台。
松毛虫常聚集在这儿晒太阳。它们晒太阳的时候，像叠罗汉似的堆成一堆，上面张着的
丝线用来减弱太阳光的强度，使它们不至于被太阳晒得过热。
松毛虫的巢里并不是一个整洁的地方，这里面满是杂物的碎屑，毛虫们蜕下来的皮、
以及其它各种垃圾，真的可以称作是“败絮其中”。
松毛虫整夜歇在巢里，早晨十点左右出来，到阳台上集合，大家堆在一起，在太阳
底下打吨。它们就这样消磨掉整个白天。它们会时不时地摇摆着头以表示它们的快乐和
舒适。到傍晚六七点钟光景，这班瞌睡虫都醒了，各自从门口回到自己家里。
它们一面走一面嘴上吐着丝。所以无论走到哪里，它们的巢总是愈变愈大，愈来愈
坚固。它们在吐着丝的时候还会把一些松叶掺杂进去加固。每天晚上总有两个小时左右
的时间做这项工作。它们早已忘记夏天了，只知道冬天快要来了，所以每一条松毛虫都
抱着愉快而紧张的心情工作着，它们似乎在说：
“松树在寒风里摇摆着它那带霜的枝桠的时候，我们将彼此拥抱着睡在这温暖的巢
里！多么幸福啊！让我们满怀希望，为将来的幸福努力工作吧！”
不错，亲爱的毛毛虫们，我们人类也和你们一样，为了求得未来的平静和舒适而孜
孜不倦地劳动。让我们怀着希望努力工作吧！你们为你们的冬眠而工作，它能使你们从
幼虫变为蛾；我们为我们最后的安息而工作，它能消灭生命，同时创造出新的生命。让
我们一起努力工作吧！
做完了一天的工作，就是它们的用餐时间了。它们都从巢里钻出来，爬到巢下面的
针叶上去用餐。它们都穿着红色的外衣，一堆堆地停在绿色的针叶上，树枝都被它们压
得微微向下弯了。多么美妙的一副图画啊！这些食客们都静静地安详地咬着松叶，它们
那宽大的黑色的额头在我的灯笼下发着光。它们都要吃到深夜才肯罢休。回到巢里后还
要继续工作一会儿。当最后一批松毛虫进巢的时候，大约已是深夜一二点钟了。
松毛虫所吃的松叶通常只有三种，如果拿其它的常绿树的叶子给它们吃，即使那些
叶子的香味足以引起食欲，可松毛虫是宁可饿死也不愿尝一下的。这似乎没什么好说的，
松毛虫的胃和人的胃有着相同的特点。
松毛虫们在松树上走来走去的时候，随路吐着丝，织着丝带，回去的时候就依照丝
带所指引的路线。有时候它们找不到自己的丝带而找了别的松毛虫的丝带，那样它就会
走入一个陌生的巢里。但是没有关系，巢里的主人和这不速之客之间丝毫不会引起争执。
大家似乎都习以为常，平静得跟什么事都没有发生一样。到了睡觉的时候，大家也就像
兄弟一般睡在一起了，谁都没有一点生疏的感觉。不论是主人还是客人，大家都依旧在
限定的时间里工作，使它们的巢更大、更厚。由于这类意外的事情常有发生，所以有几
个巢总能接纳“外来人员”为自己的巢添砖加瓦，它们的巢就显得比其它的巢大了不少。
“人人为我，我为人人”是它们的信条，每一条毛毛虫都尽力地吐着丝，使巢增大增厚，
不管那是自己的巢还是别人的巢。事实上，正是因为这样才扩大了总体上的劳动成果。
如果每个松毛虫都只筑自己的巢，宁死也不愿替别家卖命，结果会怎样？我敢说，一定
会一事无成，谁也造不了又大又厚的巢。因此它们是几百几百地一起工作的，每一条小
小的松毛虫，都尽了它自己应尽的一份力量，这样团结一致才造就了一个个属于大家的
堡垒，一个又大又厚又暖和的大棉袋。每条松毛虫为自己工作的过程也是为其它松毛虫
工作的过程，而其它松毛虫也相当于都在为它工作。多么幸福的松毛虫啊，它们不知道
什么私有财产和一切争斗的根源。

毛虫队
有一个老故事，说是有一只羊，被人从船上扔到了海里，于是其余的羊也跟着跳下
海去。“因为羊有一种天性，那就是它们永远要跟着头一只羊，不管走到哪里。就因为
这，亚里士多德曾批评羊是世界上最愚蠢、最可笑的动物。”那个讲故事的人这样说。
松毛虫也具有这种天性，而且比羊还要强烈。第一只到什么地方去，其余的都会依
次跟着去，排成一条整齐的队伍，中间不留一点空隙。它们总是排成单行，后一只的须
触到前一只的尾。为首的那只，无论它怎样打转和歪歪斜斜地走，后面的都会照它的样
子做，无一例外。第一只毛毛虫一面走一面吐出一根丝，第二只毛虫踏着第一只松毛虫
吐出的丝前进，同时自己也吐出一条丝加在第一条丝上，后面的毛毛虫都依次效仿，所
以当队伍走完后，就有一条很宽的丝带在太阳下放着耀眼的光彩。这是一种很奢侈的筑
路方法。我们人类筑路的时候，用碎石铺在路上，然后用极重的蒸汽滚筒将它们压平，
又粗又硬但非常简便。而松毛虫，却用柔软的缎子来筑路，又软又滑但花费也大。
这样的奢侈有什么意义吗？它们为什么不能像别的虫子那样免掉这种豪华的设备，
简朴地过一生呢？我替它们总结出两条理由：松毛虫出去觅食的时间是在晚上，而它们
必须经过曲曲折折的道路。它们要从一根树枝爬到另一根树枝上，要从针叶尖上爬到细
枝上，再从细枝爬到粗枝上。如果它们没有留下丝线作路标，那么它们很难找回自己的
家，这是最基本的一条理由。
有时候，在白天它们也要排着队作长距离的远征，可能经过三十码左右的长距离。
它们这次可不是去找食物，而是去旅行，去看看世界，或者去找一个地方，作为它们将
来蛰伏的场所。因为在变成蛾子之前，它们还要经过一个蛰伏期。在作这样长途旅行的
时候，丝线这样的路标是不可缺少的。
在树上找食物的时候，它们或许是分散在各处，或许是集体活动，反正只要有丝线
作路标，它们就可以整齐一致地回到巢里。要集合的时候，大家就依照着丝线的路径，
从四面八方匆匆聚集到大队伍中来。所以这丝带不仅仅是一条路，而且是使一个大团体
中各个分子行动一致的一条绳索。这便是第二个理由。
每一队总有一个领头的松毛虫，无论是长的队还是短的队。它为什么能做领袖则完
全出自偶然，没有谁指定，也没有公众选举，今天你做，明天它做，没有一定的规则。，
毛虫队里发生的每一次变故常常会导致次序的重新排列。比如说，如果队伍突然在行进
过程中散乱了，那么重新排好队后，可能是另一只松毛虫成了领袖。尽管每一位“领袖”
都是暂时的、随机的，但一旦作了领袖，它就摆出领袖的样子，承担起一个领袖应尽的
责任。当其余的松毛虫都紧紧地跟着队伍前进的时候，这位领袖趁队伍调整的间隙摇摆
着自己的上身，好像在做什么运动。又好像在调整自己——毕竟，从平民到领袖，可是
一个不小的飞跃，它得明确自己的责任，不能和刚才一样，只需跟在别人后面就行了，
当它自己前进的同时，它就不停地探头探脑地寻找路径。它真是在察看地势吗？它是不
是要选一个最好的地方？还是它突然找不到引路的丝线，所以犯了疑？看着它那又黑又
亮，活像一滴柏油似的小脑袋，我实在很难推测它真的在想什么？我只能根据它的一举
一动，作一些简单的联想。我想它的这些动作是帮助它辨出哪些地方粗糙，哪些地方光
滑，哪些地方有尘埃，哪些地方走不过去。当然，最主要的是辨出那条丝带朝着哪个方
向延伸。
松毛虫的队伍长短不一，相差悬殊，我所看到的最长的队伍有十二码或十三码，其
中包含二百多只松毛虫，排’成极为精致的波纹形的曲线，浩浩荡荡的，最短的队伍一
共只有两条松毛虫，它们仍然遵从原则，一个紧跟在另一只的后面。
有一次我决定要和我养在松树上的松毛虫开一次玩笑，我要用它们的丝替它们铺一
条路，让它们依照我所设想的路线走。既然它们只会不假思索地跟着别人走，那么如果
我把这路线设计成一个既没有始点也没有终点的圆，它们会不会在这条路上不停地打转
转呢？
一个偶然的发现帮助我实现了这个计划。在我的院子里有几个栽棕树的大花盆，盆
的圆周大约有一码半长。松毛虫们平时很喜欢爬到盆口的边沿，而那边沿恰好是一个现
成的圆周。
有一天，我看到很大一群毛虫爬到花盆上，渐渐地来到它们最为得意的盆沿上。慢
慢地，这一队毛虫陆陆续续到达了盆沿，在盆沿上前进着。我等待并期盼着队伍形成一
个封闭的环，也就是说，等第一只毛虫绕过一因而回到它出发的地方。一刻钟之后，这
个目的达到了。现在有整整一圈的松毛虫在绕着盆沿走了。’第二步工作是，必须把还
要上来的松毛虫赶开，否则它们会提醒原来盆沿上的那围虫走错了路线，从而扰乱实验。
要使它们不走上盆沿，必须把从地上到花盆间的丝拿走。于是我就把还要继续上去的毛
虫拨开，然后用刷子把丝线轻轻刷去，这相当于截断了它们的通道。这样下面的虫子再
也上不去，上面的再也找不到回去的路。这一切准备就绪后，我们就可以看到一幕有趣
的景象在眼前展开了：
一群毛虫在花盆沿上一圈一圈地转着，现在它们中间已经没有领袖了。因为这是一
个封闭的圆周，不分起点和终点，谁都可以算领袖，谁又都不是领袖，可它们自己并不
知道这一点。
丝织和轨道越来越粗了，因为每条松毛虫都不断地把自己的丝加上去。除了这条圆
周路之外，再也没有别的什么叉路了，看样子它们会这样无止境地一圈一圈绕着走，直
到累死为止？
旧派的学者都喜欢引用这样一个故事：“有一头驴子，它被安放在两捆干草中间，
结果它竟然饿死了。因为它决定不出应该先吃哪一捆。”其实现实中的驴子不比别的动
物愚蠢，它舍不得放弃任何一捆的时候，会把两捆一起吃掉。我的毛虫会不会表现得聪
明一点呢？它们会离开这封闭的路线吗？我想它们一定会的。我安慰自己说：
“这队伍可能会继续走一段时间，一个钟头或两个钟头吧。然后，到某个时刻，毛
毛虫自己就会发现这个错误，离开那个可怕的骗人的圈子，找到一条下来的路。”
而事实上，我那乐观的设想错了，我太高估了我的毛毛虫们了。如果说这些毛虫会
不顾饥饿，不顾自己一直回不到巢，只要没有东西阻挠它们，它们就会一直在那儿打圈
子，那么它们就蠢得令人难以置信了。然而，事实上，它们的确有这么蠢。
松毛虫们继续着它们的行进，接连走了好几个钟头。到了黄昏时分，队伍就走走停
停，它们走累了。当天气逐渐转冷时，它们也逐渐放慢了行进的速度。到了晚上十点钟
左右，它们继续在走，但脚步明显慢了下来，好像只是懒洋洋地摇摆着身体。进餐的时
候到了，别的毛虫都成群结队地走出来吃松叶。可是花盆上的虫子们还在坚持不懈地走。
它们一定以为马上可以到目的地和同伴们一起进晚餐了。走了十个钟头，它们一定又累
又饿，食欲极好。一棵松树离它们不过几寸远，它们只要从花盆上下来，就可以到达松
树，美美地吃上一顿松叶了。但这些可怜的家伙已经成了自己吐的丝的奴隶了，它们实
在离不开它，它们一定像看到了海市蜃楼一样，总以为马上可以到达目的地，而事实上
还远着呢！十点半的时候，我终于没有耐心了，离开它们去睡我的觉。我想在晚上的时
候它们可能清醒些。可是第二天早晨，等我再去看它们的时候，它们还是像昨天那样排
着队，但队伍是停着的。晚上太冷了，它们都蜷起身子取暖，停止了前进。等空气渐渐
暖和起来后，它们恢复了知觉，又开始在那儿兜圈子了。
第三天，一切还都像第二天一样。这天夜里非常冷，可怜的毛虫又受了一夜的苦。
我发现它们在花盆沿分成两堆，谁也不想再排队。它们彼此紧紧地挨在一起，为的是可
以暖和些。现在它们分成了两队，按理说每队该有一个自己的领袖了，可以不必跟着别
人走，各自开辟一条生路了。我真为它们感到高兴。看到它们那又黑又大的脑袋迷茫地
向左右试探的样子，我想不久以后它们就可以摆脱这个可怕的圈子了。可是不久我发现
自己又错了。当这两支分开的队伍相逢的时候，又合成一个封闭的圆圈，于是它们又开
始了整天兜圈子，丝毫没有意识到错过了一个绝佳的逃生机会。
后来的一个晚上还是很冷。这些松毛虫又都挤成了一堆，有许多毛虫被挤到丝织轨
道的两边，第二天一觉醒来，发现自己在轨道外面，就跟着轨道外的一个领袖走，这个
领袖正在往花盆里面爬。这队离开轨道的冒险家一共有七位，而其余的毛虫并没有注意
它们，仍然在兜圈子。
到达花盆里的毛虫发现那里并没有食物，于是只好垂头丧气地依照丝线指示的原路
回到了队伍里，冒险失败了。如果当初选择的冒险道路是朝着花盆外面而不是里面的活，
情形就截然不同了。
一天又过去了，这以后又过了一天。第六天是很暖和的。我发现有几个勇敢的领袖，
它们热得实在受不住了，于是用后脚站在花盆最外的边沿上，做着要向空中跳出去的姿
势。最后，其中的一只决定冒一次险，它从花盆沿上溜下来，可是还没到一半，它的勇
气便消失了，又回到花盆上，和同胞们共甘苦。这时盆沿上的毛虫队已不再是一个完整
的圆圈，而是在某处断开了。也正是因为有了一个唯一的领袖，才有了一条新的出路。
两天以后，也就是这个实验的第八天，由于新道路的开辟，它们已开始从盆沿上往下爬，
到日落的时候，最后一只松毛虫也回到了盆脚下的巢里。
我计算了一下，它们一共走了四十八个小时。绕着圆圈走过的路程在四分之一公里
以上。只有在晚上寒冷的时候，队伍才没有了秩序，使它们离开轨道，几乎安全到达家
里。可怜无知的松毛虫啊！有人总喜欢说动物是有理解力的，可是在它们身上，我实在
看不出这个优点。不过，它们最终还是回到了家，而没有活活饿死在花盆沿上，说明它
们还是有点头脑的。

松毛虫能预测气候
在正月里，松毛虫会脱第二次皮。它不再像以前那么美丽了，不过有失也有得，它
添了一种很有用的器官。现在它背部中央的毛变成暗淡的红色了。由于中央还夹杂着白
色的长毛，所以看上去颜色更淡了。这件褪了色的衣服有一个特点，那就是在背上有八
条裂缝，像口子一般，可以随毛虫的意图自由开闭。当这种裂缝开着的时候，我们可以
看到每只口子里有一个小小的“瘤”。这玩意儿非常的灵敏，稍稍有一些动静它就消失
了。这些特别的口子和“瘤”有什么用处呢？当然不是用来呼吸的，因为没有一种动物
——即便是一条松毛虫，也不会从背上呼吸的。让我们来想想松毛虫的习性，或许我们
可以发现这些器官的作用。
冬天和晚上的时候，是松毛虫们最活跃的时候，但是如果北风刮得太猛烈的话，天
气冷得太厉害，而且会下雨下雪或是雾厚得结成了冰屑，在这样的天气里，松毛虫总会
谨慎地呆在家里，躲在那雨水不能穿透的帐篷下面。
松毛虫们最怕坏天气，一滴雨就能使它们发抖，一片雪花就能惹起它们的怒火。如
果能预先料到这种坏天气。那么对松毛虫的日常生活是非常有意义的。在黑夜里，这样
一支庞大的队伍到相当远的地方去觅食，如果遇到坏天气，那实在是一件危险的事。如
果突然遭到风雨的袭击，那么松毛虫就要遭殃了，而这样的不幸在坏的季节里是常常会
发生的。可松毛虫们自有办法。让我来告诉你它们是怎样预测天气的吧。
有一天，我的几个朋友，和我一起到院子里看毛虫队的夜游。我们等到九点钟，就
进入到院子里。可是……可是……这是怎么了？巢外一只毛虫都没有！就在昨天晚上和
前天晚上还有许多毛虫出来呢，今天怎么会一只都没有了？它们都上哪儿去了？是集体
出游吗？还是遇到了灭顶之灾？我们等到十点、十一点，一直到半夜。失望之余，我只
得送我的朋友走了。
第二天，我发现那天晚上竟然下了雨，直到早晨还继续下着，而且山上还有积雪。
我脑子里突然闪过一个念头，是不是毛虫对天气的变化比我们谁都灵敏呢？它们昨晚没
有出来，是不是因为早已预料到天气要变坏，所以不愿意出来冒险？一定是这样的！我
为自己的想法暗暗喝彩，不过我想我还得仔细观察它们。
我发现每当报纸上预告气压来临的时候，比如说暴风雨将要来临的时候，我的松毛
虫总躲在巢里。虽然它们的巢暴露在坏天气中，可风啊、雨啊、雪啊、寒冷啊，都不能
影响它们。有时候它们能预报雨天以后的风暴。它们这种推测天气的天赋，不久就得到
我们全家的承认和信任。每当我们要进城去买东西的时候，前一天晚上总要先去征求一
下松毛虫们的意见，我们第二天去还是不去，完全取决于这个晚上松毛虫的举动，它成
了我们家的“小小气象预报员”。
所以，想到它的小孔，我推测松毛虫的第二套服装似乎给了它一个预测天气的本领。
这种本领很可能是与那些能自由开闭的口子息息相关。它们时时张开，取一些空气作为
样品，放到里面检验一番，如果从这空气里测出将有暴风雨来临，便立刻发出警告。

松蛾
三月到来的时候，松毛虫们纷纷离开巢所在的那棵松树，作最后一次旅行。三月二
十日那天，我花了整整一个早晨，观察了一队三码长，包括一百多只毛虫在内的毛虫队。
它们衣服的颜色已经很淡了。队伍很艰难地徐徐地前进着，爬过高低不平的地面后，就
分成了两队，成为两支互不相关的队伍，各分东西。
它们目前有极为重要的事情要做。队伍行进了两小时光景，到达一个墙角下，那里
的泥土又松又软，极容易钻洞。为首的那条松毛虫一面探测，一面稍稍地挖一下泥土，
似乎在测定泥土的性质。其余的松毛虫对领袖百分之一百的服从，因此只是盲目地跟从
着它，全盘接受领袖的一切决定，也不管自己喜欢不喜欢。最后，领头的松毛虫终于找
到了一处它自己挺喜欢的地方，于是停下脚步。接着其余的松毛虫都走出队伍，成为乱
哄哄的一群虫子，仿佛接到了“自由活动”的命令，再也不要规规矩矩地排队了。所有
的虫子的背部都杂乱地摇摆着，所有的脚都不停地靶着，所有的嘴巴都挖着泥土，渐渐
地它们终于挖出了安葬自己的洞。到某个时候，打过地道的泥土裂开了，就把它们埋在
里面。于是一切都又恢复平静了。现在，毛虫们是葬在离地面三寸的地方，准备着织它
们的茧子。
两星期后，我往地面下挖土，又找到了它们。它们被包在小小的白色丝袋里，丝袋
外面还沾染着泥土。有时候，由于泥土土质的关系，它们甚至能把自己埋到九寸以下的
深处。
可是那柔软的、翅膀脆弱而触须柔软的蛾子是怎么从下面上来到达地面的呢？它一
直要到七八月才出来。那时候，由于风吹雨打，日晒雨淋，泥土早已变得很硬了。没有
一只蛾子能够冲出那坚硬的泥土，除非它有特殊的工具，并且它的身体形状必须很简单。
我弄了一些茧子放到实验室的试管里，以便看得更仔细些。我发现松娥在钻出茧子的时
候，有一个蓄势待发的姿势，就像短跑运动员起跑前的下蹲姿势一样。它们把它美丽的
衣服卷成一捆，自己缩成一个圆底的圆柱形，它的翅膀紧贴在脚前，像一条围巾一般，
它的触须还没有张开，于是把它们弯向后方，紧贴在身体的两旁。它身上的毛发向后躺
平，只有腿是可以自由活动的，为的是可以帮助身体钻出泥土。
虽然有了这些准备，但对于挖洞来说，还远远不够，它们还有更厉害的法宝呢！如
果你用指尖在它头上摸一下，你就会发现有几道很深的皱纹。我把它放在放大镜下，发
现那是很硬的鳞片。在额头中部顶上的鳞片是所有鳞片中最硬的。这多象一个回旋钻的
钻头呀。在我的试管里，我看到蛾子用头轻轻地这边撞撞，那边碰碰，想把沙块钻穿。
到第二天，它们就能钻出一条十寸长的隧道通到地面上来了。
最后，蛾子终于到达了泥土外面，只见它缓缓地展开它的翅膀，伸展它的触须，蓬
松一下它的毛发。现在它已完全打扮好了，完全是一只漂亮成熟又自由自在的蛾子了。
尽管它不是所有蛾子中最美丽的一种，但它的确已经够漂亮了。你看，它的前翅是灰色
的，上面嵌着几条棕色的曲线，后翅是白色的，腹部盖着淡红色的绒毛。颈部围着小小
的鳞片，又因为这些鳞片挤得很紧密，所以看上去就像是一整片，非常像一套华丽的盔
甲。
关于这鳞片，还有些极为有趣的事情。如果我们用针尖去刺激这些鳞片，无论我们
的动作多么轻微，立刻会有无数的鳞片飞扬起来。这种鳞片就是松蛾用来做盛卵的小筒
用的，我们在这一章的开头已经讲过了。
------------------








卷心菜毛虫



卷心菜几乎可以说是我们所有的蔬菜中最为古老的一种，我们知道古时候的人就已
经开始吃它了。而实际上在人类开始吃它之前，它已经在地球上存在了很久很久，所以
我们实在是无法知道究竟是什么时候，它出现了，又是什么时候人类第一次种植它们，
用的又是什么方法。植物学家告诉我们，它最初是一种长茎、小叶、长在滨海悬崖的野
生植物。历史对于这类细小的事情的记载是不愿多花笔墨的。它所歌颂的，是那些夺去
千万人生命的战场，它觉得那一片使人类生生不息的土地是没有研究价值的。它详细列
举各国国王的嗜好和怪癖，却不能告诉我们小麦的起源！但愿将来的历史记载会改变它
的作风。
我们对于卷心菜知道得实在太少了，那实在有点可惜，它的确算得上一种很贵重的
东西。因为它拥有许多有趣的故事。不仅是人类，就是别的动物也都与它有千丝万缕的
联系。其中有一种普通大白蝴蝶的毛虫，就是靠卷心菜生长的。它们吃卷心菜皮及其一
切和卷心菜相似的植物叶子，像花椰菜、白菜芽、大头菜，以及瑞典萝卜等，似乎生来
就与这种样子的莱类有不解之缘。
它们还吃其它一些和卷心菜同类的植物。它们都属于十字花科——植物学家们这样
称呼它们，因为它们的花有四瓣，排成十字形。白蝴蝶的卵一般只产在这类植物上。可
是它们怎么知道这是十字花科植物呢？它们又没有学过植物学。这倒是个谜。我研究植
物和花草已有五十多年，但如果要我判定一种没有开花的植物是不是属于十字花科，我
只能去查书。现在我不需要去一一查书了，我会根据白蝴蝶留下的记号作出判断——我
是很信任它的。
白蝴蝶每年要成熟两次。一次是在四五月里，一次是在十月，这正是我们这里卷心
菜成熟的时候。白蝴蝶的日历恰巧和园丁的日历一样。当我们有卷心菜吃的时候，白蝴
蝶也快要出来了。
白蝴蝶的卵是淡橘黄色的，聚成一片，有时候产在叶子朝阳的一面，有时候产在叶
子背着阳光的一面。大约一星期后，卵就变成了毛虫，毛虫出来后第一件事就是把这卵
壳吃掉。我不止一次看到幼虫自己会把卵壳吃掉，不知道这是什么意思。我的推测是这
样的：卷心菜的叶片上有蜡，滑得很，为了要使自己走路的时候不至于滑倒，它必须弄
一些细丝来攀缠住自己的脚，而要做出丝来，需要一种特殊的食物。所以它要把卵壳吃
掉，因为那是一种和丝性质相似的物质，在这初生的小虫胃里，它比较容易转化成小虫
所需要的丝。
不久，小虫就要尝尝绿色植物了。卷心菜的灾难也就此开始了。它们的胃口多好啊！
我从一颗最大的卷心菜上采来一大把叶子去喂我养在实验室的一群幼虫，可是两个小时
后，除了叶子中央粗大的叶脉之外，已经什么都不剩了。照这样的速度吃起来，这一片
卷心菜田没多少日子就会被吃完了。
这些贪吃的小毛虫，除了偶尔有一些伸胳膊挪腿的休息动作外，什么都不做，就知
道吃。当几只毛虫并排地在一起吃叶子的时候，你有时候可以看见它们的头一起活泼地
抬起来，又一起活泼地低下去。就这样一次一次重复着做，动作非常整齐，好像普鲁士
士兵在操练一样。我不知道它们这种动作是什么含义，是表示它们在必要的时候有作战
能力呢，还是表示它们在阳光下吃食物很快乐？总之，在它们成为极肥的毛虫之前，这
是它们唯一的练习。
吃了整整一个月之后，它们终于吃够了。于是就开始往各个方向爬。一面爬，一面
把前身仰起，作出在空中探索的样子，似乎是在做伸展运动，为了帮助消化和吸收吧。
现在气候已经开始转冷了，所以我把我的毛虫客人们都安置在花房里，让花房的门开着。
可是，令我惊讶的是，有一天我发现，这群毛虫都不见了。
后来我在附近各处的墙脚下发现了它们。那里离花房差不多有三十码的距离。它们
都栖在屋檐下，那里可以作为它们冬季的居所了。卷心菜的毛虫长得非常壮实健康，应
该不十分怕冷。
就在这居所里，它们织起茧子，变成蛹。来年春天，就有蛾从这里飞出来了。
听着这卷心菜毛虫的故事，我们也许会感到非常有趣。可是如果我们任凭它大量繁
殖，那么我们很快就会没有卷心菜吃了。所以当我们听说有一种昆虫，专门猎取卷心莱
毛虫，我们并不感到痛惜。因为这样可以使它们不至于繁殖得太快。如果卷心菜毛虫是
我们的敌人，那么那种卷心菜的敌人就是我们的朋友了。但它们长得那样细小，又都喜
欢埋头默默无闻地工作，使得园丁们非但不认识它，甚至连听都没听说过它，即使他偶
然看到它在它所保护的植物周围徘徊，他也决不会注意它，更不会想到它对自己会有那
么大的贡献。
我现在要给这小小的侏儒们一些应得的奖赏。
因为它长得细小，所以科学家们称它为“小侏儒”，那么让我也这么称呼它吧，我
实在不知道它还有什么别的好听一点的名字。
它是怎样工作的呢？让我们来看看吧。春季，如果我们走到菜园里去，一定可以看
见，在墙上或篱笆脚下的枯草上，有许多黄色的小茧子，聚集成一堆一堆的，每堆有一
个榛仁那么大。每一堆的旁边都有一条毛虫，有时候是死的，看上去大都很不完整，这
些小茧子就是“小侏儒”的工作成果，它们是吃了可怜的毛虫之后才长大的，那些毛虫
的残尸，也是“小侏儒”们剥下的。
这种“小侏儒”比幼虫还要小。当卷心菜毛虫在菜上产下橘黄色的卵后，“小侏儒”
的蛾就立刻赶去，靠着自己坚硬的钢毛的帮助，把自己的卵产在卷心菜毛虫的卵膜表面
上。一只毛虫的卵里，往往可以有好几个“小侏儒”，跑去产卵。照它们卵的大小来看，
一只毛虫差不多相当于六十五只“小侏儒”。
当这毛虫长大后，它似乎并不感到痛苦。它照常吃着菜叶，照常出去游历，寻找适
宜做茧子的场所。它甚至还能开展工作，但是它显得非常萎靡、非常无力，经常无精打
彩的，渐渐地消瘦下去，最后，终于死去。那是当然的，有那么一大群“小侏儒”在它
身上吸血呢！毛虫们尽职地活着，直到身体里的“小侏儒”准备出来的时候。它们从毛
虫的身体里出来后就开始织茧，最后，变成蛾，破茧而出。
------------------








孔雀蛾



孔雀蛾是一种长得很漂亮的蛾。它们中最大的来自欧洲，全身披着红棕色的绒毛，
脖子上有一个白色的领结，翅膀上洒着灰色和褐色的小点儿。横贯中间的是一条淡淡的
锯齿形的线，翅膀周围有一圈灰白色的边，中央有一个大眼睛，有黑得发亮的瞳孔和许
多色彩镶成的眼帘，包括黑色、白色、栗色和紫色的弧形线条。这种蛾是由一种长得极
为漂亮的毛虫变来的，它们的身体以黄色为底色，上面嵌着蓝色的珠子。它们靠吃杏叶
为生。
五月六日的早晨，在我的昆虫实验室里的桌子上，我看着一只雌的孔雀蛾从茧子里
钻出来。我马上把它罩在一个金属丝做的钟罩里。我这么做没有别的什么目的，只是一
种习惯而已。我总是喜欢搜集一些新鲜的事物，把它们放到透明的钟罩里细细欣赏。
后来我很为自己的这种方法庆幸。因为我获得了意想不到的收获，在晚上九点钟左
右，当大家都准备上床睡觉的时候，隔壁的房间里突然发出很大的声响。
小保罗衣服都没穿好，在屋里奔来跑去，疯狂地跳着、顿着足、敲着椅子。我听到
他在叫我：
“快来快来！”他喊道，“快来看这些蛾子，像鸟一样大，满房间都是！”
我赶紧跑进去一看，孩子的话一点儿也不夸张。房间里的确充满了那种大蛾子，已
经有四只被捉住关在笼子里了，其余的拍打着翅膀在天花板下面翱翔。
看到这情形，我立即想起那只早上被我关起来的囚徒。
“快穿好衣服”，我对儿子说，“把鸟笼放下，跟我来。我们立刻就要看到更有趣
的事情了。”
我们立刻下楼，来到我的书房，那在整个房子的右侧。我发现厨房里的仆人已被这
突然发生的事件吓慌了，她用她的围裙扑打着这些大蛾，起初她还以为它们是蝙蝠呢。
这样看来，孔雀蛾们已经占据了我家里的每一部分，惊动了家里的每一个人。
我们点着蜡烛走进书房，书房的一扇窗开着。我们看到了难忘的一幕情景：那些大
蛾子轻轻地拍着翅膀，绕着那钟罩飞来飞去。一会儿飞上，一会儿飞下，一会儿飞出去，
一会儿又飞回来，一会儿冲到天花板上，一会儿又俯冲下来。它们向蜡烛扑来，用翅膀
把它扑灭。它们停在我们的肩上，扯我们的衣服，咬我们的脸。小保罗紧紧地握着我的
手，努力保持镇定。
一共有多少蛾子？这个房间里大约有二十只，加上别的房间里的，至少在四十只以
上。四十个情人来向这位那天早晨才出生的新娘致敬——这位关在象牙塔里的公主！
在那一个星期里，每天晚上这些大蛾总要来朝见它们美丽的公主。那时候正是暴风
雨的季节，晚上黑得伸手不见五指。我们的屋子又被遮蔽在许多大树后面，很难找到。
它们经过这么黑暗和艰难的路程，历尽困苦来见它们的女王。
在这样恶劣的天气条件下，连那凶狠强壮的猫头鹰都不敢轻易离开巢，可孔雀蛾却
能果断地飞出来，而且不受树枝的阻挡，顺利到达目的地。它们是那样的无畏，那样的
执着，以至于到达目的地的时候，它身上没有一个地方被刮伤，哪怕是细微的小伤口也
没有。这个黑夜对它们来说，如同大白天一般。
孔雀蛾一生中唯一的目的就是找配偶，为了这一目标，它们继承了一种很特别的天
赋：不管路途多么远，路上怎样黑暗，途中有多少障碍，它总能找到它的对象。在它们
的一生中大概有两三个晚上它们可以每晚花费几个小时去找它们的对象。如果在这期间
它们找不到对象。那么它的一生也将结束了。
孔雀蛾不懂得吃。当许多别的蛾成群结队地在花园里飞来飞去吮吸蜜汁的时候，它
从不会想到吃东西这回事。这样，它的寿命当然是不会长的了，只不过是两三天的时间，
只来得及找一个伴侣而已。
------------------








找枯露菌的甲虫



在讲到甲虫之前，让我先来讲一下我的狗朋友，它会找枯露菌。所谓枯露菌，指的
是一种长在地底下的蘑菇。狗常常被用来做这种工作。我的狗的运气极好，有好几次跟
着一只在这方面极有经验的狗一同出去工作。而那只狗，那位我急于见识一下的找蘑菇
专家，其外貌实在没有什么可取的地方：一只极为普通的狗，态度平静而从容，又丑又
不整洁。总之，是一种你绝对不会让它歇在你的火炉边的狗，可它的的确确是一个名符
其实的找蘑菇专家。许多道理和人类世界中的一样：天才和贫穷总是连在一起的。
这只狗的主人，是村里有名的枯露菌商。他起初怀疑我要窥探他的秘密从而和他进
行商业竞争，后来有人告诉他我只是要采集地下植物的标本，要借他的狗用一用，他相
信了，于是就允许我和他一同出发去工作。
我们有言在先，规定任何一方都不可以干涉狗的行为，而且只要它发现一种菌类，
不管那是人们喜欢吃的蘑菇还是其它不可以吃的蘑菇，它都理应得到一片面包作为酬劳。
另外，它的主人不能禁止他的狗到它喜欢的地方去，即使他深知那个地方的蘑菇绝对卖
不出去。因为对我的研究课题而言，这种蘑菇是否可吃并不重要，我的目的和枯露菌商
人的有所不同。
遵循这个原则，这次远征获得了极大的成功。这只忙碌的狗一路上踱着慢慢的步子
走着，用鼻子嗅着。每走几步，它就要停下来，用鼻子测试着泥土。它用鼻子扒几下土，
然后泰然地望着主人，似乎在说：
“就在这里了！就在这里了！我以我的名义担保！这里有蘑菇！”
它说的果然没错。主人依着它所指示的方向掘下去。万一主人的铲掘得偏了，它会
赶紧发出一声鼻音，提醒主人如何把铲子放到正确的部位。这样掘下去没有一次落空。
狗的鼻子果然名不虚传，从不说谎，不过它指示给我们的，包括各种地下菌类：大的小
的，有气味的，没气味的……当我收集着这些蘑菇的时候，我非常惊奇，这里面几乎包
括了附近一带所有的地下蘑菇品种。
是不是那种我们常说的嗅觉在帮助狗找寻呢？我不相信，如果完全靠嗅觉，它决不
能找出这许多气味完全不同的菌类来。它一定有一种我们所没有的感觉。通常我们用人
类的标准去推测一切事物的时候，往往容易犯错。在这个世界上，有许多种感觉是我们
人类所不知道的。而这种感觉，在昆虫中更加明显。
我们现在就来讲讲找蘑菇的甲虫吧。
这是一种美丽的甲虫，小小的，黑黑的，有一个白绒肚皮，形状是圆的，像一粒樱
桃的核，当它用翅膀的边缘擦着腹部的时候，就会发出一种柔软的“唧唧”声，就像小
鸟看见母亲带着食物回来时所发出的声音一样。雄的甲虫头上还长着一个美丽的角。
我是在一个长满蘑菇的松树林里发现这种甲虫的。那是一个美丽的地方，在秋季气
候温和的日子里，我们全家都喜欢到那儿去玩。那个地方几乎什么都有：有用细枝做成
的老喜鹊的巢；有饶舌的柽鸟吃饱后在树上互相嬉戏、追逐；有兔子突然从树丛里跳出
来，翘着它们的短尾巴；那里还有可爱的小河，可以给孩子们筑隧道，这种河很容易堆
成一排小屋，我们用草盖在屋顶上，算是草屋，一段芦草插上屋顶便成为烟囱。微风轻
轻地掠过松针，发出轻轻的叹息声，我们就在这美妙的音乐中开始了我们的午餐。
的确，对于小孩子们来说，这是个真正的乐园。即使是成人，也会喜欢这个地方。
我到这里来最大的乐趣便是守候那些找蘑菇的甲虫们，它们的洞到处都可以看到，而且
门是开着的，不过在洞口堆着一堆疏松的泥土。洞大约有几寸深，一直向下，而且往往
筑在比较松的泥土中。当我用小刀一直挖下去的时候，我总是发现这种洞是空的，甲虫
们已经乘着夜色离开这里了，它在这里做完了它的工作后，便迁到别处去了。这种甲虫
是个流浪者，并且是个夜行客，随便什么时候，它想离开这个洞的时候，它就能很容易
地在别处另筑新巢。有时候我侥幸能在洞底发现甲虫，但永远只有一个，或是雌的或是
雄的，从不会成对。看来这种洞并不是一个家庭的所在地，而是专门给独身的甲虫住的。
你看，这洞里的甲虫正在啃着一个小蘑菇，已经吃完了一部分。它虽然已经累了，
但仍旧紧紧地抱着它，它是决不肯轻易放弃这个蘑菇的，这是它的宝贝，是它一生中的
最爱，从周围许多吃剩的碎片来看，这只甲虫已经吃得饱饱的了。
当我从它手中夺过这宝物的时候，我发现这是一种很小的地下菌，跟枯露菌很相像。
这个事实似乎可以解释甲虫的习惯和它常要换新居的理由。让我们想象一下吧，在
静静的黄昏中，这个小旅行家便从它的洞里慢慢地跟着步走出来，一边快活地唱着歌，
一边悠闲地散着步。它仔细地检查着土地，探究这地底下所埋的东西，就像狗找枯露菌
一样。它的嗅觉告诉它哪个地方有菌，只不过盖着几寸泥土而已；那个地方虽然泥土肥
沃，但地底下决不会有菌类。当它判定在某一点下面有菌的时候，便一直往下挖，结果
总能得到它的食物。它挖的洞也成了它的临时宿舍，在食物没有吃完之前，它是不会离
开洞的，它会在自己掘的洞底快活地吃着，管它洞门是开着的还是关着的。
等到洞里的食物都吃完后，它就要搬家了。它会在别处找一个适当的地方，再掘下
去，然后住一阵子，吃一阵子，等到新屋里的食物吃完了，它就再搬一次家。在整个秋
季到来年的春季——菌类的生长季节里，它就这样游历着，“打一枪换一个地方”，从
一个洞搬到另一个洞，很辛苦又很洒脱。
甲虫所猎取的菌并没有特殊的气味，那么它是怎么能从地面上检查出地底下菌类的
存在呢？它是聪明的甲虫，它自有办法。可是我们人类就望尘莫及了，哪怕是“千里眼”
或是“顺风耳”，也无法说出隐藏在地底下的秘密。
------------------








爱好昆虫的孩子



现在，有许多人总喜欢把一切人的品格、才能、爱好等归于遗传。也就是说承认人
类及一切动物的智慧都是从祖先那儿得来的。我并不完全同意这种观点。我现在就用我
自己的故事来证明我那喜爱昆虫的嗜好并不是从哪个先辈身上继承下来的。
我的外祖父和外祖母从来没有对昆虫产生过丝毫的兴趣和好感。关于我的外祖父，
我不大知道，我只知道他曾经历过相当苦难的日子。我敢说，如果要说他曾经和昆虫发
生过关系的话，那就是他曾一脚把它踩死。外祖母是不识字的文盲，每天为琐碎的家务
所累，没有什么闲情雅致去欣赏一些风花雪月的故事，对于科学或昆虫当然更不会产生
兴趣。当她蹲在水龙头下洗菜的时候，偶尔会发现菜叶上有一条毛虫，她会立刻把这又
讨厌又可恶的东西打掉。
关于我的祖父母，我知道的比较详细。因为我小时候，我的父母穷得无法养活我，
所以在五六岁的时候，我就跟着祖父母一同生活了。祖父母的家在偏僻的乡村里，他们
靠着几亩薄田维持生计。他们不识字，一生中从没有模过课本。祖父对于牛和羊知道得
很多，可是除此之外，其它的便什么也不知道了。如果他知道在将来他家里的一个人花
费了许多时间去研究那些渺小的、微不足道的昆虫，他会多么地吃惊啊！如果他再知道
那个疯狂的人正是坐在他旁边的小孙子，他将一定会愤怒地给我一巴掌的。
“哼，把时间和力气花费在这种没出息的东西上！”他一定会怒吼。
祖母是一个可爱的人，她整天忙着洗衣服，照顾孩子、烧饭、纺纱、看小鸭、做乳
酪和奶油，以及其它一些家务，一心为这个家操劳。
有时候，在晚上，当我们都坐在火炉边的时候，她就会常常讲一些狼的故事给我们
听。我很想见一见这匹狼，这位在一切故事里使人心惊肉跳的英雄，可是我从来没有亲
眼见过一次。亲爱的祖母，我是始终深深地感激着您的。在您的膝上，我第一次得到了
温柔的安慰，使痛苦和忧伤得到缓解。你遗传给了我强壮的体质和爱好工作的品格，可
是你的确没有给我爱好昆虫的天性。
我自己的父母都是不爱好昆虫的。母亲没有受过教育，父亲小时候虽然进过学校，
稍稍能读能写，可是为了生活整天忙得不可开交，再也没有时间顾及到别的事情了，更
谈不上爱好昆虫了。有一次当他看到我把一只虫子钉在软木上的时候，他狠狠地打了我
一拳，这就是我从他那里得到的鼓励。
尽管如此，从幼年的时候开始，我就喜欢观察和怀疑一切事物。每次忆起童年，我
总会想起一件难忘的往事，现在说起来还觉得很有趣。在我五六岁的时候，有一天我光
着脚丫子站在我们的田地前面的荒地上，粗糙的石子刺痛了我。我记得我有一块用绳子
系在腰间的手帕——很惭愧，我那时常常遗失手帕，然后用袖子代替它，所以不得不把
宝贵的手帕系在腰上。
我把脸转向太阳，那眩目的光辉使我心醉。这种光辉对我的吸引力相当于光对于任
何一只蛾子的吸引力甚至还要大得多，当我这样站着的时候，我的脑海里就突然冒出一
个问题：我究竟在用哪个器官来欣赏这灿烂的光辉？是嘴巴？还是眼睛？请读者千万不
要见笑，这的确算得上一种科学的怀疑。我把嘴也张得大大的，又把眼睛闭起来，光明
消失了；我张开眼睛闭上嘴巴：光明又出现了。这样反复试验了几次，结果都是一样。
于是我的问题被我自己解决了：我确定我看太阳用的是眼睛，后来我才知道这种方法叫
“演绎法”。这是一个多么伟大的发现啊！晚上我兴奋地把这件事告诉大家。对于我这
种幼稚和天真，只有祖母慈祥地微笑着，其余的人都大笑不止。
另外一次是在黑夜的树林里，有一种断断续续的叮当声大大地引起了我的注意。这
种声音显得分外优美而柔和。在寂静的夜里，是谁在发出这种声音？是不是巢里的小鸟
在叫？还是小虫子们在开演唱会呢？
“哦，我们快去看看吧，那很可能是一只狼。狼的确是在这种时候出声的，”同行
的人对我说，“我们一起走，但不要走得太远，声音就是从那一堆黑沉沉的木头后面发
出来的。”
我站在那里守候了许多时候，什么也没有。后来树林中发出一个轻微的响声，仿佛
是谁动了一下，接着那叮当声也消失了。第二天，第三天，我再去守候，不发现真相决
不罢休。我这种不屈不挠的精神终于获得了回报。嘿！终于抓到它了，这一个音乐家已
经在我的股掌之间了。它不是一只鸟，而是一只蚱蜢，我的同伴曾告诉我它的后腿非常
鲜美。这就是我守候了那么久所得到的微乎其微的回报。不过我所得意的，倒不是那两
只像虾肉一样鲜美的大胆，而是我又学到了一种知识，而且，这知识是我亲自通过努力
得来的。现在，从我个人的观察来看，我知道蚱蜢是会唱歌的。我没有把这发现告诉别
人，为的是怕再像上次看太阳的事情那样遭别人的嘲笑。
哦，我们屋子旁边的花长得多么美丽啊！它们好像张着一个个彩色的大眼睛向我甜
甜地笑。后来，在那个地方，我又看到一堆堆又大又红的樱桃。我尝了尝，滋味也不过
如此，没有像看上去的那么诱人，而且没有核，这究竟是些什么樱桃呢？夏天将要结束
的时候，祖父拿着铁锹来，把这块土地的泥土从底下翻起来，从地底下掘出了许许多多
圆圆的根。我认得那种根，在我们的屋子里面有许许多多，我们时时把它们放在煤炉上
煨着吃。那就是马铃薯。普通得不能再普通的马铃薯。我的探索一下戛然而止，不过，
那紫色的花和红色的果子被永远地留在了我的记忆里。
我利用自己这双对于动植物特别机警的眼睛，独自观察着一切惊异的事物。尽管那
时候我只有六岁，在别人看来什么也不懂。我研究花，研究虫子；我观察着，怀疑着；
不是受到了遗传的影响，而是受到了好奇心的驱使和对大自然的热爱。
不久我又回到了我父亲的屋子里。那时候我已经七岁，到了必须进学校的年龄了。
可我并不觉得学校生活比我以前那种自由自在地沉浸在大自然中的生活更有意思。我的
教父就是老师。那间我坐在里面学习字母的屋子，我该称它做什么呢？的确很难有一个
恰如其分的名字，因为那屋子用处太多了。它既是学校，又是厨房；既是卧室，又是餐
厅；既是鸡窝，又是猪圈。在那种时代，谁也不会梦想有王宫般美丽堂皇的学校，无论
什么破棚子都可以被认为是最理想的学校。
在这间屋子里，一张很宽的梯子通到楼上去，梯子脚边的一个凹形的房间里有一张
大床。楼上究竟有什么东西我不大知道。有时候我们会看见老师从楼上捧来一捆干草给
驴子吃；有时候也会看见他从楼上提着一篮马铃薯下来，交给师母去煮猪食。我猜这一
定是一间堆放物品的屋子，是人和畜牲共同的储藏室。
让我们回过头来讲那间做教室的房间吧。这间屋子里唯一的一扇窗，是一扇朝南的
窗，又小又矮。当你的头碰着窗顶的时候，你的肩膀同时地碰到了窗栏。这个透着阳光
的窗户是这个屋子里惟一有生气的地方，它俯视着这个村庄的大部分。从窗口往外望，
你会发现这是一个散落在斜坡上的村落。窗口下面是老师的小桌子。
对面墙上有一个壁龛，里面放着一个盛满水的发亮的铜壶，孩子们口渴的时候可以
信手从这里倒杯水解渴。在壁龛的顶端有几个架子，上面放着闪闪发光的碗，那些碗只
有在举办盛会时才拿出来用。
在光线所能射到的墙壁上，到处挂着各种色彩不协调的图画，最远的那垛墙边有一
只大壁炉，左右是用木石筑成的，上面放着塞了糠的垫褥。由两块滑动的板充当门。如
果想独自静静地躺下睡觉，你可以把门关起来。这两张床是给主人和主妇睡的。无论北
风在黑暗的谷口怎样怒吼，无论雪花在外面如何打转，他们一定在这里面睡得很舒服。
其余的地方就放着一些零碎的杂物：一条三脚凳，一只挂在壁上的盐罐，一个铁铲，重
得需要两只手一起使劲才拿得动，最后还有那风箱，就像我祖父家里的那个一样，风一
吹，炉里的木块和树枝就烧起来。我们如果要享受火炉的温暖，每人每天早上就得带一
块小柴来。
可是炉子并不是为我们生的，主要是为那三只煮猪食的锅子。老师和师母总是挑一
个最舒适的位于坐下，其余的人却围着那大锅子，围成一个半圆形。那锅里不住地冒着
热气，发出呼呼的声音。我们中胆子比较大的人会趁着先生看不见的时候用小刀挑一个
煮熟了的马铃薯，夹在他的面包里吃。我不得不承认，如果说我们在学校里做了一些工
作的话，那就是我们吃得很多。在写字的时候剥着栗子或咬着面包，似乎已经成为改不
了的习惯了。
至于我们这些年纪较小的学生，除了享受满口含着食物读书的乐趣外，还有两件快
乐的事情，在我看来不见得比栗子的味道差。我们的教室后门外就是庭院。在那里，一
群小鸡围着母鸡在扒土，小猪们自由自在地打着滚。有时候，我们中会有人偷偷地溜出
去，回来的时候故意不把门关上，于是马铃薯的香味便一阵阵地飘到门外去。外面的小
猪闻到这香味，一个个循着香味接通跑来。我的长凳子，是年纪最小的学生们坐的，恰
巧是靠着墙壁，在铜壶的下方，也正是小猪们的必经之路。我每次都能看见小猪们快步
地跑着，一边大声地呼叫着，摇着它们的小尾巴。它们用身体蹭我们的腿，把又冷又红
的鼻子拱到我们的手掌里面找吃剩的面包屑。它们那细小的圆溜溜的眼珠子望着我们，
似乎在问我们口袋里还有没有干栗子给它们吃。它们这样东窜西闻了一圈之后，就会被
怒气冲冲的老师挥着手帕赶回院子里去了。
接着就是母鸡带着它的小鸡雏们来看我们了。我们每个人都会热情地剥一些面包来
招待这些可爱的小客人，然后美滋滋地欣赏它们吃东西的样子。
在这样的一个学校里，我们能学到些什么呢？每一个年纪较小的学生手里都有——
或者说，假定他们都有一本灰纸订成的小书，上面印着字母，封面上画着一只鸽子，确
切地说，那只是一只很像鸽子的动物。封面上有一个十字架，是用字母按照一定的顺序
排出来的。老师可能觉得这本书挺管用的，因此把书发给我们，并且解释给我们听。就
因为这样，老师总是被那些年纪较大的学生们缠着，没有工夫顾及我们这些小不点儿。
他还是把书也发给我们，不过其作用只是为了让我们看上去更像学生而已。于是我们这
些小不点儿就自己坐在长凳上读书，同时请旁边的大孩子教我们——如果他能认得一两
个字母的话。我们的学习常常被一些无足轻重的小事打断，一会儿老师和师母去看锅里
的马铃薯了，一会儿小猪的同伴们叫唤着进来，一会儿又是一群小鸡忙不迭地奔进来，
我们这样忙里偷闲地看一会儿书，实在是学不到什么东西，得不到什么知识的。
大孩子们常常练习写字。他们的位置比较优越，能够借着从那狭小的窗洞透进来的
光，并且他们前面还有那张全屋唯一的桌子，学校里什么设备都没有，甚至连墨水都没
有一滴，所以每个人上学的时候都得自己带上全套的文具。那时候的墨水是装在一个长
条的纸板匣里面的，里面分成两格，上面一格放鹅毛管做的笔，下面一格是一个小小的
墨水池，里面盛着墨水，那时候的墨水是用烟灰和醋合成的。
我们的老师最伟大的工作就是修笔——然后在某一页的顶上写一行字母或是单字，
至于他写什么字，依照各个学生的要求而定。看，当老师写字的时候，他的手腕抖动得
多厉害啊！他用小拇指拈着纸，做好奋笔疾书的准备。忽然，老师的手开始运动了，在
纸上飞着、打着转。看啊，笔尖所到之处展开了一条花边，里面有圆圈、有螺旋、有花
体字、有张着翅膀的鸟……我们要他画什么他就画什么。只要你喜欢，什么都有。这些
画都是用红墨水画出来的。就是这样一只笔创造了一个个奇迹。面对这一个个奇迹，我
们都惊得目瞪口呆。
我们在学校里读些什么呢？大概是法文吧，常常从圣经上记载的历史中选一两段来
读读。拉丁语倒是学得比较多些，为的是使我们能够准确地唱赞美诗。
历史、地理呢？谁也没听到过这两个名词。地球是方的还是圆的，对我们来说有什
么不同呢？方也罢，圆也罢，反正从地里长出东西来是同样的不容易啊！
语法呢？我们的老师从来不拿这个问题去为难自己，我们当然更不会了。
数学呢？是的，我们的确学了一点，不过还不配用那么堂皇的名字，我们一直称它
为“算术”。
在星期六的晚上，通常总是用“算术”的仪式来结束这一星期。最优秀的学生先站
起来把乘法口诀表背诵一遍，然后全班，包括最小的学生，依着他的样子齐声合背一遍。
我们的声音很响亮，把偶而跑进屋来想觅一点食的鸡和猪都吓跑了。
别人都说我们的老师是个很能干的人，能把学校管理得很好。的确，他不是一个等
闲之辈，但他也确实不能称作一个好老师，因为缺少一样东西——时间。他替一个出门
的地主保管着财产；他照顾着一个极大鸽棚；他还负责指挥干草、苹果、栗子和燕麦的
收获，在夏天，我们常常帮着他干活。在那个时候，上课才是一件有趣的事，因为我们
常在干草堆上上课，有时候甚至就利用上课的时候清除鸽棚，或是消灭那些雨天从墙脚
爬出来的蜗牛。这对我来说，倒是正中下怀。
我们的老师还是个剃头匠。他那双灵巧的手替我们的抄写本装饰“花边”，也为地
方上的大人物剃头，像市长、牧师和公证人等等。
我们的老师又是个打钟的能手。每逢有婚礼或洗礼的时候，他总要到教堂里去打钟
——那时我们的功课当然要暂告停止。暴风雨来临的时候，又可以给我们一天的休假，
因为那时候必须用钟声来驱除雷电和冰雹，我们的老师责无旁贷地去敲钟了。
我们的老师还是唱诗班里的一员。我们的老师还管着村里教堂顶上的钟。那是他最
引以为豪的工作。只须对着太阳一望，他便可以说出一个准确的时间，然后他爬到教堂
顶上尖尖的阁楼里，打开一个大匣子，让自己置身于一堆齿轮和发条中间。这些东西的
秘密，除了他之外，没有第二个人知道。
这样一个学校，这样一个老师，对于我那尚未充分表现的特点，将有什么影响呢？
我那热爱昆虫的个性，几乎不得不渐渐地枯萎以至永远消失了。但是，事实上，这种个
性的种子有着很强的活力，它永远在我的血液里流动，从来没有离开过我。它能够随时
激发出来或找到滋生的养料，无时无刻不体现出来，甚至在我的教科书的封面上，也能
显而易见地看出书的主人的爱好——那里有着一只色彩配合得不很协调的鸽子，它对于
我来说，比书本里的ABC有意思得多。它的圆眼睛似乎在冲着我笑，它那翅膀我已一根一
根地数过共有多少羽毛，那些羽毛告诉我怎样飞上天空，翱翔在美丽的云朵里。这只鸽
子带着我飞到毛榉树上，我看到那些透着光泽的树干高高地矗立在长满苔藓的泥土上在
泥土上。在长着许多白色的蘑菇，看上去好像是过路的母鸡产下的蛋。这只鸽子又带我
到积雪的山顶上，在那里，鸟类用它们的红脚踏出了星形的足迹。这个鸽子是我的好伙
伴、好朋友，它减轻了我整天背字母的压力。应该谢谢它，有了它作伴，我才能静静地
坐在长凳上等候放学。
露天学校有着更大的诱惑力。当老师带着我们去消灭黄杨树下的蜗牛的时候，我却
常常阴奉阳违，不忍心杀害那些小生命。当我捉到了满手的蜗牛时，我的脚步便迟缓起
来了。它们是多么美丽啊！只要我愿意，我能捉到各种颜色的蜗牛：黄色的、淡红色的、
白色的、褐色的……上面都有深色的螺旋纹。我挑了一些最美丽的塞满衣兜，以便空闲
的时候拿出来看看。
在帮先生晒干草的日子里，我又认识了青蛙。它用自己作诱饵，引诱着河边巢里的
虾出来；在赤杨树上，我捉到了青甲虫，它的美丽使天空都为之逊色；我采下水仙花，
并且学会了用舌尖从它花冠的裂缝处吸取小滴的蜜汁，我也体验到太用力吸花蜜所导致
的头痛，不过这种不舒服与那美丽的白色花朵所带给我的赏心悦目的感觉相比，实在是
太微不足道了。我还记得这种花的漏斗的颈部有一圈美丽的红色，像挂了一串红项链。
在收集胡桃的时候，我在一块荒芜的草地上找到了蝗虫，它们的翅膀张得像一把扇
子，有红色的也有蓝色的，让人眼花缭乱。无论在什么地方，我都能源源不断地得到精
神食粮，自得其乐。我对于动植物的爱好也自然有增无减，日益弥深。
最后，这种爱好促进了我对字母的认识。由于我太喜欢封面上的鸽子，早把封面后
的字母丢到九霄云外去了。所以我的认识程度一直停留在初级阶段。一个偶然的念头使
我的父亲把我从学校里领回家去，这才是我真正读书的开始。这回读的是印得很大的字，
花了三角半钱买来的。那上面划着许多五彩的格子，每一格里画着一种动物，这些动物
就用它们的名字和第一个字母来教我认ABC。第一个就是“驴”。在法文中，它的名字是
Ane，于是我认识了A；牛的名字叫Boeuf，它教我认识了B；Canara是鸭子，于是我认识
了C；Dinod是火鸡，它教我认识了D。其余的字母也是如此这般让我认识的。当然，有几
格印得很不清楚，像那教我认得H、K和E的河马(Hippopotamus)、雨燕(Kamichi)和瘤牛
(zebu)之类，不过没什么大碍。我进步得很快，不到几天工夫，居然能很有兴趣地读那
本鸽子封面的书了。我已经被启发了，接着便懂得语法了，这激起了我对学习的浓厚兴
趣，我的父母都为我的进步感到惊异。现在我能够解释那惊人的进步的原因了：那些图
画把我引入到一群动物中，这恰巧投合了我的兴趣。我心爱的动物们开始教我念书，而
以后，动物永远成为我学习研究的对象。
后来，好运第二次降临到我身上。为了让我用功读书，我得到了一本廉价的《拉封
丹寓言》，里面有许多插图，虽然又小又不准确，可是看起来的确很有趣。这里有乌鸦、
喜鹊、青蛙、兔子、驴子、猫和狗；这些都是我所熟悉的东西，这里面的动物会走路会
讲话，因此大大激起了我的兴趣。至于了解这本书究竟讲了些什么，那是另一回事了。
不过不要担心，我试着把一个个音节连起来，慢慢地你就知道全篇的意思了。于是拉封
丹也成为我的朋友了。
十岁的时候，我已是路德士书院(RodezColloge)的学生了。我在那里成绩很好，尤
其是作文和翻译两课都能得到很高的分数。在那种古典派的气氛中，我们听到了许多神
话故事。那些故事都是很吸引人的。可是在崇拜那些英雄之余，我不会忘记趁着星期天
去看看莲香花和水仙花有没有在草地上出现；梅花雀有没有在榆树丝里孵卵；金虫是不
是在摇摆于微风中的白杨树上跳跃，无论如何，我是不能忘记它们的！
可是，忽然恶运又降临了：饥饿威胁着我们一家。父母再也没有钱供我念书了。我
不得不离开学校。生命几乎变得像地狱一样可怕。我什么都不想，只盼望能快快熬过这
段时期！
在这些悲惨的日子里，我对于昆虫的偏爱应该暂时搁在一边了吧？就像我的先辈那
样，为生计所累。但是，事实并非如此，我仍然常常能够回忆起那只第一次遇到的金虫：
它那触须上的羽毛，它那美丽的花色——褐色底子上嵌着白点——这些好像是那种凄惨
晦暗的日子里的一道闪亮的阳光，照亮并温暖了我悲伤的心。
总而言之，好运不会抛弃勇敢的人。后来我又进了在伏克罗斯(Voncluse)的初级师
范学校，在那里我能免费分到食物，尽管只是干栗子和豌豆而已，校长是位极有见识的
人，他不久便信任了我，并且给了我完全的自由。他说只要我能应付学校里的课程，我
几乎可以随心所欲做自己喜欢的事情。而当时我的程度比同班的同学要稍高一些，于是
我就利用比别人多的空闲时间来增加自己对动植物的认识。当周围的同学们都在订正背
书的错误时，我却可以在书桌的角落里观察夹竹桃的果子、金鱼草种子的壳，还有黄蜂
的刺和地甲虫们的翅膀。
我对于自然科学的兴趣，就这样慢慢地滋长起来了。在那时候，生物学是被一般学
者所轻看的学科，学校方面所承认的必修课程是拉丁文、希腊文和数学。
于是我竭尽全力地去研究高等数学。这是一种艰难的奋斗，没有老师的指导，碰到
疑难问题，往往好几天得不到解决，可我一直坚持不懈地学着，从未想过半途而废，而
终于有所成就。后来我用同样的方法自学了物理学，用一套我自己制造的简陋的仪器来
做各种实验。我违背了自己的志愿，把我的生物学书籍一直埋在箱底。
毕业后，我被派到埃杰克索书院(AjaccioCollege)去教物理和化学。那个地方离大
海不远，这对我的诱惑力实在太大了。那包蕴着无数新奇事物的海洋，那海滩上美丽的
贝壳，还有番石榴树、杨梅树和其它一些树，都足够让我研究好半天的。这乐园里美丽
的东西比起那些三角、几何定理来，吸引力大得多了。可是我努力控制着自己。我把我
的课余时间分成两部分：大部分时间用来研究数学；小部分的时间用来研究植物和搜寻
海洋里丰富的宝藏。
我们谁都不能预测未来。回顾我的一生，数学，我年轻时花费了那么多时间和精力
去钻研，结果对我却没有丝毫的用处；而动物，我竭力想方没法地回避它，在我的老年
生活中，它却成了我的慰藉。
在埃杰克索，我碰到两位著名的科学家：瑞昆(Rrguien)和莫昆·坦顿(MoquinTand
on)，瑞昆是一位著名的植物学家；而莫昆·坦顿教了我植物学的第一课。那时他因为没
有旅馆住而寄住在我的房子里。在他离开的前一天，他对我说，“你对贝壳很感兴趣，
这当然很好。不过这样还远远不够。你应当知道动物本身的组织结构，让我来指给你看
吧！这会使你对动物的认识提高到一个新的水平。”
他拿起一把很锋利的剪刀和一对针，把一个蜗牛放在一个盛水的碟子里，开始解剖
给我看。他一边解剖，一边一步步地把各部分器官解释给我听。这就是我一生中所得到
的最难以忘怀的一堂生物课，从此，当我观察动物时，不再仅仅局限在表面上了。
现在我应该把自己的故事结束了。从我的故事里可以看出，早在幼年时期，我就有
着对大自然的偏爱。而且我具有善于观察的天赋。为什么我有这种天赋？怎样才会有？
我自己也说不清楚。
无论是人还是动物，都有一种特殊的天赋：一个孩子可能有音乐的天赋，一个孩子
可能在雕塑方面很有天赋，而另一个孩子可能是速算的天才。昆虫也是这样，一种蜜蜂
生来就会剪叶子，另一种蜜蜂会造泥屋，而蜘蛛则会织网。为什么它们有这种才能？天
生就有的，除此之外就没有什么理由可解释了。在人类生活中，我们称这样的人才为
“天才”；在昆虫中，我们称这样的本领为“本能”。本能，其实就是动物的天才。
------------------








条纹蜘蛛



不管是谁，大概都不会喜欢冬季。在这个季节里，许多虫子都在冬眠。不过这并不
说明你没有什么有虫子可观察了。这时候如果有一个观察者在阳光所能照到的沙地里寻
找，或是搬开地下的石头，或是在树林里搜索，他总能找到一种非常有趣的东西，那是
一件真正的艺术品。那些有幸欣赏到这艺术作品的人真是幸福。在一年将要结束的时候，
发现这种艺术品的喜悦使我忘记了一切不快，忘记了一天比一天更糟的气候。
如果有人在野草丛里或柳树丛里搜索的话，我祝福他能找到一种神秘的东西：这是
条纹蜘蛛的巢。正像我眼前所呈现的一样。
无论从举止还是从颜色上讲，条纹蜘蛛是我所知道的蜘蛛中最完美的一种。在它那
胖胖的像榛仔仁一般大小的身体上，有着黄、黑、银三色相间的条纹，所以它的名字叫
“条纹蜘蛛”。它们的八只脚环绕在身体周围，好像车轮的辐条。
几乎什么小虫子它都爱吃。不管那是蝗虫跳跃的地方还是苍蜂盘旋的地方，是蜻蜓
跳舞的地方还是蝴蝶飞翔的地方。只要它能找到攀网的地方，它就会立刻织起网来。它
常常把网横跨在小溪的两岸，因为那种地方猎物比较丰盛。有时候它也在长着小草的斜
坡上或榆树林里织网，因为那里是蚱蜢的乐园。
它捕获猎物的武器便是那张大网，网的周围攀在附近的树枝上。它的网和别种蜘蛛
的网差不多：放射形的蛛丝从中央向四周扩散，然后在这上面连续地盘上一圈圈的螺线，
从中央一直到边缘。整张网做得非常大，而且整齐对称。
在网的下半部，有一根又粗又宽的带子，从中心开始沿着辐一曲一折，直到边缘，
这是它的作品的标记，也是它在作品中的一种签名。同时这种粗的折线也能增加网的坚
固性。
网需要做得很牢固，因为有时候猎物的份量很重，它们一挣扎，很可能会把网撑破。
而蜘蛛自己不会选择或捕捉猎物，所以只能不断地改进自己的大网以捕获更多的猎物。
它静静地坐在网的中央，把八只脚撑开，为的是能感觉到网的每一个方向的动静。摆好
阵势后，它就等候着，看命运会赐予它什么：有时候是那种微弱到无力控制自己飞行的
小虫；有时候是那种强大而鲁莽的昆虫，在做高速飞行的时候一头撞在网上，有时候它
好几天一无所获，也有时候它的食物会丰盛得好几天都吃不完。
蝗虫，尤其是一种火蝗，它控制不了自己腿部的肌肉，于是常常跌进网中。你可能
会想，蜘蛛的网一定受不住蝗虫的冲撞，因为蝗虫的个儿头要比蜘蛛大得多，只要它用
脚一蹬，立刻就可以把网蹬出一个大洞，然后逃之天天。其实情况并不是这样的，如果
在第一下挣扎之后不能逃出的话，那么它就再也没有逃生的希望了。
条纹蜘蛛并不急于吃掉蝗虫，而是用它全部的丝囊同时射出丝花，再用后腿把射出
来的丝花捆起来。它的丝囊是制造丝的器官，上面有细孔，像喷水壶的莲蓬头一般。它
的后腿比其余的腿要长，而且能张得很开，所以射出的丝能分散得很开。这样，它从腿
间射出来的丝已不是一条条单独的丝了，而是一片丝，像一把云做的扇子，有着虹霓一
般的色彩。然后它就用两条后腿很快地交替着把这种薄片，或者说是裹尸布吧，一片片
地向蝗虫撤去，就这样把蝗虫完全缠住了。
这不由得让我想起了古时候的角斗士。每逢要和强大的野兽角斗的时候，他们总是
把一个网放在自己的左肩上，当野兽扑过来时，他右手一挥，就能敏捷地把网撒开，就
像能干的渔夫撒网捕鱼那样，把野兽困在网里，再加上三叉戟一刺，就把它的性命结果
了。
蜘蛛用的也是这种方法。而且它还有一个绝招是人类所没有的：它可以把自己制造
的丝制的锁链绵绵不断地缠到蝗虫身上，一副不够，第二副立即跟着抛上来，第三副、
第四副……直到它所有的丝用完为止，而人类的网只有一副，即使有很多的话，也不能
这么迅速地接连抛出去。
当那白丝网里的囚徒决定放弃抵抗、坐以待毙的时候，蜘蛛便得意洋洋地向它走过
去，它有一个比角斗士的三叉戟还厉害的武器，那就是它的毒牙。它用它的毒牙咬住蝗
虫，美滋滋地饱餐一顿，然后回到网中央，继续等待下一个自己送上门来的猎物。

蜘蛛的巢
蜘蛛在母性方面的表露甚至比猎取食物时所显示的天才更令人叹服。它的巢是一个
丝织的袋，它的卵就产在这个袋里。它这个巢要比鸟类的巢神秘，形状像一个倒置的气
球，大小和鸽蛋差不多，底部宽大，顶部狭小，顶部是削平的，围着一圈扇蛤形的边。
整个看来，这是一个用几根丝支持着的蛋形的物体。
巢的顶部是凹形的，上面像盖着一个丝盖碗。巢的其它部分都包着一层又厚又细嫩
的白缎子，点缀着一些丝带和一些褐色或黑色的花纹。我们立刻可以猜到这一层的白缎
子的作用，它是防水的，雨水或露水都不能浸透它。
为了防止里面的卵被冻坏，仅仅使巢远离地面或藏在枯草丛里是远远不够的，还必
须有一些专门的保暖设备。让我们用剪刀把包在外面的这层防雨缎子剪开来看看。在这
下面我们发现了一层红色的丝。这层丝不是像通常那样的纤维状，而是很蓬松的一束。
这种物质，比天鹅的绒毛还要软，比冬天的火炉还要暖和，它是未来的小蜘蛛们的安乐
床。小蜘蛛们在这张舒适的床上就不会受到寒冷空气的侵袭了。
在巢的中央有一个锤子一样的袋子，袋子的底部是圆的，顶部是方的，有一个柔嫩
的盖子盖在上面。这个袋子是用非常细软的缎子做成的，里面就藏着蜘蛛的卵。蜘蛛的
卵是一种极小的橘黄色的颗粒，聚集在一块儿，拼成一颗豌豆大小的圆球。这些是蜘蛛
的宝贝，母蜘蛛必须保护着它们不受冷空气的侵袭。
那么蜘蛛是怎样造就这样精致的袋子的呢？让我们来看看它做袋子时的情形吧！它
做袋子的时候，慢慢地绕着圈子，同时放出一根丝，它的后腿把丝拉出来叠在上一个圈
的丝上面，就这样一圈圈地加上去，就织成了一个小袋子。袋子与巢之间用丝线连着，
这样使袋口可以张开。袋的大小恰好能装下全部的卵而不留一点空隙，也不知道蜘蛛妈
妈如何能掌握得那么精确。
产完卵后，蜘蛛的丝囊又要开始运作了。但这次工作和以前不同。只见它先把身体
放下，接触到某一点，然后把身体抬起来，再放下，接触到另一点，就这样一会儿在这，
一会儿在那，一会儿上，一会儿下，毫无规则，同时它的后脚拉扯着放出来的丝。这种
工作的结果，不是织出一块美丽的绸缎，而是造就一张杂乱无章、错综复杂的网。
接着它射出一种红棕色的丝，这种丝非常细软。它用后腿把丝压严实，包在巢的外
面。
然后它再一次变换材料，又放出白色的丝，包在巢的外侧，使巢的外面又多了一层
白色的外套。而且，这时候巢已经像个小气球了，上端小，下端大，接着它再放出各种
颜色不同的丝，赤色、褐色、灰色、黑色……让你目不暇接，它就用这种华丽的丝线来
装饰它的巢。直到这一步结束，整个工作才算大功告成了。
蜘蛛开着一个多么神奇的纱厂啊！靠着这个简单而永恒的工厂——它可以交替做着
搓绳、纺线、织布、织丝带等各种工作，而这里面的全部机器只是它的后腿和丝囊。它
是怎样随心所欲地变换“工种”的呢？它又是怎样随心所欲地抽出自已想要的颜色的丝
呢？我只能看到这些结果，却不知道其中的奥妙。
建巢的工作完成后，蜘蛛就头也不回地跨着慢步走开了。再也不会回来，不是它狠
心，而是它真的不需要再操心了。时间和阳光会帮助它孵卵的，而且，它也没有精力再
操心了。在替它的孩子做巢的时候，它已经把所有的丝都用光了，再也没有丝给自己张
网捕食了。况且它自己也已经没有食欲了。衰老和疲惫使它在世界上苟延残喘了几天后
安详地死去了。这便是我那匣子里的蜘蛛一生的终结，也是所有树丛里的蜘蛛的必然归
宿。

条纹蜘蛛的家族
你还记得那小小的巢里的橘黄色的卵吧？那些美丽的卵的总数有五颗之多。你还记
得它们是被密封在白缎子做的巢里的吧？那么当里面的小东西要跑出来，又冲不破白丝
做的墙的时候该怎么办呢？当时它们的母亲又不在身旁，不能帮助它们冲破丝袋，它们
是用什么办法来解决这个问题的呢？
动物在许多地方和植物有类似之处。蜘蛛的巢在我看来相当于植物的果实，只不过
它里面包含的不是种子而是卵而已，自己不能动，但它们的种子却能在很远的地方生根
发芽。因为植物有许多传播种子的方法，把它们送到四面八方：凤仙花的果实成熟的时
候，只要受到轻轻的碰触，便会裂成五瓣，每一瓣各自蜷缩起来，把种子弹到很远的地
方；还有一种很轻的种子，像蒲公英的种子，长着羽毛，风一吹就能把它们带到很远的
地方；榆树的种子是嵌在一张又宽又轻的扇子里的；槭树的种子成对地搭配，好像一双
张开的翅；柽树的种子形状像船桨，风能够让它飞到极远的地方……这些种子随遇而安，
落到不知什么地方就可以安家落户，开始下一轮的生死循环。
和植物一样，动物也有各种千奇百怪的方法凭借大自然的力量让它们的种族散布在
各地。你可以从条纹蜘蛛的身上略知一二。
三月间，正是蜘蛛的卵开始孵化的时候。如果我们用剪刀把蜘蛛的巢剪开，就可以
看到有些卵已经变成小蜘蛛爬到中央那个袋子的外面，有些仍旧是橘黄色的卵。这些刚
刚拥有生命、乳臭末干的小蜘蛛还没有披上像它们的母亲身上那样美丽的条纹衣服，它
们的背部是淡黄色的，腹部是棕色的，它们要在袋子的外面，巢的里面，呆上整整四个
月。在这段时期里，它们的身体渐渐地变得强壮丰满起来，和其它动物不同的是，它们
是在巢里而不是在外面的大天地中逐渐变为成年的蜘蛛的。
到了六七月里，这些小蜘蛛急于要冲出来了。可是它们无法在那坚硬的巢壁上挖洞。
那孩子怎么办呢？不用担心，那巢自己会裂开的，就像成熟种子的果皮一样，自动地把
后代送出来。它们一出巢，就各自爬到附近的树枝上，同时放出极为轻巧的丝来，这些
丝在空中飘浮的时候，会把它们牵引到别的地方去。关于小蜘蛛这种奇特的飞行工具，
我们在下一篇还会详细讲到。
------------------








狼蛛



蜘蛛有一个很坏的名声：大多数人都认为它是一种可怕的动物，一看到它就想把它
一脚踩死，这可能和蜘蛛狰狞的外表有关。不过一个仔细的观察家会知道，它是一个十
分勤奋的劳动者，是一个天才的纺织家，也是一个狡猾的猎人，并且在其它方面也很有
意思。所以，即使不从科学的角度看，蜘蛛也是一种值得研究的动物。但大家都说它有
毒，这便是它最大的罪名，也是大家都惧怕它的原因。不错，它的确有两颗毒牙，可以
立刻致它的猎物于死地。如果仅从这一点出发，我们的确可以说它是可怕的动物，可是
毒死一只小虫子和谋害一个人是两件迥然不同的事情。不管蜘蛛能怎样迅速地结束一只
小虫子的生命，对于人类来说，都不会有比蚊子的一刺更可怕的后果了。所以，我可以
大胆的说，大部分的蜘蛛都是无辜的，它们莫名其妙地被冤枉了。
不过，有少数种类的蜘蛛的确是有毒的。据意大利人说，狼蛛的一刺能使人痉挛而
疯狂地跳舞。要治疗这种病，除了音乐之外，再也没有别的灵丹妙药了。并且只有固定
的几首曲子治疗这种病特别灵验。这种传说听起来有点可笑，但仔细一想也有一定道理。
狼蛛的刺或许能刺激神经而使被刺的人失去常态，只有音乐能使他们镇定而恢复常态，
而剧烈地跳舞能使被刺中的人出汗，因而把毒驱赶出来。
在我们这一带，有最厉害的黑肚狼蛛，从它们身上可以得知蜘蛛的毒性有多大。我
家里养了几只狼蛛，让我把它介绍给你，并告诉你它是怎样捕食的吧！
这种狼蛛的腹部长着黑色的绒毛和褐色的条纹，腿部有一圈圈灰色和白色的斑纹。
它最喜欢住在长着百里香的干燥沙地上。我那块荒地，刚好符合这个要求，这种蜘蛛的
穴大约有二十个以上。我每次经过洞边，向里面张望的时候，总可以看到四只大眼睛。
这位隐士的四个望远镜像金钢钻一般闪着光，在地底下的四只小眼睛就不容易看到了。
狼蛛的居所大约有一尺深，一寸宽，是它们用自己的毒牙挖成的，刚刚挖的时候是
笔直的，以后才渐渐地打弯。洞的边缘有一堵矮墙，是用稻草和各种废料的碎片甚至是
一些小石头筑成的，看上去有些简陋，不仔细看还看不出来。有时候这种围墙有一寸高，
有时候却仅仅是地面上隆起的一道边。
我打算捉一只狼蛛。于是我在洞口舞动一根小穗，模仿蜜蜂的嗡嗡声。我想狼蛛听
到这声音会以为是猎物自投罗网，马上会冲出来。可是我的计划失败了。那狼蛛倒的确
往上爬了一些，想试探这到底是什么东西发出的声音，但它立刻嗅出这不是猎物而是一
个陷阱，于是一动不动地停在半途，坚决不肯出来，只是充满戒心地望着洞外。
看来要捉到这只狡猾的狼蛛，唯一的办法就是用活的蜜蜂作诱饵。于是我找了一只
瓶子，瓶子的口和洞口一样大。我把一只土蜂装在瓶子里，然后把瓶口罩在洞口上。这
强大的土蜂起先只是嗡嗡直叫，歇斯底里地撞击着这玻璃囚室，拼命想冲出这可恶的地
方。当它发现有一个洞口和自己的洞口很像的时候，便毫不犹豫地飞进去了。它实在是
愚蠢得很，走了那么一条自取灭亡的路。当它飞下去的时候，那狼蛛也正在匆匆忙忙往
上赶，于是它们在洞的拐弯处相撞了。不久我就听到了里面传来一阵死亡时的惨叫——
那只可怜的土蜂！这以后便是一段很长的沉默。我把瓶子移开，用一把钳子到洞里去探
索。我把那土蜂拖出来，它已经死了，正像刚才我所想象的那样。一幕悲剧早已在洞里
发生了。这狼蛛突然被夺走了从天而降的猎物，愣了一下，实在舍不得放弃这肥美的猎
物，急急地跟上来，于是猎物和打猎的都出洞了，我赶紧趁机用石子把洞口塞住。这狼
蛛被突如其来的变化惊呆了，一下子变得很胆怯，在那里犹豫着，不知该怎么办才好，
根本没有勇气逃走。不到一秒钟功夫，我便毫不费力地用一根草把它拔进一个纸袋里。
我就用这样的办法诱它出洞，然后捉拿归案。不久我的实验室里就有了一群狼蛛。
我用土蜂去引诱它，不仅仅是为了捉它，而且还想看看它怎样猎食。我知道它是那
种每天都要吃新鲜食物的昆虫，而不是像甲虫那样吃母亲为自己储藏的食物，或者像黄
蜂那样有奇特的麻醉术可以将猎物的新鲜程度保持到两星期之后。它是一个凶残的屠夫，
一捉到食物就将其活活地杀死，当场吃掉。
狼蛛得到它的猎物确实也不容易，也须冒很大的风险。那有着强有力的牙齿的蚱蜢
和带着毒刺的蜂随时都可能飞进它的洞去。说到武器，这两方不相上下。究竟谁更胜一
筹呢？狼蛛除了它的毒牙外没有别的武器，它不能像条纹蜘蛛那样放出丝来捆住敌人。
它唯一的办法就是扑在敌人身上，立刻把它杀死。它必须把毒牙刺入敌人最致命的地方。
虽然它的毒牙很厉害，可我不相信它在任何地方轻轻一刺而不是刺中要害就能取了敌人
的性命。

与木匠蜂作战
我已经讲过狼蛛生擒土蜂的故事，可这还不能使我满足，我还想看看它与别种昆虫
作战的情形。于是我替它挑了一种最强大的敌手，那就是木匠蜂。这种蜂周身长着黑绒
毛，翅膀上嵌着紫线，差不多有一寸长。它的刺很厉害，被它刺了以后很痛，而且会肿
起一块，很久以后才消失。我之所以知道这些，是因为曾经身受其害，被它刺过。它的
确是值得狼蛛去决一胜负的劲敌。
我捉了几只木匠蜂，把它们分别装在瓶子里。又挑了一只又大又凶猛并且饿得正慌
的狼蛛，我把瓶口罩在那只穷凶极恶的狼蛛的洞口上，那木匠蜂在玻璃囚室里发出激烈
的嗡嗡声，好像知道死期临头似的。狼蛛被惊动了从洞里爬了出来，半个身子探出洞外，
它看着眼前的景象，不敢贸然行动，只是静静地等候着。我也耐心地等候着。一刻钟过
去了，半个小时过去了，什么事都没有发生，狼蛛居然又若无其事地回到洞里去了。大
概它觉得不对头，冒然去捕食的话太危险了。我照这个样子又试探了其它几只狼蛛，我
不信每一只狼蛛都会这样面对丰盛的美食而无动于衷，于是继续一个一个的试探着，都
是这个样子，总对“天上掉下的猎物”怀有戒心。
最后，我终于成功了。有一只狼蛛猛烈地从洞里冲出来，无疑，它一定饿疯了。就
在一眨眼间，恶斗结束了，强壮的木匠蜂已经死了。凶手把毒牙刺到它身体的哪个部位
呢？是在它的头部后面。狼蛛的毒牙还咬在那里，我怀疑它真具有这种知识：它能不偏
不倚正好咬在唯一能致其于死的地方，也就是它的俘虏的神经中枢。
我做了好几次试验，发现狼蛛总是能在转眼之间干净利落地把敌人干掉，并且作战
手段都很相似。现在我明白了为什么在前几次试验中，狼蛛会只看着洞口的猎物，却迟
迟不敢出击。它的犹豫是有道理的。像这样强大的昆虫，它不能冒失鲁莽地去捉，万一
它没有击中其要害的话，那它自己就完蛋了。因为如果蜂没有被击中要害的话，至少还
可活上几个小时，在这几个小时里，它有充分的时间来回击敌人。狼蛛很知道这一点，
所以它要守在安全的洞里，等待机会，直到等到那大蜂正面对着它，头部极易被击中的
时候，它才立刻冲出去，否则决不用自己的生命去冒险。

狼蛛的毒素
让我来告诉你，狼蛛的毒素是一种多么厉害的暗器。
我做了一次试验，让一只狼蛛去咬一只羽毛刚长好的将要出巢的幼小的麻雀。麻雀
受伤了，一滴血流了出来，伤口被一个红圈圈着，一会儿又变成了紫色，而且这条腿已
经不能用了，使不上劲。小麻雀只能用单腿跳着。除此之外它好像也没什么痛苦，胃口
也很好。我的女儿同情地把苍蝇、面包和杏酱喂给它吃，这可怜的小麻雀作了我的实验
品。但我相信它不久以后一定会痊愈，很快就能恢复自由——这也是我们一家共同的愿
望和推测。十二个小时后，我们对它的伤情仍然挺乐观的。它仍然好好地吃东西，喂得
迟了它还要发脾气。可是两天以后，它不再吃东西了，羽毛零乱，身体缩成一个小球，
有时候一动不动，有时候发出一阵痉挛。我的女儿怜爱地把它捧在手里，呵着气使它温
暖。可是它痉挛得越来越厉害，次数越来越多，最后，它终于离开了这个世界。
那天的晚餐席上透着一股寒气。我从一家人的目光中看出他们对我的这种试验的无
声的抗议和责备。我知道他们一定认为我太残忍了。大家都为这只不幸的小麻雀的死而
悲伤。我自己也很懊悔：我所要知道的只是很小的一个问题，却付出了那么大的代价。
尽管如此，我还是鼓起勇气试验一只鼹鼠，它是在偷田里的莴苣时被我们捉住的，
所以即使它死于非命也不足为惜。我把它关在笼子里，用各种甲虫、蚱蜢喂它，它大口
大口贪婪地吃着，被我养得胖胖的，健康极了。
我让一只狼蛛去咬它的鼻尖。被咬过之后，它不住地用它的宽爪子挠抓着鼻子。因
为它的鼻子开始慢慢地腐烂了。从这时开始，这只大鼹鼠食欲渐渐不振，什么也不想吃，
行动迟钝，我能看出它浑身难受。到第二个晚上，它已经完全不吃东西了。大约在被咬
后三十六小时，它终于死了。笼里还剩着许多的昆虫没有被吃掉，证明它不是被饿死的，
而是被毒死的。
所以狼蛛的毒牙不止能结束昆虫的性命，对一些稍大一点的小动物来说，也是危险
可怕的。它可以致麻雀于死地，也可以使鼹鼠毙命，尽管后者的体积要比它大得多。虽
然后来我再没有做过类似的试验，但我可以说，我们千万要小心戒备，不要被它咬到，
这实在不是一件可以拿来试验的事。
现在，我们试着把这种杀死昆虫的蜘蛛和麻醉昆虫的黄蜂比较一下。蜘蛛，因为它
自己靠新鲜的猎物生活，所以它咬昆虫头部的神经中枢，使它立刻死去；而黄蜂，它要
保持食物的新鲜，为它的幼虫提供食物，因此它刺在猎物的另一个神经中枢上，使它失
去了动弹的能力。相同的是，它们都喜欢吃新鲜的食物，用的武器都是毒刺。
没有谁教它们怎样根据自己的需要分别用不同的方法对待猎物，它们生来就明白这
一点。这使我们相信冥冥之中，世界上的确有着一位万能的神在主宰着昆虫，也统治着
人类世界。

狼蛛猎食
我在实验室的泥盆里，养了好几只狼蛛。从它们那里，我看到狼蛛猎食时的详细情
形。这些做了我的俘虏的狼蛛的确很健壮。它们的身体藏在洞里，脑袋探出洞口，玻璃
般的眼睛向四周张望，腿缩在一起，作着准备跳跃的姿势，它就这样在阳光下静静地守
候着，一两个小时，不知不觉就过去了。
如果它看到一只可作猎物的昆虫在旁边经过，它就会像箭一般地跳出来，狠狠地用
它的毒牙打在猎物的头部，然后露出满意又快乐的神情，那些倒霉的蝗虫、蜻蜓和其它
许多昆虫还没有明白过来是怎么回事，就做了它的盘中美餐。它拖着猎物很快地回到洞
里，也许它觉得在自己家里用餐比较舒适吧。它的技巧以及敏捷的身手令人叹为观止。
如果猎物离它不太远，它纵身一跃就可以扑到，很少有失手的时候。但如果猎物在
很远的地方，它就会放弃，决不会特意跑出来穷追不舍。看来它不是一个贪得无厌的家
伙，不会落得一个“鸟为食亡”的下场。
从这一点可以看出狼蛛是很有耐性，也很有理性的。因为在洞里没有任何帮助它猎
食的设备，它必须始终傻傻地守候着。如果是没有恒心和耐心的虫子，一定不会这样持
之以恒，肯定没多久就退回到洞里去睡大觉了。可狼蛛不是这种昆虫。它确信，猎物今
天不来，明天一定会来；明天不来，将来也总有一天会来。在这块土地上，蝗虫、蜻蜓
之类多得很，并且它们又总是那么不谨慎，总有机会刚好跳到狼蛛近旁。所以狼蛛只需
等待时候一到，它就立刻窜上去捉住猎物，将其杀死。或是当场吃掉，或者拖回去以后
吃。
虽然狼蛛很多时候都是“等而无获”，但它的确不大会受到饥饿的威胁，因为它有
一个能节制的胃。它可以在很长一段时间内不吃东西而不感到饥饿。比如我那实验室里
的狼蛛，有时候我会连续一个星期忘了喂食，但它们看上去照样气色很好。在饿了很长
一段时间后，它们并不见得憔悴，只是变得极其贪婪，就像狠一样。
在狼蛛还年幼的时候，它还没有一个藏身的洞，不能躲在洞里“守洞待虫”，不过
它有另外一种觅食的方法。那时它也有一个灰色的身体，像别的大狼蛛一样，就是没有
黑绒腰裙——那个要到结婚年龄时才能拥有。它在草丛里徘徊着，这是真正的打猎。当
小狼蛛看到一种它想吃的猎物，就冲过去蛮横地把它赶出巢，然后紧迫不舍，那亡命者
正预备起飞逃走，可是往往来不及了——小狼蛛已经扑上去把它逮住了。
我喜欢欣赏我那实验室里的小狼蛛捕捉苍蝇时那种敏捷的动作。苍蝇虽然常常歇在
两寸高的草上，可是只要狼蛛猛然一跃，就能把它捉住。猫捉老鼠都没有那么敏捷。
但是这只是狼蛛小时候的故事，因为它们身体比较轻巧，行动不受任何限制，可以
随心所欲。以后它们要带着卵跑，不能任意地东跳西窜了。所以它就先替自己挖个洞，
整天在洞口守候着，这便是成年蜘蛛的猎食方式。

狼蛛的卵袋
假如你听到这可怕的狼蛛怎样爱护自己的家庭的故事，你一定会在惊异之余改变对
它的看法。
在八月的一个清晨，我发现一只狼蛛在地上织一个丝网，大小和一个手掌差不多。
这个网很粗糙，样子也不美观，但是很坚固。这就是它将要工作的场所，这网能使它的
巢和沙地隔绝。在这网上，它用最好的白丝织成一片大约有一个硬币大小的席子，它把
席子的边缘加厚，直到这席子变成一个碗的形状，周围圈着一条又宽又平的边，它在这
网里产了卵，再用丝把它们盖好，这样我们从外面看，只看到一个圆球放在一条丝毯上。
然后它就用腿把那些攀在圆席上的丝一根根抽去，然后把圆席卷上来，盖在球上，
然后它再用牙齿拉，用扫帚般的腿扫，直到它把藏卵的袋从丝网上拉下来为止，这可是
一项费神费力的工作。
这袋子是个白色的丝球，摸上去又软又粘，大小像一颗樱桃。如果你仔细观察，那
么你会发现在袋的中央有一圈水平的折痕，那里面可以插一根针而不致于把袋子刺破。
这条折纹就是那圆席的边。圆席包住了袋子的下半部，上半部是小狼蛛出来的地方。除
了母蜘蛛在产好卵后铺的丝以外，再也没有别的遮蔽物了。袋子里除了卵以外，也没有
别的东西，不像条纹蜘蛛那样，里面衬着柔软的垫褥和绒毛。狼蛛不必担心气候对卵的
影响，因为在冬天来临之前，狼蛛的卵早已孵化了。
母蛛整个早晨都在忙着编织袋子。现在它累了。它紧紧抱着它那宝贝小球，静静地
休息着，生伯一不留神就把宝贝丢了。第二天早晨，我再看到它的时候，它已经把这小
球挂到它身后的丝囊上了。
差不多有三个多星期，它总是拖着那沉重的袋子。不管是爬到洞口的矮墙上的时候，
还是在遭到了危险急急退入地洞的时候，或者是在地面上散步的时候，它从来不肯放下
它的宝贝的小袋。如果有什么意外的事情使这个小袋子脱离它的怀抱，它会立刻疯狂地
扑上去，紧紧地抱住它，并准备好反击那抢它宝贝的敌人。接着它很快地把小球挂到丝
囊上，很不安地带着它匆匆离开这个是非之地。
在夏天将要结束的那几天里，每天早晨，太阳已经把土地烤得很热的时候，狼蛛就
要带着它的小球从洞底爬出洞口，静静地趴着。初夏的时候，它们也常常在太阳高挂的
时候爬到洞口，沐浴着阳光小睡。不过现在，它们这么做完全是为了另外一个目的。以
前狼蛛爬到洞口的阳光里是为了自己，它躺在矮墙上，前半身伸出洞外，后半身藏在洞
里。它让太阳光照到眼睛上，而身体仍在黑暗中；现在它带着小球，晒太阳的姿势刚好
相反：前半身在洞里，后半身在洞外。它用后腿把装着卵的白球举到洞口，轻轻地转动
着它，让每一部分都能受到阳光的沐浴。这样足足晒了半天，直到太阳落山。它的耐心
实在令人感动，而且它不是一天两天这样做，而是在三四个星期内天天这样做。鸟类把
胸伏在卵上，它的胸能像火炉一样供给卵充分的热量；狼蛛把它的卵放在太阳底下，直
接利用这个天然的大火炉。

狼蛛的幼儿
在九月初的时候，小狼蛛要准备出巢了。这时小球会沿着折痕裂开。它是怎么裂开
的呢？会不会是母蛛觉察到里面有动静，所以在一个适当的时候把它打开了？这也是可
能的。但从另一方面看，也可能是那小球到了一定时间自己裂开的，就像条纹蜘蛛的袋
子一样。条纹蜘蛛出巢的时候，它们的母亲早已过世多时了。所以只有靠巢的自动裂开，
孩子们才能出来。
这些小狼蛛出来以后，就爬到母亲的背上，紧紧地挤着，大约有二百只之多，像一
块树皮似的包在母蛛身上。至于那袋，在孵化工作完毕的时候就从丝囊上脱落下来，被
抛在一边当垃圾了。
这些小狼蛛都很乖，它们不乱动，也不会为了自己挤上去而把别人推开。它们只是
静静地歇着。它们在干什么呢？它们是让母亲背着它们到处去逛。而它们的母亲，不管
是在洞底沉思，还是爬出洞外去晒太阳，总是背着一大堆孩子一起跑，它从不会把这件
沉重的外衣甩掉，直到好季节的降临。
这些小狼蛛在母亲背上吃些什么呢？照我看来，它们什么也没吃。我看不出它们长
大，它们最后离开母亲的时候，和它们刚从卵里出来的时候大小完全一样。
在坏的季节里，狼蛛母亲自己也吃得很少。如果我捉一只蝗虫去喂它，常会过了很
久以后它才开口。为了保持元气，它有时候不得不出来觅食，当然，它还是背着它的孩
子。
如果在三月里，当我去观察那些被风雨或霜雪侵蚀过的狼蛛的洞穴的时候，总可以
发现母蛛在洞里，仍是充满活力的样子，背上还是背满了小狼蛛。也就是说，母蛛背着
小蛛们活动，至少要经过五六个月。著名的美洲背负专家——鼹鼠，它也不过把孩子们
背上几个星期就把它们送走了，和狼蛛比起来，真是小巫见大巫了。
背着小狼蛛出征是很危险的，这些小东西常常会被路上的草拨到地上。如果有一支
小狼蛛跌落到地上，它将会遭遇什么命运呢？它的母亲会不会想到它，帮它爬上来呢？
绝对不会。一只母蛛需要照顾几百只小蛛，每只小蛛只能分得极少的一点爱。所以不管
是一只、几只或是全部小狼蛛从它背上摔下来，它也决不为它们费心。它不会让孩子们
依靠别人的帮助解决难题，它只是静静地等着，等它们自己去解决困难，事实上这困难
并不是不能解决，而且往往解决得很迅速、很利落。
我用一只笔把我实验室中的一个母狼蛛背上的小蛛刮下，母亲一点儿也不显得惊慌，
也不准备帮助它的孩子，继续若无其事地往前走。那些落地的小东西在沙地上爬了一会
儿，不久就都攀住了它们的母亲身体的一部分：有的在这里攀住了一只脚，有的在那里
攀住一只脚。好在它们的母亲有不少脚，而且撑得很开，在地面上摆出一个圆，小蛛们
就沿着这些柱子往上爬，不一会儿，这群小蛛又像原来那样聚在母亲背上了。没有一只
会漏掉。在这样的情况下，小狼蛛很会自己照顾自己，母亲从不需为它们的跌下而费心。
在母蛛背着小蛛的七个月里，它究竟喂不喂它们吃东西呢？当它猎取了食物后，是
不是邀孩子们共同享受呢？起初我以为一定是这样的，所以我特别留心母蛛吃东西时的
情形，想看看它怎样把食物分给那么多的孩子们。通常母蛛总是在洞里吃东西，不过有
时候偶然也到门口就着新鲜空气用餐。只有在这时候我才有机会，看到这样的情形：当
母亲吃东西的时候，小蛛们并不下来吃，连一点要爬下来分享美餐的意思都没有。好像
丝毫不觉得食物诱人一样，它们的母亲也不客气，没给它们留下任何食物。母亲在那儿
吃着，孩子们在那儿张望着——不，确切的说，它们仍然伏在妈妈的背上，似乎根本不
知道“吃东西”是怎样一种概念。在它们的母亲狼吞虎咽的时候，它们安安静静地呆在
那儿，一点儿也不觉得馋。
那么，在爬在母亲背上整整七个月的时间里，它们靠什么吸取能量、维持生命呢？
你或许会猜想它们不会是从母亲的皮肤上吸取养料的吧？我发现并不是这样的。因为据
我观察，它们从来没有把嘴巴贴在母亲的身上吮吸。而那母蛛，也并不见得瘦削和衰老，
它还是和以往一样神采奕奕，而且比以前更胖了。
那么又要问了，它们这些小蛛靠什么维持生命呢？一定不是以前在卵里吸收的养料。
以前那些微不足道的养料。别说是不能帮它们造出丝来，连维持它们的小生命都很困难。
在小蛛的身体里一定有着另外一种能量。
如果它们不动，我们很容易理解为什么它们不需要食物，因为完全的静止就相当于
没有生命。但是这些小蛛，虽然它们常常安静地歇在母亲背上，但它们时刻都在准备运
动。当它从母亲这个“婴儿车”上跌落下来的时候，它们得立刻爬起来抓住母亲的一条
腿，爬回原处；即使停在原地，它也得保持平衡；它还必须伸直小肢去搭在别的小蛛身
上，才能稳稳地趴在母亲背上。所以，实际上绝对的静止是不可能的。
从生理学角度看，我们知道每一块肌肉的运动都需要消耗能量。动物和机器一样，
用得久了会造成磨损，因此需常常修理更新。运动所消耗的能量，必须从别的地方得到
补偿。我们可以把动物的身体和火车头相比。当火车头不停地工作的时候，它的活塞、
杠杆、车轮以及蒸汽导管都在不断地磨损，铁匠和机械师随时都在修理和添加些新材料，
就好像供给它食物，让它产生新的力量一样。但是即使机器各部分都很完美，火车头还
是不能开动。一直要等到火炉里有了煤，燃起了火，然后才能开动。这煤就是产生能量
的“食物”，就是它让机器动起来的。
动物也是这样。有能量才能运动。小动物在胚胎时期，从母亲的胎盘里或者卵里吸
取养料，那是一种制造纤维素的养料，它使小动物的身体长大长坚固，并且补偿一些不
足的地方。但是，除此之外，必须有产生热量的食物，才能使小动物跑、跳、游泳、飞
跃，或是作其它各种运动。任何运动都少不了能量。
再讲这些小狼蛛，它们在离开母亲的背之前，并不曾长大。七个月的小蛛和刚刚出
生的小蛛完全一样大。卵供给了足够的养料，为它们的体质打下了一个良好的基础。但
它们后来不再长大，因此也不再需要吸收制造纤维的养料。这一点我们是能够理解的。
但它们是在运动的呀！并且运动得很敏捷。它们从哪里取得产生能量的食物呢？
我们可以这样想：煤——那供给火车头动能的食物究竟是什么呢？那是许多许多年
代以前的树埋在地下，它们的叶子吸收了阳光。所以煤其实就是贮存起来的阳光，火车
头吸收了煤燃烧提供的能量，也就是相当于吸收了太阳光的能量。
血肉之躯的动物也是这样，不管它是吃什么别的动物或植物以维持生命，大家最终
都是靠着太阳的能量生存的。那种热能量贮藏在草里、果子里、种子里和一切可作为食
物的东西里。太阳是宇宙的灵魂，是能量的最高赐予者，没有太阳，就没有地球上的生
命。
那么除了吃进食物，然后经过胃的消化作用变成能量以外，太阳光能不能直接射入
动物的身体，产生活力，就像蓄电池充电那样？为什么不能直接靠阳光生存呢？我们吃
的果子中除了阳光外，还有别的什么物质吗？
化学家告诉我们，将来我们可以靠一种人工的食物来维持生命。那时候所有的田庄
将被工厂和实验室所取代，化学家们的工作就是配置产生纤维的食物和产生能的食物，
物理学家们也靠着一些精巧仪器的帮助，每天把太阳能注进我们的身体，供给我们运动
所需的能量。那样我们就能不吃东西而维持生命。不吃饭而是吃太阳的光线，你能想象
吗？那将是一个多么美妙而有趣的世界！
这是我们的梦想，它能实现吗？这个问题倒是很值得科学家们研究的。

小蛛的飞逸
到三月底的时候，母蛛就常常蹲在洞口的矮墙上。这是小蛛们与母亲告别的时候了。
作母亲的仿佛早已料到这么一天，完全任凭它们自由地离去。对于小蛛们以后的命运，
它再也不需要负责了。
在一个天气很好的日子里，它们决定在那天最热的一段时间里分离。小蛛们三五成
群地爬下母亲的身体。看上去丝毫没有依依惜别之情，它们在地上爬了一会儿后，便用
惊人的速度爬到我的实验室里的架子上。它们的母亲喜欢住在地下，它们却喜欢往高处
爬。架子上恰好有一个竖起来的环，它们就顺着环很快地爬了上去。就在这上面，它们
快活地纺着丝，搓着疏松的绳子。它们的腿不住地往空中伸展，我知道这是什么意思：
它们还想往上爬，孩子长大了，一心想走四方闯天下，离家越远越好。
于是我又在环上插了一根树枝。它们立刻又爬了上去，一直爬到树枝的梢上。在那
里，它们又放出丝来，攀在周围的东西上，搭成吊桥。它们就在吊桥上来来去去，忙碌
地奔波着，看它们的样子似乎还不满足，还想一个劲儿往上爬。
我又在架子上插了一根几尺高的芦梗，顶端还伸展着细枝。那些小蛛立刻又迫不及
待地爬了上去，一直到达细枝的梢上。在那儿，它们又乐此不疲地放出丝、搭成吊桥。
不过这次的丝很长很细，几乎是飘浮在空中的，轻轻吹口气就能把它吹得剧烈地抖动起
来，所以那些小蛛在微风中好像在空中跳舞一般。这种丝我们平时很难看见，除非刚好
有阳光照在丝上，才能隐隐约约看到它。
忽然一阵微风把丝吹断了。断了的一头在空中飘扬着。再看这些小蛛，它们吊在丝
上荡来荡去，等着风停；如果风大的话，可能把它们吹到很远的地方，使它们重新登陆，
到一个陌生的地方。
这种情形又要维持好多天。如果在阴天，它们会保持静止，动都不想动，因为没有
阳光供给能量，它们不能随心所欲地活动。
最后，这个庞大的大家庭消失了。这些小蛛纷纷被飘浮的丝带到各个地方。原来背
着一群孩子的荣耀的母蛛变成了孤老。一下子失去那么多孩子，它看来似乎并不悲痛。
它更加精神焕发地到处觅食，因为这时候它背上再也没有厚厚的负担了，轻松了不少，
反而显得年轻了。不久以后它就要做祖母，以后还要做曾祖母，因为一只狼蛛可以活上
好几年呢。
从这一家狼蛛中，我们可以看到，有一种本能，很快地赋予小蛛，不久又很快地而
且是永远地消失。那就是攀高的本能。它们的母亲不知道自己的孩子曾有这样的本事，
孩子们自己不久以后也会彻底地忘记。它们到了陆地上，做了许多天流浪儿之后，便要
开始挖洞了。这时候，它们中间谁也不会梦想爬上一颗草梗的顶端。可那刚刚离开母蛛
的小蛛的确是那样迅速、那样容易地爬到高处，在它生命的转折之处，它曾是一个满怀
激情的攀登大师。我们现在知道了它这样做的目的：在很高的地方，它可以攀一根长长
的丝。那根长丝在空中飘荡着，风一吹，就能使它们飘荡到远方去。我们人类有飞机，
它们也有它们的飞行工具。在需要的时候，它替自己制造这种工具，等到旅行结束，它
也就把它忘记了。
------------------








克鲁蜀蜘蛛



克鲁蜀蜘蛛是一个极为聪明、灵巧的纺织家，而且就一只蜘蛛而言，克鲁蜀蜘蛛算
是很漂亮的了。它这名字是取自古希腊三位命运女神中的一位，也是最年幼的一位，她
是掌管纺线杆的，从她那里纺出了万物各自不同的命运。克鲁蜀蜘蛛能为自己纺出最精
美的丝，克鲁蜀女神却不能为我们纺出幸福的命运和舒适的生活，这实在是一件令世人
遗憾的事！
如果我们想认识克鲁蜀蜘蛛，我们必须到橄榄地的岩石的斜坡上。在被太阳灼得又
热又亮的地方，让我们把一些不大不小又扁平的石块翻起来——最好是翻开那些小石堆，
那是牧童堆起来做凳子用的，这种凳子尽管简陋，但深受牧童们的喜爱，因为他们可以
坐在上面看守山底下的羊群，边休息边工作，不亦乐乎。而这些地方往往是克鲁蜀蜘蛛
最爱呆的地方，如果你开始找不着它们，不要灰心，克鲁蜀蜘蛛在这个世界上的确是很
稀少的，而且并不是每个地方都适合它们的。所以找不着它们一点儿也不奇怪，如果我
们运气不坏的话，在我们翻起的石块下面，就会发现一个样子很特别的东西：形状好像
一个翻转的穹形屋顶，大概有半个梅子那么大，外面挂着一些小贝壳，一些泥土和干了
的虫子。
穹形顶的边缘有十二个尖尖的扇蛤，向各方伸展着，固定在石头上。显然，这就是
克鲁蜀蜘蛛的宫殿！
而入口处在什么地方呢？尽管周围有许多拱，但那些拱都开在屋顶的上部，没有一
个可以通到屋子里去。可这屋子的主人总得出来，出来后当然还得再回去呀，它究竟是
从哪里进出的呢？一根稻草会告诉我们一切秘密。
如果我们用一根稻草往拱形的开口处插进去，我们可以发现这些拱门里面都是反锁
着的，关得严严实实的。但是如果你把稻草很小心地用力插进去，你就会发现，其中必
定有一个拱门，它的边缘会裂成嘴唇般的两片。这便是门了，它有弹性，自己会关闭。
当克鲁蜀蜘蛛遇到危险的时候，它会飞快地跑回自己的家，用脚爪一触门，门就开
了，等它进去后，门又自动关闭了。如果需要的话，它还可以将门反锁，尽管那所谓的
“锁”也只不过是用几根丝线做的，起不了多大作用。但从外边看它跟别的拱完全一样，
可以起到迷惑的作用，敌人只能看见它很快地消失了，却不知道它究竟是从哪里逃遁的。
现在让我们打开它的门到里面看一看。啊！那么的富丽堂皇！我马上想起一个关于
娇贵的公主的神话故事，说她是如此的娇贵，只要她的垫褥底上有一片折皱的玫瑰花的
叶子，她就会睡得很不舒服。看到克鲁蜀蜘蛛的家，你会觉得它比那位公主还要难侍候。
它的床比天鹅绒还软，比夏天的云还白。床上有绒毯，有被子，都非常软。克鲁蜀蜘蛛
就安居在这绒毯和被子之间。它长着一双短短的腿，穿着黑色的衣服，背上还有五个黄
色的徽章。
要在这屋子里过上舒适的生活，必须有一个条件：那就是屋子必须建筑得很坚固，
尤其是在遭到大风大雨的时候。只要我们仔细观察的话，就可以发现克鲁蜀蜘蛛是如何
做到这一点的：支持着整个屋子的许多拱门都是固定在石头上的。我们可以看到在接触
点有一缕缕的长线，沿着石面伸展开去。我用尺子量了量，每一根足足有九尺长。原来
这屋子由许多“链条”攀着，就像阿拉伯人的帐幕用许多绳子攀着一样。所以显得格外
牢固。
还有一件小事挺值得我们注意，虽然它的屋子里那么整洁、华丽，可是屋子外面却
堆满了垃圾：有泥土、有腐烂了的木屑，还有一粒粒脏兮兮的砂石，有时候还有更脏的
东西，比如风干了的甲虫尸体、干足虫破碎的尸体，还有蜗牛的壳等等，一片狼藉，一
踏糊涂，而且都已被太阳晒得发白了。
克鲁蜀蜘蛛没有设陷阱的本领，所以它完全靠吃那种在石堆里跳来跳去的虫子维持
生计。哪一只冒失的虫子跳过它的居所，就会被它逮个正着，捉来饱餐一顿。至于那些
风干的尸体，它并不把它们丢开，而是挂在墙壁的四周，这似乎在炫耀自己的捕猎经历。
它到底为什么要把尸体挂起来呢？那些蜗牛壳大部分是空的，有时候里面也有活着
而且没有受伤的蜗牛，可是这些蜗牛躲在壳的深处，蜘蛛根本不能接触到它，把它挂起
来又有什么用呢？它不能把那么硬的壳打破，也无法从开口处把蜗牛揪出来，它居然还
收集了那么多蜗牛，它到底想干什么呢？
后来我们想到了简单的平衡问题。普通的家蛛在墙角张了网，为了要使网保持一定
的形状而不受风吹雨打的影响，就常常把石灰或泥嵌入网中。克鲁蜀蜘蛛在自己的屋里
挂上重物虽然也是为了使屋子稳固，只是不是这种原理，而是它懂得如果把重物挂在屋
子里很低的地方，屋子四周就会平衡：在低的地方增加重量就可以把重心降低，使平衡
更加稳定。这些重物主要是昆虫的尸体或空壳，因为这类东西举手可得，不需到远处去
找，全都有现成的。
现在，你们或许要问：它在那软软的屋子里做什么呢？醉生梦死还是辛勤劳动呢？
据我所知，它什么也不做。它的肚子已经装得饱饱的了。它舒适地摊开了脚，躺在软软
的绒毯上，什么都不用做，什么都不用想，只是静静地聆听着地球旋转的声音。这不是
睡，也不是醒，而是一种半梦半醒的状态，这时候它除了快乐以外，什么感觉都没有。
让我们想象一下：如果你辛苦了一天后很舒适地躺在床上，彻底地放松，到快要入睡的
时候，将会无忧无虑地享受到这生活中最美满的一刻。克鲁蜀蜘蛛似乎也知道这美妙的
一刻，并且比我们更会享受这一刻的到来。
------------------








迷宫蛛



会结网的蜘蛛称得上是个纺织能手，它们用蛛网来猎取自投罗网的小虫子们，可谓
“坐享其成，得来全不费功夫”。还有许多其它种类的蜘蛛，它们用许多别的聪明的方
法猎取食物，同样可以以逸待劳，大获丰收。其中有几种在这方面很有造诣，几乎所有
的有关昆虫的书都会把它们列举出来。
那是一种黑色的蜘蛛，也有人叫它美洲狼蛛，它们是住在洞里的，就像我以前讲到
的欧洲狼蛛一样。但是它们的洞穴比欧洲狼蛛的洞穴要完备精细得多。欧洲狼蛛的洞口
只有一圈矮墙，用小石子、丝和废料堆成的，而美洲狼蛛的洞口上有一扇活动门，是由
一块圆板、一个槽和一个栓子做成的。当一只狼蛛回家的时候，门便会落进槽里，自动
把门关了。如果有谁在门外想把它掀起来的话，狼蛛只要用两只爪把柱子抵住，门就紧
紧关闭住，不会受外面的影响。
另外一种是水蛛。它能替自己做一只性能很好的潜水袋，里面贮藏着空气。它在这
里面等待猎物经过，同时也可以说是在避暑。在太阳像大火炉一样的日子里，这地方的
确是一个舒适凉爽的避暑胜地。人类中也有人尝试用最硬的石块或大理石在水下造房子。
不知大家有没有听说过泰比利斯，他是罗马的一个暴君，他生前曾经叫人为他造了一座
水下宫殿，供自己寻欢作乐。不过到现在这个宫殿只给人们留下一点回忆和感慨，而狼
蛛的水晶宫，却是永远灿烂辉煌的。
如果我有机会观察一下这些水蛛的话，我一定能在它们的生命史上添上一些未经记
载的事实。但是现在我不得不放弃这个想法。因为我们这一带没有水蛛。至于那美洲狼
蛛，我也只有在路旁看到过一次。而那时候我恰巧有别的事情要去办理，没有时间去看
它。错失这个良机后，后来也就一直没见到过它。
但是，并不是稀罕的虫子才值得研究。普通的虫子，如果好好地研究起来，也能发
现许多有趣的事情。我对迷宫蛛的接触机会极多，对它也很感兴趣，所以对它作了一番
研究，我觉得是很有收获的。在七月的清晨，太阳还没有焦灼着人的头颈的时候，每星
期我总要去树林里看几次迷宫蛛。孩子们也都跟着我去，每人还带上一个橘子，以供解
渴之用。
走进树林，不久，我们就发现许多很高的丝质建筑物，丝线上还串着不少露珠，在
太阳光的照射下闪闪发光，好像皇宫里的稀世珍宝一般。孩子们被这个美丽的“灯架”
惊呆了，几乎忘记了他们的橘子。我们的蜘蛛的迷宫真算得上一个奇观！
经过太阳半小时的照射，魔幻般的珍珠随着露水一起消失了。现在可以来专心观察
它的网了。在那丛蔷蔽花的上方张着一张网，大概有一块手帕那么大，周围有许多线把
它攀到附近的矮树丛中，使它能够在空中固定住，中间这张网看起来犹如一层又轻又软
的纱。
网的四周是平的，渐渐向中央凹，到了最中间便变成一根管子，大约有八九寸深，
一直通到叶丛中。
蜘蛛就坐在管子的进口处。它对着我们坐着，一点儿也不惊慌。它的身体是灰色的，
胸部有两条很阔的黑带，腹部有两条细带，由白条和褐色的斑点相间排列而成。在它的
尾部，有一种“双尾”，这在普通蜘蛛中是很少见的。
我猜想在管子的底部，一定有一个垫得软软的小房间，作为迷宫蛛空闲时候的休息
室。可事实上那里并没有什么小房间，只有一个像门一样的东西，一直是开着的，它在
外面遇到危险的时候，可以直接逃回来。
上面那个网由于用许多丝线攀到附近的树枝上，所以看上去活像一艘暴风雨下抛锚
的船。这些充当铁索的丝线中，有长的，也有短的；有垂直的，也有倾斜的；有紧张的，
也有松弛的；有笔直的，也有弯曲的，都杂乱地交叉在三尺以上的高处。这确实可以算
是一个迷宫，除了最强大的虫子外，谁都无法打破它，逃脱它的束缚。
迷宫蛛不像别的蜘蛛那样可以用粘性的网作为陷阱，它的丝是没有粘性的，它的网
妙就妙在它的迷乱。你看那只小蝗虫，它刚刚在网上落脚，便由于网摇曳不定，根本设
法让自己站稳。一下子陷了下去，它开始焦躁地挣扎，可是越挣扎陷得越深，好像掉进
了可怕的深渊一样。蜘蛛呆在管底静静地张望着，看着那倒霉的小蝗虫垂死挣扎，它知
道，这个猎物马上会落到网的中央，成为它的盘中美餐。
果然，一切都在蜘蛛的意料之中。它不慌不忙地扑到猎物上，慢慢地一口一口地吮
吸着它的血，一副得意洋洋的样子。至于那蝗虫，在蜘蛛咬它第一口的时候就死了——
蜘蛛的毒液使它一命呜呼。接下来蜘蛛就要从容地来吃完它了，而对于这只蝗虫来说，
这远比半死不活或者活活被蜘蛛撕成碎片要舒服多了。
到快要产卵的时候，迷宫蛛就要搬家了。尽管它的网还是完好无损，但它必须忍痛
割爱。它不得不舍弃它，而且以后也不再回来了。它必须去完成它的使命，一心一意去
筑巢了。它把巢做在什么地方呢？迷宫蛛自己当然知道得很清楚，而我，却一点头绪都
没有，实在猜想不出它会把巢造在哪儿。我花了整个早晨在树林中各个地方搜索。工夫
不负有心人，我终于发现了它的秘密。
在离网相当远的一个树丛里，它造好了它的巢。那里堆着一堆枯柴，草率而杂乱地
缠在一起，显得有点脏。就在这个简陋的盖子下，有一个做得比较细致、精巧的丝囊，
里面就是迷宫蛛的卵。
看到它的巢那么简陋，我不禁有些失望。但是后来我想到了，这一定是因为环境不
够好。你想，在这样一个密密的树丛里，一堆枯枝枯叶中，哪有条件让它做精致的活呢
7为了要证明我的推想没有错，我带了六只快要产卵的迷宫蛛到我家里，放在实验室的一
个铁笼子里面，然后把铁笼子竖在一个盛沙的泥盘子里，又在泥盘中央插了一根百里香
的小树枝，使每一个巢有攀附的地方。一切准备就绪后，现在就让它们大显身手吧。
这个实验获得了极大的成功。七月底的时候，我得到了六只雪白的、外观富丽精致
的丝囊。迷宫蛛在这样一个舒适的环境里工作，活干得自然细致了许多。让我来尽情地
观察吧！这个巢是一个由白纱编织而成的卵形的囊，有一个鸡蛋那么大。内部的构造也
很迷乱，和它的网差不多——看来这种建筑风格在它的脑子里已经根深蒂固了，所以无
论在什么场所，在什么条件下，它造的建筑物都是那样的迷乱无章。
这个布满丝的迷宫还是一个守卫室。在这乳白色的半透明的丝墙里面还装着一个卵
囊，它的形状有点像那些代表某一等级的骑士的星形勋章。这是一个很大的灰白色的丝
袋，周围筑着圆柱子，使它能够固定在巢的中央。这种圆柱都是中间细，两头粗，总共
大约有十个，在卵室的周围构成一个白色的围廊。母蛛在这个围廊里徘徊着，一会儿在
这儿停住，一会儿又在那儿停住，时时聆听着卵囊里的动静，活像一个马上要做父亲的
人在产室外面焦急地等待着孩子的第一声啼哭。这样一个卵巢里面，大概藏着一百颗左
右淡黄色的卵。
轻轻移去外面的白丝墙后，可以看到里面，还有一层泥墙，那是丝线夹杂着小碎石
做成的。可是这些小沙子怎么到丝墙里面去的呢？是跟着雨水渗进去的吗？不对，因为
外面的丝墙上白得没有一丝斑点，更不用说什么水迹了，看来决不是从这墙上渗进去的。
到后来，我才发现这是母蛛自己搬进去的，它为了怕卵受到寄生虫的侵犯，所以特地把
砂粒掺在丝线里面做成一培坚固的墙。
这丝墙里面还有一个丝囊，那才是盛卵的囊。我打开的这个巢里面的卵已经孵化了，
所以我能看到许多弱小的小蜘蛛在囊里快乐地爬来爬去。
但是，再回过头来看看那母蛛，它为什么要舍弃那张还完好无损的网，而把巢筑到
那么远的地方呢？它的舍近求远自然有它的道理，你还记得它的网的样子吧？在它的网
的上方，有一个错综复杂的迷宫，高高地露在树叶丛的外面，这是一个巨大的陷阱，同
时也是一个很醒目的标志，它的敌人——寄生虫轻而易举地就能看到这个迷宫，然后循
着它再找到迷宫蛛的巢——如果这巢离那醒目的网很近的话，那么寄生虫会不费吹灰之
力就把它的巢找到。寄生虫的卑鄙手段上面已经详细地报述过了，提防寄生虫入侵是每
一个母亲为了保护下一代所必须做的一项重要事业。况且这种迷宫蛛的死敌寄生虫专门
吃新生的卵，如果找到迷宫蛛的巢，会毫不客气地把它毁灭。所以聪明而尽责的迷宫蛛
就趁着夜色到各处去察看地形，找一个最安全的地方作为未来家族的安乐窝，至于那个
地方美观不美观，环境怎么样，倒是次要考虑的了。那沿着地面生长的矮矮的荆棘丛，
它们的叶子在冬天里也不会脱落，而且它们还能钩住附近的枯叶，对迷宫蛛来说，不能
说不是一个理想的居处，还有那又矮又细的迷迭香丛，也是迷宫蛛爱做巢的地方。在这
种地方，我常常能够找到不少迷宫蛛的巢。
有许多蜘蛛产卵以后就永远离开自己的巢了。可是迷宫蛛和蟹蛛一样，会一直紧紧
地守着巢。不过和蟹蛛不同的是，它不会像蟹蛛那样绝食，以致日益消瘦下去，它会照
常捕蝗虫吃。它用一团纷乱错杂的丝。筑起了一个捕虫箱，继续补充营养。
当它不捕食时，那就像我们所看到的那样，在走廊里踱来踱去，侧耳倾听四面八方
的动静。如果我用一根稻草在巢的某一处拨一下，它就会立即冲出来查看个究竟。就是
用这种警惕的办法，它尽心尽责地保护着自己的未成年的孩子们。
迷宫蛛产了卵后胃口还那么好，表示它还要继续工作。因为昆虫不像人类，有时候
吃东西仅仅是因为嘴馋。它们吃东西就是为了工作。
可是产完卵后，它这一生中最伟大的任务已经完成了，它还要作什么工作呢？经过
我细细地探究后，我才发现它所要做的工作是什么。大约又花了一个月左右的工夫，它
继续在巢的墙上添着丝。这墙最初是透明的，现在却变得又厚又不透明了。这就是它之
所以还要大吃特吃的原因：为了充实它的丝腺来为它的巢造一垛厚墙。
大约在九月中旬，小蜘蛛们从巢里出来了。但是它们并不离开巢，它们要在这温软
舒适的巢里过冬。母蜘蛛继续看护着它们，继续纺着丝线。不过岁月无情，它一天比一
天迟钝了。它的食量也渐渐地小起来。有时候我特意放几条蝗虫到它的陷阱里去喂它，
它也显得无动于衷，一口也不想吃。虽然这样，它还能维持四五个星期的生命，在它离
开这个世界之前，它继续一步不离地守着这巢，每次一听到巢里新生的小蛛在活泼地爬
来爬去，它便感到无限的满足和快慰。最后，到十月底的时候，它用最后一点力气替孩
子们咬破巢后，便精疲力尽地死去。它已尽了一个最慈爱的母亲所应尽的责任，它无愧
于它的孩子，无愧于这个世界。至于以后的事；它便托付给造物主了。到了来年的春天，
小蜘蛛们从它们舒适的屋里走出来，然后像蟹蛛那样，靠着它们的飞行工具——游丝，
飘散到各地去了。它们的母亲在天之灵，看着自己的孩子一个个各自成家立业了，也该
欣慰了。
------------------








蛛网的建筑



即使在最小的花园里，也能看到园蛛的踪迹。它们都算得上是天才的纺织家。
如果我们在黄昏的时候散步，我们可以从一丛迷迭香里寻找蛛丝马迹。我们所观察
的蜘蛛往往爬行得很慢，所以我们应该索性坐在矮树丛里看。那里的光线比较充足。让
我们再来给自己加一个头衔，叫做“蛛网观察家”吧！世界上很少有人从事这种职业，
而且我们也不用指望从这行业上嫌点钱。但是，不要计较这些，我们将得到许多有趣的
知识。从某种意义上讲，这比从事任何一个职业要有意思得多。
我所观察的都是些小蜘蛛。它们比成年的蜘蛛要小得多。而且它们都是在白天工作，
甚至是在太阳底下工作的，尽管它们的母亲只有在黑夜里才开始纺织。当到每年一定的
月份的时候，蜘蛛们便在太阳下山前两小时左右开始它们的工作了。
这些小蛛都离开了它们白天的居所，各自选定地盘，开始纺线。有的在这边，有的
在那边，谁也不打扰谁。我们可以任意地拣一只小蛛来观察。
让我们就在这只小蛛面前停下吧。它正在打基础呢。它在迷迭香的花上爬来爬去，
从一根枝端爬到另一根枝端忙忙碌碌的，它所攀到的枝大约都是十八寸距离之内的。太
远的它就无能为力了。渐渐地它开始用自己梳子似的后腿把丝从身体上拉出来，放在某
个地方作为基底，然后漫无规则地一会儿爬上，一会儿爬下，这样奔忙了一阵子后，结
果就构成了一个丝架子。这种不规则的结构正是它所需要的。这是一个垂直的扁平的
“地基”。正是因为它是错综交叉的，因此这个“地基”很牢固。
后来它在架子的表面横过一根特殊的丝，别小看这根细丝，那是一个坚固的网的基
础。这根线的中央有一个白点，那是一个丝垫子。
现在是它做捕虫网的时候了。它先从中心的白点沿着横线爬，很快就爬到架子的边
缘，然后以同样快的速度回到中心，再从中心出发以同样的方式爬到架子边缘，就这样
一会儿上，一会儿下，一会儿左，一会儿右。每爬一次便拉成一个半径，或者说，做成
一根辐。不一会儿，便这儿那儿地做成了许多辐，不过次序很乱。
无论谁，如果看到它已完成的网是那么地整洁而有规则，一定会以为它做辐的时候
也是按着次序一根根地织过去，然而恰恰相反，它从不按照次序做，但是它知道怎样使
成果更完美。在同一个方向安置了几根辐后，它就很快地往另一个方向再补上几条，从
不偏爱某个方向，它这样突然地变换方向是有道理的：如果它先把某一边的辐都安置好，
那么这些辐的重量，会使网的中心向这边偏移从而使网扭曲，变成很不规则的形状。所
以它在一边安放了几根辐后，立刻又要到另一边去，为的是时刻保持网的平衡。
你们一定不会相信，像这样毫无次序又是时时间断的工作会造出一个整齐的网。可
是事实确实如此，造好的辐与辐之间的距离都相等，而且形成一个很完整的圆。不同的
蜘蛛网的辐的数目也不同，角蛛的网有二十一根辐，条纹蜘蛛有三十二根，而丝光蛛有
四十二根。这种数目并不是绝对不变的，但是基本上是不变的，因此你可以根据蛛网上
的辐条数目来判定这是哪种蜘蛛的网。
想想看，我们中间谁能做到这一点：不用仪器，不经过练习，而能随手把一个圆等
分？但是蜘蛛可以，尽管它身上背着一个很重的袋子，脚踩在软软的丝垫上，那些垫还
随风飘荡，摇曳不定，它居然能够不加思索地将一个圆极为精细地等分。它的工作看上
去杂乱无序，完全不合乎几何学的原理，但它能从不规则的工作中得出有规则的成果来。
我们都对这个事实感到惊异。它怎么能用那么特别的方法完成这么困难的工作呢？这一
点我至今还在怀疑。
安排辐的工作完毕后，蜘蛛就回到中央的丝垫上。然后从这一点出发，踏着辐绕螺
旋形的圈子。它现在正在做一种极精致的工作。它用极细的线在辐上排下密密的线圈。
这是网的中心，让我们把它叫作“休息室”吧。越往外它就用越粗的线绕。圈与圈之间
的距离也比以前大。绕了一会，它离中心已经很远了，每经过一次辐，它就把丝绕在辐
上粘住。最后，它在“地基”的下边结束了它的工作。圈与圈之间的平均距离大约有三
分之一寸左右。
这些螺旋形的线圈并不是曲线。在蜘蛛的工作中没有曲线，只有直线和折线。这线
圈其实是辐与辐之间的横档所连成的。
以前所做的只能算作是一个支架，现在它将要在这上面做更为精致的工作。这一次
它从边缘向中心绕。而且圈与圈之间排得很紧，所以圈数也很多。
这种工作的详细情形很不容易看清，因为它的动作极为快捷而且振动得很厉害，包
括一连串的跳跃、摇摆和弯曲，使人看得眼花缭乱。如果分解它们的动作，可以看到它
的其中两条腿不停地动着，一条腿把丝拖出来传给另外一条腿，另一条腿就把这丝安在
辐上。由于丝本身有粘性，所以很容易在横档和丝接触的地方把新技出来的丝粘上去。
蜘蛛不停地绕着圈，一面绕一面把丝粘在辐上。它到达了那个被我们称作“休息室”
的边缘了。于是它立刻结束了它的绕线运动。以后它就会把中央的丝垫子吃掉。它这么
做是为了节约材料，它下一次织网的时候就可以把吃下的丝再纺出来用了。
有两种蜘蛛，也就是条纹蛛和丝光蛛，做好了网后，还会在网的下部边缘的中心织
一条很阔的锯齿形的丝带作为标记。有时候，它们还在这一条丝带的封面，就是网的上
部边缘到中心之间再织一条较短的丝带，以表明这是它们的作品，著作权不容侵犯。

粘性的网
蛛网中用来作螺旋圈的丝是一种极为精致的东西，它和那种用来做辐和“地基”的
丝不同。它在阳光中闪闪发光，看上去像一条编成的丝带。我取了一些丝回家，放在显
微镜下看，竞发现了惊人的奇迹：
那根细线本来就细得几乎连肉眼都看不出来，但它居然还是由几根更细的线缠合而
成的，好像大将军剑柄上的链条一般。更使人惊异的是，这种线还是空心的，空的地方
藏着极为浓厚的粘液，就和粘稠的胶液一样，我甚至可以看到它从线的一端滴出来。这
种粘液能从线壁渗出来，使线的表面有粘性。我用一个小试验去测试它到底有多大粘性：
我用一片小草去碰它，立刻就被粘住了。现在我们可以知道，园蛛捕捉猎物靠的并不是
围追堵截。而是完全靠它粘性的网，它几乎能粘住所有的猎物。可是又有一个问题出来
了：蜘蛛自己为什么不会被粘住呢？
我想其中一个原因是，它的大部分时间被用来坐在网中央的休息室里，而那里的丝
完全没有粘性。不过这个说法不能自圆其说，它无法一辈子坐在网中央不动，有时候，
猎物在网的边缘被粘住了。它必须很快地赶过去放出丝来缠住它，在经过自己那充满粘
性的网时，它怎么防止自己不被粘住呢？是不是它脚上有什么东西使它能在粘性的网上
轻易地滑过呢？它是不是涂了什么油在脚上？因为大家都知道，要使表面物体不粘，涂
油是最佳的办法。
为了证明我的怀疑，我从一只活的蜘蛛身上切下一条腿，在二硫化碳里浸了一个小
时，再用一个也在二硫化碳里浸过的刷子把这条腿小心地洗一下。二硫化碳是能溶解脂
肪的，所以如果腿上有油的话，这一洗就会完全洗掉了。现在我再把这条腿放到蛛网上，
它被牢牢地粘住了！由此我们知道，蜘蛛在自己身上，涂上了一层特别的“油”，这样
它能在网上自由地走动而不被粘住。但它又不愿老停在粘性的螺旋圈上，因为这种“油”
是有限的，会越用越少。所以它大部分时间呆在自己的“休息室”里。
从实验中我们得知这蛛网中的螺旋线是很容易吸收水分的。因为这个，当空气突然
变得潮湿的时候，它们就停止织网工作，只把架子、辐和“休息室”做好，因为这些都
不受水分的影响。至于那螺旋线的部分，它们是不会轻易做上去的，因为如果它吸收过
多的水分，以后就不能充分地吸水解潮了。有了这螺旋线，在极热的天气里，蛛网也不
会变得干燥易断，因为它能尽量地吸收空气中的水分以保持它的弹性并增加它的粘性。
哪一个捕鸟者在做网的时候，在艺术上和技术上能比得上蜘蛛呢？而蜘蛛织这么精致的
网只是为了捕一只小虫！真是有点大材小用了！
同时蜘蛛还是一个热忱积极的劳动者。我曾计算过，角蛛每做一个网需制造大约二
十码长的丝，至于那更精巧的丝，光蛛就得造出三十码，在这两个月中，我的角蛛邻居
几乎每夜都要修补它的网。这样，在这个时期中，它就得从它娇小瘦弱的身体上绵绵不
断抽出这种管状的，富有弹性的丝。
我们不禁要怀疑，它小小的身体怎么能产出那么多丝？它怎么能把这些丝搓成管状，
又怎么能在里面灌上粘液呢？它又怎能有时制出普通的丝，有时造出云朵状的丝花来垫
巢，最后还能制出黑色的丝带来装饰巢呢？这些问题一直在我的脑子里盘绕，并使我百
思而不得其解。
------------------








蜘蛛的几何学



当我们观察着园蛛，尤其是丝光蛛和条纹蛛的网时，我们会发现它的网并不是杂乱
无章的，那些辐排得很均匀，每对相邻的辐所交成的角都是相等的；虽然辐的数目对不
同的蜘蛛而言是各不相同的，可这个规律适用于各种蜘蛛。
我们已经知道，蜘蛛织网的方式很特别，它把网分成若干等份，同一类蜘蛛所分的
份数是相同的。当它安置辐的时候，我们只见它向各个方向乱跳，似乎毫无规则，但是
这种无规则的工作的结果是造成一个规则而美丽的网，像教堂中的玫瑰窗一般。即使他
用了圆规、尺子之类的工具。没有一个设计家能画出一个比这更规范的网来。
我们可以看到，在同一个扇形里，所有的弦，也就是那构成螺旋形线圈的横辐，都
是互相平行的，并且越靠近中心，这种弦之间的距离就越远。每一根弦和支持它的两根
辐交成四个角，一边的两个是钝角，另一边的两个是锐角。而同一扇形中的弦和辐所交
成的钝角和锐角正好各自相等——因为这些弦都是平行的。
不但如此，凭我们的观察，这些相等的锐角和钝角，又和别的扇形中的锐角和钝角
分别相等，所以，总的看来，这螺旋形的线圈包括一组组的横档以及一组组和辐交成相
等的角。
这种特性使我们想到数学家们所称的“对数螺线”。这种曲线在科学领域是很著名
的。对数螺线是一根无止尽的螺线，它永远向着极绕，越绕越靠近极，但又永远不能到
达极。即使用最精密的仪器，我们也看不到一根完全的对数螺线。这种图形只存在科学
家的假想中，可令人惊讶的是小小的蜘蛛也知道这线，它就是依照这种曲线的法则来绕
它网上的螺线的，而且做得很精确。
这螺旋线还有一个特点。如果你用一根有弹性的线绕成一个对数螺线的图形，再把
这根线放开来，然后拉紧放开的那部分，那么线的运动的一端就会划成一个和原来的对
数螺线完全相似的螺线，只是变换了一下位置。这个定理是一位名叫杰克斯·勃诺利的
数学教授发现的，他死后，后人把这条定理刻在他的墓碑上，算是他一生中最为光荣的
事迹之一。
那么，难道有着这些特性的对数螺线只是几何学家的一个梦想吗？这真的仅仅是一
个梦、一个谜吗？那么它究竟有什么用呢？
它确实广泛的巧合，总之它是普遍存在的，有许多动物的建筑都采取这一结构。有
一种蜗牛的壳就是依照对数螺线构造的。世界上第一只蜗牛知道了对数螺线，然后用它
来造壳，一直到现在，壳的样子还没变过。
在壳类的化石中，这种螺线的例子还有很多。现在，在南海，我们还可以找到一种
太古时代的生物的后代，那就是鹦鹉螺。它们还是很坚贞地守着祖传的老法则，它们的
壳和世界初始时它们的老祖宗的壳完全一样。也就是说，它们的壳仍然是依照对数螺线
设计的。并没有因时间的流逝而改变，就是在我们的死水池里，也有一种螺，它也有一
个螺线壳，普通的蜗牛壳也是属于这一构造。
可是这些动物是从哪里学到这种高深的数学知识的呢？又是怎样把这些知识应用于
实际的呢？有这样一种说法，说蜗牛是从蠕虫进化来的。某一天，蠕虫被太阳晒得舒服
极了，无意识地揪住自己的尾巴玩弄起来，便把它绞成螺旋形取乐。突然它发现这样很
舒服，于是常常这么做。久而久之便成了螺旋形的了，做螺旋形的壳的计划，就是从这
时候产生的。
但是蜘蛛呢？它从哪里得到这个概念呢？因为它和蠕虫没有什么关系。然而它却很
熟悉对数螺线，而且能够简单地运用到它的网中。蜗牛的壳要造好几年，所以它能做得
很精致，但蛛网差不多只用一个小时就造成了，所以它只能做出这种曲线的一个轮廊，
尽管不精确，但这确实是算得上一个螺旋曲线。是什么东西在指引着它呢？除了天生的
技巧外，什么都没有。天生的技巧能使动物控制自己的工作，正像植物的花瓣和小蕊的
排列法，它们天生就是这样的。没有人教它们怎么做，而事实上，它们也只能作这么一
种，蜘蛛自己不知不觉地在练习高等几何学，靠着它生来就有的本领很自然地工作着。
我们抛出一个石子，让它落到地上，这石子在空间的路线是一种特殊的曲线。树上
的枯叶被风吹下来落到地上，所经过的路程也是这种形状的曲线。科学家称这种曲线为
抛物线。
几何学家对这曲线作了进一步的研究，他们假想这曲线在一根无限长的直线上滚动，
那么它的焦点将要划出怎样一道轨迹呢？答案是：垂曲线。这要用一个很复杂的代数式
来表示。如果要用数字来表示的话，这个数字的值约等于这样一串数字1＋1/1＋1/1*2＋
1/1*2*3＋1/1*2*3*4＋……的和。
几何学家不喜欢用这么一长串数字来表示，所以就用“e”来代表这个数。e是一个
无限不循环小数，数学中常常用到它。
这种线是不是一种理论上的假想呢？并不，你到处可以看到垂曲线的图形：当一根
弹性线的两端固定，而中间松驰的时候，它就形成了一条垂曲线；当船的帆被风吹着的
时候，就会弯曲成垂曲线的图形；这些寻常的图形中都包含着“e”的秘密。一根无足轻
重的线，竟包含着这么多深奥的科学！我们暂且别惊讶。一根一端固定的线的摇摆，一
滴露水从草叶上落下来，一阵微风在水面拂起了微波，这些看上去稀松平常、极为平凡
的事，如果从数学的角度去研究的话，就变得非常复杂了。
我们人类的数学测量方法是聪明的。但我们对发明这些方法的人，不必过分地佩服。
因为和那些小动物的工作比起来，这些繁重的公式和理论显得又慢又复杂。难道将来我
们想不出一个更简单的形式，并使它运用到实际生活中吗？难道人类的智慧还不足以让
我们不依赖这种复杂的公式吗？我相信，越是高深的道理，其表现形式越应该简单而朴
实。
在这里，我们这个魔术般的“e”字又在蜘蛛网上被发现了。在一个有雾的早晨，这
粘性的线上排了许多小小的露珠。它的重量把蛛网的丝压得弯下来，于是构成了许多垂
曲线，像许多透明的宝石串成的链子。太阳一出来，这一串珠子就发出彩虹一般美丽的
光彩。好像一串金钢钻。“e”这个数目，就包蕴在这光明灿烂的链子里。望着这美丽的
链子，你会发现科学之美、自然之美和探究之美。
几何学，这研究空间的和谐的科学几乎统治着自然界的一切。在铁杉果的鳞片的排
列中以及蛛网的线条排列中，我们能找到它；在蜗牛的螺线中，我们能找到它；在行星
的轨道上，我们也能找到它，它无处不在，无时不在，在原子的世界里，在广大的宇宙
中，它的足迹遍布天下。
这种自然的几何学告诉我们，宇宙间有一位万能的几何学家，他已经用它神奇的工
具测量过宇宙间所有的东西。所以万事万物都有一定的规律。我觉得用这个假设来解释
鹦鹉螺和蛛网的对数螺线，似乎比蠕虫绞尾巴而造成螺线的说法更恰当。
------------------








蜘蛛的电报线



在六种园蛛中，通常歇在网中央的只有两种，那就是条纹蜘蛛和丝光蜘蛛。它们即
使受到烈日的焦灼，也决不会轻易稍离开网去阴凉处歇一会儿。至于其它蜘蛛，它们一
律不在白天出现。它们自有办法使工作和休息两个互不相误，在离开它们的网不远的地
方，有一个隐蔽的场所，是用叶片和线卷成的。白天它们就躲在这里面，静静地，让自
己深深地陷入沉思中。
这阳光明媚的白天虽然使蜘蛛们头晕目眩，却也是其它昆虫最活跃的时候：蝗虫们
更活泼地跳着，蜻蜓们更快活地飞舞着。所以正是蜘蛛们捕食的好时机，那富有粘性的
网虽然晚上是蜘蛛的居所，白天还是一个大陷阱，如果有一些又粗心又愚蠢的昆虫碰到
网上，被粘住了，躲在别处的蜘蛛是否会知道呢？不要为蜘蛛会错失良机而担心，只要
网上一有动静，它便会闪电般地冲过来。它是怎么知道网上发生的事的呢？让我来解释
吧。
使它知道网上有猎物的是网的振动，而不是它自己的眼睛。为了证明这一点，我把
一只死蝗虫轻轻地放到有好几只蜘蛛的网上，并且放在它们看得见的地方。有几只蜘蛛
是在网中，有几只是躲在隐蔽处，可是它们似乎都不知道网上有了猎物。后来我把蝗虫
放到了它们面前，它们还是一动不动。它们似乎瞎了，什么也看不见。于是我用一根长
草拨动那死蝗，让它动起来，同时使网振动起来。
结果证明：停在网中的条纹蛛和丝光蛛飞速赶到蝗虫身边；其它隐藏在树叶里的蜘
蛛也飞快地赶来，好像平时捉活虫一般，熟练地放出丝来把死蝗虫捆了又捆，缠了缠，
丝毫不怀疑自己是不是在浪费宝贵的丝线，由这个实验可见，蜘蛛什么时候出来攻击猎
物，完全要看网什么时候振动。
如果我们仔细观察那些白天隐居的蜘蛛们的网，我们可以看到从网中心有一根丝一
直通到它隐居的地方，这根线的长短大约有二十二寸；不过角蛛的网有些不同，因为它
们是隐居在高高的树上的，所以它的这根丝一般有八九尺长。
这条斜线还是一座桥梁，靠着它，蜘蛛才能匆匆地从隐居的地方赶到网中，等它在
网中央的工作完毕后，又沿着它回到隐居的地方，不过这并不就是这根线的全部效用。
如果它的作用仅仅在于这些的话，那么这根线应该从网的顶端引到蜘蛛的隐居处就可以
了。因为这可以减小坡度，缩短距离。
这根线之所以要从网的中心引出是因为中心是所有的辐的出发点和连接点，每一根
辐的振动，对中心都有直接的影响。一只虫子在网的任何一部分挣扎，都能把振动直接
传导到中央这根线上。所以蜘蛛躲在远远的隐蔽处，就可以从这根线上得到猎物落网的
消息。这根斜线不但是一座桥梁，并且是一种信号工具，是一根电报线。
年轻的蜘蛛都很活泼，它们都不懂得接电报线的技术。只有那些老蜘蛛们，当它们
坐在绿色的帐幕里默默地沉思或是安详地假寐的时候，它们会留心着电报线发出的信号，
从而得知在远处发生的动静。
长时间的守候是辛苦的，为了减轻工作的压力和好好休息。同时又丝毫不放松对网
上发生的情况的警觉，蜘蛛总是把腿搁在电报线上。这里有一个真实的故事可以证明这
一点。
我曾经打到一只在两棵相距一码的常青树间结了一张网的角蛛。太阳照得丝网闪闪
发光，它的主人早已在天亮之前藏到居所里去了。如果你沿着电报线找过去，就很容易
找到它的居所。那是一个用枯叶和丝做成的圆屋顶。造得很深，蜘蛛的身体几乎全部隐
藏在里面，用后端身体堵住进口。
它的前半身埋在它的居所里，所以，它当然看不到网上的动静了——即使它有一双
敏锐的眼睛也未必看得见，何况它其实是个半瞎子呢！那么在阳光灿烂的白天，它是不
是就放弃捕食了呢？让我们再看看吧。
你瞧，它的一条后腿忽然伸出叶屋，后腿的顶端连着一根丝线，而那线正是电报线
的另一个端点！我敢说，无论是谁，如果没有看见过蜘蛛的这手绝活，即把手(即它的脚
端)放在电报接收器上的姿势，他就不会知道动物表现自己智慧的最有趣的一个例子。让
猎物在这张网上出现吧，让这位假寐的猎手感觉到电报传来的信号吧！我故意放了一只
蝗虫在网上——以后呢？一切都像我预料的那样，虫子的振动带动网的振动，网的振动
又通过丝线——“电报线”传导到守株待兔的蜘蛛的脚上。蜘蛛它为得到食物而满足，
而我比它更满意：因为我学到了我想学的东西。
还有一点值得讨论的地方。那蛛网常常要被风吹动，那么电报线是不是不能区分网
的振动是来自猎物的来临还是风的吹动呢？事实上，当风吹动引起电报线晃动的时候，
在居所里闭目养神的蜘蛛并不行动，它似乎对这种假信号不屑一顾。所以这根电报线的
另外一个神奇之处在于，它像一台电话，就像我们人类的电话一样，能够传来各种真实
声音。蜘蛛用一个脚趾接着电话线，用腿听着信号，还能分辨出囚徒挣扎的信号和风吹
动所发出的假信号。
------------------








蟹蛛



前面我们讲到过的条纹蜘蛛虽然工作很勤快，为了替它的卵造一个安乐窝，一直孜
孜不倦地废寝忘食地工作着。可是到了后来，它却不能再顾到它的家了。为什么呢？因
为它寿命太短。它在第一个寒流到达之时就要死了。而它的卵要过了冬天才能孵化。它
不得不丢下它的巢。如果小宝宝在母亲还在世的时候能出世，我相信母蛛对小蛛的细心
呵护不会亚于鸟类。另外一种蜘蛛证明了我的推测；它是一种不会织网的蜘蛛：只是等
着猎物跑近它才去捉，而且它是横着走路的，有点儿像螃蟹，所以叫蟹蛛。
这种蜘蛛不会用网猎取食物，它的捕食方法是：埋伏在花的后面等猎物经过，然后
上去在它颈部轻轻一刺，你别小看这轻轻的一刺，这能致它的猎物于死地。我所观察到
的这群蟹蛛尤其喜欢捕食蜜蜂。
蜜蜂采花蜜的时候是很专心致志的，什么都不会想不会开小差。它用舌头舔着花蜜，
然后挑选一个能采到许多花蜜的花蕊上，一心一意地开始工作。当它正埋头苦干的时候，
蟹蛛早就虎视耽耽从隐藏的地方偷偷地、悄悄地爬出来，走到蜜蜂背后，越走越近，然
后一下子冲上去，在蜜蜂颈背上的某一点刺了一下。蜜蜂无论怎么挣扎也摆脱不了那一
刺。
这一刺可不是随随便便出手的。它刚好刺在蜜蜂颈部的神经中枢上。蜜蜂的神经中
枢被麻痹了，腿也开始硬化，不能动弹了。一秒钟内，一个小生命就宣告结束了。蟹蛛
这个凶手快乐而满足地吸着它的血，吸完后抹抹嘴巴残酷地把它的尸体的遗骸丢在一边。
可是，我们将会看到这个残暴的刽子手在家里是个非常慈爱温和的母亲。就像“食
人怪”一样，虽然吃起别人的孩子来毫不留情，却十分疼爱自己的孩子。在饥饿的逼迫
下，我们也是一样的，人和畜牲都是食人怪。
蟹蛛虽然是个杀蜂不见血的凶手，但你又不得不承认，它也是一只非常漂亮的小东
西。虽然它们的身材并不十分好，像是一个雕在石基上的又矮的又胖的锥体，在其中的
一边还有一块小小的隆起的肉，好像骆驼的驼蜂一样，但是它们的皮肤比任何绸缎都要
好看，有的是乳白色的，有的是柠檬色的。它们中间有些特别漂亮：腿上有着粉红色的
环，背上镶着深红的花纹，有时候在胸的左边或右边还有一条淡绿色的带子，这身打扮
虽然不像条纹蛛那么富丽，但是由于它的肚子不那么松驰，花纹又细致，色彩鲜艳又搭
配协调，所以，看起来倒反而比条纹蛛的衣服典雅、高贵。人们见了别种蜘蛛都敬而远
之，但对美丽的蟹蛛却怎么也怕不起来，因为它实在长得太漂亮，太可爱了，如果它们
是一些不会动的小东西，大家一定会对它们爱不释手。

蟹蛛的巢
在筑巢方面，蟹蛛的高超手艺并不亚于它在觅食时的技艺。
有一次我在一株水蜡树上找到它，当时它正在一丛花的中间筑巢，它织着一只白色
的丝袋，形状像一个顶针：这个白色的丝袋就是它的卵的安乐窝，袋口上还盖着一个又
圆又扁的绒毛盖子。
在屋顶的上部有一个用绒线张成的圆顶，里面还夹杂着一些凋谢了的花瓣，这就是
它的了望台。从外面到了望台上，有一个开口作为通道。
就在这了望台上，蟹蛛像一名尽心尽责的卫兵一样，天天守在这里。自从产了卵之
后，它比以前消瘦多了，差不多完全失去了以前那朝气蓬勃的样子，它全神贯注地守在
这了望台上，一有风吹草动就会全身紧张，埋没入战备状态，然后从那儿走出来，挥着
一条腿威吓来惊扰它的不速之客，它激动地做着手势，叫它赶紧滚开，否则后果自负。
它那狰狞的样子和激动的动作的确把那些或怀有恶意或无辜的外来者吓了一跳，把那些
鬼鬼祟祟的家伙赶走以后，它才心满意足地回到自己的岗位上。
那么它又在那丝和花瓣做成的穹顶下做什么呢？原来它在舒展着身体来遮蔽它宝贝
的卵。尽管此时它已经非常瘦小孱弱了仿佛一阵风就能把它卷走，它已忘记了饮食，为
了守望工作不被影响，它现在已抛弃了睡眠，不再去捕蜜蜂，吸它们的血充饥，它只是
静静地坐在自己的卵上。
我不由想到母鸡。母鸡孵蛋的时候也是这个样子，不同的是母鸡的身体可以提供热
量，当它孵蛋的时候，它身上的暖气传导到卵上唤醒了生命的种子。而对于蜘蛛而言，
太阳提供的热量已经足够了。蜘蛛母亲不需要再提供热量了。事实上，它也没有什么力
量提供热能了。因为这个不同点，我们不能称蜘蛛对小蛛的守候为“孵育下一代”。
两三个星期后，母蛛因为滴水未进，一点没吃，而愈来愈来瘦了。可它的守望工作
却丝毫不见有松懈。在临死之前，它还有什么愿望呢？它似乎一直在等待什么，这等待
使它苦苦地支撑着自己，用它的精神撑起早已没有活力的身体。它究竟在等什么呢？是
什么值得它这样用生命去苦苦等待呢？后来我才知道它是在等它的孩子们出来，这个垂
死的母亲还能为孩子们尽一点力。
条纹蛛的孩子们在离开那气球形的巢之前早已成为孤儿，没有人来帮它们打破巢，
它们自己又没有能力破巢而出，只有等巢自动裂开才能把小蛛们送到四面八方。它们出
来时根本连自己的母亲是谁都不知道。而蟹蛛的巢封闭得很严密，又不会自动裂开，顶
上的盖也不会自动升起，那么小蛛是怎么出来的呢？等小蛛孵出后我们会发现在盖的边
缘有一个小洞。这个洞在以前的时候是没有的，显然是谁暗中帮助小蜘，为它们在盖子
上咬了一个小孔，便于让它们钻出来。可是又是谁悄悄地在那儿开了一个洞呢？
袋子的四壁又厚又粗，微弱的小蛛们自己是决不能把它抓破的，其实这洞是它们那
奄奄一息的母亲打的。当它感觉到袋子里的小生命不耐烦的骚动的时候，它知道孩子们
急于想出来，于是就用全身的力气在袋壁上打了一个洞。母蛛虽然衰弱得可以随时死去，
但为了为它的家庭尽最后一份力，它一直顽强地支撑了五六个星期。然后把全身的力量
积聚到一点上爆发出来打这个洞，这个任务完成之后，它便安然去世了。它死的时候非
常平静，脸上带着安祥的神情，胸前紧紧抱着那已没有用处的巢，慢慢地缩成一个僵硬
的尸体。多么伟大的母亲啊！母蚁的牺牲精神令人感动，可是和蟹蛛相比，似乎还略逊
一筹。
在七月里，我的实验室中的小蟹蛛从卵里出来了。我知道它们有攀绳的嗜好，所以
我把一捆细树枝插在它们的笼上。果然，它们立刻沿着铁笼很快地爬到树枝的顶端，又
很快地用交叉的丝线织成互相交错的网，这便是它们的空中沙发。它们安静地在这沙发
上休息了几天，后来它们就开始搭起吊桥来。
我把爬着许多小蛛的树枝拿到窗口的一张桌子上，然后把窗户打开。不久小蛛们便
开始纺线做它们的飞行工具了，不过它们做得很慢，因为它们总是三心二意的，一会儿
爬到树枝下面，一会儿又回到顶上，好像不知道自己要干什么，又不知道该怎么干。
照这种速度，它们在那儿忙乎了半天也没什么成果，它们都急于要飞出去，可是就
是没胆量。在十一点钟的时候，我把载着它们的树枝拿到窗栏上，使让太阳照射到它们
身体上。几分钟以后，太阳的光和热射入它们的身体后积聚起来，成为一个小发动机，
驱使小蛛们纷纷活跃起来。只见它们的动作越来越快，越来越敏捷，都一个劲地往树枝
的顶上爬去，尽管我不能确切地看到它纺着线往空中飞去，但我很相信它们此刻正在树
梢上飞快地纺线，蓄势待发呢！
有三四只蟹蛛同时出发了，但各走各的路，自向不同的方向，其余的也纷纷爬到顶
上，后面拖着丝。突然起了一阵风，啊呀，那些蜘蛛那样的轻巧，它们编的丝又那么细，
风会把它们卷走吗？
我仔细看了看，风的确猝然把细丝扯断了，小蛛们顺着风在空中飘荡了一会儿便随
着它们的降落伞——断丝飘走了。我望着它们离去的背影，直到它在我的视野里消失。
它越飞越远，离出发点有四十尺远了。在又黑又暗的柏树叶丛中，它犹如一颗闪亮的明
星。它越飞越高，越飞越远，终于看不见了。其余的小蛛也接着飞出去，有些飞得很高，
有些飞得很低。有的往这边，有的往那边，最终都找到了自己的安身立命之处。
这时候，所有的小蛛都准备起飞了。而现在已不是开始的时候那样三三两两地飞出，
而是呈放射线状一队一队地飞出了，也许几个先锋的英雄行为感染激励了它们。不久它
们就陆续安全着陆了，有的在远处，有的在近处，这个简单的降落伞成功地完成了它的
使命。
关于它们以后故事，我就不知道了。在它们还没有力量刺蜜蜂的时候，它们怎么捕
食小虫子呢？小虫子和小蜘蛛争斗的话，谁又会占上风呢？它会耍什么阴谋吗？它会受
哪些天敌的威胁呢？我都不得而知。不过等到明年夏天，我们可以看到它已经长得很肥
很大，纷纷躲在花丛里偷袭那些勤劳采蜜的蜜蜂了。

------------------











昆虫记 

作者：J·H·法布尔[法] 


论祖传



人人都有自己的才能和自己的性格。有的时候这种性格看起来好像是从我们的祖先
那里遗传下来的，然而要想再追究这些性格是来源于何处，却又是一件非常非常困难的
事情。
例如，有一天看到一个牧童，他正低声地数着一颗颗小石子，计算这些小石子的总
数，把这当做一种消遣，于是他长大后竟然成了十分著名的教授，最后，他也许可以成
为数学家。另外又有一个孩子，他的年龄比起别的小孩子们也大不了多少，别的孩子们
只注意玩闹的事情，然而他却不和别的小孩子们在一起玩儿，而是整日幻想一种乐器的
声音，于是当他独自一人的时候，竟听到一种神秘的合奏曲子了。可见这个小孩是很有
音乐天才的。第三个小孩，长得又小又瘦，年龄也很小，也许他吃面包和果酱时，还会
不小心涂到脸上，但他竟然有他独自的爱好——喜欢雕塑粘土，制成各种各样的小模型，
这些小模型被他雕塑得各具形态。如果这个小孩子运气好的话，他将来总有一天会成为
一名著名的雕刻家的。
我知道，在背后议论别人的私事，是十分让人讨厌的一种行为，但是我想也许大家
能允许我来讲一番，并借这个机会来介绍我自己和我的研究。
在我很小很小的时候，我已经有一种与自然界的事物接近的感觉。如果你认为我的
这种喜欢观察植物和昆虫的性格是从我的祖先那里遗传下来的，那简直是一个天大的笑
话，因为，我的祖先们都是没有受过教育的乡下佬，对其他的东西都一无所知。他们唯
一知道和关心的，就是他们自己养的牛和羊。在我的祖父辈之中，只有一个人翻过书本
儿，甚至就连他对于字母的拼法在我看来也是十分不可信的。至于如果要说到我曾经受
过什么专门的训练，那就更谈不上了，从小就没有老师教过我，更没有指导者，而且也
常常没有什么书可看。不过，我只是朝着我眼前的一个目标不停地走，这个目标就是有
朝一日在昆虫的历史上，多少加上几页我对昆虫的见解。
回忆过去，在很多年以前，那时候我还是一个不懂事的小孩子，那时我才刚刚学会
认字母，然而，我对于当时我那种初次学习的勇气和决心，至今都感到非常骄傲。
我记得很清楚的一次经历是我第一次去寻找鸟巢和第一次去采集野菌的情景，当时
那种高兴的心情真令我直到今天还难以忘怀。
记得有一天，我去攀登离我家很近的一座山。在这座山顶上，有一片很早就引起我
浓厚兴趣的树林，从我家的小窗子里看出去，可以看见这些树木朝天立着，在风中摇摆，
在雪里弯腰，我很早就想能有机会跑到这些树林那儿去看一看了。这一次的爬山，爬了
好长的时间，而我的腿又很短，所以爬的速度十分缓慢，草坡十分陡峭，就跟屋顶一样。
忽然，在我的脚下，我发现了一只十分可爱的小鸟。我猜想这只小鸟一定是从它藏
身的大石头上飞下来的。不到一会儿工夫，我就发现了这只小鸟的巢。这个鸟巢是用干
草和羽毛做成的，而且里面还排列着六个蛋。这些蛋具有美丽的纯蓝色，而且十分光亮，
这是我第一次找到鸟巢，是小鸟们带给我许多的快乐中的第一次。我简直高兴极了，于
是我伏在草地上，十分认真地观察它。
这时候，母鸟十分焦急的在石上飞来飞去，而且还“塔克！塔克！”地叫着，表现
出一种十分不安的样子。我当时年龄还太小，甚至还不能懂得它为什么那么痛苦，当时
我心里想出了一个计划，我首先带回去一只蓝色的蛋，作为纪念品。然后，过两星期后
再来，趁着这些小鸟还不能飞的时候，将它们拿走。我还算幸运，当我把蓝鸟蛋放在青
苔上，小心翼翼地走回家时，恰巧遇见了一位牧师。
他说“呵！一个萨克锡柯拉的蛋！你是从哪里捡到这只蛋的？”
我告诉他前前后后捡蛋的经历，并且说：“我打算再回去拿走其余的蛋，不过要等
到当新生出的小鸟们刚长出羽毛的时候。”
“哎，不许你那样做！”牧师叫了起来；“你不可以那么残忍，去抢那可怜母鸟的
孩子。现在你要做一个好孩子，答应我从此以后再也不要碰那个鸟巢。”
从这一番谈话当中，我懂得了两件事。第一件，偷鸟蛋是件残忍的事。第二件，鸟
兽同人类一样，它们各自都有各自的名字的。
于是我自己问自己道：“在树林里的，在草原上的，我的许多朋友，它们是叫什么
名字呢？萨克锡柯拉的意思是什么呢？”
几年以后，我才晓得萨克锡柯拉的意思是岩石中的居住者，那种下蓝色蛋的鸟是一
种被称为石鸟的鸟儿。
有一条小河沿着我们的村子旁边悄悄地流过，在河的对岸，有一座树林，全是光滑
笔直的树木，就像高高耸立的柱子一般，而且地上铺满了青苔。
在这座树林里，我第一次采集到了野菌。这野菌的形状，猛一眼看上去，就好像是
母鸡生在青苔上的蛋一样。还有许多别的种类的野菌形状不一，颜色也各不相同。有的
形状长得像小铃儿，有的形状长得像灯泡，有的形状像茶杯，还有些是破的，它们会流
出像牛奶一样的泪，有些当我踩到它们的时候，变成蓝蓝的颜色了。其中，有一种最稀
奇的，长得像梨一样，它们顶上有一个圆孔，大概是一种烟筒吧。我用指头在下面一戳，
会有一簇烟从烟筒里面喷出来，我把它们装满了好大一袋子，等到心情好的时候，我就
把它们弄得冒烟，直到后来它们缩成一种像火绒一样的东西为止。
在这以后，我又好几次回到这片有趣的树林。我在乌鸦队里，研究真菌学的初步功
课，通过这种采集所得到的一切，是呆在房子里不可能获得的。
在这种一边观察自然与一边做试验的方法相结合的情况之下，我的所有功课，除两
门课，差不多都学过了。我从别人那里，只学过两种科学性质的功课，而且在我的一生
中，也只有这两种：一种是解剖学，一种是化学。
第一种是我得力于造诣很深的自然科学家摩根·斯东，他教我如何在盛水的盆中看
蜗牛的内部结构。这门功课的时间很短，但是能学到很多东西。
我初次学习化学时，运气就比较差了。在一次实验中，玻璃瓶爆炸，使多数同学受
了伤，有一个人眼睛险些儿瞎了，老师的衣服也被烧成了碎片，教室的墙上沾污了许多
斑点。后来，我重新回到这间教室时，已经不是学生而是教师了，墙上的斑点却还留在
那里。这一次，我至少学到了一件事，就是以后我每做一种试验，总是让我的学生们离
开远一点。
我有一个最大的愿望，就是想在野外建立一个试验室。当时我还处于在为每天的面
包问题而发愁的生活状况下，这真是一件不容易办到的事情！我几乎四十年来都有这种
梦想，想拥有一块小小的土地，把土地的四面围起来，让它成为我私人所有的土地；寂
寞、荒凉、太阳曝晒、长满荆草，这些都是为黄蜂和蜜蜂所喜好的环境条件。在这里，
没有烦扰，我可以与我的朋友们，如猎蜂手，用一种难解的语言相互问答，这当中就包
含了不少观察与试验呢。
在这里，也没有长的旅行和远足，以至于白白浪费了时间与精力，这样我就可以时
时留心我的昆虫们了！
最后，我实现了我的愿望。在一个小村落的幽静之处，我得到了一小块土地。这是
一块哈麻司，这个名字是给我们洽布罗温司的一块不能耕种，而且有许多石子的地方起
的。那里除了一些百里香，很少有植物能够生长起来。如果花费功夫耕耘，是可以长出
东西的，可是实在又不值得。不过到了春天会有些羊群从那里走过，如果碰巧当时下点
雨，也是可以生长一些小草的。
然而，我自己专有的哈麻司，却有一些掺着石子的红土，并且曾经被人粗粗地耕种
过了。有人告诉我说，在这块地上生长过葡萄树，于是我心里真有几分懊恼，因为原来
的植物已经被人用二脚叉弄掉了，现在已经没有百里香了。百里香对于我也许有用，因
为可以用来做黄蜂和蜜蜂的猎场，所以我不得已又把它们重新种植起来。
这里长满了偃卧草、刺桐花、以及西班牙的牡莉植物——那是长满了橙黄色的花，
并且有硬爪般的花序的植物。在这些上面，盖着一层伊利里亚的棉蓟，它那耸然直立的
树枝干，有时长到六尺高，而且末梢还长着大大的粉红球，还带有小刺，真是武装齐备，
使得采集植物的人不知应从哪里下手摘取才好。在它们当中，有穗形的矢车菊，长了好
长一排钩子，悬钩子的嫩芽爬到了地上。假使你不穿上高筒皮鞋，就来到有这么多刺的
树林里，你就要因为你的粗心而受到惩罚了。
这就是我四十年来拼命奋斗得来的属于我的乐园啊！
在我的这个稀奇而又冷清的王国里，是无数蜜蜂和黄蜂的快乐的猎场，我从来没有
在单独的一块地方，看见过这么多的昆虫。各种生意都以这块地为中心，来了猎取各种
野味的猎人、泥土匠、纺织工人、切叶者、纸板制造者，同时也有石膏工人在拌和泥灰，
木匠在钻木头，矿工在掘地下隧道，以及牛的大肠膜(用来隔开金箔)工人，各种各样的
人都有。
快看啊！这里有一种会缝纫的蜜蜂。它剥下开有黄花底的刺桐的网状线，采集了一
团填充的东西，很骄傲地用它的腮(即颚)带走了。它准备到地下，用采来的这团东西储
藏蜜和卵。那里是一群切叶蜂，在它们的身躯下面，带着黑色的，白色的，或者血红色
的，切割用的毛刷，它们打算到邻近的小树林中，把树叶子割成圆形的小片用来包裹它
们的收获品。这里又是一群穿着黑丝绒衣的泥水匠蜂，它们是做水泥与沙石工作的。在
我的哈麻司里我们很容易在石头上发现它们工作用的工具。另外，这有一种野蜂，它把
窝巢藏在空蜗牛壳的盘梯里。还有一种，把它的蛴螬安置在干燥的悬钩子的秆子的木髓
里。第三种，利用干芦苇的沟道做它的家。至于第四种，住在泥水匠蜂的空隧道中，而
且连租金都用不着付。还有的蜜蜂生着角，有些蜜蜂后腿头上长着刷子，这些都是用来
收割的。
我的哈麻司的墙壁建筑好了，到处可以看到成堆成堆的石子和细沙，这些全是建筑
工人们堆弃下来的，并且不久就被各种住户给霸占了。泥水匠蜂选了个石头的缝隙，用
来做它们睡眠的地方。若是有凶悍的蜥蜴，一不小心压到它们的时候，它们就会去攻击
人和狗。它们挑选了一个洞穴，伏在那里等待路过的蜣螂。黑耳毛的鸫鸟，穿着白黑相
间的衣裳，看上去好像是黑衣僧，坐在石头顶上唱简单的歌曲。那些藏有天蓝色的小蛋
的鸟巢，会在石堆的什么地方才能找到呢？当石头被人搬动的时候，在石头里面生活的
那些小黑衣僧自然也一块儿被移动了。我对这些小黑衣僧感到十分惋惜，因为它们是很
可爱的小邻居。至于那个蜥蜴，我可不觉得它可爱，所以对于它的离开，我心里没有丝
毫的惋惜之情。
在沙土堆里，还隐藏了掘地蜂和猎蜂的群落，令我感到遗憾的是，这些可怜的掘地
蜂和猎蜂们后来无情地被建筑工人给无辜地驱逐走了。但是仍然还有一些猎户们留着，
它们成天忙忙碌碌，寻找小毛虫。还有一种长得很大的黄蜂，竟然胆大包天地敢去捕捉
毒蜘蛛，在哈麻司的泥土里，有许多这种相当利害的蜘蛛居住着。而且你可以看到，还
有强悍勇猛的蚂蚁，它们派遣出一个兵营的力量，排着长长的队伍，向战场出发，去猎
取它们强大的俘虏。
此外，在屋子附近的树林里面，住满了各种鸟雀。它们之中有的是唱歌鸟，有的是
绿莺，有的是麻雀，还有猫头鹰。在这片树林里有一个小池塘，池中住满了青蛙，五月
份到来的时候，它们就组成振耳欲聋的乐队。在居民之中，最最勇敢的要数黄蜂了，它
竟不经允许地霸占了我的屋子。在我的屋子门口，还居住着白腰蜂。每次当我要走进屋
子里的时候，我必须十分小心，不然就会踩到它们，破坏了它们开矿的工作。在关闭的
窗户里，泥水匠蜂在软沙石的墙上建筑土巢。我在窗户的木框上一不小心留下的小孔，
被它们利用来做门户。在百叶窗的边线上，少数几只迷了路的泥水匠蜂建筑起了蜂巢。
午饭时候一到，这些黄蜂就翩然来访，它们的目的，当然是想看看我的葡萄成熟了没有。
这些昆虫全都是我的伙伴，我的亲爱的小动物们，我从前和现在所熟识的朋友们，
它们全都住在这里，它们每天打猎，建筑窝巢，以及养活它们的家族。而且，假如我打
算移动一下住处，大山离我很近，到处都是野草莓树、岩蔷薇和石楠植物，黄蜂与蜜蜂
都是喜欢聚集在那里的。我有很多理由，使我为了乡村而逃避都市，来到西内南，做些
除杂草和灌溉莴苣的事情。
------------------








神秘的池塘



当我面对池塘，凝视着它的时候，我可从来都不觉得厌倦。在这个绿色的小小世界
里，不知道会有多少忙碌的小生命生生不息。在充满泥泞的池边，随处可见一堆堆黑色
的小蝌蚪在暖和的池水中嬉戏着，追逐着；有着红色肚皮的蝾螈也把它的宽尾巴像舵一
样地摇摆着，并缓缓地前进；在那芦苇草丛中，我们还可以找到一群群石蚕的幼虫，它
们各自将身体隐匿在一个枯枝做的小鞘中——这个小鞘是用来作防御天敌和各种各样意
想不到的灾难用的。
在池塘的深处，水甲虫在活泼地跳跃着，它的前翅的尖端带着一个气泡，这个气泡
是帮助它呼吸用的。它的胸下有一片胸翼，在阳光下闪闪发光，像佩带在一个威武的大
将军胸前的一块闪着银光的胸甲。在水面上，我们可以看到一堆闪着亮光的“蚌蛛”在
打着转，欢快地扭动着，不对，那不是“蚌蛛”，其实那是豉虫们在开舞会呢！离这儿
不远的地方，有一队池鳐正在向这边游来，它们那傍击式的泳姿，就像裁缝手中的缝针
那样迅速而有力。
在这个地方你还会见到水蝎，只见它交叉着两肢，在水面上悠闲地做出一副仰泳的
姿势，那神态，仿佛它是天底下最伟大的游泳好手。还有那蜻蜓的幼虫，穿着沾满泥巴
的外套，身体的后部有一个漏斗，每当它以极高的速度把漏斗里的水挤压出来的时候，
借着水的反作用力，它的身体就会以同样的高速冲向前方。
在池塘的底下，躺着许多沉静又稳重的贝壳动物。有时候，小小的田螺们会沿着池
底轻轻地、缓缓地爬到岸边，小心翼翼地慢慢张开它们沉沉的盖子，眨巴着眼睛，好奇
地展望这个美丽的水中乐园，同时又尽情地呼吸一些陆上空气；水蛭们伏在它们的征服
物上，不停地扭动着它们的身躯，一副得意洋洋的样子；成千上万的孑孓在水中有节奏
地一扭一曲，不久的将来它们会变成蚊子，成为人人喊打的坏蛋。
乍一看，这是一个停滞不动的池塘，虽然它的直径不超过几尺，可是在阳光的孕育
下，它却犹如一个辽阔神秘而又丰富多彩的世界。它多能打动和引发一个孩子的好奇心
啊！让我来告诉你，在我的记忆中的第一个池塘怎样深深地吸引了我，激发起我的好奇
心。
我小的时候，家里很穷。除了我妈妈继承的一所房子和一块小小的荒芜的园子之外，
几乎什么也没有了。“我们将怎么生活下去呢？”这个严重的问题，常常会挂在我爸爸
妈妈的嘴边。
你听说过“大拇指”的故事吗？那个“大拇指”藏在他父亲的矮凳子下，偷听他父
亲和母亲所说的一些关于生活窘迫的对话。我就很像那个“大拇指”。但是我没有像他
那样，可以藏在凳子底下，我是伏在桌子上一面假装睡着了，一面偷听他们的谈话。幸
运的是，我所听到的，并不像“大拇指”的父亲所说的那种使人心寒的话，相反地，那
是一个美妙的计划。我听了以后，心中涌起一阵难以形容的快乐和欣慰。
“如果我们来养一群小鸭”，妈妈说，“将来一定可以换得不少钱。我们可以买些
油脂回来，让亨利天天照料它们，把它们喂得肥肥的。”
“太好了！”父亲高兴地说道，“让我们来试试吧。”
那天晚上，我做了一个美妙的梦。我和一群可爱的小鸭子们一起漫步到池畔，它们
都穿着鲜黄色的衣裳，活泼地在水中打闹、洗澡。我在旁边微笑地看着它们洗澡，耐心
地等它们洗痛快，然后带着它们慢悠悠地走回家。半路上，我发现其中有一只小鸭累了，
就小心翼翼地把它捧起来放在篮子里面，让它甜甜地睡觉。
没想到就在两个月之后，我的美梦就实现了：我们家里养了二十四只毛茸茸的小鸭
子。鸭子自己不会孵蛋，常常由母鸡来孵。可怜的老母鸡分不出孵的是自己的亲骨肉还
是别家的“野孩子”，只要是那圆溜溜、和鸡蛋差不多样子的蛋，它都很乐意去孵，并
把孵出来的小生物当作自己的亲生孩子来对待。负责孵育我们家的小鸭的是两只黑母鸡，
其中一只是我们自己家的，而另一只是向邻居借来的。
我们家的那只黑母鸡，每天陪着小鸭们玩，不厌其烦地和它们做游戏玩耍，让它们
快乐健康地长大。我往一只木桶里盛了些水，大约有两寸高，这个木桶就成了小鸭们的
游泳池。只要是晴朗的日子，小鸭们总是一边沐浴着温暖的阳光，一边在木桶里洗澡嬉
戏，显得无比的美满、和谐和舒适，令旁边的黑母鸡羡慕不已。
两星期以后，这只小小的木桶渐渐地不能满足小鸭们的要求了。它们需要大量的水，
这样它们才能在里面自由自在地翻身跳跃，它们还需要许多小虾米、小螃蟹、小虫子之
类作为它们的食物。而这些食物通常大量地蕴藏在互相缠绕的水草中，等候着它们自己
去猎取。现在我感觉到取水是个大问题。因为我们家住在山上，而从山脚下带大量的水
上来是非常困难的。尤其是在夏天，我们自己都不能痛快地喝水，哪里还顾得了那些小
鸭呢？
虽然在我们家附近也有一口井，可那是一口半枯的井，每天要供四五家邻居轮流使
用，还有学校里的校长先生养的那头驴子，它总是贪得无厌地对着那井水大口大口地喝
水，那口井很快就被喝干了。直到整整一昼夜之后，才看见有井水渐渐地升起来，恢复
到原来的样子。在这么艰难的水荒中，我们可怜的小鸭子自然就没有自由嬉水的份了。
不过，在那山脚下，有一条潺潺的小溪。那倒是小鸭们的天然乐园。可是从我们家
到那小溪，必须穿过一条村里的小路，可是我们不能走那条小路，因为在那条路上我们
很可能会碰到几只凶恶的猫和狗，它们会毫不犹豫地冲散小鸭们的队伍，使我没法把它
们重新聚拢在一起。于是，我只得另谋出路。我想起在离山不远的地方，有一块很大的
草地和一个很不小的池塘。那是一个很荒凉很偏僻的地方，没有什么猫狗的打扰，的确
可以成为小鸭们的乐园。
我第一天做牧童，心中又快活又自在。不过有件事很令我难受，那就是赤裸裸的双
脚，渐渐地起泡了。因为跑了太多的路，我又不能把箱子里那双鞋子拿出来穿。那双宝
贵的鞋子只有在过节的时候才能穿。我赤裸的脚不停地在乱石杂草中奔跑，伤口越来越
大了。
小鸭们的脚似乎也受不了这么折腾，因为它们的蹼还没有完全长成，还远不够坚硬。
当它们走在这么崎呕的山路上时，不时地发出“呷呷——”的叫声，似乎是在请求我允
许它们休息一下。每当这个时候，我也只得满足它们的要求，招呼它们在树荫下歇歇脚，
否则恐怕它们再也没有力气走完剩下的路了。
我们终于到达了目的地。那池水浅浅的，温温的，水中露出的土丘就好像是一个个
小小的岛屿。小鸭们一到那儿就飞奔过去忙碌地在岸上寻找食物。吃饱喝足后，它们会
下到水里去洗澡。洗澡的时候，它们常常会把身体倒竖起来，前身埋在水里，尾巴指向
空中，仿佛在跳水中芭蕾。我美滋滋地欣赏着小鸭们优美的动作，看累了，就看看水中
别的景物。
那是什么？在泥土上，我看到有几段互相缠绕着的绳子又粗又松，黑沉沉的，像熏
满了的烟灰。如果你看到它，可能会以为它是从什么袜子上拆下来的绒线。于是我想：
可能是哪位牧羊女在水边编一只黑色的绒线袜子，突然发现某些地方漏了几针，不能往
下编了，埋怨了一阵子后，就决心全部拆掉，重新开始，而在她拆得不耐烦的时候，就
索性把这编坏的部分全丢在水里。这个推测看起来合情合理。
我走过去，想拾一段放到手掌里仔细观察，没想到这玩意儿又粘又滑，一下子就从
我的手指缝里滑走了。我花费了好大的劲，就是捉不住它，并且有几段绳子的结突然散
了，从里面跑出一颗颗小珠子，只有针尖般大小，后面拖着一条扁平的尾巴，我一下子
认出它们了，那是我所熟悉的一种动物的幼虫。它就是青蛙的幼虫——蝌蚪。
在这里我还看到了许多别的生物。其中有一种不停地在水面上打旋，它的黑色的背
部在阳光下发着亮光。每当我伸手去捉它们的时候，它们似乎早就预料到危险来临似的，
不等我碰它们，就逃得无影无踪了。我本想捉几个放到碗里面仔细研究，可惜就是捉不
到它们。
看啦！在那池水深处，有一团绿绿的、浓浓的水草，我轻轻拨开一束水草，看到立
刻有许多水珠争先恐后地浮到水面聚成一个大大的水泡。我想在这厚厚的水草底下一定
藏着什么奇怪的生物。我继续往下探索，看到了许多贝壳像豆子一样扁平，周围冒着几
个涡圈；有一种小虫看上去像戴了羽毛；还有一种小生物舞动着柔软的鳍片，像穿着华
丽的裙子在跳舞。我也不知道它们为什么这样不停地游来游去，也不晓得它们叫什么，
我只能出神地对着这个神秘、玄妙的水池，浮想联翩。
池水通过小小的渠道缓缓地流入附近的田地，那儿长着几棵赤杨，我又在那儿发现
了美丽的生物，那是一只甲虫，像核桃那么大，身上带着一些蓝色。那蓝色是如此的赏
心悦目，使我联想起了那天堂里美丽的天使，她的衣服一定也是这种美丽的蓝色。我怀
着虔诚的心情轻轻地捉起它，把它放进了一个空的蜗牛壳，用叶子把它塞好。我要把它
带回家中，细细欣赏一番。
接着我的注意力又被别的东西吸引住了。清澈又凉爽的泉水源源不断地从岩石上流
下来，不停地滋润着这个池塘。泉水先流到一个小小的潭里，然后汇成一条小溪。我看
着看着就突发奇想，觉得这样让溪水默默地流过就太可惜了。可以把它看作一个小小的
瀑布，去推动一个磨。于是，我就开始着手做一个小磨。我用稻草做成轴，用两个小石
块支着它，不一会儿就完工了。这个磨子做得很成功，只可惜当时没有小伙伴和我一起
玩，只有几只小鸭来欣赏我的杰作。
这个小小的成功大大地激发了我的创造欲望，一发不可收拾。我又计划筑一个小水
坝，那里有许多乱石可以利用，我耐心地挑选着可以用来筑坝的石块，挑着挑着，我忽
然发现了一个奇迹，它使我再也无心去继续建造水坝的事了。
当我打开一个大石头时，有一个小拳头那么大的窟窿，从窟窿里面发出一簇簇光环，
好像是一簇簇钻石的小面在阳光照耀下闪着耀眼的光，又好像是教堂里彩灯上垂下来的
一串串晶莹剔透的珠子。
多么灿烂而美丽的东西啊。它使我想起孩子们躺在打禾场的干草上所讲的神龙传奇
的故事。神龙是地下宝库的守护者，它们守护着不计其数的奇珍异宝。现在在我眼前闪
光的这些东西，会不会就是神话中所说的皇冠和首饰呢？难道它们就是蕴藏在这些砖石
中吗？在这些破碎的砖石中，我可以搜集到许多发光的碎石，这些都是神龙赐给我的珍
宝啊！我仿佛觉得神龙在召唤我，要给我数不清的金子。在潺潺的泉水下，我看见许多
金色的颗粒，它们都粘在一片细砂上。我俯下身子仔细观察，发现这些金粒在阳光下随
着泉水打着转，这真是金子吗？真是那可以用来制造二十法郎金币的金子吗？对一个贫
穷的家庭来说，这金币是多么的宝贵啊！
我轻轻地拣起一些细砂，放在手掌中。这发光的金粒数量很多，但是颗粒却很小，
得把麦杆用唾沫浸湿了，才能用来沾住它们。我不得不放弃这项麻烦的工作。我想一定
有一大块一大块的金子深藏在山石中，可以等到以后我来把山炸毁了再说，这些小金粒
太微不足道了，我才不去拣它们呢！
我继续把砖石打碎，看看里面还有什么，可是这下我看到的不是珠宝，我看到有一
条小虫从碎片里爬出来，它的身体是螺旋形的，带着一节一节的疤痕，像一条蜗牛在雨
天的古墙里蜿蜒着爬到墙外，那有节疤的地方显得格外沧桑和强壮。我不知道它们是怎
样钻进这些砖石内部的，也不知道它们钻进去干嘛。
为了纪念我发现的“宝藏”，再加上好奇心的驱使，我把砖石装在口袋里，塞得满
满的。这时候，天快黑了，小鸭们也吃饱了，于是我对它们说：“来，跟着我，我们得
回家了。”
我的脑海里装满了幻想，脚跟的疼痛早已忘记了。
在回去的路上，我尽情地想着我的蓝衣甲虫，像蜗牛一样的甲虫，还有那些神龙所
赐的宝物。可是一踏进家门，我就回过神来，父母的反应令我一下子很失望。他们看见
了我那膨胀的衣袋里面尽是一些没有用处的砖石，我的衣服也快被砖石撑破了。
“小鬼，我叫你看鸭子，你却自顾自地去玩耍，你捡那么多砖石回来，是不是还嫌
我们家周围的石头不够多啊？赶紧把这些东西扔出去！”父亲冲着我吼道。
我只好遵照父亲的命令，把我的那些珍宝、金粒、羊角的化石和天蓝色的甲虫统统
抛在门外的废石堆里，母亲看着我，无奈地叹了口气。
“孩子，你真让我为难。如果你带些青菜回来，我倒也不会责备你，那些东西至少
可以喂喂兔子，可这种碎石，只会把你的衣服撑破，这种毒虫只会把你的手刺伤，它们
究竟能给你什么好处呢？蠢货！准是什么东西把你迷住了！”
可怜的母亲，她说得不错，的确有一种东西把我迷住了——那是大自然的魔力。几
年后，我知道了那个池塘边的“钻石”其实是岩石的晶体；所谓的“金粒”，原来也不
过是云母而己，它们并不是什么神龙赐给我的宝物。尽管如此，对于我，那个池塘始终
保持着它的诱惑力，因为它充满了神秘，那些东西在我看来，其魅力远胜于钻石和黄金。

玻璃池塘

你有一处建在房子里面的小池塘吗？在那个小池塘里，你可以随时观察水中生物生
活的每一个片断。它没有像户外的池塘那么大，也没有太多的生物，可这些恰恰又为观
察提供了有利条件。除此之外，还不会有行人来打扰你专注的观察。其实这并不是什么
天方夜谈，这是很容易实现的。
我的户内池塘是在铁匠和木匠的合作下造成的：先用铁条做好池架，把它装在木头
做的基座上面。池上面盖着一块可以活动的木板，下面的池底是铁做的，底上有一个排
水的小洞。池的四周镶着玻璃。这是一个设计得相当不错的玻璃池，就放在我的窗口，
它的体积大约有十到十二加仑。
我先往池里放进一些滑腻腻的硬块。那是一种分量很重的东西，表面长着许多小孔，
看上去很像珊瑚礁。硬块上面盖着许多绿绿的绒毛般的苔藓，这苔藓能够使水保持清洁，
为什么呢？让我们来看一看吧。
动物在水池里和我们在空气中一样，要吸入新鲜的气味，同时，排出废气(二氧化碳)。
这些废气是不适宜人呼吸的。而植物刚好相反，它们吸入二氧化碳。所以池中的水草就
吸收这种不可以呼吸的废气，经过一番工作后，释放出可以供动物呼吸的氧气。
如果你在充满阳光的池边站一会儿，你就能观察到这种变化，在有水草的珊瑚礁上，
那一点点发亮的闪烁的星光，好像是绿苗遍地的草坪上点缀着的零零碎碎的珍珠。这些
珍珠不断地消逝，又接连不断地出现，它们会倏然在水面上飞散开来，好像水底下发生
了小小的爆炸，冒出一串串的气泡。
水草分解了水中的二氧化碳，得到炭元素，炭可以用来制造淀粉。淀粉是生物细胞
所不可缺少的东西。营养物水草所吐出来的废气是新鲜的氧气。这些氧气一部分溶解在
水中，供给水中的生物呼吸，一部分离开水面跑到空气中。你在外面看到的像珍珠一样
的气泡就是氧气！
我注视着池水中的气泡，作了一番遐想：在许多许多年以前，陆地刚刚脱离了海洋，
那时草是第一棵植物，它吐出第一口氧气，供给生物呼吸。于是各种各样的动物相继出
现了，而且一代一代繁衍、变化下去，一直形成今天的生物世界。我的玻璃池塘似乎在
告诉我一个行星航行在没有氧气的空间里的故事。
------------------








石蚕



我往我的玻璃池塘里放进一些小小的水生动物，它们叫石蚕。确切地说，它们是石
蚕蛾的幼虫，平时很巧妙地隐藏在一个个枯枝做的小鞘中。
石蚕原本是生长在泥潭沼泽中的芦苇丛里的。在许多时候，它依附在芦苇的断枝上，
随芦苇在水中漂泊。那小鞘就是它的活动房子，也可以说是它旅行时随身带的简易房子。
这活动房子其实可以算得上是一个很精巧的编织艺术品，它的材料是由那种被水浸
透后剥蚀、脱落下来的植物的根皮组成的。在筑巢的时候，石蚕用牙齿把这种根皮撕成
粗细适宜的纤维，然后把这些纤维巧妙地编成一个大小适中的小鞘，使它的身体能够恰
好藏在里面。有时候它也会利用极小的贝壳七拼八凑地拼成一个小鞘，就好像一件小小
的百纳衣；有时候，它也用米粒堆积起来。布置成一个象牙塔似的窝，这算是它最华丽
的住宅了。

暴徒的袭击

石蚕的小鞘不但是它的寓所，同时还是它的防御工具。我曾在我的玻璃池塘里看到
一幕有趣的战争，鲜明地证实了那个其貌不扬的小鞘的作用。
玻璃池塘的水中原本潜伏着一打水甲虫，它们游泳的姿态激起了我极大的兴趣。有
一天，我无意中撒下两把石蚕，正好被潜在石块旁的水甲虫看见了，它们立刻游到水面
上，迅速地抓住了石蚕的小鞘，里面的石蚕感觉到此次攻击来势凶猛，不易抵抗，就想
出了金蝉脱壳的妙计，不慌不忙地从小鞘里溜出来，一眨眼间就逃得无影无踪了。
野蛮的水甲虫还在继续凶狠地撕扯着小鞘，直到知道早已失去了想要的食物，受了
石蚕的骗，这才显出懊恼沮丧的神情，无限留恋又无可奈何地把空鞘丢下，去别处觅食
了。
可怜的水甲虫啊！它们永远也不会知道聪明的石蚕早已逃到石底下，重新建造它的
新鞘，准备着你们的下一次袭击了。

潜水艇——石蚕

石蚕靠着它们的小鞘在水中任意遨游，它们好像是一队潜水艇，一会儿上升，一会
儿下降，一会儿又神奇地停留在水中央。它们还能靠着那舵的摆动随意控制航行的方向。
我不由想到了木筏，石蚕的小鞘是不是有木筏那样的结构，或是有类似于浮囊作用
的装备，使它们不致于下沉呢？
我将石蚕的小鞘剥去，把它们分别放在水上。结果小鞘和石蚕都往下沉。这是为什
么呢？
原来，当石蚕在水底休息时，它把整个身子都塞在小鞘里。当它想浮到水面上时，
它先拖带着小鞘爬上芦梗，然后把前身伸出鞘外。这时的小鞘的后部就留出一段空隙，
石蚕靠着这一段空隙便可以顺利往上浮。就好像装了一个活塞，向外拉时就跟针筒里空
气柱的道理一样。这一段装着空气的鞘就像轮船上的救生圈一样，靠着里面的浮力，使
石蚕不致于下沉。所以石蚕不必牢牢地粘附在芦苇枝或水草上，它尽可以浮到水面上接
触阳光，也可以在水底尽情遨游。
不过，石蚕并不是十分擅长游泳的水手，它转身或拐弯的动作看上去很笨拙。这是
因为它只靠着那伸在鞘外的一段身体作为舵桨，再也没有别的辅助工具了，当它享受了
足够的阳光后，它就缩回前身，排出空气，渐渐向下沉落了。
我们人类有潜水艇，石蚕也有这样一个小小的潜水艇。它们能自由地升降，或者停
留在水中央——那就是当它们在慢慢地排出鞘内的空气的时候。虽然它们不懂人类博大
精深的物理学，可这只小小的鞘造得这样的完美，这样的精巧，完全是靠它们的本能。
大自然所支配的一切，永远是那么巧妙和谐。
------------------








蜣螂



一、圆球
蜣螂第一次被人们谈到，是在过去的六七千年以前。古代埃及的农民，在春天灌溉
农田的时候，常常看见一种肥肥的黑色的昆虫从他们身边经过，忙碌地向后推着一个圆
球似的东西。他们当然很惊讶地注意到了这个奇形怪状的旋转物体，像今日布罗温司的
农民那样。
从前埃及人想象这个圆球是地球的模型，蜣螂的动作与天上星球的运转相合。他们
以为这种甲虫具有这样多的天文学知识，因而是很神圣的，所以他们叫它“神圣的甲虫”。
同时他们又认为，甲虫抛在地上滚的球体，里面装的是卵子，小甲虫是从那里出来的。
但是事实上，这仅是它的食物储藏室而已。里面并没有卵子。
这圆球并不是什么可口的食品。因为甲虫的工作，是从土面上收集污物，这个球就
是它把路上与野外的垃圾，很仔细地搓卷起来形成的。
做成这个球的方法是这样的：在它扁平的头的前边，长着六只牙齿，它们排列成半
圆形，像一种弯形的钉把，用来掘割东西。甲虫用它们抛开它所不要的东西，收集起它
所选拣好的食物。它的弓形的前腿也是很有用的工具，因为它们非常的坚固，而且在外
端也长有五颗锯齿。所以，如果需要很大的力量去搬动一些障碍物，甲虫就利用它的臂。
它左右转动它有齿的臂，用一种有力的扫除法，扫出一块小小的面积。于是，在那堆集
起了它所耙集来的材料。然后，再放到四支后爪之间去推。这些腿是长而细的，特别是
最后的一对，形状略弯曲，前端还有尖的爪子。甲虫再用这后腿将材料压在身体下，搓
动、旋转，使它成为一个圆球形。一会儿，一粒小丸就增到胡桃那么大，不久又大到像
苹果一样。我曾见到有些贪吃的家伙，把圆球做到拳头那么大。
食物的圆球做成后，必须搬到适当的地方去。于是甲虫就开始旅行了。它用后腿抓
紧这个球，再用前腿行走，头向下俯着，臀部举起，向后退着走。把在后面堆着的物件，
轮流向左右推动。谁都以为它要拣一条平坦或不很倾斜的路走。但事实并非如此！它总
是走险峻的斜坡，攀登那些简直不可能上去的地方。这固执的家伙，偏要走这条路。这
个球，非常的重，一步一步艰苦的推上，万分留心，到了相当的高度，而且它常还是退
着走的。只要有一些不慎重的动作，劳力就全白费了：球滚落下去，连甲虫也被拖下来
了。再爬上去，结果再掉下来。它这样一回又一回地向上爬，一点儿小故障，就会前功
尽弃，一根草根能把它绊倒，一块滑石会使它失足。球和甲虫都跌下来，混在一起，有
时经过一、二十次的继续努力，才得到最后的成功。有时直到它的努力成为绝望，才会
跑回去另找平坦的路。
有的时候，蜣螂好像是一个善于合作的动物，而这种事情是常常发生的。当一个甲
虫的球已经做成，它离开它的同类，把收获品向后推动。一个将要开始工作的邻居，看
到这种情况，会忽然抛下工作，跑到这个滚动的球边上来，帮球主人一臂之力。它的帮
助当然是值得欢迎的。但它并不是真正的伙伴，而是一个强盗。要知道自己做成圆球是
需要苦工和忍耐力的！而偷一个已经做成的，或者到邻居家去吃顿饭，那就容易多了。
有的贼甲虫，用很狡猾的手段，有的简直施用武力呢！
有时候，一个盗贼从上面飞下来，猛地将球主人击倒。然后它自己蹲在球上，前腿
靠近胸口，静待抢夺的事情发生，预备互相争斗。如果球主人起来抢球，这个强盗就给
它一拳，从后面打下去。于是主人又爬起来，推摇这个球，球滚动了。强盗也许因此滚
落。那末，接着就是一场角力比赛。两个甲虫互相扯扭着，腿与腿相绞，关节与关节相
缠，它们角质的甲壳互相冲撞，摩擦，发出金属互相摩擦的声音，胜利的甲虫爬到球顶
上，贼甲虫失败几回被驱逐后，只有跑开去重新做自己的小弹丸。有几回，我看见第三
个甲虫出现，像强盗一样抢劫这个球。
但也有时候，贼竟会牺牲一些时间，利用狡猾的手段来行骗。它假装帮助这个被驱
者搬动食物，经过生满百里香的沙地，经过有深车轮印和险峻的地方，但实际上它用的
力却很少，它做的大多只是坐在球顶上观光，到了适宜于收藏的地点，主人就开始用它
边缘锐利的头，有齿的腿向下开掘，把沙土抛向后方，而这贼却抱住那球假装死了。土
穴越掘越深，工作的甲虫看不见了。即使有时它到地面上来看一看，球旁睡着的甲虫一
动不动，觉得很安心。但是主人离开的时间久了，那贼就乘这个机会，很快的将球推走，
同小偷怕被人捉住一样快。假使主人追上了它——这种偷盗行为被发现了——它就赶快
变更位置，看起来好像它是无辜的，因为球向斜坡滚下去了，它仅是想止住它啊！于是
两个“伙伴”又将球搬回，好像什么事情都没有发生一样。
假使那贼安然逃走了，主人艰苦做起来的东西，只有自认倒霉。它揩揩颊部，吸点
空气，飞走，重新另起炉灶。我颇羡慕而且嫉妒它这种百折不挠的品质。
最后，它的食品才平安的储藏好了。储藏室是在软土或沙土上掘成的土穴。做的如
拳头般大小，有短道通往地面，宽度恰好可以容纳圆球。食物推进去，它就坐在里面，
进出口用一些废物塞起来，圆球刚好塞满一屋子，肴馔从地面上一直堆到天花板。在食
物与墙壁之间留下一个很窄的小道，设筵人就坐在这里，至多两个，通常只是自己一个。
神圣甲虫昼夜宴饮，差不多一个礼拜或两个礼拜，没有一刻停止过。

二、梨
我已经说过，古代埃及人以为神圣甲虫的卵，是在我刚才叙述的圆球当中的。这个
已经我证明不是如此。关于甲虫被放卵的真实情形，有一天碰巧被我发现了。
我认识一个牧羊的小孩子，他在空闲的时候，常来帮助我。有一次，在六月的一个
礼拜日，他到我这里来，手里拿着一个奇怪的东西，看起来好像一只小梨，但已经失掉
新鲜的颜色，因腐朽而变成褐色。但摸上去很坚固，样子很好看，虽然原料似乎并没有
经过精细的筛选。他告诉我，这里面一定有一个卵，因为有一个同样的梨，掘地时被偶
然弄碎，里面藏有一粒像麦子一样大小的白色的卵。
第二天早晨，天色才刚刚亮的时候，我就同这位牧童出去考察这个事实。
一个神圣甲虫的地穴不久就被找到了，或者你也知道，它的土穴上面，总会有一堆
新鲜的泥土积在上面。我的同伴用我的小刀铲向地下拼命的掘，我则伏在地上，因为这
样容易看见有什么东西被掘出来。一个洞穴掘开，在潮湿的泥土里，我发现了一个精制
的梨。我真是不会忘记这是我第一次看见，一个母甲虫的奇异的工作呢！当挖掘古代埃
及遗物的时候，如果我发现这神圣甲虫是用翡翠雕刻的，我的兴奋却也不见得更大呢。
我们继续搜寻，于是发现了第二个土穴。这次母甲虫在梨的旁边，而且紧紧抱着这
只梨。这当然是在它未离开以前，完工毕事的举动，用不着怀疑，这个梨就是蜣螂的卵
子了。在这一个夏季，我至少发现了一百个这样的卵子。
像球一样的梨，是用人们丢弃在原野上的废物做成的，但是原料要比较精细些，为
的是给蛴螬预备好食物。当它从卵里跑出来的时候，还不能自己寻找食物，所以母亲将
它包在最适宜的食物里，它可以立刻大吃起来，不至于挨饿。
卵是被放在梨的比较狭窄的一端的。每个有生命的种子，无论植物或动物，都是需
要空气的，就是鸟蛋的壳上也分布着无数个小孔。假如蜣螂的卵是在梨的最后部分，它
就闷死了，因为这里的材料粘得很紧，还包有硬壳。所以母甲虫预备下一间精制透气的
小空间，薄薄的墙壁，给它的小蛴螬居住，在它生命最初的时候，甚至在梨的中央，也
有少许空气，当这些已经不够供给柔弱的小蛴螬消耗，它要到中央去吃食，已经很强壮，
能够自己支配一些空气了。
当然，梨子大的一头，包上硬壳子，也是有很好的理由的。蜣螂的地穴是极热的，
有时候温度竟达到沸点。这种食物，经过三四个礼拜之后，就会干燥，不能吃了。如果
第一餐不是柔软的食物，而是石子一般硬得可怕的东西，这可怜的幼虫就会因为没有东
西吃，而饿死了。在八月的时候，我就找到了许多这样的牺牲者，这苦东西烤在一个封
闭的炉内，要减少这种危险，母甲虫就拼命用它强健而肥胖的前臂，压那梨子的外层，
把它压成保护的硬皮，如同栗子的硬壳，用以抵抗外面的热度。在酷热的暑天，管家婆
会把面包摆在闭紧的锅里，保持它的新鲜。而昆虫也有自己的方法，实现同样的目的：
用压力打成锅子的样子来保藏家族的面包。
我曾经观察过甲虫在巢里工作，所以知道它是怎样做梨子的。
它收集建筑用的材料，把自己关闭在地下，可以专心从事当前的任务，这材料大概
是由两种方法得来的。照常例，在天然环境下，甲虫用常法搓成一个球推向适应的地点。
当推行的时候，表面已稍微有些坚硬，并且粘上了一些泥土和细沙，这在后来是很多见
的，不只在离收集材料很近的地方，可以寻找到用来储藏的场所，在这种情况下，它的
工作不过是捆扎材料，运进洞而已。后来的工作，却尤其显得稀奇。有一天，我见它把
一块不成形的材料隐藏到地穴中去了。第二天，我到达它的工作场地时，发现这位艺术
家正在工作，那块不成形的材料已成功的变成了一个梨，外形已经完全具备，而且是很
精致的做好了。
梨紧贴着地板的部分，已经敷上了细沙。其余的部分，也已磨光得像玻璃一样，这
表明它还没有把梨子细细的滚过，不过是塑成形状罢了。
它塑造这只梨时，用大足轻轻敲击，如同先前在日光下塑造圆球一样。
在我自己的工作室里，用大口玻璃瓶装满泥土，为母甲虫做成人工的地穴，并留下
一个小孔以便观察它的动作，因此它工作的各项程序我都可以看得见。
甲虫开始是做一个完整的球，然后环绕着梨做成一道圆环，加上压力，直至圆环成
为一条深沟，做成一个瓶颈似的样子。这样，球的一端就做出了一个凸起。在凸起的中
央，再加压力，做成一个火山口，即凹穴，边缘是很厚的，凹穴渐深，边缘也渐薄，最
后形成一个袋。它把袋的内部磨光，把卵产在当中，包袋的口上，即梨的尾端，再用一
束纤维塞住。
用这样粗糙的塞子封口是有理由的，别的部分甲虫都用腿重重的拍过，只有这里不
拍。因为卵的层端朝着封口，假如塞子重压深入，蛴螬就会感到痛苦。所以甲虫把口塞
住，却不把塞子撞下去。

三、甲虫的生长
甲虫在梨里面产卵约一个星期或十天之后，卵就孵化成蛴螬了，它毫不迟疑地开始
吃四周的墙壁，它聪明异常，因为它总是朝厚的方向去吃，不致把梨弄出小孔，使自己
从空隙里掉出来。不久它就变得很肥胖了，不过样子实在很难看，背上隆起，皮肤透明，
假如你拿它来朝着光亮看，能看见它的内部器官。如果是古代埃及人有机会看见这肥白
的蛴螬，在这种发育的状态之下，他们是不会猜想到将来甲虫会具有的那些庄严和美观
了。
当第一次脱皮时，这个小昆虫还未长成完全的甲虫，虽然全部甲虫的形状，已经能
辨别出来了。很少有昆虫能比这个小动物更美丽，翼盘在中央，像折叠的宽阔领带，前
臂位于头部之下。半透明的黄色如蜜的色彩，看来真如琥珀雕成的一般。它差不多有四
个星期保持这个状态，到后来，重新再脱掉一层皮。
这时候它的颜色是红白色，在变成檀木的黑色之前，它是要换好几回衣服的，颜色
渐黑，硬度渐强，直到披上角质的甲胄，才是完全长成的甲虫。
这些时候，它是在地底下梨形的巢穴里居住着的。它很渴望冲开硬壳的甲巢，跑到
日光里来。但它能否成功，是要依靠环境而定的。
它准备出来的时期，通常是在八月份。八月的天气，照例是一年之中最干燥而且最
炎热的。所以，如果没有雨水来软一软泥土，要想冲开硬壳，打破墙壁，仅凭这只昆虫
的力量，是办不到的，它是没有法子打破这坚固的墙壁的。因为最柔软的材料，也会变
成一种不能通过的坚壁，烧在夏天的火炉里，早已成为硬砖头了。
当然，我也曾做过这种试验，将干硬壳放在一个盒子里，保持其干燥，或早或迟，
听见盒子里有一种尖锐的摩擦声，这是囚徒用它们头上和前足的耙在那里刮墙壁，过了
两三天，似乎并没有什么进展。
于是我加入一些助力给它们中的一对，用小刀戳开一个墙眼，但这两个小动物也并
没有比其余的更有进步。
不到两星期，所有的壳内都沉寂了。这些用尽力量的囚徒，已经死了。
于是我又拿了一些同从前一样硬的壳，用湿布裹起来，放在瓶里，用木塞塞好，等
湿气浸透，才将里面的潮布拿开，重新放到瓶子里。这次试验完全成功，壳被潮湿浸软
后，遂被囚徒冲破。它勇敢地用腿支持身体，把背部当作一条杠杆，认准一点顶和撞，
最后，墙壁破裂成碎片。在每次试验中，甲虫都能从中解放出来。
在天然环境下，这些壳在地下的时候，情形也是一样的。当土壤被八月的太阳烤干，
硬得像砖头一样，这些昆虫要逃出牢狱，就不可能了。但偶尔下过一阵雨，硬壳回复从
前的松软，它们再用腿挣扎，用背推撞，这样就能得到自由。
刚出来的时候，它并不关心食物。这时它所最需要的，是享受日光。跑到太阳里，
一动不动地取暖。
一会儿，它就要吃了。没有人教它，它也会做，像它的前辈一样，去做一个食物的
球，也去掘一个储藏所，储藏食物，一点不用学习，它就完全会从事它的工作。
------------------








蝉



一、蝉和蚁
我们大多数人对于蝉的歌声，总是不大熟悉的，因为它是住在生有洋橄榄树的地方，
但是凡读过拉封敦的寓言的人，大概都记得蝉曾受过蚂蚁的嘲笑吧。虽然拉封敦并不是
谈到这个故事的第一人。
故事上说：整个夏天，蝉不做一点事情，只是终日唱歌，而蚂蚁则忙于储藏食物。
冬天来了，蝉为饥饿所驱，只有跑到它的邻居那里借一些粮食。结果他遭到了难堪的待
遇。
骄傲的蚂蚁问道：“你夏天为什么不收集一点儿食物呢？”蝉回答道：“夏天我歌
唱太忙了。”
“你唱歌吗？”蚂蚁不客气地回答：“好啊，那么你现在可以跳舞了”，然后它就
转身不理它了。
但在这个寓言中的昆虫，并不一定就是蝉，拉封敦所想的恐怕是螽斯，而英国常常
把螽斯译为蝉。
就是在我们村庄里，也没有一个农夫，会如此没常识地想象冬天会有蝉的存在。差
不多每个耕地的人，都熟悉这种昆虫的蛴螬，天气渐冷的时候，他们堆起洋橄榄树根的
泥土，随时可以掘出这些蛴螬。至少有十次以上，他见过这种蛴螬从土穴中爬出，紧紧
握住树枝，背上裂开，脱去它的皮，变成一只蝉。
这个寓言是造谣，蝉并不是乞丐，虽然它需要邻居们很多的照应。每到夏天，它成
阵地来到我的门外唱歌，在两棵高大筱悬木的绿荫中，从日出到日落，那粗鲁的乐声吵
得我头脑昏昏。这种振耳欲聋的合奏，这种无休无止的鼓噪，使人任何思想都想不出来
了。
有的时候，蝉与蚁也确实打一些交道，但是它们与前面寓言中所说的刚刚相反。蝉
并不靠别人生活。它从不到蚂蚁门前去求食，相反的倒是蚂蚁为饥饿所驱乞求哀恳这位
歌唱家。我不是说哀恳吗？这句话，还不确切，它是厚着脸皮去抢劫的。
七月时节，当我们这里的昆虫，为口渴所苦，失望地在已经枯萎的花上，跑来跑去
寻找饮料时，蝉则依然很舒服，不觉得痛苦。用它突出的嘴——一个精巧的吸管，尖利
如锥子，收藏在胸部——刺穿饮之不竭的圆桶。它坐在树的枝头，不停的唱歌，只要钻
通柔滑的树皮，里面有的是汁液，吸管插进桶孔，它就可饮个饱了。
如果稍许等一下，我们也许就可以看到它遭受到的意外的烦扰。因为邻近很多口渴
的昆虫，立刻发现了蝉的井里流出的浆汁，跑去舔食。这些昆虫大都是黄蜂、苍蝇、蛆
蜕、玫瑰虫等，而最多的却是蚂蚁。
身材小的想要到达这个井边，就偷偷从蝉的身底爬过，而主人却很大方地抬起身子，
让它们过去。大的昆虫，抢到一口，就赶紧跑开，走到邻近的枝头，当它再转回头来时，
胆子比从前变大来了，它忽然就成了强盗，想把蝉从井边赶走。
最坏的罪犯，要算蚂蚁了。我曾见过它们咬紧蝉的腿尖，拖住它的翅膀，爬上它的
后背，甚至有一次一个凶悍的强徒，竟当着我的面，抓住蝉的吸管，想把它拉掉。
最后，麻烦越来越多，无可奈何，这位歌唱家不得已抛开自己所做的井，悄然逃走
了。于是蚂蚁的目的达到，占有了这个井。不过这个井也干得很快，浆汁立刻被吃光了。
于是它再找机会去抢劫别的井，以图第二次的痛饮。
你看，真正的事实，不是与那个寓言相反吗？蚂蚁是顽强的乞丐，而勤苦的生产者
却是蝉呢！

二、蝉的地穴
我有很好的环境可以研究蝉的习惯，因为我是与它同住的。七月初，它就占据了靠
我屋子门前的那棵树。我是屋里的主人，门外就它是最高的统治者，不过它的统治无论
怎样总是不会让人觉得舒服。
蝉初次被发现是在夏至。在行人很多，有太阳光照着的道路上，有好些圆孔，与地
面相平，大小约如人的手指。在这些圆孔中，蝉的蛴螬从地底爬出来，在地面上变成完
全的蝉。它们喜欢特别干燥而阳光充沛的地方。因为蛴螬有一种有力的工具，能够刺透
焙过的泥土与沙石。
当我考察它们的储藏室时，我是用手斧来开掘的。
最使人注意的，就是这个约一寸口径的圆孔，四边一点尘埃都没有，也没有泥土堆
积在外面。大多数的掘地昆虫，例如金蜣，在它的窝巢外面总有一座土堆。蝉则不同，
是由于它们工作方法的不同。金蜣的工作是在洞口开始，所以把掘出来的废料堆积在地
面；但蝉蛴螬是从地底上来的。最后的工作，才是开辟门口的生路，因为当初并没有门，
所以它不是在门口堆积尘土的。
蝉的隧道大都是深达十五至十六寸，一直通行无阻，下面的部分较宽，但是在底端
却完全关闭起来。在做隧道时，泥土搬移到哪里去了呢？为什么墙壁不会崩裂下来呢？
谁都以为蝉是用了有爪的腿爬上爬下的，而这样却会将泥土弄塌了，把自己房子塞住。
其实，它的举措简直像矿工或是铁路工程师一样。矿工用支柱支持隧道，铁路工程
师利用砖墙使地道坚固。蝉的聪明同他们一样，它在隧道的墙上涂上水泥。这种粘液是
藏在它身子里的，用它来做灰泥，地穴常常建筑在含有汁液的植物须上的，它可以从这
些根须取得汁液。
能够很容易的在穴道内爬上爬下，对于它是很重要的，因为当它爬出去到日光下的
时候，它必须知道外面的气候如何。所以它要工作好几个星期，甚至一个月，才做成一
道坚固的墙壁，适宜于它上下爬行。在隧道的顶端，它留着手指厚的一层土，用以保护
并抵御外面空气的变化，直到最后的一霎那。只要有一些好天气的消息，它就爬上来，
利用顶上的薄盖，以便测知气候的状况。
假使它估计到外面有雨或风暴——当纤弱的蛴螬脱皮的时候，这是一件最重要的事
情——它就小心谨慎地溜到隧道底下。但是如果气候看来很温暖，它就用爪击碎天花板，
爬到地面上来了。
在他肿大的身体里面，有一种液汁，可以利用它避免穴里面的尘土。当它掘土的时
候，将液汁倒在泥土上，使它成为泥浆。于是墙壁就更加柔软了。蛴螬再用它肥重的身
体压上去，便把烂泥挤进干土的缝隙里。因此，当它在顶端出口处被发现时，身上常有
许多湿点。
蝉的蛴螬，初次出现在地面上时，常常在附近徘徊，寻找适当的地点脱掉身上的皮
——一棵小矮树，一丛百里香，一片野草叶，或者一枝灌木枝——找到后，它就爬上去，
用前足的爪紧紧地握住，丝毫不动。
于是它外层的皮开始由背上裂开，里面露出淡绿色的蝉。当时头先出来，接着是吸
管和前腿，最后是后腿与翅膀。此时，除掉身体的最后尖端，身体己完全蜕出了。
然后，它会表演一种奇怪的体操，身体腾起在空中，只有一点固着在旧皮上，翻转
身体，使头向下，花纹满布的翼，向外伸直，竭力张开。于是用一种差不多看不清的动
作，又尽力将身体翻上来，并且前爪钩住它的空皮，用这种运动，把身体的尖端从鞘中
脱出，全部的过程大约需要半个小时。
在短时期内，这个刚被释放的蝉，还不十分强壮。它那柔软的身体，在还没具有足
够的力气和漂亮的颜色以前，必须在日光和空气中好好地沐浴。它只用前爪挂在已脱下
的壳上，摇摆于微风中，依然很脆弱，依然是绿色的。直到棕色的色彩出现，才同平常
的蝉一样。假定它在早晨九点钟取得树枝，大概在十二点半，弃下它的皮飞去。那壳有
时挂在枝上有一两月之久。

三、蝉的音乐
蝉是非常喜欢唱歌的。它翼后的空腔里带有一种像钹一样的乐器。它还不满足，还
要在胸部安置一种响板，以增加声音的强度。的确，有种蝉，为了满足音乐的嗜好，牺
牲了很多。因为有这种巨大的响板，使得生命器官都无处安置，只得把它们压紧到身体
最小的角落里。当然了，要热心委身于音乐，那么只有缩小内部的器官，来安置乐器了。
但是不幸得很，它这样喜欢的音乐，对于别人，却完全不能引起兴趣。就是我也还
没有发现它唱歌的目的。通常的猜想以为它是在叫喊同伴，然而事实明显，这个意见是
错误的。
蝉与我比邻相守，到现在已有十五年了，每个夏天差不多有两个月之久，它们总不
离我的视线，而歌声也不离我的耳畔。我通常都看见它们在筱悬木的柔枝上，排成一列，
歌唱者和它的伴侣比肩而坐。吸管插到树皮里，动也不动地狂饮，夕阳西下，它们就沿
着树枝用慢而且稳的脚步，寻找温暖的地方。无论在饮水或行动时，它们从未停止过歌
唱。
所以这样看起来，它们并不是叫喊同伴，你想想看，如果你的同伴在你面前，你大
概不会费掉整月的功夫叫喊他们吧！
其实，照我想，便是蝉自己也听不见所唱的歌曲。不过是想用这种强硬的方法，强
迫他人去听而已。
它有非常清晰的视觉。它的五只眼睛，会告诉它左右以及上方有什么事情发生，只
要看到有谁跑来，它会立刻停止歌唱，悄然飞去。然而喧哗却不足以惊扰它。你尽管站
在它的背后讲话，吹哨子、拍手、撞石子。就是比这种声音更轻微，要是一只雀儿，虽
然没有看见你，应当早已惊慌得飞走了。这镇静的蝉却仍然继续发声，好像没事儿人一
样。
有一回，我借来两枝乡下人办喜事用的土铳，里面装满火药，就是最重要的喜庆事
也只要用这么多。我将它放在门外的筱悬木树下。我们很小心的把窗打开，以防玻璃被
震破。在头顶树枝上的蝉，看不见下面在于什么。
我们六个人等在下面，热心倾听头顶上的乐队会受到什么影响。“碰！”枪放出去，
声如霹雷。
一点没有受到影响，它仍然继续歌唱。它既没有表现出一点儿惊慌扰乱之状，声音
的质与量也没有一点轻微的改变。第二枪和第一枪一样，也没有发生影响。
我想，经过这次试验，我们可以确定，蝉是听不见的，好像一个极聋的聋子，它对
自己所发的声音是一点也感觉不到的！

四、蝉的卵
普通的蝉喜欢把卵产在干的细枝上，它选择最小的枝，粗细大都在枯草与铅笔之间。
这些小枝干，垂下的很少，常常向上翘起，并且差不多已经枯死了。
蝉找到适当的细树枝，即用胸部尖利的工具，把它刺上一排小孔——这样的孔好像
是用针斜刺下去的，把纤维撕裂，使其微微挑起。如果它不被打扰与损害，在一根枯枝
上，常常被刺成三十或四十个孔。
它的卵就产在这些小孔里，这些小穴是一种狭窄的小径，一个个的斜下去。每个小
穴内，普通的约有十个卵，所以总数约有三百或四百个。
这是一个蝉的很好的家族。然而它所以产这么多卵，其理由是为防御一种特别的危
险，必须要生产出大量的蛴螬，预备将会被毁坏掉一部分。经过多次的观察，我才知道
这种危险是什么。就是一种极小的蚋，拿它们的大小相比较，蝉简直是庞然大物呢！
蚋和蝉一样，也有穿刺工具，位于身体下面靠近中部的地方，伸出来时和身体成直
角。蝉卵刚产出，蚋立刻就会把它毁坏。这真是蝉的家族中的灾祸！大怪物只须一踏，
就可轧扁它们，然而它们竟镇静异常，毫无顾忌，置身于大怪物之前，真令人惊讶之至。
我曾见过三个蚋顺序地排列着，同时预备掠夺一个倒楣的蝉。
蝉刚装满一个小穴的卵，移到稍高处，另外做穴时，蚋立刻就会到那里去，虽然蝉
的爪可以够得着它，然而它却镇静而无恐，像在自己家里一样，它们在蝉卵之上，加刺
一个孔，将自己的卵产进去。蝉飞回去时，它的孔穴内，多数已加进了别人的卵，这些
冒充的家伙能把蝉的卵毁坏掉。这种成熟得很快的蛴螬——每个小穴内一个——即以蝉
卵为食，代替了蝉的家族。
几世纪的经验，这可怜的蝉的母亲仍一无所知。它的大而锐利的眼睛，并非看不见
这些可怕的恶人，鼓翼其旁。它当然知道有其它昆虫跟在后面，然而它仍然不为所动，
宁肯让自己做牺牲。它要轧碎这些坏种子是非常容易的，不过它竟不改变原来的本能，
解救它的家族，以免遭破坏。
从放大镜里，我曾见过蝉卵的孵化过程。开始很像极小的鱼，眼睛大而黑，身体下
面有一种鳍状物。由两个前腿连在一起组成。这种鳍有些运动力，可以帮助蛴螬冲出壳
外。并且帮它走出有纤维的树枝，而这恰恰是比较困难的事情。
鱼形蛴螬到穴外后，立刻把皮脱去。但脱下的皮会形成一种线，蛴螬依靠它附着在
树枝上。它在未落地以前，就在这里进行日光浴，用腿踢着，试试它的精力，有时则又
懒洋洋地在绳端摇摆。
等到触须自由了，可以左右挥动，腿可以伸缩，在前面的能够张合其爪，身体悬挂
着，只要有一点微风，就摇摆不定，在空气中翻跟斗。我所看到的昆虫中再没有比这个
更为奇观的了。
不久，它就落到地面上来。这个像跳蚤一般大小的小动物，在它的绳索上摇荡，以
防在硬地面上摔伤。身体渐渐地在空气中变硬。现在它开始该投入到严肃的实际生活中
去了。
此时，它当前仍有着千重危险。只要有一点儿风，就能把它吹到硬的岩石上，或车
辙的污水中，或不毛的黄沙上，或粘土上，硬得它不能钻下去。
这个弱小的动物，有如此迫切的需要藏身，所以必须立刻钻到地底下寻觅藏身之所。
天气是冷起来了，迟缓一些就有死亡的危险。它不得不四处寻找软土，毫无疑问，它们
之中有许多在没有找到合适的地方之前就死去了。
最后，它寻找到适当的地点，用前足的钩爬挖掘地面。从放大镜中，我看见它挥动
斧头向下掘，并将士抛出地面。几分钟后，土穴完成，这个小生物钻下去，埋藏了自己，
此后就再也看不见了。
未长成的蝉的地下生活，至今还是未发现的秘密，我们所知道的，只是它未成长爬
到地面上来以前，地下生活经过了许多时间而已，它的地下生活大概是四年。此后，日
光中的歌唱不到五个星期。
四年黑暗的苦工，一月日光中的享乐，这就是蝉的生活，我们不应厌恶它歌声中的
烦吵浮夸。因为它掘土四年，现在忽然穿起漂亮的衣服，长起与飞鸟可以匹敌的翅膀，
在温暖的日光中沐浴着。那种钹的声音能高到足以歌颂它的快乐，如此难得，而又如此
短暂。
------------------








泥水匠蜂



一、选择造屋的地点
有很多种昆虫都非常喜欢在我们的屋子旁边建筑它们的巢穴，在这些昆虫中最能够
引起人们兴趣的，要首推那种叫舍腰蜂的动物了。为什么呢？主要原因在于，舍腰蜂有
着十分美丽而动人的身材，非常聪明的头脑，还有一点应该注意的就是它那种非常奇怪
的窠巢。但是，知道舍腰蜂这种小昆虫的人却是很少的。甚至有的时候，它们住在某一
家人的火炉的旁边，但是，这户人家都对这个小邻居竟然一无所知。为什么呢？主要是
由于它那种天生下来就具备了的，安静，而且平和的本性。的确，这个小东西居住得十
分隐避，很难引起人们的注意。因此，连它自己的主人都不知道它就住在自己的家里，
算得上是自家成员之一。然而，讨厌吵闹，而且特别怕麻烦的人类，和这些隐避性很强
的小动物相比，要想使它出名，倒是件很容易就能达到的事情。现在，就让我来把这个
谦逊的、默默无闻的小动物，从不知名中提拔出来吧！
舍腰蜂是一种非常怕冷的动物。它搭建起自己的帐蓬，在那帮助橄榄树茁壮成长，
鼓励着蝉儿纵悦高歌的太阳光下建筑自己的安乐之居。甚至有的时候，为了它们整个家
族的需要，为了让大家都觉得比在阳光下更加温暖舒适一些，它们常常找到我们人类的
门上，要求和我们一起作伴。不用敲开人们的大门，询问一下主人是否同意它们和大家
同住在一个屋檐下，便自作主张，举家迁移进来，并且定居下来享受生活。舍腰蜂平常
的居所，主要是一些农夫们的单独的茅舍。在那茅屋的门外，大部分都生长着一些高大
挺拔的无花果树。这些果树的树荫遮盖着一口小小的水井。舍腰蜂在具体确定它的住所
的时候，主要会选择一个能够暴露在夏日里的炎热之下的地点，并且，如果有可能的话，
最好能够有一只大一点儿的火炉，还要有一些能够燃烧使用的柴火，这些条件对于舍腰
蜂而言都是必要的，不可缺少的。这是由它的天性所决定的。到了寒冷的冬天的夜晚，
火炉中喷射出来的温暖无比的火焰，对于它的选择，有着十分重要的影响力。因此，每
当看到从烟筒里面出来的黑炭，舍腰蜂就会欣喜若狂，因为它们知道那里便是一个可以
考虑选择的地方。因为，那里将会提供给它所必需的温暖与安逸。但是，相反的，要是
烟筒里面并没有什么黑炭的话，那么它是绝对不会信任这种地方的，也绝对不会选择这
样的地方来建筑自己的家。因为舍腰蜂会利用它的头脑作出判断，这问屋子里的主人们
一定是在里面忍受着饥寒交迫的悲惨境域。
在七、八月里的大暑天中，这位小客人，忽然出现了。它在找寻着适合它做窠的地
点。舍腰蜂一点儿也不为这间屋子里面的一切喧闹行为所惊动和扰乱。而住在屋子里的
人们也一点儿都注意不到它。他们互相都没有注意到，因此也就互无干扰了。舍腰蜂只
不过在有的时候，利用它那尖锐的目光，有的时候，又利用它那灵敏十足的触须，视察
一下已经变得乌黑的天花板、木缝、烟筒等。但是，特别受到它关注的是火炉的旁边。
这是它从不轻易放过的地方。甚至，它连烟筒内部都要仔仔细细地视察一遍。它可是一
种细致入微的小动物，一旦视察工作完毕，并且已经决定了建巢的地点以后，它们便立
即飞走了。然后，不久就会带着少量的泥土又飞回来，开始建筑它的房子的底层了。于
是，筑造家园的工作便正式破土动工了。
舍腰蜂所选择的地点各不相同，也是非常奇怪的一个特点。炉子内部的温度最适合
那些小蜂了，因此，舍腰蜂所中意的地点，至少得是烟筒内部的两侧，其高度大约是二
十寸或者差不多的地方。不过，尽管这个地点可以说是一个非常舒服的藏身之妙处，但
是，世上没有十分完善的东西，它也有相当不少的缺点。由于巢是建在烟筒的内部的，
那么自然便会有烟在里面。如果烟要是喷到蜂巢上面，那么，巢中的舍腰蜂就会被“污
染”了，会被弄成棕色的或者是黑色的，就好像烟筒里被熏过的砖石一样。假使火炉里
的火焰烧不到蜂巢，那还不是一件最要紧的事。最重要的事是小黄蜂有可能会被闷死在
粘土罐子里。不过，不用替它们担心，它们的母亲似乎早就已经知道这些事情了，因为
这位母亲总是把它自己的家族安排在烟筒的适当位置上。它们选定的位置非常宽大。在
那个地方，除了烟灰以外，其它的东西都是很难到达的。
虽然舍腰蜂样样都当心，时刻都仔细、谨慎。但是“智者千虑，必有一失”。它如
此地认真，但还是有一件很危险的事情在等待着它们。这件事有的时候会发生，那就是
当舍腰蜂正在建造它的房屋的时候，如果在这个关健时刻，有一阵蒸汽或者是烟幕的侵
扰，那么，它刚刚造成的一半的房子，便不得不半途而废。于是，它们要么停工一些时
候，要么就全日停工不干。特别是在这家的主人在煮、洗衣服的日子里，这种事情发生
的可能性最大，危险性也最大。一天从早到晚，大盆子里不停地滚沸着，炉灶里的烟灰、
大盆和木桶里面的大量蒸汽，一起混合成为浓厚的云雾。这给蜂巢造成了严重的威胁。
这个时候舍腰蜂就会面临着家毁人亡的危险。
我以前曾经听别人说过，河鸟在回巢的时候，总是要飞过水坝下的大瀑布。这一点
听起来会让人觉得河鸟已经算得上是一种相当有勇气、有胆量的小动物了。但是，与之
相比的舍腰蜂也毫不示弱，甚至，它的勇敢已经超过河鸟。它在回巢的时候，牙齿间总
是要含着一块用于建造它的巢穴的泥土。要想到达它的施工工地，它要从浓厚的烟灰的
云雾中穿越过去。但是，那层烟幕简直太厚重了，舍腰蜂冲进去以后，就完全都看不见
它那小小的身影了。虽然看不见它那小小的躯体，但是能够听见一阵不太规则的呜呜的
声音。这是什么声音呢？这不是别的什么声音，这是它在一边工作，一边低唱的歌声。
因此，我们可以断定，舍腰蜂肯定还呆在里面，而且它很快乐，高高兴兴地从事着它的
本职工作，不知劳苦地建筑着它自己的住所。看得出来，它对自己的劳动很满意，也很
乐意从事这项工作，在这层厚厚的云雾里，它很神秘地进行着它自己的工作。忽然，低
低的劳动之歌停止了。不一会儿它飞出来了，从那层充满神秘色彩的浓雾里飞出来了，
它安然无恙，什么伤也没有得。毕竟这是它的本能嘛！差不多每天它都要经过很多次这
种十分危险的事情，直到它把巢最终建好，把食物都储藏好，最后把自家的大门关上为
止。然后，它才休息一下。这个小东西为了自己的家园也真够不辞辛苦的了！
每一次，只有我一个人能够看到舍腰蜂在我的炉灶里不停地忙碌着，建造住所，储
备食物。这大概是因为我比较细心。记得我第一次看到它们的时候，是有一天我在煮、
洗衣服的时候。本来，那个时候，我是在爱维浓(Avignon)学院里教书的。那天，时间已
经将近两点钟了，几分钟之内，外面就会敲鼓催促我去给羊毛工人们做演讲了。就在这
个时候，忽然，我看见了一个非常奇怪而且轻灵的小昆虫。它从由木桶里升腾起来的蒸
汽中穿飞出来。这只小动物的身体很有意思，当中的部分非常的瘦小，但是后部却是非
常肥大的，而这两个部分之间，竟然是由一根长线连接起来的。多么奇妙的小东西啊！
这个小东西就是舍腰蜂，这是我第一次没有用观察的眼光来看它。于是，便有了第一印
象。
在初次相识之后，我对家里的这个小客人一直报有非常浓厚的兴趣。我非常热心地
希望能和这个小不点儿客人互相熟识，作一些交流。于是，我便嘱咐我的家人，在我不
在家的时候，不要去主动打扰它们，破坏它们的正常生活。瞧，我多么注意保护这个没
有受到邀请的不速之客呀！事情发展的良好态势已然胜过了我所希望的那样。当我回到
家里的时候，发现它一点儿也没有受到什么打扰，而且一个个都安然无恙。它仍然呆在
蒸汽的后面，努力地进行着它自己的工作，为自己的家而辛苦。由于我想要观察一下舍
腰蜂的建筑以及它的建筑才能，还有它的食物的性质，以及幼小的黄蜂的进化及其生长
过程等等，因此，我把炉灶中的火焰给弄灭了。我这样做的目的主要是减少烟灰的量。
差不多将近两小时，我非常仔细地注视着它。
但是，从这以后，不知道是什么原因，差不多将近四十年来，我的屋子里，再也没
有这样小的客人光临了，一点儿也见不到它们的踪影了。有关舍腰蜂的进一步的知识，
我还是从我的邻居家的炉灶旁边的蜂巢里得出来的。
通过细心观察我发现，在这个小小的动物身上，有一种十分孤癖的流浪的习性。这
一点使得它和其它大多数黄蜂，以及蜜蜂是不尽相同的。一般情况下，它总是选择好一
个地点，自己筑起一个显得特别孤独的巢穴。同时，在舍腰蜂自己养活自己的地方，是
很少能见到它自己家族的成员及亲属的。在距离我们城南不远的地方，经常可以看到这
种小动物。但是，这个小东西，宁愿挑选农民那充满烟灰的屋子里的炉灶来筑造自己的
小家，也不喜欢那些城镇居民的雪白的别墅里的炉灶。我所到过的任何地方所看到的舍
腰蜂，都没有像我们村里这么多的。与此同时，我们村里的屋子都很有特点。我们村上
的茅屋都有一定的倾斜性，而且茅屋都被日光晒成了黄色，这使得它们看上去都很有特
色。
事实是很明显的，泥水匠蜂选择烟筒做为自己的住所。这一点是不用置疑的了。但
是，它之所以为自己选择这样一个地方，倒并不是意味着它贪图安逸与享乐。因为，很
显然，选择这样的地方可不是什么特舒服的地方。这种地方更需要这种小动物加倍地努
力，能够具备更好的才能。而且，在这种地方工作，是有很大危险性的。因为时常有险
情发生，需要冒一定的危险，甚至是生命的危险。从这一点来看，说它选择烟筒建巢是
为了自己的安逸，那可真的要大大的冤枉了我们这位小客人了。它选择这样的地点来筑
巢建穴，主要意图还完全是为了它的整个家族来考虑的，而并非出于私利。它不希望只
是自己舒服就可以了，应该是大家共同享福，共同舒适，那才是它们真正要达到的目的。
因而可以说，舍腰蜂还是一种比较热爱家庭的动物，家庭责任感很强。当然了，舍腰蜂
选择烟筒还有一个很重要的原因，那就是舍腰蜂及它的家族成员对温度的要求比较高，
这主要是由于本能的原因，它们的住所必须建在很温暖的地方，这一点和其它的黄蜂、
蜜蜂是很不相同的。
我记得有一次去一家丝厂，在那里我见到过一个舍腰蜂的巢。它把自己的巢建在机
房里，并且为自己选择了刚好是在大锅炉的上面的天花板上的一个地方。看来，它真是
慧眼独具啊！它为自己选择的这个地点，整个一年，无论寒暑，也无论春夏秋冬的变迁，
温度计所显示出的温度，总是不变的一百二十度，只是要除去晚上的时间，还有那些放
假的日子。很显然，在这些日子里，锅炉里并没有加热，所以，温度当然会随之有所变
化的。这个事实很明显地告诉我们，这个小小的动物对温度真是要求很高啊！而且，它
也是个非常会为自己挑选地点的家伙。
还有，在乡下的那些蒸酒的屋子里，我也曾经不止一次地看到过许多舍腰蜂的巢穴。
而且，凡是那些可以选择的、方便它们安居与行动的地方，都已经被它们占满了。甚至，
连那些帐薄堆积的地方，都被它们占据下来了。蒸酒房里的温度，和刚才提到的丝厂里
的温度相差得不是太多。大约有一百一十三度左右。这些温度字再次告诉我们，这种泥
水匠蜂甚至足可以在那种使油棕树生长的热度下生存。
这样看来，锅，还有炉灶，当然也就很自然地成了舍腰蜂最理想的家和首选对象了。
但是，除了这些首屈一指的地方以外，舍腰蜂也不厌弃一些其它可以选择的地点。它非
常希望居住在任何可以让它觉得舒适、安逸的角落里面。比如说，在养花房里，在厨房
的天花板上，可关闭窗户的凹进去的地方，还有就是茅舍中卧室的墙上等等。至于建造
自己窠巢的地基，这一点，它并不放在心上。为什么呢？因为，在平常它的多孔的巢穴，
一般都是建筑在石壁或者是木头上的。这些地方相对而言，还是比较坚实的。因而，它
们似乎并不是很关心房屋的基础。不过，也有的时候我曾经看到过它把自己的巢筑在葫
芦的内部，或者在皮帽子里，砖的缝隙之中，或者是装麦子用的空袋子里，还有的时候，
它建巢在铅管里面。
记得有一次，我在接近学院的一个农夫的家里所看到的事情，更让人觉得特别的新
奇。在这个农夫的家里，有一个特别宽大的炉灶。在装有这个炉灶的大房间里，在炉灶
上的一排锅里，正煮着农工们要喝的汤，还有一些供牲畜们食用的东西。过了一会儿，
工人们都从田地里收工回家了。累了一天，他们的肚子肯定饿坏了。回来后，他们便迫
不及待地，不声不响地，在一边非常迅速地吞食着他们的食品及汤。他们为了要享受这
休工用饭的大约半小时的舒适，干脆摘下了戴在头上妨碍吃饭的帽子，随后也脱去了他
们的上衣，随手把它们挂在一个木钉的上面。这吃饭的时间，对于农工们而言，虽然是
短暂的，但是，要是让泥水匠蜂去占据工人们刚刚脱下的衣物，却又是绰绰有余的了。
在这些衣物中、草帽里边，被它们视为最合适的地方，它们抢先去占领它。那些上衣的
褶缝，则被视为最佳的地点。与此同时，泥水匠蜂的建筑工作也就马上破土动工。这时，
一个工人已经吃完了他的饭，从饭桌旁边站了起来，抖了抖他自己的衣服。另外一个人
也站起来，走了过来，摘下自己的草帽，也抖了一下。这样几下抖动便去掉了舍腰蜂刚
刚初具规模的窠巢，就是在这个时候，在这么短暂的时间里，它的蜂巢居然已经有一个
橡树果子那样大了，真让人始料不及。它们可真是一些让人惊奇的小动物。
那个农夫家，有一位专门烹调食物的女人。她对于泥水匠蜂这种动物可是一点儿好
感也没有。她报怨说这些可恶的小东西常常跑出来，弄脏了许多的东西。天花板、墙壁，
还有烟筒上，经常被涂满了泥，非常烦人，打扫起来很费力气的。但是，在衣服和窗幔
上，情况就大不相同了。这个女人每天都会用一根竹子，使劲地敲打窗幔，以保持它的
清洁。所以，在这些地方情况会稍好一些，略微干净一些。但是，驱逐这些扰人的小动
物是多么地不容易啊！赶走了一次，第二天早晨它又会一样地跑回来来做巢。它可真是
个执着的小家伙，总是不厌其烦地从事着它的本能工作。

二、它的建筑物
事实上，我也非常同情这个农家厨役，很能理解她的烦恼。但是，我同时感到遗憾
的是，我不能代替她的位置。对此，我无能为力。如果，我能够以某种力量，使得这种
小动物安安静静地固定地在某一稳定的地点建屋居住，那该有多好啊，我肯定会特别高
兴的。这样一来既便它把家俱弄满了泥土，那也是不碍事的！我更希望能够知道它的那
种巢的命运。如果这个巢是做在不太稳固的东西上，比如，在衣服上，或是在窗幔上，
那么它们该怎么办呢？
泥水匠蜂的窠巢是利用硬的灰泥制做而成的。一般它的巢都围绕在树枝的四周。由
于是灰泥组成的，所以它就能够非常坚固地附着在上面。但是，泥水匠蜂的窠巢，只是
用泥土做成的，没有加水泥，或者是其它什么更能让它坚固的基础。那么，它怎么解决
这些问题呢？
建筑上的材料，并没有什么特殊的。只是潮湿的泥土，从那种湿地上取来的。因此，
河边的粘土是最合适的选择。但是，在我们这样一个多沙石的村庄里面，河道非常的少。
然而在我自己的小园子里，我在种植蔬菜的区域里，挖掘了一些小沟渠，以便更好地种
植。因此，有的时候，有一点儿水，便会整天在沟里流。因而，这里便经常会有舍腰蜂
的身影出没。它们在这里选择适宜的泥土，于是在无事可做的时候，我就可以观察这些
建筑家了。这里倒是一个很好的观察地点。
临近沟渠的时候，它当然就会注意到这件可喜的事情，于是就匆匆忙忙地跑过来取
水边这一点点十分宝贵的泥土。它们不肯轻意放过这没有湿气的时节极为珍稀的发现。
那么它们是怎样掘取这里的泥土的呢？它们用下颚刮取沟渠旁边那层表面光滑的泥，足
直立起来，双翼还振动着，把它那黑色的身体抬举得相当的高。我的管家妇在这泥土的
旁边做工。她把她的裙子非常小心谨慎地提起来，以免弄脏了。但是事实上，却很少能
够不沾上污渍。可是这样一群不停地搬取着泥土的黄蜂，原本应该是很脏的，但是事实
上它们的身上竟然连一点儿泥迹都没有。之所以会这样，它们自然有它们自己聪明的办
法。它们会把身子提起来，这样就能使它们全身上下一点儿泥污也沾染不上。除去它们
的足尖以及用于工作的下颚之外，其它的地方都看不到什么泥迹之类的脏东西。
这样，用不了多长时间，一个泥球就制做成功了。差不多能有豌豆那么大。然后，
泥水匠蜂会用牙齿把它衔住，飞回去，在它自己的建筑物上再增加上一层。这项工作完
成以后，它歇也不歇一下，便继续投入新的工作之中。接着飞回来，再做第二个泥球。
在一天中，天气最为炎热的时候，只要那片泥土未干，仍然是潮湿的，那么，泥水匠蜂
的工作就会不停地坚持下去。
除了我这园中的小小的沟渠边这片潮湿的泥土以外，在村子里，最好的地点，就算
是村里的人牵着驴子去饮水的那片泉水旁边了。在这个地方，无论什么时候都有潮湿的
黑色的烂泥。哪怕是那种最热的太阳，最强烈的风，都不可能把这片泥土吹干。这种泥
泞不堪的地方，对于那些走路的人来说，是非常不方便的，也是极不受欢迎的。然而，
舍腰蜂却不一样的。它非常喜欢到这个地方来，因为这里的泥土质量非常好，它也很喜
欢在驴子的蹄旁做小泥丸。每次它都会有丰富的收获。
和泥水匠蜂这位粘土建筑家不一样，黄蜂并不把泥土先做成水泥，它就这样把现成
的泥土拿走，直接应用于建筑。所以，黄蜂的巢建造得很不结实，更不稳定，完全不能
抵挡气候的千变万化。只要有一点儿水滴落上去，蜂巢就会变软，变成了和原来一样的
泥土。要是有一阵狂风大雨的话，它的巢穴就会被打成泥浆。这主要是因为，这种蜂巢
实际上只不过是由干了的烂泥做成的，一旦浸了水以后，就会马上变成和原来一样的软
泥，自然巢穴也就不复存在了，还须再次辛苦地重建家园。
事实是很显然的，即便是幼小的舍腰蜂一点儿也不惧怕寒冷，不怕雨水把蜂巢打得
粉碎，那蜂巢也必须建在避雨的地方。这就是为什么这种小动物喜欢选择人类居住的屋
子，特别是选择温暖的烟筒里面来建筑自己的住所的缘故。看来，安全是很重要的。
在最后一项装饰工作——那遮盖起它辛苦制造的建筑的各层——还没有完全成功之
前，舍腰蜂的窠巢确实具有一种非常自然的美感。它有一些小窠穴，有的时候它们互相
并列成一排，那种形状有一点儿像口琴。不过，那些小巢穴，还是以那种互相堆叠起来
成层的居多。有的时候，数一下有十五个小巢穴；有的时候，有十个；有时，又减少至
三、四个，甚至仅有一个。
舍腰蜂的巢穴的形状和一个圆筒子差不多。它的口稍微有点儿大，底部又稍小一些。
大的有一寸多长，半寸多宽，蜂巢有一个非常别致的表面，它是经过了非常仔细的粉饰
而形成的。在这个表面上，有一列线状的凸起围绕在它的四周，就好像金线带子上的线
一样。每一条线，就是建筑物上的一层。这些线的形状，是由于用泥土盖起每一层已经
造好的巢穴而显露出来的。数一数它们，就可以知道，在黄蜂建筑它的时候，来回旅行
了一共有多少次。它们通常是十五到二十层之间。每一个巢穴，这位辛辛苦苦的不辞辛
劳的建筑家在建筑它时，大概须用二十次来往返复搬运材料。可见，它们有多少勤劳！
蜂窠的口当然是朝着上面的。如果一个罐子的口是朝下的，那么，它还能盛下什么
东西呢？当然什么也盛不下了。道理也便就在这里。黄蜂的窠穴，也并不是什么特殊的
东西，不过就像一个罐子而已，其中预备盛储的食物便是：一堆小蜘蛛。
这些巢穴，一一建造好了以后，黄蜂便往里面塞满了蜘蛛。等它们自己产下卵以后，
便把它们全部封闭好。但是，这时候，它依然保存着美观的外表。这种外表一直要保持
到黄蜂认为巢穴的数量已经足够多了的时候为止。于是，黄蜂会把整个巢穴的四周，再
堆上一层泥土，以便使它能够更加坚固一些，从而可以起到保护的作用。这一次黄蜂在
工作时，也不进行什么周密的计算了。因此，它做得特别不精巧，更不像从前做巢那样，
铺加以相当的修饰之物。黄蜂能带回多少泥土，就往上面堆积多少泥土。只要能够堆积
得上去就可以了，再没有更多的修补、装璜的动作了。泥土一旦取了回来，便堆放到原
来的巢穴上去。然后，就那么非常漫不经心地轻轻地敲几下，使这些泥土可以铺开。这
一层包裹物质，一下子把建筑物的美观统统都掩盖住了。这最后一道工序完成以后，蜂
巢的最后形状就形成了。此时此刻的蜂巢就好像是一堆泥，一堆人们抛掷到墙壁上的泥。

三、它的食物
现在，我们都已经很清楚这个装食物的罐子是怎样形成的了。接下来，我们必须知
道的是，在这个罐子里边，究竟都隐藏了一些什么东西。
幼小的舍腰蜂，是以各种各样的蜘蛛作为食物的。甚至，在同一窠巢中，其食品的
形状各个都不相同，因为，各种各样的蜘蛛，都可以充当食品，只是个头不要过大。否
则就装不到罐子里去了。在幼蜂的各种食品中，那种后背上有三个交叉着的白点的十字
蜘蛛，是最为常见的美味佳肴。这其中的理由，我觉得应该是很简单的。因为，黄蜂不
是那种跑到离家很远的地方去千里迢迢地捕猎食物的动物。它只不过经常在住所的附近
地区游猎而已。而在它的住宅的近区内，这种有交叉纹的蜘蛛是最容易寻找得到的。
对于幼蜂而言，那种生长着毒爪的蜘蛛，要算是最最危险的野味儿了。假使蜘蛛的
身体特别的大，就需要黄蜂拥有更大的勇气和更多的技艺，才能够征服它。这可不是一
件容易的事情！而且，蜂巢的地方太小，也盛不下这么大的一个东西。所以，黄蜂只得
放弃猎取大个儿的蜘蛛，不去干这种费时、费力、又不讨好的傻事。还是更实际些吧。
于是，它只得选择去猎取那些较小一些的蜘蛛为食。如果，它偶然会碰上一群可以猎食
的蜘蛛，那么它总是很聪明，从来也不贪多，只选择其中最小的那一个。但是，虽然个
儿头都是较小，但它的俘虏的身材还是差别比较大。因此，大小的不同，就会影响到数
目的不同。在这个巢穴里面，盛有一打蜘蛛，而在另外一个巢穴里面，只藏着五个或者
六个蜘蛛。
黄蜂专选那些个儿小的蜘蛛，还有一个理由，那就是，在它还没有把猎物装入它的
巢穴里之前，它先得把那个蜘蛛杀死。它所要采取的行动，有以下几步：它先是突然一
下子落到蜘蛛的身上，以快取胜，差不多连翅膀都还没来得及停下来，就要把这个小蜘
蛛带走。其它的昆虫所采用的什么麻醉的方法等等，这个小动物可是一点儿也不知道。
这个小小的食物，一旦被储藏起来，就很容易变坏。幸好这个蜘蛛的个子小，一顿就可
以把它全部吃掉。要是换了一只大一些的蜘蛛，一顿是不可能吃完的，只能分成几次吃。
这样的话，这个蜘蛛是一定要腐烂的。这样烂了的食品就会毒害窠巢里其它的幼虫，这
对整个家族是不利的。
我经常能够看到，黄蜂的卵并不是放在蜂案的上面。而是在蜂案里面储藏着的第一
个蜘蛛的身上。差不多都是这样的，完全没有什么例外。黄蜂都是把第一个被捉到的蜘
蛛放在最下层，然后把卵放到它的上面的，再把别的蜘蛛放在顶上。用了这种聪明的办
法以后，小幼虫就能先吃掉那些比较陈旧的死蜘蛛，然后再吃那些比较新鲜的。这样一
来，蜂案里面储藏的食物也就没有什么时间足以变坏了。这不失为一种很安全的办法。
蜂的卵总是放在蜘蛛的身上的某一部分的。蜂卵的包含头的一端，放在靠近蜘蛛最
肥的地方。这对于幼虫是很好的。因为，一经孵化以后，幼虫就可以直接吃到最柔软、
最可口和最有营养的食物了。因此，这是一个很聪明的主意。应该说，大自然赋予了黄
蜂一种相当巧妙的天性。这样的一个有经济头脑的动物，一口食物也不浪费掉。等到它
完全吃光这个蜘蛛的时候，一堆蜘蛛什么也剩不下来了。这种大嚼的生活要经过八天到
十天之久。
在一顿美餐之后，蛴螬就开始做它的茧了。那是一种纯洁的白丝袋，异常而又精致。
还有一些东西，能够使这个幼虫的丝袋更加坚实。这些东西，可以用作保护之用。于是，
蛴螬就又从它身体里生出一种像漆一样的流质。这种流质慢慢地浸入丝的网眼里，然后
会渐渐地变硬，成为一种很光亮的保护漆。此时，幼虫又会在它正在做的茧下面增加一
个硬的填充物，使得一切都十分妥当。
这一项工作完成以后，这个茧呈现出琥珀的黄颜色，很容易让人联想到那种洋葱头
的外皮。因为，它和洋葱头有着同样细致的组织，同样的颜色，同样的透明感，而且，
它和洋葱头一样，如果用指头摸一摸，便会立即发出沙沙的响声，完整的昆虫就从这个
黄茧里孵化出来。早一点或是迟一点，这要随气候的变化而变化，各有不同。
当黄蜂在蜂巢中把东西储藏好以后，如果我们打算和它开一次玩笑的话，就立刻会
显露出黄蜂的本能是如何的机械了。
在它辛辛苦苦地把它自己的巢穴做好以后，便带回了它的第一个蜘蛛。黄蜂会马上
把它拖进巢里，然后收藏起来，立刻，又在它的身体的最肥大的部位产下一个卵。做好
了这一切以后，它便又飞了出去，继续它的第二次野外旅行和捕猎。当它不在家里的时
候，我从它的巢穴里，把那只死蜘蛛连同那个卵一起都取走了。就算和这只黄蜂开个小
小的玩笑吧。不知道它会有什么样的反应。
我们很自然地就会想到，如果这个小动物稍微有一点儿头脑的话，那么，这个蜘蛛
和卵的失踪，它是一定能够发觉得到的，而且应该会感到奇怪的！蜂卵虽然是小的，但
是，它是被放在那个大的蜘蛛的身体上的。那么，当我们的这个小东西回来以后，发现
巢穴里面是空的，它会怎么做呢？将有什么举动呢？它将很有理智地行动，再产下一个
卵，以补偿它所失去的那一个吗？事实上这些都不是，它的举动是非常不合情理的。
现在，这个小东西所做的事情，只不过是又带回了一只蜘蛛，非常坦然地再次把它
放到那巢穴里边去。对于其它的事情一律不理睬，就好像并没有发生过什么意外一样。
似乎它根本就没有看到自己的孩子已经丢失了。那只刚刚捕获的蜘蛛也已经丢了。它没
有发现这一切的不幸，也并没有表现出什么吃惊、诧异、着急、不知所措之类的失意的
表情。这以后，居然若无其事地一只又一只地盲目地往窠里继续传带着蜘蛛。每当它把
巢里的猎物和卵都安排妥当了以后，便又飞了出去，继续盲目地执着地奋斗着。每次在
它飞出去的时候，我都会把这些蜘蛛和蜂卵悄悄地拿出来。因此，它每一次游猎回来，
储藏室里实际上总是空着的。就这样，它十分固执而徒劳地忙碌了整整两天时间。它一
心打算要使劲努力，无论如何也要争取装满这个不知为什么永远也装不满的食物瓶子。
我呢，也和它一样，不屈不挠地坚持了有两天的工夫。一次又一次耐心地把巢穴里的蜘
蛛和卵取出来。想要看看这个执着的小傻瓜究竟要等何时才能终结它这种看起来毫无意
义的工作。当这个傻乎乎的小动物完成了它的第二十次任务的时候，也就是到了第二十
次的收获物送来的时候，这位辛苦多时的猎人大概以为这罐子已经装够了——或许也是
因为这么多次的旅行，疲倦了——于是，它便自认为非常小心而且谨慎地把自己的巢穴
封锁了起来，然而，实际上，里面却完全是空的！什么东西都没有。它忙碌了这么久，
事实上它根本意识不到这一点，真是让人可怜啊！
在任何情况下，昆虫的智慧都是非常有限的。这一点是毫无疑问的，无论是哪一种
临时的困难，昆虫，这种小小的的动物，都是无力加以很好而且迅速的解决的。无论是
哪一个种类的昆虫，都同样的不能对抗。这一点，我可以列举出一大堆的例子来，证明
昆虫是一种完全没有理解能力的动物。当然，同时，它也是一种不具有意识的动物。虽
然它们的工作是那么异常的完备。经过长时间的经验和观察，使我不能不断定它们的劳
动，既不是自动的，也不是有意识的。它们的建筑、纺织、打猎、杀害，以及麻醉它们
的捕获物，都和消化食物或是分泌毒汁一样，其方法和目的完全都是不自知的。所以，
我相信这样一点，即这些动物对于它们所具有的特殊的才能，完全是莫名其妙的，既不
知也不觉。
动物的本能是不能改变的。经验不能指导它们；时间也不能使它们的无意识有一丝
一毫的觉醒。如果它们只有单纯的本能，那么，它们便没有能力去应付大千世界，应付
大自然环境的变化。环境是要经常有所变化的，意外的事情有很多，也时常会发生，正
因为如此，昆虫需要具备一种特殊的能力，来教导它，从而让它们自己能够清楚什么是
应该接受的，什么又是应该拒绝的。它需要某种指导。这种指导，它当然是具备的，不
过，智慧这样一个名词，似乎太精细了一点，在这里是不适用的。于是，我预备叫它为
辨别力。
那么，昆虫，能够意识到它自己的行动吗？能，但同时也不能。如果它的行动是由
于它所拥有的本能而引起的，那么它就不能知道自己的行动。如果它的行动是由于辨别
力而产生的结果，那么，它就能意识到。
比如，舍腰蜂利用软土来建造巢穴，这一点就是它的本能。它常常是如此建造巢穴
的，从一生下来就会。既不是时间，也不是生活的奋斗与激励，能够使得它模仿泥水匠
蜂，用那种细沙的水泥去建造它自己的巢，这并不是它的本能。
黄蜂的这个泥巢，一定要建在一种隐避之处，以便抵抗自然风雨的侵袭。在最初的
时候，大概那种石头下面可以隐匿的地方就能够被认为是相当合适了。但是，当它发现
还有更好的其它的地方可以选择时，它便会立刻去占据下来，然后搬到人家的屋子里边
去住。那么，这一种就属于辨别力了。
黄蜂利用蜘蛛作为它的子女的食物，这就是它本能的一种。表现没有其它的任何方
法，能够让这只黄蜂明白，小蟋蟀也是一样的好，和蜘蛛一样可以当作食物。不过假设
那种长有交叉白点的蜘蛛缺少了，那它也不会让它的宝宝挨饿的。它会选择其它类型的
蜘蛛，将其捕捉回来，给它的子女吃。那么，这种就是辨别力。
在这种辨别力的性质之下，隐伏了昆虫将来进步的可能性。

四、它的来源
舍腰蜂又给我们带来了另外一个问题。它寻找着我们房子里的火炉的热量。这是因
为它的巢穴是用软土建筑起来的，潮湿会把它给弄成泥浆而无法居住。所以，基于上述
原因，它必须要有一个干燥的隐避场所。因此，热量，也是黄蜂生活中所必要的。
那么，它是不是一个侨民呢？或许它是从海边被卷过来的？它是从有枣椰树的陆地
来到生长洋橄榄的陆地的？如果真的是这样的话，那么它们也就自然地会觉得我们这个
地方的太阳不够温暖，也就必须要寻找一些地方，比如说火炉，作为人工取暖的地方了。
这样就可以解释它的习性了，为什么它能和别的种类的黄蜂有如此大的差别。而且这种
蜂都是避人的。
在它还没有到我们这里来做客以前，它的生活是什么样子的呢？在没有房屋以前，
它住在什么样的地方呢？没有烟筒的时候，它把蛴螬隐藏在哪里呢？
也许，当古代山上的居民用燧石做武器，剥掉羊皮做衣服，用树枝和泥土造屋子的
时候，这些屋子也已经早就有舍腰蜂的足迹了。也许，它们的巢就建筑在一个破盆里面，
那是我们的祖先用手指取粘土制做成的。或者它就在狼皮及熊皮做的衣服的褶缝里边筑
巢。我感到非常奇怪，当它们在用树枝和粘土造成的粗糙的壁上做巢的时候，它们是否
选择那些靠近烟筒的地点呢？这些烟筒，虽然和我们现在所使用的烟筒不同，但是，在
不得已的时候，那些烟筒也是可以利用的。
如果说，舍腰蜂在古代的时候，的确和那些最古老的人们共同在这个地方居住过，
那么它所经历和见到的进步，就真的是不少了，而且，它所得到的文明的利益也真正是
不少了。它已经把人类不断增进的幸福转变成为自己的了。当人类社会，发明出在房屋
的屋顶上铺上天花板的法子，想出在烟筒上加上管子的主意以后，我们便可以想象得到，
这个怕冷的动物就会悄悄地对自己自言自语道：
“这是如何的舒适啊！让我们在这里撑开帐篷吧！”
但是，我们还应该追究得更远一些。在小屋没有出现以前，在壁龛也不常见之前，
甚至是在人类还没有出现之前，舍腰蜂又是在哪里造房子的呢？这个问题，当然不是一
个孤立的问题。我们还可以提出这样的问题：燕子和麻雀在没有窗子、烟筒等东西以前，
它们又是在哪里筑巢的呢？
燕子、麻雀、舍腰蜂是在人类出现以前就已经存在了的动物。显然，它们的工作是
不能依靠人类劳作的。当这里还没有人类的时候，它们各自必定已经具有了高超的建筑
技艺了。
三四十年来，我都常常问我自己，在那个时候舍腰蜂住在哪里的问题。
在我们的屋子外面，我们找不到它们的窠巢的痕迹。在房子外边，在空旷的广场、
在荒丘的草地里，我们都没有找到舍腰蜂的住处。
但是，最后，我长时间的研究结果表明，一个帮助我的机会出现了。
在西往南的采石场上，有许多碎石头子和很多的废弃物，堆积在这里有很长时间了。
据说已经有几百年的时间了。在这个乱石堆上，沉积了几个世纪的污泥，几百年的风风
雨雨，将这些乱石堆摆在人们面前。田鼠也在那里生活着。在我寻找这些宝藏的时候，
我有三次发现了在乱石堆中的舍腰蜂的巢。
这三个巢与在我们屋子里发现的完完全全一个样，材料当然也是泥土，而用以保护
的外壳，也是相同的泥土。
这个地方的危险性，并没有促使这位建筑家丝毫的一点儿进步。我们有时——不过
很少——看到舍腰蜂的巢筑在石堆里和不靠着地的平滑的石头下面。
在它们还没有侵入我们的屋子以前，它们的窠巢一定是建筑在这类地方的。
然而，这三个巢的形状，是很凄惨的，湿气已经把它们给侵蚀坏了，茧子也被弄得
粉碎了。周围也没有厚厚的土保护着它，它们的幼虫也已经牺牲了——已经被田鼠或别
的动物吃光了。
这个荒废的景象，使我惊疑的来到我邻居的屋外。是否能够真为舍腰蜂建巢的地点，
挑选一个适当的位置呢？事实很显然，母蜂不愿意这么做，并且也不至于被驱逐到这么
绝望的地步。同时，如果气候使它不能从事它祖先的生活方式，那么，我想，我可以断
言，它就是一个侨民。它很可能是从遥远的异国他乡，侨居到这个地方来的侨民。也可
能是另一种移民，那种背井离乡的移民。也可能是难民，是为了生计，不得不远走他乡，
被其他地方收养的难民。
事实的确如此，它是从炎热的、干燥的、缺水的、沙漠式的地方来的，在它们那里，
雨水不多，雪简直是没有的。
我相信，舍腰蜂是从非洲来的。
很久以前，它经过了西班牙，又经过了意大利，来到了我们这里，它可以说的上是
千里迢迢，也可以说它是不远万里、不辞辛苦地到我们这里来。
它不会越过长着洋橄榄树的地带，再往北去。它的祖籍是非洲，而现在它又归入了
我们布曼温司。
在非洲，据说它常把巢穴建筑在石头的下面，而在马来群岛，听说也有它们的同族、
同宗，它们是住在屋子里的。
从世界的这一边，来到世界的那一边，从世界的南边来到世界的北边，从地球的南
边——非洲，来到地球的北边——欧洲！最后又来到马来群岛。它的嗜好都是一个样的：
蜘蛛、泥巢，还有人类的屋顶。
假如我是在马来群岛，我一定要翻开乱石堆，翻找它居住的巢穴。这时，我会很高
兴地在一块平滑的石头下面，发现它的巢穴，发现它的住所——原来它的位置，就在这
些石头的下面。
------------------








螳螂



一、打猎
在南方有一种昆虫，与蝉一样，很能引起人的兴趣，但不怎么出名，因为它不能唱
歌。如果它也有一种钹，它的声誉，应比有名的音乐家要大得多，因为它在形状上与习
惯上都十分的不平常。它将是一名出色的乐手。
多年以前，在古希腊时期，这种昆虫叫做螳螂，或先知者。农夫们看见它半身直起，
立在太阳灼烧的青草上，态度很庄严，宽阔的、轻纱般的薄翼，如面膜似的拖曳着，前
腿形状如臂，伸向半空，好像是在祈祷，在无知识的农夫看来，它好像是一个女尼，所
以后来，就有人称呼它为祈祷的螳螂了。
这个错误再大没有了！那种貌似真诚的态度是骗人的，高举着的似乎是在祈祷的手
臂，其实是最可怕的利刃，无论什么东西经过它的身边，它便立刻原形毕露，用它的凶
器加以捕杀。它真是凶猛如饿虎，残忍如妖魔，它是专食活的动物的。看来，在它温柔
的面纱下，隐藏着十分吓人的杀气。
如果单从外表上看来，它并不令人生畏，相反，看上去它相当美丽，它有纤细而优
雅的姿态，淡绿的体色，轻薄如纱的长翼。颈部是柔软的，头可以朝任何方向自由转动。
只有这种昆虫能向各个方向凝视，真可谓是眼观六路。它甚至还有一个面孔。这一切都
构成了这样一个小动物的温柔。
螳螂天生就有着一副娴美而且优雅的身材。不仅如此，它还拥有另外一种独特的东
西，那便是生长在它的前足上的那对极具杀伤力，并且极富进攻性的冲杀、防御的武器。
而它的这种身材和它这对武器之间的差异，简直是太大了，太明显了，真让人难以相信，
它是一种温存与残忍并存的小动物。
见过螳螂的人，都会十分清楚地发现，它的纤细的腰部非常的长。不光是很长，还
特别的有力呢。与它的长腰相比，螳螂的大腿要更长一些。而且，它的大腿下面还生长
着两排十分锋利的像锯齿一样的东西。在这两排尖利的锯齿的后面，还生长着一些大齿，
一共有三个。总之，螳螂的大腿简直就是两排刀口的锯齿。当螳螂想要把腿折叠起来的
时候，它就可以把两条腿分别收放在这两排锯齿的中间，这样是很安全的，不至于自己
伤到自己。
如果说螳螂的大腿像是两排刀口的锯齿的话，那么它的小腿可以说是两排刀口的锯
子。生长在小腿上的锯齿要比长在大腿上的多很多。而且，小腿上的锯齿和大腿上的有
一些不太相同的地方。小腿锯齿的末端还生长着尖而锐的很硬的钩子，这些小钩子就像
金针一样。除此以外，锯齿上还长着一把有着双面刃的刀，就好像那种成弯曲状的修理
各种花枝用的剪刀一样。
对于这些小硬钩，我有着许多不堪回首的记忆。每次想到它们，都有一种难受的感
觉。记得从前曾经有过许多次这样的经历。在我到野外去捕捉螳螂的时候，经常遭到这
个小动物的强有力的自我保护与还击，总是捉它不成，反过来倒中了这个小东西的十分
厉害的“暗器”，被它抓住了手。而且，它总是抓得很牢，不轻易松开，让我自己无法
从中解脱出来，只有想其他的方法，请求别的人前来相助，帮我摆脱它的纠缠。所以，
在我们这种地方，或许再也没有什么其它的昆虫比这种小小的螳螂更难以对付，更难以
捕捉的了。螳螂身上的武器、暗器很多，因此，它在遇到危险的时候，可以选择多种方
法来自我保护。比如，它有如针的硬钩，可以用镰钩去钩你的手指；它长有锯齿般的尖
刺，可以用它来扎、刺你的手；它还有一对锋利无比、而且十分健壮的大钳子。这对大
钳子对你的手有相当的威力，当它挟住你的手时，那滋味儿可不太好受啊！综上所述，
这种种有杀伤力的方法，让你很难对付它。要想活捉这个小动物，还真得动一番脑筋，
费一番周折呢！否则，捉住它将是不可能的。这个小东西不知要比人类小多少倍，但却
能威胁住人类。
平时，在它休息、不活动的时候，这个异常勇猛的捕捉其它昆虫的机器，只是将身
体蜷缩在胸坎处，看上去，似乎特别的平和，不至于有那么大的攻击性，甚至会让你觉
得，这个小动物简直是一只热爱祈祷的温和的小昆虫。但是，它可不总是这样的，否则
的话，它身上具备的那些进攻、防卫的武器也就派不上什么用场了。只要是有其它的昆
虫从它们的身边经过，无论是什么样的昆虫，也无论它们是无意路过，还是有意地侵袭，
螳螂的那副祈祷和平的相貌便会一下子烟消云散了。这个刚才还是蜷缩着休息的小动物，
立刻便伸展开它身体的三节，于是，那个可怜的路过者，还没有完全反应过来，便己糊
里糊涂地成了螳螂利钩之下的俘虏了。它被重压在螳螂的两排锯齿之间，移动不得。然
后，螳螂很有力地把钳子夹紧，一切战斗就都结束了。无论是蝗虫，还是蚱蜢，或者甚
至是其它更加强壮的昆虫，都无法逃脱这四排锋利的锯齿的宰割。于是，一旦被捉，只
好束手就擒了。它可真是个了不得的杀虫机器。
假如你想到原野里面去详尽地研究，观察螳螂的习性，那几乎是不可能的。因此，
也就不得不把螳螂拿到室内来进行观察、分析和研究。如果把螳螂放在一个用铜丝盖住
的盆里面，再往盆里加上一些沙子，那么，这只螳螂将会生活得十分快乐和满意。我所
要做的，只是提供给它充足而又新鲜的食物就可以了。有了它必须的食品，它会生活得
更满意。因为我想要做一些试验，测量一下螳螂的筋力究竟能够有多大，所以，我不仅
仅是提供一些活的蝗虫或者是活的蚱蜢给螳螂吃，同时，还必须供给它一些最大个儿的
蜘蛛，以使它的身体更加强壮。至于我的观察、研究，以下便是在我做了上述工作以后，
所观察到的情形。
有这样一只不知危险、无所畏惧的灰颜色的蝗虫，朝着那只螳螂迎面跳了过去。后
者，也就是那只螳螂，立刻表现出异常愤怒的态度，接着，反应十分迅速地做出了一种
让人感到特别诧异的姿势，使得那只本来什么也不怕的小蝗虫，此时此刻也充满了恐惧
感。螳螂表现出来的这种奇怪的面像，我敢肯定，你从来也没有见到过。螳螂把它的翅
膀极度地张开，它的翅竖了起来，并且直立得就好像船帆一样。翅膀竖在它的后背上，
螳螂将身体的上端弯曲起来，样子很像一根弯曲着手柄的拐杖，并且不时地上下起落着。
不光是动作奇特，与此同时，它还会发出一种声音。那声音特别像毒蛇喷吐气息时发出
的声响。螳螂把自己的整个身体全都放置在后足的上面。显然，它已经摆出了一副时刻
迎接挑战的姿态。因为，螳螂已经把身体的前半部完全都竖起来了，那对随时准备东挡
西杀的前臂也早已张了开来，露出了那种黑白相间的斑点。这样一种姿势，谁能说不是
随时备战的姿势呢？
螳螂在做出这种令谁都惊奇的姿势之后，一动不动，眼睛瞄准它的敌人，死死盯住
它的俘虏，准备随时上阵，迎接激烈的战斗。哪怕那只蝗虫轻轻地、稍微移动一点位置，
螳螂都会马上转动一下它的头，目光始终不离开蝗虫。螳螂这种死死的盯人战术，其目
的是很明显的，主要就是利用对方的惧怕心理，再继续把更大的惊恐纳入这个不久以后
就将成为牺牲者的对手心灵深处，造成“火上加油”的效果，给对手施加更重的压力。
螳螂希望在战斗未打响之前，就能让面前的敌人因恐惧心理而陷于不利地位，达到使其
不战自败的目的。因此，螳螂现在需要虚张声势一番，假装什么凶猛的怪物的架势，利
用心理战术，和面前的敌人进行周旋。螳螂真是个心理专家啊！
看起来，螳螂的这个精心安排设计的作战计划是完全成功的。那个开始天不怕、地
不怕的小蝗虫果然中了螳螂的妙计，真的是把它当成什么凶猛的怪物了。当蝗虫看到螳
螂的这副奇怪的样子以后，当时就有些吓呆了，紧紧地注视着面前的这个怪里怪气的家
伙，一动也不动，在没有弄清来者是谁之前，它是不敢轻易地向对方发起什么攻势的。
这样一来，一向擅于蹦来跳去的蝗虫，现在，竟然一下子不知所措了，甚至连马上跳起
来逃跑也想不起来了。已经慌了神儿的蝗虫，完全把“三十六计，走为上策”这一招儿
忘到脑后去了。可怜的小蝗虫害怕极了，怯生生地伏在原地，不敢发出半点声响。生怕
稍不留神，便会命丧黄泉，在它最害怕的时候，它甚至莫明其妙地向前移动，靠近了螳
螂。它居然如此地恐慌，到了自己要去送死的地步。看来螳螂的心理战术是完全成功了。
当那个可怜的蝗虫移动到螳螂刚好可以碰到它的时候，螳螂就毫不客气，一点儿也
不留情地立刻动用它的武器，用它那有力的“掌”重重的地击打那个可怜虫，再用那两
条锯子用力地把它压紧。于是，那个小俘虏无论怎样顽强抵抗，也无济于事了。接下来，
这个残暴的魔鬼胜利者便开始咀嚼它的战利品了。它肯定是会感到十分得意的。就这样，
像秋风扫落叶一样地对待敌人，是螳螂永不改变的信条。
在蜘蛛捕捉食物、降服敌人的时候，它通常采取的办法是：首先，一上来便先发制
人，猛烈地刺击敌人的颈部，让它中毒。这样做的好处是对手中了毒，自然也就没有了
力气，也就不能继续抵抗防卫了。先下手为强嘛！与此相同的，螳螂在攻击蝗虫的时候，
也是首先重重地、不留情面地击打对方的颈部。受了一顿狂轰乱炸的痛捶之后，再加上
先前万分的恐惧，蝗虫的运转能力逐渐下降，动作慢慢地迟缓下来。也许是已经被打蒙
了的原因吧。这种办法既有效又非常的实用。螳螂就是利用这种办法，屡屡取得战斗的
胜利。无论是杀伤并食用和它一样大小的动物，还是对付比自己还要大一些的昆虫，这
种办法都是十分有效的。不过，最让人感到奇怪的，就是这么一只小个儿的昆虫，竟然
是一种十分贪吃的动物，能吃掉这么多的食物。
那些爱掘地的黄蜂们，算得上是螳螂的美餐之一了，因此常常受到螳螂的光顾。螳
螂经常出没于黄蜂的地穴附近。因此，在黄蜂的窠巢近区看到螳螂的身影屡屡出现，便
不足为奇了。螳螂总是埋伏在蜂窠的周围，等待时机，特别是那种能获得双重报酬的好
机会。为什么说是双重报酬呢？原来，有的时候，螳螂等待的不仅仅是黄蜂本身，因为
黄蜂自己的身上常常也会携带一些属于它自己的俘虏。这样一来，对于螳螂而言，不就
是双份的俘虏，双重报酬了吗？不过，螳螂并不总是这么走运的，也有不太幸运的时候。
有时，它也会常常什么都等不到，竟无功而返。主要原因是，黄蜂已经有所疑虑，从而
有所戒备了，方让螳螂失望而归。但是，也有个别掉以轻心者虽已发觉但仍不当心的，
被螳螂看准时机，一举将其抓获。这些命运悲惨的黄蜂为什么会遭到螳螂的毒手呢？因
为，有一些刚从外面回家的黄蜂，它们振翅飞来，有一些粗心大意，对早已埋伏起来的
敌人毫无戒备。当突然发觉大敌当前时，会被猛地吓了一跳，心里会稍稍迟疑一下，飞
行速度忽然减慢下来。但是，就在这千钧一发的关键时刻，螳螂的行动简直是迅雷不及
掩耳。于是，黄蜂一瞬间便坠入那个两排锯齿的捕捉器中——即螳螂的前臂和上臂的锯
齿之中了。螳螂就是这样出其不备，以快致胜的。接下来，那个不幸的牺牲者就会被胜
利者一口一口地蚕食掉。又成了螳螂的一顿美餐。
记得有一次，我曾看见过这样有趣的一幕。有一只黄蜂，刚刚俘获了一只蜜蜂，并
把它带回到自己的储藏室里，正在享用这只蜜蜂体内的蜜汁。不料，正在它吃得高兴的
时候，遭到了一只凶悍的螳螂的突然袭击。它无力还击，便束手就擒了。这只黄蜂正在
吃蜜蜂的嗉袋里储藏的蜜，但是螳螂的双锯，在不经意中，竟然有力地夹在了它的身上。
可是，就是在这种被俘虏的关键时刻，无论怎样的惊吓、恐怖和痛苦，竟然不能让这只
贪吃的小动物停止继续吸食蜜蜂体内的蜜汁。它依然在甜食着那芬香诱人的蜜汁。这真
是太奇异了，真是人为财死，鸟为食亡啊！
螳螂，这样一种凶狠恶毒、有如魔鬼一般的小动物，它的食物的范围并不仅仅局限
于其它种类的所有昆虫。螳螂的气概虽然特别神圣，但是，或许你想不到，因为这实在
是让人不可思议。事实上，螳螂还是一种自食其同类的动物呢。也就是说，螳螂是会吃
螳螂的，吃掉自己的兄弟姐妹。而且，在它吃的时候，面不改色，心不跳，十分泰然自
若，那副样子，简直和它吃蝗虫，吃蚱蜢的时候一模一样，仿佛这是天经地义的事情。
并且，与此同时，围绕在食同类的螳螂旁边围观的观众们，也没有任何反应，没有任何
抵抗的行动。不仅如此，这些观众还纷纷跃跃欲试，时刻准备着，一旦有了机会，它们
也会做同样的事情，也同样地毫不在乎，仿佛顺理成章似的。然而在事实上，螳螂甚至
还具有食用它丈夫的习性。这可真让人吃惊！在吃它的丈夫的时候，雌性的螳螂会咬住
它丈夫的头颈，然后一口一口地吃下去。最后，剩余下来的只是它丈夫的两片薄薄的翅
膀而已。这真人难以置信。
螳螂真的是比狼还要狠毒十倍啊！听说，即便是狼，也不吃它们的同类。那么，螳
螂真的是很可怕的动物了！

二、它的巢
虽然我们的螳螂是如此的凶猛而又可怕，它身上有那么多的杀伤性很强的武器，还
有那么凶恶的捕食方法，甚至它居然要以自己的同类为食。尽管如此，螳螂也和人类是
一样的，不光有缺点和不足之处，还拥有很多自己的优点。比如，螳螂能够建造十分精
美的巢穴，这便是螳螂众多优点中很突出的一个。
螳螂建造的窠巢，在有太阳光照耀的地方随处都可以找得到。比如，石头堆里，木
头块下，树枝上，枯草丛里，一块砖头底下，一条破布下，或者是旧皮鞋的破皮子上面
等等。总之一句话，在任何东西上，只要那个东西上有凸凹不平的表面，都可以作为非
常坚固的地基。螳螂就是利用这样的地基建巢的。
螳螂的巢，大小约有一两寸长，不足一寸宽。巢的颜色是金黄色的，样子很像一粒
麦子。这种巢是由一种泡沫很多的物质做成的。但是，不久以后，这种多沫的物质就逐
渐变成固体了，而且慢慢地变硬了。如果燃烧一下这种物质，便会产生出一种像燃烧丝
质品一样的气味儿。螳螂巢的形状各不相同。这主要是因为巢所附着的地点不同，因而
巢随着地形的变化而变化，会有不同形状的巢存在。但是，不管巢的形状多么干变万化，
它的表面总是凸起的。这一点是不变的。
整个的螳螂巢，大概可以分成三个部分。其中的一部分是由一种小片做成的，并且
排列成双行，前后相互覆盖着，就好像屋顶上的瓦片一样。这种小片的边沿，有两行缺
口，是用来做门路的。在小螳螂孵化的时候，就是从这个地方跑出来的。至于其它部分
的墙壁，全都是不能穿过的。
螳螂的卵在巢穴里面堆积成好几层。其中每一层，卵的头都是向着门口的。前面我
已经提到过了，那道门有两行，分成左、右两边。所以，在这些幼虫中，有一半是从左
边的门出来的，其余的则从右边的门出来。
有这样一个事实是值得注意的，那就是母螳螂在建造这个十分精致的巢穴的时候，
也正是它产卵的时候。在这个时候，从母螳螂的身体里，会排泻出一种非常有粘性的物
质。这种物质和毛虫排泻出来的丝液很相象。这种物质在排泻出来以后，将与空气互相
混合在一起，于是就会变成泡沫。然后，母螳螂会用身体末端的小杓，把它打起泡沫来。
这种动作，特别像我们用叉子搅打鸡蛋蛋白一样。打起来的泡沫是灰白色的，与肥皂沫
十分相似。开始的时候，泡沫是有粘性的。但是过了几分钟以后，粘性的泡沫就变成了
固体。
母螳螂就是在这种泡沫的海洋中产卵，繁衍后代的，每当它产下一层卵以后，它就
会往卵上覆盖上一层这样的泡沫。于是，很快地，这层泡沫就将变成固体了。
在新建的巢穴的门外面，有一层材料，把这个巢穴封了起来。看上去，这层材料和
其它的材料并不一样——那是一层多孔、纯洁无光的粉白状的材料。这与螳螂巢内部，
其它部分的灰白颜色是完全不一样的。这就好像面包师们把蛋白、糖和面粉搅合在一起，
用来作饼干外衣的混合物一样。这样一种雪白色的外壳，是很容易破碎的，也很容易脱
落下来。当这层外壳脱落下来的时候，螳螂巢的门口，完全裸露在外，可以看出来，门
的中间装着两行板片。不久以后，风吹雨打会把它侵蚀，将它剥成小片。于是这些小片
会逐渐地脱落下去。所以，在旧巢上，就看不见它的痕迹了。
至于这两种材料，虽然它们从外表上看来，一点儿也不一样，但是实际上，它们的
质地是完全一样的。它们只不过是同样原质的东西的两种不同的表现形式罢了。螳螂用
它身上的杓打扫着泡沫的表面，然后，撇掉表面上的浮皮，使其形成一条带子，覆盖在
巢穴的背面。这看起来，就像那种冰霜的带一样。因此，这种物质实际上仅仅是粘性物
质的最薄、最轻的那一部分。它看上去之所以会比较白一些，主要是因为它的泡沫比较
细巧，光的反射力比较强罢了。
这可真是一个非常奇异的操作方式。它相当有自己的一套方法，可以很迅速、很自
然地做成一种角质的物质。于是，螳螂的第一批卵就生产在这种物质上面了。
螳螂真是一种很能干的动物，也是一种很有建筑才能的动物。产卵时，它排泻出用
于保护的泡沫，制造出柔软得像糖一样的包被物，同时，它还能制做出一种遮盖用的薄
片，以及通行用的小道。而在进行这一切工作的时候，螳螂都只是在巢的根脚处站立着，
一动也不动，用不着移动身体，就在它背后建筑起一座了不起的建筑物，而它自己对这
个建筑物连看都不看一眼。它那粗壮而有力的大腿，在这件事的整个过程中，竟然没有
用武之地，发挥不了什么作用，一点儿都没有做什么。这所有一切的繁杂工作，完全都
是这部小机器自己完成的。
母亲的工作成功完成以后，就洗手不干了，放开一切，跑走了。我总是对它抱着一
线希望，盼望着它有朝一日能够回来看一下，以便表示一些它对整个家族的生产地的爱
护和关切之情。但是，我的这个希望总也得不到实现。看起来很显然，这一点对于它竟
然没有一点儿兴趣了。它真的是一去不回头了。
所以，根据这一事实，我便得出了这个结论：螳螂都是些没有心肝的东西，尽干一
些残忍、恶毒达到极点的事情。比如，它以自己的丈夫作为美餐，而且，它居然还会抛
弃它自己的子女，弃家出走且永不归还。
螳螂卵的孵化，通常都是在有太阳光的地方进行的，而且，大约是在六月中旬，上
午十点钟的时候。
在前面的文章里，我已经告诉过大家了，在这个螳螂巢里，只有一小部分可以为螳
螂幼虫当作出路。这一部分指的就是窠巢里面那一带鳞片的地方。再仔细地观察一下，
你就会发现在每一个鳞片的下面，都可以看见一个物体，而且稍微有一点儿透明的小块
儿。在这个小块儿的后面，紧接着的就是两个大大的黑点。那不是别的什么东西，就是
那个可爱的小动物的一对小眼睛了。幼小的螳螂蛴螬，静静地伏卧在那个薄薄的片下面。
如果仔细地看一下，就会发现它现在差不多已经有将近一半的身体解放了出来。下面，
再看看这个小东西的身体是什么样的吧。它身体的颜色主要是黄色，又带有一些红的颜
色。它长了一个十分肥胖而且很大的脑袋。从这个幼虫的外面的皮肤来看，非常容易地
就能分辨出它的那对特别大的眼睛来。幼虫的小嘴是贴在它的胸部的，腿又是和它的腹
部紧紧相贴的。这只小幼虫，从它的外形上看，除了它那些和腹部紧贴着的腿以外，其
它部分都能够让人联想到另一种动物的状态，那就是那种刚刚才离开巢穴的蝉的最初期
的状态。这种状态和此时的螳螂幼虫像极了。
和蝉一样，为了图方便，但是更重要的是为了安全起见，幼小的螳螂刚一降临到这
个世界上来，的确有穿上一层结实的外套的必要。要是幼虫打算从巢穴中非常狭小而又
弯曲的那条小道里爬出来，如果它想要完全地把自己的小腿伸展开来，那将是不太可能
的一件事了。这主要是因为，如果它完全伸展开身体，高高翘起那尚还缺乏力量的用来
杀戮敌人的长矛，然后再竖立起它那十分灵敏的触须，那么这样一来，它自己就完全把
自己的道路给阻挡住了，以至它无法前行，而根本不可能从通道中爬行出来。正因为如
此，这个小动物，在它刚刚降临到这个世界上的时候，它是被团团包裹在一个襁褓之中
的，那种形状就好像一只小船一般。
在小幼虫刚刚降生，出现在巢中的薄片下面不久以后，它的头便逐渐地变大，一直
膨胀到形状像一粒水泡一样为止。这个有力气的小生命，在出生后不久，就开始靠自己
的力量努力生存。它一刻也不停地一推一缩地努力地解放着自己的躯体。就这样，每每
做一次动作的时候，它的脑袋就要稍稍变大一些。最终的结果是，它的胸部的外皮终于
破裂了。于是，它便更加努力，“乘胜追击”。它摆动得更加剧烈了，也更加快了。它
挣扎着，用尽浑身解数，不停歇地弯曲扭动着它那副小小的躯干。看来，它是义无反顾
地下定决心要挣脱掉这一件外衣的束缚，想马上看到外面的大千世界究竟是个什么样子。
渐渐地，首先得到解放的是它的腿和触须。然后，它继续不懈地努力。又进行了几次摆
动与挣扎以后，终于，它的目的和企图就完全实现了。
有几百只小螳螂，它们同时团团地拥挤在不太宽敞的巢穴之中，这场景，倒真的算
得上是一种不可多得的奇观呢！当巢中的螳螂幼虫还没有集体打破外衣，还没有集体冲
出襁褓，变成螳螂的形态之前，首先暴露出它的那双小眼睛。我们很少会见到哪一个单
独的小动物独自行动。情况恰恰相反，就好像存在什么统一行动的信号一样，每当这信
号传达出来的时候，速度非常之快，几乎所有的卵差不多在同一时刻孵化出来，一起打
破它们的外衣，从硬壳中抽出身体来。因此，也就是在一刹那之间，螳螂巢穴的中部顿
时如同召开大会一样，无数个幼虫一下子集合起来，挤满了这个不太大的地方。它们近
乎狂热一般地爬动着，似乎很兴奋、很急切地要马上脱掉这件困扰它们生活的讨厌的外
衣。在这之后，它们或者是不小心跌落，或者是使劲地爬行到巢穴附近的其它的枝叶上
面去。再过几天以后，就会在巢穴中又发现一群幼虫，它们同样要进行与前辈们相同的
工作，直到它们全都孵化出来。于是，繁衍就这样不停地继续下去。
然而，有一点非常不幸！这些可怜的小幼虫竟然孵化到了一个布满了危险与恐怖的
世界上来，也许它们自己还并不清楚明白这一点。曾经有过好多次，我在门外边的围墙
内，或者是在树林中的那些幽静的地方，看到螳螂的卵在孵化，一个个小幼虫破壳而出。
我总有一种美好的愿望，希望能够尽自己微薄的力量，好好地保护这些可爱的小生命，
让它们能够平平安安而且快快乐乐地生活在这个世界上。但是，很不幸，这种愿望总是
会成为泡影。已经至少有二十次了，(实际上比这要多得多)我总是看到那种非常残暴的
景象，总是亲眼目睹那令人恐惧的一幕。这些还不知道什么叫危险的小幼虫，在它们乳
臭未干的时候，便惨遭杀戮，还没来得及体验一下生活，体会一下生命的宝贵，就已经
结束了年幼的生命，真是可怜啊！螳螂虽然产下了许多卵，但是事实上，它并没有生产
极大数目的卵。至少，它所产下的那些卵，还不足以抵御那些早已在巢穴门口埋伏多时，
一旦幼虫出现，便会不失时机地加以杀戮的强大的敌人。
对于螳螂幼虫而言，它们的最具杀伤力的天敌，要算是蚂蚁了。几乎每一天，我都
会有意无意地看到，一只只蚂蚁不厌其烦地光临到螳螂巢穴的旁边，非常耐心，而且信
心十足地等待时机的成熟，以便立即采取先下手为强的行动。我一看到它们，就千方百
计地帮着螳螂驱赶它们。可是，无济于事。我的能力经常都驱逐不了它们。因为，它们
常常是先人一步，率先占据有利的位置。看来，蚂蚁的时间观念还是很强的，但是，虽
然它们早早就静候在大门之外，可它们却很难深入到巢穴的内部去。这主要是因为，螳
螂巢穴的四围有一层硬硬的厚壁，这便形成了十分坚固的壁垒，因而，蚂蚁对此束手无
策。它们的智慧还不足以使它们想出冲破这一层屏障的办法。不过，它们总是埋伏在巢
穴的门口，静候着它们的俘虏。
因此，螳螂幼虫的处境实在是非常危险的。只要它一不小心跨出自家大门一步，那
么，马上就会坠入深渊，葬送了自己的生命。因为守候在巢边的蚂蚁是不会轻易放过任
何一顿美餐的。一旦有猎物探出头来，便立刻将其擒住，然后再扯掉幼虫身上的外衣，
将其毫不客气地切成碎片。在这场战斗中，你可以看到，那些只能利用随意的乱摆来进
行自我保护的小动物，和那些前来俘虏食品的非常凶猛、残忍的大队的强盗们展开激烈
的拼杀。小动物们尽管非常弱小，但是仍然坚持着、挣扎着，不放弃对求生的渴望。但
是，这种挣扎与那些凶恶之众相比，显得多么可怜呢！用不了多长时间，也就是一小会
儿的工夫，这场充满血腥的大屠杀便宣告终结了。残杀过后，剩余下来的，只不过是碰
巧有幸能够逃脱敌人的恶爪的少数几个幸存者而已。其它的小生命，都已经变成了蚂蚁
口中之食了。就这样，一个原本的人丁兴旺的家族就衰败了。
这是非常奇怪的事情。前面我们曾提到过，螳螂是一种十分凶残的动物。它不仅以
锋利的杀伤性武器去攻击其它的动物，以猎取食物，而且还居然以自己的同类为食，并
且在食用自己的同胞骨肉时，竞然还那样心安理得坦坦荡荡。然而，就是这种可以被视
为昆虫中的灾害的螳螂，如今，在它刚刚拥有生命的初期，本身也要牺牲在昆虫中个儿
头最小的蚂蚁的魔爪下，这难道不奇妙吗！大自然造物真是让人不可思议啊！这种小小
的恶魔，眼睁睁地看着它自己的家族被这样毁掉，它自己的兄弟姐妹被这么一群小小的
侏儒所欺凌，所吞食，却对此束手无策。只能傻傻地目送亲人们远离这个充满危险的世
界。
不过，这样的情形并不会持续相当长的时期。因为，遭到不测的只是那些刚刚问世，
刚刚从卵中孵化出来的幼虫而己。但是，当这些幼虫开始和空气相接触以后，用不了多
长时间，便会马上变得非常的强壮。这样一来，渐渐地，它自己就具备了能够自我保护
的能力了，再也不是那些任人宰割的可怜虫了！
在它长大一些以后，情况就大不相同了。它从蚂蚁的群里，快速地走过去，它所经
过的地方，原来任意行凶的敌人们都纷纷跌倒下来，再也不敢去攻击和欺负这个已经长
大了的“弱者”了。螳螂在行进的时候，把它的前臂放置在胸前，作出一副自卫的警戒
状态。它那种骄傲的态度和不可小视的神气，早已经把这群小小的蚂蚁吓倒了。它们再
也不敢轻举妄动了，有些甚至已经望风而逃了。
但是，事实上，螳螂的敌人，不只是这些小个子的蚂蚁，还有许多其它的敌人。这
些天敌可不是那么容易就能吓倒的。比如说，那种居住在墙壁上面的小形的、灰色的蜥
蜴，就很难对付。对于小小螳螂的自卫和恐吓的姿势，它是全然不在意的。小蜥蜴进攻
螳螂的方法主要是用它的舌尖，一个一个地舐起那些刚刚幸运地逃出蚂蚁虎口的小昆虫。
虽然一个小螳螂还不能填满蜥蜴的嘴，但是，从它的面目表情便可以很清楚地看出来，
那味道却是非常之好。看来，它相当满意。每吃掉一个，蜥蜴的眼皮总是要微微一闭，
这的确是一种极端满足的表现。然而，对于那些年轻的、仍不走运的少年螳螂而言，它
们真可谓“才出龙潭，又入虎穴”啊！
不仅仅是在卵孵化出来以后是如此的危险，甚至就是在卵还没有发育出来以前，它
们就已经处于万分危险之中了。有这样一只小个儿的野蜂(challis)，它随身携带着一种
刺针，其尖利的程度，足可以刺透螳螂的由泡沫硬化以后而形成的巢穴，这样一来，螳
螂的血统，就如同蝉的子孙后代一样，遭受到相同的命运。这样一位外来的客人，并没
有受到谁的邀请，就在螳螂的巢穴中擅自决定产下自己的卵。它的卵的孵化也要比这巢
穴的主人的卵提前一步。于是，螳螂的卵就会顺其自然地受到侵略者的骚扰，被侵略者
吞食掉。比如说螳螂产下一千枚卵，那么，最后剩留下来的，没有遭受噩运而被残酷地
毁灭了的，大概也就只有一对而已了。
这样一来，便形成了下面这条生物链。螳螂以蝗虫为食，蚂蚁又会吃掉螳螂，而蚂
蚁又是鸡的食品。但是，等到了秋天的时候，鸡长大了，长肥了，我又会把鸡做成佳看
吃掉，这可真有趣！
或许螳螂、昨晚、蚂蚁，甚至是其它个儿头更小一些的动物，食用之后都可以增加
人类的脑力。它们采用一种非常奇妙但又见不到的方法，提供给我们的大脑某种有益的
物质。然后，作为我们人类思想之灯的油料。它们的精力慢慢地发达起来，然后贮蓄起
来，并且一点一点地传送到我们身体的各个部位，流进我们的血脉里。它们滋养着我们
身上的不足之处。我们就是生存在它们的死亡之上的。世界本来就是一个永无穷尽的循
环着的圆环。各种物质完结以后，在此基础上，各种物质又纷纷重新开始一切；从某种
意义上讲，各种物质的死，就是各种物质的生。这是一个十分深刻的哲学道理。
很多年前，人们总是习惯性地把螳螂的巢看做是一种充满迷信的东西。在布罗温司
这个地方，螳螂的巢，被人们视为医治冻疮的一种丹灵妙药。大多数人拿一个螳螂的巢，
然后把它劈开成两半，挤出里面的浆汁来，涂抹在痛楚的部位。农村里的人常说，螳螂
巢的功效，就仿佛有什么神奇的魔力一样。然而，我自己从没有感到它有什么功效。
与此同时，也有一些人盛传，认为螳螂巢医治牙痛非常有效。假如你有了它，也就
用不着再怕什么牙痛了。一般情况下，妇女们常常在月夜里到野外去收集它，然后，很
小心地收藏在杯碗橱子的角落里，或者是把它们缝在一个袋子里面，仔细珍藏起来。如
果附近的邻居们，谁要是患了牙痛病的话，就会跑过来，借用它。妇女们管它叫做“铁
格奴(Tigno)”。
如果是脸肿了的病人，他们会说：“请你借给我一些铁格奴，好吗？我现在痛得厉
害呢！”另外的一个就会赶快放下手里的针线活儿，拿出这个宝贝东西来。
而她会很慎重地对朋友说：“你随便做什么都可以，但是不要摘掉它。我只有这么
一个了，而且，现在又是没有月亮的时候！”
没有想到，农民们的这种心理上的简单而幼稚的反应，竟然被十九世纪的一位英国
医生兼科学家所超越。他曾经告诉过我们如此荒唐可笑的事情。他说在那个时候，如果
一个小孩子在树林里迷了路，他可以询问螳螂，让它指点道路。并且，他还说道：“螳
螂会伸出它的一足，指引给他正确的道路，而且很少或是从不会出错的。”
------------------








蜜蜂、猫和红蚂蚁



我希望能够了解更多的关于蜜蜂的故事。我曾听人说起过蜜蜂有辨认方向的能力，
无论它被抛弃到哪里，它总是可以自己回到原处。于是我想亲自试一试。
有一天，我在屋檐下的蜂窝里捉了四十只蜜蜂，叫我的小女儿爱格兰等在屋檐下，
然后我把蜜蜂放在纸袋里，带着它们走了二里半路，接着打开纸袋，把它们抛弃在那里，
看有没有蜜蜂飞回来。
为了区分飞到我家屋檐下的蜜蜂是否是被我扔到远处的那群，我在那群被抛弃的蜜
蜂的背上做了白色的记号。在这过程中，我的手不可避免地被刺了好几口，但我一直坚
持着，有时候竟然忘记了自己的痛，只是紧紧地按住那蜜蜂，把工作做完，结果有二十
多只损伤了，当我打开纸袋时，那些被闷了好久的蜜蜂一拥而出地向四面飞散，好像在
区分该从哪个方向回家一样。
放走蜜蜂的时候，空中吹起了微风。蜜蜂们飞得很低，几乎要触到地面，大概这样
可以减少风的阻力，可是我想，它们飞得这样低，怎么可以眺望到它们遥远的家园呢？
在回家的路上，我想到它们面临的恶劣环境，心里推测它们一定都我不到回家的方
向了。可是没等我跨进家门，爱格兰就冲过来，她的脸红红的，看上去很激动。她冲着
我喊道：
“有两只蜜蜂回来了！在两点四十分的时候到达巢里，还带来了满身的花粉。”
我放蜜蜂的时间是两点整。也就是说，在三刻钟左右的时间里，那两只小蜜蜂飞了
二里半路，这还不包括采花粉的时间。
那天天快黑的时候，我们还没见到其它蜜蜂回来。可是第二天当我检查蜂巢时，又
看见了十五只背上有白色记号的蜜蜂回到巢里了。这样，二十只中有十七只蜜蜂没有迷
失方向，它们准确无误地回到了家，尽管空中吹着逆向的风，尽管沿途尽是一些陌生的
景物。但它们确确实实地回来了。也许是因为它们怀念着巢中的小宝贝和丰富的蜂蜜。
凭借这种强烈的本能，它们回来了。是的，这不是一种超常的记忆力，而是一种不可解
释的本能，而这种本能正是我们人类所缺少的。

猫
我一直没有相信过这样一种说法，即猫也和蜜蜂一样，能够认识自己的归途。直到
有一天我家的猫的确这样做了，我才不得不相信这一事实。
有一天，我在花园里看见一只并不漂亮的小猫，薄薄的毛皮下显露着一节一节的脊
背，瘦骨嶙峋的。那时我的孩子们还都很小，他们很怜惜这只小猫，常塞给它一些面包，
一片一片还都涂上了牛乳。小猫很高兴地吃了好几片，然后就走了。尽管我们一直在它
后面温和地叫着它，“咪咪，咪咪——”，它还是无怨无悔地走了。可是隔了一会儿，
小猫又饿了。它从墙头上爬下来，又美美地吃了几片。孩子们怜惜地爱抚着它瘦弱的身
躯，眼里充满了同情。
我和孩子们作了一次谈话，我们达成一致，决定驯养它。后来，它果然不负众望，
长成一只小小的“美洲虎”——红红的毛，黑色的斑纹，虎头虎脑的，还有锋利的爪子。
它的小名叫做“阿虎”。后来阿虎有了伴侣，她也是从别处流浪而来的。他们俩后来生
了一大堆小阿虎。不管我家有什么变迁，我一直收养着它们，大约有二十多年了。
第一次搬家时，我们很为它们担忧，假如遗弃这些我们所宠爱的猫，它们将再度遭
受流浪的生活。可是如果把它们带上的话，雌猫和小猫们还能稳住气，保持安静，可两
只大雄猫——一只老阿虎，一只小阿虎在旅途上是一定不会安静的。最后我们决定这样：
把老阿虎带走，把小阿虎留在此地，替它另外找一个家。
我的朋友劳乐博士愿意收留小阿虎。于是某天晚上，我们把这只猫装在篮子里，送
到他家去。我们回来后在晚餐席上谈论起这只猫，说它运气真不坏，找到了一户人家。
正说着，突然一个东西从窗口跳进来。我们都吓了一跳，仔细一看，这团狼狈不堪的东
西快活而亲切地用身体在我们的腿上蹭着，这正是那只被送掉的小阿虎。
第二天，我们听到了关于它的故事：它刚到劳乐博士家里，就被锁在一间卧室里。
当它发现自己已在一个陌生的地方做了囚犯时，它就发狂一般地乱跳。一会儿跳到家俱
上，一会儿跳到壁炉架上，撞着玻璃窗，似乎要把每一样东西都撞坏。劳乐夫人被这个
小疯子吓坏了，赶紧打开窗子，于是它就从窗口里跳了出来。几分钟之后，它就回到了
原来的家。这可不是容易的事啊，它几乎是从村庄的一端奔到另一端，它必须经过许多
错综复杂的街道，其间可能遭遇到几千次的危险，或是碰到顽皮的孩子，或是碰到凶恶
的狗，还有好几座桥，我们的猫不愿意绕着圈子去过桥，它决定拣取一条最短的路径，
于是它就勇敢地跳入水中——他那湿透了的毛告诉了我们一切。
我很可怜这只小猫，它对它的家是如此的忠心。我们都同意带它一起走，正当我们
担心它在路上会不安分的时候，这个难题竟自动解决了。几天之后，我们发现它已经僵
硬地躺在花园里的矮树下。有人已经替我把它毒死了。是谁干的呢？这种举动可不会是
出自好意！
还有那只老阿虎。当我们离开老屋的时候，却怎么也找不到它了。于是我们另外给
车夫两块钱，请他负责找那老阿虎，无论什么时候找到它，都要把它带到新家这边来。
当车夫带着最后一车家俱来的时候，他把老阿虎带来了。他把它藏在自己的座位底下。
当我打开这活动囚箱，看到这前两天就被关进去的囚徒的时候，我真不能相信它就是我
的老阿虎了。
它跑出来的时候，活像一只可怕的野兽，它的脚爪不停在张舞着，口里挂着口水，
嘴唇上沾满了白沫，眼睛充满了血，毛已经倒竖起来，完全已经没有了原来的阿虎的神
态和风采。难道它发疯了吗？我仔细把它察看了一番。我终于明白了，它没有疯，只是
被吓着了。可能是车夫捉他的时候把它吓坏了，也可能是长途的旅行把它折磨得精疲力
尽。我不能确定到底是什么原因，但显而易见的是，它的性格大变，它不再口中常常念
念有词，不再用身体擦我们的腿了，只有一副粗暴的表情和深沉的忧郁。慈爱的抚慰也
不能消除它的苦痛了。终于有一天，我们发现它死了，躺在火炉前的一堆灰上，忧郁和
衰老结束了它的生命。如果它精力还足够的话，它会不会回到我们的老房子去呢？我不
敢断定。但是，这样一个小生灵，因为衰老的体力不允许它回到老家，终于得了思乡病，
忧郁而死，这总是一件令人感慨的事吧！
当我们第二次搬家的时候，阿虎的家族己完全换了一批了：老的死了，新的生出来
了。其中有一只成年的小阿虎，长得酷像它的先辈。也只有它会在搬家的时候增加我们
的麻烦。至于那些小猫咪和母亲们，是很容易制服的。只要把它们放在一只篮子里就行
了，小阿虎却得被单独放在另一只篮子里，以免它把大家都闹得不太平。这样一路上总
算相安无事。到了新居后，我们先把母猫们抱出篮子。它们一出篮子，就开始审视和检
阅新屋，一间一间地看过去，靠着它们粉红色的鼻子，它们嗅出了那些熟悉的家俱的气
味。它们找到了自己的桌子、椅子和铺位，可是周围的环境确实变了。它们惊奇地发出
微微的“喵喵”声，眼睛里时时闪着怀疑的目光。我们疼爱地抚摸着它们，给它们一盆
盆牛奶，让它们尽情享用。第二天它们就像在自己家里一样习惯了。
可是轮到我们的小阿虎，情形却完全不同了。我们把它放到阁楼上，让它渐渐习惯
新环境，那儿有好多空屋可以让它自由地游玩。我们轮流陪着它，给它加倍的食物，并
时时刻刻把其余的猫也捉上去和它作伴。我们想让它知道，它并不是独自一个在这新屋
里。我们想尽了一切办法，让它忘掉原来的家。果然，它似乎真的忘记了。每当我们抚
摸它的时候，它显得非常温和驯良，一叫它，它就会“咪咪”地叫着过来，还把背弓起
来。这样关了一个星期，我们觉得应该恢复它的自由了，于是把它从阁楼上放了出来。
它走进了厨房，和别的猫一同站在桌子边。后来它又走进了花园。我的女儿爱格兰紧紧
地盯着它看它有什么异样的举动，只见它做出一副非常天真的样子，东张张，西望望，
最后仍回到屋里。太好了，小阿虎再也不会出逃了。
第二天，当我们唤它的时候，任凭我们叫了多少声“咪咪咪咪——”，就是没有它
的影子！我们到处找，呼唤它，丝毫没有结果。骗子！骗子！我们上了它的当！它还是
走了，我说它是回到老家去了。可是家里其他人都不相信。
我的两个女儿为此特意回了一次老家。正如我说的那样，她们在那里找到了小阿虎。
她们把它装在篮子里又带了回来，虽然天气很干燥，也没有泥浆，可它的爪子上和腹部
都沾满了沙泥，无疑它一定是渡过河回老家去的，当它穿过田野的时候，泥土就粘在了
它湿漉漉的毛上，而我们的新屋，距离原来的老家，足足有四里半呢！
我们把这个逃犯关在阁楼上，整整两个星期之后，再放它出来。可是还不到一天工
夫，它又跑回去了，对于它的前途，我们只能听天由命了。后来有一位老屋的邻居来看
我们，说起小阿虎，说他有一次看到我们的猫口里叼着一只野兔，躲在篱笆下。是啊，
再也没有人喂给它食物了，它得用自己的力量去寻找食物。后来我就再也没有听到过它
的消息了。它的结局一定是挺悲惨的。它变成了强盗，当然要遭到强盗的命运。
这些真实的故事证明了猫和泥匠蜂一样，有着辨别方向的本领。鸽子也是这样，当
它们被送到几百里以外的时候；它们还能回来找到自己的老巢。还有燕子，还有许多别
的鸟也是这样。让我们再回到昆虫的问题上吧。蚂蚁和蜜蜂是最相似的一对昆虫，我很
想知道它们是不是像蜜蜂一样有着辨别方向的本领。

红蚂蚁
在一块废墟上，有一处地方是红蚂蚁的山寨。红蚂蚁是一种既不会抚育儿女也不会
出去寻找食物的蚂蚁，它们为了生存，只好用不道德的办法去掠夺黑蚂蚁的儿女，把它
们养在自己家里，将来这些被它们占为已有的蚂蚁就永远沦为了奴隶。
夏天的下午，我时常看见红蚂蚁出征的队伍，这队伍大约有五、六码长。当它们看
见有黑蚂蚁的巢穴时，前面的队伍出现一阵忙乱。几只间谍似的蚂蚁先离开了队伍往前
走。一队的蚂蚁仍旧列着队伍蜿蜒不停地前进，有时候有条不紊地穿过小径，有时在荒
草的枯叶中若隐若现。
最后，它们终于找到了黑蚂蚁的巢穴，就长驱直入地进入到小蚂蚁的卧室里，把它
们抱出了巢。在巢内，红蚂蚁和黑蚂蚁有过一番激烈的厮杀，最终黑蚂蚁败下阵来，无
可奈何地让强盗们把自己的孩子抢走。
我再讲一下它们一路上怎样回去的情形吧。
有一天我看见一队出征的蚂蚁沿着池边前进，那时天刮着大风，许多蚂蚁被吹落了，
白白地做了鱼的美餐。这一次鱼又多吃了一批意外的食物——黑蚂蚁的婴儿。显然蚂蚁
不会像蜜蜂那样，会选择另一条路回家，它们只会沿着原路回家。
我不能把整个下午都消耗在蚂蚁身上，所以我叫小孙女拉茜帮我监视它们。她喜欢
听蚂蚁的故事，也曾亲眼看到红蚂蚁的战争，她很高兴接受我的嘱托。凡是天气不错的
日子里，小拉茜总是蹲在园子里，瞪着小眼睛往地上张望。
有一天，我在书房里听到拉茜的声音，“快来快来！红蚂蚁已经走到黑蚂蚁的家里
去了！”
“你知道它们走的是哪条路吗？”
“是的，我已经做了记号。”
“什么记号？你怎么做的？”
“我沿路撒了小石子。”
我急忙跑到园子里，拉茜说得没错。红蚂蚁们正沿着那一条白色的石子路凯旋呢！
我取了一张叶子，截走几只蚂蚁，放到别处。这几只就这样迷了路，其它的，凭着它们
的记忆力顺着原路回去了。这证明它们并不是像蜂那样，直接辨认回家的方向，而是凭
着对沿途景物的记忆找到回家的路的。所以即使它们出征的路程很长，需要几天几夜，
但只要沿途不发生变化，它们也照旧回得来。
------------------








开隧道的矿蜂



矿蜂是细长形的蜜蜂，它们的身材大小不同，大的比黄蜂还大，小的比苍蝇还小。
但是它们有一个共同的特征，那就是它的腹部的底端有一条明显的沟，沟里藏着一根刺，
遇到敌人来侵犯时，这根刺可以沿着沟来回地移动，以保护自己。我这里要讲的是关于
矿蜂中的一种有红色斑纹的蜂。雌蜂的斑纹是很美丽夺目的，细长的腹部被黑色和褐色
的条纹环绕着。至于它的身材，大约和黄蜂差不多。
它的巢往往建在结实的泥土里面，因为那里没有崩溃的危险。比如，我们家院子里
那条平坦的小道就是它们最理想的屋基。每到春天，它们就成群结队地来到这个地方安
营扎寨。每群数量不一，最大的差不多有上百只黄蜂。这地方简直成了它们的大都市。
每只蜜蜂都有自己单独的一个房间。这个房间除了它自己以外，谁也不可以进去。
如果有哪只不识趣的蜜蜂想闯进别人的房间，那么主人就会毫不客气地给它一剑。因此，
大家都各自守着自己的家，谁也不冒犯谁，这个小小的社会充满了和平的气氛。
一到四月，它们的工作就不知不觉地开始了。唯一可以显而易见地证明它们在工作
的，是那一堆堆新鲜的小土山。至于那些劳动者，我们外人是很少有机会看到的。它们
通常是在坑的底下忙碌着，有时在这边，有时在那边。我们在外面可以看到，那小土堆
渐渐地有了动静，先是顶部开始动，接着有东西从顶上沿着斜坡滚下来，一个劳动者捧
着满怀的废物，把它们从土堆顶端的开口处抛到外面来，而它们自己却用不着出来。
五月到了，太阳和鲜花带来了欢乐。四月的矿工们，这时已经演变成勤劳的采蜜者
了。我们常常看到它们满身披着黄色的尘土停在土堆上，而那些土堆现在已变得像一只
倒扣着的碗了，那碗底上的洞就是它们的入口了。
它们的地下建筑离地面最近的部分是一根几乎垂直的轴，大约有一支铅笔那么粗，
在地面下约有六寸到十二寸深，这个部分就算是走廊了。
在走廊的下面，就是一个小小的巢。每个小巢大概有四分之三寸长，呈椭圆形。那
些小巢有一个公共的走廊通到地上。
每一个小巢内部都修聋得很光滑，很精致。我们可以看出一个个淡淡的六角形的印
子，这就是它们作最后一次工程时留下的痕迹。它们用什么工具来完成这么精细的工作
呢？是它们的舌头。
我曾经试图往巢里面灌水，看看会有什么后果，可是水一点儿也流不到巢里去。这
是因为斑纹蜂在巢上涂了一层唾液，这层唾液像油纸一样包住了巢，在下雨的日子里，
巢里的小蜜蜂就再也不会被弄湿了。
斑纹蜂一般在三、四月里筑巢。那时候天气不大好，地面上也缺少花草。它们在地
下工作，用它的嘴和四肢代替铁锹和把子。当它们把一堆堆的泥粒带到地面上后，巢就
渐渐地做成了。最后用它的铲子——舌头，涂上一层唾液。当快乐的五月到来时，地下
的工作已经完毕，那和照的阳光和灿烂的鲜花也已经开始向它们招手了。
田野里到处可以看到蒲公英、野蔷薇、雏菊花等，在花丛里尽是些忙忙碌碌的蜜蜂。
它们带上花蜜和花粉后，就兴高采烈地回去了。它们一回到自己的城市里，就会立即改
变飞行方式，它们很低地盘旋着，好像对这么多外观酷似的地穴产生了迟疑，不知道哪
个才是自己真正的家。但是没过一会儿，它们就各自认清了自己的记号，很快地，准确
无误地钻了进去。
斑纹蜂也像其它蜜蜂一样，每次采蜜回来，先把尾部塞入小巢，刷下花粉，然后一
转身，再把头部钻入小巢，把花蜜洒在花粉上，这样就把劳动成果储藏起来了。虽然每
一次采的花蜜和花粉都微乎其微，但经过多次的采运，积少成多，小巢内已经变得很满
了。接着斑纹蜂就开始动手制造一个个“小面包”，“小面包”是我给那些精巧的食物
起的名字。
斑纹蜂，开始为它未来的子女们预备食品了。她把花粉和花蜜搓成一粒粒豌豆大小
的“小面包。”这种“小面包”和我们吃的小面包大不一样：它的外面是甜甜的蜜质，
里面充满了干的花粉，这些花粉不甜，没有味道。这外面的花蜜是小蜜蜂早期的食物，
里面的花粉则是小蜜蜂后期的食物。
斑纹蜂做完了食物，就开始产卵。它不像别的蜜蜂产了卵后就把小巢封起来，它还
要继续去采蜜，并且看护它的小宝宝。
小蜜蜂在母亲的精心养护和照看之下渐渐长大了。当它们作茧化蛹的时候，斑纹蜂
就用泥把所有的小巢都封好。在它完成这项工作以后，也到了该休息的时候了。
如果没有什么意外发生的话，在短短的两个月之后，小蜜蜂就能像它们的妈妈一样
去花丛中玩耍了。

温厚长者和小强盗
可是斑纹蜂的家并不像想象中那样安逸，在它们周围埋伏着有许多凶恶的强盗。其
中有一种蚊子，虽然小得微不足道，却是矿蜂的劲敌。
这种蚊子是什么样的呢？它的身体不到五分之一寸长，眼睛是红黑色的；脸是白色
的；胸甲是黑银灰色的，上面有五排微小的黑点儿，长着许多刚毛；腹部是灰色的；腿
是黑色的，像一个又凶恶又奸诈的杀手。
在我所观察到的这一群蜂的活动范围内，就有许多这样的蚊子。这些蚊子在太阳底
下时能找到一个隐蔽的地方潜伏起来。等到斑纹蜂携带着许多花粉过来时，蚊子就紧紧
地跟在它后面，跟着打转、飞舞。忽然，斑纹蜂俯身一冲，冲进自己的屋子。立刻，蚊
子也跟着在洞口停下，头向着洞口，就这样等了几秒钟，蚊子纹丝不动。
它们常常这样面对着面，彼此只隔一个手指那么宽的距离僵持着。但彼此都显得十
分镇定。斑纹蜂这温厚的长者，只要它愿意，它完全有能力把门口那个破坏它家庭的小
强盗打倒，它可以用嘴把它咬阵，可以用刺把它刺得遍体鳞伤，可它并没有这么做。它
任凭那小强盗安然地埋伏在那里。至于那小强盗呢？虽然有强大的对手在它眼前虎视耽
耽，而那可恶的小蚊子尽管知道斑纹蜂只要举手之劳就可以把它撕碎，可它丝毫没有恐
惧的样子。
不久斑纹蜂就飞走了。蚊子便开始行动了。它飞快地进入了巢中，像回到自己的家
里那样不客气。现在它可以在这储藏着许多粮食的小巢里胡作非为了。因为这些巢都还
没有封好。它从从容容地选好一个巢，把自己的卵产在那个巢里。在主人回来之前，它
是安全的，谁也不会来打扰它，而在主人回来之时，它早已完成任务，拍拍屁股逃之夭
夭了。它会再在附近找一处藏身之处，等候着第二次盗窃的机会。
几个星期后，让我们再来看看斑纹蜂藏在巢里的花粉团吧，我们将发现这些花粉团
已被吃得狼藉一片。在藏着花粉的小巢里，我们会看到几条尖嘴的小虫在蠕动着——它
们就是蚊子的小宝宝，在它们中间，我们有时候也会发现几条斑纹蜂的幼虫——它们本
该是这房子的真正的主人，却已经饿得很瘦很瘦了。那帮贪吃的入侵者剥夺了原该属于
它们的一切。这可怜的小东西渐渐地衰弱，渐渐地萎缩，最后竟完全消失了。那凶恶的
蚊子的幼虫就一口一口把这尸体也吞下去了。
小蜜蜂的母亲虽然常常来探望自己的孩子，可是它似乎并没有意识到巢里已经发生
了天翻地覆的变化。它从不会把这陌生的幼虫杀掉，也不会毫不犹豫地把它们抛出门外，
它只知道巢里躺着它亲爱的小宝贝。它认真地小心地把巢封好，好像自己的孩子正在里
面睡觉一样。其实，那时巢里已经什么都没有了，连那蚊子的宝宝也早已趁机飞走了。
多么可怜的母亲啊！

老门警
斑纹蜂的家里如果没有碰到意外，也就是说没有像刚才我说的那样被蚊子所偷袭，
那么它们大约应有十个姐妹。为了节约时间和劳动力，它们不再另外挖隧道，只要把它
们的母亲遗留下来的老屋拿过来继续用就是了。大家都客客气气地从同一个门口进出，
各自做着自己的工作，互不打扰。不过在走廊的尽头，它们有各自的家，每一个家包括
一群小屋，那就是它们自己挖的，不过那走廊是公用的。
让我们来看看它们是怎样来来去去地忙碌的吧。当一只采完花蜜的蜜蜂从田里回来
的时候，它的腿上都沾满了花粉。如果那时门正好开着，它就会立刻一头钻进去。因为
它忙得很，根本没有空闲时间在门口徘徊。有时候会有几只蜜蜂同时到达门口的情况，
可那隧道的宽度又不允许两只蜂并肩而行，尤其是在大家都满载花粉的时候，只要轻轻
一触就会把花粉都掉到地上，半天的辛勤劳动就都白废了。于是它们定了一个规矩：靠
近洞口的一个赶紧先进去，其余的依次在旁边排着队等候。第一个进去后，第二个很快
地跟上，接着是第三个，第四个，第五个……大家都排着队很有秩序地进去。
有时候也会碰到这样的情况：一只蜂刚要出来，而另一只正要进去。在这种情况下，
那只要进去的蜂会很客气地让到一边，让里面的那只蜜蜂先出来，每只蜜蜂在自己的同
类面前，都表现得非常有风度，有礼貌。有一次我看到一只蜂已经从走廊到达洞口，马
上要出来了，忽然，它又退了回去，把走廊让给刚从外面回来的蜂。多有趣啊！这种互
助的精神实在令人佩服，有了这样一种精神，它们的工作才能很快地进行。
让我们把眼睛睁大些仔细地观察，还有比这更有趣的事呢！当一只蜜蜂从花田里采
了花粉回到洞口的时候，我们可以看到一块堵住洞口的活门忽然落下，开出一条通路来。
当外来的蜂进去以后，这活门又升上来把洞口堵住。同样，当里面的蜜蜂要出来的时候，
这活门也是先降下，等里面的蜜蜂飞出去后，又升上来关好。
这个像针筒的活塞一般忽上忽下的东西究竟是什么呢？这是一只蜂，是这所房子的
门卫。它用它的大头顶住了洞口。当这所房子的居民要进进出出的时候，它就把“门栓”
一拔，也就是说，它立刻退到一边，那儿的隧道特别宽大，可以容得下两只蜂。当别的
蜜蜂都通过了，这“门警”又上来用头顶住洞口。它一动不动地守着门，那样的尽心尽
责，除非它不得不去驱除一些不知好歹的不速之客，否则它是不会擅自离开岗位的。
当这位门警偶尔走出洞口的时候，让我们趁机仔仔细细地看看它吧。我们发现它和
其它蜂一样，不过它的头长得很扁，它的衣服是深黑色的，并且有着一条条的纹路。身
上的绒毛已经看不出来了，它本该有的那种美丽的红棕色的花纹也没有了。这一套破碎
的衣服似乎告诉了我们一切。
这一只用自己的身躯顶住门口充当老门警的蜜蜂看起来比谁都显得沧桑和年老。事
实上它正是这所屋子的建筑者，现在的工蜂的母亲，现在的幼虫的祖母。就在三个月之
前，它还挺年轻的，那时候它正在独自辛辛苦苦地建筑这座房子。现在它算是告老退休
了——不，这不是退休，它还要发挥它的余热，用它的全力来保护着这个家呢。
你还记得那多疑的小山羊的故事吗。它从门缝里往外张望一下，然后对门外的狼说，
“你是我们的妈妈吗？请你把白腿伸给我看，如果你的腿是黑色的，我们就不开门。”
我们这位老祖母的警惕绝不亚于那小山羊。它对每一位来客说：
“把你的蜜蜂黑脚伸给我看，否则我就不让你进来。”
只有当它认出这是它家的一员时，它才会开门，否则它是决不会让任何外客进入到
它家里去的。
你看！在洞旁走过一只蚂蚁，它是一个大胆的冒险家。它很想知道这个散发着一阵
阵蜂蜜香味的地方究竟是怎样的。
“滚开！”老蜜蜂摇了摇头说道。
蚂蚁被它吓了一跳，悻悻地走开了。也幸亏它走开了。如果它仍逗留在蜂房旁的话，
老蜜蜂就要离开它的岗位，飞过去不客气地追击它了。
也有一种不擅长挖隧道的蜜蜂，也就是樵叶蜂，它要寻找人家从前挖掘好的隧道。
斑纹蜂的隧道对它再适合不过了。那些以前受蚊子偷袭，被蚊子占据的斑纹蜂的巢一直
是空着的。因为蚊子让它们家绝了后，整个家都已经败落了。于是樵叶蜂就可以顺理成
章地占据这个空巢，来个废物利用了。为了找到这样的空巢，以便于让它们放那些用枯
叶做成的蜜罐，这帮樵叶蜂常到我的这种斑纹蜂的领地里来巡视。有时候它似乎找到目
标了，可还没等它的脚站稳，它的嗡嗡声已引起了门警的注意。门警立刻冲出洞来，在
门口作了几下手势，告诉它这洞早就有主人了。樵叶蜂明白了它的意思，立即飞到别处
去找房子了。
有时候没等门警出来，樵叶蜂已经迫不及待地把头伸了进去。于是做门警的老祖母
立刻把头顶上来塞住通路，并且发出一个并不十分严厉的信号，以示警告，樵叶蜂立即
明白了这屋子的所有权，很快就离开了。
有一种“小贼”，它是樵叶蜂的寄生虫，有时候会受到斑纹蜂的教训。有次我亲眼
看到它受了一顿重罚。这鲁莽的东西一闻进隧道便为非作歹，以为自己进了樵叶蜂的家
了。可是不一会儿，它立刻发现自己犯了一个大错误，它闯进的是斑纹蜂的家。它碰到
了守门的老祖母，受了一顿严厉的惩罚。于是它急急忙忙地往外逃。同样，其它野心勃
勃又没有头脑的傻瓜，如果想闯进斑纹蜂的家，毫无疑问它将受到同样的待遇。
有时候守门的蜜蜂也会和另外一位老祖母发生争执。七月中旬，是蜜蜂们最忙的时
候。这时候我们会看到两种迥然不同的蜂群：那就是老蜜蜂和年轻的母蜂。年轻的母蜂
又漂亮又灵敏，忙忙碌碌地从花间飞到巢里，又从巢里飞向花间。而那些老蜂，失去了
青春，失去了活力。只是从一个洞口踱到另一个洞口。看上去就像迷失了方向找不到自
己的家。这些流浪者究竟是谁呢？它们就是那些受了可恶的小强盗的蒙骗而失去家庭的
老蜜蜂。当初夏来临的时候，老蜜蜂终于发现从自己的巢里钻出来的是可恶的蚊子，这
才恍然大悟、痛心疾首。可是这已经太晚了，它已经变成了无家可归的孤老。它只好委
屈地离开自己的老家，到别处去另谋生路了。看看哪一家需要一个管家或是需要一个门
警。可是那些幸福的家庭早已有了自己的祖母来打点一切了。而这些老祖母往往对外来
找工作，抢自己饭碗的老蜜蜂心存敌意，往往会给它一个不客气的答复。的确，一个家
只需要一个门警就早已足够了。来了两个的话，反而把那原本就不宽敞的走廊给堵住了。
有些时候，两个老祖母之间真的会发生一阵恶斗。当流浪的老蜜蜂停在别家门口的
时候，这家的看门老祖母一方面紧紧守着门，一方面张牙舞爪地向外来的老蜂挑战，而
胜的那一方，往往是那身心疲惫，悲伤孱弱的老孤蜂。
这些无家可归的老蜜蜂后来怎样了呢？它们一天一天地衰老下去，渐渐数目也少了
起来，最后全部绝迹了。它们有的是被那些灰色的小蜥蜴吃掉了，有的是饿死的，有的
是老死了，还有的是万念俱灰，心力衰竭而死。
至于那守门的老祖母，它似乎从来不休息，在清晨天气还很凉快的时候，它已经到
达它的岗位，到了中午，正是工蜂们采蜜工作最忙的时候，许许多多蜜蜂从洞口飞进飞
出，它仍旧守护在那里；到了下午，外边很热，工蜂都不去采蜜，留在家里建造新的巢，
这时候，老祖母仍旧在上面守着门。在这种闷热的时候，它连瞌睡都不打一下，它不能
打瞌睡，这个家的安全都靠它了。
到了晚上，甚至是深夜，别的蜜蜂都休息了，它还像白天一样忙，防备着夜里的盗
贼。
在它小心的守护下，整个蜂巢的安全可以一直持续到五月以后。如果那蚊子来抢巢，
让它来吧，老祖母会立即冲出去和它拼个你死我活。但它们不会来。因为在明年冬天到
来之前，它们还是躲在茧子里的蛹。
虽然没有蚊子，其它的寄生虫类也不少。它们也很可能来侵犯蜂巢。但是，奇怪的
是，我天天认真观察那个蜂巢，却从没有在它的附近发现什么蜂类的敌人。整个夏天它
都那么安静而平和。可见那些暴徒己深知老祖母的厉害。同时也可见老祖母是如何的警
觉了。
------------------








萤



一、它的外科器具
在昆虫的各种类型中，很少有能够发光的。但其中有一种是以发光而出名的。这个
稀奇的小动物的尾巴上像挂了一盏灯似的，用来表达它对快乐生活的美好祝愿。即便是
我们不曾与它相识，不曾见过它在黑夜中从草丛上飞过，不曾见过它从圆月上落下来。
就像一点点火星儿一样，那么，至少从它的名字上，我们可以多少对它有一些了解。古
代的时候，希腊人曾经把它叫做亮尾巴，是很形象的一个名字。现代，科学家们则给它
起了一个新的名字，叫做萤火虫(Iamqyris)。
事实上，萤无论如何也不是蠕虫，即便是从它的外表上来看，它也不能算作是蠕虫。
它生长着六只短短的腿，而且，它能够知道如何去利用这些短足。从某种意义上讲，它
可真算得上一位真正的闲游家。雄性的萤到了它发育完全的时候，会生长出翅盖，像真
的甲虫一样。不过，事实上，它也就是甲虫的一种。然而，雌性的萤却不怎么引人注意，
它对于飞行的快乐，却是一无所知的。真的有些可怜，它不能完全懂得世界上还有自由
飞行这种快乐可以享受。它终身都处于幼虫的状态，也就是说处于一种不完全的形状，
似乎永远也长不大。可是，即便在这种状态下，使用“蠕虫”这个名词也是很不贴切，
不得当的。我们法国人，经常用“像蠕虫一样的粉光”一语来表示这些没有一点儿保护
和遮掩的动物，不过现在，萤却是有衣服的。可以说，外皮就是它的衣服，它用自己的
外皮来保护自己。而且，它的外皮还具有很丰富的颜色呢。它全身是黑棕色的，但是胸
部有一些微红。在它的身体的每一节的边沿部位，还装饰着一些粉红色的斑点。像这样
色彩丰富的衣服，蠕虫是不会穿的！
尽管如此，我们还是继续把它叫做发光的蠕虫。因为，它的这个名字是全世界的人
早就已熟知的了。
萤，有两个最有意思的特点：第一、就是它获取食物的方法；第二、就是它的尾巴
上有灯。
有一位著名的法国研究食物的科学家曾经说过：“告诉我，你吃的是什么东西，那
么我就会告诉你，你究竟是什么东西。”
同样的问题，都应该对任何昆虫提出，我们想要研究的东西就是这些昆虫们的生活
习性——因为，有关昆虫的食品供给方面的知识，是动物生活中最主要的问题，就是我
们人类常说的“民以食为天”，它也就成了我们应该重点研究的问题之一。虽然，从萤
的外表来看，它似乎是一个纯洁善良而可爱的小动物，但是，事实上，它却是一个凶猛
无比的食肉动物。它是一个善于猎取山珍野味的猎人，而且，它的捕猎方法是十分凶恶
的。看来，它的外表也像其他一些昆虫一样具有一定的欺骗性。通常，它的俘虏对象主
要是一些蜗牛。这个事实，早就被人们认识到了。而人们所不知道的，鲜为人知的，只
是萤的那种有些稀奇古怪的捕捉食品的方法。这个方法，我在其他的地方还没有看到过
相同的例子，可见其独特性并非一般。
萤的这种捕食方法，具体情况如下所述：在它开始捉食它的俘虏以前，它总是先要
给它打一针麻醉药，使这个小猎物失去知觉，从而也就失去了防卫抵抗的能力，以便它
捕捉并食用。这就好比我们人类在动手术之前，在病床上，先接受麻醉，从而渐渐失去
知觉而不感到疼痛一样。在一般情况下，萤所猎取的食物，都是一些很小很小的蜗牛。
很少能捕捉到比樱桃大的蜗牛。在气候非常炎热的时候，就会在路旁边的枯草或者是麦
根上，聚集着大群的蜗牛，像集体纳凉一般。也许是酷热难耐的原因，它们一动也不动
地群伏在那些地方。生怕动一动，就会觉得热气逼人。它们就是这样静止着，懒洋洋地
经过炎热的夏天的。在这些地方，我常常会看到，一些萤在咀嚼它们那已经失去知觉的
俘虏。萤就是在这些摇摆不定的物体上把它们麻醉了的。
除了上述路边的枯草、麦根等地方以外，萤也常常选择一些其他可以获得食物的地
方出没或停留。比如说，它常到一些又凉又潮湿的阴暗的沟渠附近去蹓跶。因为在这些
地方，经常会有一些杂草丛生。在那里经常可以找到大量的蜗牛。这可是难得的美餐啊！
饱饱地吃上几顿是没问题的。通常在这些地方，萤把它们的俘虏在地上杀死，就像人们
说的就地处决一样，干净利落地结束战斗。然后，获得丰厚的战利品。在我自家的屋子
里，我也可以经常制造出这样的地方，来吸引萤到这里来捕食。因此，待在家里，我便
能制造出一个战场，并且，更便了我非常仔细地观察它们的一举一动。
那么，下面我就来叙述一下这种奇怪的情形吧。我拿了一个大玻璃瓶，然后再找一
点儿小草，把草放到大玻璃瓶子里面，再往里边放进几只萤，还有一些蜗牛。我取的蜗
牛，大小比较适中，不算特别的大，也不算特别的小。这一切准备工作就绪以后，我们
所需要继续进行的工作，就是等待，而且，必须要耐心地等待。不过，最为重要的一点
是必须十分留心，时刻注视着玻璃瓶中发生的一切动静，哪怕是最微小的动作，也不能
轻易放过。因为，整个事情的发生，是在非常不经意的时候，而且，时间持续也不长，
几乎就是一会儿的时间。所以，必须目不转睛地紧紧盯住瓶中的这些生灵。
不一会儿，玻璃瓶中就有事情要发生了。萤已经开始注意到它的牺牲品了。看起来，
蜗牛对于萤而言，有极强的，难以抗拒的吸引力。按照通常情况下，蜗牛的习性，除去
外套膜的边缘的地方，它的身体会微微露出一点儿以外，其余的躯体全部都隐藏在它的
家中——即它背上的壳子里面，可能它觉得这样会更安全一些。于是，这位猎人跃跃欲
试，准备发起总攻了。它先做的事情，就是把自己身上随身携带着的兵器迅速地抽出来。
这件兵器是何等地细小啊，要是没有放大镜的帮助，简直是一点儿也看不出来。萤的身
上长有两片颚，它们分别弯曲起来，再合拢到一起，就形成了一把钩子，一把尖利、细
小，像一根毛发一样的钩子。如果把它放到显微镜下面观察，就可以发现，在这把钩子
上有一条沟槽。如此而已，这件武器并没有什么其他更特别的地方。然而，这可是一件
有用的兵器，是可以致对手于死地的夺命宝刀。
这个小小的昆虫，正是利用这样一件兵器，在蜗牛的外膜上，不停地、反复地刺击。
但是，萤所表现出来的态度很平和，神情也很温和，并不像恶狠般凶猛，乍一看起来，
好像并不是猎人在捕猎食物，在咬它的俘虏，倒好像是两个动物在亲呢接吻一般。当小
孩子在一起互相戏弄对方的时候，他们常常用两个手指头，互相握住对方的皮肤，轻轻
地揉搓。这种动作，一般情况下，我们常用“扭”这个字眼儿来表示。因为，事实上，
这种动作近乎相互搔痒，而并不是那样重重地打。现在，就让我们来使用“扭”这个字
吧。一提到动物，除去那些最简单平实的字以外，通常在言语中使用的字，可以说，好
多都没有用处。那么，我们就可以说，萤是在“扭”动蜗牛，这大概更贴切一些。
萤在扭动蜗牛时，颇有它自己的方法。你会看到它一点儿也不着急，不慌不忙，很
有章法。它每扭动一下对方，总是要停下来一小会儿。仿佛是要审看一下，这一次扭动
产生了何等效果一般。萤扭动的次数并不是很多，顶多有五六次足以。就这么几下，就
能让蜗牛动弹不得，失去了一切知觉而不醒人世。再后来，也就是在萤开始吃战利品的
时候，还要再扭上几下。看起来，这几下扭动更至关重要。但是，至于究竟萤为什么要
如此这般行事，我就不能确定其真正的原因了。确实在最初的时候，只要轻轻地几下，
就足以使蜗牛慢慢地不能感觉一切了。那么，它为什么在食用前，还要来上几下呢？我
不得而知，至今仍是个迷。由于萤的动作非常迅速而敏捷，如同闪电一般快，就已经将
毒汁从沟槽中传送到蜗牛的身上了。这只是一个瞬间的动作，要非常仔细地观察才能觉
察到。
当然，有一点是不用怀疑的，那就是，在萤对蜗牛进行刺击时，蜗牛一点也不会感
觉到痛苦。关于这一点，我曾经做过一次小小的试验。在一只萤进攻一只蜗牛的时候，
当萤只扭了四五次以后，我马上迅速地把那只受了毒汁破害的蜗牛拿开。然后，用一根
很小很小的针去刺激这可怜虫的皮肤。那一点儿被刺伤了的肉，竟然一点儿也没有收缩
的迹象。这就已经很清楚地表明，此时此刻，这只蜗牛已经一点儿活气也没有了。它是
不会感觉到痛苦的。因为它已经沉浸到另一个世界里去了。还有一次，我非常偶然地看
到一只可怜的蜗牛遭受到萤的攻击。当时，这只蜗牛正在向前自由自在地爬行着。它的
足慢慢地蠕动着，触角也伸得很长。忽然，由于一下子的刺激和兴奋，这只蜗牛自己乱
动了几下。但是，马上这一切就静止了下来，它的足也不再向前慢爬了。整个身体的前
部也全然失去了它刚才的那种温文尔雅的曲线。它的长长的触角也变得软了，不再向上
伸展着了，而是拖垂到下边来，就像一根已经坏了的手杖一样，再也感受不到什么东西
了。从这种种表面的现象来看，这只蜗牛已经死了，已经真的到另一个世界里去了。
然而，实际上，这只蜗牛并没有真正悲惨地死去。我完全有办法，能让它重新活过
来，我可以给它第二次生命的机会。就在这个可怜的、假死的蜗牛既不生，又不死的两
三天内，我每天都坚持给它洗浴，清洁身体，特别是伤口。就在几天以后，奇迹出现了。
这只被萤无情地伤害得很惨重的、几乎一命呜呼了的家伙，恢复到了以前的状态中，它
已经能够自由地爬来爬去了。而且，它的知觉也已经恢复正常了。因为当我用小针刺击
它的肉时，它立刻就会有反应，小小的躯体马上就会缩到背壳里去藏了起来，这充分说
明它已经恢复知觉了，就像什么都不曾发生过一样。它是完全可以爬行了，那对长长的
触角重新又伸展开来，好像并没有发生过什么特别意外的事情一样。而且，它还精神倍
增。在它失去知觉的日日夜夜里，它仿佛进入了一种什么都不知晓的沉醉的状态，一切
都惊动不了它，而现在则大不一样了。它醒了，而且完全苏醒了过来，从死亡中复活了，
奇迹般地逃离魔爪，获得了第二次生命。
在人类社会的科学中，人们已经发明创造了在外科手术时不会让人感觉到疼痛的方
法，并且这种方法在医学实践上，已经被证明是非常成功的了。然而，在人类还没有找
到这种方法之前，萤以及很多其他种类的动物，就已经通过实践，实行了好几个世纪之
久了。所不同的是，外科医生在手术前，让我们闻以太或者是其他麻醉剂。而那些昆虫
采用的方法是，利用它们天生就长着的毒牙，向别的动物注射极小量的特别的毒药，以
达到让别的动物失去知觉的目的。
当我们偶然想到蜗牛具有那样温柔、平和而无害的天性，可萤却要采取向它注射毒
汁以麻醉它的特别方法来制服它，并且以它为食物，似乎总有些奇怪的感觉。但是，我
想我可以知道萤利用这种方法猎取蜗牛的种种鲜为人知的理由。
假如蜗牛只是在地上爬行，甚至是蜷缩在自己的壳子里，那么，对于它的种种攻击，
是一点儿也不困难的。原因是，蜗牛背上背的壳儿上并没有任何遮盖的东西加以保护。
而且，蜗牛身体的前部是完全暴露在外面的，也是毫无遮挡的。但是，实际上，事情并
不这么简单。蜗牛不仅仅单单地在上爬行，它经常是置身在较高的而且不稳定的地方。
比如，它喜欢趴在草杆的顶上，或者是待在很光滑的石面上。它贴身在这些地方，就无
需什么其他的保护了。因为，这些地点本身就为它提供了再好不过的天然保护。其原因
是，当蜗牛把自己的身体紧紧地依附在这些东西上时，这些东西就起到了盖子的作用。
于是，便能起到很好的保护作用了。不过，只要稍微一不留神，有一点儿没有遮盖严密，
一旦被萤发现，它的钩子可一点儿也不讲情面。只要有机可乘，它是一定要钻空子的。
总之，萤的钩子总会有办法可以触及到蜗牛的身体。然后，一下子钩住，释放出毒液，
蜗牛便会失去知觉。萤就可以安安稳稳地找个地方坐下来，享用它的美餐了。
不过，蜗牛深居高处，也是有一定危险的。当它爬在草杆上时，很容易掉下来。哪
怕稍微有一点儿扭动，或者是挣扎，蜗牛就会移动它的身体。一旦蜗牛落到地面上，那
么，萤就不得不去选择一个更好的猎物。所以，在萤捕捉蜗牛时，必须要使它没有丝毫
的痛苦感，失去知觉，让它动晃不得，从而也就不能轻易地逃跑了。因此，萤在进攻蜗
牛时，一定要采取触得很轻微的方法，以免惊动了这只蜗牛，让它从摇摆不定的草杆上
摇落下来。否则的话，可就前功尽弃，白费了一番心思，让到手的食物给溜走了。因此，
我想，萤之所以具备这样的外科器具，理由不过如此而已吧！除此以外，我再也想不到
其他的理由了。

二、蔷薇花饰物
萤不仅仅是在草木的枝干上结束战斗，使它的俘虏逐渐失去知觉，而且，也在这种
存在一定危险性的地方，就地把它解决处理掉，也就是要把它给完全吃掉。所以，萤的
食品的获得可并不是件简单容易的事呢！
那么，萤在吃蜗牛时，又是采用怎样奇特方法呢？它真的是在食用它吗？是不是要
先把蜗牛分割成一片一片的，或者是割成一些小碎片或碎粒什么的，然后再去慢慢地、
细细地咀嚼品味它呢？我猜想，它并不是以这样普通的方式食用它的。因为，我从来也
没有在这些动物体内，找到过任何这种小粒的食物。这就证明萤的吃，并不是通常的狭
义上的吃的意思，它只不过是以另一种方式，来解决问题罢了。具体方法是这样的，它
要将蜗牛先制造成非常稀薄的肉粥，然后才开始饮用。就像蝇吃小的幼虫一样，它能够
在还没有吃之前，先把它弄成流质，然后再痛快地享用。
更为具体的情形和做法是这样的。萤先使蜗牛失去知觉，无论蜗牛的身体大小如何。
在开始的时候，总是常常只有一只的四分之一大小。客人们也二三两两地跑过来了。它
们和主人毫无争吵，全部聚集到一起，准备和主人一起分享食物。过了两三天以后，如
果把蜗牛的身体翻转过来，把它的面孔朝下面放置，那样，它体内盛的东西，就会像锅
里的羹一样流出来。这个时候，萤的膳食已经结束了。它所饮下的只不过是一些其他动
物已经吃剩下的东西。因而，一只蜗牛被众虫同时分享了。
事实是很显然的。和前面我们已经看到过的“扭”的动作相似，它们经过几次轻轻
的咬，蜗牛的肉就已经变成了肉粥。然后，许多客人一起跑过来共同享用。很随意地，
每一位客人都一口一口地把它吃掉。而且，每一位客人都利用自己的一种消化素把它做
成汤。能够应用这样一种方法，说明萤的嘴是非常柔软的。萤在用毒牙给蜗牛注射毒药
的同时，也会注入——引进其他的无物质到蜗牛的体内，以便蜗牛身上固体的肉能够变
成流质。这样一来，这种流质很适合萤那柔软的嘴，使它吃得更加方便自如。
蜗牛被关闭在我的玻璃瓶里，虽然有的时候，它所处的地位不是特别稳固，但是，
它还是非常仔细小心的。有的时候，蜗牛爬到了瓶子的顶部，而那顶口是用玻璃片盖住
的。于是，它为了能在那里停留得更加稳固、踏实一些，它就利用那自己随身携带着的
粘性液体，粘在那个玻璃片上。这样一来，的确是非常稳定安全的。不过一定要多用一
些粘液，不然的话，哪怕稍微少用了一点儿粘液，都将是十分危险的。即便是微微地动
一点儿，也足以使它的壳脱离那个玻璃片，掉到瓶子底下去。
萤常常要利用一种爬行器——为了弥补它自己腿部，以及足部力量的不足——爬到
瓶子的顶部去，先仔细地观察一下蜗牛的动静，然后，作一下判断、和选择，寻找可以
下来的地方。然后，就那么迅速地轻轻一咬，就足以使对手失去知觉了。这一切都发生
在一瞬间。于是，一点儿也不拖延，萤开始抓紧时间来制造它的美味佳肴——肉粥，以
准备作为数日内的食品。
当萤一阵风卷残云以后，便吃得很饱了。剩下的蜗牛壳也就完全空了。但是，这个
空壳依然是粘在玻璃片上的，并没有脱落到瓶底上来。而且，壳的位置也一点儿都没有
改变，这都是粘液作用的结果。那个牺牲了的隐居者一点儿也不加以反抗，就这样静悄
悄地，不知不觉地任人宰割，最终，变成了别人嘴里营养丰富、美不胜收的大餐。就在
它那受到攻击的原地，逐渐流干了身体的全部，成了一个空空如也的壳儿。这种详细的
情形，向我们表明了这样一个事实，即萤的这种麻醉式的咬伤，是何等的有效。因此，
可以说，萤处理蜗牛的方法是十分的巧妙的。
萤要想顺利地完成自己的任务，实现自己的目的，比如，爬到悬在半空中的玻璃片
上去，或者是爬行到草杆上去，必须要具备一种特别的爬行足或其他什么有利的器官，
以便使它自己不至于在还未触及到猎物时，就先从高空跌落下来，从而半途而废。显然
它现有的笨拙的足是不够用的，这就决定它需要辅助的东西。
把一只萤放到放大镜下面进行仔细的观察研究，我们就可以很容易地发现，在萤的
身上，的确生长着这种特别的器官。大自然在创造它的时候很公平，非常细心，并没有
忘掉赐给它必要的工具。在萤的身体下面，接近它尾巴的地方，有一块白点，通过放大
镜可以清楚地看到。这主要是由一些一打以上的短小的细管，或者是指头组合而成的。
有的时候，这些东西合拢在一起形成为一团，而有的时候，它们则张开，成蔷薇花的形
状。就是这精细的结构，这些隆起来的指头，帮助了萤，使得它能够牢牢地吸伏在非常
光滑的表面上，与此同时，还可以帮助它向前爬行。如果萤想使自己紧紧地吸到玻璃片
上，或者是草杆上，那么，它就会放开那些指头，让蔷薇花绽放开来。在支撑物上，这
些指头放开得很大。萤就利用它自己自然的粘力而牢固地附着在那些它想停留的支撑物
体上。而且，当萤想在它所待的地方爬行时，它便让那些指头相互交错地一张一缩。这
样一来，萤就可以在看起来很危险的地方自由地爬行了。
那些长在萤身上的，构成蔷薇花形的指头，是不长节的，但是，它们每一个都可以
向各个不同的方向随意地转动。事实上，与其说它们像是指头，倒不如说它们更像一根
根细细的管子。因为，这一个比喻要更加合适，贴切一些。要是说它们像指头的话，它
们却并不能拿起什么东西。它们只能是利用其粘附力而附着在其他东西上。它们的作用
很大，除掉粘附，以及在危险处爬行这两大功能外，它还具有第三种功能，那就是它们
能当海绵以及刷子使用。在萤饱餐一顿以后，当它休息的时候，它便会利用这种自动的
小刷子，在头上、身上到处进行扫刷和清洁工作，这样既方便，又卫生，它之所以能够
如此自如地利用身体的这一器官，主要是因为那刺有着很好的柔韧性，使用起来相当便
利。在它饱餐之后，舒舒服服地休息一下，再用刷子一点一点，从身体的这一端刷到另
外一端，而且非常仔细、认真，几乎哪个部位都不会被遗漏掉。可以说，它是一种非常
爱清洁，注意文明修身的小动物。从它那副神采奕奕，得意洋洋，而舒服的表情来判断，
这个小动物对这个清理个人卫生的事情还是非常重视的，也非常的有兴趣去做的。刚一
开始的时候，我们当然会产生某种疑问：为什么这个小东西在拂拭自己的时候，是如此
的专心致志，而且如此地当心呢？答案是显而易见的。把一只蜗牛做成一顿肉粥，而且
花费了很多心思，用很多天的工夫去饮食它，肯定会把自己的身体弄得出奇的肮脏，那
么，认认真真地在饱餐之后，把自己的身体好好清洗一番，让自己焕然一新，是很有必
要的。

三、它的灯
如果，我们的萤，除了利用那种类似于接吻一样的动作——轻轻地扭动几下，来施
行麻醉术以外，就再也不具备什么其他的才能了，那么它的名声就不会有如此之大了，
以至于所有的人都知道它的大名。因此，它必定还具有一些其他的特殊本领，与其他动
物区别开来，比如特异功能什么的。那究竟它还有什么样的奇特本领呢？
众所周知，它的身上还带有一盏灯。它会在自己的身上点燃这盏灯。在黑夜中为自
己留一盏灯，照耀着自己行进的路程。这就是它成名的最重要的原因之一了。
雌性萤那个发光的器官，生长在它身体最后三节的地方。在前两节中的每一节下面
发出光来，形成了宽宽的节形。而位于第三节的发光部位比前两节要小得多，只是有两
个小小的点，发出的光亮可以从背面透射出来，因而在这个小昆虫的上下面都可以看得
见光。从这些宽带和小点上，发出的光是微微带蓝色的、很明亮的光。
而雄性的萤则不一样，它与雌萤相比，只有雌黄那些灯中的小灯，也就是说，只有
尾部最后一节处的两个小点。而这两个小点，在萤类的全族之中，差不多全都具备，从
萤还处于幼小的蛴螬的时代开始，就已经具备这两个用于发光的小点。此后，随着萤的
成长，它们也随着身体的生长不断地长大。在萤的一生中都不改变。这两个小点，经常
是无论在身体的上面，还是下面，都可以看见。但是雌萤所特有的那两条宽带子则不同，
它只能在下面发光的。这就是雄、雌的主要区别之一。
我曾经在我的显微镜下，观察过这两条发光的带子。在萤的皮上，有一种白颜色的
涂料，形成了很细很细的粒形物质。于是，光就是发源于这个地方。在这些物质的附近，
更是分布着一种非常奇特的器官，它们都有短干，上面还生长着很多细枝。这种枝干散
布在发光物体上面，有时还深入其中。
我很清楚地知道，光亮是产生于萤的呼吸器官的。世界上有一些物质，当它和空气
相混合以后，立即便会发出亮光，有的时候，甚至还会燃烧，产生火焰。此等物质，被
人们称为“可燃物”。而那种和空气相混合便能发光或者产生火焰的作用，则通常被人
们称之为“氧化作用”。萤能够发光，便是这种氧化作用的一个很好的例证与说明。萤
的灯就是氧化的结果。那种形如白色涂料的物质，就是经过氧化作用以后，剩下的余物。
氧化作用所需要的空气，是由连接着萤的呼吸器官的细细的小管提供的。至于那种发光
的物质的性质，至今尚无人知晓其答案。
但是，另外有一个问题，我们是知道得比较详细的。我们清楚地知道，萤完全有能
力调节它随身携带的亮光。也就是说，它可以随意地将自己身上的光放大一些，或者是
调暗一些，或者是干脆熄灭它。
那么，这个聪明的小动物，究竞是怎样行动才达到它调节自身光亮的目的呢？经过
观察我了解到，如果萤身上的细管里面流入的空气量增加了，那么它发出来的光亮度就
会变得更强一些；要是哪天萤不高兴了，把气管里面的空气的输送停止下来，那么，光
的亮度自然就会变得很微弱，甚至是熄灭了。
一些外界的刺激，将会对气管产生影响。这盏精致的小灯——萤身后最后一节上的
两个小点，哪怕只有一点点的侵扰，立刻就会熄灭。这一点我深有体会，每次当我想要
捕捉那些十分幼稚可爱的小动物的时候，它们总是爱和我玩捉迷藏的游戏。我明明就在
刚才，清清楚楚地看见它在草丛里发光，并且飞旋着，但是，只要我的脚步稍微有一点
儿不经意，发出一点儿声响，或者是我不知不觉地触动了旁边的一些枝条，那个光亮立
刻就会消失掉，这个昆虫自然也就不见了。我也就失去了捕捉对象，又浪费了一次机会。
然而，雌萤的光带，即便是受到了极大的惊吓与扰动，都不会产生多么大的影响。
比如说，把一个雌萤放在一个铁丝做的笼子里面，空气是完全可以流通的。然后，我们
再在铁笼子旁边放上一枪。就是这样爆烈的声音，竟然也毫无结果。萤似乎什么也没有
听到，或是听到了，也置之不理。它的光亮依然如故，丝毫变化都没有。于是，我又换
了一种方法试探。我取了一个树枝，而且还把冷水洒到它们的身上去，但是，这种种方
法都失败了。各种刺激居然都不奏效。没有一盏灯会熄灭，顶多是把光亮稍微停一下。
但是，这种情况是很少发生的。然后，我又拿了我的一个烟斗，往铁笼子里吹进一阵烟
去。这一吹，那光亮停止的时间长久了一些。还有一些竟然停熄掉了。但是，即刻之间
便又点着了。等到烟雾全部散去以后，那光亮便又像刚才一样明亮了。假如把它们拿在
手掌上，然后轻轻地一捏。只要你捏得不是特别的重，那么，它们的光亮并不会减少得
很多。总之，到目前为止，我们根本就没有什么办法，能让它们全体熄灭光亮。
从各个方面来看，毫无疑问，萤千真万确地能够控制并且调节它自己的发光器官，
随意地使它更明亮，或更微弱，或熄灭。不过，在某一种我们还不知道的环境之下，它
也会失去它这种自我调节的能力。如果我们从它发光的地方，割下一片皮来，把它放在
玻璃瓶或管子里面，虽然并没有像在活着的萤体上那么明亮耀眼，但是，它也还是能够
从容地发出亮光的。因为，对于发光的物质而言，是并不需要什么生命来支持的。原因
在于，能够发光的外皮，直接和空气相接触而起作用。因此，气管中氧气的流通也就不
必要了。就是在那种含有空气的水中，这层外皮发出的光也和在空气中发出的光同样明
亮。如果是在那种已经煮沸过的水里，由于空气已经被“驱逐”出来了，于是，发出的
光就会渐渐地熄灭了。再没有更好的证据来证明萤的光是氧化作用的结果了。
萤发出来的光，是白色而且平静的。另外，它的光对于人的眼睛一点儿也不刺激很
柔和。这种光看过以后，便会很自然地让人联想到，它们简直就像那种从月亮里面掉落
下来的一朵朵可爱的洁白的小花朵，充满诗情画意的温馨。虽然这种光亮十分灿烂，但
是同时它又是很微弱的。假使在黑暗之中，我们捉住一只细小的萤，然后把萤的光向一
行油印的字上照过去，于是我们便会很容易地辨别出一个一个的字母，甚至也可分辨出
不是很长的词来。不过，超过了这份光亮所涉及到的比较狭小的范围以外，那就什么都
看不清楚了。不过，这样的灯，这样吝啬的光亮，不久就会令读书人厌倦的。
但是，这些能够发出光亮的小动物，这些本该是心中一片光明的小昆虫，在事实上
却是一群心理很黑暗的家伙。它们对于整个家族的感情是完全不存在的。家庭对于它们
而言，是无足轻重的。柔情对于它们也是没有丝毫实际意义的。它们能够随处地产卵。
有的时候，产在地面上；有的时候，产在草叶上。无论何时何地，它都可以随意散播自
己的子孙后代。真可谓四处闯荡，四海为家，随遇而安。而且，在它们产下卵以后，就
再也不去注意它们了，随它们自生自灭，自然生长去了。
从生到死，萤总是放着亮光的。甚至连它的卵也是要发光的。幼虫也是如此。当寒
冷的气候马上就要降临的时候，幼虫就会立刻钻到地面下边去，但是并不钻得很深。假
如我从地面下，把它轻轻地掘起来，它的小灯仍然还是亮着的。就是在土壤的下面，它
的小灯还是点着的，永远为自己留一盏希望的灯！
------------------








被管虫



一、衣冠齐整的毛虫
当春天来临的时候，只要长着一对眼睛，就可以看清楚世界上任何东西的人，在破
旧的墙壁和尘土飞扬的大路上，或者是在那些空旷的土地上，都能够发现一种比较奇怪
的小东西。
那是一个小小的柴束，不知因为什么，它能自己自由自在地行动，一跳一跳地向前
走动。没有生命的东西变成了有生命的东西，不会活动的居然能够跳动了。这究竟是怎
么一回事呢？
这一点的确非常稀奇，而且很令人感到奇怪。不过如果我们靠近些仔细地看一看，
很快就能解开这个谜了。
在那些会动的柴束中，有一条特别漂亮、特别好看的毛虫。在它的身上装饰着白色
和黑色的条纹。大概它是正在寻找自己的食物。也许它是正在寻找一个可以让它安全地
化成蛾的适当的地点。对于它这些让人猜测不适的动作，通过以后它自己的所做所为就
能了解清楚了。
它很懦怯的朝前方急切地行走着，它总是穿着树枝做成的奇异的服装，完全把自己
的身体遮挡住了。只有除掉头和长有六只短足的前部暴露在外面。
它只要受到一点小小的惊动，就会本能地隐藏到这层壳里去，而且一动也不动了。
生怕一不小心被其他的东西侵害了，这显然是它的一种自我保护的本能。
这就是一束柴枝会走动的秘密的答案，它就是柴把毛虫，它是属于被管虫一类的。
为了防御气候的变化，这个既非常害怕寒冷而且又全身裸体的被管虫，建筑起了一
个属于它自己的很轻便的，又很舒服的隐避的场所，一个能够移动的安全的茅草屋。
在它还没有变成蛾的时候，一刻也不敢贸然离开这间茅屋。这确实要比那种装有轮
盘的草屋要好一些，它完全像是由一种特殊的材料制做而成的隐士们穿的保护外衣。
邓内白，山谷里的农夫，穿着一种用兰草带子紧紧地扎住的外衣，而且是羊皮的，
它是皮板朝里，羊皮朝外的。特别是居住在深山里的农夫，尤其是黄土高坡上的农夫，
这种穿着打扮更是常见，他们的头上还要系一条白色的羊肚毛巾。相比较而言，被管虫
的外衣，比这种打扮还要草率一些，因为它们只是拿一个简简单单的柴枝做成一件再朴
素不过的外衣而已，没有任何过多的装饰物品。可见，它们是多么不拘于小节啊！四月
里，在我们家的作坊上面有很多昆虫，在这些墙上能够发现很多的被管虫，它们都向我
提供了十分详尽的常识，如果它是在蛰伏的状态下，这就表示了它们不久就要变成蛾子
了。这是一个最好的机会，它使我能够直接地仔细地观察一下它的柴草的外衣。
这些外衣形状都是一个样子的，真的很像一个纺锤，大约有一寸半那么长。那位于
前端的细枝是固定的，而末端则是分散开的，它们就是这样排列着的，要是没有什么其
他更好的可以当作保护的地方，那么这里就是可以抵挡日光与雨水侵袭的避难场所了。
在并不认识它以前，乍一看上去，它真像一捆普通的草束。不过只是用草束这两个
字还并不能来正确地形容它的样子，因为麦茎实在是很少见得到的。
它的这件外衣的主要的材料是那些光滑的、柔韧的、富有木髓的小枝和小叶，其次
则是那些草叶和柏树的鳞片枝等，最后如果材料不够用了，就采用那些干叶的碎片和碎
枝。
总之，小毛虫遇到什么就使用什么，只要它是轻巧的、柔韧的、光滑的、干燥的、
大小适当的就可以了。所以，它的要求还不算苛刻。
它所利用的材料完全都是依照其原有的形状，一点儿都不加以改变。也就是说既保
持原有材料的性质，又保持原有材料的形状。
一些过长的材料，它也不修整一下，并使其成为适合的、适当的长度。造屋顶的板
条也直接被它拉过来使用。它的工作只不过是把前面固定了就行了。这在它是很简单易
行的。
因为要是想让旅行中的毛虫可以自由地行动，特别是在它装上新枝的时候，仍然能
够使它的头和足可以自如地活动，这个匣子的前部必须用一种特别的方法装置而成。仅
仅是用树枝装饰成的匣子对它而言是不适用的，理由很简单，因为它的枝特别长而且还
很硬实，这就大大防碍了这位勤劳的工人的工作，使它不能正常地尽职尽责。
它所需要的，是必须拥有一个柔软的前部，使得它可以向任何方向自由地转来转去，
从而可以很高兴地完成本职工作。
所以那些硬树枝，在离开毛虫前部相当远的地方，就中止了，取而代之的是一种领
圈，那里的丝带只是用一种碎木屑来衬托，这样一来，也就增加了材料的强度和韧性，
从而不妨碍毛虫的弯曲性。这样一个能够使毛虫自由行动和弯曲的领圈是非常重要，而
且绝对不可缺少的。以至于无论它的做法有怎样的不同，而所有的被管虫都要用到它。
在柴束前部，那张装得下可以自由转动的头部，触摸起来让人觉得很柔软，它的内
部是用纯丝织成的网，外面包裹着绒状的木屑。这种木屑，是毛虫在割碎那些干草的时
候得到的。
我把草匣的外层轻轻地剥掉，将它撕碎，就会发现里面有很多极细的枝干，我曾经
仔细地数过，大概有八十多个呢。在这里面，从靠近毛虫的这一端到那一端，我又发现
了同样的内衣，在把它的外衣打开以前，只有中部与前端是可以看见的，而现在则可以
看到全部了。这种内衣全都是由坚韧的丝做成的，这种丝的韧性是很强的，人用手拉都
不能把它拉断。这是一种光滑的组织，其内部是美丽的白颜色，外部是褐色的而且是有
皱纹的，还有细碎的木屑分散地装饰在上面。
于是，我要看看毛虫是如何做成这件精巧的外衣的了。这件外衣内外共有三层，它
们互相按一定次序叠加在一起。第一层是极细的绫子，它可以和毛虫的皮肤直接相接触；
第二层是粉碎的木屑，用来保护衣服上的丝，并使之坚韧；最后一层是小树枝做成的外
壳。
虽然各种被管虫全都穿上了这种三层的衣服，不过各个种族的外壳却有所不同。比
如，有一种，六月底我在靠近屋子旁边的尘土飞扬的大路上遇见的，它的壳无论从形式
还是从做法上来看，比前边提到过的那一种都要更加高明一些。它外面的厚披是用很多
片材料制做而成的，比如那种空心树杆的断片，细麦杆的小片，还有那些青草的碎叶等。
在壳的前部，简直找不到一点儿枯叶的痕迹。我先前所说的那一种，是常常有的，但那
足以妨碍其美观。在它的背部，也没有什么长的突出物，长出外皮之后，除去颈部的领
圈之外，这个毛虫的全身都武装在那个用细杆做成的壳里面。总体上的差别并不是很大，
不过最显著的一点差异就是它有比较完整的美观。
还有一种身材比较小，衣服穿得比较简易一些的被管虫，在冬天快要结束的时候，
在墙上或树上，在树皮多皱的老树上，比如洋橄榄树，或榆树上，常常可以发现它的踪
迹。当然在其他的地方也会见到。它的壳非常小，常常还不到一寸的五分之二长。它随
意地拾起一些干草，然后把它们平行地粘合起来，除去丝质的内壳以外，这就是它的全
身衣服的材料。
衣服要穿的更经济，更便宜，而且看上去更漂亮、更美丽，那是难度更大的事情了。

二、良母
如果我们在四月的时候捉几条幼小的被管虫，把它们放在铁丝罩子里面，关于它们
的一些事实，我们可以看得更多一点，也可以观察的更清楚一点了。
这时它们中的多数还是处在蛹的时代，等待着有朝一日变成蛾子。但是它们并不都
是那么安份守己地静静地待着，有的比较活跃好动一些，它们会很自豪地慢慢地爬到铁
丝格子上去。在那里，它们会用一种丝质的小垫子，把自身的身体固定好，无论是对它
们而言还是对我而言，都要耐心地等待几个星期，然后，才会有一些事情要发生。
到了六月底的时候，雄性的幼虫从它的壳里跑出来了，它已经不再是什么毛虫了，
而已经变成蛾子了。
这个壳，即一束细杆，你应当记得，它有两个出口，一个在前面，另一个在后面。
前面的一个，是这个毛虫很谨慎、很当心地制做的，是永远封闭着的，因为毛虫要利用
这一端钉在支持物上，以便使蛹得以固定在上面。因此，孵化的蛾必须从后面的口钻出
来，在毛虫还没有变化成蛾子之前，要在壳内先转一个身。然后，才会慢慢地出来。
虽然雄蛾只穿着一件十分简单的黄灰色的衣服，只有苍蝇差不多大小的翼翅，然而，
它却是异常漂亮的。它们长有羽毛状的触须，翼边还挂着细须头。
至于雌蛾，则很少能够在一些比较显眼暴露的地方捕捉到，而且，它们当中的大多
数是很难见得到的。
比别的昆虫迟几天以后，它才会从壳里姗姗来迟地钻出来，其形状简直是难看到了
极点，这个怪物也就是雌蛾。当你刚刚见到它的时候，甚至会惊吓地叫起来。也许它的
样子会吓你一跳的。没有一个人能够马上就看惯眼前这个凄惨的情形。它的难看程度并
不比那些毛虫差些。它没有长翅膀，什么都没有，在它背的中央，连毛也没有，光秃秃、
圆溜溜的。人们简直难得睁开眼睛看它一眼。在它圆圆的有装饰的体端，戴有一顶灰白
色的小帽子，第一节上，在背部的中央，长着一个大大的、长方形的黑斑点——这便是
它身体惟一的装饰物，母被管虫放弃了蛾类所有的一切美丽。这就是雌蛾，这个怪物般
的雌蛾形像。
当它离开它的蛹壳的时候，就在里面生卵。于是，母亲的茅屋(即她的大衣)就留传
给它的后代子孙了。她的卵产的很多，所以这产卵的时间也很长，要经过三十个小时以
上。
产完卵后，它将门关闭起来，使其免受外来的一些侵扰，从而获得一种安全感。为
了达到这个目的，某种填塞物是必要的。于是这位溺爱的母亲，在她一贫如洗、穷困潦
倒的情况下就只能利用它仅有的衣服了。也就是说要利用戴在它体端的那顶丝绒帽子，
塞住门口，以保母子平安，安然无恙地生活。
最后，它所做的还不限于此，它还要拿自己的身体来做屏障。经过一次激烈的震动
以后，它死在这个新屋的门前，留在那里慢慢地干掉，即在死后，它还依然留守在阵地，
为了下一代，死了也甘心。别看它外表上看起来丑陋不堪，但实际上它的内心、它的精
神是很伟大的。
假如破开外面的壳，我们可以看到那里面储存有蛹的外衣，除去前面蛾子钻出来的
地方留下的孔以外，一点也没有受到损坏，雄蛾要从这个狭小的隧道中出来的时候，会
感觉到它的翼和羽毛是很笨重的负担，而且对它形成了一定的阻力。
因此，当毛虫还处在蛹的时代时，就拼命地朝门口奔跑出将近一半的旅程来。最后
终于成功地撞出琥珀色的外衣来，在它的前面，出现了一块开阔的场所，可以允许它自
由地飞行了。
但是，母蛹不长翼，也不生羽毛，用不着经过这种艰难的步骤。
她的圆筒形的身体是裸露出来的，和毛虫没有多少区别。所以可以容许她在狭小的
隧道中爬出爬进，一点儿困难也没有。因此它把外衣抛弃在后面——抛在壳里面，作为
盖着茅草的屋顶。
同时，还有一种非常深谋远虑的举动，足以表现出她对于卵的命运有极其深切的关
心，事实上它们已经好像是被装在桶里面了，在它脱下的羊皮纸状的袋子里，母蛾已经
非常有办法地把卵产在里面了，直到把它装满为止，但是仅仅把它的房子与丝绒帽子遗
传给子孙，这并不能让它感到满足，最后的举动，还要把自己的皮也奉献出来留给子孙
后代，在它身上，“可怜天下父母心”这句话得到了最好的体现。
我想方便地观察这件事情的过程，于是我曾有一次从柴草的外壳里捡来一只装满卵
的蛹袋，并把它放在玻璃管中。在七月的第一个星期里，我忽然发现我竟然拥有了一个
被管虫的兴旺的大家族。它们孵化的速度是如此之快，差不多有四十多只以上的新生的
毛虫，竟在我没有看见的时候，在我还没来得及注意的时候，统统都穿上衣服了。
它们穿的衣服特别像波斯人戴的头巾，由光亮的白绒制做而成，说的普通一点，通
俗一些，就像一种白棉的礼帽，只是没有帽缨子。
不过说起来很奇怪，它们的这项帽子不是戴在头顶上的，而是从尾部一直披到前面
来的，它们在这玻璃管里非常得意地跑来跑去，因为这是属于它们自己的广大的屋子啊！
因此，我就想要看一看这顶帽子，究竟是由哪种材料做成的，织造的初步手续又是什么
样的。
幸运得很，蛹袋是不大会变空的。在里面，我又找到了它们第二个大家族，其数目
和先前跑出去的差不多。大概总有五打或六打的卵在里面。
我把那些已经穿好衣服的毛虫拿走，只留下这些裸露着身体的新客房在玻璃管里面，
它们有鲜红的头部，身体的其余部分全都是灰白色的，全身还不足一寸的二十五分之一
长。
我等待的时间并不长久，从第二天开始，这些小动物，慢慢地，成群结队地，开始
离开它们的蛹袋，用不着把这些摇篮弄破，只从它们母亲在当中弄破的口中出来就行了。
虽然它们都有洋葱头般的、漂亮的琥珀色，但是，没有一个把它拿来用做衣服的材
料，也没有一个利用那些柔软摇床的毛绒，谁都可能会以为这种材料可以做成这些伯冷
的动物的毛毯，但是事实上没有一个小动物去利用它一下。
它们一起冲到柴枝壳粗糙的外面，那是我故意为它们留下来的，而且直接靠近那个
装有卵的蛹袋，于是这些小动物们开始感觉到它们面临的情况有些不对头。于是便产生
了一种迫切感。
在你还未进入世界去打猎的时候，首先要做的是必须穿好自己的衣服，这一点对于
这些小动物们同样是适用的。它们也一样地焦急，恨不得马上攻破这个令人厌倦的陈旧
的老壳，赶紧穿上准备好的安全的外衣。
它们之中有的注意到了已经咬裂开的细枝，撕下那柔软的洁白的内层，有的很大胆，
深入到空茎的隧道，在黑暗中努力收集一些材料，它们的勇敢当然会有所报酬的，它们
得到了极其优等的材料，用这些织成雪白的衣服，还有一些毛虫加入了一些它们所选择
的东西，制做成了杂色的衣服，于是雪白的颜色给黑的微粒玷污了。
小毛虫制做衣服的工具就是它们的大头，其形状很像一把剪刀，并且它还长有五个
坚硬的利齿，这把剪刀的刀口靠得很紧凑。虽然它实际上很小，但它却很锋利，刀很快，
能夹住也能剪断各种纤维。
把它放在显微镜下可以清楚地观察到，小毛虫的这把剪刀竟然是有机械的、正确的，
而且是强有力的奇异标本。
如果羊也具备这样的工具的话，并且与它的身体成一定的比例，那么羊也就可以不
光吃草而也能吃树干了。由此可见，小毛虫的头可不能等闲视之啊！
观察这些被管虫的幼虫，制造棉花一样的灰白色的礼帽，这一点很能够启发人们的
智慧。无论是它们工作的行程，或是它们所应用的方法，都有很多的事情值得人们注意
一下。它们太微小了，它们也太纤弱了。当我用放大镜观察它时，必须非常小心，非常
仔细，既不敢使劲呼吸、喘粗气，也不敢大声说话，哪怕稍有一点不小心，就会惊扰了
它们，也许会把它们移动了位置，或者也可能是把它给一口气吹跑了。
别看这个小东西是如此微小，但是，它可是一位有着高超的制造毛毯技术的专家，
这个刚刚生下来一小会儿的小孤儿，竟然天生的知道怎样从它母亲留给它的旧衣服上裁
剪下自己的衣服来。它所采用的方法，我现在可以告诉人们，不过在此之前，我必须先
交代一点关于它的死去的母亲的事情。
我已经说过铺在蛹袋里的毛绒被，它很像一只鸭绒的床铺，软软乎乎，舒舒服服的，
小毛虫钻出卵以后，就睡在这张床上面休息一会儿，从而取得适当的温暖，并为到外面
的世界中去工作做好准备。
野鸭会脱下身上的绒毛，用它为子孙后代做成一张华丽舒适的床。母兔则会剪下身
上那些最柔软的毛，为它的新出生的儿女做成一张温暖的垫褥。母的被管虫也做着同样
的事情。看来，天下的母亲总还是有一定的共性的，这种共性也是它们的本能所决定的，
那就是无私地疼爱自己的儿女。
母亲会用一块柔软的充塞物，给小毛虫做成温暖的外衣，这材料非常的精细而且美
观。从显微镜下仔细地观察，可以看到上面有一点一点的鳞状片体，这就是它为小儿女
们制做衣服的最好的呢绒材料。小幼虫不久就会在壳里出现，因此要给它们准备好一个
温暖的屋子，让它们可以在里面自由地游戏玩耍。在它们还没有进入到广大的世界里去
之前，可以在里面修养，积蓄力量。所以母蛾像母兔、母鸭一样从身上取下毛来，为儿
女不辞辛劳地建造一片美好的天地。
这大概是以一种非常机械的方式进行的，好像是连续不断地磨擦墙壁而且并不是有
意识的有心的举动一样，然而的确没有理由向我们证实确是如此。甚至连最蠢笨的母亲
也有它自己的先见之明。这位看上去似乎有毛病的蛾子翻来覆去地打着滚，在狭窄的通
道中跑来跑去，想方设法地把自己身上的毛弄下来，给它的家族制做舒适的床铺。
有些书上说，小被管虫自从有了生命以后，就会吃掉它们的母亲。事实上，我却始
终也没有看到过这种情形发生。而且也不知道这个说法是怎样传说起来的。事实上，它
已经为它的家族奉献、牺牲了那么多，最后自己只留下干干的、薄薄的一个条，还不够
许多小子孙们的一口食物。实际上我的小被管虫们，它们是不吃母亲的。我看到的是它
们自从穿上衣服以后，一直到自己开始吃食的时候，没有一个曾经咬到自己的已死的母
亲的身上。

三、聪明的裁缝
现在我要详细地讲一讲这些小幼虫的衣服了。
卵的孵化是在七月初开始的，小幼虫的头部和身体的上部呈现出鲜明的黑色，下面
的两节，是带棕色的，其他部分都是灰灰的琥珀色。它们是一些十分精锐的小生物，跑
来跑去的脚步是很短小的，而且也是很快的。
它们从孵化地点的袋里钻出来以后，有一段时间，它们仍然需要待在从它们的母亲
身上取来的绒毛堆里。这里要比它们钻出来的那个袋子里更加空旷舒适一些。它们待在
绒毛堆里，有些在休息，有些十分忙乱，有的比较心急一些的已经开始练习行走了。它
们全体在离开外壳以前，都在修身养性，增强体质，以迎接未知世界风雨的洗礼。
在这个看上去比较奢华的地点，它们却并不留恋。等到它们的精力逐渐充沛起来，
就纷纷爬出来散布在壳上面。随后积极的工作就开始了，逐渐将自己穿着打扮起来。食
物问题以后才会想起来解决，目前却只有穿衣服是最要紧的事情，看来这些小家伙把脸
面上的事看得很重。
蒙坦穿上他父亲从前曾经穿过的衣服时，常常说：“我穿起我父亲的衣服了。”如
今，幼被管虫同样地穿起自己母亲的衣服(这同样必须记清，不是它身上的皮，而是它的
衣服)。它们从树枝的外壳，也就是我有时称做屋子，有时称做衣服的那种东西，剥取下
一些适当的材料，然后开始利用这些材料，给自己做衣服。它所用的材料都是小枝中的
木髓，特别是裂开的几枝，主要是因为它的髓是更容易取到的原故。
它们制做衣服的方法倒是非常值得注意的。这个小动物所采用的方法，真是出乎我
们人类的想象力，它是那样的灵巧，那样的细致，那样的精心，这种填塞物都被弄成极
其微小的圆球。那么这些小圆球是怎样连接在一起的呢？这位小裁缝需要一种支持物，
作为一个基础。而这个支持物又不能是从毛虫自己的身体上得来的。这个困难，并不能
难倒这些聪明的小家伙的；它们把小圆球聚集起来弄成一堆，然后依次用丝将它们一个
个绑起来。于是，困难就这样被克服了。你已经知道了，毛虫是能从自己身上吐出丝来
的，就像蜘蛛能吐丝织网一样。采用这种方法，把圆球或微粒连接在同一根丝上，做成
一种十分好看的花环，等到足够长了以后，这个花环就围绕在这个小动物的腰间，留出
六只脚，以便行动自由，末梢再用丝捆住，于是就形成了一根圈带，围绕在这个小幼虫
的身上。
这个圈带就是所有工作的起点和幼虫所需的支持物，完成第一道工序以后，小幼虫
再用大腮从壳上取下树心，固定上去，使它增长增大，于是就形成了一件完全的外衣。
这些碎树心或圆球，有时被放置在顶上，有时又被放在底下或旁边，不过通常都是放在
前边的时候居多。没有其他的设计，要比这个花环的做法更好了。外衣刚一做出来的时
候，是平的，后来把它扣住以后就像带子，圈在小毛虫的身体上。
最初工作的起点已经完成了，然后它会继续纺织下去。于是，那个最初的圈带逐渐
成为披肩、背心和短衫，后来成为长袍，几个小时以后，就完全变成一件雪白的崭新的
大衣了。
还要感谢它的母亲的关心，小幼虫得以免去光着身子到处跑来跑去的危险。假如它
不放弃那个旧的壳，那么，它们要想获得新的衣服将有很大的困难呀，因为草束和有心
髓的枝杆不是随处都可以找到的呢！然而，除非它们曝露而死，看来迟早它们总会找到
它们穿的衣服的，因为它们能利用随便什么材料，只要能找得到，什么都行。在玻璃管
中，我对于这些新生的小幼虫也曾做过好几回这样的试验。
从一种蒲公英的茎里，它毫不犹豫的挖出雪白的心髓，然后将它做成洁净的长袍子，
比它的母亲遗留给它的旧衣服所做成的要精致得多。有时还有更好的衣服，是用一种特
殊植物的心髓织造而成的。这一回的衣服上面饰有细点，像一粒粒的结晶块，或白糖的
颗粒。这可真正算是我们裁缝制作家的杰出作品了。
第二种材料，是我提供给它们的。那是一张吸墨纸，同样的，我的小幼虫也毫不犹
豫地割碎其表面，用它做成一件纸衣服，它们对这种新奇的材料非常高兴，也非常感兴
趣。当我再给它们提供那种原来的柴壳当做服装的材料时，它们竟然不予理睬，弃而不
顾，选取这种吸墨纸来继续做它们的衣服。
对于别的小幼虫，我什么东西也没有提供给它们，然而它们并没有就此失败。它们
非常聪明，采用了另一种方法，急急地去割碎那个瓶塞，使其成为小碎块，然后将这些
小碎块割成极其微小的颗粒，好像它们和它们的祖先也曾经利用过这种材料一样，因为
看上去这些小幼虫对这些材料并不陌生。这种稀奇的材料，也许毛虫们从来没有利用过，
然而它们把这些材料拿来做成衣服，竟然与其他材料做成的毫无差别。这些小幼虫的所
做所为真是让人感到惊奇！
从而我已经知道了它们能够接受干而轻的植物材料了，于是我决定换一种方法做试
验。用动物与矿物的材料来试试，我割下一片大孔雀蛾的翅膀，把两个裸体的小毛虫放
在上面。它们两个先是迟疑了好长时间。然后其中的一个就决心要利用这块奇怪的地毯，
一天的工夫都不到，它就穿起了它亲手用大孔雀蛾的鳞片做成的灰色的绒衣了。
第二回，我又拿来一些软的石块，其柔软的程度，只要轻轻一碰，就能破碎到如同
蝴蝶翼上的粉粒。在这种材料上，我放了四个需要衣服的毛虫。有一个很快就决定把自
己打扮起来，开始为自己缝制衣服。它的金属的衣服，有彩虹一样发出各种颜色的亮光，
闪烁在小毛虫的外壳上。这当然是很贵重，而且非常华丽的，只不过有点太笨重了。在
这样一个金属物的重压之下，小毛虫的行走变得非常辛苦，非常缓慢。不过，东罗马的
皇帝在国家有重大仪式的时候，也得如此呢！
为了满足本能上的迫切需要，幼小的毛虫也不顾忌这种蠢笨的行动了。穿衣服的需
要太迫切了，与其光着身子还不如纺织一些矿物好一些。爱美之心虫也有之，它也愿意
把自己打扮得漂漂亮亮的。吃的东西对于它并没有像穿的东西那样重要，只顾穿衣打扮，
外表好看，是这些小毛虫的共性与天性。假如先将它关起来两天，然后再换去它的衣服，
将它放在它喜欢吃的食物面前，比如一片山柳菊的叶子，它一定先做一件衣服，这是必
然的，因为一件衣服穿在身上后，它才会放心地去满足它的饮食需要。
它们对于衣服如此需要，并不是因为有特别寒冷的感觉，而是因为这种毛虫的先见。
别的毛虫在冬天，都是把自己隐藏在厚厚的树叶里，有的藏在地下的案穴里避寒，有的
在树枝的裂缝里，这是怕寒的毛虫。但是，我们所说的被管虫却安然地暴露在空气当中。
它不怕寒，也不怕冷，它从有生之日起，就学会了怎样预防冬季的寒冷。
受到秋天细雨的威胁以后，它又开始做外层的柴壳，开始时做得很草率、很不用心，
参差不齐的草茎和一片片的枯叶，混杂在一起，没有次序地缀在颈部后面的衬衣上，头
部必须仍然是柔软的，可让毛虫向任何方向自由转动。这些不整齐的第一批材料，并不
妨害建筑物后来的整齐。当这件长袍在前面增长起来的时候，那些材料便被甩到后边去
了。
经过一段时候以后，碎叶渐渐的加长，并且小毛虫也更细心地选择材料。各种材料
都被它直排的铺下去。它铺置草茎时的敏捷与精巧，真令人大吃一惊。人们不仅惊异地
发现小毛虫的动作如此之快，如此之轻巧，而且做的还很认真实在，铺垫的如此舒适，
这是一些大的昆虫都无法比拟的。真的不能小看它呀！
它将这些东西放在它的腮和脚之间，不停地搓卷，然后用下腮很紧的把它们含住，
在末端削去少许，立即贴在长袍的尾端。它的这种做法或许是要使丝线能粘的更坚固、
更结实些，和铅管工匠在铅管接合的尾梢锉去一点的意思是一个样的。
于是，在还没有放到背上以前，小毛虫用腮的力气，将草管竖起来，并且在空中舞
动它，吐丝口就立即开始工作，将它粘在适当的地方。于是，毛虫也不再摸索行动，也
不再移动，一切手续都已完成了。等到寒冷的气候来临的时候，保护自己的、温暖的外
壳已经作好了，所以，它可以安心地过自己的日子了。
不过这衣服内部的丝毡，并不很厚实，但能使它感到很舒服安逸。等到春天来临以
后，它可以利用闲暇的时间，加以改良，使它又厚又密，而且变得很柔软。就是我们拿
去它的外壳，它也不再重新制造了，它只管在衬衣上加上新层，甚至到不能再加为止。
这件长袍非常柔软，宽松而且多皱，又舒适、又美观。它既没有保护，也没有隐避之所，
然而它以为这并不要紧。做木工的时候已经一去不复返了，该是装饰室内的时候了，它
只一心一意地装饰它的室内，填充房子——即衬它的长袍，而房子已经没有了。它将要
凄惨地死去，被蚂蚁咬得粉碎，成为蚂蚁的一顿美餐。这就是本能过分顽固的结果呢！
------------------








樵叶蜂



如果你在园子里漫步，会发现丁香花或玫瑰花的叶子上，有一些精致的小洞，有圆
形的，也有椭圆形的，好像是被谁用巧妙的手法剪过了一般。有些叶子的洞实在太多了，
最后只剩下了叶脉。这是谁干的呢？它们又为什么要这样做呢？是因为好吃，还是好玩
呢？这些事都是樵叶蜂干的，它们用嘴巴作剪刀，靠眼睛和身体的转动，剪下了小叶片。
它们这么做，既不是觉得好吃，也不是为了好玩，而是这些剪下来的小叶片在它们的生
活中实在太重要了。它们把这许多小叶片凑成一个个针箍形的小袋，袋里可以储藏蜂蜜
和卵。每一个樵叶蜂的巢都有一打针箍形的袋，那些袋一个个地重叠在一起。
我们常看到的那种樵叶蜂是白色的，带着条纹。它常常寄居在蚯蚓的地道里，如果
你走到泥滩边，蹲下身子细细搜索，不难发现这样的地道。樵叶蜂并不利用地道的全部
作自己的居所，因为地道的深处往往又阴暗又潮湿，而且不适合排泄废物，有时还会遭
受昆虫的暗袭。所以它用靠近地面七八寸长的那段作自己的居所。
樵叶蜂一生中会碰到许多天敌，那地道毕竟不是一个安全坚固的防御工事，它用什
么办法来加强对自己家园的护卫呢？你瞧，那些剪下来的碎叶又派上大用场了。它用剪
下的许多零零碎碎的小叶片，把深处给堵塞了。这些用来堵塞的小叶片，都是樵叶蜂漫
不经心地从叶子上剪下来的。所以看上去非常零碎，不像筑巢用的那些小叶片一样整齐。
在樵叶蜂的防御工事之上有一叠小巢，大约有五六个。这些小巢由樵叶蜂所剪的小
叶片筑成，这些筑巢用的小叶片比那些做防卫工事的碎片，要求要高得多，它们必须是
大小相当、形状整齐的碎叶，圆形叶片用来作巢盖，椭圆形叶片用来做底和边缘。
樵叶蜂的小叶片，都是用它那把小刀——嘴角剪成的。为了适应巢的各部分的要求，
它往往用这把剪刀剪出大小不同的小叶片。对于巢的底部，它往往精心去设计，一点儿
也不含糊。如果一张较大的叶片不能完全吻合地道的截面的话，它会用两三张较小的椭
圆的叶片凑成一个巢底，一直到紧密地与地道截面吻合为止，决不留一点空隙。
做巢盖子的是一张正圆形的叶片。它好像是用圆规精确地规划过的，可以完美元缝
地盖在小巢上。
一连串的小巢做成后，樵叶蜂就着手剪许多大小不等的叶片，搓成一个栓塞把地道
塞好。
最值得我们思量的是，樵叶蜂没有任何可以用来当模子用的工具。它怎么能够剪下
这么多精确的叶子呢？它有可以依照的模型吗？还是它有什么特殊仪器可以测量呢？有
人推测，樵叶蜂的身体可以当作圆规来使用，一端固定住，即尾部固定在叶片某一点上；
另一端，也就是它们的头部，像圆规的脚一样在叶片上转动。这样就可以剪下一个标准
的圆。就像我们的手臂那样，固定肩肿挥动起来就是一个圆形。但是我们的手臂不会像
樵叶蜂那样巧妙又精确地画出大小一样的圆圈，这些用来做巢的盖子的圆叶片，恰好天
衣无缝地盖在巢上，非常完美。而小巢在地道的下面，它们不知道随时测量小巢的大小，
它们只靠摸索得到的感觉，来决定这只小巢所需要的叶盖大小。
圆形的叶片，不能剪得太大或太小。太大了盖不下，太小了会跌落在小巢内，使卵
活活闷死。你不用担心樵叶蜂的技术，它能很熟练地从叶子上剪下符合要求的叶片，虽
然没有什么模子，却是那么精确。樵叶蜂为什么有这么深厚的几何学基础呢？
冬天的一个晚上，我们都围着炉子坐着。我想起樵叶蜂剪叶片的事情，于是我设计
了一个小实验。
“明天是赶集的日子，你们中有人得出去采办这个星期所要用的东西。我们厨房里
一只天天要用的罐子的盖子被猫打破了。我要求他买一只盖子回来，不大不小，恰好能
盖上我们那个罐子。在去买之前，我们允许他仔细地把那罐子的大小估计一下，但不可
以用任何东西来测量，然后明天到集市上，凭记忆力选择合适的盖子。”大家听了都面
露难色，谁也不敢立即站出来承担这项任务。
的确，这似乎是一件很难的事。可是樵叶蜂的工作比我们这事更难以估测，它没有
看到自己的巢盖，根本没有这样一个印象；它也不能像我们选择盖子似的在摊贩的一大
堆罐中，靠着互相比较来选择一个最合适的盖子。对樵叶蜂来说，它必须在离家很远的
地方，毫不犹豫地剪下一片大圆叶，使它恰好能做巢的盖子。我们觉得很难的事，对它
来说像小孩游戏一样稀松平常。我们如果不用测量工具的话，比如绳子之类，或一个模
型或是一个图样，我们就很难选择一个大小适宜的盖子。可樵叶蜂什么都不需要。对于
如何治家，它们的确比我们聪明得多。
在实用几何学问题上，椎叶蜂的确胜过我们。当我看到樵叶蜂的巢和盖子，再观察
了其他昆虫在“科技”方面创造的奇迹——那些都不是我们的结构学所能解释的，我不
得不承认我们的科学还远不及它们。
------------------








采棉蜂和采脂蜂



我们知道，有许多蜜蜂像樵叶蜂一样自己不会筑巢，只会借居别的动物遗留或抛弃
的巢作自己的栖身之所。有的蜜蜂会借居泥匠蜂的故居，有的会借居于蚯蚓的地道中或
蜗牛的空壳里，有的会占据矿蜂曾经盘踞过的树枝，还有的会搬进掘地蜂曾经居住过的
砂坑。在这些借居它屋的蜜蜂中有一种采棉蜂，它的借居方式尤其奇特。它在芦枝上做
一个棉袋，这个棉袋便成了它的绝佳的睡袋；还有一种叫采脂蜂，它在蜗牛的空壳里塞
上树胶和树脂，经过一番装修，就可以当房间用了。
泥匠蜂很匆忙地用泥土筑成了“水泥巢”，就算大功告成了；木匠蜂在枯木上钻了
一个九英寸深的孔也开始心满意足地过日子了。尽管它们的家很粗糙，它们还是以采蜜
产卵为第一重要的大事，没有时间去精心装修它们的居室，屋子只要能够遮风挡雨就行
了；而另几类蜜蜂可算得上是装饰艺术大师，像樵叶蜂在蚯蚓的地道中做一串盖着叶片
的小巢，像采棉蜂在芦枝中做一个小小的精致的棉袋，使原来的地道和芦枝中别有一番
风情，令人不由得拍案叫绝。
看到那一个个洁白细致的小棉袋，我们可以知道采棉蜂是不适宜做掘土的工作的，
它们只能做这种装修工作。棉袋做得很长也很白，尤其是在没有灌入蜜糖的时候，看起
来像一件轻盈精致的艺术品。我想没有一个鸟巢可以像采棉蜂做的棉袋那样清洁、精巧
的。它是怎样把一个个棉花小球集中起来，拼成一个针箍形的袋子的呢？它也没有其他
特殊的工具，只有和泥匠蜂、樵叶蜂一样的灵巧的嘴，但它们的工作无论从方式上还是
从成果上看，都截然不同。
我们很难看清楚采棉蜂在芦枝内工作的情形，它们通常在毛蕊花、蓟花、鸢尾草上
采棉花，那些棉花早已没有水分了，所以将来不会出现难看的水痕。
它是这样工作的：它先停在植物的干枝上，用嘴巴撕去外表的皮，采到足够的棉花
后，用后足把棉花压到胸部，成为一个小球，等到小球有一粒豌豆那么大的时候，它再
把小球放到嘴里，衔着它飞走了。如果我们有耐心等待的话，将会看到它一次次地回到
同一棵植物上采棉，直到它的棉袋做完。
采棉蜂会把采到的棉花分成不同的等级，以适应袋中各个部分不同的需要。有一点
它们很像鸟类。鸟类为使自己的巢结实一些，会用硬硬的树枝卷成架子；又为了要使巢
温暖舒适些，而且宜于孵育小鸟，会用不同的羽毛填满巢的底部。采棉蜂也是这样做它
的巢，它用最细的棉絮衬在巢的内部，入口处用坚硬的树枝或叶片做“门”和“窗。”
我看不到采棉蜂在树枝上做巢的情形，但我却看到了它怎样做“塞子”，这个“塞
子”其实就是它的巢的“屋顶”。它用后足把棉花撕开并铺开，同时用嘴巴把棉花内的
硬块撕松，然后一层一层地叠起来，并用它的额头把它压结实。这是一种很粗的工作。
推想起来，它做别的部分的精细工作时，大概也是用这种办法。
有几只采棉蜂在做好屋顶后，怕不可靠，还要把树枝间的空隙填起来。它们利用了
所有能够得到的材料：小粒的沙土、一撮泥、几片木屑、一小块水泥，或是各种植物的
断枝碎屑。这巢的确是一个坚固的防御工事，任何敌人都无法攻进去。
采棉蜂藏在它巢内的蜂蜜是一种淡黄色的胶状颗粒，所以它们不会从棉袋里渗出来。
它的卵就产在这蜜上。不久，幼虫孵出来了。它们刚睁开眼睛，就发现食物早巳准备好
了，就把头钻进花蜜里，大口大口地吃着，吃得很香，也渐渐变得很肥。现在我们已经
可以不去照看它了。因为我们知道，不久它就会织起一个茧子，然后变成一只像它们母
亲那样的采棉蜂。
另外有一种蜜蜂，它们也是利用人家现成的房子，稍作改造变为自己的居住之处，
那就是采脂蜂。在矿石附近的石堆上，常常可以看到坐着吃各种硬壳果的蜗牛。它们吃
完后就跑了，石堆上留下一堆空壳。在这中间我们很可能找到几只塞着树脂的空壳，那
就是采脂蜂的巢了。竹蜂也利用蜗牛壳做巢，不过它们是用泥土做填充物的。
关于采脂蜂巢内的情形我们很难知道。因为它的巢总是做在蜗牛壳的螺旋的末端，
离壳口有很长的距离，从外面根本看不到里面的构造。我拿起一只壳照了照，看上去挺
透明的，也就是说这是只空壳，以后很可能被某个采脂蜂看中，在此安家落户，于是我
把它放回原处，让它作为将来的采脂蜂的巢。我又换一只照照，结果发现第二节是不透
明的，看来这里面一定有些东西。是什么呢？是下雨时冲进去的泥土？还是死了的蜗牛？
我不能确定。于是我在末端的壳上弄一个小洞，我看见了一层发亮的树脂，上面还嵌着
沙粒，一切都真相大白了：我得到的正是采脂蜂的巢。
采脂蜂往往在蜗牛壳中选择大小适宜的一节作它的巢。在大的壳中，它的巢就在壳
的末端。在小的壳中，它的巢就筑在靠进壳口的地方。它常常用细砂嵌在树胶上做成有
图案的薄膜。起初我也不知道这就是树胶。这是一种黄色半透明的东西，很脆，能溶解
在酒精中，燃烧的时候有烟，并且有一股强烈的树脂气味。你可以根据这些特点，判断
出来脂蜂用的是树干里流出来的树脂。
在用树脂和砂粒做成的盖子下，还有第二道防线，用砂粒、细枝等做的壁垒，这些
东西把壳的空隙都填得严严实实的。采棉蜂也有着类似的防御工程。不过，采脂蜂这种
工程只有在大的壳中才有，因为大的壳中空隙较多。在小的壳中，如果它的巢离入口处
不远，那它就用不着筑第二道防线了。
在第二道防线后面就是小房间了。在采脂蜂所选定的一节壳的末尾，共有两间小屋，
前屋较大，有一只雄蜂，后室较小，有一只雌蜂——采脂蜂的雄蜂比雌蜂要大。有一件
事科学家们至今仍无法解释，那就是母蜂怎能预先知道它所产的卵将是雌的还是雄的呢？
也就是它们怎么保证产在前屋的卵将来是只雄蜂，而产在后屋的卵一定会变成雌蜂呢？
有时候，采脂蜂筑巢的时候，一个小小的疏忽会造成下一代的一个大悲剧。让我们
来看看这只倒霉的采脂蜂吧！它选择了一只大的壳，把巢筑在壳的末端，但是从入口处
到巢的一段空间它忘记用废料来填充。前面我们提到过有一种竹蜂也是把巢筑在蜗牛壳
里的，它往往不知道这壳的底部已有了主人，一看到这个壳里还有一段空隙，就把巢筑
在这段空间里，并且用厚厚的泥土层把入口处封好。七月来了，悲剧就开始了。后面采
脂蜂巢里的蜂已经长大，它们咬破了胶膜，冲破了防线，想解放自己。可是，它们的通
路早已被一个陌生的家庭堵住了。它们试图通知那些邻居，让它们暂时让一让，可是无
论它们怎么闹，外屋那邻居始终没有动静。是不是它们故意装作听不见呢？不是的，竹
蜂的幼虫此时还正在孕育中，至少要到明年春天才能长成呢！难怪它们一直无动于衷。
采脂蜂无法冲破泥土的防线，一切都完了。它们只能让自己活活地饿死在洞里。这只能
怪那粗心的母亲，如果它们早能料到这一点，那么这悲剧也就不会发生了。如果那粗心
的母亲得知是自己活活杀死了孩子们，不知道该有多恨自己！不幸的遭遇并不能使采脂
蜂的后代学乖，事实上，时时有采脂蜂犯这样的错误，这与科学家所说的“动物不断地
从自己的错误和经验中学习和改进”的理论不符合。不过也难怪，你想，那些被关在壳
里的小蜂们永远地埋在了里面，没有一个能生还，这件事也随着小蜂们的死去而永远地
埋在了泥土里，成了无人能知的千古奇冤，更不用说让采脂蜂的后代吸取教训了。
------------------








西班牙犀头的自制



我希望你还记得神圣甲虫，它消耗掉它的时间，做成即可以当食物，又可以当梨形
窝巢的基础的圆球。
我已经指出，这种形状对于小甲虫的利处和害处，因为圆形是顶好的形状，可以保
存好食物使其不干也不硬。
经过长时间的观察这种甲虫的工作，我开始怀疑我极力赞扬它的本能，或许是我估
计错误了。它们是否真的关心它们的小幼虫，并且替它们预备下最柔软最合适的食物呢？
甲虫做球是它们自己的职业啊！它要继续在地底做球不是很奇怪吗？一个动物生着长而
弯的腿，用它把球在地上滚来滚去是很便利的。无论在那里，自然要从事自己所喜欢的
职业。自己想干的工作，就一定要干好，只有这样才能在自然界中求生存，才能在大自
然中繁衍后代，一代一代地生存下去。
它并不顾及它自己的幼虫，或许它做成梨形的外壳这件事仅仅是碰巧了而已。
为了要圆满地解决这个疑难问题，我还观察过一种清道的甲虫，在它的日常工作中，
它非常不熟悉做球这种工作。可是，到了产卵期，它突然改变了以往的习惯，将自己储
存的所有食物都统统做成圆圆的一个团。这一点表明这不仅仅是习惯而已，而是真的关
心它的幼虫，因而选择圆形的球做为它的窠巢。
如今，在我的住所附近，就有这样一种甲虫。它是甲虫中最漂亮的，个子最大的。
虽然不如神圣甲虫那么魁伟，它的名字就是——西班牙犀头。
它最显著、最特别的地方，就是它胸部的陡坡和头上长的角。
这种甲虫是圆的，而且很短，当然也就不适合做神圣甲虫所做的那些运动。它的腿
不足以供做球使用。稍有一点点惊扰，它的腿就本能地卷缩在自己身体的下面，它不像
一个勇敢者，也不像神圣甲虫那样，有一个勇敢者的气魄。
它们一点也不像搓滚弹丸的工具，它们那种发育不全的形象，表明它们缺乏挖掘性，
这足以使我们清楚它是不能带着一个滚动的圆球走路的。
的确，犀头的性格很不活泼。有一次，在夜里，或在黄昏的月光下，它寻找到食物，
就在原来的地点挖开一个洞穴。它的这种挖掘草率的很，其最大的也只能藏下一个苹果。
在这里，它逐渐堆下刚刚才找来的食品和食料，至少一直要堆积到洞穴的门口。
它的大量的食物要堆积为不成形的一大堆，这就足以证明这个犀头的贪食、贪吃和
馋嘴了。食物能够吃多长的时间，它自身也就在这地底下待多长时间，一直待到吃完所
存的食物为止。
等它把所有存储的食物全都吃完以后，它的食品仓库空了，它这才又重新跑出来，
再去寻找新鲜的食物，然后再另挖掘一个洞穴，重复它那种存了吃，吃了再出来找的周
期性运动。
实实在在地说，它只不过是一个清道夫，是一个肥料的收集者而已。总之，它没有
什么特别的本事，是一个平庸之虫。
对于搓捏圆球的技术，它明显表现出特别的外行。而且，它的短而笨的腿，也极其
的不适合干这种技术性的工作。
在五六月之间，产卵的时候到了，这个昆虫则变成了非常擅长于选择最柔软的材料，
选择最舒适的环境，为它顺利产卵打下一个良好环境的能手了。
它开始为它的家族制做食物，只要在一个地方找到，如果它认为是最好的，它立刻
就把它们埋在地下，它从不旅行，从不搬运，从不做任何添加配制工作，也从不进行再
加工。
然而，我看到这个洞穴，比它自己吃食的临时的洞穴，挖掘得更宽大一些，而且建
筑得也比较精细。
我觉得在这种野外的环境里，要想仔细观察犀头的一些生活习惯，以及它的生长过
程，是非常不容易的，所以后来我就将它放到我的昆虫屋里面，这样，我可以更加认真、
更加仔细地观察。这为我自己提供了许多的方便。
起初，这个可怜的昆虫，因为被我俘虏了，所以有一些胆怯，它可能认为大难即将
来临。当它做好了洞穴以后，自己出入洞穴时，也还是提心吊胆，唯恐自己被再次伤害
了一样。然而从这以后，它也就逐渐的胆壮起来，在一夜之间，将我提供给它的食物全
部储存起来了。
在一个星期快要过去的时候，我掘起昆虫屋中的泥土。我发现，我见过的它储存食
物的洞穴显现出来了，这是一个很大的厅堂，一个很大的仓库。它的屋顶并不很整齐，
四壁也是很普通的，地板差不多是平平坦坦的。
在一个角上，找个圆孔，从这里一直通往倾斜的走廊，这个走廊一直通到土面上。
这个房子——这个昆虫的别墅——用新鲜的泥土掘成的一个大洞。它的墙壁，曾经被很
仔细的压过，很认真地装饰过。这也就足以抵抗我在做试验时所引起的地震了。并且很
容易就能看到这个昆虫以及它所有的技能，它不遗余力，用尽所有的掘地力量，来做一
个永久的家。可是它的餐室却仅仅是一个土穴，墙壁做的也不那么坚固。
当它从事这个大型建筑的建设的时候，我想，它的丈夫，或者是它的伴侣一定会来
帮助它的，至少我常常看见它和它的丈夫一同待在一个洞穴里。我也相信这个帮助会使
它的妻子更加勤快，丈夫和妻子可以一起收集并储存食物。因为夫妻二人同做一件事情，
同干一件工作，自然要快得多，至少比一个人干事要快得多，但是等到屋子里储备满了，
足够它生活以后，它的丈夫也就隐退了。这位丈夫就跑回到土面上来，到别的地方去安
身了。它对这个家庭应做的工作，应尽的职责也就全部结束了，尽到了一个丈夫应尽的
职责，就此结束了对这个家庭的义务。
那么，在许多食物放下去的土屋中，我所看到的是什么样的呢？是一大堆小土块，
互相堆叠在一起吗？但是，一点儿都不对。实际上不是想象中的样子。我只看到单独的
一个很大的土块，除掉一条小路以外，储存食物的那一个屋子，全都被塞满了。
这种圆堆块没有一定的形状，有的大小像吐绶鸡的蛋，有的像普通的洋葱头。有的
是差不多是完整的圆形。这使我想起了荷兰的那种圆形硬酪。有的是圆形而上部微微有
点突起。然而，无论是哪一种，其表面都是很光滑的，呈现出精致的曲线。
这位母亲，不辞辛苦地一次一次地带去很多很多的材料，收集在一起并搓成一个大
团。它的做法是，捣碎这许多的小堆，将它们合在一起，并把它们揉合起来，同时也踩
踏它们。有好几回我都曾经见到它在这个巨大的球顶上。当然，这个球要比神圣的甲虫
做的那个大得多，两个互相比较一下，后者只不过是个小小的弹丸而已。它也有时在约
四寸直径的凸面上徘徊，它敲它、拍它、打它、揉它、含它，使它变得坚固而且平坦。
我只有一次见过如此新奇的景观，而且只有一次。这是多么难得的机会啊！但是当它一
见到我的时候，立刻就滚到弯曲的斜坡下不见了。它发现，它的所做所为已被人注意到
了，完全暴露身份和目标，所以它就逃之大吉了。
我得力于一排墨纸盖住的玻璃瓶，在这里我发现了许许多多有趣的事情。
第一我发现了这个大球的雕饰过程——常常是很整齐的，无论其倾斜程度的差异如
何——这并不是由于搓滚的方法而形成的。
事实上我已经知道，这么大的体积决不能滚进这个差不多已经被塞满了的洞里去。
而且这个昆虫的力量也不足以移动这么大堆的东西。
我每次到瓶边观察时，所得到的证据都是一个样子的，我常常看到母虫爬到球顶上，
看看这里、又看看那里，看看这边、又看看那边，它轻轻地敲，轻轻地拍，尽量使之光
滑，似乎没有见过它有想移动这个球的意思。
事实明确地证实，制球是并不采用搓滚的方法的。
最后已经准备好了。这就像面包工人将面粉团分成许多许多的小块，每一块将来都
将成为面包。这犀头甲虫也是一样的做法。它用头部锋利的边缘，及前爪的利齿，划开
圆形的裂口，从大块上随意割下小小的一块来。在做这次工作的时候，一点犹豫都没有，
也不重复改做一下。它从不在这里加上一点，或者在那里去掉一点。直接了当，只要一
次切割，它就得到适当的一块了。
其次，就是如何使球有一定的形状。它竭力将球抱在那双短臂之间，叫人看起来它
很不适于做这项工作的，只用压力把它做成圆块。它很庄严、很正重地，在不成形的一
块食物上爬上爬下，向左爬，向右爬，向前爬，向后爬，不停地爬，耐心地一再触摸，
最后经过二十四小时以上的工作，终于有棱有角的东西变圆了，像成熟的梅子一样大小。
在它狭小的技术操作室里，简直就没有什么余地可以自由地转动一下了。这位又矮
又胖的艺术家完成此项工作，竞然会没有动摇它的底面。但是经过相当的时间与耐力以
后，它竟然做成了确实适当的圆球。从如此笨拙的工具与有限的地位而论，这看来似乎
是不可能的。
它亲切有味地用足摩擦圆球的表面，再经过很长的时间，最后它终于满意了。然后，
它爬到圆顶上面，慢慢地压，压出一个浅浅的穴来，就在这个盆样的孔穴里它产下一个
卵。
于是，它非常当心，非常精细地把这个盒子的边缘合拢起来，以遮盖它产下的那个
卵，再把边缘挤向顶上，使之略略尖细而突出。最后，这个球就做成椭圆形的了。
这个昆虫于是又开始从事第二个小块的工作，制造的方法完全相同。余下的，又重
新做第三个乃至第四个，你当然记得，神圣甲虫用很熟悉的方法只做一个梨形的巢。
它的洞穴中隐藏着三四个蛋形的球，一个紧靠着一个，而且组合、排列都很有规则，
细小的一端全都朝着上面。
它经过长期的工作以后，谁都要以为它也像神圣的甲虫一样，跑出来寻找自己的食
物去了。然而事实却不是这样的，它没有那样做。它没有跑出去，也没有去寻找食物，
而是在那里一动不动地守着，并且自打它钻入地下以后，它一点食物也没有吃过，它像
宇宙间任何母亲一样，一样的无私，一样的奉献。对自己的子女只有爱护、关怀与牺牲。
它没有，而且也不肯去碰一碰那为自己的子女预备下的食物。它宁愿自己挨饿，宁
愿自己受痛苦，也不愿意自己的小幼虫将来感受到一点儿痛苦，这是多么了不起的奉献
精神啊：在昆虫的世界里也充分体现了母爱是最伟大的。
它不出去的目的，当然是为了看守这几个为子女建筑下的摇篮。因为这是这个家族
生存的基本条件之一。这是它们的房子，是它们的小别墅，是它们生活在世上的惟一栖
身的地方。因而要仔细地看护它。
神圣的甲虫的梨正是因为母亲的离开，而遭到损坏的，当母亲离开不久，梨就已破
裂开了。经过一个相当长的时期以后，就不成形状了，就这样，一个家被毁掉了。
但是这个甲虫的蛋，可以保存完好，并长时间地保存，因为它有母亲的关心爱护，
母亲的一份责任感，才使它们的蛋完好地保存下来。
它从这一个跑到那一个上，再从那一个跑到另一个上，看看它们，听听它们，唯恐
它们有什么闪失，受到了什么外来的侵害。就像人类母亲对自己怀里的婴儿一样，关怀
得无微不至。这小甲虫真是一个好母亲。
它修补这一处，然后又修补那一处，生怕它的小幼虫受到什么干扰，受到外来的欺
辱。我们的眼睛看不出什么不足的地方，它虽然很笨拙而且有角，有足，但是在黑暗中
竟然比我们的视觉在日光中还要灵敏，还要看的清楚，这一点我们可以感觉的到。只要
有细微的破裂，它立刻就会跑过去，赶紧地修补一下，惟恐空气会透进去，干掉它的卵。
它在摇篮当中狭窄的过道里跑出跑进，为的是保护它的卵，它仔细观察，认真巡视，
假如我们打扰它，破坏它正常的生活，它就立刻用体尖抵住翼尖壳的边缘，做出柔软的
沙沙之声，如同和平的鸣声，又像发出强烈的抗议一般。
它就是这样，辛辛苦苦地关心着它的摇篮。有时候它实在困了，也会在旁边睡上一
小会儿，但时间不会太长的，只是打一会儿盹而已，决不会高枕无忧的睡上一大觉。这
位母亲就是这样在看守它的卵，为它的后代做出无私的奉献，为儿女操碎了一颗心。
犀头在地下室中，有着一个昆虫所稀有的特点，那就是照顾自己家庭的快乐。这是
多么伟大的母爱呀！这是一个奉献者的自豪。
它在自己弄下的缺口处。听见它的幼虫在壳内爬动，争取自由。当这个小囚犯，伸
直了腿，弯曲了腰，想推开压在自己头上的天花板时，它的母亲会意识到，小幼虫一天
天长大了，要独立生活了，该自己去世界上闯荡一番了。这位小幼虫自己出来，感受自
由与生命的美好。
即然有建造修理的本领，为什么不能打碎它呢？然而我不能做出肯定的回答，因为
我没有见到过这种事情发生。或许可以说这个母虫，被关在无法逃脱的玻璃瓶子里，所
以它一直守在巢中，因为它没有任何行动的自由。不过，假使如此，它对摩擦工作与长
时间的观察难道不感到焦急吗？这个工作显然对于它很自然，形成了它已经习惯了的一
部分生活了。
假如它急切的想恢复自由，它当然要在瓶中爬上爬下，毫无休止地忙碌。但是，我
只看见它常常是很平静的，也很安心地待在它的圆球旁。
为了要得到确切的第一手资料，为了得到确切的真象，所以我随时去察看玻璃瓶中
的现象如何。
如果它要休息，它可以任意的钻入沙土中，到处都可以隐藏它的身体，如果需要饮
食，也可以出来取得新鲜食物，然而既不是休息，也不是日光与饮食，可以便它离开它
自己的家族片刻。它只坐镇在那里，直到最后一个圆球破裂开，我常见它总是坐在摇篮
旁边的，那分安静，那分重担在肩的责任感很让我感动。
大概有四个月的时间，它不吃任何食物，它已不像最初为了照顾家族时那么贪嘴了，
而这时它竟然对于长时间的坐守，有非常惊人的自制力了。
母鸡伏在它的蛋上，忘记饮食数星期以后，自己的蛋才变成小鸡，然而犀头却要忘
记饮食达到一年的三分之一那么久。
夏天过去了。人类和牲畜都很希望下几场雨，终于下来了，地上积了很深的水。
于是，在我们布罗温司酷热干燥、生命不安的夏季过后，我们有凉爽的气候来使它
复活了。
石南开放了它的红色钟形的花，海葱绽放穗状的花朵，草莓树的珊瑚色果子也已经
开始变软了，神圣甲虫和犀头也裂开外层的包壳，跑到地面上来，享受一下一年来最后
这几天的好天气了。
刚刚解放出来的犀头家族，与它们的母亲一起，逐渐地来到地面。大概有三四个，
最多的是五个。
公的犀头生有比较长的角，很容易就能分辨出来。
母的犀头与母亲则很难分别。因此它们之间，很容易混淆。
不久，又有一种突然的改变发生了。从前牺牲一切的母亲，现在对于家族的利益，
已不再那么关心了。
自此它们各自开始管理自己的家和自己的利益了。它们彼此之间也就不相互照应了。
目前虽然母甲虫对家族漠不关心，但我们都不能因此而忘记它四个月来辛辛苦苦的
看护，除掉蜜蜂、黄蜂、蚂蚁等外来的干涉和侵犯。自己能养儿育女，关心它们的健康，
直到长成之后，据我所知，再没有别的昆虫能够做到这些了。
它独自一个毫无外来帮助，为每个孩子预备摇篮似的食物，并且尽心修补，以防止
其破裂，使摇篮十分安全。这是一个母亲无私的奉献。
它的情感如此的浓厚与执着，使它失掉了一切的欲望和饮食的需要。
在洞穴的黑暗里看护它的骨肉达到四个月之久。细心地看护着它的卵。
它在的子女们未得到解放出来之前，它决不恢复户外的快乐生活。
我们竟从田野中愚蠢的清道夫身上，看到最深切的关于母性本能的例子，不禁对这
种小昆虫产生了无限的敬意。
------------------








两种稀奇的蚱蜢



一、恩布沙
海是生物最初出现的地方，至今还存在许多种奇形怪状的动物，让人们无法统计出
它们的具体数目，也分不清它们的具体种类。这些动物界原始的模型，保存在海洋的深
处。这就是我们常说的，海洋是人类无价的宝库，它是人类生存的重要条件之一。
但是，在陆地上，从前的奇形动物，差不多都已经灭绝了，只有少数的还遗留下来，
能留到现在的大多都是一些昆虫类的动物。其中之一就是那种祈祷的螳螂，关于它特有
的形状和习性，我已经在前文对你们说过了。别一种则是恩布沙。
这种昆虫，在它的幼虫时代，大概要算布罗温司省内最怪的动物了。它是一种细长，
摇摆不定的奇形的昆虫。它的形状和任何昆虫都不一样，没有看惯的人，决不敢用手指
去碰触它。我的近邻的小孩，看了这个奇怪的昆虫以后，看到它这个奇异的模样，留下
了很深的印象，他们叫它为“小鬼”。他们想象它和妖法魔鬼等等多少有些关系。从春
季到五月，或是到秋天，有时在有阳光和温暖的冬天，可以遇见它们，虽然从不集成大
群。
荒地上坚韧的草丛，可以受到日光照耀，并且有石头可以遮风的矮丛树，都是畏寒
的恩布沙最喜欢的住宅。
我要尽我一切的可能告诉你们，它看起来像什么样子。它身体的尾部常常向背上卷
起，曲向背上，形成一个钩的形状，身体的下面，即钩的上面，铺垫着许多叶状的鳞片，
并排列成三行。
这个钩架在四只长而细的，形如高跷的腿上；每只足的大腿和小腿连接之处，有一
个弯的、突出的刀片，这个刀片与屠夫切肉常用的那种刀片相仿。
在高跷或四足蹬上的钩的前面，有很长而且很直的胸部突起。形状圆而且很细，像
一根草一样，草干的末梢，有猎狩的工具，是完全类似螳螂的那种猎具。
这里有比较尖利的鱼叉，还有一个残酷的老虎钳，生长着如锯子似的牙齿。上臂做
成的钳口中间有一道沟，两边各有五只长长的钉，当中也有小锯齿。臂做成的钳口也有
同样的沟，但是锯齿比较细巧，比较密一些，而且很整齐。
在它休息的时候，前臂的锯齿嵌在上臂的沟里。它的整体就像一架可以加工的机器，
有锯齿、有老虎钳、有沟、有道，如果这部机器再稍微大一点，那它就成了一部令人可
畏的刑具了。
它的头部也和这种机器相辅相承。这是一个多么怪异的头啊！尖形的面孔，卷曲而
长的胡须，巨大而且突出的眼睛，在它们中间还有短剑的锋口；在前额，有一种从未见
过的东西——一种高的僧帽一样的东西，一种向前突出的精美的头饰，向左向右分开，
形成尖起的翅膀。
为什么这个“小鬼”要这样像古代占卜家一样戴着奇形怪状的尖帽子呢？它的用途
在不久以后我们就会知道的。
在这个时候，这动物的颜色是普通的，大抵为灰色，待发育以后，就会变为装饰着
灰绿、白、与粉红的条纹。
如果你在丛林中遇见这个奇怪的东西，它在四只长足上动荡，头部向着你不停地摇
摆，转动它的僧帽，凝视着你的眉头。
在它的尖脸上，你似乎可以看到要遭受危险的形象。但是，如果你想要捉到它，这
种恐吓姿势，马上就会不见了。
它高举的胸部就会低下去，竭力用大步逃之夭夭，并且它的武器会帮助它握着小树
枝。假如你有比较熟练的眼光，它就很容易被捉住，关在铁丝笼子里。
起初，我不知道应该如何喂养它们。我的“小鬼”又很小，最多只有一两个月大。
我捉大小适宜的蝗虫给它们吃，我选取了其中最小的一些喂给它吃。
“小鬼”不但不要它们，而且还惧怕它们，无论那个无思想的蝗虫怎样很温和地靠
近它，都会受到很坏的待遇。
尖帽子低下来，愤怒的一捅，使蝗虫滚跌开去。
因此可知，这个魔术家的帽子实际上是自卫的武器。雄羊用它的前额来冲撞，和它
的对手进行搏斗，同样的，思布沙也在用它的僧帽来和它的对手进行抵抗。
第二次，我喂给它一个活的苍蝇，这个恩布沙立即就接受了它，把它当成一次酒席
上的佳肴。当苍蝇走近它的时候，早己守候着的恩布沙掉转了它的头，弯曲了胸部，给
苍蝇猛然一叉，把它夹在两条锯子之间。就连老猫扑捉老鼠也没有这样的迅速。
我惊奇地发现，一只苍蝇不仅可供给它一餐，而且足够整日食用，甚至可以连着吃
上几天。这种相貌凶恶的昆虫，竞然是极其节食的动物。
我开始以为它们是一个个的魔鬼，但是，后来发现它们的食量像病人一样少。经过
一个时期以后就连小蝇也不能引诱它们了。在冬天的几个月里，它完全是断食的。到了
春天，才又准备吃一些小量的米蝶和蝗虫。它们总在颈部攻击俘虏，如螳螂一般。
幼小的恩布沙，被关在笼子里时，有一种非常特殊的习性。
在铁丝笼里，它的态度从最初一直到最后，都是一样的，而且是一种顶奇怪的态度。
它用它那四只后足的爪，紧握着铁丝倒悬着，纹丝不动，活像一只倒挂在横杠上的小金
丝猴一样，它的背部向下，整个的身体就挂在那四个点上。如果它想移动一下，前面的
鱼叉就会张开，向外伸展开去，然后，紧握住另一根铁丝，朝怀里拉过来。
用这种方法将这个昆虫在铁丝上拽动，仍然是背朝下的，于是鱼叉两口合拢，缩回
来放在胸前。
这种倒悬的位置，对于我们而言一定会很难受的，也是很不容易做到的，要是人很
可能就会得病的，要么是高血压，要么是脑出血。但是，恩布沙保持这样的姿势的时间
并不短，它在铁丝笼里，可以持续十个月以上，竟然毫无改变。
苍蝇在天花板上，确实也是这种姿势的，但是它有休息的时间，它累了就要休息一
会，养足了精神以后，再做这种动作。它在空中飞动，用平常的习惯走路，沐浴在阳光
中。
恩布沙则完全相反，它保持这种奇怪的姿势，达到十个月以上，绝不休息。它悬挂
在铁丝网上，背部朝下，猎取、吃食、消化、睡眠、经过昆虫生活所有的经历，直至最
后死亡。它爬上去时年纪还很轻，而落下来的时候，已经是年老的尸首了。
它这个习惯的动作，应该注意的是只有处在俘囚期的时候才会如此，并不是这种昆
虫天生的、固有的习惯。因为在户外，除去很少的时候，它站在草上时是背脊向上的，
并不是倒悬着的。
和这种行为相似的，我还知道另外一个稀奇的例子，比起这个还要特别一些。这就
是一种黄蜂和蜜蜂，在夜晚休息时的姿态。有一种特别的黄蜂——生有红色的前脚的
“泥蜂”八月底的时候在我的花园里非常之多，它们很喜欢在薄荷草上睡眠。在傍晚薄
暮时，特别是在窒闷的日子里，暴风雨正在酝酿，大风大雨即将来临的时候，可是，我
们却能见到一个奇怪的睡眠者——仍然在那里安祥地熟睡着。
大概在晚上休息时，它的睡眠姿态没有比这个更奇怪的了。当你见到它以后就会觉
得特别的稀奇古怪了。它用颚咬入薄荷草的茎内，方的茎比圆的茎更能握的牢固一些，
它只用嘴咬住它，身体却笔直地横在空中，腿折叠着，它和树干成直角，这昆虫把全身
的重量，完完全全的放置在它的大腮上。
泥蜂利用它强有力的颚这样睡觉，身体伸展在空中。如果按动物的这种情形来推测，
我们从前对于休息的固有观念就要被推翻了。
任凭风暴狂欢，树枝摇摆，这位睡眠者并不被这摇晃的吊床所烦扰，至多是在某个
时候用前足抵住这摇动的枝干罢了。也许黄蜂的颚像鸟类的足趾一般，具有极强的把握
力，比风的力量还要强大许多。
尽管如此，有好几种黄蜂和蜜蜂都是采用这种奇怪的姿势来睡眠的——用大腮咬住
枝干，身体伸直，腿缩着。
大约在五月中旬，那时候恩布沙已经发育完整了。它的体态和服饰比螳螂更引入注
目。它还保留着一点幼稚时代的怪相——垂直的胸部，膝上的武器和它身体下面的三行
鳞片。但是它现在已经不能卷成钩子，它现在看起来也文雅多了：大型灰绿色的翅膀，
粉红色的肩头，矫捷的飞翔，下面的身体装饰着白色和绿色的条纹。
雄的恩布沙，是一个花花公子，和有些蛾类相似，更是夸张地用羽毛状的触须修饰
着自己。
在春天，农人们遇见恩布沙的时候，他们总以为是看到了螳螂——这个秋天的女儿
了。
它们外表很相像，以致人们都怀疑它们的习性也是一样的。因为外观一样，又都是
昆虫类的动物，所以人们没有认真仔细观察，也没有考察过它们的行动坐卧，所以就猜
测它们的生活习惯是一样的。
但是，事实上因为它的那种异常的甲胄，会使人们想到恩布沙的生活方式甚至比螳
螂要凶狠得多。但是，这种想法却错了，这个误解对恩布沙是不公平的，无调查研究的
结论是靠不住的。
尽管它们都具有一种作战的姿态，但是，恩布沙却是一个比较和平友好的动物呢！
它不是一个好斗好战的恶劣的凶手。
把它们关在铁丝罩里，无论是半打(一打是十二只，半打是六只)或者只有一对，它
们没有一刻忘掉柔和的态度。它们之间都是和平友好，互利相处的。
甚至到发育完成的时候，它们几个也是互相体谅，互相谦让，互不侵犯的。它们吃
的东西比较少，每天的食物只有两三只苍蝇就足够了。
食量大的小动物，当然是好争斗的。吃得饱的动物，把争斗当做一种消化食物的手
段，同时也是一种健身的方式。争强好胜，事事不让人，从来不吃亏，这是典型的弱肉
强食者的特点，它从来就是见便宜就占，见利益就争，见好事就抢。螳螂一见到蝗虫立
刻就会兴奋起来，于是战争就不可避免地开始了。螳螂立刻就扑向蝗虫，但是蝗虫也不
示弱，两者你争我斗，蝗虫用利齿欲扑向螳螂，但螳螂用它尖利的双夹给蝗虫以有力的
反扑。你争我斗的场面，十分精彩。
但是，节食的恩布沙，是个和平的使者，它从不和邻居们争斗，也从不用做鬼的形
状，去恐吓外来者。它也从不像螳螂那样，和邻居们争夺地盘。它从不突然张开翅膀，
也不像毒蛇那样作喷气、吐舌状。它从来也不吃掉自己的兄弟姐妹。更不像螳螂那样，
吞食自己的丈夫。这种惨无人道的事情，它是从来不做的。
这两种昆虫的器官，是完全一样的。所以这种性格上的不同，与身体的形状无关，
与其外表也无关。或许可以说是由于食物的差异而造成的。
无论是人还是动物，淳朴的生活总可以使性格变得温和一些，随和一些。这些都可
以营造一个和平共处的好环境。但是，自奉太厚了，就要开始残忍起来。贪食者吃肉又
饮酒——这是野性勃发的普遍原因——从不能像自制的隐士一样温和平静。它是吃些面
包，在牛奶里浸浸，这样简单的生活。它是一个普普通通的昆虫，它是平和、温柔、和
善的。而螳螂则是十足的贪食者。
虽然我的解释已经很清楚明白了，但是还有人可能会提出更深一层的问题。
这两种昆虫有完全相同的形状，想来一定也有同样的生活需要，而为什么，一种如
此的贪食，而另一种又如此有节制呢？它们在态度方面，如同别的昆虫已经告诉我们的
一样，嗜好和习性并不完全取决于自身的形状，以及身体结构，而是在决定物质的定律
方面，还有决定本能的定律存在。

二、白面孔螽斯
在我所居住的区域里的螽斯是白面孔的。无论在其善于歌唱，还是在其庄严的丰彩
上，它都可以算得上是蚱蜢类中的首领。它生有灰色的身体，一对强有力的大腮，以及
宽阔的象牙色的面孔。
如果要想捕捉它，这并不是什么难做到的事，也并不烦人。在夏天最炎热的时候，
我们常可以见到它在长长的草上来回跳跃。特别是在岩石下面，那里有松树生长着。
希腊字Dectikog(即白面孔螽斯、Decticns的语源)的意义是咬，喜欢咬。白面孔螽
斯因此取了这个名字。
它确实是善于咬的昆虫。假如有一种强壮的蚱蜢抓住了你的指头，你可是要当心一
点儿，它会把你的指头咬出血来，咬得你生疼，甚至有时疼痛难忍。它那强有力的颚仿
佛是凶猛的武器。当我要捕捉它时，我必须非常小心提防它，否则随时都有被它咬伤的
危险和被它咬破的可能。它那两颊突出的大型肌肉，显然是用来切碎它捕捉的、硬皮的
捕获物时用的。
把白面孔螽斯关在我的笼子里，我发现蝗虫蚱蜢等任何新鲜的肉食，都符合它们的
需要。特别是那种长着兰色翅膀的蝗虫，尤其适合它的嗜好。
当把食物放进笼子里时，常常会引起一阵骚动。特别是在它们饿极了的时候，它们
一步一步的很笨重地向前突进。因为受长颈的阻碍，它不能很敏捷地行动。有些蝗虫立
刻就被捉住，有的乱飞，乱蹦，乱跳，有的急了跳到笼子的顶上，逃出这螽斯所能捉捕
到的范围之外。因为它的身体很笨重，不能爬得那么高。不过蝗虫也只能是延长它们自
己的生命而已，最终也无法逃脱被白面螽斯蚕食的恶运。它们或因疲倦、或因被下面的
绿色食物所引诱，纷纷从上面跑下来，于是立刻就会被螽斯所捕获，成为其口中之美食。
这种螽斯，虽然智力很低下，然而却会用一种科学的杀戳方法。如同我们在别的地
方见到的一样，它常常先刺捕猎物的颈部，然后再咬住主宰它运动的神经，使它立刻失
去抵抗的能力。和其他肉食动物一样如哺乳动物虎、猎豹等等，它们都是先将所捕捉的
猎物的喉头管咬住，使其停止呼吸，丧失反抗力后，再一点点地享用它的肉体。
这是一种很聪明的方法，因为蝗虫是很难杀死的。有时虽然蝗虫的头已经掉下来了，
但它的躯体依然还能够跳动不已。我曾经见过几只蝗虫，已经被吃掉一半了，还不断地
乱跳，居然被它逃走了。
因它嗜好蝗虫，以及有些对于未成熟的谷类有害的种族，所以这类螽斯多一些，对
于农业也许有相当的益处。
不过现在它对于土地上保存果实的帮助，是非常薄弱的。它带给我们的主要的兴趣，
事实上是那些远古遗留下来的纪念物。它留给我们一些现今已经不用了的习性。
我应该谢谢白面孔螽斯，使我再次知道了关于幼小螽斯的一两件事情。
它产下的卵，并不和蝗虫，螳螂一样，把它们装在硬沫做成的桶里，它也不像蝉那
样，将它们产在树枝的洞穴里。
这种螽斯将卵像植物种子一般，种植在土壤里。母的白面螽斯身体的尾部有一种器
官，可以帮助它在土面上掘下一个小小的洞穴。在这个洞穴内，产下若干个卵，将洞穴
四周的土弄松一些，用这种器具，将土推入洞中，就像我们用手杖将土填入洞穴一样。
用这样一种方法，它将这个小土井盖好，再将上面的土弄平整。
然后，它到附近的地方散一会儿步，以作消遣和放松。用不了多长时间，它就会回
到先前产卵的那个地方，靠近原来的地点——这是它记得很清楚的——又重新开始工作。
如果我们注意观察它一个小时，就可以看到这种全部的动作，不下五次以上，连附
近的散步也包括在内。它产卵的地点，常是靠得很近的。
各种工作都已经完成以后，我察看这种小穴。只有卵放在那里，没有小室或壳来保
护它们。通常约有六十个，颜色大部分是紫灰色的，形状如同棱一样。
我开始观察螽斯的工作，就想看看它的卵子孵化的情形，于是在八月底的时候，我
取来很多的卵，放在一个里面铺有一层沙土的玻璃瓶子中。它们在里面度过八个月的时
间，感受不到气候变化带来的痛苦：没有风暴，没有大雨，没有大雪，也没有它们在户
外必须经受到的、过度炎热的光照和日晒。
六月来临时，瓶中的卵，还没有表现出开始孵化的征兆。和九个月以前，我刚把它
们取来的时候一样，即不发皱，也不变色，反而表现出极其健康的外观。在六月里，小
螽斯在原野里经常可以遇到了，有的，甚至已发育的很大了。因此我很怀疑，究竟是什
么理由使它迟延下来的。
于是，就产生了一种意见，这种螽斯的卵，如同植物一样，被种在土地里，是毫无
保护地暴露在雨雪之中的。
在我瓶子里的卵，在比较干燥的状况下，度过了一年的三分之二的时间。因为它们
本来是像植物种子一样散播着的。它的孵化大概也需要潮，需要适合它的一切孵化条件，
如同种子发芽时需要潮湿一样。这时，我开始决定要试一试。
我将从前取来的卵，分出一部分，放在我的玻璃管里，在它们上面，薄薄的加上一
层细细的潮湿的沙子。然后把玻璃管用湿棉花塞好，以保持里面的湿度。无论谁看见我
的试验，都会以为我是那种在试验种子的植物学家。
我的希望可以实现了。在温暖的，潮湿的环境之下，卵不久就表示出要孵化的迹象，
它们渐渐的，一点点地涨大，壳显然就要分裂开了。我花费了两个星期的功夫，每个小
时我都很认真仔细、不知疲倦地守候着它，想看看小螽斯跑出卵来的情形，以解决遗留
在我心中很长时间的疑问。
那个疑问是这样的。这种螽斯，按照惯例，是埋在土下边约一寸深的地方，现在这
个新生的小螽斯，夏初时在草地上跳跃，发育得完全一样，长有一对很长的触须，细得
如同发丝一般；并且身后生有两条十分异常的腿——像两条跳跃用的支撑杆，对于走路
是很不方便的障碍。
我很想知道，这个柔弱的小动物，携带着这样笨重的行李，当它到地面上来时，其
间所有的工作，是怎样进行的呢？它用什么东西从土中开出一条小道路来呢？它有遇到
一粒小沙就会折断的触角，少许的力量就会断脱的长腿，这个小动物是显然不可能从土
坑中解放出来的。
我已经告诉过你们：蝉和螳螂，一个从它的枝头、一个从它的巢出来时，穿有一种
保护物，就像一件大衣一样。
我想，这个小螽斯，从沙土里钻出来的时候，一定也有比出生以后，在草间跳跃时
所穿的还要简单而且又紧又窄的衣服，作为一种保护。
我的估计并没有错。这时候，白面孔螽斯，和别的昆虫一样，的确穿有一件保护外
衣。这个细小的，肉白色的，小动物，已经长在一个鞘里了，六个足平置胸前，向后伸
直。
为了让出来时比较容易一些，它的大腿绑在身旁；另一半不太方便的器官——触须
——一动也不动地压在包袋里面。
它的颈弯向胸部。大的黑点——是它的眼睛，那毫无生气而且十分肿大的面孔，使
人以为那是盔帽。颈部则因头弯曲的关系，十分开阔。它的筋脉同时微微地跳动着，时
张时合，因为有了这种突出的、可以跳动的筋脉，新生的螽斯的头部才能自由转动。依
赖颈部推动潮湿的沙土，挖掘出一个小洞穴。于是筋脉张开，成为球状，紧塞在洞里，
在它的幼虫移动它的背，并推土时，可以有足够的力量。
如此，进一步的步骤已经成功了，球泡的每一次涨起，对于小螽斯在洞中的爬动，
都是很有帮助的。
看到这个柔软的小动物，身上还是没有什么颜色，移动着它那膨胀的颈部，攒掘土
壁，真是可怜。
它的肌肉还没有达到强健的时候，这真无益于与硬石的斗争啊！不过经过不懈的奋
斗它却居然获得了最终的成功。
一天早晨，这块地方，已经做成了小小的孔道，不是直的，约有一寸深，宽阔得像
一根柴草。一般用这样的方法，这个疲倦的昆虫终于可以达到地面上了。
在还没有完全脱离土壤以前，这位奋斗者也要休息一会儿，以恢复它这次旅行后的
精力。再作一次最后的拼搏，竭力膨胀头后面突出的筋脉，以突破那个保护它已经很久
的鞘。这个动物就这样将外衣抛弃了。
于是，这是一个幼小的螽斯了，它还是灰色的，但是，第二天就渐渐变黑了，同发
育完全的螽斯比较起来简直是成了一个黑奴了。不过它成熟时的象牙面孔是天生的，在
大腿之下，有一条窄窄的白斑纹。
在我面前发育的螽斯啊！在你面前展开的生命是太凶险了。
你的许多亲属们，在尚没有得到自由之前，就因疲倦而死去了。在我的玻璃管中，
我看到了好多螽斯因受到沙粒的阻碍而放弃了尚未成功的奋斗。
它的身上长有一种绒毛，欲将它的尸体包裹起来。如果我不去帮助它，到地面上来
的旅行会更加危险，因为屋子外面的泥土更加粗糙，已经被太阳晒硬了。
这个有白条纹的黑鬼，在我给它的莴苣菜叶上咬啮，在我给它居住的笼子里跳跃着，
我可以很容易地豢养它。
不过它已不能再提供给我更多的知识了，所以，我就恢复了它的自由，以报答它教
给我的那些知识，我送给它这个房子——玻璃管，还有花园里的那些蝗虫。
因为它教给我蚱蜢在离开产卵的地点时，穿着一件临时的保护衣服，将那些最笨、
最重的部份，如它的长腿和它的触角等等，全都包在鞘里。它又告诉我这种略微伸缩、
干尸状的动物，为了它旅行的方便，它的头颈上生有一种瘤，或者说是颤动的泡口——
是一种原来就生成的机器，在我最初观察螽斯的时候，我并没有看见它用它作为走路的
帮助。
------------------








黄蜂



一、它们的聪明和愚笨
在九月里的一天，我和我的小儿子保罗跑出去，想去瞧一瞧黄蜂的巢。
小保罗的眼力非常好，再加上特别集中的注意力，这些都有助于我们的观察很好地
进行。我们两个饶有兴趣地欣赏着小径两旁的风景。
忽然，小保罗指着不远的地方，冲着我喊了起来：“看！一个黄蜂的巢。就在那边，
一个黄蜂的巢，比什么都要更清楚呢！”果然，在大约二十码以外的地方，小保罗看见
一种运动得非常快的东西，一个一个地从地面上飞跃起来，立即迅速地飞去，好像那些
草丛里面隐避着小小的即将爆发的火山口，马上要将它们一个个喷出来一般。
我们小心谨慎地慢慢地跑近那个地点，生怕一不小心，惊动了这些凶猛的动物，引
起它们对我们的注意和攻击，那样的话，后果可是不堪设想的。
在这些小动物们的住所的门边，有一个圆圆的裂口。口的大小大约可容下人的大拇
指。同居一室者来来去去，进进出出，摩肩接踵地向相反的方向飞去飞回，不停地忙碌
着。
突然，“噗”的一声，我不觉吃了一惊，但是马上又醒悟过来了。我忽然想起现在
我们正处于一个很不安全的时刻。要是我们太靠近去观察它们的行踪，就会引起不良的
后果。因为，这样的不速之客会让它们感到不安，会激怒这些容易发脾气的战士来袭击
我们。因此，我们不敢再多观察了。再观察下去就意味着要“牺牲”更多的东西了。
我和小保罗记住了那个地点，以便日落后再来观察。到了夕阳西下的时候，这个巢
里的居住者，全体都应该从野外回家了。那样，我们就可以更好地观察了。
当一个人决定要征服黄蜂的巢时，如果他的这一举动，没有经过谨慎而细致的思考
的话，那么这种行动简直就是一种冒险的事情。半品脱的石油，九寸长的空芦管，一块
有相当坚实度的粘土，这些构成了我的全部武器装备。还有一点必须提到的是，以前的
几次小小的观察研究，稍稍积累了一点儿成功的经验。这所有的一切物品与经验对我而
言，是最简单，同时也是再好不过了。
有一种方法对我是至关重要的，那就是窒息的方法。除非我打算用我所不能够忍受
的牺牲的方法。否则，我必须掌握窒息的方法。当瑞木特要把一个活的黄蜂的巢放在玻
璃匣子内，观察里面的同居者的习性的时候，他不是亲自行动，而是选择了另一种方法，
雇佣了一个帮手，协助他进行实验。这个帮手经常从事这种痛苦不堪的工作。为了获得
优厚的报酬，情愿牺牲自己的皮肤，为科学家们提供有偿性服务。但是，我可是打算牺
牲自己的皮肤的。
在还没有挖出我所要的蜂巢之前，我仔细思考了两次。然后，才开始我的计划。我
首先将蜂巢里的居民窒闷住，死了的黄蜂就不能刺人了。这是一个残忍的方法，但也是
一个十分安全的方法，可以让我不至于身处危险之中。我采用的是石油，因为它的刺激
作用不至于过于猛烈。
因为我要做一次观察，所以我希望能留下一部分不死的黄蜂，否则的话老是观察死
了的对象，就前功尽弃了。现在的问题只是在于如何把石油倒进有蜂巢的穴里去。蜂巢
穴的出入孔道大约有九寸长，而且差不多和地面是平行的，一直通到地底下的窠巢。假
如把石油直接倒入隧道的口上，这便是一个大大的错误，而且将会带来极其严重的不良
后果。为什么呢？主要原因是，这样少量的石油，会被泥土吸收进去，而无法到达地下
的窠巢。这样一来，到了第二天，当我们凭着想象，以为这时挖掘、凿开窠巢一定是很
安全的时候，就会遇到很大的危险。我们就会碰到一群火上浇油般的黄蜂，在我们的铁
铲下回旋，从而对我们造成一定的威胁。
早已准备好的九寸长的空芦管可以阻止这一不幸事件的发生。把这根空芦管插进差
不多九寸长的隧道里面的时候，就形成了一根自动引水管。于是石油可以顺着导管流人
土穴中，一点儿也不会漏掉，而且，速度还很快。然后，我们再用一块事先已经捏好的
泥土，像瓶塞子一样，塞住出入的孔道口，断绝这些黄蜂的后路。我们所要做的工作就
到此为止了，剩下的就只是等待了。
当我们准备做这项工作的时候，是昏暗的月夜，正是九点钟，小保罗和我一起出去。
我们只带了一盏灯，还有一篮子需要用到的工具。当时，远远的还可以听见农家的狗还
在互相吠叫着，猫头鹰在橄榄树的高枝上叫着，蟋蟀在浓密的草丛中不停地奏着动听的
音乐。小保罗和我则在谈论着昆虫。他热切而好奇地向我提出很多问题。为了不让他失
望，我将我所知道的一切告诉了他，帮助他学习，以丰富他的知识，满足的他的兴趣。
这样一个快乐的猎取黄蜂的夜晚，让我们忘记了睡眠和被黄蜂攻击时的痛苦。
将芦管插入土穴中是一件非常精巧的工作，需要一些技巧。因为孔道的方向是无从
知晓的，需要颇费一番猜疑和试探。而且有时候，黄蜂保卫室里的门卫会突然警觉地飞
出来，毫不客气地攻击正在进行这项工作而且没有防备的人的手掌。为了防止这种措手
不及的不幸事情发生，我和小保罗中的一个人，在一旁守卫，时刻警惕着，并用手帕不
停地驱赶着进攻的敌人。这样一来，即使最后有一个人的手上不幸被命中，隆起了一块，
就是很疼痛，也算是一个理想的，不很大的代价，尚可以忍受。
在石油流入土穴中以后，我们便听到地下传来的众蜂惊人的喧哗声。然后，很快地，
我们用湿泥将孔道封闭起来，一次一次地用脚踏实，使封口坚不可摧，从而使它们无路
可逃。现在，没有什么其他的事可以做了。于是，我和小保罗就跑回去睡觉休息了。
第二天清晨，我们带了一把锄头和一把铁铲，重新又回到了老地方。早一点儿去比
较好些，因为可能有很多黄蜂夜里是在外面游荡的，它们有可能在我们挖土的时候飞回
来，这就糟糕了，因为这对我们又将是一种威胁。另外，清晨的冷气，可以多少削弱一
些它们的凶恶和威风。
在孔道之前，芦管依然还插在那边，我和小保罗挖了一条壕沟，宽度刚好能容下我
们俩，行动很方便。于是，我们从沟道的两边开始挖。很小心地一片一片铲去。后来，
挖了差不多有二十寸深，蜂巢便暴露出来了。它吊在土穴的屋脊当中，一点儿也没有被
损坏，完好地吊在那里，这真让我们感到高兴。
这真是一个壮观美丽的建筑啊！它大得简直像一个大南瓜。除去顶上的一部分以外，
各方面全都是悬空的，顶上生长有很多的根，其中多数是茅草根，穿透了很深的“墙壁”
进入墙内，和蜂巢结在一起，非常坚实。如果那地方的土地是软的，它的形状就成圆形，
各部分都会同样的坚固。如果那地方的土地是沙砾的，那黄蜂掘凿时就会遇到一定的阻
碍，蜂巢的形状就会随之有所变化，至少会不那么整齐。
在低巢和地下室的旁边，常常留有手掌宽的一块空隙，这块面积是宽阔的街道。这
些建筑者，在这里可以行动自由，继续不停地进行它们各自的工作，用它们自己的双手，
使它们的窠巢更大更坚固。通向外面的那条孔道，也通向这里。在蜂巢的下面，还有一
块更大一些的空隙，其形状是圆的，就如同一个大圆盆，在蜂巢扩建新房时，可以增大
其体积。这个空穴，还有另外一个用途，那就是盛废弃物品的垃圾箱。看来这里的基本
建设还是较为齐全的。
这个地穴是黄蜂们用自己的“双手”亲自挖掘出来的。关于这一点，是用不着怀疑
的。因为如此之大、如此整齐的洞穴，在自然界是没有现成的。当初，第一个开辟这个
巢的黄蜂，也许是利用了鼹鼠所做的洞穴，加以借用，以图开始创建的便利。可是，筑
巢的绝大部分工作都是黄蜂亲自操作的。然而，事实上，并没有一些挖出的泥土堆积在
蜂巢的大门之外。那么，黄蜂们挖出的泥土被搬运到哪里去了呢？答案是：它们已经被
弃散在不引人注意的广阔的野外去了。有成千上万的黄蜂参与挖掘这个壮丽的建筑物，
必要的时候，还要将它扩大。这千百万只黄蜂，飞到外面来的时候，每一个身上都附带
着一粒土屑，抛散在离开窠巢很远的各处土地上。因此，挖出的泥土的痕迹一点儿也看
不到了。所以，蜂巢看上去像一片净土一样。
黄蜂的巢是用一种薄而柔韧的材料做成的。这种材料是木头的碎粒，很像一种棕色
的纸。它的上面有一条条的带，其颜色视所用木头的不同而不同。如果蜂巢是用整张
“纸”做的，就可以稍稍抵御寒冷，起到保暖的作用。但是，黄蜂就像做气球的人一样，
它们懂得温度可以利用各层外壳中所含有的空气来保持。所以，黄蜂把它们的低巢做成
宽的鳞片的形状，一片一片松松的铺起来，显出很多的层次来，整个蜂巢形成一种粗粗
的毛毯状，厚厚的，而且多孔，其内部含有大量的空气。这样一来，外壳里的温度，在
天气很热时，一定是很高的。
大黄蜂——黄蜂们的领导，在一样的原则下，建筑它自己的巢。在杨柳的树孔中，
或者是在空的壳层里，它用木头的碎片，做成脆弱的黄色的纸板。它就利用这种材料来
包裹它自己的窠。一层一层相互地重叠起来，就像个凸起的大鳞片一样，可以想象这有
多么保暖！这个大鳞片的中间有充分的空隙，空气停留在里边也不流动。
黄蜂们的动作常常与物理学和几何学的定理相吻合。它们可以利用空气——这个不
良导体来保持它们家里的温度。它们早在人类还未曾想到做毛毯之前就已经做出来了，
而且技艺还很高，它们在建筑窠巢的外墙时，只要极小的外围，就足以造出很多的房间，
它们的小房间也同样如此，其面积与材料都非常经济。
然而，尽管这些建筑家们有如此之聪明智慧，但是，令我们感到奇怪的是，当它们
遇到最小的困难时，居然会束手无策，愚笨无比。一方面，它们得益于大自然的本能指
导它们像科学家一样地行动；而另一方面，很显然它们完全不具备反省的能力，其智力
是相当低下的。关于这一个事实，我已用各种各样的试验加以证明了。
黄蜂碰巧将自己的房子安置在我家花园的路旁边，于是，我便可以利用一个玻璃罩
来做试验了。在原野里的时候，我无法利用这种器具，因为乡下的小孩子们很快就会把
它打破，而破坏了我精心准备的试验。有一天晚上，天已经黑了，黄蜂也已经回家了。
我弄平了泥土，放上一个玻璃罩罩住黄蜂的洞口。第二天早晨，黄蜂们习惯性地开始工
作。当它们发觉自己的飞行受到阻止时，它们是否能够在玻璃罩的边下挖掘出另外一条
道路呢？是不是这些能够掘出广阔洞穴的刚强的动物知道只要创造一条很短的地道，便
可以让它们重获自由呢？这便是我们的问题关键之所在了。那么，结果如何呢？
第二天早晨，我看到温暖耀眼的阳光已经落在玻璃罩上了。这些工作者们已经成群
地由地下上来，急于要出去寻觅它们的食物。但是，它们一次又一次地撞在透明的“墙
壁”上跌落下来，重新又上来。就这样，成群地团团飞转不停地尝试，丝毫不想放弃。
其中有一些，舞跳得疲倦了，脾气暴躁地乱走一阵，然后重新又回到住宅里去了。有一
些，当太阳更加炽热的时候，代替前者来乱撞。就这样轮换着倒班。但是，最终没有一
只黄蜂大智大勇，能够伸出手足，到玻璃罩四周的边沿下边抓、挖泥土，开辟新的谋生
之路。这就说明它们是不能设法逃脱的。它们的智慧是多么有限啊。
这个时候，有少数在外面过夜的黄蜂，从原野归来了。它们围绕着玻璃罩盘旋飞舞，
一直迟疑徘徊，不知如何是好。有一个带头决定往玻璃罩的下边去挖，其他的黄蜂也随
着学它的样子。于是，大家齐动手，很快地，一条新的通路很容易地开辟出来了。它们
也就跑了进去，终于到家了。于是，我用土将这条新辟之路堵住。假设从里面能够看出
这条狭窄的通路，当然可以帮助罩内的黄蜂轻而易举地逃走。我很愿意让这些囚徒通过
自己的观察和努力争得自由的光荣，享受阳光沐浴的欢乐，领会大自然的优美。
无论黄蜂的理解能力是如何的差劲。我想它们的逃脱，现在应该是可能实现的了。
那些刚刚进去的黄蜂当然会指引一下路径，它们会指教其他的黄蜂向玻璃罩下边挖，以
便尽快地逃离牢笼。
然而，事实却并不那么乐观。我非常失望，可爱的黄蜂们居然没有一点儿要从经验
和实例上仿效学习的企图。在那个玻璃罩里，一点儿没有要继续挖掘出逃之路——地道
的迹象。这些小昆虫们只是依旧团团乱飞，毫无计划，毫无目的，它们只是盲目地乱碰
乱撞，挤作一团不知究竞发生了什么意外。每天都有很多可怜的黄蜂死于饥饿和炎热之
下。一个星期以后，很遗憾，没有一个活的黄蜂能够侥幸存活下来，全军覆没了。一堆
死尸铺在地面上，其状况尤其惨烈。
从原野里返回的黄蜂们可以另辟新路，毫不费力地回到自己的家中。其原因是，从
泥土外面可以嗅知它们的家，并去寻找它。这是黄蜂自然本能想方设法投入家的怀抱的
一种表现，或者说是它们的一种防御方法。这是不需要任何思想和解释的。自从小小的
黄蜂初次降临到这个世界上的时候开始，地面上的一切阻碍，对于每一个黄蜂而言，都
已经很熟悉了。
但是，对于那些不幸被罩在玻璃罩里的黄蜂，就没有这种本能来帮助它们逃离险境
了。它们的目的是明确的，单一的。它们想到阳光里面去，到野外去觅食。它们被罩在
玻璃罩里，在这个透明的牢狱中，能够看到日光，它们便被蒙骗了，以为自己的目的已
经完全达到。虽然它们几经努力，一往无前，继续不断地和玻璃罩相抗衡、相碰撞，心
中抱有无限希望，想朝着日光，飞得再远一点儿。以便能觅到急需的食物。可是事实上
那是无用的。在它们以往的经历中，没有任何经验和实践指导它们遇到这种情况时，应
该如何行事。于是它们走投无路，别无选择，只能盲目地固守着它们生来就惯有的老习
性，从而生的希望越来越小，而逐渐将自己推向无奈的死亡。

二、它们的几种习性
假如我们掀开蜂巢的厚包，便可以看到里面隐藏有许多的蜂房，那好几层的小房间，
上下排列着，中间用稳固坚实的柱子紧密连在一起，层数是不一定的。在一定季节的后
期，大概是十层，或者是更多一些。各个小房间的口都是向下的。在这个种类看起来很
奇怪的小世界里，幼蜂无论是睡眠还是饮食，都是脑袋朝着下边生长的，即倒挂着的。
这一层一层的楼即蜂房层，有广大的空间把它们分隔开。在外壳与蜂房之间，有一
条门路与各个部分是相通的。经常有许多的守护者进进出出，负责照顾蜂巢中的幼虫。
在外壳的一边，矗立着这个丰富多彩的都市的大门，一个没有经过什么过多装饰的裂口，
隐藏在被包着的薄鳞片中。直面对着这个大门的，就是那从地穴深处直通到外面的大千
世界的隧道的进出口。
在黄蜂的社会中，生活着数量众多的黄蜂。它们的全部生命完全都投入到不辞劳苦
的工作之中。它们的主要职责就是，当人口不断增加的时候，就不停地扩建蜂巢，以便
新的公民居住。尽管它们并没有自己的幼虫，可它们呵护巢内的幼虫，却是极小心勤勉，
无微不至的。
为了能观察到它们的工作状况，以及快到冬天的时候它们之中会有什么事情发生，
我在十月里，把少许巢的小片放在盖子下面，里面居住着很多的卵和幼虫，并且还有一
百个以上的工蜂在细心地看护着它们。
为了便于进行观察，我将蜂房分隔开来，让小房间的口朝着上面，然后并排放着。
这样颠倒的排列，看起来似乎并没有使我的这些囚徒们烦恼，它们很快地就从被打扰的
情形下适应过来，恢复了原来的空间状态，重新开始忙碌而辛勤地工作，似乎从来没有
什么事情发生过一样。
事实上，它们当然需要再建筑一点东西。所以，我便选择了一块软木头送给它们，
并且用蜂蜜来喂养它们，满足它们的需要。用一个拿铁丝盖着的大泥锅来代替隐藏蜂巢
的土穴。再盖上一个可以移动的纸板做的圆顶形的东西，使得内部相当黑暗。当然，当
我需要亮一些时就把它移开。
黄蜂继续进行它们自己的日常工作，就好像从来没有受到过任何的扰乱一样。工蜂
们一面照料着蜂巢中的蜂宝宝，与此同时，又要照顾好它们自己的房子。它们一起努力
加油，开始慢慢地筑起一道新的铜墙铁壁。这墙壁围绕着它们最封闭的蜂房。看起来，
它们似乎是打算重新再建筑一个新的外壳，作为新的外壳，来代替那个被我用铁铲毁坏
了的旧外壳。但是，这些工蜂们并不是简单地修修补补，它们是从被我破坏了的那个地
方开始它们的工作。它们很快就筑成了一个弧形的纸鳞片似的房顶，然后，用它遮盖住
大约三分之一的蜂房。如果这个小蜂巢不曾遭到我的破坏的话，那么这些工蜂们搭建起
的这个屋顶足可以连接到外壳呢。它们亲手做成的一个房顶，还不够大，只能遮盖住整
个小房间的一部分而已。
至于我事先为它们精心准备好的那块软木头，它们根本不予理睬，甚至连碰都不曾
碰上一下，仿佛它根本就不存在一样。或许这种“新型”的材料，对于黄蜂而言，用起
来很不方便。它们宁愿放弃它，而继续选用那已经废弃不用了的旧巢，这样更加方便，
而且更加得心应手一些。因为在这些旧的小巢内，不必辛辛苦苦地重新制做纤维，因为
它们是已经做好了的，方便实用。只需拿来主义即可。而且，它们也不用浪费很多的唾
液，只需相当少的唾液，再用它们的大腮仔细咀嚼几下，然后便形成了上等质地的浆糊，
这是相当好的建筑材料。
下一步，它们一起把不居住的小房间统统毁得粉碎。然后，利用这些碎物，做成一
种似天篷一样的东西。如果有必要的话，它们也会再次利用同样的方法，筑造出新的小
房间，以便居住，活动之用。
与它们齐心协力筑造屋顶的工作相比较，更加有趣味的要算是喂养蛴螬幼虫了。刚
才还是一个个粗暴刚强，卖力气的战士，这会儿就摇身一变，成了温柔、体贴的小保姆。
看到这些，谁也不会感到厌倦和反感的。一下子，充满了战斗气息的军营一样的窠巢，
立刻变成了温馨的育婴室了。真是妙趣横生啊！
喂养好可爱的、柔弱的小宝宝，可是需要相当的耐心与细致的。假如我们只将注意
力集中到一个正在忙碌工作的黄蜂身上，我们就可以清楚地观察到，在它的嗉囊里，充
满了蜜汁。它停在一个小房间的前面，它的样子特别有意思，它把它小小的头慢慢地伸
到洞口里面去，然后再用它的触须的尖儿去轻轻地碰一碰里面的一个小幼虫。那个小宝
宝慢慢地清醒了过来，似乎看到了那个黄蜂递送进来的触须，于是向它微微地张开小嘴。
它的样子，特别像一只刚刚初生不久，羽毛尚未丰满起来，乳臭未干的小鸟，正在向着
刚刚辛辛苦苦为它觅食而归的妈妈伸出小嘴，急切地索要食品一般，不觉得让人感到一
阵温馨。
不一会儿，这个刚刚从梦中苏醒过来的小宝宝，将它的小脑袋摇来摆去的，渴望着
能够马上探索到它急切需要得到的食物，这可以算是它的本能天性了。然而它又是盲目
地探寻着，一次次试探着外面的黄蜂为它们提供的食物。可以想象小宝宝的急切心情网，
终于两张小嘴接触到了。一滴浆汁从“小保姆”的嘴里流出来，流了进那个被看护者的
小嘴里。仅仅这一点点就足够一个小宝宝享用了。现在，该轮到第二个黄蜂婴儿进食了。
于是，这个小保姆又赶快马不停蹄地跑到别处去，继续履行它神圣的职责。
小宝宝们通过口对口地交接食物后，享受到大部分的蜜汁。但是，进食并没有完全
告终，它们还没有享用完呢。因为，在喂食的时候，幼虫的胸部会暂时膨胀起来，其作
用就如同一块围嘴或餐巾纸一样，从嘴里流出来的东西全都滴落在它上面。这样等保姆
走后，小宝宝们就会在它们自己的颈根上舐来舐去，吮吸着滴在胸部的蜜汁，尽情地享
受着美味的食物，不浪费一点儿食物。大部分的蜜汁咽下之后，幼虫胸部的鼓胀便会自
然而然地消失掉。然后，幼虫会稍微往蜂巢里缩进去一点，继续回到它的甜蜜的梦乡里。
当黄蜂在我的笼子里喂养小宝宝时，小幼虫们的头是朝上的，从它们的小嘴里遗漏
出来的东西，自然会滴落在它们的围嘴上面。至于在蜂巢里喂养它们的时候，它们的小
脑袋则是朝下的，但是，我并不心存怀疑，那便是在这样头向下的位置上，小幼虫的餐
巾纸仍能发挥作用，而且，功效是一样的。这是因为幼虫在蜂巢中时，它的头不是直的，
而是略微有一点弯度的。因此，它们嘴里溢出来的蜜汁很可能是堆积在那块小小的围嘴
上。而且，溢出的蜜汁是非常粘稠的，很快就会粘在围嘴上。与此同时，细心的小保姆
就是再放下一部分食品在这个地方，也是有可能的。所以，无论小幼虫的头是朝下的还
是朝上的。无论那块围嘴是在嘴的上边，还是在嘴的下边，这都不会阻碍围嘴充分发挥
它的作用。其主要原因是这种食品非常有粘性，可以牢牢地附着在围嘴上。因而可以说，
这块小小的围嘴简直就是一个又方便又及时的小碟子，它可以减少喂食工作的困难，避
免许多不必要的麻烦。为我们的小保姆们提供方便，使得它们又省力又省时，而且，它
还可以使得小幼虫们能够舒适、宁静地享用它们的美味佳看，直至一饱口福，满足为止。
还有一个好处，那就是不致让小宝宝们吃得太饱，撑坏了小肚皮而夭折。
如果是在野外，置身于大自然中，每当一年快要结束时，也是果品数量非常少的时
候，有些青黄不接。在这种情况下，大多数的小保姆挑选其它的食物来继续喂养小幼虫。
它们大多选择苍蝇，先将它们一一切碎，然后再喂给小幼虫们食用。但是，在我为它们
制作的笼子中，一概不选择其他的东西作为幼虫的食品，我只为它们提供充足的有营养
的蜜汁。
吃了这些蜜汁以后，所有的看护者和被看护者似乎都变得精力旺盛起来。而且，一
旦有什么不速之客突然闯进蜂房里，进行袭击侵略，那么它们将很不幸地立刻被处以死
刑。显然，黄蜂分明是一种不好客的生灵，从不厚待宾客，更不允许其他动物随意侵扰
自己的家园。即使是那种叫拖足蜂的蜂，形状和颜色与黄蜂是极其相像的，如果它们稍
一走近黄蜂，来分享它们的蜜汁，那么它们的这种企图很快就会破灭，马上被觉察出来，
于是，黄蜂们群起而攻之，直到致其于死地为止。拖足蜂的外貌并不能欺骗黄蜂那敏感
的目光，如果拖足蜂反应不迅速，没有及时退避，那么就会大难临头引来杀身之祸，被
黄蜂残酷地处死。所以，擅自闯入黄蜂的巢，实在不是一件明智的事情。即使来访的客
人的外表与它们极为相似，工作也与它们大同小异，几乎可以说是团体中的一份子，都
是绝对不行的。黄蜂是不会轻易地放过任何不请自来，不知趣的所谓客人的。因此，其
它动物还是退避三舍，回避为佳。
我已经一而再、再而三地看到黄蜂对于客人们的野蛮的待遇，假如闯入境内的不速
之客，是个相当有杀伤力，而且凶猛无比的家伙，当它受到群攻而牺牲后，其尸首便会
马上被众蜂拖到蜂巢以外，抛弃在下面的垃圾堆里边。但是，黄蜂似乎不会轻易地动用
它那有毒的短剑来攻击其他的动物，还是比较手下留情的。如果我把一个锯蝇的幼虫抛
到黄蜂群里面，对于这条绿黑色的小龙一样的侵入者，黄蜂们表示出很大的兴趣，它们
一定是感到很奇怪。接下来，它们便向它发起进攻，把它弄伤，但是并不利用它们带毒
的针去刺伤它。然后众蜂齐用力，要把它拖出巢去。与此同时，这条“小龙”也不服输，
不断进行抵抗，用它的钩子钩住蜂房。有时利用它的前足，也有时利用它的后足。然而，
最终这条可怜的“小龙”还是因为伤势太重，而且它还很软弱，最终被有力的黄蜂拉了
出来。这条“小龙”很惨，小小的身体上充满了血迹，被一直拖到垃圾堆上去。黄蜂们
驱赶这样一条并无什么力气的可怜虫可并不轻松，耗费了足有两个小时的时间呢！
如果，与此相反，我放的不是一个弱小的幼虫，而放一个住在樱桃树孔里的一种相
对比较魁伟的幼虫在蜂巢里面，结果就不同了。立刻会有五六只黄蜂拥上来，纷纷用有
毒的针去刺它的身体。不大一会儿，差不多几分钟以后，这只较强壮有力的幼虫终于也
难逃恶运，一命呜呼了。但是，接下来便产生出了一个问题。这具笨重的尸体，很难把
它搬运到巢外去。所以，黄蜂们觉得动它不成，便选择了其它的方法，比如吃掉它，或
者，至少要设法使它的体重有所减少。因此，它们便一直吃它，直到吃剩下的那部分可
以被拖动为止。然后，还是要把它拖到外面去，抛弃掉。

三、它们悲惨的结果
有了如此凶猛而又残酷的方法来抵御外来闯入者的入侵，还有如此巧妙而又温柔的
喂食方法，我制作的笼子里的小幼虫们一天一天茁壮成长着，黄蜂的家族日益兴旺起来。
不过，当然也存在例外的现象。黄蜂的窠巢里，也有一些非常柔弱、不走运的小幼虫，
它们长大成人，还未经历世间的风雨，沐浴阳光的温暖，便早早地夭折而去了。
我通过观察，发现了那些柔弱的病者，亲眼目睹它们不能继续享用蜜汁，不能进食，
渐渐地，一点点憔悴下去，衰弱下去。那些小保姆们早已比我更清楚地知晓了这一切。
它们十分无奈地把头轻轻弯下来，朝着那些可怜的患病者，用触须很小心地去试听一下，
最后得出结论，证明这些病者的确是不可医治，无法挽救了。于是，慢慢地，这个弱小
的生命逐步走向生命的尽头，快到死的程度了。最终，被毫不怜惜地从小房间里拖到蜂
巢的外面去。在充满野蛮气息的黄蜂的社会里，久病者不过仅仅是一块没有用处的垃圾
而已，越是赶紧拖出去越好，否则的话，就有蔓延传染的可能。对于黄蜂而言，那将是
很可怕的事情。但是这还不是最坏的可能。因为，随着冬天渐渐来临，黄蜂们大都已经
预感到它们将来的命运。它们深知，末日就在眼前了。
十一月里，非常寒冷的夜，使得蜂巢的内部起了变化。大搞基础建设的热情逐渐衰
退了。到储蜜的地方，从事储蜜工作的黄蜂不再频繁地到那里去了。整个家庭，所有黄
蜂全都逐渐地放任自流了。幼虫由于饥饿而大张着它们的小嘴，然而，等到的只不过是
非常迟缓的救济品，或者干脆没有小保姆愿意光临到这里来给它们喂食。深深的惆怅牢
牢占据了那些小保姆的心灵，它们从前的那份工作热情也不见了，最终竟转化为厌恶。
它们知道，再过不久的时间，一切就将变成不可能了。那么，小保姆还有什么存在价值
吗？还会有什么好处吗？当然，在看护蜂的心中，答案是否定的。于是，饥饿的时候来
临了。噩运降临到小幼虫的头上，它们悲惨而孤独地死去。因此，从前的那些温柔体贴
的小保姆一变而成为不可思议的凶残的刽子手了。
那些小保姆会对自己说：“我们没有必要留下许许多多的孤儿。不久以后，等我们
都离开了这里以后，还能有谁来照顾这些可怜的后代呢？没有。既然是这样的结果，那
还不如让我们亲手来把这些卵和小幼虫统统杀死。这样一个十分残暴的结果，总比那种
慢慢被饥饿煎熬而死要强得多，长痛不如短痛嘛！”
接下来的一幕闹剧，便是一场凶残的大屠杀行动。黄蜂们残忍地咬住了小幼虫颈项
的后面，然后粗暴地把它们一个个从小房间里拖出来，拉到蜂巢的外面去，抛到外面土
穴底下的垃圾堆里，其情景简直是惨不忍睹！
那些小保姆，也就是工蜂，在把幼虫从小房间中强行拖拉出来时，那种情形之残酷，
就好像这些幼虫都是一些从外面来的生客一般，或者是一群已经死掉了的尸体。它们野
蛮地拖着小幼虫的尸体，并且还要将它们的尸体扯碎。至于那些小卵，则会被工蜂们撕
扯开来，最后把它们吃掉。
在此之后，这些小保姆，即刽子手，毫无生气地留着它们自己的生命。一天一天地，
我带着无比的惊奇，注视着我的这些昆虫的最终的结局。非常出乎我的意料，这些工蜂
忽然间都死掉了。它们跑到上面，跌倒下来，仰卧着，从此再也没有爬起来，就如同触
了电一般。它们也有它们自己的生命周期。它们被时间这个无情无义的毒品毒死了。就
算是一只钟表内的机器，当它的发条被放开到最后一圈时，也是会如此的。
工蜂是老了！然而，母蜂是蜂巢中最迟生出来的，它们既年轻，又强壮。所以，当
严冬降临，威胁到它们时，它们还仍有能力来抵挡一阵。至于那些末日已经临近的，很
容易地就能从它们的外表的病态上分辨出来。在它们的背上，是有尘土沾附着的。在它
们尚健壮，还年青的时候，它们一旦发现有尘土附着在身上，就会不停地拂拭，把它们
黑色、黄色的外衣清洁得十分光亮。然而，当它们有病时，也就无心注意卫生清洁了。
因为已经无暇顾及了。它们或是停留在阳光底下，一动也不动，或者很迟缓地踱来踱去。
它们已经不再拂拭它们的衣裳了，因为这已不再重要了，也没有任何意义了。
这种对装束的不在意，就是一种不祥的征兆。过了两三天以后，这个身上带有尘土
的动物，便最后一次离开它自己的巢穴。它跑出来，主要是打算再最后享受一点日光的
温暖。忽然，它跌倒在地上，一动也不动，再也不能够重新爬起来了。它尽量避免自己
死在它所热爱和生存的巢里。这是因为，在黄蜂中，有一种不成文的“法律”规定，那
就是巢里是要绝对保持干净整洁的。这个生命即将结束的黄蜂，要自行解决它自己的葬
礼。它把自已跌落在土穴下面的坑里。由于要保持清洁卫生，这些苦行主义者，不愿意
自己死在蜂房里。至于那些剩余下来的，还没有死去的黄蜂，它们仍然要保留这种习惯，
直到它们最终的结局为止。这形成了一种不曾被摒弃的法律条文。无论黄蜂世界中的人
口如何增加，或是减少，这一传统总是要保持遵守的。
我的笼子里，一天天地空起来了。虽然这个屋子仍然是暖和的，而且里面还储备有
很多的蜜汁，供剩下来的那些健康者食用。但是，到了圣诞节的时候，仅仅剩下了约一
打的雌蜂。到一月六日，连最后剩余下来的黄蜂也全都死掉了。
那么，这种死亡是从哪里来的呢？让我的黄蜂统统都倒毙了。它们并没有受过饿，
也没有挨过冻，更没有经历过离家的痛苦。那么，它们究意是为了什么而死的呢？
我们不应该归罪于囚禁，即便是在野外，也会发生同样的事情。在十二月末的时候，
我曾到野外去观察过很多的蜂巢，都曾经发生过同样的情况。大量数目的黄蜂，必须要
死亡，这并不是因为碰到了什么意外情况，也并不是因为疾病的担扰，或是因为某种气
候的摧残影响，而是由于一种不可逃脱的命运，这种命运摧残着它们，这和鼓舞着它们
生活下去的力量是一样有力的。不过，它们这样的生命，对于我们人类倒是很有好处的。
一只母黄蜂可以创造出一个拥有三万居民的城市。假如全体黄蜂都存活下来，那么，可
想而知，这将是一场多么大的灾难啊！若是那样的话，黄蜂就可以在野外构造自己的王
国，并且称王施虐了。
到了后期，蜂巢自己会毁灭的。一种将来会变成形状平庸的蛾子的毛虫，一种赤色
的小甲虫，还有一种身着鳞状的金丝绒外衣的小幼虫，它们都是有可能攻击，毁灭蜂巢
的小动物。它们会利用锋利的牙齿，咬碎一层层小巢的地板，使得整个蜂巢内的所有住
房全部崩塌毁坏。最后，剩下来的只有几把尘土和几片棕色的纸片。到了第二年春天到
来的时候，黄蜂们便又可以废物利用，白手起家，发挥大自然在建筑房屋方面赋予它们
的高度的灵性和悟性，建造起属于它们自己的新家园。新的结构精巧而且十分坚固的城
池，其中居住着约有三万居民——一个庞大的家族。它们将一切从零开始。它们将继续
繁衍后代，喂养小宝宝，继续抵御外来的侵略，与大自然抗争，为自己的安全而战斗，
为蜂巢内部生活的快乐而贡献自己的一份力量。生命不息奋斗不止！
------------------








蛴螬的冒险



一、蜂螨
围绕着卡本托拉斯(Carpenras)乡下沙土地的高堤一带，是黄蜂和蜜蜂最喜欢光临的
地方了。它们为什么会如此喜欢这个地方呢？究其原因，主要是因为这一地区的阳光非
常充足，而且这一带还非常容易开凿，很适合黄蜂和蜜蜂在这里安居乐业。在五月份这
样的天气，主要有两种蜜蜂特别的多。它们都是泥水匠蜂，是地下的一个个小屋的建造
者。其中的一种蜜蜂，它们在自己的住宅门口，建筑起一道自认为固若金汤的防御用的
壁垒——一个土筒。它的里面留有空白，而且整个筒是呈弧形的。筒的长和宽就像人的
一个手指头一样。有时候，会有很多蜜蜂飞到这一带来定居，当它们发现了这一个个斜
形的土手指的装饰以后，谁都会感到奇怪，不知道这是一些什么东西。
还有另外一种蜜蜂，也就是我们大家经常能够见到的，它们的名字叫做掘地蜂。它
们的走廊的外口没有什么手指形的防御壁垒，而是直接暴露在外面的。旧墙的石头之间
的缝隙之中，废弃的房舍，或者是沙石上头显露的表面，这些地方都非常适合于掘地蜂
的工作。但是，最最理想，最最适宜的地方，它们常常成群结伙地奔赴的地方，要算是
那些地面上突起的，朝着南方的直路。因为我经常可以看到它们开凿低路的处所。
这里的面积有好多码宽呢，而且墙上常常穿有很多很多的小孔，以致于这块地看起
来，呈海绵的形状。这些小小的洞孔，大概是用锥子戳出来的，因为它们是那么的整齐，
每一个孔穴都与盘曲的走廊相通相连，差不多都有四五寸深。蜂巢是在这底下的。如果
我们打算观察一下这种蜜蜂的工作情况，那么我们一定要在五月的下半月到它们的工作
场上来，但是千万要注意，必须保持一定的距离，要离开相当的远。这主要是出于安全
考虑。于是，我们会发现它们一群一群地汇合在一起，喧哗着，并且，众蜂齐努力，以
一种让人惊讶的毅力，从事着关于食物和蜂巢的各项工作。
但是，我来到这个被掘地蜂占领了的地方，次数最多的，要算是在八、九月间了，
这个时候正好是快乐、自由的夏天休假的时期。在这样的季节里，在靠近掘地蜂窠巢的
地方，都显得非常宁静。一切的工作都早已进行完毕了。在缝隙之中，有很多的蜘蛛拥
挤地呆在那里面，或者有丝管子伸入到蜜蜂的走廊里。从前住满了蜂，到处都熙熙攘攘、
热热闹闹的，现在仿佛变成了凄惨、悲凉的废墟一般。这其中的理由，我们谁也无从知
晓。距离大地表面约有数寸深的下面，有成千的幼虫被封闭在它们的土室之中。它们全
都静静地等候着春天的来临。当然，这些如此柔弱而又不能自己起来保护自己的幼虫，
是这样的肥胖，这样的有吸引力。一定足以引诱某种寄生者，或者招来某种正处在饥饿
之中到处寻觅食品的外来昆虫。这件事是非常值得研究注意一下的。
有两个事实，引起了我的留心。有一些非常丑陋的苍蝇，它们身上的颜色是半黑半
白的，这些苍蝇慢慢地从一个洞穴飞到另外一个洞穴里。它们这样飞来飞去的目的是要
表明它们在那些地方产卵。其中，有一些卵是挂在网上的，都早已干枯而死了。而在其
它的地方，比如，在堤上的蜘蛛网上，也挂了许多某种甲虫——蜂螨的尸体。在这些尸
体当中，有雌的也有雄的。不过，仍然还有少数是有生命的。雌性的甲虫，一定是伸入
到了蜂的住宅里面，而且，毫无疑问，它们一定是在蜂的窠巢中产下自己的卵。
如果我们慢慢地，稍稍掘开堤的表面，我们就会惊奇地发现更多有趣的东西。在八
月之初的时候，我们看到的是：顶上有一层的小房间，它们的样子和底下的蜂巢相比，
大不一样，相差甚远。之所以有这种区别，主要是因为这是由两种不一样的蜂建造而成
的。其中有一种是已经在前面提到过的掘地蜂，另外一种，有一个很动听的名字，叫竹
蜂。
掘地蜂组成了一只先锋队。挖掘地道的工作完全由它们负责承包下来。它们懂得，
必须选择适宜的地方来建造它们自己的住所。然后，如果它们无论是因为什么样的事情
发生，都会离开它们辛苦建筑起来的外部的小房间，那么，竹蜂就会紧随它们之后跑进
来，占据这一方难得的宝地。竹蜂就将利用很粗糙的土壁，把走廊分割成大小并不完全
相等的、毫无艺术特色的许多的小房间。这便是它们所能设计出的惟一的建筑构思了。
由此可见，它们是多么投机取巧，而且还很缺少艺术的灵感。
掘地蜂建造的窠巢，却做得非常整洁，而且还进行了非常精心别致的粉饰和装修。
所以，我们可以认为它们从事的工作是颇具艺术性的，它们自身具有高超的艺术创造才
能。它们很会利用适当的土壤，把窠巢构造得连任何一个普通的敌害都无法轻易地入侵。
也正是因为这个原因，这种蜜蜂的幼虫是不会做茧的。它们只是“赤身裸体”地躺在温
暖的小房间中享福，那里面光滑得如同粉饰过一样。
然而，竹蜂的小房间里却不一样了。那里是需要一定的东西来加以保护的。原因就
在于，竹蜂的窠巢是建筑在土壤的表面上的，做得非常的草率肤浅，而且只有相当薄的
墙壁做保垒。因此，和掘地蜂的幼虫不同，竹蜂的幼虫是包在非常坚固的厚厚的虫茧里
的。这样一来，有两点好处。一方面，厚厚的茧可以保护幼虫不致于和草率而建的巢里
的墙壁相碰撞而受到不必要的伤害。另一方面，也可以使得小幼虫能躲过闯进来的仇敌
的爪牙，不致于还在襁褓之中，就遭到不测而夭折离世。
在这样的堤上，居住着两种不同的蜜蜂。我们很容易就可以分辨出哪一种蜂巢属于
哪一种蜜蜂。很显然，在掘地蜂的窠巢里，隐藏着“一丝不挂”的赤裸小幼虫；而在竹
蜂的窠巢中，则有用坚实的茧包裹着的小幼虫。
同时，这两种不同的蜜蜂，都各自有它们特殊的寄生者，或是不速之客。竹蜂的寄
生者，是那种身上黑白相间的蝇。总是能够在蜂巢隧道的门口发现这种蝇。它们闯进到
窠巢中，然后产下一些自己的卵。掘地蜂的寄生者是蜂螨。我们经常可以在堤面上发现
很多这种甲虫的尸首。
如果我把竹蜂的小房间拿开，我便可以观察到掘地蜂的家了。在一些小房间中居住
着正在成长之中的昆虫。还有一部分小房间中，住满了掘地蜂的幼虫。也有一些小房间
中，大多数是藏着一个蛋形的壳。这种壳分成了好几节，上面还有突出来的呼吸孔。这
种壳特别的薄，而且还很脆，非常易碎。它的颜色是琥珀色的，非常透明。因此，从外
边看，可以很清楚地看到，里面有一个已经发育完全的蜂螨在挣扎着，好像极其渴望自
由，希望能早日从里面解放出来。
那么，这个很奇特的壳到底是个什么东西呢？看起来，它并不太像某一种甲虫的壳。
这个寄生者，是怎样来到这个蜂巢里面的呢？
从它的地理位置上来看，简直是不能侵入的。而且，若是使用放大镜进行仔细观察，
也看不出什么受过伤害的痕迹。经过了三年之久的周密而细致的观察，我终于寻找到了
这些问题的答案。于是，在我记录的昆虫的生活史上，便又增加了最为奇怪有趣的一页。
下面就是我所研究出来的一些结果。
蜂螨，即便是在它发育完整的时期内，也只不过有一两天的寿命而已，它的全部生
命，是在掘地蜂的门口度过的。而这短暂的生命，除去要繁殖子孙后代以外；其余的什
么也没有了。
蜂螨也具备其它的动物所有的消化器官，但是，它究竟要不要吃食物呢？我表示怀
疑。对于雌甲虫而言，它惟一的愿望，便是要产下它的小宝宝。等这件大事做完以后，
它便寿终正寝，放心地离开这个世界了。那么，雄性又怎么样呢？它们要在这种土穴上
伏上一两天之后，也同样命归九泉了。这也就为我们提供了一个问题的答案，为什么在
蜂的住宅旁边的那片蜘蛛网上，悬挂着那么多莫名其妙的尸首。现在可知道它们的来源
了。
人们乍一看，肯定都会以为这种甲虫在它产卵的时候，一定是要一个小房间一个小
房间地全都跑遍，在每一个蜜蜂的幼虫身上，都要产下一个卵。可是，事实并非如此，
在我观察的过程中，我曾经在蜜蜂的隧道里面仔仔细细地搜寻过，最后发现，蜂螨只将
所有的卵产在蜂巢的门口里边，积累成一堆，距离门口差不多有一到两寸远的地方。这
些卵全部都是白颜色的，其形状呈蛋形。它们的体积都很小，互相之间轻轻地粘连在一
起。至于它们到底有多少数目，暂时算它们一共有两千多个吧，我觉得这个数目还不能
算是过高的估计。
这一事实和任何人对于它必然要想象到的相反，它们并不把卵产在蜂巢的里面，而
仅仅将它们产在蜜蜂住宅的门口之内，而且还堆成一小堆。不只如此，它们的母亲也不
储备布置一些起保护作用的东西，也不考虑为它们防御冬天的寒冷，也不替它们关上这
扇进出孔道的大门，以便抵御前来侵略、攻击、伤害它们的成千上万的敌人。总之它们
的母亲在把它们产下之后，便抛开它们，让它们独自在这个世界上闯荡。这其中的原因
是：在冬日的严寒还不曾到来之前，这条开着口的隧道，便遭受着蜘蛛及其它更为凶悍
的侵略者们的践踏，被它们所侵占，因而那些可怜的卵也就成了侵犯者的可口美餐了。
为了能看得更清楚仔细一些，我把若干的卵放在一个盒子里面。大约到了九月，它
们还没有孵化出来的时候，我想象着，它们会立刻就跑开去，到处寻找掘地蜂的小房间。
然而，事实告诉我，我完全估计错了。这一群幼小的蛴螬——小小的黑色动物，还不到
一寸的二十五分之一长——虽然它们拥有强有力的健壮的腿，但竟然利用不上。它们并
不跑散开，而是非常混乱地相处在一起，和脱下来的卵壳混杂在一起生活。于是，我在
它们面前故意悄悄地放了一块带有蜂巢的土块，想看看它们会采取怎样的行动，可是结
果却无济于事，一点儿也不能诱惑这些小动物移动一丝一毫。要是我采取行动，强行把
其中的几个挪开一些，它们便会立即又跑回去，继续躲在其它的同伴里面，和它们混居
生活在一起。
最后，在冬天的时候，我跑到了卡本托拉斯的野外，到那里去观察掘地蜂居住的地
方。我想通过这样的观察来确定是否在自然的状态之下，蜂螨的幼虫们仍然也同样是如
此，即在孵化之后，不分散开居住，而是混杂地在一起生活。到底是不是一样的呢？是
的，完全一样。在野外的情况与我的小盒子里的情况真的一模一样。我看到那些在野外
的蜂螨的幼虫也同样是累积成一堆，并且也是和它们的卵壳混住在一起的。
到现在为止，我还不能回答这样的疑问：蜂螨究竟是怎么进到蜜蜂的小房间里面来
的呢？还有它们又是怎样走进另一种并不属于自己的壳里去的呢？

二、第一次的冒险
观察过幼小的蜂螨的外表以后，我便立刻就能感觉到，它们的生活习性一定是非常
特殊的，也一定挺有意思的。
经过仔细的观察，我发现很难使得蜂螨在很一般的平面上边轻轻移动一下。在蜂螨
的幼虫所居住的地方，很显然，它们要冒着一定的跌落下去的危险。怎样才能防备这种
危险的事情发生呢？这个问题对于蜂螨幼虫而言，是轻而易举地就能解决掉的。因为它
们天生就长着一对非常强有力的大腮，弯曲而且尖利；它们还生有强壮的腿，以及成为
能够活动的爪；还生长有很多的硬毛和尖尖的针；并且，它们生来就有一对坚硬的长钉，
有着锋利而且很坚硬的尖子，其形状和样子都很像一种犁头，它可以牢牢地刺入任何光
滑的土里。还不止这些呢。除了上面提到的这些器官以外，它还可以吐出一种粘性很强
的液汁，即便是没有其它的任何东西存在，单单是这种液汁，也是可以把它紧紧地粘住，
不致于滑下去的。可见，它的自我保护能力是很强的。
我曾经一再绞尽脑汁，具思苦想着一个问题，究竟存在什么样的因素，使得这些幼
小的蛴螬决定要居住在这里呢？可是，我怎么也想不出答案来。于是，我便只有非常急
切地等待着大自然的气候能够尽快转暖一些，以便能很快地找到答案。
到了四月底的时候，被我禁闭在牢笼中的蛴螬，以前本来一直是躺着不动的，躲避
在像海绵一样的卵壳堆里边睡觉。现在则不同了，它们忽然活动起来了。开始时，它们
在度过严冬的盒子里，到处爬走着。它们急急匆匆的动作，以及它们那不知疲倦的精气
神儿，都表明它们似乎正在寻觅一些东西一些它们急切需要的东西。看起来，这些东西
自然是它们的食物了。因为，这些蜂螨的幼虫们是在九月底进行孵化的，一直到现在，
四月底。虽然它们总是处在麻木不仁的状态之中，但是差不多足足有七个月的时间，没
有获取一点儿有营养的东西来强壮身体。从孵化的时候开始，虽然这些小动物是具有生
命的“精灵”，但是它们就像是注定被判了七个月的徒刑一样，什么事情也不能做，只
能保持着一种姿势。
与此同时，当我看到它们一个个如此兴奋充满激情的时候，我便自然而然地猜想到，
驱使这些有生命的小动物如此忙忙碌碌地工作的原因，一定是饥饿，也只有饥饿才能让
它们本能地不辞辛苦。
这些匆忙寻找食物的小动物，它们真正需要的食物不过是蜂巢中的储藏品罢了。有
什么理由可以这样说呢？因为到了后期的时候，我们是在这些蜂巢中找到那些蜂螨的。
现在这些储藏品不仅限于蜜蜂的幼虫食用，也供蜂螨们分享了。
我所提供给它们的，是里面藏着蜜蜂幼虫的蜂巢。我甚至把蜂螨直接地放到蜂巢里
边去。总之，我利用各种东西，采用各种方法，希望能引起它们的食欲。但是，事实上，
我的努力仍然是一点儿结果也没有。于是，我故意用了一种方法，利用蜂蜜进行试探。
为了能够找到储藏着蜜汁的蜂巢，我花去了五月份的大部分时间。
在找到了我所希望的蜂巢以后，我把其中蜜蜂的幼虫拿了出去，然后再把蜂螨的幼
虫放到蜂巢中储备的蜂蜜里。然而事实再一次打击了我。简直没有任何其它的试验比这
一次失败得更为严重了。幼虫们根本就不去饮食那些蜜汁，更糟糕的是，它们反而被这
种粘性的东西粘住了，以至于在其中被闷死了，这太出乎我的意料了！
于是，我非常失望地说道：“我提供给你们现成的蜂窠，幼虫，还有蜜汁，难道这
些还不足够吗！你们还不满足吗？那么，你们这些丑恶的小东西，到底需要些什么东西
呢？”
不过，最终还是让我发现了它们真正需要的东西了。原来，它们并不需要什么特殊
的东西。它们是要掘地蜂亲自把它们带到蜂巢里边去。
我在前面早就已经提到过，当四月来临的时候，在蜂巢的门内居住的一堆幼虫，已
经开始表现出一点儿活动的迹象了，它们蠢蠢欲动。仅仅几天以后，它们便已经不在那
个地方停留了。真是非常怪异的小动物。它们牢牢地，死不放手地攀附在蜜蜂的毛上，
于是，便被带到了野外去，甚至已经被带到很遥远的地方了。
当掘地蜂经过蜂巢的门口的时候，无论它是要出远门，还是刚从远游中归来，睡在
门口，已经等待许久的蜂螨的幼虫，便会立刻爬到蜜蜂的身上去。它们爬进掘地蜂的绒
毛里面，抓得十分紧，无论这只掘地蜂要飞到多么遥远的地方去，它们一点儿也不担心
自己有跌落到地上去的危险。因为它们抓得太紧了。之所以要采用这样的方法，它们惟
一的目的就是想借助蜜蜂强壮的身体，将它们带到那些储有丰富的蜜的巢里去。
当一个人第一次发现这种情形的时候，他一定会以为这种喜欢冒风险的小幼虫，可
能要在蜜蜂的身上先寻觅到一些食物。但是，事实并非如此，蜂螨的幼虫伏卧在蜂的细
细的绒毛里面，和蜜蜂的身体是呈直角状的。它的头朝向里面，尾巴却朝向外面，呆在
靠近蜜蜂肩头的地方。它们选择好了地点以后，便不再随便地移来移去。如果，它们真
的打算在蜜蜂的身上吃些什么东西的话，那么它们应该是这里那里地到处跑动，寻找一
下哪一部分是最鲜嫩的。然而，事实并不是这样。它们总是要固着在蜜蜂身上，并且是
最硬的那部分，一般是在靠近蜜蜂翅膀下面的部位，有时也附着在蜜蜂的头上。它们攀
住一根毛以后，纹丝不动。因此，在我看来，事实总归是事实，是不以人的意志为转移
的。这些小甲虫之所以如此附着在蜜蜂的身体上，它们的目的仅仅是打算让蜜蜂把它们
带到即将要建造起来的蜂巢里去。
不过，在飞行的时候，这位未来的寄生者，必须要紧紧抓牢它的主人的毛才行。无
论蜜蜂是正在花叶中穿梭飞行时，如何地急速，还是它在向窠巢里飞的时候，如何地摩
擦，甚至无论它是正在用足清洁它的身体的时候，小幼虫都必须抓得很紧才行，这样才
能确保安全。
不久以前，我们曾经怀疑过，究竟是什么东西，可以使蜂螨的幼虫依附在蜜蜂的身
上呢？现在已经知道答案了，那便是生长在蜜蜂身上的绒毛。
现在，我们可以知道长在蜂螨身上的那两根大钉有什么用途了。它们合拢起来，便
可以紧紧握住蜜蜂身上的毛，比起那些最精密的人工的钳子来，还要更加精密得多。
同时，我们也可以知道那些粘液的价值了。它能帮助这个小动物更加牢固地伏在蜜
蜂的身上。而且，我们也可以了解幼虫足上长着的尖针和硬毛的作用了。它们都是用来
插入蜜蜂的软毛里，使得它本身的地位更加稳固的。
我们越是想这些当幼虫爬在平面上时似乎毫无用武之地的“组成设备”，越是会感
觉到对于这些“机器”的惊异。当这个柔弱的小动物，在它冒着危险去周游大干世界的
时候，竟然能够利用如此多的器具，防止它从蜜蜂的身上跌落下来，是多么的奇妙啊！

三、第二次的冒险
在五月二十一日这一天，我到卡本托拉斯去，想看一看蜂螨进入蜂巢时的门路。
这件工作很不容易做，需要用尽全力去完成。在野外广大的地面上，有一群蜜蜂，
像是受了日光的刺激似的，正在那边疯狂乱舞着。就在我正用缭乱的眼光观察它们的动
作时，忽然在狂乱的蜂群中间响起了一种单调而可伯的喧哗声。然后，就像闪电一样迅
猛，掘地蜂飞身而起，到处去寻找食物。与此同时，另外一群成千的蜜蜂正飞回家来。
它们身上或是携带着已经采好了的蜜汁，或者带回了用来建造它们的蜂巢用的泥土。
在那个时候，我已经具备了一些关于这类昆虫的知识，了解了一些有关它们的习性。
我以为，无论是谁有意无意地闯入它们的群里，或者只是轻轻地碰一碰它们的住宅，那
么它马上就会遭到成千锥子的狂刺而身亡。有一次，我去观察大黄蜂的蜂房，由于距离
太近了，立刻就起了一阵恐惧的颤抖，那种感觉真的是一辈子也忘不了！
然而，不管如何困难，要想知道我所渴望知道的事情，我就必须进入到这种可怕的
蜂群里，而且，必须在那里站上几个钟头，必要的时候，甚至是一整天的时间。我必须
盯着它们的工作，目不转睛，把放大镜拿在手上，站在它们当中一动也不动，观察着蜂
巢里将会发生什么样的事情。与此同时，面套、手套等等，以及其它各类遮盖保护的东
西全都不能使用，其原因是，我的手指和眼睛一定要完全不受任何防碍。其余的一切都
不管，即便是我离开蜂巢时，脸上被刺得都快让人认不出来了，也不能带各种遮盖的东
西。
那一天，我决定要解决那个困扰了我很长时间的问题了。
我用我的网子捉住了几只掘地蜂，这竟然让我感到十分的满意。因为这几只蜂的身
上都栖息着蜂螨的幼虫，这也正如我一直所希望的那样。
我先把衣服扣紧，然后，突入这群蜜蜂的中心。我拿了锄头锄了几下，然后取下一
块泥。令我感到非常诧异的，就是我居然一点儿也没有受到攻击和伤害。
在我第二次开始的时候，花的时间比第一次还要更长一些，但是，仍然是同样的结
果。我并没有受一点儿伤，也并没有一个蜜蜂利用它的尖针来刺击我一下。这以后，我
也就没什么可担心害怕的了。于是，我就大胆地长时间地停留在蜂巢前面，揭起土块，
拿出里面的蜂蜜，赶走其中的蜜蜂。在这一过程中，始终没有引起比那种喧哗更为可怕
的事情发生。为什么呢？这主要是因为掘地蜂是一种比较爱好和平的动物。每当它们的
巢内部被扰乱的时候，它们便会马上离开老家，转移地点，躲避到其它的地方去。即便
是有的时候受了一点儿伤，它们也不会使用它们的尖针，只有当它们被人捉住的时候，
它们才会用一下。
不过我得向这个泥水匠蜂的缺乏勇气表示谢意。虽然我并没有进行一点儿防御，但
是，居然能够在这些喧闹的蜂群中，在一块石头上，安安静静地坐着，并且还能随意地
观察它们的巢达到几个小时之久，却并没有被刺击过一针。这时，有一些乡下人经过这
个地方，看见我居然很安静地坐在蜂群之中。于是，他们便问我是不是对它们施加了什
么魔法。
就是这样，我观察了很多的蜜蜂蜂巢。其中有些蜂巢还是敞开着的，里面多少储备
着一些蜜汁。还有一些蜂巢已经用土掩盖了起来。而里面的东西，是大不一样的。有的
时候，我看到的是蜜蜂的幼虫；有的时候，我又看到其它种类的稍为肥大的幼虫；另外
一些时候，还能看见一个卵飘浮在蜜汁的表面上。这个卵呈非常美丽的白颜色，它是圆
柱形的，而且稍微有一点儿弯曲，差不多有一寸的五分之二或者六分之一长，这就是掘
地蜂的卵。
在少数的小房间中，我看到这种虫卵浮在蜂蜜的表面上，在其它的许多小房间中，
我看到的更多的就是幼小的蜂螨的蛴螬，它们伏卧在蜜蜂的卵上，就好像是伏在一种木
筏上一样。它的形状和大小都和刚刚孵化出来的时候是一样的，在这个蜜蜂巢里，敌人
已经卧在家门口了。
它是在什么时候并且是用什么方法进去的呢？在很多小房间中，经过仔细观察，我
简直找不出一点儿它们可以闯进去的缝隙。因为这些小房间全都封闭得很严密；所以我
推测这位寄生者一定是在储藏蜂蜜的仓库还没有关上大门之前就进去了。可是另一方面，
我看到门没有关上的小房间里面藏满了蜂蜜，但是没有发现有卵浮在那上面，也从没有
蜂螨的幼虫在里面留宿。所以，这些幼虫一定是在蜜蜂产卵的时候，或者是后来蜜蜂封
门的时候进去的。我凭借我的一点儿经验断定，幼虫在进入小房间的时候，一定是在蜜
蜂产卵在蜜上时的一瞬间。
如果我拿了一个里面装满了蜂蜜，表面上还浮着一个卵的小房间来；然后再拿上几
只蜂螨的幼虫，把它们一起放到玻璃罩里面进行观察。然而它们却很少会跑到蜂巢里边
去，它们也不能够安然地跑到“木筏”上边去！围绕着这个“木筏”的蜂蜜看来对它们
而言是太危险了。即使有那么一两只幼虫碰巧跑近了这个蜜湖，那么它们一看到这粘性
很大的东西，或者稍一涉足其中，便马上会千方百计地设法逃离这个危险的地方。可是，
经常有一些不太走运的幼虫，一不小心，就跌落到蜂的窠巢里面，很不幸运地被闷死。
所以，我们根据这一点便可以断定，蜂螨的幼虫是绝对不会离开蜜蜂的毛的，特别是在
蜜蜂呆在小房间里或靠近小房间的时候，更要牢牢地依附于蜜蜂的身体。这是因为，只
要和蜂蜜的表面稍有一点儿接触，就可以使得柔弱的小幼虫窒息而死亡。
我们必须记清楚这一点，那就是幼小的蜂螨是在封闭的小房间中发现的，而且它们
一定是呆在蜜蜂的卵上面的。这个小小的卵，不仅可以给这个小小的动物当做一个安全
的木筏，以便它能漂浮在这个可怕的、不可信任的蜜湖中，而且，这个小小的卵也就成
为了幼虫的第一顿美味佳看。
然而，要想到达这只漂浮在蜜湖中心的，而且又将成为它的食品的木筏，这只蜂螨
的幼小的蛴螬必须要避免与蜂蜜接触。否则的话后果会不堪设想的。要想达到目的，做
成这件事情，只有一个可以选择的方法。这个聪明的小幼虫，趁着蜜蜂还在产卵的空当
儿，从它的身上迅速地一下子滑落到了那个卵上。这样一来，目的达到了。幼虫便和卵
一起作伴，共同浮在蜜上了。因为这只由蜜蜂产下的卵太小了，不能同时乘载超过一个
以上的幼虫。因此，我们在一个蜂室里面，只能看到一个蜂螨的幼虫。
蜂螨幼虫的这种动作，在我们人类看来，好像是异常有灵性似的。但是，如果我们
再继续研究昆虫的话，那么，它们还将为我们提供更多的这样有灵感的例子。
可以说，当蜜蜂产下卵以后，把卵放在蜜汁上时，同时也就把它们的小天敌——蜂
螨的幼虫一起放到了小房间里面，然后，蜜蜂要非常仔细地用土把小房间的门给密封起
来。于是，一切需要它做的工作都做完了。然后，第二个小房间是做在第一个小房间的
旁边的。大概也要经历和前面相同的过程。就照这样，不停地继续下去。一直要继续到
隐蔽在蜜蜂的绒毛中的寄生者统统安居下来，这才算结束。
现在，让我们抛开这个实际上有些苦恼的母亲，不管它在做着怎样无结果的工作，
把我们的注意力转移一下，转到这些用聪明的方法最终得到膳宿的蜂螨的幼虫身上，看
看它对我们的试验会做出什么样的反应。
让我们想一想看，如果我们一旦将一只有蜂螨的幼虫在小房间上面的盖子拿下来，
那么会有什么样的事情将要发生呢？
卵还是十分完好的，一点儿也没有受到破坏。但是，好景不长。不久以后，蜂螨幼
虫的破坏工作便开始了。我们可以观察到，幼虫朝着一个长有小黑点的白卵上跑去。最
后，它忽然停下来了，由于它有六只脚，所以身体可以停得很稳。然后，它利用长在自
己身上的大腮的尖钩咬住了那个卵身上的薄皮儿，用尽自己浑身的力气，猛烈地拉扯着，
直到那个卵被它拉破为止。于是，卵里面的东西便都流了出来。那只得胜的幼虫，见了
这种东西非常满意，立刻高兴地把它享用光。这个小小的寄生虫，一生第一次使用它的
大腿的尖钩，原来是在拉破蜂卵的时候。
蜂螨的幼虫真是天生聪明啊！想出这样奇妙的方法来。利用了这种巧妙的方法，幼
虫便可以在它寄生的蜂巢的小房间中为所欲为，毫无顾及了。它可以任意地享用蜜汁。
这是因为，在蜜蜂的幼虫在孵化过程中，也是需要蜜汁来增加营养的但是，在孵化中，
蜂卵吸收的一点点东西，是不能在日后提供给两者一起享用的。因此只要蜂螨的幼虫在
拉扯卵皮的时候，越快越用力就越好。这样一来，“僧多粥少”的这个困难就不存在了。
蜂螨的幼虫之所以要破坏蜜蜂的卵，还有另外一个重要的原因。那就是因为蜂卵有
一种特殊的滋味。这种很吸引蜂螨幼虫的滋味，驱使着小幼虫在它的第一餐就要享用这
个香甜可口的小卵。这个小幼虫，在把卵撕破的初期，饮食到的是从卵里流出来的诱人
的浆汁。接连过了好几天以后，幼虫继续加油努力，把卵的裂口撕扯得更大一些这样一
来，幼虫就可以得寸进尺了。它将继续享受卵内部的流质，直到满足为止。
在幼虫吸食蜜蜂卵的过程中，储备在蜂卵周围的甜美的蜜汁，却一点儿也诱惑不了
贪吃的蜂螨幼虫，它理都不理睬一下，也不去碰它们一下。因此，可以这样讲，蜜蜂的
卵对于蜂螨幼虫而言，是绝对重要的，它是幼虫的必需食品。因而小小的蜂卵，不仅仅
可以当作蜂螨幼虫的一叶小舟，使得它在蜜湖中安全地行进，更重要的是，它还是幼虫
相当有营养的食品，为幼虫的茁壮成长提供条件。
等到整整一个星期以后，这个可怜的小蜂卵已不再剩下什么东西了，只留下了一个
空空无物的干壳。一个生命就这样悄悄地结束了。而到了这个时候，蜂螨幼虫的第一顿
大餐也已经享用完毕。幼虫也茁壮成长，差不多有原来的两倍大了。它的形状也发生了
一些不大不小的变化。它的背部裂开了，形成了自己的第二种形状，长成了一只简单的
甲虫。小幼虫从那个裂缝中解脱出来，然后落到蜂蜜上。从它身上脱下来的那个壳，还
依然停留在原来的那个小“木筏”上面。但是，在不久以后，它们都被掩没在蜜浪之中
了。
此时此刻，蜂螨幼虫的历史便画上了一个圆满的句号！
------------------








蟋蟀



一、家政
居住在草地里的蟋蟀，差不多和蝉是一样有名气的。它们在有数的几种模范式的昆
虫中，表现是相当不错的。它之所以如此名声在外，主要是因为它的住所，还有它出色
的歌唱才华。只占有这其中的一项，是不足以让它们成就如此大的名气的。一位动物故
事学家拉封丹，对于它只谈了简单的几句，仿佛并没有注意到这种小动物的天才与名气。
另外，还有一位法国寓言作家曾经写过一篇关于蟋蟀的寓言故事，但是很可惜，太
缺乏真实性和含蓄一些的幽默感。而且，这位寓言作家在这个蟋蟀的故事中写到：蟋蟀
并不满意，在叹息它自己的命运！事实可以证明，这是一个多么错误的观点。因为，无
论是什么样的人，只要曾经亲自研究过蟋蟀，观察过它们的生活情况，哪怕仅仅是一点
表面上的观察与研究，都会感觉到蟋蟀对于自己的住所，以及它们天生的歌唱才能，是
非常满意而又愉快的。是的，这两点所给它们带来的名气真的足以让它们感到庆幸了。
在这个故事的结尾处，他承认了蟋蟀的这种满足感。他写到：

“我的舒适的小家庭，是个快乐的地方，
如果你想要快乐的生活，就隐居在这里面吧！”

在我的一位朋友所做的一首诗中，给了我另一种感觉。我觉得这首诗所要表达的更
具有真实性，更加有力地表现出蟋蟀对于生活的热爱。
下面就是我的朋友写的这首诗：

曾经有个故事是讲述动物的，
一只可怜的蟋蟀跑出来，
到它的门边，
在金黄色的阳光下取暖，
看见了一只趾高气扬的蝴蝶儿。

她飞舞着，
后面拖着那骄傲的尾巴，
半月形的兰色花纹，
轻轻快快地排成长列，
深黄的星点与黑色的长带，
骄傲的飞行者轻轻地拂过。

隐士说道：飞走吧，
整天到你们的花里去徘徊吧，
不论菊花白，
玫瑰红，
都不足与我低凹的家庭相比。

突然，
来了一阵风暴，
雨水擒住了飞行者，
她的破碎的丝绒衣服上染上了污点儿，
她的翅膀被涂满了烂泥。

蟋蟀藏匿着，
淋不到雨，
用冷静的眼睛看着，
发出歌声。
风暴的威严对于它毫不相关，
狂风暴雨从它的身边无碍地过去。

远离这世界吧！
不要过分享受它的快乐与繁华，
一个低凹的家庭，
安逸而宁静，
至少可以给你以不须忧虑的时光。

从这首诗里，我们可以认识一下可爱的蟋蟀了。
我经常可以在蟋蟀住宅的门口看到它们正在卷动着它们的触须，以便使它们的身体
的前面能够凉快一些，后面能更加暖和一些。它们一点儿也不妒嫉那些在空中翩翩起舞
的各种各样的花蝴蝶。相反的，蟋蟀反倒有些怜惜它们了。它们的那种怜悯的态度，就
好像我们常看到的一样，那种有家庭的人，能体会到有家的欢乐的人，每当讲到那些无
家可归，孤苦伶仃的人时，都会流露出一样的怜悯之情。蟋蟀也从来不诉苦、不悲观，
它一向是很乐观的、很积极向上的，它对于自己拥有的房屋，以及它的那把简单的小提
琴，都相当的满意和欣慰。从某种意义上，可以这样说，蟋蟀是个地道正宗的哲学家。
它似乎清楚地懂得世间万事的虚无漂渺，并且还能够感觉到那种躲避开盲目地、疯狂地
追求快乐的人的扰乱的好处。
对了，这样来描写我们的蟋蟀，无论如何，总应该是正确的。不过，仍然需要用几
行文字，以便把蟋蟀的优点公之于众。自从那个动物故事学家拉封丹，忽略了它们以后，
蟋蟀已经等待了很长的时间了，等待着人们对它加以描述，加以介绍，加以重视。它们
的朋友——人类忽略了它们。
对于我，一个自然学者而言，前面提到的两篇寓言中，最为重要的一点，乃是蟋蟀
的窠穴，教训便建筑这在上面。
寓言作家在诗中谈到了蟋蟀的舒适的隐居地点；而拉封丹，也赞美了它的在他看来
是低下的家庭。所以，从这一点讲，最能引起人们注意的，毫无疑问，就是蟋蟀的住宅。
它的住宅，甚至吸引了诗人的目光来观察它们，尽管他们常常很少能做注意到真正存在
的事物。
确实，在建造窠穴以及家庭方面，蟋蟀可以算是超群出众的了。在各种各样的昆虫
之中，只有蟋蟀在长大之后，拥有固定的家庭，这也算是它辛苦工作的一种报酬吧！在
一年之中最坏的时节，大多数其它种类的昆虫，都只是在一个临时的隐避所里暂且躲避
身形，躲避自然界的风风雨雨。因此，它们的隐避场所得来的方便，在放弃它的时候，
也并不会觉得可惜。
这些昆虫在很多时候，也会制造出一些让人感到惊奇的东西，以便安置它们自己的
家。比如，棉花袋子，用各种树叶制做而成的篮子，还有那种水泥制成的塔等等。有很
多的昆虫，它们长期在埋伏地点伏着，等待着时机，以捕获自己等待已久的猎物。例如，
虎甲虫。它常常挖掘出一个垂直的洞，然后，利用它自己平坦的、青铜颜色的小脑袋，
塞住它的洞口。如果一旦有其它种类的昆虫涉足到这个具有迷惑性的、诱捕它们的大门
上时，那么，虎甲虫就会立刻行动，毫不留情地掀起门的一面来捕捉它。于是，这位很
不走运的过客，就这样落入虎甲虫精心伪装起来的陷阱里，不见踪影了。
另外一个例子，是蚁狮。它会在沙子上面，做成一个倾斜的隧道。这里的牺牲者是
蚂蚁。蚂蚁一旦误入歧途，便会从这个斜坡上不由自主地滑下去，然后，马上就会被一
阵乱石击死。这条隧道中守候猎物的猎者，把颈部做成了一种石弩。
但是，上面提到的例子统统都只是一种临时性的避难所或是陷阱而已，实在不是什
么长久之计。
经过辛辛苦苦的劳作构造出的家，昆虫住在里面，无论是朝气蓬勃、生机盎然的春
天，或者是在寒风刺骨、漫天雪飘的冬令时节，都让昆虫无比地依赖，不想迁移到其它
的任何地方去居住。这样一个真正的居住之所，是为了安全以及舒适而建筑的，是从长
远的角度考虑的，而并不是像前面所提到的那样，家是为了狩猎而建的。或是所谓的
“育儿院”之类的延期行为。那么，只有蟋蟀的家是为了安全和温馨而建造的了。在一
些有阳光的草坡上，蟋蟀就是这个隐逸者的场院的所有者。正当其它的或许正在过着孤
独流浪的生活，或许是卧在露天地里，或许是埋伏在枯树叶、石头和老树的树皮底下的
昆虫正为没有一个稳定的家庭而烦恼时，蟋蟀却成了大自然中的一个拥有固定居所的优
越的居民。由此可见，它是有远见意识的。
要想做成一个稳固的住宅，并不那么简单不过，现在对于蟋蟀、兔子，最后是人类，
已经不再是什么大问题了。在与我的住地相距不太远的地方，有狐狸和獾猪的洞穴，它
们绝大部分只是由不太整齐的岩石构建而成的，而且一看就知道这些洞穴都很少被修整
过。对于这类动物而言，只要能有个洞，暂且偷生，“寒窑虽破能避风雨”也就可以了。
相比之下，兔子要比它们更聪明一些。如果，有些地方没有任何天然的洞穴可以供兔子
们居住，以便躲避外界所有的侵袭与烦扰，那么，它们就会到处寻找自己喜欢的地点进
行挖掘。
然而，蟋蟀则要比它们中的任何一位更聪明得多。在选择住所时，它常常轻视那些
偶然碰到的以天然的隐避场所为家的种类。它总是非常慎重地为自己选择一个最佳的家
庭住址。它们很愿意挑选那些排水条件优良，并且有充足而温暖的阳光照射的地方。凡
是这样的地方。都被视为佳地，要优先考虑选取。蟋蟀宁可放弃那种现成的天然而成的
洞穴，因为，这些洞都不合适，而且它们都建造得十分草率，没有安全保障。有时，其
它条件也很差。总之这种洞不是首选对象。蟋蟀要求自己的别墅每一点都必须是自己亲
手挖掘而成的，从它的大厅一直到卧室，无一例外。
除去人类以外，至今我还没有发现哪种动物的建筑技术要比蟋蟀更加高超。即便是
人类，在混合沙石与灰泥使之凝固，以及用粘土涂抹墙壁的方法尚未发明之前，也不过
是以岩洞为隐避场所，和野兽进行战斗和大自然进行博击。那么，为什么这样一种非常
特殊的本能，大自然单单把它赋予了这种动物呢？最为低下的动物，但却可以居住得非
常完美和舒适。它拥有自己的一个家，有很多被文明的人类所不知晓的优点：它拥有安
全可靠的躺避隐藏的场所；它有享受不尽的舒适感，同时，在属于它自己的家的附近地
区，谁都不可能居住下来，成为它们的邻居。除了我们人类以外，没有谁可以与蟋蟀相
比。
令人感到不解和迷惑的是，这样一种小动物，它怎么会拥有这样的才能呢？难道说，
大自然偏向它们，赐予了它们某种特别的工具吗？当然，答案是否定的。蟋蟀，它可不
是什么掘凿技术方面的一流专家。实际上，人们也仅仅是因为看到蟋蟀工作时的工具非
常柔弱，所以才对蟋蟀有这样的工作结果，建造出这样的住宅感到十分惊奇的。
那么，是不是因为蟋蟀的皮肤过于柔嫩，经不起风雨的考验，才需要这样一个稳固
的住宅呢？答案仍然是否定的。因为，在它的同类兄弟姐妹中，也有和它一样，有柔美
的、感觉十分灵敏的皮肤，但是，它们并不害怕在露天底下呆着，并不怕暴露于大自然
之中。
那么，它建筑它那平安舒适的住所的高超才能，是不是由于它的身体结构上的原因
呢？它到底有没有进行这项工作的特殊器官呢？答案又是否定的。在我住所的附近地区，
分别生活着三种不同的蟋蟀。这三种蟋蟀，无论是外表、颜色，还是身体的构造，和一
般田野里的蟋蟀是非常相像的。在开始时，刚一看到它们，经常就把它们当成田野中的
蟋蟀。然而，就是这些由一个模子刻出来的同类，竟然没有一个晓得究竟怎样才能为自
己挖掘一个安全的住所。其中，有一只身上长有斑点的蟋蟀，它只是把家安置在潮湿地
方的草堆里边；还有一只十分孤独的蟋蟀，它自个儿在园丁们翻土时弄起的土块上，寂
寞地跳来跳去，像一个流浪汉一样；而更有甚者，如波尔多蟋蟀，甚至毫无顾忌，毫不
恐惧地闯到了我们的屋子里来，真是不请自来的客人，不顾主人的意愿。从八月份到九
月份，它独自呆在那些既昏暗又特别寒冷的地方，小心翼翼地唱着歌。
如果再继续前面已经提到过的那些问题，将是毫无意义的。因为那些问题的答案统
统都是否定的。蟋蟀自然形成的本能，从来也不为我们提供有关答案的原因所在。如果
寄希望于从蟋蟀的体态，身体结构，或是工作时所利用的工具上来寻找答案，来解释那
些答案，同样是不可能的。长在昆虫身上的所有的东西，没有什么能够提供给我们一些
满意的解释与答案，或者是能够让我们知晓一些原因，给不了我们任何有力的帮助。
在这四种相互类似的蟋蟀中，只有一种能够挖掘洞穴。于是，我们可以得知，蟋蟀
本能的由来，我们尚不可得知。
难道会有谁不晓得蟋蟀的家吗？哪一个人在他还是小孩子的时候，没有到过这位隐
士的房屋之前去观察过呢？无论你是怎样的小心，脚步是如何地轻巧，这个小小的动物
总能发觉，总能感觉到你的来访。然后，它立刻警觉起来，并且有所反映，马上躲到更
加隐避的地方去。而当你好不容易才接近这些动物的定居地时，此时此刻，这座住宅的
门前已经是空空如也了，很让人失望。
我想，凡是有过如此经历的人，谁都会知道，如何把这些隐匿者从躲藏处诱惑出来
的方法。你可以拿起一根草，把它放到蟋蟀的洞穴里去，轻轻地转动几下。这样一来，
小蟋蟀肯定会认为地面上发生了什么事情。于是，这只已经被搔痒了，而且已经有些恼
怒了的蟋蟀，将从后面的房间跑上来。然后，停留在过道中，迟疑着，同时，鼓动着它
的细细的触须认真而警觉地打探着外面的一切动静。然后，它才渐渐地跑到有亮光的地
方来，只要这个小东西一跑到外面来，便是自投罗网，很容易就会被人捉到。因为，前
面发生的一系列事情，已经把我们这只可怜的小动物的简单的小小头脑给弄迷糊了，毕
竞它的智力水平是何等低下啊！假如这一次，小蟋蟀逃脱掉了，那么，它将会很疑虑，
很机警，时刻提高它的警惕性，不肯再轻易地冒险，从躲避的地方跑出来。在这种情况
下，就不得不选择其它的应付手段了。比如，你可以利用一杯水，把蟋蟀从洞穴中冲出
来。
想起我们的孩童时代，那个时候真的是值得人怀念与羡慕。我们跑到草地里去，到
处捉拿蟋蟀这种昆虫。捉到以后，就把它们带回家里。就把它们放在笼子里供养。采来
一些新鲜的莴苣叶子来养活它们。这真是一种莫大的童趣啊！
现在，回过头来谈谈我这里的情况吧。为了能够更好地研究它们，我到处搜寻着它
们的窠穴。孩童时代发生的事情，就仿佛昨天刚刚发生过一般。当我的另一个小同伴—
—小保罗——一个在利用草须方面，可以称为专家级的孩子，在很长时间地实施他的战
略战术之后，忽然，他十分激动而兴奋地叫起来：“我捉住它了！我捉住它了！一只可
爱的小蟋蟀！”
“动作快一点儿，”我对小保罗说道，“我这里有一个袋子。我的小战俘，你快快
跳进去吧，你可以在袋子里面安心居住。里面有充足的饮食。不过，有个条件，那就是，
你可一定不要让我们失望啊！你一定要赶快告诉我们一些事情，一些我们渴望知道而且
正在苦苦寻觅的答案。而这些事情中，需要你做的头一件便是：把你的家给我看一看。”

二、它的住屋
在那些青青的草丛之中，不注意的话，就会不为人知地隐藏着一个有一定倾斜度的
隧道。在这里，即便是下了一场滂沱的暴雨，也会立刻就干了的。这个隐蔽的隧道，最
多不过有九寸深的样子，宽度也就像人的一个手指头那样。隧道按照地形的情况和性质，
或是弯曲，或是垂直。差不多如同定律一样，总是要有一叶草把这间住屋半遮掩起来，
其作用是很明显的，如同一所罩壁一样，把进出洞穴的孔道遮避在黑暗之中。蟋蟀在出
来吃周围的青草的时候，决不会去碰一下这一片草。那微斜的门口，仔细用扫帚打扫干
净，收拾得很宽敝。这里就是它们的一座平台，每当四周的事物都很宁静的时候，蟋蟀
就会悠闲自在地聚集在这里，开始弹奏它的四弦提琴了。多么温馨的促夏消暑音乐啊！
屋子的内部并不奢华，有暴露的，但是并不粗糙的墙。房子的住户有很多空闲的时
间去修整太粗糙的地方。隧道的底部就是卧室，这里比别的地方修饰的略微精细些，并
且宽敞些。大体上说，这是个很简单的住所，非常清洁，也不潮湿，一切都符合卫生标
准。从另一方面来说，假如我们考虑到蟋蟀用来掘土的工具十分简单，那么可以说这真
是一个伟大的工程了。如果想要知道它是怎样做的，它是什么时候开始这么大的工程的，
我们一定要回溯到蟋蟀刚刚下卵的时候。
蟋蟀像黑螽斯一样，只把卵产在土里，深约四分之三寸，它把它们排列成群，总数
大约有五百到六百个。这卵真是一种惊人的机器。孵化以后，它看起来很像一只灰白色
的长瓶子，瓶顶上有一个因而整齐的孔。孔边上有一顶小帽子，像一个盖子一样。去掉
盖子的原因，并不是蛴螬在里面不停地冲撞，把盖子弄破了，而是因为有一种环绕着的
线——一种抵抗力很弱的线，它自己会自动裂开。
卵产下两个星期以后，前端出现两个大的蛴螬，是一个呆在襁褓中的蛴螬，穿着紧
紧的衣服，还不能完全辨别出来。你应当记得，螽斯也以同样的方法孵化，当它来到地
面上时，也一样穿着一件保护身体的紧紧的外衣。蟋蟀和螽斯是同类动物，虽然事实上
并不需要，但它也穿着一件同样的制服。螽斯的卵留在地下有八个月之久，它要想从地
底下出来必须同已经变硬了的土壤搏斗一番，因此需要一件长衣保护它的长腿。但是蟋
蟀整体上比较短粗，而且卵在地下也不过几天，它出来时无非只要穿过粉状的泥土就可
以了。用不着和土地相抗挣。因为这些理由，它不需要外衣，于是它就把这件外衣抛弃
在后面的壳里了。
当它脱去襁褓时，蟋蟀的身体差不多完全是灰白色的，它开始和眼前的泥土战斗了。
它用它的大腮将一些毫无抵抗力的泥土咬出来，然后把它们打扫在一旁或干脆踢到后面
去，它很快就可以在土面上享受着阳光，并冒着和它的同类相冲突的危险开始生活，它
是这样弱小的一个可怜虫，还没有跳蚤大呢！
二十四小时以后，它变成了一个小黑虫，这时它的黑檀色足以和发育完全的蟋蟀相
媲美，它全部的灰白色到最后只留下来一条围绕着胸部的白肩带，它身上生有两个黑色
的点。在这两点中上面的一点，就在长瓶的头上，你可以看见一条环绕着的，薄薄的、
突起的线。壳子将来就在这条线上裂开。因为卵是透明的，我们可以看见这个小动物身
上长着的节。现在是应该注意的时候了，特别是在早上的时候。
好运气是关爱带来的，如果我们不断地到卵旁边去看，我们会得到报酬的。在突起
的线的四周，壳的抵抗力会渐渐消失，卵的一端逐渐分裂开，被里面的小动物的头部锥
动，它升起来，落在一旁，像小香水瓶的盖子一样，战俘就从瓶子里跳了出来。
当它出去以后，卵壳还是长形的，光滑、完整、洁白，馆子似的盖子挂在口上的一
端。鸡卵破裂，就是小鸡用嘴尖上的小硬瘤撞破的；蟋蟀的卵做得更加巧妙，和象牙盒
子相似，能把盖子打开。它的头顶，已经足可以做这件工作了。
我们上面说过，盖子去掉以后，一个幼小的蟋蟀跳出来，这句话还不十分精确。它
是非常灵敏和活泼的，不时用长的而且经常颤动的触须打探四周发生的情况，并且很性
急地跑来跳去。当有一天，它长胖了，不能如此放肆了，那才真有些滑稽呢！
现在我们要看一看母蟋蟀为什么要产下这么多的卵。这是因为多数的小动物是要被
处以死刑的。它们常遭到别的动物大量的残忍的大屠杀，特别是小形的灰螅蜴和蚂蚁的
杀害。蚂蚁这种讨厌的流寇，常常不留一只蟋蟀在我们的花园里。它一口就能咬住这可
怜的小动物，然后狼吞虎咽地将它们吞咽下去。
唉，这个可恨的恶人，请想想看，我们还将蚂蚁放在比较高级的昆虫当中，还为它
写了很多的书，更对它大加赞美。称赞之声，不绝于耳。自然学者对它们很推祟，而且
其名誉日益增加。这样看来，动物和人一样，引起人们注意的最绝妙的方法，就是损害
别人。
那些从事十分有益处的清洁工作的甲虫，并不能引来人们的注意与称赞，甚至无人
去理睬它们；而吃人血的蚊虫，却是每个人都知道的；同时人们也知道那些带着毒剑，
暴躁而又虚夸的黄蜂，以及专做坏事的蚂蚁。后者在我们南方的村庄中，常常会跑到人
们的家里面弄坏椽子，而且它们在做这些坏事时，还像品尝无花果一样高兴。
我花园里的蟋蟀，已经完全被蚂蚁残杀殆尽，这就使得我不得不跑到外面的地方去
寻找它们。八月里在落叶下，那里的草还没有完全被太阳晒枯干，我看到幼小的蟋蟀，
已经长得比较大了，全身已经都是黑色了，白肩带的痕迹一点也没有存留下来，在这个
时期，它的生活是流浪式的，一片枯叶，一块扁石头，已经足够它去应付大千世界中的
一些事情了。
许多从蚂蚁口中逃脱而残生的蟋蟀，现在又作了黄蜂的牺牲品。它们猎取这些旅行
者，然后把它们埋在地下。其实只要蟋蟀提前几个星期做好防护工作，它们就没有这种
危险了。但是它们从来也没想到过这点，总是死守着旧习惯，仿佛视死如归的样子。
一直要到十月末，寒气开始袭人时，蟋蟀才开始动手建造自己的巢穴。如果以我们
对养在笼子里的蟋蟀的观察来判断，这项工作是很简单的。挖穴并不在裸露的地面上进
行，而是常常在莴苣叶——残留下来的食物——掩盖的地点。或者是其它的能代替草叶
的东西，似乎为了使它的住宅秘密起见，这些掩盖物是不可缺少的。
这位矿工用它的前足扒着土地，并用大腮的钳子，咬去较大的石块。我看到它用强
有力的后足蹬踏着土地，后腿上长有两排锯齿式的东西。同时，我也看到它清扫尘土将
其推到后面，把它倾斜地铺开。这样，就可以知道蟋蟀挖掘巢穴的全部方法了。
工作开始做的很快。在我笼子里的土中，它钻在下面一呆就是两个小时，而且隔一
小会儿，它就会到进出口的地方来。但是它常常是向着后面的，不停地打扫着尘土。如
果它感到劳累了，它可以在还没完成的家门口休息一会儿，头朝着外面，触须特别无力
地摆动，一副倦怠的样子。不久它又钻进去，用钳子和耙继续劳作。后来，休息的时间
渐渐加长，这使我感到有些不耐烦了。
这项工作最重要的部分已经完成了。洞口已经有两寸多深了，足够满足一时之需。
余下的事情，可以慢慢地做，今天做一点，明天再做一点，这个洞可以随天气的变冷和
蟋蟀身体的长大而加大加深。如果冬天的天气比较暖和，太阳照射到住宅的门口，仍然
还可以看见蟋蟀从洞穴里面抛散出泥土来。在春天尽情享乐的天气里，这住宅的修理工
作仍然继续不已。改良和装饰的工作，总是经常地不停歇地在做着，直到主人死去。
四月的月底，蟋蟀开始唱歌，最初是一种生疏而又羞涩的独唱，不久，就合成在一
起形成美妙的奏乐，每块泥土都夸赞它是非常善于演奏动听的音乐的乐者。我乐意将它
置于春天的歌唱者之魁首。在我们的荒废了的土地上，在百里香和欧薄荷繁盛的开花时
节，百灵鸟如火箭般飞起来，打开喉咙纵情歌唱，将优美的歌声，从天空散布到地上。
而呆在下面的蟋蟀，它们也禁不住吸引，放声高歌一曲，以求与相知者相应和。它们的
歌声单调而又无艺术感，但它的这种艺术感和它生命复苏的单调喜悦相协调，这是一种
警醒的歌颂，为萌芽的种子和初生的叶片所了解、所体味。对于这种二人合奏的乐曲，
我们应该判定蟋蟀是优秀中的胜者。它的数目和不间断的音节足以使它当之无愧。百灵
鸟的歌声停止以后，在这些田野上，生长着青灰色的欧薄荷，这些在日光下摇摆着芳香
的批评家，仍然能够享受到这样朴实的歌唱家的一曲赞美之歌，从而伴它们度过每一刻
寂寞的时光。多么有益的伴侣啊！它给大自然以美好的回报。

三、它的乐器
为了科学的研究，我们可以很坦率地对蟋蟀说道：“把你的乐器给我们看看。”像
各种有价值的东西一样，它是非常简单的。它和螽斯的乐器很相像，根据同样的原理，
它不过是一只弓，弓上有一只钩子，以及一种振动膜。右翼鞘遮盖着左翼鞘，差不多完
全遮盖着，只除去后面和转折包在体侧的一部分，这种样式和我们原先看到的蚱蜢、螽
斯，及其同类相反。蟋蟀是右边的盖着左边的，而蚱蜢等，是左边的盖着右边的。
两个翼鞘的构造是完全一样的。知道一个也就知道另一个了。它们分别平铺在蟋蟀
的身上。在旁边，突然斜下成直角，紧裹在身上，上面还长有细脉。
如果你把两个翼鞘揭开，然后朝着亮光仔细地留意，你可以看到它是极其淡的淡红
色，除去两个连接着的地方以外，前面是一个大的三角形，后面是一个小的椭圆，上面
生长有模糊的皱纹，这两个地方就是它的发声器官，这里的皮是透明的，比其它的地方
要更加紧密些，只是略带一些烟灰色。
在前一部分的后端边隙的空隙中有五条或是六条黑色的条纹，看来好像梯子的台阶。
它们能互相磨擦，从而增加与下面弓的接触点的数目，以增强其振动。
在下面，围绕着空隙的两条脉线中的一条，呈肋状。切成钩的样子的就是弓，它长
着约一百五十个三角形的齿，整齐得几乎符合几何学的规律。
这的确可以说是一件非常精致的乐器。弓上的一百五十个齿，嵌在对面翼鞘的梯级
里面，使四个发声器同时振动，下面的一对直接摩擦，上面的一对是摆动摩擦的器具，
它只用其中的四只发音器就能将音乐传到数百码以外的地方，可以想象这声音是如何的
急促啊！
它的声音可以与蝉的清澈的鸣叫相抗衡，并且没有后者粗糙的声音。比较来说，蟋
蟀的叫声要更好一些，这是因为它知道怎样调节它的曲调。蟋蟀的翼鞘向着两个不同的
方向伸出，所以非常开阔。这就形成了制音器，如果把它放低一点，那么就能改变其发
出声音的强度。根据它们与蟋蟀柔软的身体接触程度的不同，可以让它一会儿能发出柔
和的低声的吟唱，一会儿又发出极高亢的声调。
蟋蟀身上两个翼盘完全相似，这一点是非常值得注意的。我可以清楚地看到上面弓
的作用，和四个发音地方的动作。但下面的那一个，即左翼的弓又有什么样的用处呢？
它并不被放置在任何东西上，没有东西接触着同样装饰着齿的钩子。它是完全没有用处
的，除非能将两部分器具调换一下位置，那下面的可以放到上面去。如果这件事可以办
到的话，那么它的器具的功用还是和以前相同，只不过这一次是利用它现在没有用到的
那只弓演奏了。下面的胡琴弓变成上面的，但是所演奏出来的调子还是一样的。
最初我以为蟋蟀是两只弓都是有用的，至少它们中有些是用左面那一只的。但是观
察的结果恰恰与我的想象相反。我所观察过的蟋蟀(数目很多)都是右翼鞘盖在左翼鞘上
的，没有一只例外。
我甚至用人为的方法来做这件事情。我非常轻巧地，用我的钳子，使蟋蟀的左翼鞘
放在右翼鞘上，决不碰破一点儿皮。只要有一点技巧和耐心，这件事情是容易做到的。
事情的各方面都得很好，肩上没有脱落，翼膜也没有皱褶。
我很希望蟋蟀在这种状态下仍然可以尽情歌唱，但不久我就失望了。它开始回复到
原来的状态。我一而再而三地摆弄了好几回，但是蟋蟀的顽固终于还是战胜了我的摆布。
后来我想这种试验应该在翼鞘还是新的、软的时候进行，即在蛴螬刚刚蜕去皮的时
候。我得到刚刚脱化的一只幼虫，在这个时候，它未来的翼和翼鞘形状就像四个极小的
薄片，它短小的形状和向着不同方向平铺的样子，使我想到面包师穿的那种短马甲，这
蛴螬不久就在我的面前，脱去了这层衣服。
小蟋蟀的翼鞘一点一点长大，渐渐变大，这时还看不出哪一扇翼鞘盖在上面。后来
两边接近了，再过几分钟，右边的马上就要盖到左边的上面去了。于是这时是我加以干
涉的时候了。
我用一根草轻轻地调整其鞘的位置，使左边的翼鞘盖到右边的上面。蟋蟀虽然有些
反抗，但是最终我还是成功了。左边的翼鞘稍稍推向前方，虽然只有一点点。于是我放
下它，翼鞘逐渐在变换位置的情况下长大。蟋蟀逐渐向左边发展了。我很希望它使用它
的家族从未用过的左琴弓来演奏出一曲同样美妙动人的乐曲。
第三天，它就开始了。先听到几声摩擦的声音，好像机器的齿轮还没有切合好，正
在把它调整一样。然后调子开始了，还是它那种固有的音调。
唉，我过于信任我破坏自然规律的行为了。我以为已造就了一位新式的奏乐师，然
而我一无所获。蟋蟀仍然拉它右面的琴弓，而且常常如此拉。它因拼命努力，想把我颠
倒放置的翼鞘放在原来的位置，导致肩膀脱臼，现在它已经经过自己的几番努力与挣扎，
把本来应该在上面的翼鞘又放回了原来的位置上，应该放在下面的仍放在下面。我想把
它做成左手的演奏者的方法是缺乏科学性的。它以它的行动来嘲笑我的做法，最终，它
的一生还是以右手琴师的身份度过的。
乐器已讲得够多了，让我们来欣赏一下它的音乐吧！蟋蟀是在它自家的门口唱歌的，
在温暖的阳光下面，从不躲在屋里自的欣赏。翼鞘发出“克利克利”柔和的振动声。音
调圆满，非常响亮、明朗而精美，而且延长之处仿佛无休止一样。整个春天寂寞的闲暇
就这样消遣过去了。这位隐士最初的歌唱是为了让自己过得更快乐些。它在歌颂照在它
身上的阳光，供给它食物的青草，给它居住的平安隐避之所。它的弓的第一目的，是歌
颂它生存的快乐，表达它对大自然恩赐的谢意。
到了后来，它不再以自我为中心了，它逐渐为它的伴侣而弹奏。但是据实说来，它
的这种关心并没收到感谢的回报，因为到后来它和它的伴侣争斗的很凶，除非它逃走，
否则它的伴侣会把它弄成残废，甚至将吃掉它一部分的肢体。不过无论如何，它不久总
要死的，就是它逃脱了好争斗的伴侣，在六月里它也是要死亡的。听说喜欢听音乐的希
腊人常将它养在笼子里，好听它们的歌唱。然而我不信这回事，至少是表示怀疑。第一，
它发出的略带烦嚣的声音，如果靠近听久了，耳朵是受不了的，希腊人的听觉恐怕不见
得爱听这种粗糙的，来自田野间的音乐吧！
第二，蝉是不能养在笼子里面的，除非我们连洋橄榄或榛系木一齐都罩在里面。但
是只要关一天，就会使这喜欢高飞的昆虫厌倦而死。
将蟋蟀错误地作为蝉，好像将蝉错误地当作蚱蜢一样，并不是不可能的。如果如此
形容蟋蟀，那么是有一定道理的。它被关起来是很快乐的，并不烦恼。它长住在家里的
生活使它能够被饲养，它是很容易满足的。只要它每天有莴苣叶子吃，就是关在不及拳
头大的笼子里，它也能生活的很快乐，不住地叫。雅典小孩子挂在窗口笼子里养的，不
就是它吗？
布罗温司的小孩子，以及南方各处的小孩子们，都有同样的嗜好。至于在城里，蟋
蟀更成为孩子们的珍贵财产了。这种昆虫在主人那里受到各种恩宠，享受到各种美味佳
肴。同时，它们也以自己特有的方式来回报好心的主人，为他们不时地唱起乡下的快乐
之歌。因此它的死能使全家人都感到悲哀，足可以说明它与人类的关系是多么亲密了。
我们附近的其它三种蟋蟀，都有同样的乐器，不过细微处稍有一些不同。它们的歌
唱在各方面都很像，不过它们身体的大小各有不同。波尔多蟋蟀，有时候到我家厨房的
黑暗处来，是蟋蟀一族中最小的，它的歌声也很细微，必须要侧耳静听才能听得见。
田野里的蟋蟀，在春天有太阳的时候歌唱，在夏天的晚上，我们则听到意大利蟋蟀
的声音了。它是个瘦弱的昆虫，颜色十分浅淡，差不多呈白色，似乎和它夜间行动的习
惯相吻合。如果你将它放在手指中，你就会怕把它捏扁。它喜欢呆在高高的空气中，在
各种灌木里，或者是比较高的草上，很少爬下地面来。在七月到十月这些炎热的夜晚，
它甜蜜的歌声，从太阳落山起，继续至半夜也不停止。
布罗温司的人都熟悉它的歌声，最小的灌木叶下也有它的乐队。很柔和很慢的“格
里里，格里里”的声音，加以轻微的颤音，格外有意思。如果没有什么事打扰它，这种
声将会一直持续并不改变，但是只要有一点儿声响，它就变成迷人的歌者了。你本来听
见它在你面前很靠近的地方，但是忽然你听起来，它已在十五码以外的地方了。但是如
果你向着这个声音走过去，它却并不在那里，声音还是从原来的地方传过来的。其实，
也并不是这样的。这声音是从左面，还是从后面传来的呢？一个人完全被搞糊涂了，简
直辨别不出歌声发出的地点了。
这种距离不定的幻声，是由两种方法造成的。声音的高低与抑扬，根据下翼鞘被弓
压迫的部位而不同，同时，它们也受翼鞘位置的影响。如果要发较高的声音，翼鞘就会
抬举得很高；如果要发较低的声音，翼鞘就低下来一点。淡色的蟋蟀会迷惑来捕捉它的
人，用它颤动板的边缘压住柔软的身体，以此将来者搞昏。
在我所知道的昆虫中，没有什么其它的歌声比它更动人、更清晰的了。在八月夜深
人静的晚上，可以听到它。我常常俯卧在我哈麻司里迷迭香旁边的草地上，静静地欣赏
这种悦耳的音乐。那种感觉真是十分的惬意。
意大利蟋蟀聚集在我的小花园中，在每一株开着红花的野玫瑰上，都有它的歌颂者，
欧薄荷上也有很多。野草莓树、小松树，也都变成了音乐场所。并且它的声音十分清澈，
富有美感，特别动人。所以在这个世界中，从每棵小树到每根树枝上，都飘出颂扬生存
的快乐之歌。简直就是一曲动物之中的“欢乐颂”！
高高的在我头顶上，天鹅飞翔于银河之间，而在地面上，围绕着我的，有昆虫快乐
的音乐，时起时息。微小的生命，诉说它的快乐，使我忘记了星辰的美景，我已然完全
陶醉于动听的音乐世界之中了。那些天眼，向下看着我，静静的，冷冷的，但一点也不
能打动我内在的心弦。为什么呢？因为它们缺少一个大的秘密——生命。确实，我们的
理智告诉我们：那些被太阳晒热的地方，同我们的一样，不过终究说来，这种信念也等
于一种猜想，这不是一件确实无疑的事。
在你的同伴里，相反的啊，我的蟋蟀，我感到生命的活力，这是我们土地的灵魂，
这就是为什么我不看天上的星辰，而将注意力集中于你们的夜歌的原因了。一个活着的
微点——最小最小的生命的一粒，它的快乐和痛苦，比无限大的物质，更能引起我的无
限兴趣，更让我无比地热爱你们！
------------------








娇小的赤条蜂



细细的腰，玲珑的身材，腹部分成两节，下面大，上面小，中间好像是用一根细线
连起来，黑色的肚皮上面围着一丝红色的腰带：这就是赤条蜂。
赤条蜂的巢穴是建筑在疏松的极容易钻通的泥土里的。小路的两旁，太阳照耀着的
泥滩上，那些地方的草长得很稀疏，都是黄蜂最理想的住所。在春季，四月初的时候，
我们总可以在这样的地方找到它们。
赤条蜂通常在泥土里筑一个垂直的洞，好像一口井，口径只有鹅毛管那么粗，约有
二寸深，洞底是一个孤立的小房间，专为产卵用的。黄蜂建巢的时候，总是静静地、慢
慢地工作着，丝毫没有什么热烈或兴奋的样子。像别的蜂一样，它用前足作耙，用嘴巴
作挖掘的工具。有时候我们就可以听到，从洞底发出一声尖利、刺耳的摩擦声，这是因
为它遇到了一颗极不容易搬去的沙粒引起翅膀和全身剧烈振动的缘故。每隔短短的十几
分钟，我们就可以看到赤条蜂在洞口出现，嘴里衔着一些垃圾或是一颗沙粒。它总把这
种垃圾丢到几寸以外的地方，这样可以保持自己的居所和周围环境的整洁干净。
有几颗沙粒会被区别对待。赤条蜂们会对它们进行特殊的处理，使它们免遭被远远
抛出去的命运。这些沙粒被赤条蜂们堆在洞的附近，将来会另有重大的用途。当赤条蜂
把洞完全挖好了，它就在这小沙滩上察看有没有适合它需要的沙粒。如果没有，它就到
附近去找，直到找到为止。它需要的是一粒扁平的，比它的洞口稍大一些的沙粒，它可
以把这个沙粒盖在洞口，做成一扇门。第二天它从外面猎取一条毛毛虫回来，就不慌不
忙地把门打开，把猎物拖进去。这门看起来和其它沙粒完全一样，谁也不会想到它底下
会藏着食物，藏着一只赤条蜂的家，只有它自己才能辨别出它的家。它打开门，不紧不
但地把猎物放到洞底后，就开始在上面产卵，然后再用它以前藏在附近的沙粒把洞口堵
住。这听起来似乎有点像《阿里巴巴与四十大盗》中“芝麻开门”的故事。
赤条蜂所猎取的食物是一种灰蛾的幼虫。这种虫大部分都是生活在地底下的，赤条
蜂又是怎么把它捉到的呢？让我们来看看吧。有一天当我散步回来的时候，正好看到一
只赤条蜂在一丛百里香底下忙碌着。我立刻在它附近的地上躺下。我的出现并没有把它
吓走。它先飞到我的衣袖上停了一会儿，断定我不会伤害它之后，又飞回到百里香丛中
去了。从过去的经验我知道这意味着什么：它忙得很，没有时间来考虑我这个不速之客。
赤条蜂把百里香根部的泥土挖去，又把周围的小草拔掉，然后把头钻进它挖松的土
块里。它匆匆忙忙地从这里飞到那里，向每一条裂缝里张望。而不是在为自己筑巢，而
是在寻找地底下的食物，活像一只猎狗在寻找洞里的野兔一般。
灰蛾的幼虫觉察到了上面的动静，决定离开自己的巢，爬到地面上来看看到底发生
了什么事。这一念之差就决定了它的命运。那赤条蜂是早已准备就绪，就等着它的出现
了。果然，灰蛾的幼虫一露出地面，赤条蜂就冲出去一把将它抓住了，然后伏在它的背
上，像一个熟练的外科医生一样，不慌不忙地用刺把毛虫的每一节都刺一下。从前到后
一节一节地往下刺，一点儿也不遗漏。它那熟练的动作，让人想起游刃有余的屠夫。
赤条蜂的技巧令科学家们都自愧不如。它可以靠观察去推断人类所从不知道的事情。
它很熟悉它的俘虏的神经系统，它知道往哪些神经中枢上扎刺，可以使它的俘虏神经麻
木而不致于死亡。它从哪里学到的这种知识呢？我们人类有学校，有老师，有各种书籍，
我们通过这些逐渐地积累知识懂得大自然的许多奥妙。可是赤条蜂是怎么知道这些复杂
的知识的呢？而且不用练习就掌握了那么熟练的技术，难道是在它们出生前，冥冥之中
就有神灵赋予它们这种本领了吗？大自然多神奇啊，当我们孜孜不倦地探索它的秘密时，
它早就有条不紊地安排好了一切！
我还要告诉你另一幕赤条蜂和毛毛虫大战的故事，这也是我亲眼看见的。那是在五
月里，我看着赤条蜂在一条光洁的路旁为它的巢作最后一步的清除工作。它的几码外的
地方已经麻醉好了一条毛毛虫，当它清除好那条街道并且把洞开得足够大后，它就出去
搬毛毛虫了。
它很容易地找到了那条毛毛虫，那条毛毛虫就躺在地上。可是糟糕的是，蚂蚁也正
在猎取那条毛毛虫。赤条蜂不愿意和蚂蚁分享这条毛毛虫，可是要把蚂蚁赶掉，也不是
一件容易的事，再三考虑之后，它认为自己的能力实在有限，还是不要作无谓的牺牲吧。
于是它决定放弃这条毛毛虫，再到别处去寻找食物。
它在离巢大约十尺以内的地方，一步一步慢慢地走着，察看着泥土，不时地用它那
弯着的触须，在地面上挥动，像一名执着的士兵用探雷针寻找着地雷。在烈日的下面，
我观察了它整整三个钟头！要找到一条毛毛虫是多么困难啊，尤其是在急需的时候。
即使是对人，这也是一件困难的工作。我一心要帮助它，替它找到一条毛毛虫，因
为我想看它怎样麻痹毛毛虫。
于是我就想起我的老朋友法维，他是我的园丁，正在那里照料花园，于是我把他招
呼过来。
“快来快来！我想要几条灰色毛毛虫！”我把事情飞快地向他解释了一下。他明白
了，马上去找虫子。他挖掘着莴苣的根，耙着草莓里的泥，察看着鸢尾草丛的边缘。我
非常信任他，他的眼力和聪明，因为那么多年以来，大家都认为他是一个出色的园丁。
过了好久，也没见法维拿毛毛虫过来。
“喂，法维，毛毛虫呢？”
“先生，我一条也找不着。”
“怎么会？！把你们所有的人都叫过来！克兰亚、爱格兰，你们都来！到莴苣田里
来！帮我找毛毛虫！”
于是全家都出动了，每个人都很努力，可是毫无结果。三个钟头过去了。我们中间
谁也没找着毛毛虫。
赤条蜂也没有找着毛毛虫，它已经很疲倦了，我看到它很果断地在地面上有裂缝的
地方寻找。它尽着它最大的努力寻找，甚至把杏核般大的泥块搬开。可是不久它又离开
了这些地方。于是我开始怀疑，赤条蜂捕获不到猎物，不是因为找不到毛毛虫，而是因
为虽然它知道毛毛虫在哪儿，却没办法捉到它们，可能因为毛毛虫早有防备，把巢挖得
很深，而赤条蜂没有能力把虫子从地底下挖出来。我真愚昧，为什么早没有想到这一点？
难道这样一个经验丰富的猎取家会盲目地浪费精力吗？当然不会。
此刻赤条蜂又在挖另一个地方了，可是不久它又放弃了，正像它所尝试过的许多地
方一样。我决心要帮它的忙，于是我就继续它的工作，用小刀朝那儿挖下去。可是什么
也没有，于是我也放弃了那块地方。
可是不久赤条蜂又回来了，在我挖过的地方继续往下挖。我明白它的意思了，我为
它创造了条件，重新激起它对这个地方的信心。
“滚开，你这个笨头笨脑的家伙！”赤条蜂似乎在说“让我来告诉你这里到底有没
有毛毛虫！”
于是我按照赤条蜂指引的方向挖下去，果然挖出了一条毛毛虫。太好了！聪明的赤
条蜂！你没有辜负我对你的信任！
照这种办法，我挖到了第二条毛毛虫，不久，第三条、第四条都被我挖到了。我发
现赤条蜂所挖掘的地方，都是光秃秃的几个月前曾经翻松过的泥土。除此之外，再也没
有什么记号可以显示出毛毛虫的所在地了。好了，法维、克兰亚、爱格兰和你们其余所
有的人，你们还有什么话说？你们找了三个钟头都没找到一条毛毛虫，而这只聪明的赤
条蜂，却提供给了我足够的毛毛虫，同时，我也为自己对赤条蜂的信任和了解而沾沾自
喜。是啊，我能够懂得它的心思，能够和它密切配合，互补长短，那一堆丰盛的“战果”
就是我们之间天衣无缝的完美合作的最好证明。

袭击毛毛虫
我把第五条毛毛虫留给赤条蜂，当时我正躺在地上，和这位屠夫靠得很近，所以没
有一个细小的动作能逃过我的眼睛。现在我要把我眼前所发生的情景一段一段地记下来。
1．赤条蜂用它的嘴巴夹住了毛毛虫的颈部，毛毛虫剧烈地挣扎着，扭动着身体。赤
条蜂却不慌不忙，自己让到一边，以避免剧烈的冲撞。它的刺扎在毛毛虫的头和第一节
之间的关节上，那是毛毛虫的皮最嫩的地方。这是最性命攸关的一下，这一下可以使毛
毛虫完全受赤条蜂的控制。
2．赤条蜂突然离开毛毛虫，躺倒在地上，剧烈地扭动着，不停地打着滚，抖动着足，
拍打着翅膀，像是在垂死挣扎。我以为它也被毛毛虫扎了一下，受了致命的伤。看着它
的生命就要这样结束，我对它充满了无限的同情。可是它突然又恢复正常了，扇扇翅膀，
理理须发，又活灵活现地回到猎物旁。刚才看到的那一幕，其实正是它庆祝胜利的表现，
而不是像我想的那样受了伤。
3．赤条蜂抓住了毛毛虫的背部，抓的部位比第一次稍微低些，然后开始用刺扎它身
体的第二个体节，仍旧是刺在下方。它一节一节地往下刺。头三节上有脚，接着两节没
有脚，再以后四节又有脚，不过那不是真正的脚，充其量也只能算是沈状突起物。一共
有九节。但早在赤条蜂第一针刺下后，毛毛虫已经没有多大的抵抗力了。
4．最后赤条蜂把钳子般的嘴巴张到最大的宽度，钳住了毛毛虫的头，有节奏地轻轻
压榨它，但尽量不使它受伤。每压一次，赤条蜂就要停一下，看看毛毛虫有什么反应。
这样一停、一等、一压，循环往复地进行着。这种控制大脑的手术不能做得太猛烈，否
则毛毛虫很可能会死掉。说来很奇怪，赤条蜂为什么并不想让它死掉呢？
现在外科医生的手术已经结束，毛毛虫瘫趴在地上。它不会动了，几乎没有生命，
只有一息尚存。它任凭被赤条蜂拖到洞里，不做也不能做丝毫的反抗。当赤条蜂把卵产
在它身上后，它也没有能力伤害在它身上成长的赤条蜂的幼虫。这就是赤条蜂所做的麻
醉工作的目的：它是在为未来的婴儿预备食物。它把毛毛虫拖到洞里以后，就在它身上
产一个卵。等到幼虫从卵里孵化出来，就可以把毛毛虫当作食物。想想看，如果毛毛虫
还能动弹会有什么后果？只要它轻轻一转身，就能轻而易举地把赤条蜂的卵压破！当然，
毛毛虫是不会动了。可是它又不能完全死掉，因为如果它死了，尸体很快就会腐烂，不
适宜做赤条蜂幼虫的食物了。所以赤条蜂用它的毒刺刺进毛毛虫的每一节神经中枢，使
它失去运动的能力，半死不活地苟且残喘下去，自动地为幼虫将来的食物“保鲜”。赤
条蜂想得多周到啊！不过，等你看到它把猎物拉回家的过程，你会发现它对事物考虑的
周到程度还远不止这些。它想到毛毛虫的头部还没有受伤，嘴巴还能动。当它被赤条蜂
拖着走的时候，它能够咬住地上的草，从而阻碍赤条蜂继续把它往前拖，所以赤条蜂还
得想办法把毛毛虫的头部也麻痹掉。这次它不再用它的毒刺，因为那会致毛毛虫于死地。
它连续不断地压榨和摩擦毛毛虫的头部，这种方法实在是十分恰到好处，毛毛虫很快便
失去了知觉，它被折腾晕了。
虽然我们很佩服赤条蜂的技巧，但同时也不禁为毛毛虫们捏把汗。毕竟，被赤条蜂
弄得求生不得，求死不能是一件很残忍的事，不过如果我们是农夫，我们就不会对毛毛
虫深表同情了。它是农作物和花草最可怕的敌人！它们白天弯曲着身子在洞里睡大觉，
晚上爬出来为非作歹，咬植物的根和茎。无论是观赏用的花草还是食用的蔬菜，都是它
们的美餐。有时候，一棵幼苗无缘无故地枯萎了，如果你把它轻轻地拔起来，就会发现
它的根部受了伤。原来，晚上的时候，恶贯满盈的毛毛虫到过这儿，它那剪刀一样的嘴
巴把这棵幼苗咬伤了。它和另一种白色的毛毛虫(金虫的幼虫)一样的坏。只要它到了甜
菜园里，这菜园就要蒙受巨大的损失。它的的确确是一种害虫。赤条蜂把它杀死，是在
为民除害。我们就完全没有必要对毛毛虫产生同情了。
------------------








西西斯



我希望你们听过了关于清道的甲虫做球的奇怪的事情，还不至于厌倦。我已经告诉
过你们神圣甲虫和西班牙的犀头，现在我想再讲一些这种动物的其它种类。在昆虫的世
界里，我们遇到过许多模范母亲，现在只是为了好玩，来注意一回好的父亲吧！
除非在高等动物中，好的父亲是很少见的。在这方面，鸟类是优秀的，而人类最能
尽这种义务。低级动物当中，父亲对家族中的事情是漠不关心的。很少有昆虫是这种定
律的例外。这种无情，在高级动物的世界中是要被厌恶的，而且它们幼小的动物不需要
长时间的看护。而对于昆虫的父亲来说这是可以原谅的。因为只要有个适当的地点，新
生昆虫就可以十分健康地成长，很可能无需帮助而得到食物。例如粉蝶为了种族的安全，
只要把卵产在菜叶上，父亲的责任心又有什么用呢？母亲有利用植物的本能，是不需要
帮助的。在产卵的时候，也是不需要父亲在一边保护的。
许多昆虫都采用一种简单的养育法。即它们先要找一个餐室，当作幼虫卵化后的家，
或者先找一个地方，使幼虫自己能觅到适当的食物，在出生后食用。在这种情况下，它
们是不需要父亲的，所以父亲通常到死都没有给它的后代成长工作以丝毫的帮助。
然而事情也不是常常以这种原始的方式进行的。有些种族为它们的家庭预备下妆奁，
作它们将来的食宿。蜜蜂和黄蜂特别善于营造小巢，例如口袋、小瓶等，并在里面装满
蜜，它们还十分善于造筑土穴，储藏着野味，给蛴螬做食物。
这种伟大的建筑巢穴和收集食物的工作，要花去它全部的生命，而这工作却是母亲
单独一人做的。这工作消磨它的时间，耗去它的生命。父亲则沉醉于日光下，懒惰地站
在工作场之外，只是看着它的勤劳的伴侣在从事艰苦的工作。
为什么它不帮助一下呢？事实上它从没有帮助过。为什么它不学学燕子夫妻，它们
都带一些草和一些泥土到巢里，还带一些小虫给小鸟。而雄性昆虫一点也没做那种事。
也许它借口比较软弱，无以作辩解。这是个无聊的议论。因为在叶子上割下一块，从植
物上摘下一些棉花，从泥土中收集一点水泥，完全是它的力量所能做到的。它很可以像
工人一样地帮助雌虫，它很适于收集一些材料，再由更智慧的雌虫建筑起来。它不做的
真正原因，只是因为它不愿做而已。
多数从事劳作的昆虫，竟然都不知道做父亲的责任，这是很令男人们感到奇怪的。
谁都为了幼虫发展最高才能的需要而努力，但这些父亲们仍然愚钝如蝴蝶，对于家族是
很少费力的。我们每一次都不能回答下面的问题：为什么这种昆虫，有这个特别的本能，
而别的昆虫就没有呢？
当我们看到清道夫甲虫有这种高贵的品质而收蜜的昆虫却没有时，我们非常的惊奇
和难以理解，好多种清道夫甲虫善于负起家政的重任，并知道两人共同工作的价值。例
如蜣螂夫妻，它们共同预备蛴螬的食物，父亲帮助它的伴侣在制造腊肠般的食物时，助
以强有力的轧榨工作。
它们就是形成家族共同劳作的习惯的最好的榜样，在普遍的自私的情形中，是最稀
罕的一个例外。
关于这件事，经过我长期的研究，在这个例子之外，我又可以增加三个另外的例子，
全都是清道夫甲虫合作的事实。
这三个中的一个是锡赛弗斯，它是搓丸药者中最小最勤劳的一个。它在它们当中最
活泼、最灵敏，并且毫不介意在危险的道路上倾倒和翻跟斗，在那里它固执地爬起来，
但又重新倒下去。正是因为那些狂乱的体操，所以人们给它起了一个名字，叫西西弗。
我想你们总该知道，一个可怜的人变得很著名，一定要经过很多很艰难的奋斗。它
被迫把一块大石头滚上高山，每次好容易到山顶时，那石头又滑落，滚到山脚下。我很
喜欢这个神话，这是我们当中许多人的历史，就拿我自己来说，刻苦的爬峻峭的山坡已
经五十多年了，把我的精力浪费在为了安全地得到每天的面包的挣扎里。面包一经滑去，
就滚下去，落到深渊里，很难摆稳。
我现在所谈及的西西弗，就不知道有这种困难，它在陡峭的山坡上毫无挂念地滚着
粮食，有时供给它自己，有时供给它的子女。在我们这些地方，它是很少见的，如果没
有我们从前几次提起过的助手，我也没有地方得到这么多观察物来研究。
我的小儿子保罗，年龄才七岁。他是我猎取昆虫的热心的同伴，而且比任何同龄的
小孩，更清楚地知道蝉、蝗虫、蟋蟀的秘密，尤其是清道夫甲虫。他锐利的眼光能在二
十步以外，辨别出地上隆起的土堆，哪一个是甲虫的巢穴，哪一个不是。他的灵敏的耳
朵可以听到螽斯细微的歌声，这是我完全办不到的。他帮助我看和听，我则把意见给他
用以作交换，他是很注意接受我的意见的。
小保罗有他自己养虫子的笼子，神圣的甲虫在里面做巢。他自己的花园，和手帕差
不多大小，能在里面种些豆子，但他常常将它们挖起来，看看小根长了一点没有。他的
林地上，有四株小槲树，只有手掌那么高，一旁还连着槲树子，在供给它养料。这是研
究昆虫之余极好的休息，对于昆虫研究的进步是毫无妨碍的。
五月将近的时候，有一天保罗和我起得很早，因为太早了，出去时连早饭都没有吃，
我们在山脚下的草场上，在羊群曾经走过的地方寻找。在这里，我们找着了西西弗，保
罗非常热心地搜索，不久我们就得到了足足好几对，收获真是不少。
使它们安居下来所需要的是一个铁丝的罩子，沙土的床，以及食物的供给——为了
这个我们也变成清道夫了。这些动物是很小的，还不及樱桃核大。形状也很奇怪！一个
短而肥的身体，后部是尖的，足很长，伸开来和蜘蛛的脚很像。后足更长，呈弯曲形，
挖土和搓小球时最有用。
不久，建设家族的时候到了。父亲和母亲同样热心地从事着搓卷、搬动和贮藏食物
的工作，都是为了它们的子女。它们利用前足的刀子，随意的从食物上割下小块来。夫
妻俩一同工作，一次次地抚拍和挤压，做成了一粒豌豆大的球。
和在神圣甲虫的工作场里一样，它们把圆球做成正确的圆形，是用不着机械的力量
来滚这球的。材料在没有移动之前，甚至在没有拾起之前，就已经做成圆形了。现在我
们又有了一个圆形学家，善于制作制造和保藏食物的最好的样式——圆形。
球不久就制造成功了。现在须用力地滚动，使它具有一层硬壳，保护里面柔软的物
质，使它不致变得太干燥。我们可以从大一些的身段上辨别出在前面全幅武装的母亲。
它将长长的后足放在地上，前足放在球上，将球向自己的身边拉，向后退着走。父亲处
在相反的方位，头向着下面，在后面推，这与神圣甲虫两个在一起工作时的方法相同，
不过目的是两样的，西西弗夫妻是为蛴螬搬运食物，而大的滚梨者(即神圣甲虫)则是准
备食物为自己在地下大嚼用的。
这一对在地面上走过，它们没有固定的目标，只是一直走下去，不管横在路中央的
障碍物。这样倒退着走，障碍固然是免不了的，但是即使看到了，它们也不会绕过障碍
走。它甚至做顽固的尝试，想爬过我的铁丝笼子。这是一种费时而且不可能完成的工作，
母亲的后足抓住铁丝网将球向它拉过来，然后用前足抱住它，把球抱在空中。父亲觉得
无物可推就抱住了球，伏在上面，把它身体的重量，加在球上，不再费什么力气了。这
种努力维持下去，未免太难了。于是球和骑在上面的昆虫，滚成一团，掉落到地上。母
亲从上面惊异地看着下面，不久就下来了，扶好这个球，重新做这个不可能成功的尝试。
一再的跌落之后，才放弃攀爬这个铁丝网。
就是在平地运输的时候，也不是完全没有困难的。差不多每分钟都会碰到隆起的石
头堆，货物就会翻倒。正在奋力推的昆虫也翻倒了，仰卧着把脚乱踢。不过这只是小事
情，很小很小的事情。西西弗是常常翻倒的，它并不在意。甚至有人也许以为它是喜欢
这样的。然而无论如何，球是变硬了，而且相当的坚固。跌倒、颠簸等都是程序单上的
一部分。这种疯狂的跳跃往往要继续几个小时！
最后母亲以为工作已经完毕，跑到附近找个适当的地点贮存球。父亲留守，蹲在实
物的上面。如果它的伴侣离开太久，它就用它高举的后足灵活地搓球，用以解闷儿，它
处置它珍贵的小球时，如同演戏者处置他的球一样。它用它变形的腿实验那个球是否完
整。那种高举的样子，无论谁看了，都不会怀疑它生活得很满足——父亲会保障它子女
将来的幸福的满足。
它好像是在说：“我搓成的这块圆球，是我给我的儿子们做面包的！”
它高高举起那个球，给每个人看看，这个是它工作的成果。这时候，母亲已经找到
了埋藏的地方，开始的一小部分工作已经做了，已经开始做下一个浅穴了。可以将球推
进浅穴。守卫的父亲一刻也不离开，母亲则用足和头挖土。不久，地穴已足够容纳那个
球了。它始终坚决地把球靠近自己，它一定感到在洞穴做成以前，一定要前后左右地把
它摇动摇动，以免寄生物的侵害。如果把它放在洞穴边上，一直到这个家完成，它害怕
会有什么不幸的事发生。因为有很多蚊蝇和别的动物，会出其不意地来攫取，因此不能
不格外当心。
于是圆球已经有一半放在还没有完成的土穴里了。母亲在下面，用足把球抱住往下
拉，父亲在上面，轻轻地往下放，而且还要注意落下去的泥土会不会把穴堵住，一切都
进行得很顺利。掘凿进行着，继续往下放球，非常的小心，一个虫往下拉，一个管理着
落下去的速度，并清除那些阻碍工作的事物。经过再进一步的努力，球和矿工都到地下
去了。以后所要做的事，是把从前做好的事再做一回，而且我们必须再等半天或几个小
时。
如果我们仔细等待，我们可以见到父亲又单独到地面上来了，它蹲在靠近土穴的沙
土。母亲为了尽它的伴侣不能帮助它的责任，常常要到第二天才出现。最后它也出来了，
父亲才离开它打瞌睡的地点，同它一道走。这对重新联合在一起的夫妇，又回到它们从
前找到食物的地方，休息一会，又收集起材料来。于是它们俩又重新工作，又一起塑模
型，运输和储藏球。
我对于这种恒心很是佩服。然而我不敢公然宣布，这是甲虫确定的习惯。无疑，有
许多甲虫是轻浮的，没有恒心的。但不要紧，我所看见的这点，关于西西弗爱护家庭的
习惯，已经使我看重了它们。
该是我们查看土巢的时候了。它并不很深，我看到墙边有一个小空隙，宽度足以让
母亲在球旁转动。寝室很小，这告诉我们父亲是不能在那里留很久的。当工作室准备好
了的时候，它一定要跑出去，请女雕刻家来继续工作。
地窖中只储藏着一只球，一件艺术的杰作。和神圣甲虫的梨形状相同，不过小得多。
因为小，球表面的光滑和圆形的标准，更加令人吃惊，最宽的地方，直径也只有一寸的
二分之一到四分之三。
另外还有一次对西西弗的观察。在我的铁丝笼下有六对，它们做了五十七个梨，每
个当中都有一颗卵——每一对平均有九个以上的蛴螬。神圣甲虫远不及这个数目。什么
原因使它产下这么多的后代呢？我看只有一个理由，就是父亲和母亲共同工作，一个家
族的负担，一人的精力不足以应付，两人分担起来就不觉得太重了。
------------------








捕蝇蜂



你已经知道了赤条蜂和黄蜂怎样麻痹毛毛虫或蟋蟀来喂自己的孩子，然后怎样封闭
洞口，离开巢飞到别处去。不过并不是每一种蜂都是这样生活的，现在你将要听到另一
种蜂，它们每天用新鲜的食物喂它的孩子，这就是捕蝇蜂。
这种蜜蜂喜欢在明亮的阳光下和蔚蓝的天空中选择最轻最松的泥土做它的巢。我有
时候会在一片没有树荫的广场上观察它们。天气很热，要避免烈日的煎熬只有躺在小沙
堆后面，把头钻进兔子洞，或是为自己预备一把大伞。我就采取了后一种办法，如果大
家愿意在七月快要结束的时候来和我一同坐在这样的大伞下，那么他(或她)也可以和我
一起饱饱眼福：
一只捕蝇蜂突然飞来，毫不犹豫地停在某个地方，这地方在我看来和别处没什么不
同。它前足上长着一排排的硬毛，会使你想起一把扫帚，一个刷子或一个钉锚。它用前
脚工作，用四只后脚支持着自己的身体。它先把沙耙起，然后向后拂去，它的动作非常
快，使这些连续不断的沙子看上去像不住地流水一样流到七八寸以外的地方。这种沙粒
的飞射要维持五分钟或十分钟左右。
和这些沙粒堆在一起的，还有木屑、腐烂的叶片和其它废料的碎屑。捕蝇蜂把这些
垃圾一一用嘴搬掉，这就是它工作的目的。它要使它家门前的沙都是又轻又细的“高级
沙粒”，没有任何粗重的杂质。这样，当它为孩子们捕了蝇回来的时候，它就可以很容
易地打开一条通路，把猎物带到洞里去。这种清洁工作，它总是在空闲的时候做。譬如，
在猎物已经储藏了许多，足够它的孩子们吃一段时间的时候，它没有必要再出去觅食，
那么它就会抽出时间来清除垃圾，像一个出色的家庭主妇那样。我们可以看出，在它认
真勤勉地工作的时候，它显得非常快乐和满足。也许这正是一个做母亲的看到孩子们在
自己盖的屋子里成长起来所涌上心头的喜悦吧！
如果我们用一把小刀，在母蜂所刮的沙地上挖下去，我们首先会发现一条隧道，有
一个手指那么粗，或许有八寸到十二寸那么长，接着就是一个小屋。小屋的大小足以容
纳三个胡桃，可是到目前为止，这里面只有一个蝇和一只白色的小卵，那就是捕蝇蜂的
卵。大约二十四个小时之后，这个卵就能孵化成一条小虫，小虫出来后靠吃母亲为它准
备好的死蝇长大。
大约两三天之后，捕蝇蜂的幼虫就快把那死蝇吃完了。这时的母蜂离家并不很远，
你可以看到它有时从花蕊里吸几口蜜汁充饥，有时快乐地坐在火热的沙地上——它是在
看守自己的家。它会常常在家门口耙去一些沙，然后，又飞走了，过了一段时间再来。
可是不管它在外面呆多久，它总不会忘记估算一下它家里的小屋中的食物还能维持多久，
作为一个母亲，它的本能会告诉它什么时候它孩子的食物快吃完了，于是它就回到自己
的巢里。至于这巢，上面已经提到过，在外面看来和其它沙地是一样的，没有明显的洞
口或什么标记，可是它自己却清清楚楚地知道它的巢在哪一点。它每次回来探望孩子，
总不忘带些丰盛的礼物回来，这次它带回一只大蝇，它把蝇送进地下的家后，自己又出
来，直到需要它再送第三只蝇的时候再下去，这中间隔的时间是很短的，因为幼虫的胃
口始终很好，如果母蜂稍有懈怠的话，它的孩子就要挨饿了。
这样维持了两个星期，幼虫不停地在成长。食物的需求量也越来越大，母蜂也不断
地送食物进来。在第二个星期末，幼虫已经长得很肥胖了。母亲加倍努力地寻找食物以
供养这老也吃不饱的孩子，直到它完全长大，不再需要别人给它准备食物为止。有一次
我算了一下，一条幼虫长大过程中所需要吃的蝇加起来达八十二只之多。
我有时候怀疑，这种蜂为什么不像其它蜂那样预先储藏好食物，把洞封好，自己也
就可以离开，何必老耐着性子守在洞口呢？可能是因为它捕回的死蝇不能藏得太久的缘
故吧。可是它又为什么不像黄翅蜂一样把蝇麻痹，而是把它杀死呢？我推测可能是因为
蝇毕竟和毛毛虫、蟋蟀不大一样，它是那样的轻，那样的软，放不了多久就会缩得没有
了。所以这东西必须吃新鲜的，否则就没有什么价值了。另一个原因可能是蝇非常灵敏，
必须擒得快，不像那呆头呆脑的毛毛虫和庞然大物似的蟋蟀，目标明显，动作又不灵便，
让母蜂有充分的时间去麻痹它们。捕蝇蜂在必要的时候，须随时应用它的爪子、嘴巴或
刺，这样捕捉来的蝇当然不能随心所欲地让它半死不活了，要么让蝇逃掉了，要么捉个
死蝇。母蜂当然选择后者。
要观察到捕蝇蜂袭击苍蝇可不是件容易的事。因为它总是在离巢很远的地方捕捉的。
可是“踏破铁蹄无觅处，’得来全不费功夫”，有一次，我就无意之中看到了这精彩的
一幕，饱了眼福。那天我张着伞坐在烈日下。享受着伞的阴影的不只我一个，还有各种
马蝇也躲在我的伞下休息。它们平静地歇在张着的伞顶上。我在伞下没有事情做，就欣
赏着它们大大的金色眼睛来消磨时间。那些眼睛在我的伞下闪闪发光，好像一颗颗宝石。
有时候伞的某一部分被晒得太热了，它们就不得不转移阵地，移到没有被太阳光晒到的
那部分。我很喜欢看它们这种严肃的动作。
有一天，我正在伞下打瞌睡，突然，“梆！”的一声，张着的伞忽然像皮鼓似地被
击了一下。
“怎么回事？”我一下子清醒了许多。
大概是一颗榆树的果子掉到伞上了吧！我想。
可是“梆——梆——梆！”一声接着一声地传来。是哪个爱搞恶作剧的家伙把种子
或石子往我伞上扔？我离开了我的伞荫，四处巡视了一下，什么也没有。那声音又响起
来了，我抬头往伞顶一望，原来如此！我终于明白是怎么回事了。附近的捕蝇蜂发现我
这里有这么多肥美的食物，都飞过来捕取猎物，一切都像我希望的那样进行着，我只需
静静地坐着观看便是了。
每隔十五分钟左右，就有一只捕蝇蜂飞进来，直向伞顶冲去，发出一声重击。于是
战争就在伞顶上展开了。那是多么精彩和紧张啊，大家打得难分难解，不分上下，使你
辨不清谁是袭击者谁是自卫者。不过这种争执并没有维持很久。不一会儿，捕蝇蜂就用
双腿夹着它的俘虏飞走了。奇怪的是，这愚蠢的蝇群还不肯离开这危险的场所——的确，
外面实在太热了，与其被晒死，还不如在里面“今朝有酒今朝醉”，先尽情享受再说。
现在，让我来观察这只带着战利品回去的蜂吧。当它接近自己家的时候，突然发出
一种尖锐的嗡嗡声，听来颇有点凄凉的意味，好像十分不安。这声音一直继续着，直到
它降落到地面上为止。它先在地面上方盘旋了一会，然后又小心翼翼地降落，如果它那
敏锐的眼睛发现了一些什么不正常的情形，它就要降低下落的速度，在上面盘旋几秒钟，
飞上去又飞下来，然后像一支箭一般地飞开去了。稍过一会儿，我们就知道它为什么犹
豫不决了。不到一会儿，它又回来了。这次它先在高处巡视一遍，然后慢慢降落到地上
某一点——这一点在我看来实在没什么特别之处。
我想它大概是随便降落在这一点上的，降落之后，它还得慢慢地寻找自己巢的入口。
可是后来我发现自己又低估了捕蝇蜂。它恰好不偏不倚地降落在自己的巢上。它把前面
的沙扒开一些，再用头一顶，便顺利地拖着它的猎物进巢了。它进去后，旁边的沙粒立
刻又堆上洞口把它堵住。这和我从前所看到的无数次捕蝇蜂回巢的情形一样。我常常惊
异于蜂类为什么毫不犹豫地找到它的巢的入口，虽然那入口处和旁边的地方完全一样，
没有任何可以辨别的记号。
捕蝇蜂回巢的时候，并不是每次要在空中盘旋许多时候，它之所以要这样不停盘旋
是因为它看到自己的巢被一种巨大的危险所笼罩。它那种凄凉的嗡嗡声是表示它内心的
忧愁和恐惧。在没有危险的时候，它绝对不会发出这种声音。那么它的敌人是谁呢？原
来是一只小小的蝇，外表看上去十分的软弱无能。而这捕蝇蜂，它虽是蝇类的天敌，大
马蝇的刽子手，但当它发现自己被这种小蝇所监视的时候，它竟然会吓得不敢进洞去。
而事实上那只小蝇小得像一个不够它的幼虫吃一顿的侏儒。
这情形似乎像猫怕老鼠一样让人费解。为什么捕蝇蜂不冲下去把这个无耻的小蝇赶
走呢？这个我无法解释。也许这其貌不扬的小蝇自有它的厉害之处，在茫茫宇宙中占有
相当的地位，就像许多凶猛的动物一样。大自然的规则常常是我们人类不能了解的。
我以后还会讲到这种蝇把卵产在捕蝇蜂放在巢内的猎物上。它的幼虫孵出来以后就
掠夺了捕蝇蜂幼虫的养料。如果食物不够的话，它们还会毫不手软地把捕蝇蜂的幼虫当
作美食吃掉。所以，它决不是一种微不足道的小蝇，而是一个无情的杀手。捕蝇蜂那么
害怕它，不是没有道理的。那么这种小蝇是怎么把卵产在捕蝇蜂的卵上的呢？这非常值
得研究。
它从不走近捕蝇蜂的巢，只是耐着性子等着捕蝇蜂拖着丰美的猎物回来。当捕蝇蜂
半个身体钻进洞穴的时候，它就冲下去附在那只死的马蝇身上，当捕蝇蜂艰难地拖着马
蝇进洞的时候，这种小蝇就以惊人的速度在马蝇身上产下一个卵，有时候会连续地产下
两三个。捕蝇蜂从前半身钻进洞到完全把猎物拖进洞，前后只不过是一瞬间的功夫，可
就是在这一瞬间的时间内，小蝇已经完成了它的任务。现在它可以在洞旁的阳光里蹲下，
为它的第二次偷袭作准备了。
平时总有三四个这样的小蝇同时出现在一个巢的附近。对于进巢的入口它们往往知
道得比较清楚。它们那暗红的肤色、大而红的眼睛、以及它们惊人的耐心，常常使我想
到绑票的情形。那些歹徒穿着黑色的衣服，头上包着红布，静静地躲在阴暗的角落里，
等候机会来拦住过路的客人。
那可怜的捕蝇蜂是因为看到了这种歹徒在家门口才犯踌躇的。它知道那帮歹徒一定
会干坏事的。但是最后它还是飞进了自己的家。于是这些小歹徒似的蝇们便飞起来紧紧
跟着它，它向前，它们也向前；它后退，它们也后退；它无法使它们离开。最后它终于
撑不住了，不得不歇歇脚，那些小歹徒也跟着歇下来，但仍然虎视耽耽地跟在背后。于
是捕蝇蜂又飞起来，带着一声愤怒的呜咽声。这些无耻的小歹徒仍然厚着脸皮紧迫不舍。
捕蝇蜂只好另外想了一个办法，另找了一条路，以比较高的速度飞行，希望它的敌人跟
不上它而最终迷失方向。没想到那伙小歹徒早已料到这一招，折回到洞口等它回来。果
然，不一会儿，以为已经摆脱了危险的捕蝇蜂们回来了，这帮小歹徒赶紧起身直追。母
蜂的耐心已经没有了，最后终于被它们找着了产卵的机会。
好在我们刚才所讲的那只捕蝇蜂没有遭到这种不幸，所以让我们来结束这一章吧。
蜂的幼虫吃着母亲留给它们的粮食，慢慢地长大。过了两个星期，它就开始做茧了。可
是在它身体内没有足够的丝，所以它必须掺入沙粒以增加它的硬度。它把残余的食物堆
积到小屋的一角，先把地面扫清，然后在墙和墙之间搭起白色的美丽的丝来。它先把丝
攀成一个网，然后开始第二步工作。
它在网的中央做一个吊床。这吊床好像一个袋子，一端封闭，一端留有小孔。捕蝇
蜂的幼虫半个身体伸在床外，用嘴巴一粒一粒地挑选沙粒，太大的沙粒它看不上眼，会
一下子把它丢开。选好后，它再把沙一粒一粒地衔进去，很均匀地铺在吊床袋的四周，
就像泥水匠把石子嵌入灰泥一样。
到现在为止，茧子的一端还是开着口的，它必须把它封上。它用丝织成一顶帽子，
大小恰巧能盖住茧子的开口处，在这上面，也嵌进一粒一粒的沙。现在茧子可以说做完
了。不过捕蝇蜂在茧里还要做一番修葺工作。它在墙上涂一层浆液，为的是避免让自己
柔嫩的皮肤被沙粒擦伤或者蹭破。在这之后它就可以安心地睡大觉了。不久它将变成一
只成年的捕蝇蜂，就像它的母亲一样。
------------------








寄生虫



在八九月里，我们应该到光秃秃的、被太阳灼得发烫的山峡边去看看，让我们找一
个正对太阳的斜坡，那儿往往热得烫手，因为太阳已经把它快烤焦了。恰恰是这种温度
像火炉一般的地方，正是我们观察的目标。因为就是在这种地方，我们可以得到很大的
收获。这一带热土，往往是黄蜂和蜜蜂的乐土。它们往往在地下的土堆里忙着料理食物
——这里堆上一堆象鼻虫、蝗虫或蜘蛛，那里一组组分列着蝇类和毛毛虫类，还有的正
在把蜜贮藏在皮袋里、土罐里、棉袋里或是树叶编的瓮里。
在这些默默地埋头苦干的蜜蜂和黄蜂中间，还夹杂着一些别的虫，那些我们称之为
寄生虫。它们匆匆忙忙地从这个家赶到那个家，耐心地躲在门口守候着，你别以为它们
是在拜访好友，它们这些鬼鬼祟祟的行为决不是出于好意，它们是要找一个机会去牺牲
别人，以便安置自己的家。
这有点类似于我们人类世界的争斗。劳苦的人们，刚刚辛辛苦苦地为儿女积蓄了一
笔财产，却碰到一些不劳而获的家伙来争夺这笔财产。有时还会发生谋杀、抢劫、绑票
之类的恶性事件，充满了罪恶和贪婪。至于劳动者的家庭，劳动者们曾为它付出了多少
心血，贮藏了多少他们自己舍不得吃的食物，最终也被那伙强盗活活吞灭了。世界上几
乎每天都有这样的事情发生，可以说，哪里有人类，哪里就有罪恶。昆虫世界也是这样，
只要存在着懒惰和无能的虫类，就会有把别人的财产占为已有的罪恶。蜜蜂的幼虫们都
被母亲安置在四周紧闭的小屋里，或呆在丝织的茧子里，为的是可以静静地睡一个长觉，
直到它们变为成虫。可是这些宏伟的蓝图往往不能实现，敌人自有办法攻进这四面不通
的堡垒。每个敌人都有它特殊的战略——那些绝妙又狠毒的技巧，你根本连想都想不到。
你看，一只奇异的虫，靠着一根针，把它自己的卵放到一条蛰伏着的幼虫旁边——这幼
虫本是这里真正的主人；或是一条极小的虫，边爬边滑地溜进了人家的巢，于是，蛰伏
着的主人格永远长睡不醒了，因为这条小虫立刻要把它吃掉了。那些手段毒辣的强盗，
毫无愧意地把人家的巢和茧子作为自己的巢和茧子，到了来年，善良的女主人已经被谋
杀，抢了巢杀了主人的恶棍倒出世了。
看看这一个，身上长着红白黑相间的条纹，形状像一只难看而多毛的蚂蚁，它一步
一步地仔细地考察着一个斜坡，巡查着每一个角落，还用它的触须在地面上试探着。你
如果看到它，一定会以为它是一只粗大强壮的蚂蚁，只不过它的服装要比普通的蚂蚁漂
亮。这是一种没有翅膀的黄蜂，它是其它许多蜂类的幼虫的天敌。它虽然没有翅膀，可
是它有一把短剑，或者说是一根利刺。只见它踯躅了一会儿，在某个地方停下来，开始
挖和扒，最后居然挖出了一个地下巢穴，就跟经验丰富的盗墓贼似的。这巢在地面上并
没有痕迹，但这家伙能看到我们人类所看不到的东西。它钻到洞里停留了一会儿，最后
又重新在洞口出现。这一去一来之间，它已经干下了无耻的勾当：它潜进了别人的茧子，
把卵产在那睡得正酣的幼虫的旁边，等它的卵孵化成幼虫，就会把茧子的主人当作丰美
的食物。
这里是另外一种虫，满身闪耀着金色的、绿色的、蓝色的和紫色的光芒。它们是昆
虫世界里的蜂雀，被称作金蜂，你看到它的模样，决不会相信它是盗贼或是搞谋杀的凶
手。可它们的确是用别的蜂的幼虫作食物的昆虫，是个罪大恶极的坏蛋。
这十恶不赦的金蜂并不懂得挖人家墙角的方法，所以只得等到母蜂回家的时候溜进
去。你看，一只半绿半粉红的金蜂大摇大摆地走进一个捕蝇蜂的巢。那时，正值母亲带
着一些新鲜的食物来看孩子们。于是，这个“侏儒”就堂而皇之地进了“巨人”的家。
它一直大摇大摆在走到洞的底端，对捕蝇蜂锐利的刺和强有力的嘴巴似乎丝毫没有惧意。
至于那母蜂，不知道是不是不了解金蜂的丑恶行径和名声，还是给吓呆了，竟任它自由
进去。来年，如果我们挖开捕蝇蜂的巢看看，就可以看到几个赤褐色的针箍形的茧子，
开口处有一个扁平的盖。在这个丝织的摇篮里，躺着的是金蜂的幼虫。至于那个一手造
就这坚固摇篮的捕蝇蜂的幼虫呢？它已完全消失了。只剩下了一些破碎的皮屑了。它是
怎么消失的？当然是被金蜂的幼虫吃掉了！
看看这个外貌漂亮而内心奸恶的金蜂，它身上穿着金青色的外衣，腹部缠着“青铜”
和“黄金”织成的袍子，尾部系着一条蓝色的丝带。当一只泥匠蜂筑好了一座弯形的巢，
把入口封闭，等里面的幼虫渐渐成长，把食物吃完后，吐着丝装饰着它的屋子的时候，
金蜂就在巢外等候机会了。一条细细的裂缝，或是水泥中的一个小孔，都足以让金蜂把
它的卵塞进泥匠蜂的巢里去。总之，到了五月底，泥匠蜂的巢里又有了一个针箍形的茧
子，从这个茧子里出来的，又是一个口边沾满无辜者的鲜血的金蜂，而泥匠蜂的幼虫，
早被金蜂当作美食吃掉了。
正如我们所知道的那样，蝇类总是扮演强盗或小偷或歹徒的角色。虽然它们看上去
很弱小，有时候甚至你用手指轻轻一撞，就可以把它们全部压死。可它们的确祸害不小。
有一种小蝇，身上长满了柔软的绒毛，娇软无比，只要你轻轻一摸就会把它压得粉身碎
骨，它们脆弱得像一丝雪片，可是当它们飞起来时有着惊人的速度。乍一看，只是一个
迅速移动的小点儿。它在空中徘徊着，翅膀震动得飞快，使你看不出它在运动，倒觉得
是静止的。好像是被一根看不见的线吊在空中。如果你稍微动一下，它就突然不见了。
你会以为它飞到别处去了，怎么找都没有。它到哪儿去了呢？其实，它哪儿都没去，它
就在你身边。当你以为它真的不见了的时候，它早就回到原来的地方了。它飞行的速度
是如此之快，使你根本看不清它运动的轨迹，那么它又在空中干什么呢？它正在打坏主
意，在等待机会把自己的卵放在别人预备好的食物上。我现在还不能断定它的幼虫所需
要的是哪一种食物：蜜、猎物，还是其它昆虫的幼虫？
有一种灰白色的小蝇，我对它比较了解，它蜷伏在日光下的沙地上，等待着抢劫的
机会。当各种蜂类猎食回来，有的衔着一只马蝇，有的衔着一只蜜蜂，有的衔着一只甲
虫，还有的衔着一只蝗虫。大家都满载而归的时候，灰蝇就上来了，一会儿向前，一会
儿向后，一会儿又打着转，总是紧跟着蜂，不让它从自己的眼皮底下溜走。当母蜂把猎
物夹在腿间拖到洞里去的时候，它们也准备行动了。就在猎物将要全部进洞的那一刻，
它们飞快地飞上去停在猎物的末端，产下了卵。就在那一眨眼的工夫里，它们以迅雷不
及掩耳之势完成了任务。母蜂还没有把猎物拖进洞的时候，猎物已带着新来的不速之客
的种子了，这些“坏种子”变成虫子后，将要把这猎物当作成长所需的食物，而让洞的
主人的孩子们活活饿死。
不过，退一步想，对于这种专门掠夺人家的食物吃人家的孩子来养活自己的蝇类，
我们也不必对它们过于指责。一个懒汉吃别人的东西，那是可耻的，我们会称他为“寄
生虫”，因为它牺牲了同类来养活自己，可昆虫从来不做这样的事情。它从来不掠取其
同类的食物，昆虫中的寄生虫掠夺的都是其它种类昆虫的食物，所以跟我们所说的“懒
汉”还是有区别的。你还记得泥匠蜂吗？没有一只泥匠蜂会去沾染一下邻居所隐藏的蜜，
除非邻居已经死了，或者已经搬到别处去很久了。其它的蜜蜂和黄蜂也一样。所以，昆
虫中的“寄生虫”要比人类中的“寄生虫”要高尚得多。
我们所说的昆虫的寄生，其实是一种“行猎”行为。例如那没有翅膀，长得跟蚂蚁
似的那种蜂，它用别的蜂的幼虫喂自己的孩子，就像别的蜂用毛毛虫、甲虫喂自己的孩
子一样。一切东西都可以成为猎手或盗贼，就看你从怎样的角度去看待它。其实，我们
人类是最大的猎手和最大的盗贼。他们偷吃了小牛的牛奶，偷吃了蜜蜂的蜂蜜，就像灰
蝇掠夺蜂类幼虫的食物一样。人类这样做是为了抚育自己的孩子。自古以来人类不也总
是想方设法地把自己的孩子拉扯大，而且往往不择手段吗——这不是很像灰蝇吗？
------------------








新陈代谢的工作者



有许多昆虫，它们在这世界上做着极有价值的工作，尽管它们从来没有因此而得到
相应的报酬和相称的头衔。当你走近一只死鼹鼠，看见蚂蚁、甲虫和蝇类聚集在它身上
的时候，你可能会全身起鸡皮疙瘩，拔腿就跑。你一定会以为它们都是可怕而肮脏的昆
虫，令人恶心。事实并不是这样的，它们正在忙碌着为这个世界做清除工作。让我们来
观察一下其中的几只蝇吧，我们就可以知道它们的所作所为是多么的有益于人类，有益
于整个自然界了。
你一定看见过碧蝇吧？也就是我们通常所说的“绿头苍蝇”。它们有着漂亮的金绿
色的外套，发着金属般的光彩，它们还有一对红色的大眼睛。
当它们嗅出在很远的地方有死动物的时候，会立即赶过去在那里产卵。几天以后，
你会惊讶地发现那动物的尸体变成了液体，里面有几千条头尖尖的小虫子，你一定会觉
得这种方法实在有点令人反胃，可是除此之外，还有什么别的更好更容易的方法消灭腐
烂发臭的动物的尸体，让它们分解成元素被泥土吸收而再为别的生物提供养料呢？是谁
能够使死动物的尸体奇迹般地消失，变成一滩液体的呢？正是碧蝇的幼虫。
如果这尸体没有经过碧蝇幼虫的处理，它也会渐渐地风干，这样的话，要经过很长
一段时间才会消失。碧蝇和其它蝇类的幼虫一样，有一种惊人的本事，那就是能使固体
物质变成液体物质。有一次我做了一个试验，把一块煮得很老的蛋白扔给碧蝇作食物，
它马上就把这块蛋白变成一滩像清水一样的液体。而这种使它能够把固体变成液体的东
西，是它嘴里吐出来的一种酵母素，就好像我们胃里的胃液能把食物消化一样。碧蝇的
幼虫就靠着这种自己亲手制作的肉汤来维持自己的生命。
其实，能做这种工作的，除了碧蝇之外，还有灰肉蝇和另一种大的肉蝇。你常常可
以看到这种蝇在玻璃窗上嗡嗡飞着。千万不要让它停在你要吃的东西上面，要不然的话，
它会使你的食物也变得充满细菌了。不过你可不必像对待蚊子一样，毫不客气地去拍死
它们，只要把它们赶出去就行了。因为在房间外面，它们可是大自然的功臣。它们以最
快的速度，用曾经活过的动物的尸体产生新的生命，它们使尸体变成一种无机物质被土
壤吸收，使我们的土壤变得肥沃，从而形成新一轮的良性循环。
------------------








松毛虫



在我那个园子里，种着几棵松树。每年毛毛虫都会到这松树上来做巢，松叶都快被
它们吃光了。为了保护我们的松树，每年冬天我不得不用长叉把它们的巢毁掉，搞得我
疲惫不堪。
你这贪吃的小毛虫，不是我不客气，是你太放肆了。如果我不赶走你，你就要喧宾
夺主了。我将再也听不到满载着针叶的松树在风中低声谈话了。不过我突然对你产生了
兴趣，所以，我要和你订一个合同，我要你把你一生的传奇故事告诉我，一年、两年，
或者更多年，直到我知道你全部的故事为止。而我呢，在这期间不来打扰你，任凭你来
占据我的松树。
订合同的结果是，不久我们就在离门不远的地方，拥有了三十几只松毛虫的巢。天
天看着这一堆毛毛虫在眼前爬来爬去，使我不禁对松毛虫的故事更有了一种急切了解的
欲望。这种松毛虫也叫作“列队虫”，因为它们总是一只跟着一只，排着队出去。
下面我开始讲它的故事：
第一，先要讲到它的卵。在八月份的前半个月，如果我们去观察松树的枝端，一定
可以看到在暗绿的松叶中，到处点缀着一个个白色的小圆柱。每一个小圆柱，就是一个
母亲所生的一簇卵。这种小圆柱好像小小的手电筒，大的约有一寸长，五分之一或六分
之一寸宽，裹在一对对松针的根部。这小筒的外貌，有点像丝织品，白里略透一点红，
小筒的上面叠着一层层鳞片，就跟屋顶上的瓦片似的。
这鳞片软得像天鹅绒，很细致地一层一层盖在筒上，做成一个屋顶，保护着筒里的
卵。没有一滴露水能透过这层屋顶渗进去。这种柔软的绒毛是哪里来的呢？是松毛虫妈
妈一点一点地铺上去的。它为了孩子牺牲了自己身上的一部分毛。它用自己的毛给它的
卵做了一件温暖的外套。
如果你用钳子把鳞片似的绒毛刮掉，那么你就可以看到盖在下面的卵了，好像一颗
颗白色珐琅质的小珠。每一个圆柱里大约有三百颗卵，都属于同一个母亲。这可真是一
个大家庭啊！它们排列得很好看，好像一颗玉蜀黍的穗。无论是谁，年老的或年幼的，
有学问的还是没文化的，看到松蛾这美丽精巧的“穗”，都会禁不住喊道：“真好看啊！”
多么光荣而伟大母亲啊！
最让我们感兴趣的东西，不是那美丽的珐琅质的小珠本身，而是那种有规则的几何
图形的排列方法。一只小小的蛾知道这精妙的几何知识，这难道不是一件令人惊讶的事
吗？但是我们愈和大自然接触，便愈会相信大自然里的一切都是按照一定的规则安排的。
比如，为什么一种花瓣的曲线有一定的规则？为什么甲虫的翅鞘上有着那么精美的花纹？
从庞然大物到微乎其微的小生命，一切都安排得这样完美，这是不是偶然的呢？似乎不
大可能吧？是谁在主宰这个世界呢？我想冥冥之中一定有一位“美”的主宰者在有条不
紊地安排着这个缤纷的世界。我只能这样解释了。
松蛾的卵在九月里孵化。在那时候，如果你把那小筒的鳞片稍稍掀起一些，就可以
看到里面有许多黑色的小头。它们在咬着，推着它们的盖子，慢慢地爬到小筒上面，它
们的身体是淡黄色的，黑色的脑袋有身体的两倍那么大。它们爬出来后，第一件事情就
是吃支持着自己的巢的那些针叶，把针叶啃完后，它们就落到附近的针叶上。常常可能
会有三四个小虫恰巧落在一起，那么，它们会自然地排成一个小队。这便是未来大军的
松毛虫雏形。如果你去逗它们玩，它们会摇摆起头部和前半身，高兴地和你打招呼。
第二步工作就是在巢的附近做一个帐篷。这帐篷其实是一个用薄绸做成的小球，由
几片叶子支持着。在一天最热的时候，它们便躲在帐里休息，到下午凉快的时候才出来
觅食。
你看松毛虫从卵里孵化出来还不到一个小时，却已经会做许多工作了：吃针叶、排
队和搭帐篷，仿佛没出娘胎就已经学会了似的。
二十四小时后，帐篷已经像一个榛仁那么大。两星期后，就有一个苹果那么大了。
不过这毕竟是一个暂时的夏令营。冬天快到的时候，它们就要造一个更大更结实的帐篷。
它们边造边吃着帐篷范围以内的针叶。也就是说，它们的帐篷同时解决了它们的吃住问
题。这的确是一个一举两得的好办法。这样它们就可以不必特意到帐外去觅食。因为它
们还很小，如果贸然跑到帐外，是很容易碰到危险的。
当它们把支持帐篷的树叶都吃完了以后，帐篷就要塌了。于是，像那些择水草而居
的阿拉伯人一样，全家会搬到一个新的地方去安居乐业。在松树的高处，它们又筑起了
一个新的帐篷。它们就这样辗转迁徒着，有时候竟能到达松树的顶端。
也就是这时候，松毛虫改变了它们的服装。它们的背上面上了六个红色的小圆斑，
小圆斑周围环绕着红色和绯红色的刚毛。红斑的中间又分布着金色的小斑。而身体两边
和腹部的毛都是白色的。
到了十一月，它们开始在松树的高处，木枝的顶端筑起冬季帐篷来。它们用丝织的
网把附近的松叶都网起来。树叶和丝合成的建筑材料能增加建筑物的坚固性。全部完工
的时候，这帐篷的大小相当于半加仑的容积，它的形状像一个蛋。巢的中央是一根乳白
色的极粗的丝带，中间还夹杂着绿色的松叶。顶上有许多圆孔，是巢的门，毛毛虫们就
从这里爬进爬出。在矗立在帐外的松叶的顶端有一个用丝线结成的网，下面是一个阳台。
松毛虫常聚集在这儿晒太阳。它们晒太阳的时候，像叠罗汉似的堆成一堆，上面张着的
丝线用来减弱太阳光的强度，使它们不至于被太阳晒得过热。
松毛虫的巢里并不是一个整洁的地方，这里面满是杂物的碎屑，毛虫们蜕下来的皮、
以及其它各种垃圾，真的可以称作是“败絮其中”。
松毛虫整夜歇在巢里，早晨十点左右出来，到阳台上集合，大家堆在一起，在太阳
底下打吨。它们就这样消磨掉整个白天。它们会时不时地摇摆着头以表示它们的快乐和
舒适。到傍晚六七点钟光景，这班瞌睡虫都醒了，各自从门口回到自己家里。
它们一面走一面嘴上吐着丝。所以无论走到哪里，它们的巢总是愈变愈大，愈来愈
坚固。它们在吐着丝的时候还会把一些松叶掺杂进去加固。每天晚上总有两个小时左右
的时间做这项工作。它们早已忘记夏天了，只知道冬天快要来了，所以每一条松毛虫都
抱着愉快而紧张的心情工作着，它们似乎在说：
“松树在寒风里摇摆着它那带霜的枝桠的时候，我们将彼此拥抱着睡在这温暖的巢
里！多么幸福啊！让我们满怀希望，为将来的幸福努力工作吧！”
不错，亲爱的毛毛虫们，我们人类也和你们一样，为了求得未来的平静和舒适而孜
孜不倦地劳动。让我们怀着希望努力工作吧！你们为你们的冬眠而工作，它能使你们从
幼虫变为蛾；我们为我们最后的安息而工作，它能消灭生命，同时创造出新的生命。让
我们一起努力工作吧！
做完了一天的工作，就是它们的用餐时间了。它们都从巢里钻出来，爬到巢下面的
针叶上去用餐。它们都穿着红色的外衣，一堆堆地停在绿色的针叶上，树枝都被它们压
得微微向下弯了。多么美妙的一副图画啊！这些食客们都静静地安详地咬着松叶，它们
那宽大的黑色的额头在我的灯笼下发着光。它们都要吃到深夜才肯罢休。回到巢里后还
要继续工作一会儿。当最后一批松毛虫进巢的时候，大约已是深夜一二点钟了。
松毛虫所吃的松叶通常只有三种，如果拿其它的常绿树的叶子给它们吃，即使那些
叶子的香味足以引起食欲，可松毛虫是宁可饿死也不愿尝一下的。这似乎没什么好说的，
松毛虫的胃和人的胃有着相同的特点。
松毛虫们在松树上走来走去的时候，随路吐着丝，织着丝带，回去的时候就依照丝
带所指引的路线。有时候它们找不到自己的丝带而找了别的松毛虫的丝带，那样它就会
走入一个陌生的巢里。但是没有关系，巢里的主人和这不速之客之间丝毫不会引起争执。
大家似乎都习以为常，平静得跟什么事都没有发生一样。到了睡觉的时候，大家也就像
兄弟一般睡在一起了，谁都没有一点生疏的感觉。不论是主人还是客人，大家都依旧在
限定的时间里工作，使它们的巢更大、更厚。由于这类意外的事情常有发生，所以有几
个巢总能接纳“外来人员”为自己的巢添砖加瓦，它们的巢就显得比其它的巢大了不少。
“人人为我，我为人人”是它们的信条，每一条毛毛虫都尽力地吐着丝，使巢增大增厚，
不管那是自己的巢还是别人的巢。事实上，正是因为这样才扩大了总体上的劳动成果。
如果每个松毛虫都只筑自己的巢，宁死也不愿替别家卖命，结果会怎样？我敢说，一定
会一事无成，谁也造不了又大又厚的巢。因此它们是几百几百地一起工作的，每一条小
小的松毛虫，都尽了它自己应尽的一份力量，这样团结一致才造就了一个个属于大家的
堡垒，一个又大又厚又暖和的大棉袋。每条松毛虫为自己工作的过程也是为其它松毛虫
工作的过程，而其它松毛虫也相当于都在为它工作。多么幸福的松毛虫啊，它们不知道
什么私有财产和一切争斗的根源。

毛虫队
有一个老故事，说是有一只羊，被人从船上扔到了海里，于是其余的羊也跟着跳下
海去。“因为羊有一种天性，那就是它们永远要跟着头一只羊，不管走到哪里。就因为
这，亚里士多德曾批评羊是世界上最愚蠢、最可笑的动物。”那个讲故事的人这样说。
松毛虫也具有这种天性，而且比羊还要强烈。第一只到什么地方去，其余的都会依
次跟着去，排成一条整齐的队伍，中间不留一点空隙。它们总是排成单行，后一只的须
触到前一只的尾。为首的那只，无论它怎样打转和歪歪斜斜地走，后面的都会照它的样
子做，无一例外。第一只毛毛虫一面走一面吐出一根丝，第二只毛虫踏着第一只松毛虫
吐出的丝前进，同时自己也吐出一条丝加在第一条丝上，后面的毛毛虫都依次效仿，所
以当队伍走完后，就有一条很宽的丝带在太阳下放着耀眼的光彩。这是一种很奢侈的筑
路方法。我们人类筑路的时候，用碎石铺在路上，然后用极重的蒸汽滚筒将它们压平，
又粗又硬但非常简便。而松毛虫，却用柔软的缎子来筑路，又软又滑但花费也大。
这样的奢侈有什么意义吗？它们为什么不能像别的虫子那样免掉这种豪华的设备，
简朴地过一生呢？我替它们总结出两条理由：松毛虫出去觅食的时间是在晚上，而它们
必须经过曲曲折折的道路。它们要从一根树枝爬到另一根树枝上，要从针叶尖上爬到细
枝上，再从细枝爬到粗枝上。如果它们没有留下丝线作路标，那么它们很难找回自己的
家，这是最基本的一条理由。
有时候，在白天它们也要排着队作长距离的远征，可能经过三十码左右的长距离。
它们这次可不是去找食物，而是去旅行，去看看世界，或者去找一个地方，作为它们将
来蛰伏的场所。因为在变成蛾子之前，它们还要经过一个蛰伏期。在作这样长途旅行的
时候，丝线这样的路标是不可缺少的。
在树上找食物的时候，它们或许是分散在各处，或许是集体活动，反正只要有丝线
作路标，它们就可以整齐一致地回到巢里。要集合的时候，大家就依照着丝线的路径，
从四面八方匆匆聚集到大队伍中来。所以这丝带不仅仅是一条路，而且是使一个大团体
中各个分子行动一致的一条绳索。这便是第二个理由。
每一队总有一个领头的松毛虫，无论是长的队还是短的队。它为什么能做领袖则完
全出自偶然，没有谁指定，也没有公众选举，今天你做，明天它做，没有一定的规则。，
毛虫队里发生的每一次变故常常会导致次序的重新排列。比如说，如果队伍突然在行进
过程中散乱了，那么重新排好队后，可能是另一只松毛虫成了领袖。尽管每一位“领袖”
都是暂时的、随机的，但一旦作了领袖，它就摆出领袖的样子，承担起一个领袖应尽的
责任。当其余的松毛虫都紧紧地跟着队伍前进的时候，这位领袖趁队伍调整的间隙摇摆
着自己的上身，好像在做什么运动。又好像在调整自己——毕竟，从平民到领袖，可是
一个不小的飞跃，它得明确自己的责任，不能和刚才一样，只需跟在别人后面就行了，
当它自己前进的同时，它就不停地探头探脑地寻找路径。它真是在察看地势吗？它是不
是要选一个最好的地方？还是它突然找不到引路的丝线，所以犯了疑？看着它那又黑又
亮，活像一滴柏油似的小脑袋，我实在很难推测它真的在想什么？我只能根据它的一举
一动，作一些简单的联想。我想它的这些动作是帮助它辨出哪些地方粗糙，哪些地方光
滑，哪些地方有尘埃，哪些地方走不过去。当然，最主要的是辨出那条丝带朝着哪个方
向延伸。
松毛虫的队伍长短不一，相差悬殊，我所看到的最长的队伍有十二码或十三码，其
中包含二百多只松毛虫，排’成极为精致的波纹形的曲线，浩浩荡荡的，最短的队伍一
共只有两条松毛虫，它们仍然遵从原则，一个紧跟在另一只的后面。
有一次我决定要和我养在松树上的松毛虫开一次玩笑，我要用它们的丝替它们铺一
条路，让它们依照我所设想的路线走。既然它们只会不假思索地跟着别人走，那么如果
我把这路线设计成一个既没有始点也没有终点的圆，它们会不会在这条路上不停地打转
转呢？
一个偶然的发现帮助我实现了这个计划。在我的院子里有几个栽棕树的大花盆，盆
的圆周大约有一码半长。松毛虫们平时很喜欢爬到盆口的边沿，而那边沿恰好是一个现
成的圆周。
有一天，我看到很大一群毛虫爬到花盆上，渐渐地来到它们最为得意的盆沿上。慢
慢地，这一队毛虫陆陆续续到达了盆沿，在盆沿上前进着。我等待并期盼着队伍形成一
个封闭的环，也就是说，等第一只毛虫绕过一因而回到它出发的地方。一刻钟之后，这
个目的达到了。现在有整整一圈的松毛虫在绕着盆沿走了。’第二步工作是，必须把还
要上来的松毛虫赶开，否则它们会提醒原来盆沿上的那围虫走错了路线，从而扰乱实验。
要使它们不走上盆沿，必须把从地上到花盆间的丝拿走。于是我就把还要继续上去的毛
虫拨开，然后用刷子把丝线轻轻刷去，这相当于截断了它们的通道。这样下面的虫子再
也上不去，上面的再也找不到回去的路。这一切准备就绪后，我们就可以看到一幕有趣
的景象在眼前展开了：
一群毛虫在花盆沿上一圈一圈地转着，现在它们中间已经没有领袖了。因为这是一
个封闭的圆周，不分起点和终点，谁都可以算领袖，谁又都不是领袖，可它们自己并不
知道这一点。
丝织和轨道越来越粗了，因为每条松毛虫都不断地把自己的丝加上去。除了这条圆
周路之外，再也没有别的什么叉路了，看样子它们会这样无止境地一圈一圈绕着走，直
到累死为止？
旧派的学者都喜欢引用这样一个故事：“有一头驴子，它被安放在两捆干草中间，
结果它竟然饿死了。因为它决定不出应该先吃哪一捆。”其实现实中的驴子不比别的动
物愚蠢，它舍不得放弃任何一捆的时候，会把两捆一起吃掉。我的毛虫会不会表现得聪
明一点呢？它们会离开这封闭的路线吗？我想它们一定会的。我安慰自己说：
“这队伍可能会继续走一段时间，一个钟头或两个钟头吧。然后，到某个时刻，毛
毛虫自己就会发现这个错误，离开那个可怕的骗人的圈子，找到一条下来的路。”
而事实上，我那乐观的设想错了，我太高估了我的毛毛虫们了。如果说这些毛虫会
不顾饥饿，不顾自己一直回不到巢，只要没有东西阻挠它们，它们就会一直在那儿打圈
子，那么它们就蠢得令人难以置信了。然而，事实上，它们的确有这么蠢。
松毛虫们继续着它们的行进，接连走了好几个钟头。到了黄昏时分，队伍就走走停
停，它们走累了。当天气逐渐转冷时，它们也逐渐放慢了行进的速度。到了晚上十点钟
左右，它们继续在走，但脚步明显慢了下来，好像只是懒洋洋地摇摆着身体。进餐的时
候到了，别的毛虫都成群结队地走出来吃松叶。可是花盆上的虫子们还在坚持不懈地走。
它们一定以为马上可以到目的地和同伴们一起进晚餐了。走了十个钟头，它们一定又累
又饿，食欲极好。一棵松树离它们不过几寸远，它们只要从花盆上下来，就可以到达松
树，美美地吃上一顿松叶了。但这些可怜的家伙已经成了自己吐的丝的奴隶了，它们实
在离不开它，它们一定像看到了海市蜃楼一样，总以为马上可以到达目的地，而事实上
还远着呢！十点半的时候，我终于没有耐心了，离开它们去睡我的觉。我想在晚上的时
候它们可能清醒些。可是第二天早晨，等我再去看它们的时候，它们还是像昨天那样排
着队，但队伍是停着的。晚上太冷了，它们都蜷起身子取暖，停止了前进。等空气渐渐
暖和起来后，它们恢复了知觉，又开始在那儿兜圈子了。
第三天，一切还都像第二天一样。这天夜里非常冷，可怜的毛虫又受了一夜的苦。
我发现它们在花盆沿分成两堆，谁也不想再排队。它们彼此紧紧地挨在一起，为的是可
以暖和些。现在它们分成了两队，按理说每队该有一个自己的领袖了，可以不必跟着别
人走，各自开辟一条生路了。我真为它们感到高兴。看到它们那又黑又大的脑袋迷茫地
向左右试探的样子，我想不久以后它们就可以摆脱这个可怕的圈子了。可是不久我发现
自己又错了。当这两支分开的队伍相逢的时候，又合成一个封闭的圆圈，于是它们又开
始了整天兜圈子，丝毫没有意识到错过了一个绝佳的逃生机会。
后来的一个晚上还是很冷。这些松毛虫又都挤成了一堆，有许多毛虫被挤到丝织轨
道的两边，第二天一觉醒来，发现自己在轨道外面，就跟着轨道外的一个领袖走，这个
领袖正在往花盆里面爬。这队离开轨道的冒险家一共有七位，而其余的毛虫并没有注意
它们，仍然在兜圈子。
到达花盆里的毛虫发现那里并没有食物，于是只好垂头丧气地依照丝线指示的原路
回到了队伍里，冒险失败了。如果当初选择的冒险道路是朝着花盆外面而不是里面的活，
情形就截然不同了。
一天又过去了，这以后又过了一天。第六天是很暖和的。我发现有几个勇敢的领袖，
它们热得实在受不住了，于是用后脚站在花盆最外的边沿上，做着要向空中跳出去的姿
势。最后，其中的一只决定冒一次险，它从花盆沿上溜下来，可是还没到一半，它的勇
气便消失了，又回到花盆上，和同胞们共甘苦。这时盆沿上的毛虫队已不再是一个完整
的圆圈，而是在某处断开了。也正是因为有了一个唯一的领袖，才有了一条新的出路。
两天以后，也就是这个实验的第八天，由于新道路的开辟，它们已开始从盆沿上往下爬，
到日落的时候，最后一只松毛虫也回到了盆脚下的巢里。
我计算了一下，它们一共走了四十八个小时。绕着圆圈走过的路程在四分之一公里
以上。只有在晚上寒冷的时候，队伍才没有了秩序，使它们离开轨道，几乎安全到达家
里。可怜无知的松毛虫啊！有人总喜欢说动物是有理解力的，可是在它们身上，我实在
看不出这个优点。不过，它们最终还是回到了家，而没有活活饿死在花盆沿上，说明它
们还是有点头脑的。

松毛虫能预测气候
在正月里，松毛虫会脱第二次皮。它不再像以前那么美丽了，不过有失也有得，它
添了一种很有用的器官。现在它背部中央的毛变成暗淡的红色了。由于中央还夹杂着白
色的长毛，所以看上去颜色更淡了。这件褪了色的衣服有一个特点，那就是在背上有八
条裂缝，像口子一般，可以随毛虫的意图自由开闭。当这种裂缝开着的时候，我们可以
看到每只口子里有一个小小的“瘤”。这玩意儿非常的灵敏，稍稍有一些动静它就消失
了。这些特别的口子和“瘤”有什么用处呢？当然不是用来呼吸的，因为没有一种动物
——即便是一条松毛虫，也不会从背上呼吸的。让我们来想想松毛虫的习性，或许我们
可以发现这些器官的作用。
冬天和晚上的时候，是松毛虫们最活跃的时候，但是如果北风刮得太猛烈的话，天
气冷得太厉害，而且会下雨下雪或是雾厚得结成了冰屑，在这样的天气里，松毛虫总会
谨慎地呆在家里，躲在那雨水不能穿透的帐篷下面。
松毛虫们最怕坏天气，一滴雨就能使它们发抖，一片雪花就能惹起它们的怒火。如
果能预先料到这种坏天气。那么对松毛虫的日常生活是非常有意义的。在黑夜里，这样
一支庞大的队伍到相当远的地方去觅食，如果遇到坏天气，那实在是一件危险的事。如
果突然遭到风雨的袭击，那么松毛虫就要遭殃了，而这样的不幸在坏的季节里是常常会
发生的。可松毛虫们自有办法。让我来告诉你它们是怎样预测天气的吧。
有一天，我的几个朋友，和我一起到院子里看毛虫队的夜游。我们等到九点钟，就
进入到院子里。可是……可是……这是怎么了？巢外一只毛虫都没有！就在昨天晚上和
前天晚上还有许多毛虫出来呢，今天怎么会一只都没有了？它们都上哪儿去了？是集体
出游吗？还是遇到了灭顶之灾？我们等到十点、十一点，一直到半夜。失望之余，我只
得送我的朋友走了。
第二天，我发现那天晚上竟然下了雨，直到早晨还继续下着，而且山上还有积雪。
我脑子里突然闪过一个念头，是不是毛虫对天气的变化比我们谁都灵敏呢？它们昨晚没
有出来，是不是因为早已预料到天气要变坏，所以不愿意出来冒险？一定是这样的！我
为自己的想法暗暗喝彩，不过我想我还得仔细观察它们。
我发现每当报纸上预告气压来临的时候，比如说暴风雨将要来临的时候，我的松毛
虫总躲在巢里。虽然它们的巢暴露在坏天气中，可风啊、雨啊、雪啊、寒冷啊，都不能
影响它们。有时候它们能预报雨天以后的风暴。它们这种推测天气的天赋，不久就得到
我们全家的承认和信任。每当我们要进城去买东西的时候，前一天晚上总要先去征求一
下松毛虫们的意见，我们第二天去还是不去，完全取决于这个晚上松毛虫的举动，它成
了我们家的“小小气象预报员”。
所以，想到它的小孔，我推测松毛虫的第二套服装似乎给了它一个预测天气的本领。
这种本领很可能是与那些能自由开闭的口子息息相关。它们时时张开，取一些空气作为
样品，放到里面检验一番，如果从这空气里测出将有暴风雨来临，便立刻发出警告。

松蛾
三月到来的时候，松毛虫们纷纷离开巢所在的那棵松树，作最后一次旅行。三月二
十日那天，我花了整整一个早晨，观察了一队三码长，包括一百多只毛虫在内的毛虫队。
它们衣服的颜色已经很淡了。队伍很艰难地徐徐地前进着，爬过高低不平的地面后，就
分成了两队，成为两支互不相关的队伍，各分东西。
它们目前有极为重要的事情要做。队伍行进了两小时光景，到达一个墙角下，那里
的泥土又松又软，极容易钻洞。为首的那条松毛虫一面探测，一面稍稍地挖一下泥土，
似乎在测定泥土的性质。其余的松毛虫对领袖百分之一百的服从，因此只是盲目地跟从
着它，全盘接受领袖的一切决定，也不管自己喜欢不喜欢。最后，领头的松毛虫终于找
到了一处它自己挺喜欢的地方，于是停下脚步。接着其余的松毛虫都走出队伍，成为乱
哄哄的一群虫子，仿佛接到了“自由活动”的命令，再也不要规规矩矩地排队了。所有
的虫子的背部都杂乱地摇摆着，所有的脚都不停地靶着，所有的嘴巴都挖着泥土，渐渐
地它们终于挖出了安葬自己的洞。到某个时候，打过地道的泥土裂开了，就把它们埋在
里面。于是一切都又恢复平静了。现在，毛虫们是葬在离地面三寸的地方，准备着织它
们的茧子。
两星期后，我往地面下挖土，又找到了它们。它们被包在小小的白色丝袋里，丝袋
外面还沾染着泥土。有时候，由于泥土土质的关系，它们甚至能把自己埋到九寸以下的
深处。
可是那柔软的、翅膀脆弱而触须柔软的蛾子是怎么从下面上来到达地面的呢？它一
直要到七八月才出来。那时候，由于风吹雨打，日晒雨淋，泥土早已变得很硬了。没有
一只蛾子能够冲出那坚硬的泥土，除非它有特殊的工具，并且它的身体形状必须很简单。
我弄了一些茧子放到实验室的试管里，以便看得更仔细些。我发现松娥在钻出茧子的时
候，有一个蓄势待发的姿势，就像短跑运动员起跑前的下蹲姿势一样。它们把它美丽的
衣服卷成一捆，自己缩成一个圆底的圆柱形，它的翅膀紧贴在脚前，像一条围巾一般，
它的触须还没有张开，于是把它们弯向后方，紧贴在身体的两旁。它身上的毛发向后躺
平，只有腿是可以自由活动的，为的是可以帮助身体钻出泥土。
虽然有了这些准备，但对于挖洞来说，还远远不够，它们还有更厉害的法宝呢！如
果你用指尖在它头上摸一下，你就会发现有几道很深的皱纹。我把它放在放大镜下，发
现那是很硬的鳞片。在额头中部顶上的鳞片是所有鳞片中最硬的。这多象一个回旋钻的
钻头呀。在我的试管里，我看到蛾子用头轻轻地这边撞撞，那边碰碰，想把沙块钻穿。
到第二天，它们就能钻出一条十寸长的隧道通到地面上来了。
最后，蛾子终于到达了泥土外面，只见它缓缓地展开它的翅膀，伸展它的触须，蓬
松一下它的毛发。现在它已完全打扮好了，完全是一只漂亮成熟又自由自在的蛾子了。
尽管它不是所有蛾子中最美丽的一种，但它的确已经够漂亮了。你看，它的前翅是灰色
的，上面嵌着几条棕色的曲线，后翅是白色的，腹部盖着淡红色的绒毛。颈部围着小小
的鳞片，又因为这些鳞片挤得很紧密，所以看上去就像是一整片，非常像一套华丽的盔
甲。
关于这鳞片，还有些极为有趣的事情。如果我们用针尖去刺激这些鳞片，无论我们
的动作多么轻微，立刻会有无数的鳞片飞扬起来。这种鳞片就是松蛾用来做盛卵的小筒
用的，我们在这一章的开头已经讲过了。
------------------








卷心菜毛虫



卷心菜几乎可以说是我们所有的蔬菜中最为古老的一种，我们知道古时候的人就已
经开始吃它了。而实际上在人类开始吃它之前，它已经在地球上存在了很久很久，所以
我们实在是无法知道究竟是什么时候，它出现了，又是什么时候人类第一次种植它们，
用的又是什么方法。植物学家告诉我们，它最初是一种长茎、小叶、长在滨海悬崖的野
生植物。历史对于这类细小的事情的记载是不愿多花笔墨的。它所歌颂的，是那些夺去
千万人生命的战场，它觉得那一片使人类生生不息的土地是没有研究价值的。它详细列
举各国国王的嗜好和怪癖，却不能告诉我们小麦的起源！但愿将来的历史记载会改变它
的作风。
我们对于卷心菜知道得实在太少了，那实在有点可惜，它的确算得上一种很贵重的
东西。因为它拥有许多有趣的故事。不仅是人类，就是别的动物也都与它有千丝万缕的
联系。其中有一种普通大白蝴蝶的毛虫，就是靠卷心菜生长的。它们吃卷心菜皮及其一
切和卷心菜相似的植物叶子，像花椰菜、白菜芽、大头菜，以及瑞典萝卜等，似乎生来
就与这种样子的莱类有不解之缘。
它们还吃其它一些和卷心菜同类的植物。它们都属于十字花科——植物学家们这样
称呼它们，因为它们的花有四瓣，排成十字形。白蝴蝶的卵一般只产在这类植物上。可
是它们怎么知道这是十字花科植物呢？它们又没有学过植物学。这倒是个谜。我研究植
物和花草已有五十多年，但如果要我判定一种没有开花的植物是不是属于十字花科，我
只能去查书。现在我不需要去一一查书了，我会根据白蝴蝶留下的记号作出判断——我
是很信任它的。
白蝴蝶每年要成熟两次。一次是在四五月里，一次是在十月，这正是我们这里卷心
菜成熟的时候。白蝴蝶的日历恰巧和园丁的日历一样。当我们有卷心菜吃的时候，白蝴
蝶也快要出来了。
白蝴蝶的卵是淡橘黄色的，聚成一片，有时候产在叶子朝阳的一面，有时候产在叶
子背着阳光的一面。大约一星期后，卵就变成了毛虫，毛虫出来后第一件事就是把这卵
壳吃掉。我不止一次看到幼虫自己会把卵壳吃掉，不知道这是什么意思。我的推测是这
样的：卷心菜的叶片上有蜡，滑得很，为了要使自己走路的时候不至于滑倒，它必须弄
一些细丝来攀缠住自己的脚，而要做出丝来，需要一种特殊的食物。所以它要把卵壳吃
掉，因为那是一种和丝性质相似的物质，在这初生的小虫胃里，它比较容易转化成小虫
所需要的丝。
不久，小虫就要尝尝绿色植物了。卷心菜的灾难也就此开始了。它们的胃口多好啊！
我从一颗最大的卷心菜上采来一大把叶子去喂我养在实验室的一群幼虫，可是两个小时
后，除了叶子中央粗大的叶脉之外，已经什么都不剩了。照这样的速度吃起来，这一片
卷心菜田没多少日子就会被吃完了。
这些贪吃的小毛虫，除了偶尔有一些伸胳膊挪腿的休息动作外，什么都不做，就知
道吃。当几只毛虫并排地在一起吃叶子的时候，你有时候可以看见它们的头一起活泼地
抬起来，又一起活泼地低下去。就这样一次一次重复着做，动作非常整齐，好像普鲁士
士兵在操练一样。我不知道它们这种动作是什么含义，是表示它们在必要的时候有作战
能力呢，还是表示它们在阳光下吃食物很快乐？总之，在它们成为极肥的毛虫之前，这
是它们唯一的练习。
吃了整整一个月之后，它们终于吃够了。于是就开始往各个方向爬。一面爬，一面
把前身仰起，作出在空中探索的样子，似乎是在做伸展运动，为了帮助消化和吸收吧。
现在气候已经开始转冷了，所以我把我的毛虫客人们都安置在花房里，让花房的门开着。
可是，令我惊讶的是，有一天我发现，这群毛虫都不见了。
后来我在附近各处的墙脚下发现了它们。那里离花房差不多有三十码的距离。它们
都栖在屋檐下，那里可以作为它们冬季的居所了。卷心菜的毛虫长得非常壮实健康，应
该不十分怕冷。
就在这居所里，它们织起茧子，变成蛹。来年春天，就有蛾从这里飞出来了。
听着这卷心菜毛虫的故事，我们也许会感到非常有趣。可是如果我们任凭它大量繁
殖，那么我们很快就会没有卷心菜吃了。所以当我们听说有一种昆虫，专门猎取卷心莱
毛虫，我们并不感到痛惜。因为这样可以使它们不至于繁殖得太快。如果卷心菜毛虫是
我们的敌人，那么那种卷心菜的敌人就是我们的朋友了。但它们长得那样细小，又都喜
欢埋头默默无闻地工作，使得园丁们非但不认识它，甚至连听都没听说过它，即使他偶
然看到它在它所保护的植物周围徘徊，他也决不会注意它，更不会想到它对自己会有那
么大的贡献。
我现在要给这小小的侏儒们一些应得的奖赏。
因为它长得细小，所以科学家们称它为“小侏儒”，那么让我也这么称呼它吧，我
实在不知道它还有什么别的好听一点的名字。
它是怎样工作的呢？让我们来看看吧。春季，如果我们走到菜园里去，一定可以看
见，在墙上或篱笆脚下的枯草上，有许多黄色的小茧子，聚集成一堆一堆的，每堆有一
个榛仁那么大。每一堆的旁边都有一条毛虫，有时候是死的，看上去大都很不完整，这
些小茧子就是“小侏儒”的工作成果，它们是吃了可怜的毛虫之后才长大的，那些毛虫
的残尸，也是“小侏儒”们剥下的。
这种“小侏儒”比幼虫还要小。当卷心菜毛虫在菜上产下橘黄色的卵后，“小侏儒”
的蛾就立刻赶去，靠着自己坚硬的钢毛的帮助，把自己的卵产在卷心菜毛虫的卵膜表面
上。一只毛虫的卵里，往往可以有好几个“小侏儒”，跑去产卵。照它们卵的大小来看，
一只毛虫差不多相当于六十五只“小侏儒”。
当这毛虫长大后，它似乎并不感到痛苦。它照常吃着菜叶，照常出去游历，寻找适
宜做茧子的场所。它甚至还能开展工作，但是它显得非常萎靡、非常无力，经常无精打
彩的，渐渐地消瘦下去，最后，终于死去。那是当然的，有那么一大群“小侏儒”在它
身上吸血呢！毛虫们尽职地活着，直到身体里的“小侏儒”准备出来的时候。它们从毛
虫的身体里出来后就开始织茧，最后，变成蛾，破茧而出。
------------------








孔雀蛾



孔雀蛾是一种长得很漂亮的蛾。它们中最大的来自欧洲，全身披着红棕色的绒毛，
脖子上有一个白色的领结，翅膀上洒着灰色和褐色的小点儿。横贯中间的是一条淡淡的
锯齿形的线，翅膀周围有一圈灰白色的边，中央有一个大眼睛，有黑得发亮的瞳孔和许
多色彩镶成的眼帘，包括黑色、白色、栗色和紫色的弧形线条。这种蛾是由一种长得极
为漂亮的毛虫变来的，它们的身体以黄色为底色，上面嵌着蓝色的珠子。它们靠吃杏叶
为生。
五月六日的早晨，在我的昆虫实验室里的桌子上，我看着一只雌的孔雀蛾从茧子里
钻出来。我马上把它罩在一个金属丝做的钟罩里。我这么做没有别的什么目的，只是一
种习惯而已。我总是喜欢搜集一些新鲜的事物，把它们放到透明的钟罩里细细欣赏。
后来我很为自己的这种方法庆幸。因为我获得了意想不到的收获，在晚上九点钟左
右，当大家都准备上床睡觉的时候，隔壁的房间里突然发出很大的声响。
小保罗衣服都没穿好，在屋里奔来跑去，疯狂地跳着、顿着足、敲着椅子。我听到
他在叫我：
“快来快来！”他喊道，“快来看这些蛾子，像鸟一样大，满房间都是！”
我赶紧跑进去一看，孩子的话一点儿也不夸张。房间里的确充满了那种大蛾子，已
经有四只被捉住关在笼子里了，其余的拍打着翅膀在天花板下面翱翔。
看到这情形，我立即想起那只早上被我关起来的囚徒。
“快穿好衣服”，我对儿子说，“把鸟笼放下，跟我来。我们立刻就要看到更有趣
的事情了。”
我们立刻下楼，来到我的书房，那在整个房子的右侧。我发现厨房里的仆人已被这
突然发生的事件吓慌了，她用她的围裙扑打着这些大蛾，起初她还以为它们是蝙蝠呢。
这样看来，孔雀蛾们已经占据了我家里的每一部分，惊动了家里的每一个人。
我们点着蜡烛走进书房，书房的一扇窗开着。我们看到了难忘的一幕情景：那些大
蛾子轻轻地拍着翅膀，绕着那钟罩飞来飞去。一会儿飞上，一会儿飞下，一会儿飞出去，
一会儿又飞回来，一会儿冲到天花板上，一会儿又俯冲下来。它们向蜡烛扑来，用翅膀
把它扑灭。它们停在我们的肩上，扯我们的衣服，咬我们的脸。小保罗紧紧地握着我的
手，努力保持镇定。
一共有多少蛾子？这个房间里大约有二十只，加上别的房间里的，至少在四十只以
上。四十个情人来向这位那天早晨才出生的新娘致敬——这位关在象牙塔里的公主！
在那一个星期里，每天晚上这些大蛾总要来朝见它们美丽的公主。那时候正是暴风
雨的季节，晚上黑得伸手不见五指。我们的屋子又被遮蔽在许多大树后面，很难找到。
它们经过这么黑暗和艰难的路程，历尽困苦来见它们的女王。
在这样恶劣的天气条件下，连那凶狠强壮的猫头鹰都不敢轻易离开巢，可孔雀蛾却
能果断地飞出来，而且不受树枝的阻挡，顺利到达目的地。它们是那样的无畏，那样的
执着，以至于到达目的地的时候，它身上没有一个地方被刮伤，哪怕是细微的小伤口也
没有。这个黑夜对它们来说，如同大白天一般。
孔雀蛾一生中唯一的目的就是找配偶，为了这一目标，它们继承了一种很特别的天
赋：不管路途多么远，路上怎样黑暗，途中有多少障碍，它总能找到它的对象。在它们
的一生中大概有两三个晚上它们可以每晚花费几个小时去找它们的对象。如果在这期间
它们找不到对象。那么它的一生也将结束了。
孔雀蛾不懂得吃。当许多别的蛾成群结队地在花园里飞来飞去吮吸蜜汁的时候，它
从不会想到吃东西这回事。这样，它的寿命当然是不会长的了，只不过是两三天的时间，
只来得及找一个伴侣而已。
------------------








找枯露菌的甲虫



在讲到甲虫之前，让我先来讲一下我的狗朋友，它会找枯露菌。所谓枯露菌，指的
是一种长在地底下的蘑菇。狗常常被用来做这种工作。我的狗的运气极好，有好几次跟
着一只在这方面极有经验的狗一同出去工作。而那只狗，那位我急于见识一下的找蘑菇
专家，其外貌实在没有什么可取的地方：一只极为普通的狗，态度平静而从容，又丑又
不整洁。总之，是一种你绝对不会让它歇在你的火炉边的狗，可它的的确确是一个名符
其实的找蘑菇专家。许多道理和人类世界中的一样：天才和贫穷总是连在一起的。
这只狗的主人，是村里有名的枯露菌商。他起初怀疑我要窥探他的秘密从而和他进
行商业竞争，后来有人告诉他我只是要采集地下植物的标本，要借他的狗用一用，他相
信了，于是就允许我和他一同出发去工作。
我们有言在先，规定任何一方都不可以干涉狗的行为，而且只要它发现一种菌类，
不管那是人们喜欢吃的蘑菇还是其它不可以吃的蘑菇，它都理应得到一片面包作为酬劳。
另外，它的主人不能禁止他的狗到它喜欢的地方去，即使他深知那个地方的蘑菇绝对卖
不出去。因为对我的研究课题而言，这种蘑菇是否可吃并不重要，我的目的和枯露菌商
人的有所不同。
遵循这个原则，这次远征获得了极大的成功。这只忙碌的狗一路上踱着慢慢的步子
走着，用鼻子嗅着。每走几步，它就要停下来，用鼻子测试着泥土。它用鼻子扒几下土，
然后泰然地望着主人，似乎在说：
“就在这里了！就在这里了！我以我的名义担保！这里有蘑菇！”
它说的果然没错。主人依着它所指示的方向掘下去。万一主人的铲掘得偏了，它会
赶紧发出一声鼻音，提醒主人如何把铲子放到正确的部位。这样掘下去没有一次落空。
狗的鼻子果然名不虚传，从不说谎，不过它指示给我们的，包括各种地下菌类：大的小
的，有气味的，没气味的……当我收集着这些蘑菇的时候，我非常惊奇，这里面几乎包
括了附近一带所有的地下蘑菇品种。
是不是那种我们常说的嗅觉在帮助狗找寻呢？我不相信，如果完全靠嗅觉，它决不
能找出这许多气味完全不同的菌类来。它一定有一种我们所没有的感觉。通常我们用人
类的标准去推测一切事物的时候，往往容易犯错。在这个世界上，有许多种感觉是我们
人类所不知道的。而这种感觉，在昆虫中更加明显。
我们现在就来讲讲找蘑菇的甲虫吧。
这是一种美丽的甲虫，小小的，黑黑的，有一个白绒肚皮，形状是圆的，像一粒樱
桃的核，当它用翅膀的边缘擦着腹部的时候，就会发出一种柔软的“唧唧”声，就像小
鸟看见母亲带着食物回来时所发出的声音一样。雄的甲虫头上还长着一个美丽的角。
我是在一个长满蘑菇的松树林里发现这种甲虫的。那是一个美丽的地方，在秋季气
候温和的日子里，我们全家都喜欢到那儿去玩。那个地方几乎什么都有：有用细枝做成
的老喜鹊的巢；有饶舌的柽鸟吃饱后在树上互相嬉戏、追逐；有兔子突然从树丛里跳出
来，翘着它们的短尾巴；那里还有可爱的小河，可以给孩子们筑隧道，这种河很容易堆
成一排小屋，我们用草盖在屋顶上，算是草屋，一段芦草插上屋顶便成为烟囱。微风轻
轻地掠过松针，发出轻轻的叹息声，我们就在这美妙的音乐中开始了我们的午餐。
的确，对于小孩子们来说，这是个真正的乐园。即使是成人，也会喜欢这个地方。
我到这里来最大的乐趣便是守候那些找蘑菇的甲虫们，它们的洞到处都可以看到，而且
门是开着的，不过在洞口堆着一堆疏松的泥土。洞大约有几寸深，一直向下，而且往往
筑在比较松的泥土中。当我用小刀一直挖下去的时候，我总是发现这种洞是空的，甲虫
们已经乘着夜色离开这里了，它在这里做完了它的工作后，便迁到别处去了。这种甲虫
是个流浪者，并且是个夜行客，随便什么时候，它想离开这个洞的时候，它就能很容易
地在别处另筑新巢。有时候我侥幸能在洞底发现甲虫，但永远只有一个，或是雌的或是
雄的，从不会成对。看来这种洞并不是一个家庭的所在地，而是专门给独身的甲虫住的。
你看，这洞里的甲虫正在啃着一个小蘑菇，已经吃完了一部分。它虽然已经累了，
但仍旧紧紧地抱着它，它是决不肯轻易放弃这个蘑菇的，这是它的宝贝，是它一生中的
最爱，从周围许多吃剩的碎片来看，这只甲虫已经吃得饱饱的了。
当我从它手中夺过这宝物的时候，我发现这是一种很小的地下菌，跟枯露菌很相像。
这个事实似乎可以解释甲虫的习惯和它常要换新居的理由。让我们想象一下吧，在
静静的黄昏中，这个小旅行家便从它的洞里慢慢地跟着步走出来，一边快活地唱着歌，
一边悠闲地散着步。它仔细地检查着土地，探究这地底下所埋的东西，就像狗找枯露菌
一样。它的嗅觉告诉它哪个地方有菌，只不过盖着几寸泥土而已；那个地方虽然泥土肥
沃，但地底下决不会有菌类。当它判定在某一点下面有菌的时候，便一直往下挖，结果
总能得到它的食物。它挖的洞也成了它的临时宿舍，在食物没有吃完之前，它是不会离
开洞的，它会在自己掘的洞底快活地吃着，管它洞门是开着的还是关着的。
等到洞里的食物都吃完后，它就要搬家了。它会在别处找一个适当的地方，再掘下
去，然后住一阵子，吃一阵子，等到新屋里的食物吃完了，它就再搬一次家。在整个秋
季到来年的春季——菌类的生长季节里，它就这样游历着，“打一枪换一个地方”，从
一个洞搬到另一个洞，很辛苦又很洒脱。
甲虫所猎取的菌并没有特殊的气味，那么它是怎么能从地面上检查出地底下菌类的
存在呢？它是聪明的甲虫，它自有办法。可是我们人类就望尘莫及了，哪怕是“千里眼”
或是“顺风耳”，也无法说出隐藏在地底下的秘密。
------------------








爱好昆虫的孩子



现在，有许多人总喜欢把一切人的品格、才能、爱好等归于遗传。也就是说承认人
类及一切动物的智慧都是从祖先那儿得来的。我并不完全同意这种观点。我现在就用我
自己的故事来证明我那喜爱昆虫的嗜好并不是从哪个先辈身上继承下来的。
我的外祖父和外祖母从来没有对昆虫产生过丝毫的兴趣和好感。关于我的外祖父，
我不大知道，我只知道他曾经历过相当苦难的日子。我敢说，如果要说他曾经和昆虫发
生过关系的话，那就是他曾一脚把它踩死。外祖母是不识字的文盲，每天为琐碎的家务
所累，没有什么闲情雅致去欣赏一些风花雪月的故事，对于科学或昆虫当然更不会产生
兴趣。当她蹲在水龙头下洗菜的时候，偶尔会发现菜叶上有一条毛虫，她会立刻把这又
讨厌又可恶的东西打掉。
关于我的祖父母，我知道的比较详细。因为我小时候，我的父母穷得无法养活我，
所以在五六岁的时候，我就跟着祖父母一同生活了。祖父母的家在偏僻的乡村里，他们
靠着几亩薄田维持生计。他们不识字，一生中从没有模过课本。祖父对于牛和羊知道得
很多，可是除此之外，其它的便什么也不知道了。如果他知道在将来他家里的一个人花
费了许多时间去研究那些渺小的、微不足道的昆虫，他会多么地吃惊啊！如果他再知道
那个疯狂的人正是坐在他旁边的小孙子，他将一定会愤怒地给我一巴掌的。
“哼，把时间和力气花费在这种没出息的东西上！”他一定会怒吼。
祖母是一个可爱的人，她整天忙着洗衣服，照顾孩子、烧饭、纺纱、看小鸭、做乳
酪和奶油，以及其它一些家务，一心为这个家操劳。
有时候，在晚上，当我们都坐在火炉边的时候，她就会常常讲一些狼的故事给我们
听。我很想见一见这匹狼，这位在一切故事里使人心惊肉跳的英雄，可是我从来没有亲
眼见过一次。亲爱的祖母，我是始终深深地感激着您的。在您的膝上，我第一次得到了
温柔的安慰，使痛苦和忧伤得到缓解。你遗传给了我强壮的体质和爱好工作的品格，可
是你的确没有给我爱好昆虫的天性。
我自己的父母都是不爱好昆虫的。母亲没有受过教育，父亲小时候虽然进过学校，
稍稍能读能写，可是为了生活整天忙得不可开交，再也没有时间顾及到别的事情了，更
谈不上爱好昆虫了。有一次当他看到我把一只虫子钉在软木上的时候，他狠狠地打了我
一拳，这就是我从他那里得到的鼓励。
尽管如此，从幼年的时候开始，我就喜欢观察和怀疑一切事物。每次忆起童年，我
总会想起一件难忘的往事，现在说起来还觉得很有趣。在我五六岁的时候，有一天我光
着脚丫子站在我们的田地前面的荒地上，粗糙的石子刺痛了我。我记得我有一块用绳子
系在腰间的手帕——很惭愧，我那时常常遗失手帕，然后用袖子代替它，所以不得不把
宝贵的手帕系在腰上。
我把脸转向太阳，那眩目的光辉使我心醉。这种光辉对我的吸引力相当于光对于任
何一只蛾子的吸引力甚至还要大得多，当我这样站着的时候，我的脑海里就突然冒出一
个问题：我究竟在用哪个器官来欣赏这灿烂的光辉？是嘴巴？还是眼睛？请读者千万不
要见笑，这的确算得上一种科学的怀疑。我把嘴也张得大大的，又把眼睛闭起来，光明
消失了；我张开眼睛闭上嘴巴：光明又出现了。这样反复试验了几次，结果都是一样。
于是我的问题被我自己解决了：我确定我看太阳用的是眼睛，后来我才知道这种方法叫
“演绎法”。这是一个多么伟大的发现啊！晚上我兴奋地把这件事告诉大家。对于我这
种幼稚和天真，只有祖母慈祥地微笑着，其余的人都大笑不止。
另外一次是在黑夜的树林里，有一种断断续续的叮当声大大地引起了我的注意。这
种声音显得分外优美而柔和。在寂静的夜里，是谁在发出这种声音？是不是巢里的小鸟
在叫？还是小虫子们在开演唱会呢？
“哦，我们快去看看吧，那很可能是一只狼。狼的确是在这种时候出声的，”同行
的人对我说，“我们一起走，但不要走得太远，声音就是从那一堆黑沉沉的木头后面发
出来的。”
我站在那里守候了许多时候，什么也没有。后来树林中发出一个轻微的响声，仿佛
是谁动了一下，接着那叮当声也消失了。第二天，第三天，我再去守候，不发现真相决
不罢休。我这种不屈不挠的精神终于获得了回报。嘿！终于抓到它了，这一个音乐家已
经在我的股掌之间了。它不是一只鸟，而是一只蚱蜢，我的同伴曾告诉我它的后腿非常
鲜美。这就是我守候了那么久所得到的微乎其微的回报。不过我所得意的，倒不是那两
只像虾肉一样鲜美的大胆，而是我又学到了一种知识，而且，这知识是我亲自通过努力
得来的。现在，从我个人的观察来看，我知道蚱蜢是会唱歌的。我没有把这发现告诉别
人，为的是怕再像上次看太阳的事情那样遭别人的嘲笑。
哦，我们屋子旁边的花长得多么美丽啊！它们好像张着一个个彩色的大眼睛向我甜
甜地笑。后来，在那个地方，我又看到一堆堆又大又红的樱桃。我尝了尝，滋味也不过
如此，没有像看上去的那么诱人，而且没有核，这究竟是些什么樱桃呢？夏天将要结束
的时候，祖父拿着铁锹来，把这块土地的泥土从底下翻起来，从地底下掘出了许许多多
圆圆的根。我认得那种根，在我们的屋子里面有许许多多，我们时时把它们放在煤炉上
煨着吃。那就是马铃薯。普通得不能再普通的马铃薯。我的探索一下戛然而止，不过，
那紫色的花和红色的果子被永远地留在了我的记忆里。
我利用自己这双对于动植物特别机警的眼睛，独自观察着一切惊异的事物。尽管那
时候我只有六岁，在别人看来什么也不懂。我研究花，研究虫子；我观察着，怀疑着；
不是受到了遗传的影响，而是受到了好奇心的驱使和对大自然的热爱。
不久我又回到了我父亲的屋子里。那时候我已经七岁，到了必须进学校的年龄了。
可我并不觉得学校生活比我以前那种自由自在地沉浸在大自然中的生活更有意思。我的
教父就是老师。那间我坐在里面学习字母的屋子，我该称它做什么呢？的确很难有一个
恰如其分的名字，因为那屋子用处太多了。它既是学校，又是厨房；既是卧室，又是餐
厅；既是鸡窝，又是猪圈。在那种时代，谁也不会梦想有王宫般美丽堂皇的学校，无论
什么破棚子都可以被认为是最理想的学校。
在这间屋子里，一张很宽的梯子通到楼上去，梯子脚边的一个凹形的房间里有一张
大床。楼上究竟有什么东西我不大知道。有时候我们会看见老师从楼上捧来一捆干草给
驴子吃；有时候也会看见他从楼上提着一篮马铃薯下来，交给师母去煮猪食。我猜这一
定是一间堆放物品的屋子，是人和畜牲共同的储藏室。
让我们回过头来讲那间做教室的房间吧。这间屋子里唯一的一扇窗，是一扇朝南的
窗，又小又矮。当你的头碰着窗顶的时候，你的肩膀同时地碰到了窗栏。这个透着阳光
的窗户是这个屋子里惟一有生气的地方，它俯视着这个村庄的大部分。从窗口往外望，
你会发现这是一个散落在斜坡上的村落。窗口下面是老师的小桌子。
对面墙上有一个壁龛，里面放着一个盛满水的发亮的铜壶，孩子们口渴的时候可以
信手从这里倒杯水解渴。在壁龛的顶端有几个架子，上面放着闪闪发光的碗，那些碗只
有在举办盛会时才拿出来用。
在光线所能射到的墙壁上，到处挂着各种色彩不协调的图画，最远的那垛墙边有一
只大壁炉，左右是用木石筑成的，上面放着塞了糠的垫褥。由两块滑动的板充当门。如
果想独自静静地躺下睡觉，你可以把门关起来。这两张床是给主人和主妇睡的。无论北
风在黑暗的谷口怎样怒吼，无论雪花在外面如何打转，他们一定在这里面睡得很舒服。
其余的地方就放着一些零碎的杂物：一条三脚凳，一只挂在壁上的盐罐，一个铁铲，重
得需要两只手一起使劲才拿得动，最后还有那风箱，就像我祖父家里的那个一样，风一
吹，炉里的木块和树枝就烧起来。我们如果要享受火炉的温暖，每人每天早上就得带一
块小柴来。
可是炉子并不是为我们生的，主要是为那三只煮猪食的锅子。老师和师母总是挑一
个最舒适的位于坐下，其余的人却围着那大锅子，围成一个半圆形。那锅里不住地冒着
热气，发出呼呼的声音。我们中胆子比较大的人会趁着先生看不见的时候用小刀挑一个
煮熟了的马铃薯，夹在他的面包里吃。我不得不承认，如果说我们在学校里做了一些工
作的话，那就是我们吃得很多。在写字的时候剥着栗子或咬着面包，似乎已经成为改不
了的习惯了。
至于我们这些年纪较小的学生，除了享受满口含着食物读书的乐趣外，还有两件快
乐的事情，在我看来不见得比栗子的味道差。我们的教室后门外就是庭院。在那里，一
群小鸡围着母鸡在扒土，小猪们自由自在地打着滚。有时候，我们中会有人偷偷地溜出
去，回来的时候故意不把门关上，于是马铃薯的香味便一阵阵地飘到门外去。外面的小
猪闻到这香味，一个个循着香味接通跑来。我的长凳子，是年纪最小的学生们坐的，恰
巧是靠着墙壁，在铜壶的下方，也正是小猪们的必经之路。我每次都能看见小猪们快步
地跑着，一边大声地呼叫着，摇着它们的小尾巴。它们用身体蹭我们的腿，把又冷又红
的鼻子拱到我们的手掌里面找吃剩的面包屑。它们那细小的圆溜溜的眼珠子望着我们，
似乎在问我们口袋里还有没有干栗子给它们吃。它们这样东窜西闻了一圈之后，就会被
怒气冲冲的老师挥着手帕赶回院子里去了。
接着就是母鸡带着它的小鸡雏们来看我们了。我们每个人都会热情地剥一些面包来
招待这些可爱的小客人，然后美滋滋地欣赏它们吃东西的样子。
在这样的一个学校里，我们能学到些什么呢？每一个年纪较小的学生手里都有——
或者说，假定他们都有一本灰纸订成的小书，上面印着字母，封面上画着一只鸽子，确
切地说，那只是一只很像鸽子的动物。封面上有一个十字架，是用字母按照一定的顺序
排出来的。老师可能觉得这本书挺管用的，因此把书发给我们，并且解释给我们听。就
因为这样，老师总是被那些年纪较大的学生们缠着，没有工夫顾及我们这些小不点儿。
他还是把书也发给我们，不过其作用只是为了让我们看上去更像学生而已。于是我们这
些小不点儿就自己坐在长凳上读书，同时请旁边的大孩子教我们——如果他能认得一两
个字母的话。我们的学习常常被一些无足轻重的小事打断，一会儿老师和师母去看锅里
的马铃薯了，一会儿小猪的同伴们叫唤着进来，一会儿又是一群小鸡忙不迭地奔进来，
我们这样忙里偷闲地看一会儿书，实在是学不到什么东西，得不到什么知识的。
大孩子们常常练习写字。他们的位置比较优越，能够借着从那狭小的窗洞透进来的
光，并且他们前面还有那张全屋唯一的桌子，学校里什么设备都没有，甚至连墨水都没
有一滴，所以每个人上学的时候都得自己带上全套的文具。那时候的墨水是装在一个长
条的纸板匣里面的，里面分成两格，上面一格放鹅毛管做的笔，下面一格是一个小小的
墨水池，里面盛着墨水，那时候的墨水是用烟灰和醋合成的。
我们的老师最伟大的工作就是修笔——然后在某一页的顶上写一行字母或是单字，
至于他写什么字，依照各个学生的要求而定。看，当老师写字的时候，他的手腕抖动得
多厉害啊！他用小拇指拈着纸，做好奋笔疾书的准备。忽然，老师的手开始运动了，在
纸上飞着、打着转。看啊，笔尖所到之处展开了一条花边，里面有圆圈、有螺旋、有花
体字、有张着翅膀的鸟……我们要他画什么他就画什么。只要你喜欢，什么都有。这些
画都是用红墨水画出来的。就是这样一只笔创造了一个个奇迹。面对这一个个奇迹，我
们都惊得目瞪口呆。
我们在学校里读些什么呢？大概是法文吧，常常从圣经上记载的历史中选一两段来
读读。拉丁语倒是学得比较多些，为的是使我们能够准确地唱赞美诗。
历史、地理呢？谁也没听到过这两个名词。地球是方的还是圆的，对我们来说有什
么不同呢？方也罢，圆也罢，反正从地里长出东西来是同样的不容易啊！
语法呢？我们的老师从来不拿这个问题去为难自己，我们当然更不会了。
数学呢？是的，我们的确学了一点，不过还不配用那么堂皇的名字，我们一直称它
为“算术”。
在星期六的晚上，通常总是用“算术”的仪式来结束这一星期。最优秀的学生先站
起来把乘法口诀表背诵一遍，然后全班，包括最小的学生，依着他的样子齐声合背一遍。
我们的声音很响亮，把偶而跑进屋来想觅一点食的鸡和猪都吓跑了。
别人都说我们的老师是个很能干的人，能把学校管理得很好。的确，他不是一个等
闲之辈，但他也确实不能称作一个好老师，因为缺少一样东西——时间。他替一个出门
的地主保管着财产；他照顾着一个极大鸽棚；他还负责指挥干草、苹果、栗子和燕麦的
收获，在夏天，我们常常帮着他干活。在那个时候，上课才是一件有趣的事，因为我们
常在干草堆上上课，有时候甚至就利用上课的时候清除鸽棚，或是消灭那些雨天从墙脚
爬出来的蜗牛。这对我来说，倒是正中下怀。
我们的老师还是个剃头匠。他那双灵巧的手替我们的抄写本装饰“花边”，也为地
方上的大人物剃头，像市长、牧师和公证人等等。
我们的老师又是个打钟的能手。每逢有婚礼或洗礼的时候，他总要到教堂里去打钟
——那时我们的功课当然要暂告停止。暴风雨来临的时候，又可以给我们一天的休假，
因为那时候必须用钟声来驱除雷电和冰雹，我们的老师责无旁贷地去敲钟了。
我们的老师还是唱诗班里的一员。我们的老师还管着村里教堂顶上的钟。那是他最
引以为豪的工作。只须对着太阳一望，他便可以说出一个准确的时间，然后他爬到教堂
顶上尖尖的阁楼里，打开一个大匣子，让自己置身于一堆齿轮和发条中间。这些东西的
秘密，除了他之外，没有第二个人知道。
这样一个学校，这样一个老师，对于我那尚未充分表现的特点，将有什么影响呢？
我那热爱昆虫的个性，几乎不得不渐渐地枯萎以至永远消失了。但是，事实上，这种个
性的种子有着很强的活力，它永远在我的血液里流动，从来没有离开过我。它能够随时
激发出来或找到滋生的养料，无时无刻不体现出来，甚至在我的教科书的封面上，也能
显而易见地看出书的主人的爱好——那里有着一只色彩配合得不很协调的鸽子，它对于
我来说，比书本里的ABC有意思得多。它的圆眼睛似乎在冲着我笑，它那翅膀我已一根一
根地数过共有多少羽毛，那些羽毛告诉我怎样飞上天空，翱翔在美丽的云朵里。这只鸽
子带着我飞到毛榉树上，我看到那些透着光泽的树干高高地矗立在长满苔藓的泥土上在
泥土上。在长着许多白色的蘑菇，看上去好像是过路的母鸡产下的蛋。这只鸽子又带我
到积雪的山顶上，在那里，鸟类用它们的红脚踏出了星形的足迹。这个鸽子是我的好伙
伴、好朋友，它减轻了我整天背字母的压力。应该谢谢它，有了它作伴，我才能静静地
坐在长凳上等候放学。
露天学校有着更大的诱惑力。当老师带着我们去消灭黄杨树下的蜗牛的时候，我却
常常阴奉阳违，不忍心杀害那些小生命。当我捉到了满手的蜗牛时，我的脚步便迟缓起
来了。它们是多么美丽啊！只要我愿意，我能捉到各种颜色的蜗牛：黄色的、淡红色的、
白色的、褐色的……上面都有深色的螺旋纹。我挑了一些最美丽的塞满衣兜，以便空闲
的时候拿出来看看。
在帮先生晒干草的日子里，我又认识了青蛙。它用自己作诱饵，引诱着河边巢里的
虾出来；在赤杨树上，我捉到了青甲虫，它的美丽使天空都为之逊色；我采下水仙花，
并且学会了用舌尖从它花冠的裂缝处吸取小滴的蜜汁，我也体验到太用力吸花蜜所导致
的头痛，不过这种不舒服与那美丽的白色花朵所带给我的赏心悦目的感觉相比，实在是
太微不足道了。我还记得这种花的漏斗的颈部有一圈美丽的红色，像挂了一串红项链。
在收集胡桃的时候，我在一块荒芜的草地上找到了蝗虫，它们的翅膀张得像一把扇
子，有红色的也有蓝色的，让人眼花缭乱。无论在什么地方，我都能源源不断地得到精
神食粮，自得其乐。我对于动植物的爱好也自然有增无减，日益弥深。
最后，这种爱好促进了我对字母的认识。由于我太喜欢封面上的鸽子，早把封面后
的字母丢到九霄云外去了。所以我的认识程度一直停留在初级阶段。一个偶然的念头使
我的父亲把我从学校里领回家去，这才是我真正读书的开始。这回读的是印得很大的字，
花了三角半钱买来的。那上面划着许多五彩的格子，每一格里画着一种动物，这些动物
就用它们的名字和第一个字母来教我认ABC。第一个就是“驴”。在法文中，它的名字是
Ane，于是我认识了A；牛的名字叫Boeuf，它教我认识了B；Canara是鸭子，于是我认识
了C；Dinod是火鸡，它教我认识了D。其余的字母也是如此这般让我认识的。当然，有几
格印得很不清楚，像那教我认得H、K和E的河马(Hippopotamus)、雨燕(Kamichi)和瘤牛
(zebu)之类，不过没什么大碍。我进步得很快，不到几天工夫，居然能很有兴趣地读那
本鸽子封面的书了。我已经被启发了，接着便懂得语法了，这激起了我对学习的浓厚兴
趣，我的父母都为我的进步感到惊异。现在我能够解释那惊人的进步的原因了：那些图
画把我引入到一群动物中，这恰巧投合了我的兴趣。我心爱的动物们开始教我念书，而
以后，动物永远成为我学习研究的对象。
后来，好运第二次降临到我身上。为了让我用功读书，我得到了一本廉价的《拉封
丹寓言》，里面有许多插图，虽然又小又不准确，可是看起来的确很有趣。这里有乌鸦、
喜鹊、青蛙、兔子、驴子、猫和狗；这些都是我所熟悉的东西，这里面的动物会走路会
讲话，因此大大激起了我的兴趣。至于了解这本书究竟讲了些什么，那是另一回事了。
不过不要担心，我试着把一个个音节连起来，慢慢地你就知道全篇的意思了。于是拉封
丹也成为我的朋友了。
十岁的时候，我已是路德士书院(RodezColloge)的学生了。我在那里成绩很好，尤
其是作文和翻译两课都能得到很高的分数。在那种古典派的气氛中，我们听到了许多神
话故事。那些故事都是很吸引人的。可是在崇拜那些英雄之余，我不会忘记趁着星期天
去看看莲香花和水仙花有没有在草地上出现；梅花雀有没有在榆树丝里孵卵；金虫是不
是在摇摆于微风中的白杨树上跳跃，无论如何，我是不能忘记它们的！
可是，忽然恶运又降临了：饥饿威胁着我们一家。父母再也没有钱供我念书了。我
不得不离开学校。生命几乎变得像地狱一样可怕。我什么都不想，只盼望能快快熬过这
段时期！
在这些悲惨的日子里，我对于昆虫的偏爱应该暂时搁在一边了吧？就像我的先辈那
样，为生计所累。但是，事实并非如此，我仍然常常能够回忆起那只第一次遇到的金虫：
它那触须上的羽毛，它那美丽的花色——褐色底子上嵌着白点——这些好像是那种凄惨
晦暗的日子里的一道闪亮的阳光，照亮并温暖了我悲伤的心。
总而言之，好运不会抛弃勇敢的人。后来我又进了在伏克罗斯(Voncluse)的初级师
范学校，在那里我能免费分到食物，尽管只是干栗子和豌豆而已，校长是位极有见识的
人，他不久便信任了我，并且给了我完全的自由。他说只要我能应付学校里的课程，我
几乎可以随心所欲做自己喜欢的事情。而当时我的程度比同班的同学要稍高一些，于是
我就利用比别人多的空闲时间来增加自己对动植物的认识。当周围的同学们都在订正背
书的错误时，我却可以在书桌的角落里观察夹竹桃的果子、金鱼草种子的壳，还有黄蜂
的刺和地甲虫们的翅膀。
我对于自然科学的兴趣，就这样慢慢地滋长起来了。在那时候，生物学是被一般学
者所轻看的学科，学校方面所承认的必修课程是拉丁文、希腊文和数学。
于是我竭尽全力地去研究高等数学。这是一种艰难的奋斗，没有老师的指导，碰到
疑难问题，往往好几天得不到解决，可我一直坚持不懈地学着，从未想过半途而废，而
终于有所成就。后来我用同样的方法自学了物理学，用一套我自己制造的简陋的仪器来
做各种实验。我违背了自己的志愿，把我的生物学书籍一直埋在箱底。
毕业后，我被派到埃杰克索书院(AjaccioCollege)去教物理和化学。那个地方离大
海不远，这对我的诱惑力实在太大了。那包蕴着无数新奇事物的海洋，那海滩上美丽的
贝壳，还有番石榴树、杨梅树和其它一些树，都足够让我研究好半天的。这乐园里美丽
的东西比起那些三角、几何定理来，吸引力大得多了。可是我努力控制着自己。我把我
的课余时间分成两部分：大部分时间用来研究数学；小部分的时间用来研究植物和搜寻
海洋里丰富的宝藏。
我们谁都不能预测未来。回顾我的一生，数学，我年轻时花费了那么多时间和精力
去钻研，结果对我却没有丝毫的用处；而动物，我竭力想方没法地回避它，在我的老年
生活中，它却成了我的慰藉。
在埃杰克索，我碰到两位著名的科学家：瑞昆(Rrguien)和莫昆·坦顿(MoquinTand
on)，瑞昆是一位著名的植物学家；而莫昆·坦顿教了我植物学的第一课。那时他因为没
有旅馆住而寄住在我的房子里。在他离开的前一天，他对我说，“你对贝壳很感兴趣，
这当然很好。不过这样还远远不够。你应当知道动物本身的组织结构，让我来指给你看
吧！这会使你对动物的认识提高到一个新的水平。”
他拿起一把很锋利的剪刀和一对针，把一个蜗牛放在一个盛水的碟子里，开始解剖
给我看。他一边解剖，一边一步步地把各部分器官解释给我听。这就是我一生中所得到
的最难以忘怀的一堂生物课，从此，当我观察动物时，不再仅仅局限在表面上了。
现在我应该把自己的故事结束了。从我的故事里可以看出，早在幼年时期，我就有
着对大自然的偏爱。而且我具有善于观察的天赋。为什么我有这种天赋？怎样才会有？
我自己也说不清楚。
无论是人还是动物，都有一种特殊的天赋：一个孩子可能有音乐的天赋，一个孩子
可能在雕塑方面很有天赋，而另一个孩子可能是速算的天才。昆虫也是这样，一种蜜蜂
生来就会剪叶子，另一种蜜蜂会造泥屋，而蜘蛛则会织网。为什么它们有这种才能？天
生就有的，除此之外就没有什么理由可解释了。在人类生活中，我们称这样的人才为
“天才”；在昆虫中，我们称这样的本领为“本能”。本能，其实就是动物的天才。
------------------








条纹蜘蛛



不管是谁，大概都不会喜欢冬季。在这个季节里，许多虫子都在冬眠。不过这并不
说明你没有什么有虫子可观察了。这时候如果有一个观察者在阳光所能照到的沙地里寻
找，或是搬开地下的石头，或是在树林里搜索，他总能找到一种非常有趣的东西，那是
一件真正的艺术品。那些有幸欣赏到这艺术作品的人真是幸福。在一年将要结束的时候，
发现这种艺术品的喜悦使我忘记了一切不快，忘记了一天比一天更糟的气候。
如果有人在野草丛里或柳树丛里搜索的话，我祝福他能找到一种神秘的东西：这是
条纹蜘蛛的巢。正像我眼前所呈现的一样。
无论从举止还是从颜色上讲，条纹蜘蛛是我所知道的蜘蛛中最完美的一种。在它那
胖胖的像榛仔仁一般大小的身体上，有着黄、黑、银三色相间的条纹，所以它的名字叫
“条纹蜘蛛”。它们的八只脚环绕在身体周围，好像车轮的辐条。
几乎什么小虫子它都爱吃。不管那是蝗虫跳跃的地方还是苍蜂盘旋的地方，是蜻蜓
跳舞的地方还是蝴蝶飞翔的地方。只要它能找到攀网的地方，它就会立刻织起网来。它
常常把网横跨在小溪的两岸，因为那种地方猎物比较丰盛。有时候它也在长着小草的斜
坡上或榆树林里织网，因为那里是蚱蜢的乐园。
它捕获猎物的武器便是那张大网，网的周围攀在附近的树枝上。它的网和别种蜘蛛
的网差不多：放射形的蛛丝从中央向四周扩散，然后在这上面连续地盘上一圈圈的螺线，
从中央一直到边缘。整张网做得非常大，而且整齐对称。
在网的下半部，有一根又粗又宽的带子，从中心开始沿着辐一曲一折，直到边缘，
这是它的作品的标记，也是它在作品中的一种签名。同时这种粗的折线也能增加网的坚
固性。
网需要做得很牢固，因为有时候猎物的份量很重，它们一挣扎，很可能会把网撑破。
而蜘蛛自己不会选择或捕捉猎物，所以只能不断地改进自己的大网以捕获更多的猎物。
它静静地坐在网的中央，把八只脚撑开，为的是能感觉到网的每一个方向的动静。摆好
阵势后，它就等候着，看命运会赐予它什么：有时候是那种微弱到无力控制自己飞行的
小虫；有时候是那种强大而鲁莽的昆虫，在做高速飞行的时候一头撞在网上，有时候它
好几天一无所获，也有时候它的食物会丰盛得好几天都吃不完。
蝗虫，尤其是一种火蝗，它控制不了自己腿部的肌肉，于是常常跌进网中。你可能
会想，蜘蛛的网一定受不住蝗虫的冲撞，因为蝗虫的个儿头要比蜘蛛大得多，只要它用
脚一蹬，立刻就可以把网蹬出一个大洞，然后逃之天天。其实情况并不是这样的，如果
在第一下挣扎之后不能逃出的话，那么它就再也没有逃生的希望了。
条纹蜘蛛并不急于吃掉蝗虫，而是用它全部的丝囊同时射出丝花，再用后腿把射出
来的丝花捆起来。它的丝囊是制造丝的器官，上面有细孔，像喷水壶的莲蓬头一般。它
的后腿比其余的腿要长，而且能张得很开，所以射出的丝能分散得很开。这样，它从腿
间射出来的丝已不是一条条单独的丝了，而是一片丝，像一把云做的扇子，有着虹霓一
般的色彩。然后它就用两条后腿很快地交替着把这种薄片，或者说是裹尸布吧，一片片
地向蝗虫撤去，就这样把蝗虫完全缠住了。
这不由得让我想起了古时候的角斗士。每逢要和强大的野兽角斗的时候，他们总是
把一个网放在自己的左肩上，当野兽扑过来时，他右手一挥，就能敏捷地把网撒开，就
像能干的渔夫撒网捕鱼那样，把野兽困在网里，再加上三叉戟一刺，就把它的性命结果
了。
蜘蛛用的也是这种方法。而且它还有一个绝招是人类所没有的：它可以把自己制造
的丝制的锁链绵绵不断地缠到蝗虫身上，一副不够，第二副立即跟着抛上来，第三副、
第四副……直到它所有的丝用完为止，而人类的网只有一副，即使有很多的话，也不能
这么迅速地接连抛出去。
当那白丝网里的囚徒决定放弃抵抗、坐以待毙的时候，蜘蛛便得意洋洋地向它走过
去，它有一个比角斗士的三叉戟还厉害的武器，那就是它的毒牙。它用它的毒牙咬住蝗
虫，美滋滋地饱餐一顿，然后回到网中央，继续等待下一个自己送上门来的猎物。

蜘蛛的巢
蜘蛛在母性方面的表露甚至比猎取食物时所显示的天才更令人叹服。它的巢是一个
丝织的袋，它的卵就产在这个袋里。它这个巢要比鸟类的巢神秘，形状像一个倒置的气
球，大小和鸽蛋差不多，底部宽大，顶部狭小，顶部是削平的，围着一圈扇蛤形的边。
整个看来，这是一个用几根丝支持着的蛋形的物体。
巢的顶部是凹形的，上面像盖着一个丝盖碗。巢的其它部分都包着一层又厚又细嫩
的白缎子，点缀着一些丝带和一些褐色或黑色的花纹。我们立刻可以猜到这一层的白缎
子的作用，它是防水的，雨水或露水都不能浸透它。
为了防止里面的卵被冻坏，仅仅使巢远离地面或藏在枯草丛里是远远不够的，还必
须有一些专门的保暖设备。让我们用剪刀把包在外面的这层防雨缎子剪开来看看。在这
下面我们发现了一层红色的丝。这层丝不是像通常那样的纤维状，而是很蓬松的一束。
这种物质，比天鹅的绒毛还要软，比冬天的火炉还要暖和，它是未来的小蜘蛛们的安乐
床。小蜘蛛们在这张舒适的床上就不会受到寒冷空气的侵袭了。
在巢的中央有一个锤子一样的袋子，袋子的底部是圆的，顶部是方的，有一个柔嫩
的盖子盖在上面。这个袋子是用非常细软的缎子做成的，里面就藏着蜘蛛的卵。蜘蛛的
卵是一种极小的橘黄色的颗粒，聚集在一块儿，拼成一颗豌豆大小的圆球。这些是蜘蛛
的宝贝，母蜘蛛必须保护着它们不受冷空气的侵袭。
那么蜘蛛是怎样造就这样精致的袋子的呢？让我们来看看它做袋子时的情形吧！它
做袋子的时候，慢慢地绕着圈子，同时放出一根丝，它的后腿把丝拉出来叠在上一个圈
的丝上面，就这样一圈圈地加上去，就织成了一个小袋子。袋子与巢之间用丝线连着，
这样使袋口可以张开。袋的大小恰好能装下全部的卵而不留一点空隙，也不知道蜘蛛妈
妈如何能掌握得那么精确。
产完卵后，蜘蛛的丝囊又要开始运作了。但这次工作和以前不同。只见它先把身体
放下，接触到某一点，然后把身体抬起来，再放下，接触到另一点，就这样一会儿在这，
一会儿在那，一会儿上，一会儿下，毫无规则，同时它的后脚拉扯着放出来的丝。这种
工作的结果，不是织出一块美丽的绸缎，而是造就一张杂乱无章、错综复杂的网。
接着它射出一种红棕色的丝，这种丝非常细软。它用后腿把丝压严实，包在巢的外
面。
然后它再一次变换材料，又放出白色的丝，包在巢的外侧，使巢的外面又多了一层
白色的外套。而且，这时候巢已经像个小气球了，上端小，下端大，接着它再放出各种
颜色不同的丝，赤色、褐色、灰色、黑色……让你目不暇接，它就用这种华丽的丝线来
装饰它的巢。直到这一步结束，整个工作才算大功告成了。
蜘蛛开着一个多么神奇的纱厂啊！靠着这个简单而永恒的工厂——它可以交替做着
搓绳、纺线、织布、织丝带等各种工作，而这里面的全部机器只是它的后腿和丝囊。它
是怎样随心所欲地变换“工种”的呢？它又是怎样随心所欲地抽出自已想要的颜色的丝
呢？我只能看到这些结果，却不知道其中的奥妙。
建巢的工作完成后，蜘蛛就头也不回地跨着慢步走开了。再也不会回来，不是它狠
心，而是它真的不需要再操心了。时间和阳光会帮助它孵卵的，而且，它也没有精力再
操心了。在替它的孩子做巢的时候，它已经把所有的丝都用光了，再也没有丝给自己张
网捕食了。况且它自己也已经没有食欲了。衰老和疲惫使它在世界上苟延残喘了几天后
安详地死去了。这便是我那匣子里的蜘蛛一生的终结，也是所有树丛里的蜘蛛的必然归
宿。

条纹蜘蛛的家族
你还记得那小小的巢里的橘黄色的卵吧？那些美丽的卵的总数有五颗之多。你还记
得它们是被密封在白缎子做的巢里的吧？那么当里面的小东西要跑出来，又冲不破白丝
做的墙的时候该怎么办呢？当时它们的母亲又不在身旁，不能帮助它们冲破丝袋，它们
是用什么办法来解决这个问题的呢？
动物在许多地方和植物有类似之处。蜘蛛的巢在我看来相当于植物的果实，只不过
它里面包含的不是种子而是卵而已，自己不能动，但它们的种子却能在很远的地方生根
发芽。因为植物有许多传播种子的方法，把它们送到四面八方：凤仙花的果实成熟的时
候，只要受到轻轻的碰触，便会裂成五瓣，每一瓣各自蜷缩起来，把种子弹到很远的地
方；还有一种很轻的种子，像蒲公英的种子，长着羽毛，风一吹就能把它们带到很远的
地方；榆树的种子是嵌在一张又宽又轻的扇子里的；槭树的种子成对地搭配，好像一双
张开的翅；柽树的种子形状像船桨，风能够让它飞到极远的地方……这些种子随遇而安，
落到不知什么地方就可以安家落户，开始下一轮的生死循环。
和植物一样，动物也有各种千奇百怪的方法凭借大自然的力量让它们的种族散布在
各地。你可以从条纹蜘蛛的身上略知一二。
三月间，正是蜘蛛的卵开始孵化的时候。如果我们用剪刀把蜘蛛的巢剪开，就可以
看到有些卵已经变成小蜘蛛爬到中央那个袋子的外面，有些仍旧是橘黄色的卵。这些刚
刚拥有生命、乳臭末干的小蜘蛛还没有披上像它们的母亲身上那样美丽的条纹衣服，它
们的背部是淡黄色的，腹部是棕色的，它们要在袋子的外面，巢的里面，呆上整整四个
月。在这段时期里，它们的身体渐渐地变得强壮丰满起来，和其它动物不同的是，它们
是在巢里而不是在外面的大天地中逐渐变为成年的蜘蛛的。
到了六七月里，这些小蜘蛛急于要冲出来了。可是它们无法在那坚硬的巢壁上挖洞。
那孩子怎么办呢？不用担心，那巢自己会裂开的，就像成熟种子的果皮一样，自动地把
后代送出来。它们一出巢，就各自爬到附近的树枝上，同时放出极为轻巧的丝来，这些
丝在空中飘浮的时候，会把它们牵引到别的地方去。关于小蜘蛛这种奇特的飞行工具，
我们在下一篇还会详细讲到。
------------------








狼蛛



蜘蛛有一个很坏的名声：大多数人都认为它是一种可怕的动物，一看到它就想把它
一脚踩死，这可能和蜘蛛狰狞的外表有关。不过一个仔细的观察家会知道，它是一个十
分勤奋的劳动者，是一个天才的纺织家，也是一个狡猾的猎人，并且在其它方面也很有
意思。所以，即使不从科学的角度看，蜘蛛也是一种值得研究的动物。但大家都说它有
毒，这便是它最大的罪名，也是大家都惧怕它的原因。不错，它的确有两颗毒牙，可以
立刻致它的猎物于死地。如果仅从这一点出发，我们的确可以说它是可怕的动物，可是
毒死一只小虫子和谋害一个人是两件迥然不同的事情。不管蜘蛛能怎样迅速地结束一只
小虫子的生命，对于人类来说，都不会有比蚊子的一刺更可怕的后果了。所以，我可以
大胆的说，大部分的蜘蛛都是无辜的，它们莫名其妙地被冤枉了。
不过，有少数种类的蜘蛛的确是有毒的。据意大利人说，狼蛛的一刺能使人痉挛而
疯狂地跳舞。要治疗这种病，除了音乐之外，再也没有别的灵丹妙药了。并且只有固定
的几首曲子治疗这种病特别灵验。这种传说听起来有点可笑，但仔细一想也有一定道理。
狼蛛的刺或许能刺激神经而使被刺的人失去常态，只有音乐能使他们镇定而恢复常态，
而剧烈地跳舞能使被刺中的人出汗，因而把毒驱赶出来。
在我们这一带，有最厉害的黑肚狼蛛，从它们身上可以得知蜘蛛的毒性有多大。我
家里养了几只狼蛛，让我把它介绍给你，并告诉你它是怎样捕食的吧！
这种狼蛛的腹部长着黑色的绒毛和褐色的条纹，腿部有一圈圈灰色和白色的斑纹。
它最喜欢住在长着百里香的干燥沙地上。我那块荒地，刚好符合这个要求，这种蜘蛛的
穴大约有二十个以上。我每次经过洞边，向里面张望的时候，总可以看到四只大眼睛。
这位隐士的四个望远镜像金钢钻一般闪着光，在地底下的四只小眼睛就不容易看到了。
狼蛛的居所大约有一尺深，一寸宽，是它们用自己的毒牙挖成的，刚刚挖的时候是
笔直的，以后才渐渐地打弯。洞的边缘有一堵矮墙，是用稻草和各种废料的碎片甚至是
一些小石头筑成的，看上去有些简陋，不仔细看还看不出来。有时候这种围墙有一寸高，
有时候却仅仅是地面上隆起的一道边。
我打算捉一只狼蛛。于是我在洞口舞动一根小穗，模仿蜜蜂的嗡嗡声。我想狼蛛听
到这声音会以为是猎物自投罗网，马上会冲出来。可是我的计划失败了。那狼蛛倒的确
往上爬了一些，想试探这到底是什么东西发出的声音，但它立刻嗅出这不是猎物而是一
个陷阱，于是一动不动地停在半途，坚决不肯出来，只是充满戒心地望着洞外。
看来要捉到这只狡猾的狼蛛，唯一的办法就是用活的蜜蜂作诱饵。于是我找了一只
瓶子，瓶子的口和洞口一样大。我把一只土蜂装在瓶子里，然后把瓶口罩在洞口上。这
强大的土蜂起先只是嗡嗡直叫，歇斯底里地撞击着这玻璃囚室，拼命想冲出这可恶的地
方。当它发现有一个洞口和自己的洞口很像的时候，便毫不犹豫地飞进去了。它实在是
愚蠢得很，走了那么一条自取灭亡的路。当它飞下去的时候，那狼蛛也正在匆匆忙忙往
上赶，于是它们在洞的拐弯处相撞了。不久我就听到了里面传来一阵死亡时的惨叫——
那只可怜的土蜂！这以后便是一段很长的沉默。我把瓶子移开，用一把钳子到洞里去探
索。我把那土蜂拖出来，它已经死了，正像刚才我所想象的那样。一幕悲剧早已在洞里
发生了。这狼蛛突然被夺走了从天而降的猎物，愣了一下，实在舍不得放弃这肥美的猎
物，急急地跟上来，于是猎物和打猎的都出洞了，我赶紧趁机用石子把洞口塞住。这狼
蛛被突如其来的变化惊呆了，一下子变得很胆怯，在那里犹豫着，不知该怎么办才好，
根本没有勇气逃走。不到一秒钟功夫，我便毫不费力地用一根草把它拔进一个纸袋里。
我就用这样的办法诱它出洞，然后捉拿归案。不久我的实验室里就有了一群狼蛛。
我用土蜂去引诱它，不仅仅是为了捉它，而且还想看看它怎样猎食。我知道它是那
种每天都要吃新鲜食物的昆虫，而不是像甲虫那样吃母亲为自己储藏的食物，或者像黄
蜂那样有奇特的麻醉术可以将猎物的新鲜程度保持到两星期之后。它是一个凶残的屠夫，
一捉到食物就将其活活地杀死，当场吃掉。
狼蛛得到它的猎物确实也不容易，也须冒很大的风险。那有着强有力的牙齿的蚱蜢
和带着毒刺的蜂随时都可能飞进它的洞去。说到武器，这两方不相上下。究竟谁更胜一
筹呢？狼蛛除了它的毒牙外没有别的武器，它不能像条纹蜘蛛那样放出丝来捆住敌人。
它唯一的办法就是扑在敌人身上，立刻把它杀死。它必须把毒牙刺入敌人最致命的地方。
虽然它的毒牙很厉害，可我不相信它在任何地方轻轻一刺而不是刺中要害就能取了敌人
的性命。

与木匠蜂作战
我已经讲过狼蛛生擒土蜂的故事，可这还不能使我满足，我还想看看它与别种昆虫
作战的情形。于是我替它挑了一种最强大的敌手，那就是木匠蜂。这种蜂周身长着黑绒
毛，翅膀上嵌着紫线，差不多有一寸长。它的刺很厉害，被它刺了以后很痛，而且会肿
起一块，很久以后才消失。我之所以知道这些，是因为曾经身受其害，被它刺过。它的
确是值得狼蛛去决一胜负的劲敌。
我捉了几只木匠蜂，把它们分别装在瓶子里。又挑了一只又大又凶猛并且饿得正慌
的狼蛛，我把瓶口罩在那只穷凶极恶的狼蛛的洞口上，那木匠蜂在玻璃囚室里发出激烈
的嗡嗡声，好像知道死期临头似的。狼蛛被惊动了从洞里爬了出来，半个身子探出洞外，
它看着眼前的景象，不敢贸然行动，只是静静地等候着。我也耐心地等候着。一刻钟过
去了，半个小时过去了，什么事都没有发生，狼蛛居然又若无其事地回到洞里去了。大
概它觉得不对头，冒然去捕食的话太危险了。我照这个样子又试探了其它几只狼蛛，我
不信每一只狼蛛都会这样面对丰盛的美食而无动于衷，于是继续一个一个的试探着，都
是这个样子，总对“天上掉下的猎物”怀有戒心。
最后，我终于成功了。有一只狼蛛猛烈地从洞里冲出来，无疑，它一定饿疯了。就
在一眨眼间，恶斗结束了，强壮的木匠蜂已经死了。凶手把毒牙刺到它身体的哪个部位
呢？是在它的头部后面。狼蛛的毒牙还咬在那里，我怀疑它真具有这种知识：它能不偏
不倚正好咬在唯一能致其于死的地方，也就是它的俘虏的神经中枢。
我做了好几次试验，发现狼蛛总是能在转眼之间干净利落地把敌人干掉，并且作战
手段都很相似。现在我明白了为什么在前几次试验中，狼蛛会只看着洞口的猎物，却迟
迟不敢出击。它的犹豫是有道理的。像这样强大的昆虫，它不能冒失鲁莽地去捉，万一
它没有击中其要害的话，那它自己就完蛋了。因为如果蜂没有被击中要害的话，至少还
可活上几个小时，在这几个小时里，它有充分的时间来回击敌人。狼蛛很知道这一点，
所以它要守在安全的洞里，等待机会，直到等到那大蜂正面对着它，头部极易被击中的
时候，它才立刻冲出去，否则决不用自己的生命去冒险。

狼蛛的毒素
让我来告诉你，狼蛛的毒素是一种多么厉害的暗器。
我做了一次试验，让一只狼蛛去咬一只羽毛刚长好的将要出巢的幼小的麻雀。麻雀
受伤了，一滴血流了出来，伤口被一个红圈圈着，一会儿又变成了紫色，而且这条腿已
经不能用了，使不上劲。小麻雀只能用单腿跳着。除此之外它好像也没什么痛苦，胃口
也很好。我的女儿同情地把苍蝇、面包和杏酱喂给它吃，这可怜的小麻雀作了我的实验
品。但我相信它不久以后一定会痊愈，很快就能恢复自由——这也是我们一家共同的愿
望和推测。十二个小时后，我们对它的伤情仍然挺乐观的。它仍然好好地吃东西，喂得
迟了它还要发脾气。可是两天以后，它不再吃东西了，羽毛零乱，身体缩成一个小球，
有时候一动不动，有时候发出一阵痉挛。我的女儿怜爱地把它捧在手里，呵着气使它温
暖。可是它痉挛得越来越厉害，次数越来越多，最后，它终于离开了这个世界。
那天的晚餐席上透着一股寒气。我从一家人的目光中看出他们对我的这种试验的无
声的抗议和责备。我知道他们一定认为我太残忍了。大家都为这只不幸的小麻雀的死而
悲伤。我自己也很懊悔：我所要知道的只是很小的一个问题，却付出了那么大的代价。
尽管如此，我还是鼓起勇气试验一只鼹鼠，它是在偷田里的莴苣时被我们捉住的，
所以即使它死于非命也不足为惜。我把它关在笼子里，用各种甲虫、蚱蜢喂它，它大口
大口贪婪地吃着，被我养得胖胖的，健康极了。
我让一只狼蛛去咬它的鼻尖。被咬过之后，它不住地用它的宽爪子挠抓着鼻子。因
为它的鼻子开始慢慢地腐烂了。从这时开始，这只大鼹鼠食欲渐渐不振，什么也不想吃，
行动迟钝，我能看出它浑身难受。到第二个晚上，它已经完全不吃东西了。大约在被咬
后三十六小时，它终于死了。笼里还剩着许多的昆虫没有被吃掉，证明它不是被饿死的，
而是被毒死的。
所以狼蛛的毒牙不止能结束昆虫的性命，对一些稍大一点的小动物来说，也是危险
可怕的。它可以致麻雀于死地，也可以使鼹鼠毙命，尽管后者的体积要比它大得多。虽
然后来我再没有做过类似的试验，但我可以说，我们千万要小心戒备，不要被它咬到，
这实在不是一件可以拿来试验的事。
现在，我们试着把这种杀死昆虫的蜘蛛和麻醉昆虫的黄蜂比较一下。蜘蛛，因为它
自己靠新鲜的猎物生活，所以它咬昆虫头部的神经中枢，使它立刻死去；而黄蜂，它要
保持食物的新鲜，为它的幼虫提供食物，因此它刺在猎物的另一个神经中枢上，使它失
去了动弹的能力。相同的是，它们都喜欢吃新鲜的食物，用的武器都是毒刺。
没有谁教它们怎样根据自己的需要分别用不同的方法对待猎物，它们生来就明白这
一点。这使我们相信冥冥之中，世界上的确有着一位万能的神在主宰着昆虫，也统治着
人类世界。

狼蛛猎食
我在实验室的泥盆里，养了好几只狼蛛。从它们那里，我看到狼蛛猎食时的详细情
形。这些做了我的俘虏的狼蛛的确很健壮。它们的身体藏在洞里，脑袋探出洞口，玻璃
般的眼睛向四周张望，腿缩在一起，作着准备跳跃的姿势，它就这样在阳光下静静地守
候着，一两个小时，不知不觉就过去了。
如果它看到一只可作猎物的昆虫在旁边经过，它就会像箭一般地跳出来，狠狠地用
它的毒牙打在猎物的头部，然后露出满意又快乐的神情，那些倒霉的蝗虫、蜻蜓和其它
许多昆虫还没有明白过来是怎么回事，就做了它的盘中美餐。它拖着猎物很快地回到洞
里，也许它觉得在自己家里用餐比较舒适吧。它的技巧以及敏捷的身手令人叹为观止。
如果猎物离它不太远，它纵身一跃就可以扑到，很少有失手的时候。但如果猎物在
很远的地方，它就会放弃，决不会特意跑出来穷追不舍。看来它不是一个贪得无厌的家
伙，不会落得一个“鸟为食亡”的下场。
从这一点可以看出狼蛛是很有耐性，也很有理性的。因为在洞里没有任何帮助它猎
食的设备，它必须始终傻傻地守候着。如果是没有恒心和耐心的虫子，一定不会这样持
之以恒，肯定没多久就退回到洞里去睡大觉了。可狼蛛不是这种昆虫。它确信，猎物今
天不来，明天一定会来；明天不来，将来也总有一天会来。在这块土地上，蝗虫、蜻蜓
之类多得很，并且它们又总是那么不谨慎，总有机会刚好跳到狼蛛近旁。所以狼蛛只需
等待时候一到，它就立刻窜上去捉住猎物，将其杀死。或是当场吃掉，或者拖回去以后
吃。
虽然狼蛛很多时候都是“等而无获”，但它的确不大会受到饥饿的威胁，因为它有
一个能节制的胃。它可以在很长一段时间内不吃东西而不感到饥饿。比如我那实验室里
的狼蛛，有时候我会连续一个星期忘了喂食，但它们看上去照样气色很好。在饿了很长
一段时间后，它们并不见得憔悴，只是变得极其贪婪，就像狠一样。
在狼蛛还年幼的时候，它还没有一个藏身的洞，不能躲在洞里“守洞待虫”，不过
它有另外一种觅食的方法。那时它也有一个灰色的身体，像别的大狼蛛一样，就是没有
黑绒腰裙——那个要到结婚年龄时才能拥有。它在草丛里徘徊着，这是真正的打猎。当
小狼蛛看到一种它想吃的猎物，就冲过去蛮横地把它赶出巢，然后紧迫不舍，那亡命者
正预备起飞逃走，可是往往来不及了——小狼蛛已经扑上去把它逮住了。
我喜欢欣赏我那实验室里的小狼蛛捕捉苍蝇时那种敏捷的动作。苍蝇虽然常常歇在
两寸高的草上，可是只要狼蛛猛然一跃，就能把它捉住。猫捉老鼠都没有那么敏捷。
但是这只是狼蛛小时候的故事，因为它们身体比较轻巧，行动不受任何限制，可以
随心所欲。以后它们要带着卵跑，不能任意地东跳西窜了。所以它就先替自己挖个洞，
整天在洞口守候着，这便是成年蜘蛛的猎食方式。

狼蛛的卵袋
假如你听到这可怕的狼蛛怎样爱护自己的家庭的故事，你一定会在惊异之余改变对
它的看法。
在八月的一个清晨，我发现一只狼蛛在地上织一个丝网，大小和一个手掌差不多。
这个网很粗糙，样子也不美观，但是很坚固。这就是它将要工作的场所，这网能使它的
巢和沙地隔绝。在这网上，它用最好的白丝织成一片大约有一个硬币大小的席子，它把
席子的边缘加厚，直到这席子变成一个碗的形状，周围圈着一条又宽又平的边，它在这
网里产了卵，再用丝把它们盖好，这样我们从外面看，只看到一个圆球放在一条丝毯上。
然后它就用腿把那些攀在圆席上的丝一根根抽去，然后把圆席卷上来，盖在球上，
然后它再用牙齿拉，用扫帚般的腿扫，直到它把藏卵的袋从丝网上拉下来为止，这可是
一项费神费力的工作。
这袋子是个白色的丝球，摸上去又软又粘，大小像一颗樱桃。如果你仔细观察，那
么你会发现在袋的中央有一圈水平的折痕，那里面可以插一根针而不致于把袋子刺破。
这条折纹就是那圆席的边。圆席包住了袋子的下半部，上半部是小狼蛛出来的地方。除
了母蜘蛛在产好卵后铺的丝以外，再也没有别的遮蔽物了。袋子里除了卵以外，也没有
别的东西，不像条纹蜘蛛那样，里面衬着柔软的垫褥和绒毛。狼蛛不必担心气候对卵的
影响，因为在冬天来临之前，狼蛛的卵早已孵化了。
母蛛整个早晨都在忙着编织袋子。现在它累了。它紧紧抱着它那宝贝小球，静静地
休息着，生伯一不留神就把宝贝丢了。第二天早晨，我再看到它的时候，它已经把这小
球挂到它身后的丝囊上了。
差不多有三个多星期，它总是拖着那沉重的袋子。不管是爬到洞口的矮墙上的时候，
还是在遭到了危险急急退入地洞的时候，或者是在地面上散步的时候，它从来不肯放下
它的宝贝的小袋。如果有什么意外的事情使这个小袋子脱离它的怀抱，它会立刻疯狂地
扑上去，紧紧地抱住它，并准备好反击那抢它宝贝的敌人。接着它很快地把小球挂到丝
囊上，很不安地带着它匆匆离开这个是非之地。
在夏天将要结束的那几天里，每天早晨，太阳已经把土地烤得很热的时候，狼蛛就
要带着它的小球从洞底爬出洞口，静静地趴着。初夏的时候，它们也常常在太阳高挂的
时候爬到洞口，沐浴着阳光小睡。不过现在，它们这么做完全是为了另外一个目的。以
前狼蛛爬到洞口的阳光里是为了自己，它躺在矮墙上，前半身伸出洞外，后半身藏在洞
里。它让太阳光照到眼睛上，而身体仍在黑暗中；现在它带着小球，晒太阳的姿势刚好
相反：前半身在洞里，后半身在洞外。它用后腿把装着卵的白球举到洞口，轻轻地转动
着它，让每一部分都能受到阳光的沐浴。这样足足晒了半天，直到太阳落山。它的耐心
实在令人感动，而且它不是一天两天这样做，而是在三四个星期内天天这样做。鸟类把
胸伏在卵上，它的胸能像火炉一样供给卵充分的热量；狼蛛把它的卵放在太阳底下，直
接利用这个天然的大火炉。

狼蛛的幼儿
在九月初的时候，小狼蛛要准备出巢了。这时小球会沿着折痕裂开。它是怎么裂开
的呢？会不会是母蛛觉察到里面有动静，所以在一个适当的时候把它打开了？这也是可
能的。但从另一方面看，也可能是那小球到了一定时间自己裂开的，就像条纹蜘蛛的袋
子一样。条纹蜘蛛出巢的时候，它们的母亲早已过世多时了。所以只有靠巢的自动裂开，
孩子们才能出来。
这些小狼蛛出来以后，就爬到母亲的背上，紧紧地挤着，大约有二百只之多，像一
块树皮似的包在母蛛身上。至于那袋，在孵化工作完毕的时候就从丝囊上脱落下来，被
抛在一边当垃圾了。
这些小狼蛛都很乖，它们不乱动，也不会为了自己挤上去而把别人推开。它们只是
静静地歇着。它们在干什么呢？它们是让母亲背着它们到处去逛。而它们的母亲，不管
是在洞底沉思，还是爬出洞外去晒太阳，总是背着一大堆孩子一起跑，它从不会把这件
沉重的外衣甩掉，直到好季节的降临。
这些小狼蛛在母亲背上吃些什么呢？照我看来，它们什么也没吃。我看不出它们长
大，它们最后离开母亲的时候，和它们刚从卵里出来的时候大小完全一样。
在坏的季节里，狼蛛母亲自己也吃得很少。如果我捉一只蝗虫去喂它，常会过了很
久以后它才开口。为了保持元气，它有时候不得不出来觅食，当然，它还是背着它的孩
子。
如果在三月里，当我去观察那些被风雨或霜雪侵蚀过的狼蛛的洞穴的时候，总可以
发现母蛛在洞里，仍是充满活力的样子，背上还是背满了小狼蛛。也就是说，母蛛背着
小蛛们活动，至少要经过五六个月。著名的美洲背负专家——鼹鼠，它也不过把孩子们
背上几个星期就把它们送走了，和狼蛛比起来，真是小巫见大巫了。
背着小狼蛛出征是很危险的，这些小东西常常会被路上的草拨到地上。如果有一支
小狼蛛跌落到地上，它将会遭遇什么命运呢？它的母亲会不会想到它，帮它爬上来呢？
绝对不会。一只母蛛需要照顾几百只小蛛，每只小蛛只能分得极少的一点爱。所以不管
是一只、几只或是全部小狼蛛从它背上摔下来，它也决不为它们费心。它不会让孩子们
依靠别人的帮助解决难题，它只是静静地等着，等它们自己去解决困难，事实上这困难
并不是不能解决，而且往往解决得很迅速、很利落。
我用一只笔把我实验室中的一个母狼蛛背上的小蛛刮下，母亲一点儿也不显得惊慌，
也不准备帮助它的孩子，继续若无其事地往前走。那些落地的小东西在沙地上爬了一会
儿，不久就都攀住了它们的母亲身体的一部分：有的在这里攀住了一只脚，有的在那里
攀住一只脚。好在它们的母亲有不少脚，而且撑得很开，在地面上摆出一个圆，小蛛们
就沿着这些柱子往上爬，不一会儿，这群小蛛又像原来那样聚在母亲背上了。没有一只
会漏掉。在这样的情况下，小狼蛛很会自己照顾自己，母亲从不需为它们的跌下而费心。
在母蛛背着小蛛的七个月里，它究竟喂不喂它们吃东西呢？当它猎取了食物后，是
不是邀孩子们共同享受呢？起初我以为一定是这样的，所以我特别留心母蛛吃东西时的
情形，想看看它怎样把食物分给那么多的孩子们。通常母蛛总是在洞里吃东西，不过有
时候偶然也到门口就着新鲜空气用餐。只有在这时候我才有机会，看到这样的情形：当
母亲吃东西的时候，小蛛们并不下来吃，连一点要爬下来分享美餐的意思都没有。好像
丝毫不觉得食物诱人一样，它们的母亲也不客气，没给它们留下任何食物。母亲在那儿
吃着，孩子们在那儿张望着——不，确切的说，它们仍然伏在妈妈的背上，似乎根本不
知道“吃东西”是怎样一种概念。在它们的母亲狼吞虎咽的时候，它们安安静静地呆在
那儿，一点儿也不觉得馋。
那么，在爬在母亲背上整整七个月的时间里，它们靠什么吸取能量、维持生命呢？
你或许会猜想它们不会是从母亲的皮肤上吸取养料的吧？我发现并不是这样的。因为据
我观察，它们从来没有把嘴巴贴在母亲的身上吮吸。而那母蛛，也并不见得瘦削和衰老，
它还是和以往一样神采奕奕，而且比以前更胖了。
那么又要问了，它们这些小蛛靠什么维持生命呢？一定不是以前在卵里吸收的养料。
以前那些微不足道的养料。别说是不能帮它们造出丝来，连维持它们的小生命都很困难。
在小蛛的身体里一定有着另外一种能量。
如果它们不动，我们很容易理解为什么它们不需要食物，因为完全的静止就相当于
没有生命。但是这些小蛛，虽然它们常常安静地歇在母亲背上，但它们时刻都在准备运
动。当它从母亲这个“婴儿车”上跌落下来的时候，它们得立刻爬起来抓住母亲的一条
腿，爬回原处；即使停在原地，它也得保持平衡；它还必须伸直小肢去搭在别的小蛛身
上，才能稳稳地趴在母亲背上。所以，实际上绝对的静止是不可能的。
从生理学角度看，我们知道每一块肌肉的运动都需要消耗能量。动物和机器一样，
用得久了会造成磨损，因此需常常修理更新。运动所消耗的能量，必须从别的地方得到
补偿。我们可以把动物的身体和火车头相比。当火车头不停地工作的时候，它的活塞、
杠杆、车轮以及蒸汽导管都在不断地磨损，铁匠和机械师随时都在修理和添加些新材料，
就好像供给它食物，让它产生新的力量一样。但是即使机器各部分都很完美，火车头还
是不能开动。一直要等到火炉里有了煤，燃起了火，然后才能开动。这煤就是产生能量
的“食物”，就是它让机器动起来的。
动物也是这样。有能量才能运动。小动物在胚胎时期，从母亲的胎盘里或者卵里吸
取养料，那是一种制造纤维素的养料，它使小动物的身体长大长坚固，并且补偿一些不
足的地方。但是，除此之外，必须有产生热量的食物，才能使小动物跑、跳、游泳、飞
跃，或是作其它各种运动。任何运动都少不了能量。
再讲这些小狼蛛，它们在离开母亲的背之前，并不曾长大。七个月的小蛛和刚刚出
生的小蛛完全一样大。卵供给了足够的养料，为它们的体质打下了一个良好的基础。但
它们后来不再长大，因此也不再需要吸收制造纤维的养料。这一点我们是能够理解的。
但它们是在运动的呀！并且运动得很敏捷。它们从哪里取得产生能量的食物呢？
我们可以这样想：煤——那供给火车头动能的食物究竟是什么呢？那是许多许多年
代以前的树埋在地下，它们的叶子吸收了阳光。所以煤其实就是贮存起来的阳光，火车
头吸收了煤燃烧提供的能量，也就是相当于吸收了太阳光的能量。
血肉之躯的动物也是这样，不管它是吃什么别的动物或植物以维持生命，大家最终
都是靠着太阳的能量生存的。那种热能量贮藏在草里、果子里、种子里和一切可作为食
物的东西里。太阳是宇宙的灵魂，是能量的最高赐予者，没有太阳，就没有地球上的生
命。
那么除了吃进食物，然后经过胃的消化作用变成能量以外，太阳光能不能直接射入
动物的身体，产生活力，就像蓄电池充电那样？为什么不能直接靠阳光生存呢？我们吃
的果子中除了阳光外，还有别的什么物质吗？
化学家告诉我们，将来我们可以靠一种人工的食物来维持生命。那时候所有的田庄
将被工厂和实验室所取代，化学家们的工作就是配置产生纤维的食物和产生能的食物，
物理学家们也靠着一些精巧仪器的帮助，每天把太阳能注进我们的身体，供给我们运动
所需的能量。那样我们就能不吃东西而维持生命。不吃饭而是吃太阳的光线，你能想象
吗？那将是一个多么美妙而有趣的世界！
这是我们的梦想，它能实现吗？这个问题倒是很值得科学家们研究的。

小蛛的飞逸
到三月底的时候，母蛛就常常蹲在洞口的矮墙上。这是小蛛们与母亲告别的时候了。
作母亲的仿佛早已料到这么一天，完全任凭它们自由地离去。对于小蛛们以后的命运，
它再也不需要负责了。
在一个天气很好的日子里，它们决定在那天最热的一段时间里分离。小蛛们三五成
群地爬下母亲的身体。看上去丝毫没有依依惜别之情，它们在地上爬了一会儿后，便用
惊人的速度爬到我的实验室里的架子上。它们的母亲喜欢住在地下，它们却喜欢往高处
爬。架子上恰好有一个竖起来的环，它们就顺着环很快地爬了上去。就在这上面，它们
快活地纺着丝，搓着疏松的绳子。它们的腿不住地往空中伸展，我知道这是什么意思：
它们还想往上爬，孩子长大了，一心想走四方闯天下，离家越远越好。
于是我又在环上插了一根树枝。它们立刻又爬了上去，一直爬到树枝的梢上。在那
里，它们又放出丝来，攀在周围的东西上，搭成吊桥。它们就在吊桥上来来去去，忙碌
地奔波着，看它们的样子似乎还不满足，还想一个劲儿往上爬。
我又在架子上插了一根几尺高的芦梗，顶端还伸展着细枝。那些小蛛立刻又迫不及
待地爬了上去，一直到达细枝的梢上。在那儿，它们又乐此不疲地放出丝、搭成吊桥。
不过这次的丝很长很细，几乎是飘浮在空中的，轻轻吹口气就能把它吹得剧烈地抖动起
来，所以那些小蛛在微风中好像在空中跳舞一般。这种丝我们平时很难看见，除非刚好
有阳光照在丝上，才能隐隐约约看到它。
忽然一阵微风把丝吹断了。断了的一头在空中飘扬着。再看这些小蛛，它们吊在丝
上荡来荡去，等着风停；如果风大的话，可能把它们吹到很远的地方，使它们重新登陆，
到一个陌生的地方。
这种情形又要维持好多天。如果在阴天，它们会保持静止，动都不想动，因为没有
阳光供给能量，它们不能随心所欲地活动。
最后，这个庞大的大家庭消失了。这些小蛛纷纷被飘浮的丝带到各个地方。原来背
着一群孩子的荣耀的母蛛变成了孤老。一下子失去那么多孩子，它看来似乎并不悲痛。
它更加精神焕发地到处觅食，因为这时候它背上再也没有厚厚的负担了，轻松了不少，
反而显得年轻了。不久以后它就要做祖母，以后还要做曾祖母，因为一只狼蛛可以活上
好几年呢。
从这一家狼蛛中，我们可以看到，有一种本能，很快地赋予小蛛，不久又很快地而
且是永远地消失。那就是攀高的本能。它们的母亲不知道自己的孩子曾有这样的本事，
孩子们自己不久以后也会彻底地忘记。它们到了陆地上，做了许多天流浪儿之后，便要
开始挖洞了。这时候，它们中间谁也不会梦想爬上一颗草梗的顶端。可那刚刚离开母蛛
的小蛛的确是那样迅速、那样容易地爬到高处，在它生命的转折之处，它曾是一个满怀
激情的攀登大师。我们现在知道了它这样做的目的：在很高的地方，它可以攀一根长长
的丝。那根长丝在空中飘荡着，风一吹，就能使它们飘荡到远方去。我们人类有飞机，
它们也有它们的飞行工具。在需要的时候，它替自己制造这种工具，等到旅行结束，它
也就把它忘记了。
------------------








克鲁蜀蜘蛛



克鲁蜀蜘蛛是一个极为聪明、灵巧的纺织家，而且就一只蜘蛛而言，克鲁蜀蜘蛛算
是很漂亮的了。它这名字是取自古希腊三位命运女神中的一位，也是最年幼的一位，她
是掌管纺线杆的，从她那里纺出了万物各自不同的命运。克鲁蜀蜘蛛能为自己纺出最精
美的丝，克鲁蜀女神却不能为我们纺出幸福的命运和舒适的生活，这实在是一件令世人
遗憾的事！
如果我们想认识克鲁蜀蜘蛛，我们必须到橄榄地的岩石的斜坡上。在被太阳灼得又
热又亮的地方，让我们把一些不大不小又扁平的石块翻起来——最好是翻开那些小石堆，
那是牧童堆起来做凳子用的，这种凳子尽管简陋，但深受牧童们的喜爱，因为他们可以
坐在上面看守山底下的羊群，边休息边工作，不亦乐乎。而这些地方往往是克鲁蜀蜘蛛
最爱呆的地方，如果你开始找不着它们，不要灰心，克鲁蜀蜘蛛在这个世界上的确是很
稀少的，而且并不是每个地方都适合它们的。所以找不着它们一点儿也不奇怪，如果我
们运气不坏的话，在我们翻起的石块下面，就会发现一个样子很特别的东西：形状好像
一个翻转的穹形屋顶，大概有半个梅子那么大，外面挂着一些小贝壳，一些泥土和干了
的虫子。
穹形顶的边缘有十二个尖尖的扇蛤，向各方伸展着，固定在石头上。显然，这就是
克鲁蜀蜘蛛的宫殿！
而入口处在什么地方呢？尽管周围有许多拱，但那些拱都开在屋顶的上部，没有一
个可以通到屋子里去。可这屋子的主人总得出来，出来后当然还得再回去呀，它究竟是
从哪里进出的呢？一根稻草会告诉我们一切秘密。
如果我们用一根稻草往拱形的开口处插进去，我们可以发现这些拱门里面都是反锁
着的，关得严严实实的。但是如果你把稻草很小心地用力插进去，你就会发现，其中必
定有一个拱门，它的边缘会裂成嘴唇般的两片。这便是门了，它有弹性，自己会关闭。
当克鲁蜀蜘蛛遇到危险的时候，它会飞快地跑回自己的家，用脚爪一触门，门就开
了，等它进去后，门又自动关闭了。如果需要的话，它还可以将门反锁，尽管那所谓的
“锁”也只不过是用几根丝线做的，起不了多大作用。但从外边看它跟别的拱完全一样，
可以起到迷惑的作用，敌人只能看见它很快地消失了，却不知道它究竟是从哪里逃遁的。
现在让我们打开它的门到里面看一看。啊！那么的富丽堂皇！我马上想起一个关于
娇贵的公主的神话故事，说她是如此的娇贵，只要她的垫褥底上有一片折皱的玫瑰花的
叶子，她就会睡得很不舒服。看到克鲁蜀蜘蛛的家，你会觉得它比那位公主还要难侍候。
它的床比天鹅绒还软，比夏天的云还白。床上有绒毯，有被子，都非常软。克鲁蜀蜘蛛
就安居在这绒毯和被子之间。它长着一双短短的腿，穿着黑色的衣服，背上还有五个黄
色的徽章。
要在这屋子里过上舒适的生活，必须有一个条件：那就是屋子必须建筑得很坚固，
尤其是在遭到大风大雨的时候。只要我们仔细观察的话，就可以发现克鲁蜀蜘蛛是如何
做到这一点的：支持着整个屋子的许多拱门都是固定在石头上的。我们可以看到在接触
点有一缕缕的长线，沿着石面伸展开去。我用尺子量了量，每一根足足有九尺长。原来
这屋子由许多“链条”攀着，就像阿拉伯人的帐幕用许多绳子攀着一样。所以显得格外
牢固。
还有一件小事挺值得我们注意，虽然它的屋子里那么整洁、华丽，可是屋子外面却
堆满了垃圾：有泥土、有腐烂了的木屑，还有一粒粒脏兮兮的砂石，有时候还有更脏的
东西，比如风干了的甲虫尸体、干足虫破碎的尸体，还有蜗牛的壳等等，一片狼藉，一
踏糊涂，而且都已被太阳晒得发白了。
克鲁蜀蜘蛛没有设陷阱的本领，所以它完全靠吃那种在石堆里跳来跳去的虫子维持
生计。哪一只冒失的虫子跳过它的居所，就会被它逮个正着，捉来饱餐一顿。至于那些
风干的尸体，它并不把它们丢开，而是挂在墙壁的四周，这似乎在炫耀自己的捕猎经历。
它到底为什么要把尸体挂起来呢？那些蜗牛壳大部分是空的，有时候里面也有活着
而且没有受伤的蜗牛，可是这些蜗牛躲在壳的深处，蜘蛛根本不能接触到它，把它挂起
来又有什么用呢？它不能把那么硬的壳打破，也无法从开口处把蜗牛揪出来，它居然还
收集了那么多蜗牛，它到底想干什么呢？
后来我们想到了简单的平衡问题。普通的家蛛在墙角张了网，为了要使网保持一定
的形状而不受风吹雨打的影响，就常常把石灰或泥嵌入网中。克鲁蜀蜘蛛在自己的屋里
挂上重物虽然也是为了使屋子稳固，只是不是这种原理，而是它懂得如果把重物挂在屋
子里很低的地方，屋子四周就会平衡：在低的地方增加重量就可以把重心降低，使平衡
更加稳定。这些重物主要是昆虫的尸体或空壳，因为这类东西举手可得，不需到远处去
找，全都有现成的。
现在，你们或许要问：它在那软软的屋子里做什么呢？醉生梦死还是辛勤劳动呢？
据我所知，它什么也不做。它的肚子已经装得饱饱的了。它舒适地摊开了脚，躺在软软
的绒毯上，什么都不用做，什么都不用想，只是静静地聆听着地球旋转的声音。这不是
睡，也不是醒，而是一种半梦半醒的状态，这时候它除了快乐以外，什么感觉都没有。
让我们想象一下：如果你辛苦了一天后很舒适地躺在床上，彻底地放松，到快要入睡的
时候，将会无忧无虑地享受到这生活中最美满的一刻。克鲁蜀蜘蛛似乎也知道这美妙的
一刻，并且比我们更会享受这一刻的到来。
------------------








迷宫蛛



会结网的蜘蛛称得上是个纺织能手，它们用蛛网来猎取自投罗网的小虫子们，可谓
“坐享其成，得来全不费功夫”。还有许多其它种类的蜘蛛，它们用许多别的聪明的方
法猎取食物，同样可以以逸待劳，大获丰收。其中有几种在这方面很有造诣，几乎所有
的有关昆虫的书都会把它们列举出来。
那是一种黑色的蜘蛛，也有人叫它美洲狼蛛，它们是住在洞里的，就像我以前讲到
的欧洲狼蛛一样。但是它们的洞穴比欧洲狼蛛的洞穴要完备精细得多。欧洲狼蛛的洞口
只有一圈矮墙，用小石子、丝和废料堆成的，而美洲狼蛛的洞口上有一扇活动门，是由
一块圆板、一个槽和一个栓子做成的。当一只狼蛛回家的时候，门便会落进槽里，自动
把门关了。如果有谁在门外想把它掀起来的话，狼蛛只要用两只爪把柱子抵住，门就紧
紧关闭住，不会受外面的影响。
另外一种是水蛛。它能替自己做一只性能很好的潜水袋，里面贮藏着空气。它在这
里面等待猎物经过，同时也可以说是在避暑。在太阳像大火炉一样的日子里，这地方的
确是一个舒适凉爽的避暑胜地。人类中也有人尝试用最硬的石块或大理石在水下造房子。
不知大家有没有听说过泰比利斯，他是罗马的一个暴君，他生前曾经叫人为他造了一座
水下宫殿，供自己寻欢作乐。不过到现在这个宫殿只给人们留下一点回忆和感慨，而狼
蛛的水晶宫，却是永远灿烂辉煌的。
如果我有机会观察一下这些水蛛的话，我一定能在它们的生命史上添上一些未经记
载的事实。但是现在我不得不放弃这个想法。因为我们这一带没有水蛛。至于那美洲狼
蛛，我也只有在路旁看到过一次。而那时候我恰巧有别的事情要去办理，没有时间去看
它。错失这个良机后，后来也就一直没见到过它。
但是，并不是稀罕的虫子才值得研究。普通的虫子，如果好好地研究起来，也能发
现许多有趣的事情。我对迷宫蛛的接触机会极多，对它也很感兴趣，所以对它作了一番
研究，我觉得是很有收获的。在七月的清晨，太阳还没有焦灼着人的头颈的时候，每星
期我总要去树林里看几次迷宫蛛。孩子们也都跟着我去，每人还带上一个橘子，以供解
渴之用。
走进树林，不久，我们就发现许多很高的丝质建筑物，丝线上还串着不少露珠，在
太阳光的照射下闪闪发光，好像皇宫里的稀世珍宝一般。孩子们被这个美丽的“灯架”
惊呆了，几乎忘记了他们的橘子。我们的蜘蛛的迷宫真算得上一个奇观！
经过太阳半小时的照射，魔幻般的珍珠随着露水一起消失了。现在可以来专心观察
它的网了。在那丛蔷蔽花的上方张着一张网，大概有一块手帕那么大，周围有许多线把
它攀到附近的矮树丛中，使它能够在空中固定住，中间这张网看起来犹如一层又轻又软
的纱。
网的四周是平的，渐渐向中央凹，到了最中间便变成一根管子，大约有八九寸深，
一直通到叶丛中。
蜘蛛就坐在管子的进口处。它对着我们坐着，一点儿也不惊慌。它的身体是灰色的，
胸部有两条很阔的黑带，腹部有两条细带，由白条和褐色的斑点相间排列而成。在它的
尾部，有一种“双尾”，这在普通蜘蛛中是很少见的。
我猜想在管子的底部，一定有一个垫得软软的小房间，作为迷宫蛛空闲时候的休息
室。可事实上那里并没有什么小房间，只有一个像门一样的东西，一直是开着的，它在
外面遇到危险的时候，可以直接逃回来。
上面那个网由于用许多丝线攀到附近的树枝上，所以看上去活像一艘暴风雨下抛锚
的船。这些充当铁索的丝线中，有长的，也有短的；有垂直的，也有倾斜的；有紧张的，
也有松弛的；有笔直的，也有弯曲的，都杂乱地交叉在三尺以上的高处。这确实可以算
是一个迷宫，除了最强大的虫子外，谁都无法打破它，逃脱它的束缚。
迷宫蛛不像别的蜘蛛那样可以用粘性的网作为陷阱，它的丝是没有粘性的，它的网
妙就妙在它的迷乱。你看那只小蝗虫，它刚刚在网上落脚，便由于网摇曳不定，根本设
法让自己站稳。一下子陷了下去，它开始焦躁地挣扎，可是越挣扎陷得越深，好像掉进
了可怕的深渊一样。蜘蛛呆在管底静静地张望着，看着那倒霉的小蝗虫垂死挣扎，它知
道，这个猎物马上会落到网的中央，成为它的盘中美餐。
果然，一切都在蜘蛛的意料之中。它不慌不忙地扑到猎物上，慢慢地一口一口地吮
吸着它的血，一副得意洋洋的样子。至于那蝗虫，在蜘蛛咬它第一口的时候就死了——
蜘蛛的毒液使它一命呜呼。接下来蜘蛛就要从容地来吃完它了，而对于这只蝗虫来说，
这远比半死不活或者活活被蜘蛛撕成碎片要舒服多了。
到快要产卵的时候，迷宫蛛就要搬家了。尽管它的网还是完好无损，但它必须忍痛
割爱。它不得不舍弃它，而且以后也不再回来了。它必须去完成它的使命，一心一意去
筑巢了。它把巢做在什么地方呢？迷宫蛛自己当然知道得很清楚，而我，却一点头绪都
没有，实在猜想不出它会把巢造在哪儿。我花了整个早晨在树林中各个地方搜索。工夫
不负有心人，我终于发现了它的秘密。
在离网相当远的一个树丛里，它造好了它的巢。那里堆着一堆枯柴，草率而杂乱地
缠在一起，显得有点脏。就在这个简陋的盖子下，有一个做得比较细致、精巧的丝囊，
里面就是迷宫蛛的卵。
看到它的巢那么简陋，我不禁有些失望。但是后来我想到了，这一定是因为环境不
够好。你想，在这样一个密密的树丛里，一堆枯枝枯叶中，哪有条件让它做精致的活呢
7为了要证明我的推想没有错，我带了六只快要产卵的迷宫蛛到我家里，放在实验室的一
个铁笼子里面，然后把铁笼子竖在一个盛沙的泥盘子里，又在泥盘中央插了一根百里香
的小树枝，使每一个巢有攀附的地方。一切准备就绪后，现在就让它们大显身手吧。
这个实验获得了极大的成功。七月底的时候，我得到了六只雪白的、外观富丽精致
的丝囊。迷宫蛛在这样一个舒适的环境里工作，活干得自然细致了许多。让我来尽情地
观察吧！这个巢是一个由白纱编织而成的卵形的囊，有一个鸡蛋那么大。内部的构造也
很迷乱，和它的网差不多——看来这种建筑风格在它的脑子里已经根深蒂固了，所以无
论在什么场所，在什么条件下，它造的建筑物都是那样的迷乱无章。
这个布满丝的迷宫还是一个守卫室。在这乳白色的半透明的丝墙里面还装着一个卵
囊，它的形状有点像那些代表某一等级的骑士的星形勋章。这是一个很大的灰白色的丝
袋，周围筑着圆柱子，使它能够固定在巢的中央。这种圆柱都是中间细，两头粗，总共
大约有十个，在卵室的周围构成一个白色的围廊。母蛛在这个围廊里徘徊着，一会儿在
这儿停住，一会儿又在那儿停住，时时聆听着卵囊里的动静，活像一个马上要做父亲的
人在产室外面焦急地等待着孩子的第一声啼哭。这样一个卵巢里面，大概藏着一百颗左
右淡黄色的卵。
轻轻移去外面的白丝墙后，可以看到里面，还有一层泥墙，那是丝线夹杂着小碎石
做成的。可是这些小沙子怎么到丝墙里面去的呢？是跟着雨水渗进去的吗？不对，因为
外面的丝墙上白得没有一丝斑点，更不用说什么水迹了，看来决不是从这墙上渗进去的。
到后来，我才发现这是母蛛自己搬进去的，它为了怕卵受到寄生虫的侵犯，所以特地把
砂粒掺在丝线里面做成一培坚固的墙。
这丝墙里面还有一个丝囊，那才是盛卵的囊。我打开的这个巢里面的卵已经孵化了，
所以我能看到许多弱小的小蜘蛛在囊里快乐地爬来爬去。
但是，再回过头来看看那母蛛，它为什么要舍弃那张还完好无损的网，而把巢筑到
那么远的地方呢？它的舍近求远自然有它的道理，你还记得它的网的样子吧？在它的网
的上方，有一个错综复杂的迷宫，高高地露在树叶丛的外面，这是一个巨大的陷阱，同
时也是一个很醒目的标志，它的敌人——寄生虫轻而易举地就能看到这个迷宫，然后循
着它再找到迷宫蛛的巢——如果这巢离那醒目的网很近的话，那么寄生虫会不费吹灰之
力就把它的巢找到。寄生虫的卑鄙手段上面已经详细地报述过了，提防寄生虫入侵是每
一个母亲为了保护下一代所必须做的一项重要事业。况且这种迷宫蛛的死敌寄生虫专门
吃新生的卵，如果找到迷宫蛛的巢，会毫不客气地把它毁灭。所以聪明而尽责的迷宫蛛
就趁着夜色到各处去察看地形，找一个最安全的地方作为未来家族的安乐窝，至于那个
地方美观不美观，环境怎么样，倒是次要考虑的了。那沿着地面生长的矮矮的荆棘丛，
它们的叶子在冬天里也不会脱落，而且它们还能钩住附近的枯叶，对迷宫蛛来说，不能
说不是一个理想的居处，还有那又矮又细的迷迭香丛，也是迷宫蛛爱做巢的地方。在这
种地方，我常常能够找到不少迷宫蛛的巢。
有许多蜘蛛产卵以后就永远离开自己的巢了。可是迷宫蛛和蟹蛛一样，会一直紧紧
地守着巢。不过和蟹蛛不同的是，它不会像蟹蛛那样绝食，以致日益消瘦下去，它会照
常捕蝗虫吃。它用一团纷乱错杂的丝。筑起了一个捕虫箱，继续补充营养。
当它不捕食时，那就像我们所看到的那样，在走廊里踱来踱去，侧耳倾听四面八方
的动静。如果我用一根稻草在巢的某一处拨一下，它就会立即冲出来查看个究竟。就是
用这种警惕的办法，它尽心尽责地保护着自己的未成年的孩子们。
迷宫蛛产了卵后胃口还那么好，表示它还要继续工作。因为昆虫不像人类，有时候
吃东西仅仅是因为嘴馋。它们吃东西就是为了工作。
可是产完卵后，它这一生中最伟大的任务已经完成了，它还要作什么工作呢？经过
我细细地探究后，我才发现它所要做的工作是什么。大约又花了一个月左右的工夫，它
继续在巢的墙上添着丝。这墙最初是透明的，现在却变得又厚又不透明了。这就是它之
所以还要大吃特吃的原因：为了充实它的丝腺来为它的巢造一垛厚墙。
大约在九月中旬，小蜘蛛们从巢里出来了。但是它们并不离开巢，它们要在这温软
舒适的巢里过冬。母蜘蛛继续看护着它们，继续纺着丝线。不过岁月无情，它一天比一
天迟钝了。它的食量也渐渐地小起来。有时候我特意放几条蝗虫到它的陷阱里去喂它，
它也显得无动于衷，一口也不想吃。虽然这样，它还能维持四五个星期的生命，在它离
开这个世界之前，它继续一步不离地守着这巢，每次一听到巢里新生的小蛛在活泼地爬
来爬去，它便感到无限的满足和快慰。最后，到十月底的时候，它用最后一点力气替孩
子们咬破巢后，便精疲力尽地死去。它已尽了一个最慈爱的母亲所应尽的责任，它无愧
于它的孩子，无愧于这个世界。至于以后的事；它便托付给造物主了。到了来年的春天，
小蜘蛛们从它们舒适的屋里走出来，然后像蟹蛛那样，靠着它们的飞行工具——游丝，
飘散到各地去了。它们的母亲在天之灵，看着自己的孩子一个个各自成家立业了，也该
欣慰了。
------------------








蛛网的建筑



即使在最小的花园里，也能看到园蛛的踪迹。它们都算得上是天才的纺织家。
如果我们在黄昏的时候散步，我们可以从一丛迷迭香里寻找蛛丝马迹。我们所观察
的蜘蛛往往爬行得很慢，所以我们应该索性坐在矮树丛里看。那里的光线比较充足。让
我们再来给自己加一个头衔，叫做“蛛网观察家”吧！世界上很少有人从事这种职业，
而且我们也不用指望从这行业上嫌点钱。但是，不要计较这些，我们将得到许多有趣的
知识。从某种意义上讲，这比从事任何一个职业要有意思得多。
我所观察的都是些小蜘蛛。它们比成年的蜘蛛要小得多。而且它们都是在白天工作，
甚至是在太阳底下工作的，尽管它们的母亲只有在黑夜里才开始纺织。当到每年一定的
月份的时候，蜘蛛们便在太阳下山前两小时左右开始它们的工作了。
这些小蛛都离开了它们白天的居所，各自选定地盘，开始纺线。有的在这边，有的
在那边，谁也不打扰谁。我们可以任意地拣一只小蛛来观察。
让我们就在这只小蛛面前停下吧。它正在打基础呢。它在迷迭香的花上爬来爬去，
从一根枝端爬到另一根枝端忙忙碌碌的，它所攀到的枝大约都是十八寸距离之内的。太
远的它就无能为力了。渐渐地它开始用自己梳子似的后腿把丝从身体上拉出来，放在某
个地方作为基底，然后漫无规则地一会儿爬上，一会儿爬下，这样奔忙了一阵子后，结
果就构成了一个丝架子。这种不规则的结构正是它所需要的。这是一个垂直的扁平的
“地基”。正是因为它是错综交叉的，因此这个“地基”很牢固。
后来它在架子的表面横过一根特殊的丝，别小看这根细丝，那是一个坚固的网的基
础。这根线的中央有一个白点，那是一个丝垫子。
现在是它做捕虫网的时候了。它先从中心的白点沿着横线爬，很快就爬到架子的边
缘，然后以同样快的速度回到中心，再从中心出发以同样的方式爬到架子边缘，就这样
一会儿上，一会儿下，一会儿左，一会儿右。每爬一次便拉成一个半径，或者说，做成
一根辐。不一会儿，便这儿那儿地做成了许多辐，不过次序很乱。
无论谁，如果看到它已完成的网是那么地整洁而有规则，一定会以为它做辐的时候
也是按着次序一根根地织过去，然而恰恰相反，它从不按照次序做，但是它知道怎样使
成果更完美。在同一个方向安置了几根辐后，它就很快地往另一个方向再补上几条，从
不偏爱某个方向，它这样突然地变换方向是有道理的：如果它先把某一边的辐都安置好，
那么这些辐的重量，会使网的中心向这边偏移从而使网扭曲，变成很不规则的形状。所
以它在一边安放了几根辐后，立刻又要到另一边去，为的是时刻保持网的平衡。
你们一定不会相信，像这样毫无次序又是时时间断的工作会造出一个整齐的网。可
是事实确实如此，造好的辐与辐之间的距离都相等，而且形成一个很完整的圆。不同的
蜘蛛网的辐的数目也不同，角蛛的网有二十一根辐，条纹蜘蛛有三十二根，而丝光蛛有
四十二根。这种数目并不是绝对不变的，但是基本上是不变的，因此你可以根据蛛网上
的辐条数目来判定这是哪种蜘蛛的网。
想想看，我们中间谁能做到这一点：不用仪器，不经过练习，而能随手把一个圆等
分？但是蜘蛛可以，尽管它身上背着一个很重的袋子，脚踩在软软的丝垫上，那些垫还
随风飘荡，摇曳不定，它居然能够不加思索地将一个圆极为精细地等分。它的工作看上
去杂乱无序，完全不合乎几何学的原理，但它能从不规则的工作中得出有规则的成果来。
我们都对这个事实感到惊异。它怎么能用那么特别的方法完成这么困难的工作呢？这一
点我至今还在怀疑。
安排辐的工作完毕后，蜘蛛就回到中央的丝垫上。然后从这一点出发，踏着辐绕螺
旋形的圈子。它现在正在做一种极精致的工作。它用极细的线在辐上排下密密的线圈。
这是网的中心，让我们把它叫作“休息室”吧。越往外它就用越粗的线绕。圈与圈之间
的距离也比以前大。绕了一会，它离中心已经很远了，每经过一次辐，它就把丝绕在辐
上粘住。最后，它在“地基”的下边结束了它的工作。圈与圈之间的平均距离大约有三
分之一寸左右。
这些螺旋形的线圈并不是曲线。在蜘蛛的工作中没有曲线，只有直线和折线。这线
圈其实是辐与辐之间的横档所连成的。
以前所做的只能算作是一个支架，现在它将要在这上面做更为精致的工作。这一次
它从边缘向中心绕。而且圈与圈之间排得很紧，所以圈数也很多。
这种工作的详细情形很不容易看清，因为它的动作极为快捷而且振动得很厉害，包
括一连串的跳跃、摇摆和弯曲，使人看得眼花缭乱。如果分解它们的动作，可以看到它
的其中两条腿不停地动着，一条腿把丝拖出来传给另外一条腿，另一条腿就把这丝安在
辐上。由于丝本身有粘性，所以很容易在横档和丝接触的地方把新技出来的丝粘上去。
蜘蛛不停地绕着圈，一面绕一面把丝粘在辐上。它到达了那个被我们称作“休息室”
的边缘了。于是它立刻结束了它的绕线运动。以后它就会把中央的丝垫子吃掉。它这么
做是为了节约材料，它下一次织网的时候就可以把吃下的丝再纺出来用了。
有两种蜘蛛，也就是条纹蛛和丝光蛛，做好了网后，还会在网的下部边缘的中心织
一条很阔的锯齿形的丝带作为标记。有时候，它们还在这一条丝带的封面，就是网的上
部边缘到中心之间再织一条较短的丝带，以表明这是它们的作品，著作权不容侵犯。

粘性的网
蛛网中用来作螺旋圈的丝是一种极为精致的东西，它和那种用来做辐和“地基”的
丝不同。它在阳光中闪闪发光，看上去像一条编成的丝带。我取了一些丝回家，放在显
微镜下看，竞发现了惊人的奇迹：
那根细线本来就细得几乎连肉眼都看不出来，但它居然还是由几根更细的线缠合而
成的，好像大将军剑柄上的链条一般。更使人惊异的是，这种线还是空心的，空的地方
藏着极为浓厚的粘液，就和粘稠的胶液一样，我甚至可以看到它从线的一端滴出来。这
种粘液能从线壁渗出来，使线的表面有粘性。我用一个小试验去测试它到底有多大粘性：
我用一片小草去碰它，立刻就被粘住了。现在我们可以知道，园蛛捕捉猎物靠的并不是
围追堵截。而是完全靠它粘性的网，它几乎能粘住所有的猎物。可是又有一个问题出来
了：蜘蛛自己为什么不会被粘住呢？
我想其中一个原因是，它的大部分时间被用来坐在网中央的休息室里，而那里的丝
完全没有粘性。不过这个说法不能自圆其说，它无法一辈子坐在网中央不动，有时候，
猎物在网的边缘被粘住了。它必须很快地赶过去放出丝来缠住它，在经过自己那充满粘
性的网时，它怎么防止自己不被粘住呢？是不是它脚上有什么东西使它能在粘性的网上
轻易地滑过呢？它是不是涂了什么油在脚上？因为大家都知道，要使表面物体不粘，涂
油是最佳的办法。
为了证明我的怀疑，我从一只活的蜘蛛身上切下一条腿，在二硫化碳里浸了一个小
时，再用一个也在二硫化碳里浸过的刷子把这条腿小心地洗一下。二硫化碳是能溶解脂
肪的，所以如果腿上有油的话，这一洗就会完全洗掉了。现在我再把这条腿放到蛛网上，
它被牢牢地粘住了！由此我们知道，蜘蛛在自己身上，涂上了一层特别的“油”，这样
它能在网上自由地走动而不被粘住。但它又不愿老停在粘性的螺旋圈上，因为这种“油”
是有限的，会越用越少。所以它大部分时间呆在自己的“休息室”里。
从实验中我们得知这蛛网中的螺旋线是很容易吸收水分的。因为这个，当空气突然
变得潮湿的时候，它们就停止织网工作，只把架子、辐和“休息室”做好，因为这些都
不受水分的影响。至于那螺旋线的部分，它们是不会轻易做上去的，因为如果它吸收过
多的水分，以后就不能充分地吸水解潮了。有了这螺旋线，在极热的天气里，蛛网也不
会变得干燥易断，因为它能尽量地吸收空气中的水分以保持它的弹性并增加它的粘性。
哪一个捕鸟者在做网的时候，在艺术上和技术上能比得上蜘蛛呢？而蜘蛛织这么精致的
网只是为了捕一只小虫！真是有点大材小用了！
同时蜘蛛还是一个热忱积极的劳动者。我曾计算过，角蛛每做一个网需制造大约二
十码长的丝，至于那更精巧的丝，光蛛就得造出三十码，在这两个月中，我的角蛛邻居
几乎每夜都要修补它的网。这样，在这个时期中，它就得从它娇小瘦弱的身体上绵绵不
断抽出这种管状的，富有弹性的丝。
我们不禁要怀疑，它小小的身体怎么能产出那么多丝？它怎么能把这些丝搓成管状，
又怎么能在里面灌上粘液呢？它又怎能有时制出普通的丝，有时造出云朵状的丝花来垫
巢，最后还能制出黑色的丝带来装饰巢呢？这些问题一直在我的脑子里盘绕，并使我百
思而不得其解。
------------------








蜘蛛的几何学



当我们观察着园蛛，尤其是丝光蛛和条纹蛛的网时，我们会发现它的网并不是杂乱
无章的，那些辐排得很均匀，每对相邻的辐所交成的角都是相等的；虽然辐的数目对不
同的蜘蛛而言是各不相同的，可这个规律适用于各种蜘蛛。
我们已经知道，蜘蛛织网的方式很特别，它把网分成若干等份，同一类蜘蛛所分的
份数是相同的。当它安置辐的时候，我们只见它向各个方向乱跳，似乎毫无规则，但是
这种无规则的工作的结果是造成一个规则而美丽的网，像教堂中的玫瑰窗一般。即使他
用了圆规、尺子之类的工具。没有一个设计家能画出一个比这更规范的网来。
我们可以看到，在同一个扇形里，所有的弦，也就是那构成螺旋形线圈的横辐，都
是互相平行的，并且越靠近中心，这种弦之间的距离就越远。每一根弦和支持它的两根
辐交成四个角，一边的两个是钝角，另一边的两个是锐角。而同一扇形中的弦和辐所交
成的钝角和锐角正好各自相等——因为这些弦都是平行的。
不但如此，凭我们的观察，这些相等的锐角和钝角，又和别的扇形中的锐角和钝角
分别相等，所以，总的看来，这螺旋形的线圈包括一组组的横档以及一组组和辐交成相
等的角。
这种特性使我们想到数学家们所称的“对数螺线”。这种曲线在科学领域是很著名
的。对数螺线是一根无止尽的螺线，它永远向着极绕，越绕越靠近极，但又永远不能到
达极。即使用最精密的仪器，我们也看不到一根完全的对数螺线。这种图形只存在科学
家的假想中，可令人惊讶的是小小的蜘蛛也知道这线，它就是依照这种曲线的法则来绕
它网上的螺线的，而且做得很精确。
这螺旋线还有一个特点。如果你用一根有弹性的线绕成一个对数螺线的图形，再把
这根线放开来，然后拉紧放开的那部分，那么线的运动的一端就会划成一个和原来的对
数螺线完全相似的螺线，只是变换了一下位置。这个定理是一位名叫杰克斯·勃诺利的
数学教授发现的，他死后，后人把这条定理刻在他的墓碑上，算是他一生中最为光荣的
事迹之一。
那么，难道有着这些特性的对数螺线只是几何学家的一个梦想吗？这真的仅仅是一
个梦、一个谜吗？那么它究竟有什么用呢？
它确实广泛的巧合，总之它是普遍存在的，有许多动物的建筑都采取这一结构。有
一种蜗牛的壳就是依照对数螺线构造的。世界上第一只蜗牛知道了对数螺线，然后用它
来造壳，一直到现在，壳的样子还没变过。
在壳类的化石中，这种螺线的例子还有很多。现在，在南海，我们还可以找到一种
太古时代的生物的后代，那就是鹦鹉螺。它们还是很坚贞地守着祖传的老法则，它们的
壳和世界初始时它们的老祖宗的壳完全一样。也就是说，它们的壳仍然是依照对数螺线
设计的。并没有因时间的流逝而改变，就是在我们的死水池里，也有一种螺，它也有一
个螺线壳，普通的蜗牛壳也是属于这一构造。
可是这些动物是从哪里学到这种高深的数学知识的呢？又是怎样把这些知识应用于
实际的呢？有这样一种说法，说蜗牛是从蠕虫进化来的。某一天，蠕虫被太阳晒得舒服
极了，无意识地揪住自己的尾巴玩弄起来，便把它绞成螺旋形取乐。突然它发现这样很
舒服，于是常常这么做。久而久之便成了螺旋形的了，做螺旋形的壳的计划，就是从这
时候产生的。
但是蜘蛛呢？它从哪里得到这个概念呢？因为它和蠕虫没有什么关系。然而它却很
熟悉对数螺线，而且能够简单地运用到它的网中。蜗牛的壳要造好几年，所以它能做得
很精致，但蛛网差不多只用一个小时就造成了，所以它只能做出这种曲线的一个轮廊，
尽管不精确，但这确实是算得上一个螺旋曲线。是什么东西在指引着它呢？除了天生的
技巧外，什么都没有。天生的技巧能使动物控制自己的工作，正像植物的花瓣和小蕊的
排列法，它们天生就是这样的。没有人教它们怎么做，而事实上，它们也只能作这么一
种，蜘蛛自己不知不觉地在练习高等几何学，靠着它生来就有的本领很自然地工作着。
我们抛出一个石子，让它落到地上，这石子在空间的路线是一种特殊的曲线。树上
的枯叶被风吹下来落到地上，所经过的路程也是这种形状的曲线。科学家称这种曲线为
抛物线。
几何学家对这曲线作了进一步的研究，他们假想这曲线在一根无限长的直线上滚动，
那么它的焦点将要划出怎样一道轨迹呢？答案是：垂曲线。这要用一个很复杂的代数式
来表示。如果要用数字来表示的话，这个数字的值约等于这样一串数字1＋1/1＋1/1*2＋
1/1*2*3＋1/1*2*3*4＋……的和。
几何学家不喜欢用这么一长串数字来表示，所以就用“e”来代表这个数。e是一个
无限不循环小数，数学中常常用到它。
这种线是不是一种理论上的假想呢？并不，你到处可以看到垂曲线的图形：当一根
弹性线的两端固定，而中间松驰的时候，它就形成了一条垂曲线；当船的帆被风吹着的
时候，就会弯曲成垂曲线的图形；这些寻常的图形中都包含着“e”的秘密。一根无足轻
重的线，竟包含着这么多深奥的科学！我们暂且别惊讶。一根一端固定的线的摇摆，一
滴露水从草叶上落下来，一阵微风在水面拂起了微波，这些看上去稀松平常、极为平凡
的事，如果从数学的角度去研究的话，就变得非常复杂了。
我们人类的数学测量方法是聪明的。但我们对发明这些方法的人，不必过分地佩服。
因为和那些小动物的工作比起来，这些繁重的公式和理论显得又慢又复杂。难道将来我
们想不出一个更简单的形式，并使它运用到实际生活中吗？难道人类的智慧还不足以让
我们不依赖这种复杂的公式吗？我相信，越是高深的道理，其表现形式越应该简单而朴
实。
在这里，我们这个魔术般的“e”字又在蜘蛛网上被发现了。在一个有雾的早晨，这
粘性的线上排了许多小小的露珠。它的重量把蛛网的丝压得弯下来，于是构成了许多垂
曲线，像许多透明的宝石串成的链子。太阳一出来，这一串珠子就发出彩虹一般美丽的
光彩。好像一串金钢钻。“e”这个数目，就包蕴在这光明灿烂的链子里。望着这美丽的
链子，你会发现科学之美、自然之美和探究之美。
几何学，这研究空间的和谐的科学几乎统治着自然界的一切。在铁杉果的鳞片的排
列中以及蛛网的线条排列中，我们能找到它；在蜗牛的螺线中，我们能找到它；在行星
的轨道上，我们也能找到它，它无处不在，无时不在，在原子的世界里，在广大的宇宙
中，它的足迹遍布天下。
这种自然的几何学告诉我们，宇宙间有一位万能的几何学家，他已经用它神奇的工
具测量过宇宙间所有的东西。所以万事万物都有一定的规律。我觉得用这个假设来解释
鹦鹉螺和蛛网的对数螺线，似乎比蠕虫绞尾巴而造成螺线的说法更恰当。
------------------








蜘蛛的电报线



在六种园蛛中，通常歇在网中央的只有两种，那就是条纹蜘蛛和丝光蜘蛛。它们即
使受到烈日的焦灼，也决不会轻易稍离开网去阴凉处歇一会儿。至于其它蜘蛛，它们一
律不在白天出现。它们自有办法使工作和休息两个互不相误，在离开它们的网不远的地
方，有一个隐蔽的场所，是用叶片和线卷成的。白天它们就躲在这里面，静静地，让自
己深深地陷入沉思中。
这阳光明媚的白天虽然使蜘蛛们头晕目眩，却也是其它昆虫最活跃的时候：蝗虫们
更活泼地跳着，蜻蜓们更快活地飞舞着。所以正是蜘蛛们捕食的好时机，那富有粘性的
网虽然晚上是蜘蛛的居所，白天还是一个大陷阱，如果有一些又粗心又愚蠢的昆虫碰到
网上，被粘住了，躲在别处的蜘蛛是否会知道呢？不要为蜘蛛会错失良机而担心，只要
网上一有动静，它便会闪电般地冲过来。它是怎么知道网上发生的事的呢？让我来解释
吧。
使它知道网上有猎物的是网的振动，而不是它自己的眼睛。为了证明这一点，我把
一只死蝗虫轻轻地放到有好几只蜘蛛的网上，并且放在它们看得见的地方。有几只蜘蛛
是在网中，有几只是躲在隐蔽处，可是它们似乎都不知道网上有了猎物。后来我把蝗虫
放到了它们面前，它们还是一动不动。它们似乎瞎了，什么也看不见。于是我用一根长
草拨动那死蝗，让它动起来，同时使网振动起来。
结果证明：停在网中的条纹蛛和丝光蛛飞速赶到蝗虫身边；其它隐藏在树叶里的蜘
蛛也飞快地赶来，好像平时捉活虫一般，熟练地放出丝来把死蝗虫捆了又捆，缠了缠，
丝毫不怀疑自己是不是在浪费宝贵的丝线，由这个实验可见，蜘蛛什么时候出来攻击猎
物，完全要看网什么时候振动。
如果我们仔细观察那些白天隐居的蜘蛛们的网，我们可以看到从网中心有一根丝一
直通到它隐居的地方，这根线的长短大约有二十二寸；不过角蛛的网有些不同，因为它
们是隐居在高高的树上的，所以它的这根丝一般有八九尺长。
这条斜线还是一座桥梁，靠着它，蜘蛛才能匆匆地从隐居的地方赶到网中，等它在
网中央的工作完毕后，又沿着它回到隐居的地方，不过这并不就是这根线的全部效用。
如果它的作用仅仅在于这些的话，那么这根线应该从网的顶端引到蜘蛛的隐居处就可以
了。因为这可以减小坡度，缩短距离。
这根线之所以要从网的中心引出是因为中心是所有的辐的出发点和连接点，每一根
辐的振动，对中心都有直接的影响。一只虫子在网的任何一部分挣扎，都能把振动直接
传导到中央这根线上。所以蜘蛛躲在远远的隐蔽处，就可以从这根线上得到猎物落网的
消息。这根斜线不但是一座桥梁，并且是一种信号工具，是一根电报线。
年轻的蜘蛛都很活泼，它们都不懂得接电报线的技术。只有那些老蜘蛛们，当它们
坐在绿色的帐幕里默默地沉思或是安详地假寐的时候，它们会留心着电报线发出的信号，
从而得知在远处发生的动静。
长时间的守候是辛苦的，为了减轻工作的压力和好好休息。同时又丝毫不放松对网
上发生的情况的警觉，蜘蛛总是把腿搁在电报线上。这里有一个真实的故事可以证明这
一点。
我曾经打到一只在两棵相距一码的常青树间结了一张网的角蛛。太阳照得丝网闪闪
发光，它的主人早已在天亮之前藏到居所里去了。如果你沿着电报线找过去，就很容易
找到它的居所。那是一个用枯叶和丝做成的圆屋顶。造得很深，蜘蛛的身体几乎全部隐
藏在里面，用后端身体堵住进口。
它的前半身埋在它的居所里，所以，它当然看不到网上的动静了——即使它有一双
敏锐的眼睛也未必看得见，何况它其实是个半瞎子呢！那么在阳光灿烂的白天，它是不
是就放弃捕食了呢？让我们再看看吧。
你瞧，它的一条后腿忽然伸出叶屋，后腿的顶端连着一根丝线，而那线正是电报线
的另一个端点！我敢说，无论是谁，如果没有看见过蜘蛛的这手绝活，即把手(即它的脚
端)放在电报接收器上的姿势，他就不会知道动物表现自己智慧的最有趣的一个例子。让
猎物在这张网上出现吧，让这位假寐的猎手感觉到电报传来的信号吧！我故意放了一只
蝗虫在网上——以后呢？一切都像我预料的那样，虫子的振动带动网的振动，网的振动
又通过丝线——“电报线”传导到守株待兔的蜘蛛的脚上。蜘蛛它为得到食物而满足，
而我比它更满意：因为我学到了我想学的东西。
还有一点值得讨论的地方。那蛛网常常要被风吹动，那么电报线是不是不能区分网
的振动是来自猎物的来临还是风的吹动呢？事实上，当风吹动引起电报线晃动的时候，
在居所里闭目养神的蜘蛛并不行动，它似乎对这种假信号不屑一顾。所以这根电报线的
另外一个神奇之处在于，它像一台电话，就像我们人类的电话一样，能够传来各种真实
声音。蜘蛛用一个脚趾接着电话线，用腿听着信号，还能分辨出囚徒挣扎的信号和风吹
动所发出的假信号。
------------------








蟹蛛



前面我们讲到过的条纹蜘蛛虽然工作很勤快，为了替它的卵造一个安乐窝，一直孜
孜不倦地废寝忘食地工作着。可是到了后来，它却不能再顾到它的家了。为什么呢？因
为它寿命太短。它在第一个寒流到达之时就要死了。而它的卵要过了冬天才能孵化。它
不得不丢下它的巢。如果小宝宝在母亲还在世的时候能出世，我相信母蛛对小蛛的细心
呵护不会亚于鸟类。另外一种蜘蛛证明了我的推测；它是一种不会织网的蜘蛛：只是等
着猎物跑近它才去捉，而且它是横着走路的，有点儿像螃蟹，所以叫蟹蛛。
这种蜘蛛不会用网猎取食物，它的捕食方法是：埋伏在花的后面等猎物经过，然后
上去在它颈部轻轻一刺，你别小看这轻轻的一刺，这能致它的猎物于死地。我所观察到
的这群蟹蛛尤其喜欢捕食蜜蜂。
蜜蜂采花蜜的时候是很专心致志的，什么都不会想不会开小差。它用舌头舔着花蜜，
然后挑选一个能采到许多花蜜的花蕊上，一心一意地开始工作。当它正埋头苦干的时候，
蟹蛛早就虎视耽耽从隐藏的地方偷偷地、悄悄地爬出来，走到蜜蜂背后，越走越近，然
后一下子冲上去，在蜜蜂颈背上的某一点刺了一下。蜜蜂无论怎么挣扎也摆脱不了那一
刺。
这一刺可不是随随便便出手的。它刚好刺在蜜蜂颈部的神经中枢上。蜜蜂的神经中
枢被麻痹了，腿也开始硬化，不能动弹了。一秒钟内，一个小生命就宣告结束了。蟹蛛
这个凶手快乐而满足地吸着它的血，吸完后抹抹嘴巴残酷地把它的尸体的遗骸丢在一边。
可是，我们将会看到这个残暴的刽子手在家里是个非常慈爱温和的母亲。就像“食
人怪”一样，虽然吃起别人的孩子来毫不留情，却十分疼爱自己的孩子。在饥饿的逼迫
下，我们也是一样的，人和畜牲都是食人怪。
蟹蛛虽然是个杀蜂不见血的凶手，但你又不得不承认，它也是一只非常漂亮的小东
西。虽然它们的身材并不十分好，像是一个雕在石基上的又矮的又胖的锥体，在其中的
一边还有一块小小的隆起的肉，好像骆驼的驼蜂一样，但是它们的皮肤比任何绸缎都要
好看，有的是乳白色的，有的是柠檬色的。它们中间有些特别漂亮：腿上有着粉红色的
环，背上镶着深红的花纹，有时候在胸的左边或右边还有一条淡绿色的带子，这身打扮
虽然不像条纹蛛那么富丽，但是由于它的肚子不那么松驰，花纹又细致，色彩鲜艳又搭
配协调，所以，看起来倒反而比条纹蛛的衣服典雅、高贵。人们见了别种蜘蛛都敬而远
之，但对美丽的蟹蛛却怎么也怕不起来，因为它实在长得太漂亮，太可爱了，如果它们
是一些不会动的小东西，大家一定会对它们爱不释手。

蟹蛛的巢
在筑巢方面，蟹蛛的高超手艺并不亚于它在觅食时的技艺。
有一次我在一株水蜡树上找到它，当时它正在一丛花的中间筑巢，它织着一只白色
的丝袋，形状像一个顶针：这个白色的丝袋就是它的卵的安乐窝，袋口上还盖着一个又
圆又扁的绒毛盖子。
在屋顶的上部有一个用绒线张成的圆顶，里面还夹杂着一些凋谢了的花瓣，这就是
它的了望台。从外面到了望台上，有一个开口作为通道。
就在这了望台上，蟹蛛像一名尽心尽责的卫兵一样，天天守在这里。自从产了卵之
后，它比以前消瘦多了，差不多完全失去了以前那朝气蓬勃的样子，它全神贯注地守在
这了望台上，一有风吹草动就会全身紧张，埋没入战备状态，然后从那儿走出来，挥着
一条腿威吓来惊扰它的不速之客，它激动地做着手势，叫它赶紧滚开，否则后果自负。
它那狰狞的样子和激动的动作的确把那些或怀有恶意或无辜的外来者吓了一跳，把那些
鬼鬼祟祟的家伙赶走以后，它才心满意足地回到自己的岗位上。
那么它又在那丝和花瓣做成的穹顶下做什么呢？原来它在舒展着身体来遮蔽它宝贝
的卵。尽管此时它已经非常瘦小孱弱了仿佛一阵风就能把它卷走，它已忘记了饮食，为
了守望工作不被影响，它现在已抛弃了睡眠，不再去捕蜜蜂，吸它们的血充饥，它只是
静静地坐在自己的卵上。
我不由想到母鸡。母鸡孵蛋的时候也是这个样子，不同的是母鸡的身体可以提供热
量，当它孵蛋的时候，它身上的暖气传导到卵上唤醒了生命的种子。而对于蜘蛛而言，
太阳提供的热量已经足够了。蜘蛛母亲不需要再提供热量了。事实上，它也没有什么力
量提供热能了。因为这个不同点，我们不能称蜘蛛对小蛛的守候为“孵育下一代”。
两三个星期后，母蛛因为滴水未进，一点没吃，而愈来愈来瘦了。可它的守望工作
却丝毫不见有松懈。在临死之前，它还有什么愿望呢？它似乎一直在等待什么，这等待
使它苦苦地支撑着自己，用它的精神撑起早已没有活力的身体。它究竟在等什么呢？是
什么值得它这样用生命去苦苦等待呢？后来我才知道它是在等它的孩子们出来，这个垂
死的母亲还能为孩子们尽一点力。
条纹蛛的孩子们在离开那气球形的巢之前早已成为孤儿，没有人来帮它们打破巢，
它们自己又没有能力破巢而出，只有等巢自动裂开才能把小蛛们送到四面八方。它们出
来时根本连自己的母亲是谁都不知道。而蟹蛛的巢封闭得很严密，又不会自动裂开，顶
上的盖也不会自动升起，那么小蛛是怎么出来的呢？等小蛛孵出后我们会发现在盖的边
缘有一个小洞。这个洞在以前的时候是没有的，显然是谁暗中帮助小蜘，为它们在盖子
上咬了一个小孔，便于让它们钻出来。可是又是谁悄悄地在那儿开了一个洞呢？
袋子的四壁又厚又粗，微弱的小蛛们自己是决不能把它抓破的，其实这洞是它们那
奄奄一息的母亲打的。当它感觉到袋子里的小生命不耐烦的骚动的时候，它知道孩子们
急于想出来，于是就用全身的力气在袋壁上打了一个洞。母蛛虽然衰弱得可以随时死去，
但为了为它的家庭尽最后一份力，它一直顽强地支撑了五六个星期。然后把全身的力量
积聚到一点上爆发出来打这个洞，这个任务完成之后，它便安然去世了。它死的时候非
常平静，脸上带着安祥的神情，胸前紧紧抱着那已没有用处的巢，慢慢地缩成一个僵硬
的尸体。多么伟大的母亲啊！母蚁的牺牲精神令人感动，可是和蟹蛛相比，似乎还略逊
一筹。
在七月里，我的实验室中的小蟹蛛从卵里出来了。我知道它们有攀绳的嗜好，所以
我把一捆细树枝插在它们的笼上。果然，它们立刻沿着铁笼很快地爬到树枝的顶端，又
很快地用交叉的丝线织成互相交错的网，这便是它们的空中沙发。它们安静地在这沙发
上休息了几天，后来它们就开始搭起吊桥来。
我把爬着许多小蛛的树枝拿到窗口的一张桌子上，然后把窗户打开。不久小蛛们便
开始纺线做它们的飞行工具了，不过它们做得很慢，因为它们总是三心二意的，一会儿
爬到树枝下面，一会儿又回到顶上，好像不知道自己要干什么，又不知道该怎么干。
照这种速度，它们在那儿忙乎了半天也没什么成果，它们都急于要飞出去，可是就
是没胆量。在十一点钟的时候，我把载着它们的树枝拿到窗栏上，使让太阳照射到它们
身体上。几分钟以后，太阳的光和热射入它们的身体后积聚起来，成为一个小发动机，
驱使小蛛们纷纷活跃起来。只见它们的动作越来越快，越来越敏捷，都一个劲地往树枝
的顶上爬去，尽管我不能确切地看到它纺着线往空中飞去，但我很相信它们此刻正在树
梢上飞快地纺线，蓄势待发呢！
有三四只蟹蛛同时出发了，但各走各的路，自向不同的方向，其余的也纷纷爬到顶
上，后面拖着丝。突然起了一阵风，啊呀，那些蜘蛛那样的轻巧，它们编的丝又那么细，
风会把它们卷走吗？
我仔细看了看，风的确猝然把细丝扯断了，小蛛们顺着风在空中飘荡了一会儿便随
着它们的降落伞——断丝飘走了。我望着它们离去的背影，直到它在我的视野里消失。
它越飞越远，离出发点有四十尺远了。在又黑又暗的柏树叶丛中，它犹如一颗闪亮的明
星。它越飞越高，越飞越远，终于看不见了。其余的小蛛也接着飞出去，有些飞得很高，
有些飞得很低。有的往这边，有的往那边，最终都找到了自己的安身立命之处。
这时候，所有的小蛛都准备起飞了。而现在已不是开始的时候那样三三两两地飞出，
而是呈放射线状一队一队地飞出了，也许几个先锋的英雄行为感染激励了它们。不久它
们就陆续安全着陆了，有的在远处，有的在近处，这个简单的降落伞成功地完成了它的
使命。
关于它们以后故事，我就不知道了。在它们还没有力量刺蜜蜂的时候，它们怎么捕
食小虫子呢？小虫子和小蜘蛛争斗的话，谁又会占上风呢？它会耍什么阴谋吗？它会受
哪些天敌的威胁呢？我都不得而知。不过等到明年夏天，我们可以看到它已经长得很肥
很大，纷纷躲在花丛里偷袭那些勤劳采蜜的蜜蜂了。

\newpage

% ========== 以下内容来自EPUB《昆虫记》（全10卷）2020版 ==========

\chapter{《昆虫记》总序}

\section*{享受生命的细节之美}

1915年9月，91岁高龄的法布尔在家人的搀扶下，坐在轮椅上最后一次巡视了他毕生钟爱的"荒石园"。在这块"矢车菊与昆虫"钟爱的土地上，法布尔用30年的时间，完成了十卷横跨科学与文学领域、史无前例的伟大经典—《昆虫记》。

法布尔全名让-亨利·卡西米尔·法布尔，1823年出生于法国南部城市普罗旺斯的一户农民家中。好奇心重的法布尔从小就表现出对自然异乎寻常的热爱和出众的观察力。15岁时，他以公费生第一名的成绩考进亚威农师范学校。毕业后，18岁的法布尔在卡本特拉开始了他的教师生涯，教授的课程便是自然科学史。在随后的几年内他继续努力自修，陆续获得文学、数学、物理学和其他自然科学的学士学位，并在1855年拿到科学博士学位。在取得博士学位后，他就决定终生致力于昆虫学的研究。

1857年，他发表了《节腹泥蜂习性观察记》。这篇论文修正了当时昆虫学祖师列翁·杜福尔的观点，赢得了法兰西研究院的赞誉，法布尔被授予实验生理学奖。其间，他曾写作出版了各种科普书籍，将科学新知与各类自然科学知识介绍给大众。执教期间，法布尔的授课方式灵活、自由，深受学生们喜爱，保守派与教会人士却因他在公开场合向妇女讲述花的生殖功能而抨击他，甚至中止了他的课程。各种流言蜚语的中伤，使法布尔心灰意冷，他辞去了学校的教职，也不得不放弃了到大学任教的愿望，举家迁往奥朗日定居，一住就是10余年。在这10余年里，法布尔完成了《昆虫记》第一卷的写作。也正是在这段时间里，他遭遇了丧子之痛，《昆虫记》第一卷末就有他怀念爱子的文句。

《昆虫记》法文原名为Souvenirs entomologiques，意为"关于昆虫学的回忆"。法布尔并不局限于传统的昆虫解剖和分类，而是直接在野地里进行实地观察，或将昆虫带回自己家中进行培养研究，并生动详尽地记录这些小生命的体貌特征、习性、喜好、生存技巧、蜕变、繁衍和死亡。细致的观察、形象的描绘，正是《昆虫记》的最大特色。《昆虫记》十卷中有许多脍炙人口的经典段落，比如对昆虫倒挂姿势的描述：

如在金属笼子里，椎头螳螂的幼虫停在一个地方后姿势始终如一，毫不改变。它用四只后爪的爪尖钩住网子，后背朝下，纹丝不动，高高挂在笼顶，四个悬点承受着整个身体的重量。倒挂栖驻姿势是如此艰难，然而苍蝇的倒挂姿势截然不同。苍蝇虽然也抓挂在天花板上，但它总要抽出时间松弛一下，随便飞一飞，操起正常姿势走一走，肚皮贴地，肢体舒展开晒晒太阳。

正是如此细致入微的描写，才会让每一个读到《昆虫记》的人都感到无比的震撼，那一只只小虫子，清晰地在我们眼前活跃着，那充满活力的自然和生命之美几乎触手可及。为了更好地对昆虫进行实地观察和研究，在第一卷《昆虫记》出版后，法布尔买下了塞里尼昂乡下的一栋房屋和一公顷的荒地，并将这块荒地命名为"荒石园"，从此他全身心投入到了对昆虫的观察与实验中，直至逝世。就是在这里，他一边进行观察和实验，一边整理前半生研究昆虫的观察笔记、实验记录和科学札记，完成了《昆虫记》后九卷的写作。

《昆虫记》被人们称为"昆虫的荷马史诗"，与哥白尼的《天体运行论》、达尔文的《物种起源》一起被认为是人们必读的科普经典。不同于一般的科普书籍或百科全书，《昆虫记》散发着浓郁的文学气息，具有一种如孩子般天真的人文品格。《昆虫记》中没有种属科目的严格分类，而是由那些或美丽或丑陋的昆虫自己出场，展示出一幅幅生动旖旎的昆虫生活风俗画。书中充满了像孩子一样对未知的渴求、对自然的热爱、对生命的敬畏，也充满了法布尔本人与自然融为一体的生命状态。他会为发现昆虫天赋本能的不同而高兴，为圣甲虫成功推动粪球而激动，在他一次次不厌其烦的描述下隐藏着睿智的思考，让我们这些现代人又亲切地找回了丧失已久的内心的纯净与静谧，唤起了我们对自然万物的热爱。自始至终，法布尔的《昆虫记》都是谦逊的、平和的，没有对生命的随意摆布，没有可笑的妄自尊大，他孜孜不倦地在"荒石园"中观察、记录了30多年，平和地向我们传达着那些关于昆虫们最朴素也最深刻的知识，也让我们在这种平和、温暖的心态下，重新思考生命的含义与价值。

《昆虫记》传入中国，始于20世纪20年代，首倡者就是著名的文学家鲁迅（周树人）和周作人两兄弟。周作人在1923年发表的短文《法布耳昆虫记》中说道：

法布耳的书中所讲的是昆虫的生活，但我们读了却觉得比看那些无聊的小说戏剧更有趣味，更有意义。他不去做解剖和分类的工作（普通的昆虫学里已经说得够多了），却用了观察与试验的方法，实地地记录昆虫的生活现象，以及本能和习性之不可思议的神妙与愚蒙。我们看了小说戏剧中所描写的同类的命运，受得深切的铭感，现在见了昆虫界的这些悲喜剧，仿佛是听说远亲—的确是很远的远亲—的消息，正是一样迫切的动心，令人想起种种事情来。他的叙述，又特别有文艺的趣味，更使他不愧有昆虫的史诗之称。戏剧家罗斯丹批评他说，"这个大科学家像哲学者一般地想，美术家一般地看，文学家一般地感受而且抒写"，实在可以说是最确切的评语。默忒林克称他为"昆虫的荷马"，也是极简明的一个别号。

周作人还曾感慨地说："羡慕有这样好书看的别国少年，也希望中国人来做这翻译编纂的事业，即使在现在的混乱秽恶之中。"鲁迅自1924年起就在收集《昆虫记》的日语译文，多次在他的文章中强调《昆虫记》的重要意义，并对《昆虫记》的翻译寄予厚望。正是周氏兄弟的这般推崇、介绍，国人才开始知道并翻译法布尔的《昆虫记》。

从世界范围来讲，《昆虫记》的翻译都是一项十分巨大的工程，最早开始翻译《昆虫记》的英、美两国都没有出版过严格意义上的全译本；而追踪欧美文化最快最多的日本倒是从一开始就尽量求全，第一套全日文译本的《昆虫记》由大杉荣在20世纪20年代翻译完成。自20世纪30年代到现在，在我国面世的《昆虫记》译本越来越多，但大多以选译或节译本为主，真正意义上的全译本寥寥可数。各版本由于编译初衷不同，虽然各有特点但也留下了不少遗憾，或为单册内容完整而删改原文，或为其文学阅读性而忽略科学严谨性，或为其全而文字不易阅读。为此我们特别推出了这套十卷全译插图版的《昆虫记》。

本套书在最大限度重现《昆虫记》原著全貌的同时兼顾原著的文学阅读性，是一套真正适合阅读的全译本《昆虫记》。本书以法文版十卷本原著为底本进行分卷、翻译，同时参照对比了美国人的分专题多卷选译本、日本大杉荣版《昆虫记》的内容，力求全尽，将一部完整的《昆虫记》奉献给读者。同时，我们的翻译工作在坚持"反映原著原貌"这一原则的基础上，特别强调了译文文字的文学性和可读性，希望通过我们的翻译，让读者充分感受到原著中朴实清新、生动活泼的语言，以及充满了盎然生机的情趣和诗意。对于原著中穿插的大量希腊神话、历史事件、《圣经》典故、拉丁文诗歌，我们也加入大量有针对性的注释，帮助读者更加方便、流畅地阅读。

《昆虫记》一书中提及了上千种昆虫，并提及了很多其他动物、植物等，内容繁多，为了让读者对书中所涉及昆虫和动植物有一个直观的认识，我们在编辑过程中特别加入了大量精确、细致的手绘插图，全十卷共计约1 800幅。这些图片有近200幅来自于《昆虫记》法文原著中的原始配图，其余近1 600幅则是我们根据文中法布尔的记述，参考实际昆虫形态设计并重新绘制的。书中的全部插图均由手工绘制，既精美又准确地描绘出了昆虫的真实形态。同时，每幅图片均配有简洁、专业的图说，与法布尔的文字意境可谓相得益彰。希望这些经过我们精心编辑的插图和图说能够帮助读者亲身融入19世纪法国南部普罗旺斯迷人的田园风光中，轻松地感受自然与生命带给我们的惊喜与感动。

《昆虫记》全套共十卷，文字逾200万字，编辑工程繁复、浩大，本套书最终能够顺利出版离不开各位参与者的共同努力。在此，特别感谢陈一青先生及其团队高质量的翻译，将这部跨越文学与科学的不朽经典完整地展现在读者面前；同时插画师郭警、陈勤对全套书约1 800幅插图进行了精心的绘制和修补，为读者还原了一个独一无二的昆虫世界；最后，在这套书的编辑过程中，得到了国内众多昆虫学家的帮助和指导。在此对所有为本书辛勤付出的同仁致以最诚挚的谢意。

\hfill\textbf{编者谨识}

\hfill 2020年1月

\newpage

% ========== 第1卷 高明的杀手 ==========

\part{昆虫记：第1卷 高明的杀手}

\begin{figure}[htbp]
    \centering
    
\includegraphics[width=0.8\textwidth]{images/00271.jpeg}
\end{figure}

\newpage

\chapter{圣甲虫}

\begin{figure}[htbp]
    \centering
    \includegraphics[width=0.8\textwidth]{images/00281.jpeg}
\end{figure}

事情的经过是这样的：我们五六个人中，我年纪最老，是他们的老师，更是他们的同伴和朋友。他们都是年轻人，充满了青春的活力与激情。我们都对自然万物充满了无限的热情与渴望。在一条盛开着接骨木和英国山楂树的山路上，树上的伞房花序散发出浓郁的苦涩香味，把金匠花金龟都陶醉了。我们沿着山路前行，一边谈天说地，一边猜测着圣甲虫是否已经莅临在安格尔多沙的高原上，并滚动着它那被古埃及人视为地球象征的粪球。我们猜测着那长着像珊瑚枝似的梭形尾巴的小蛛螺，想知道它是否游荡在山脚下的溪水中，是否在水面那绿毯般的浮萍下捉迷藏；我们猜测着小溪里那美丽的刺鱼，想知道它是否已经戴上了那蓝紫相间的结婚领带；初来乍到的燕子是否张开剪刀般的尾翼掠过草地，捕捉那正在一边跳舞一边下蛋的大蚊；长着眼状斑的蜥蜴是否从砂岩中的地穴中爬到阳光下，炫耀起它那布满蓝斑的臀部；从海上飞来的笑鸥是否成群结队地翱翔在河面上，追逐着溯罗讷河而上准备到内陆水域产卵的鱼群，并不时发出阵阵狂笑般的鸣叫……不过，我就描述到这里吧！为了简短表达，我要说，我们这些狂热地喜欢跟动物生活在一起的幼稚而纯朴的人们，将度过一个难以言喻的欢乐的上午，以此庆祝春天的万物复苏。

事情正如我们期望的那样：刺鱼已经梳妆打扮完毕，它那洁白闪亮的鳞片足以使白银都黯然失色，它甚至还在胸前抹上了最鲜艳的朱红色彩。一旦别有用心的黑色大蚂蟥靠近它，它背上、肋部的小刺就会像弹簧似的突然弹竖起来；而那个强盗在刺鱼这种坚定不可侵犯的态度面前，只得灰溜溜地钻进草堆里去了。扁卷螺、瓶螺、椎实螺，这些与世无争的软体动物也来到水面上呼吸着新鲜空气。水龟虫和它那丑陋的幼虫简直堪称是池塘里的海盗，它们扭动着脖子一窜而过，一会儿袭击袭击这个，一会儿骚扰骚扰那个；而周围那些昆虫却傻乎乎的，好像什么事儿也没发生一样。但是，我们还是先撇下平原上的水塘，去攀登那横亘在我们与高原之间的悬崖吧。在那里，绵羊在吃草，马儿好像在为筹备一场比赛而练习着奔跑。它们给快乐的食粪虫带来了美味。

把地上的粪便清除干净，这是鞘翅目食粪虫的工作，也是它们的崇高使命。食粪虫那种类齐全的工具真是令我们赞叹不已：这些工具有的用来翻动粪土，把粪土捣碎、整形；有的用来挖洞，以备将来带着战利品躲藏在这洞里。这些工具就像技术博物馆里陈列着的挖掘器械一样：有的好像是模仿人类的技艺制造的，有的甚至比人类的制造更独具匠心。说不定我们人类也可以仿效它们制造出新的器具来。

西班牙蜣螂的前额有一个强有力的角，像十字镐的镐头一样尖锐而向后翻翘。月形蜣螂除了长有这样的角以外，胸部还长着两片大大的犁桦状尖片，两个尖片之间有一根凸出的尖骨作为刮刀。生活在地中海一带的野牛粪屎蜣螂，前额有一对粗壮的岔开的角，自前胸长出的一片水平的犁铧伸到两角之间。米诺多蒂菲的前胸长着三片尖端朝前的平行尖犁，只是两旁的稍长些，中间的稍短些。牛粪屎蜣螂的工具是两个像牛角似的又弯又长的钳子，而长角粪屎蜣螂的工具是一根双刃长叉，直直地竖在扁平的头上。就连最差劲的金龟子也或者在头上或者在前胸长着一个凸出来的小硬块，而有耐心的昆虫能巧妙地使用这些圆钝的工具。所有的金龟子都配备着"铲子"——就是边缘锋利、大而扁平的头，同时它们都会像使用耙子一样，用带有锯齿的前腿将粪便耙拢在一起。

金龟子的劳动实在太肮脏了，也许是为了补偿这一缺憾吧，很多金龟子的体内都散发着麝香味，而且腹部都像磨亮的金属一样闪烁着光泽。有的粪金龟腹部发出像金铜一样的光亮，有的粪金龟腹部则闪耀着紫晶色。不过通常情况下，食粪虫牛蜣螂的腹部是黑色的。那些衣着华丽的、宛如佩戴着鲜艳首饰的食粪虫，都长在热带地区。那种生长在上埃及骆驼粪下的圣甲虫，它们浑身的绿色可与祖母绿相媲美；而圭亚那、巴西、塞内加尔的蜣螂，都闪耀着红色金属般的光泽，好似黄铜一样富丽堂皇，又像红宝石一样光彩照人。我们这里虽然没有这种宛如用粪便做成的首饰，但我们国家的食粪性甲虫的习性也同样引人注目。

那出现在一堆粪便四周的，是何等争先恐后的场面啊！从世界各地涌向加利福尼亚的冒险家们在开发砂金矿时，也没有表现出这样狂热的干劲儿！太阳刚刚升起，数以百计的大小不一、形态各异的食粪虫便已拥挤在那儿，互不相让、分秒必争地从那块大蛋糕上分着自己的一杯羹。它们有的在露天工作，梳耙着粪便的表层；有的在粪堆内部挖掘巷道，寻找更加优质的粪源；有的在底层开发，以便及时把战利品埋藏在身下的洞里；个头儿最小的只能寻找身强力壮的合作伙伴，它们总是先在一旁等待，当合作者在进行大规模发掘时，一旦有小块粪便滑落，它们就立刻走上前去将其加工成碎屑。有的初来乍到，可能实在饥肠辘辘了吧，竟然当场饱餐起来。不过大多数食粪虫都立志于积攒一笔丰厚的财产，然后躲进万无一失的隐藏所内，守着这些财产过一段衣食无忧的日子。在盛开着百里香的贫瘠的平原上，新鲜的粪便并不是随处可见的；像这样的意外收获堪称是上帝的恩赐，只有得天独厚的幸运者才有这样的福分，因此它们自然要尽最大的努力多得些财富，并小心翼翼地贮藏起来。在这之前，粪香飘扬到了方圆一千米的地方，于是所有的食粪虫都闻香而来，除了已经聚集在这块美味上的食粪虫外，路上还有很多迟到者们正或飞或跑地赶往这里呢！

这只唯恐迟到的、迈着碎步向粪块赶来的是什么虫呢？它那长腿就像被腹中的某个机械装置驱使着一样，行动起来生硬而笨拙；红棕色的触角像扇子似地张开着，透露出它垂涎欲滴的焦急心态。它终于赶到了，不甘示弱地挤了上来，挤倒了一些捷足先登者，跻身到了大餐桌跟前。这浑身黝黑、体形粗大的家伙，就是食粪虫中最鼎鼎有名的圣甲虫。现在，它已跟它的共享者们同席而坐了。它用巨大的前爪，一下一下地拍打着自己的粪球，做最后的加工。对粪球进行完这最后一层的加工后，它就可以走到一旁，安静地享受自己的劳动果实了。现在让我们看看这地道的粪球是怎样一步一步制造出来的吧！

圣甲虫的头顶上是宽大而扁平的月牙形顶壳，顶壳的前沿上排列着六个角形锯齿，这就是它们用以挖掘和切削的工具。与此同时，这带锯齿的顶壳也发挥着类似耙子的功效，能将那些难以下咽的植物纤维剔除掉，把最精美的食物统统耙过来，聚拢在一起。精选食料的工作就是这么进行的。圣甲虫如果只是为了自己吃，那么粗枝大叶地挑选一下就可以了；但是如果要制作"育儿室"——在粪球中挖一个孵卵的小洞，那就得精挑细选了。这些精明的行家是宁愿按后一种方式来行事的，于是便一丝不苟地把所有的纤维屑都剔除掉。育儿室的内层全部由粪便的精华建筑而成。这样，初生的幼虫破卵而出时就能在住所的内壁中找到精致的食物，为以后向粗糙的外层发起进攻做准备了。

圣甲虫在选择自己的食物时，毫不挑剔，只是大致挑选一番。它先用带锯齿的顶壳破土插入挖掘几下，然后漫不经心地搜索一番，把杂质剔除掉，然后再把剩下的精细的食物耙拢到一起。它在制作粪球时，强劲有力的前腿通力合作。它的前腿扁平，弯成弧状，表面纹脉凸起，上面还排列着五个坚齿。如果需要动用武力来推翻障碍，在粪团最厚的地方开辟一条道路，它就会舞动起一对齿足，左右并进，然后猛地横扫一耙，清出一个半圆形的地盘。场地清好后，前腿就开始另一项工作了，那就是一抱一抱地把顶壳耙过来的粪便聚拢到腹部下方的四只后腿之间。接下来，就轮到后腿干车工的工作了。圣甲虫的后腿，尤其是最后的两条腿十分细长，略微弯曲，呈弧形，足端长有利爪。一眼望去便会发现，这四条后腿好似两组球面圆规，两对弧形的弯腿之间环抱成一个球体，恰好可以检查和修正粪球的形状。事实就是如此，这些后腿就是用来加工粪球的。

当圣甲虫用前腿将粪便一抱一抱地聚拢在腹下四腿之间时，四条弧形后腿轻轻一压，粪便就成了圆球状，粪球也就初具雏形了。接下来，经过粗加工的粪团在两组球面圆规中间，由四只脚一边摇滚一边轻压。就这样，粪球经过腹部底下的不断旋转加工后，形状越来越完善。如果粪球的表层因黏合度差而造成粪土一片片剥落，或者某一处的粗纤维过多而难以旋转加工时，就需要前腿来修整这些有缺陷的部位：用前腿上的巨掌轻轻拍打粪球，新添一些粪土，把那些难以黏合的纤维屑裹到粪球里层去。

在炽热的阳光下，工程紧迫地进行着，这位车工干得如痴如狂，动作非常敏捷，这令我们惊叹不已。它们的工作进行得非常神速：刚刚还是一粒小粪丸，转眼间已变成核桃大的粪团了，再过一会儿又成了苹果那么大的粪球。我甚至曾经见过一些贪吃的圣甲虫们竟然把粪球做得像拳头那么大，肯定够它们吃上几天了。

储备食品制作好后，圣甲虫就要从分割蛋糕的混战中抽身而退，把食物运到安全稳妥的地方了。然而就是在运粪球的过程中，它表现出了最令人拍案叫绝的习性特征。这种食粪虫毫不迟疑地上路了，它用两条长长的后腿环抱着粪球，同时将两条后腿足端的利爪从左右两侧插入粪球，形成球体的旋转轴；中间的一对足爪按在地上，当作支架，支撑着粪球；长着锯齿的前腿交替在地面上推按粪球，充当杠杆的作用。就这样，圣甲虫带着重物，弓着身子，压低脑袋，屁股上翘，倒退着走去。它的这些一连串的肢体动作，宛如一套运行合理的机器，而在这套机器中，两条后腿是主要部件。它们始终不停地运动着，足端的利爪不断地移动位置，以调整粪球的轴心，使其保持平衡，并在两只前腿一左一右的交替推动下往前移动。这样，粪球表面的各个部位会依次与地面接触，由于整个球面滚压的力量均匀，不仅使粪球的外形更加完善，还使表面的坚硬程度也相同。

球在前进，在滚动，加油！它是一定能到达目的地的，当然，途中不会一帆风顺。你看，第一个难题出现了：它遇到了一个陡坡，沉重的粪球顺着斜坡一路滚下。圣甲虫却认准了这条路，偏要从这天然的道路上走过不可。这可是一个大胆的计划，因为只要它一步失足，或者有一颗砂粒碰到粪球，粪球就会失去平衡，它的计划就会破产，一切努力就会功亏一篑。果然，它迈错了一步，粪球滚入了路旁的沟底；圣甲虫被滑冲下来的粪球撞了个仰面朝天，六条腿在空中乱舞着。终于，它翻过身了，奔跑着追上粪球，使劲抱住，身体机器更起劲地运转起来了。当心啊，冒失鬼！为什么不顺着谷底走呢！既节省力气又不会出意外，那条路不错，十分平坦，你可以毫不费力地滚动你的粪球。可是，它偏不这么走，竟然打算重新攀上刚刚摔滑下来的陡坡。也许它应当返回高地，对此我有什么可说的呢？高度是可以加以利用的，看来在这个问题上，圣甲虫的见解比我高明。可是，你至少应该走这条小路吧，它坡度小，你准能爬上去。可它根本不屑一顾，如果附近有一道更陡峭的几乎无法攀登的斜坡，怕是这个固执的家伙都宁愿走上去！圣甲虫的工作又开始了。它步步为营、小心翼翼地倒退着走，千辛万苦地把粪球这个硕大的负载物推到了一定的高度。我们不免思索起来，究竟是靠着什么样的静力学奇迹，圣甲虫能在斜坡上抓牢这么一大团东西啊！糟了！一个动作失调，它又前功尽弃了，又连球带虫都滚了下去。它再攀登，但很快又跌了下来！再重新开始时，它做得更好了，谨慎地绕过了一根禾本植物的根。原来前几次它都是被这个植物根绊倒的。只剩一小段路就上去了，不过，它走得更慢了，非常非常慢。它似乎意识到这斜坡危机四伏，一招不慎就会全盘皆输。糟了，它一条腿在一块光滑的砾石上滑了一下，又连球带虫一道"叽里咕噜"地滚下去了。可是圣甲虫并不气馁，以百折不挠的执着精神又重新开始了。它要么十次、二十次劳而无功地攀登，直至顽强地战胜障碍；要么变得识时务些，承认那样做是白费力气，转而改从平地走。

圣甲虫搬运那珍贵的粪球时，并不总是单枪匹马地干，它经常给自己找个搭档，或者准确地说，是它的搭档主动找上门来。情况通常是这样的：粪球制作好以后，一只圣甲虫离开纷争混战的虫群和工作场地，倒退着把战利品推走。这时，旁边有只姗姗来迟、刚刚开始工作的金龟子，突然扔下自己的工作，跑向那正在滚动着的粪球，助幸运的粪球主人一臂之力，而这个粪球主人得到了帮助，当然也很乐意接受。于是两个伙伴协同合作，并竞相出力把粪球运到安全地点。它们合伙运送粪球，难道是因为它们在开采地上有心照不宣的协议或者某种平分糕点的默契吗？是否一只圣甲虫在揉制粪球时，另一只则开挖更好的粪源，从而把优质的材料提取出来添加到共有的储备食物上呢？我从来没有见过这样的合作。我看到的是每只圣甲虫都在开采地上忙着自己的事情，所以半路出现的援助者是没有任何与粪球主人分享劳动果实的权利的。

那么，这是不是异性之间的一种联合呢？是不是即将配偶成家的一雌一雄呢？有段时间，我曾这么认为。两只圣甲虫，一只在前，一只在后，怀着同样的劳动热情双双推动着沉重的粪球，这情形甚至令我想起从前的一首管风琴弹奏曲："要成家，唉！怎么办！你在前，我在后，咱们一起推酒桶。"然而，解剖的结果使我不得不放弃对这种温情脉脉的家庭牧歌的想象。雌性和雄性的圣甲虫在外貌上没有任何区别，于是我将搬运同一粪球的两只圣甲虫进行了解剖。结果，出现最多的情况是它们都是同性。

既不是家庭共同体，又不是劳动合伙人，那么它们之间的这种共事究竟是为了什么呢？答案竟然是图谋劫持。那位殷勤的搭档以助人一臂之力为幌子，却一直处心积虑地等待时机，想把粪球据为己有。在开采地上自己制作粪球既辛苦又需要耐心，如果别人做好后把它抢来，或者退一步，充当一个座上宾，这要划算得多啊！如果粪球主人稍不注意，它就携着粪球溜之大吉；如果主人寸步不离地严密监视它，最坏也就是两人共进午餐，因为它帮过忙。这样的战术真是有百利而无一弊。有的圣甲虫的确像我刚才说过的那样，殷勤地跑去帮助一个根本不需要它帮忙的伙伴。在善意援助的假象下，掩藏着极其卑鄙的贪婪野心。有的圣甲虫也许胆子更大，对自己的力量更有信心，干脆单刀直入，直截了当地把粪球抢走。

\newpage

\chapter{大笼子}

\begin{figure}[htbp]
    \centering
    \includegraphics[width=0.8\textwidth]{images/00291.jpeg}
\end{figure}

假如到书本中去搜寻关于食粪虫特别是圣甲虫习性方面的资料，我们会发现，这门科学带有鲜明的埃及法老时代的色彩，个别地方有失偏颇。资料显示，那在田野上滚动的粪球中孕育着一颗虫卵，粪球既是即将出生的幼虫的食物，也是摇篮。食粪虫父母们把粪球带到曲折坎坷的土地上滚动，是为了让它更圆润一些，如果粪球遭到撞击、颠簸而顺着斜坡滚落，父母们就把它埋在地下，交给大地这个庞大的孵化器。

对于书上记录的圣甲虫培育幼虫的方式，我十分怀疑。因为虫卵是十分娇弱的，即使包裹在柔软的物质中，也要格外小心，怎么能禁得住粪球摇篮的滚动和震荡呢？那孕育在胚胎中的小生命是柔弱的，几乎稍一碰触或者任何一个小小的闪失，都会将其扼杀，圣甲虫的父母又怎么会带着它一路颠簸跋山涉水呢？不，事情一定不是这样的！尤其是圣甲虫的母亲，它那母性的温柔是决不会容忍子女遭受雷古卢斯滚筒一样的酷刑的。但是不能只靠简单的逻辑推理就试图推翻先人的观点。为此，我曾切开数百个圣甲虫的粪球查看，还亲自到经过核实的圣甲虫的洞穴中取来粪球，然后一一打开。但是，我从未在粪球中发现虫卵和孕育虫卵的温室。我看到的全部都是匆忙中制作而成的粪球，十分粗糙，粪球里边除了食物还是食物，没有讲究的结构。圣甲虫就是因为有了这些食物，才得以闭门不出，吃几天美味盛宴，过几天安闲日子的。圣甲虫总是对别人的粪球垂涎三尺，甚至互相争抢，它们这么积极绝对不是为了争抢来一群需要照顾和养育的后代。如果说圣甲虫之间会互相盗窃虫卵，那简直太荒谬了，因为每只圣甲虫都有很强的繁衍后代的能力。因此，毫无疑问，在圣甲虫滚动着的粪球里根本没有虫卵。

为了研究圣甲虫孕育后代的问题，我做了一个伟大的尝试。我首先做了一个大笼子，在里面铺上沙土，然后捉来大约20只圣甲虫，把它们和蜣螂、圆裸胸金龟、牛粪屎蜣螂一起放到笼子里饲养，并经常为它们更换食物。可是，这竟然是我从事昆虫实验以来最失败的一次！失败就失败在了为它们更换食物上。我的邻居家养有马，可是要想得到马的粪便来喂养圣甲虫，我必须得到喂马的佣人的帮助。这位喂马的佣人起初对我的请求嗤之以鼻，直到我给了他些小银币，他才答应帮忙。这就相当于，我要为昆虫的每顿午餐都支付25个生丁。怕是圣甲虫自己也从没为午餐做过这么大一笔财政预算。可是，每天早上，约瑟夫在给马包扎之后，就会从我们两家花园之间的墙头上探过头，把手围成喇叭状放在嘴边，低声喊道："哎！哎！我去提一满钵马粪来。"做这种事儿，双方都应该小心些啊。您看，果不出所料，某一天，就在我们进行粪便交易时，被房东撞见了。他就此认定他的肥料全被偷到了墙的这一边，认定我偷盗他的卷心菜肥料来浇灌我的马鞭草和水仙。尽管我极力解释，可于事无补，而且我的理由听上去又是那么的荒唐可笑。约瑟夫遭到了一顿痛骂，主人不仅对他说三道四，还下了最后通牒，说他如果再做这样的事就解雇他，而且一再强调自己是一个说一不二的人。

这样之后，我弄了一个圆锥形的纸袋，每天拿着它悄悄地来到大路上，为我的昆虫们捡食物。多丢人啊，可我并不因此而感到羞耻。有一次我运气好，看到一只到阿维尼翁去的驴子经过我的门前，背上驮着雷纳尔堡或者巴邦塔纳的农产品，就在这时，它不失时机地留下了一团粪便。这太令人惊喜了，足够昆虫们吃上几天了，我立即将其收捡起来。总而言之，为了给我的昆虫们弄到一团粪便，我做了一切我能做的，计谋、手段，等待时机、四处奔波。如果仅仅有激情、热情和百折不挠的精神就能取得成功的话，那么，我的实验理所当然会成功，可我失败了。在狭小的空间中，它们没有进行伟大的演变，我的圣甲虫也因思乡心切而日渐憔悴，不久就抑郁而终了，连同它的秘密也一同带走了。圆裸胸金龟和牛粪屎蜣螂倒是没令我失望，为我提供了很多有价值的资料，在适当的时候我会引用到的。

我的实验过程不只是在笼子里饲养昆虫，还包括对它们进行直接的研究，可是这些研究的结果与我的期望大相径庭。正在我思量着需要一个助手的时候，一群兴高采烈地穿过高地走来的孩子映入了我的眼帘。那是一个星期四的早晨，几个孩子从邻村莱昂格尔走来，他们一边啃着苹果，一手拿着面包，把学校以及令人厌烦的功课都抛到了脑后。他们来到那光秃秃的小山丘就开始挖土，寻找驻军在射击练习时射下的铅弹。可是，他们此次远征的唯一收获也就是那几块价值1苏钱的铅块而已。在那绿茵茵的青草地上，点缀着老鹳草玫瑰红色的花瓣，青色的草、红色的花则为佩特拉阿拉伯增添了几分妖娆；在锋利的岩石间，半黑半白的大耳鹏时来时往，欢乐的鸣叫声回荡在上空；百里香花丛下的洞口旁，一只蟋蟀正奏起枯燥的交响乐，乐声弥漫了四周。

这个春天里的游玩，令孩子们兴致勃勃，更令他们自豪的是，他们拥有了一笔价值1苏钱的财富。这是捡铅弹的回报，有了这1苏钱，下个星期天，他们就可以到教堂门口的女商贩那儿买两大块薄荷糖——两大块，每块2里亚。

我走到那个最大的孩子面前，看到他充满智慧的脸庞，我不由得信心满满；其他孩子围着我们，继续啃着苹果。我说清楚我的来意以后，指着一个正在搬运粪球的圣甲虫对他们说，说不定哪个类似的粪球里就有个凹陷的小窝，窝里孕育着一个虫卵。我要他们做的事就是四处挖挖，同时留意圣甲虫的活动，找到有虫卵的粪球。为了引起他们的兴趣，使他们放弃为挣区区几里亚而挖铅弹的活动转而为我的实验服务，我答应付给他们更高的报酬：每找到一个有虫卵的粪球就付给他们1法郎——那可是一枚亮晶晶的值20个苏的硬币。当他们听说我会支付这么大一笔数目时，他们都目瞪口呆了，那天真的模样可爱极了。一块粪球竟然被我定出了这么高的价格，这可能模糊了他们头脑中的金钱概念。接着，为了证明我的诚意，我先付了订金：给他们每人几个苏。我们约定下周四在同一时间、同一地点兑现交易，到那时我会忠诚地履行诺言，对每一个得到那个珍宝的人支付应得的报酬。我说清楚以后，就让他们走了。他们一边离开一边议论纷纷地说："这可真不错，我们要争取每个人都能挣1个法郎！"孩子们对此充满了期许，不由得把掌心里的那几苏订金弄得叮当作响。子弹头被踩在了脚下，被遗忘了。远远地，我看到孩子们在平原上分散开来，开始分头寻找粪球。

约定的那一天到了，我又来到山丘上。我确信我的这个做法是相当高明的，因为那群年幼的助手们一定会在他们的同学面前炫耀这桩赚钱的交易，如果有谁不相信圣甲虫的粪球会这么值钱，他们会得意扬扬地拿出订金给他看。事情果然如此，这一次来这里的人数比上次多了很多。他们见我来了赶忙跑过来，但并没有胜利后的欢呼雀跃。我想一定是我交代的事情失败了，事情果然如我所料。他们说每天放学后都去寻找，可没找到我要的那种有虫卵的粪球。虽然他们也曾连虫带球一起捡到过，可那些粪球都仅仅是食物而已，里面并没有虫卵。听完他们的话，我又给他们鼓了鼓劲儿，约定下个星期四再在这里碰面，还是没有成功。大多数人都放弃寻找了，只有少数几个人坚持着。可他们在我的鼓励下做了最后一次努力后，仍然一无所获。最后，我给那几个坚持到底的人一些报酬，合作也就此以失败告终了。看来我只能独自进行那看似简单、实则困难的研究了。

即使在很多年以后的今天，我到适当的地点寻找，在最佳的时机观察，也没能得到一个确切而合乎逻辑的结果。于是，我只能把那些零零散散的观察结果综合起来，经过类比得出结论，填补这一方面的空白。现在，我就将我的观察以及从笼子里的圆裸胸金龟、蜣螂和牛粪屎蜣螂等食粪虫身上得到的资料结合起来，归纳总结一下。用于孕育虫卵的粪球，是不能在众目睽睽之下制造，也不能在吵吵嚷嚷的开采工地上进行的。这是一件需要耐心制作的艺术品，要在远离人群的地方专心致志、一丝不苟地进行。雌虫进入居室以后，全盘思考了一番后就开始动手做了。圣甲虫母亲在沙地里给自己挖了一个10厘米×20厘米的宽敞大厅，由一个小回廊通到外面。它把精心挑选的材料制作成球形后运到里面。当然，要往返多次才能运完，因为它运来的粪球都堆积在了大厅里，那庞大的体积是不可能一次通过入口进来的，也不可能是一次就堆积起来的。我记得我在拜访一只西班牙蜣螂时发现，它的洞里放着一个橙子那么大的粪球，可是它大厅的长廊很狭窄，仅能穿过一根手指。蜣螂既不滚动粪球，也不千里迢迢地搬运食物，而是在粪便底下直接挖一个洞，然后后退着把材料一抱一抱地拖到洞里。在粪便的下面安家，寻找起食物来既方便、容易，工作起来也更安全，也正因为如此，蜣螂养成了奢华的饮食习惯。而那些靠出苦力搬运粪球的食粪虫就从不挑三拣四。不过，圣甲虫往返两三次以后就会囤积起一大笔可观的财富，足以让西班牙蜣螂羡慕不已了。

但是，那些运进来的粪料只是胡乱地堆砌在一起，还未经过加工。现在首先要做的就是筛选材料：把最精细的放在里层，给幼虫做食物；最粗糙的则放在外层，保护整个粪球。然后，在卵虫生存的中央居室四周，将材料按照粗细程度和营养价值的高低，由内到外、由优到劣，一层层地裹好，每层材料都要十分稳固，并且要和后一层紧密贴合。最后，使最外层的纤维筋粘合镶嵌在粪球表层上，起到保护壳的作用。圣甲虫的居室内塞满了食物，几乎没有活动空间，并且昏黑一片，这样，对于动作迟滞、身体僵硬的圣甲虫来说，它如何完成那件艺术品的制作呢？一想到它们那细致的工作是由粗大的工具完成的，以及它那可以掘土甚至可以凿开岩石的多角的爪，我就不禁联想到大象做起绣花的工作。如果有人看到过这些母亲们工作的过程，就请描述一下这种堪称奇迹的技艺吧，而我则不想看到，况且也没有可能见证艺匠的制作过程。我们只能对它那伟大的作品进行一番描述！

通常情况下，产卵的粪球差不多有一个中等苹果那么大，粪球的中心有一个直径约10厘米的椭圆形小洞，卵就直直地安放在那里，十分牢固，直抵洞底。虫卵是圆柱形的，两端浑圆，颜色白中泛黄，跟麦粒差不多大，只是稍短一些。小洞的洞壁四周涂抹了一层微绿而带棕色的物质，散发着光泽，呈半流体，那是真真正正的粪糊，是幼虫出生后的第一份粮食。这种精细的食品，难道是由粪便的精华做成的？从食物的形状上看，不是这样的，而是由圣甲虫母亲消化后的某种物质制成的。看来食粪虫也有鸽子身上的柔情，鸽子会先在嗉囊里把麦粒软化，直至形成一种乳状的黏糊物质后，再喂给雏鸽吃。食粪虫则先把精选的食物吃掉，经消化形成精细的糊粥，然后吐出来涂抹在安置虫卵的小洞的洞壁上。幼虫孵化后就可以找到易消化的食物，等它的胃功能强健起来后，就能逐渐向外层越来越粗糙的食物进攻了。紧挨半流体涂层的是一层不含任何纤维屑食物的精华层，构造得匀称而坚实。再往外是有着许多植物茎的粗糙食物层，最外层的材质最粗糙，但是经过紧压、粘结后形成了十分坚硬的壳。

这清楚地反映出了圣甲虫从幼虫到成虫的饮食变化规律。当无比娇弱的小幼虫破卵而出时，就开始食用住所四周墙壁上的精细浆液。浆液虽然不多，但都是高营养物质，能使它的身体和肠胃功能迅速强壮起来。度过了新生期，圣甲虫幼虫仿佛断奶的婴儿一样，不再有糊粥的供给了，而是开始食用介于最精细的浆液和最粗糙的食物层之间的食物。这一层很厚，食物较多，能把小幼虫养得结结实实的。不过幼虫长结实以后，食物又变了，变成了那类似大麦面包一样的杂干草尖的天然粪便。幼虫的食物是多种多样的，除了它在整个生长期食用的以外，包裹在它身外的一层墙壁也是它的食物。它们的身体长大后，住所也随之扩大了，因为住所墙壁上的物质就是它们的食物。最初那个小洞上的墙壁很厚，现在已经变成了大房间，墙壁只有几毫米厚；在房屋变大的过程中，室内的居民也分别完成了幼虫、蛹或金龟子的蜕变。总之，粪球是一个坚实的硬壳，圣甲虫的生命演化过程就是在它的遮掩下秘密进行的。

可用的材料用完了，我无法继续写下去了，至于圣甲虫的身份之谜，仍旧停留在虫卵上。虽然我没亲眼见过幼虫，但是有的作家见过它们，并且在作品中做了描述。至于那些已经完成变态而仍旧生活在粪球里的、尚未开始搬运和挖掘等工作的圣甲虫，我更没见过。可它正是我梦寐以求的。我希望能对这样一只食粪虫进行观察：它住在粪球房里，刚刚完成变态，尚未开始任何工作。为什么我会有这种想法呢？我将在后面说出原因。

这些昆虫的爪子上有趾肢节以及人们所说的跗骨，它的跗骨是由一些类似人类手指指骨的精细零件组成的，最末端是带钩的指甲。每个爪有一个趾肢节。高级的鞘翅目昆虫，尤其是食粪虫，趾肢节至少有5根指骨或者关节。可是奇怪的是，金龟子完全例外，它的前爪没有跗骨，其他两对爪子都有5个关节。由此可见，金龟子是个腿脚不全的残疾，它们的前爪不像其他昆虫那样有趾肢节。这种情况在食粪虫家族中的宽胸蜣螂和粪屎蜣螂身上也能看到。对这些奇怪的情况，昆虫学上早有记载，却至今也没有合理的解释。这种昆虫的腿脚是从出生时就残疾吗？是生下来时前爪就没有趾肢节，还是刚开始工作时就意外地折断了呢？

大多数人会认为圣甲虫的肢体残缺是由劳动造成的，因为它总是在砾石中或者在布满粗纤维的粪堆里不停地搜呀、挖呀、耙呀、剁呀，而这会大大损伤它那娇嫩的跗骨。更重要的是，当它倒退着滚动粪球时，总是脑袋朝下，用前爪支撑地面。它那脆弱的趾肢节像线头那么细，却不断地碰触粗硬的大地，会造成什么后果呢？这些趾肢节一无用处，完全是负担，终究会在千百次的磨砺中消失掉，或者被压碎，或者被拔掉，或者被磨光。唉！就连人类在使用沉重的工具或搬运重物时也常有人落下残疾啊。也许金龟子就是在搬运粪球时残疾的，因为对它而言，粪球是一个沉重的负荷。断指或许就是它勤奋劳动的光荣证明吧！

但是有人听到这儿会产生疑问：如果这断残确实是因为意外和艰苦的劳动造成的，那也应当是偶然情况而不该是通例啊！一个工人或若干个工人的手被机器轧断了是常事儿，但不可能所有工人的手都断掉啊！即使对于金龟子这个粪球搬运工来说失去趾肢节是平常事，十分稀松平常，也至少会有几只幸运的或者动作灵敏的圣甲虫的前爪保留着跗骨。现在我们来看看事实到底是怎样的。我观察过许多生活在法国的金龟子，比如：普罗旺斯的圣甲虫，生活在距海较远，经常在塞特、帕拉瓦和里昂海湾的沙滩上活动的一种带斑点金龟，以及比前两种更常见，一直分布到达罗讷河谷，至少在里昂一带随处可见的宽颈金龟。我还观察过一种来自君士坦丁城郊的非洲瘢痕金龟。结果无一例外，这四种金龟子的前爪都没有跗骨，至少我观察到的情况是这样的。所以，金龟子的断指是与生俱来的，这是它的身体特征，而不是意外造成的。还有另一个理由可以进一步证明这个说法：如果昆虫前爪趾肢节是在操作失误而引起的工伤事故中断掉的，那么，那些工作更艰苦的——干挖掘工作的昆虫，更应该没有跗骨。因为跗骨是百无一用的附属品，当前爪需要出强力的时候，它就是累赘。

\newpage

\chapter{捕食吉丁的节腹泥蜂}

\begin{figure}[htbp]
    \centering
    \includegraphics[width=0.8\textwidth]{images/00296.jpeg}
\end{figure}

有些书为人们描绘了一个想象之外的世界。因为每个人的思维方式不同，所以对于有的人来说，这些作品具有划时代的意义。它为人们展现了一个全新的世界，一个需要人们充分运用聪明才智去开拓的世界；它们是引燃火炉的火源，如果没有这些星星之火，火炉中的木柴将永远无法燃烧起来。这些作品是有启发意义的，它改变了我们的思想，是创造新时代的起点，可我们往往是在无意中读到的。在不经意间，在毫不知情的情况下读到的这些文字，改变了我们的未来，指引着我们走上了命运之路。

冬天的某个夜晚，火炉散发出暖暖的热气，全家人都进入了梦乡，只有我独坐在炉边读书，把已无明日之粮的烦恼抛到了脑后。那是一个物理教师的烦恼，他取得了一个又一个学位证书，在20多年的从业生涯中取得了有目共睹的成绩，可他全家每月的生计开支只有1 600法郎，都不如一个富户的喂马工的工钱。这个时代对待教育事业的态度是多么可耻啊，1分钱都不肯多给。因为我总是独自研究，所以依据政府的规定，我成了非正式人员。当我埋头苦读时，偶然间翻到了一本昆虫学小册子，那是一本我自己都不知道从何而来的书，而它使我忘却了当教师的贫穷和苦闷。

这是一本关于捕食吉丁的膜翅目昆虫的书，出自当时的昆虫学祖师——可敬的学者列翁·杜福尔之手。其实，我并不是在这个时候喜欢上昆虫的。从童年时代起，我就对鞘翅目昆虫、蜂和蝶产生了浓厚的兴趣。我记得刚刚开始记事儿的时候，我就曾望着步甲华美的鞘翅和金凤蝶漂亮的翅膀出神。火炉里已备好木柴，缺少的是使木柴燃烧起来的火星。无意中读到的列翁·杜福尔的作品，就是这星星之火。

得到了这个新的启迪以后，我的思想就像是得到了上天的启示和指引一样。这门科学，不仅包括把美丽的鞘翅目昆虫放在软木盒里给它们命名、分类，还有很多更深层次的东西：深入探究昆虫的身体结构，研究它们的习性特征。我阅读着这本昆虫学领域的杰作，顿觉激情澎湃，我相信任何一个执着追求的人都会像我一样感觉幸运。

得到了这本书的引导后不久，我写下了一部补充列翁·杜福尔这本作品的书，这也是我的第一部昆虫学著作，荣获了法兰西学院的荣誉，被授予实验生理学奖。但是对我来说，令我最欣喜的是不久收到的一封信。这封信中对我的称赞令我深受鼓舞。这是令人敬爱的大师从朗德腹地发来的信，向我表达了他热烈的祝贺和喜悦之情，并极力鼓励我在这条路上继续走下去。直到今日，每当我想起这件事时，我昏花的老眼中还会涌出那圣洁的饱含感激之情的泪水。啊，那些美好的、对未来充满想象和企盼的日子，顿时发生了翻天覆地的变化！

我准备把促使我开始从事自己研究的那篇文章摘录下来，读者应该不会反对吧，最重要的是，了解这个摘要有助于理解我后面将要谈到的事。现在，我们就来听听大师是怎么说的吧，只是我进行了适当的精减。

在整个昆虫学史上，像我将要谈到的这样奇怪、新奇的事，我还是第一次遇见。我要谈的是一种节腹泥蜂，它总是处心积虑地用吉丁科昆虫来喂养它的子女。我的朋友，请允许我把这个膜翅目昆虫给我留下的深刻印象告诉您吧。

1839年7月，我的一个乡下朋友给我寄来了一种双面吉丁，共两只。那时我还从来没有收藏过这种昆虫，他告诉我这是一种泥蜂在运送这些美丽的鞘翅目昆虫时，把第一只吉丁扔在了他的衣服上。不一会儿，另一只泥蜂又把另一只吉丁扔在了地上。

1840年7月，身为医生的我到这位朋友家里出诊，我向他提起他送给我的昆虫并询问他当时的情况。同一个季节，相同的地点，使我产生了亲自去抓昆虫的愿望。可是那天天气有些阴冷，膜翅目昆虫通常不会出来活动。尽管如此，我们还是到花园里进行观察，因为根本找不到这种昆虫，我便开始在地上寻找这种膜翅目昆虫的窝，因为它素来以掘地见长。突然，我注意到了一个刚刚翻动过的小沙堆，就像小鼹鼠丘一样。我扒开小沙堆，发现下面有一个地下通道的口。我用铲子小心翼翼地挖开地，不一会儿就看到了零零散散的吉丁鞘翅正闪耀着光芒。不一会儿，我看到的不再是一些散落的鞘翅，而是一些残骸，是三四只吉丁共同展现着它们的"金子"和"绿宝石"。我简直不敢相信自己的眼睛了。然而，那个令我兴奋不已的发现才刚刚开始！

在我挖掘到其他地方的时候，一只膜翅目昆虫从地下飞了出来，落在我的手里。这是一种专门捕猎吉丁的昆虫，它正试图从它的猎物窝里逃走。然而我一眼就认出了这只掘地虫是节腹泥蜂，因为我和它已经是老朋友了，一生之中我曾在西班牙和圣塞维郊区找到过它不下两百次。

我的野心远远没有满足。在认识了掠夺者和被掠夺者之后，我还想找到幼虫，因为它才是这些奢华食物的最终享用者。挖完了第一个埋藏着吉丁的窝以后，我又立即挖起了新窝，更加仔细地寻找起来，结果我幸运地得到了两只幼虫，给我的挖掘行动画上了圆满的句号。在不足3个小时的时间里，我就捣毁了节腹泥蜂的3个窝，得到的完整的吉丁就有15只，还有很多很多的断臂残骸。我根据这种概率粗略算了一下，花园里应该还有25个窝，但我认为其实远不止这些，如果真是这样的话，地下可就埋藏着一大群吉丁。我苦苦思索着，为什么在这里仅仅在大蒜花上就能捉到60只节腹泥蜂呢？也许这些泥蜂的窝就在不远处，而且它们的饮食也一定非常奢华。因此我大胆地猜想，在地下的一个很小的范围内，会有几千只双面吉丁，这种猜测也是极有可能的；可是30多年来，我研究了我那个区域内的很多昆虫，在乡下却没有找到过一只吉丁。

唯有大约20年前的那一次，我在一个老橡树洞里见到了这种腹部长有鞘翅的昆虫。它的鞘翅给我带来了一线希望，它说明双面吉丁的幼虫很可能就生活在橡木里，如此一来，在一个橡树林里生活着大量的这种鞘翅目昆虫也是顺理成章的事了。由于在这里的黏土丘陵上不常见到捕食吉丁的节腹泥蜂，至少远不如在长有海松的沙地平原上更常见到它们，所以我特别想弄清楚一件事：对于这种膜翅目昆虫来说，生活在松树林里的也像生活在橡树林里的一样奢侈，用吉丁喂养它的幼虫吗？我有充足的理由证明它不是这样的，可是下面这个新发现会令您震惊：我们的节腹泥蜂在捕猎吉丁科昆虫的时候，触觉异常的敏锐。

那么我们现在就到松树林里去了解这个新发现，体验它带给我们的乐趣吧。我们检查的目的地是海松林中的园林花园。在那里，我很快就找到了节腹泥蜂的窝；这些窝全部建在主甬道上，因为那里的地面经过很多人的踩踏后变得更平整坚实了，这种掘地的膜翅目昆虫就喜欢在这样的地点建造地下住所。我大约一共检测了20个洞，累得汗流浃背，因为要找到这些窝是很辛苦的，必须挖到1法尺深的地方，在那里才能得到窝里贮存的食物。为了不损坏它们的窝，我先把一根稻草秸插进节腹泥蜂的巷道中，作为指引标杆，然后用一个方形的挖掘线在距离洞眼或标杆七八法寸处将这个地方围起来。挖地要使用园艺铲，使中间的土块跟四周的土完全隔开，这样才能把整块土挖起来，然后翻倒在地，再小心地把它捣碎。我就是这样做的。

我的朋友，我就是使用这种别具一格的探测方法，把华美的吉丁一个个摊放在我渴求的目光之下的，您或许会像我一样心花怒放吧。每当我把坑道翻个底朝天，看到里面的宝物在阳光下更加炫美多姿时，或者每当我发现那些咬住猎物的幼虫壳上镶嵌的"金子"和"蓝宝石"时，我都会不由自主地欢呼雀跃起来。我是个讲究实践的昆虫学家，唉！在之前的三四十年中，我从未见过如此动人的画面，从未有过一个日子像这样值得庆贺。可惜您没有亲眼见到，否则一定会更加高兴的。我看到这些光彩夺目的鞘翅目昆虫，看到这些吉丁仍清晰可辨，再看到节腹泥蜂把它们埋藏、存储起来，完全被它们惊人的智慧折服了，不禁赞叹起来。请相信，我挖出来的400多只昆虫全部都是吉丁科的。如此灵动的膜翅目昆虫我是不会弄错的。我们从一只如此娇小的昆虫的智慧中得到了多少收获啊！节腹泥蜂在博物学里的特殊贡献，会得到拉特雷依多高的赞誉啊！

现在，我们来了解一下节腹泥蜂是如何建造住所和捕获食物的。我在上面提到过，它的建筑地点总是选择在踩踏得紧密坚实的土地上，还需要特别强调的是，此处必须干燥、能晒到烈日。选择这样的地方是十分明智的，或者您也可以说它是出于本能，而本能得益于经验。当然，与坚实的土地相比，在一块松散的土地或者一块纯沙质的土地上挖掘要更容易些；但是在这样的地下，怎能开出一个始终敞开着的便于出入的洞口呢？怎能挖出一条两壁不会随时倒塌、不会因雨水冲刷而变形或堵塞的巷道呢？由此可见，它选择这样的地点是十分正确而且经过深思熟虑的。

我们的膜翅目掘地虫在挖掘巷道时使用的工具是它的大颚和前脚跗节，它的跗节上长着硬刺，好像耙子一样。洞穴不能太小，如果仅能容下矿工的身体是不够的，必须保证能把体形更大的猎物运进去。它们竟有如此远虑，真是令人震惊。节腹泥蜂钻进土里挖掘起来，把挖出来的土送到外面堆起来，这正是我前面所说的那个像小鼹鼠丘的土堆。巷道不是直上直下的，因为垂直的巷道在遇到大风或其他状况时，随时可能会被填满。在距洞口不远的地方，巷道有个长七八英寸的弯道。才智双全的母亲们就把后代安放在巷道尽头处。这是5个彼此相隔而独立的蜂房，排成了半圆形，蜂房的形状和大小都类似橄榄，里面光滑而稳固。每个蜂房内可存放3只吉丁，这是每只幼虫每日所需的食物。母亲在这3只吉丁中产卵，然后用土封住巷道口，让居室与外部彻底隔离，同时保证幼虫在整个孵化期内都有食物吃。

捕食吉丁的节腹泥蜂堪称是一个有勇有谋、灵活敏捷的猎手。它捉来的吉丁又干净又新鲜，甚至会让人们误以为这些鞘翅目昆虫是刚刚才在木质巷道里完成变态的。可是，以花蜜为生的节腹泥蜂，要具有多么神奇的本能，才能不辞辛苦地为它永远无法相见的孩子提供肉食，才能够飞到一棵棵陌生的树上，从树干深处搜捕那注定成为它猎物的昆虫啊？它需要具有多么超乎想象的动物学触觉，才能够在一个庞大的类别中捕猎到自己所需的某几种昆虫，而这几种昆虫又在大小、形状、颜色方面各不相同啊？我的朋友，您看，这些昆虫几乎没有相似之处：双棉芽吉丁体形修长、身色暗淡；八棉芽吉丁呈椭圆形，或蓝或绿的体色上长着两朵美丽的大黄斑；碎点吉丁的体形比双棉芽吉丁大三四倍，颜色呈金光闪闪的蓝绿金属色。

我们的吉丁杀手还有一个极其特殊的操作技巧。无论是被埋葬的吉丁还是从捕猎者手中抢来的吉丁，都已经没有了生命的迹象，当然，它们都死了。可我惊奇地发现，这些吉丁尸体在任何时候都保持着鲜亮的色泽，就连爪、触角、触须以及连接身体各部分的薄膜都十分柔软，可以弯曲。它们的身体丝毫没有损坏，也没有任何显著的伤痕。起初，我以为吉丁保持新鲜的原因有两个：埋在洞里的是因为地下比较阴凉，而且隔离了空气和光线；从捕猎者手里抢下的则是因为刚死不久！

但是请您留意一下，我在做实验时，会先把吉丁分别放进圆锥形的纸袋里，然后通常是摆弄了36个小时之后才用大头钉把它们钉起来的。您看，在炎热干燥的7月份，这些吉丁的关节却一直都能自如地弯曲。不仅如此，在这之后，我解剖了几只吉丁，发现它们的内脏都完好无损，我的解剖刀就像是插在了活的昆虫内脏上一样。我从以前的实验中得到的结论是：即使是一只相同大小的鞘翅目昆虫，在夏季死了12个小时后，它体内的器官也会或者干化或者腐烂，是无法辨别出它的形状和结构的。可是，我在被节腹泥蜂杀死的吉丁中却发现了例外，它为什么能一个星期或者两个星期都不变质发臭呢？

一只昆虫已经死了几个星期，成了一块野味，可是在最炎热的夏天，它的尸体却不会干化、腐烂，就像活着时一样新鲜。为了弄清吉丁捕猎者是如何保鲜肉类的，这位聪明的描述吉丁捕猎者的昆虫学家假想它使用了一种防腐液，这种防腐液就像人类为使解剖下来的肌体保持新鲜而使用的化学药品一样。然而即便是果真如此，也只能是这种膜翅目昆虫亲自把毒汁注入猎物体内的。它用蜇针把一小滴毒汁注入吉丁的体内，就像腌肉的盐水或者防腐液一样，把幼虫要吃的肉保存下来。可是膜翅目昆虫保鲜食物的方法比我们高明多了！我们要用盐浸泡，用烟熏，然后把食物放到密封的铁盒里，是的，这样之后食物仍然可以吃，可是远不如食物新鲜时的质量好。浸在油里的沙丁鱼罐头、荷兰的烟熏鲱鱼、盐腌日晒的鳕鱼干，这些怎么比得上刚送到厨房里活蹦乱跳的鱼呢？至于肉类，那就更糟糕了。一块鲜肉，除非经过腌制或熏干，否则哪怕只存放了很短的时间也不能食用了。今天，在经过了各种各样的无数次的失败尝试后，人们花费巨资来装备特殊船只，船上装备着强大的冷冻机器，用强冷的方法使南美潘帕斯草原上宰杀的牛羊肉免于腐烂。节腹泥蜂的方法却如此快捷、有效又不需要花钱，比我们人类高明多少倍啊！它那超群绝伦的技巧给了人类多少有益的启示啊！它用肉眼几乎看不见的一小滴毒汁，就能立即使它的猎物免于腐败。我说的是什么呀？不会腐败？远远不只是这样！它不仅使它的野味不腐化，还能使野味的关节弯曲自如，内外器官保持着像有生命时一样完整新鲜。总之，它的猎物除了一动不动像尸体外，其他的跟活着没什么区别。

这是杜福尔在面对吉丁尸体不会腐烂的奇迹时产生的想象。一种比人类科学所能生产出来的防腐液还要强千百倍的东西，似乎才能揭开这个奥秘。这位大师是精明者中的精明者，精通解剖学，他凭借放大镜和解剖刀仔细观察了整个昆虫系，任何一个细节都没有忽略。总之，在他手下，任何昆虫的组织结构都是没有秘密可言的。可是面对吉丁尸体，他也迷惑不解了，除了用一个假设来做些从表面看合情合理的解释外，再也无能为力了。请允许我反复强调一下昆虫的本能和学者的理性，以便及时揭示出昆虫独具的无与伦比的优越性。

关于以吉丁为食的节腹泥蜂的历史，我想再补充几句。专门研究它的历史学家告诉我们，这种膜翅目昆虫在朗德是随处可见的，在沃克吕兹省却难得一见。至于我，只是曾在深秋时节的阿维尼翁郊区或奥朗日和卡班特拉郊区见过几次，而且我都是在带刺茎的菊科植物的头状花序上看到的，它们也总是形单影只。在卡班特拉，海边松软的沙地有利于掘地膜翅目昆虫建筑住所，我不仅亲自挖出了很多像杜福尔描述的那种昆虫，还找到了它们的几处住所。根据茧的形状、食物以及在近处看到的膜翅目昆虫，我断定这个住所的主人就是吉丁的捕猎者。这些住所安在一种被当地称为"萨弗尔"的十分易碎的砂岩中，住所内，鞘翅目昆虫的残肢断臂随处可见：折断的鞘翅、掏空的前胸、整只腿爪，一眼就能看出来是吉丁的。而节腹泥蜂幼虫饱餐后剩下的残羹冷炙，全部成了另一种昆虫的囊中物，这就是对生吉丁，也是吉丁的一种。所以，从法国的西部到东部，从朗德省到沃克吕兹省，所有的节腹泥蜂都以吉丁为食，它们的嗜好丝毫不受地理位置的影响。它在海边沙丘的海松林中捕捉的是吉丁，在普罗旺斯橄榄树中捕猎的还是吉丁。唯一不同的是，它会因地点、天气和植被的不同而改变捕猎吉丁的种类，这些因素促使昆虫家族发生了很多变异，但是节腹泥蜂偏爱的食物自始至终都是吉丁。为什么这么奇怪呢？这也是我想加以论证的。

\newpage

\chapter{栎棘节腹泥蜂}

\begin{figure}[htbp]
    \centering
    
\includegraphics[width=0.8\textwidth]{images/00303.jpeg}
\end{figure}

吉丁捕猎者的卓越功勋占据我的脑海，我希望能亲眼看看节腹泥蜂是如何工作的，在热切的期盼下，这个机会终于来了。当然，我看到的不是杜福尔称赞的那种膜翅目昆虫，它不是以奢华的食物为餐，它吃剩下的食物残骸从地下挖出来以后，也不会像矿工在金矿内砸烂的金块碎粉一样。我看到的昆虫跟节腹泥蜂属同一种类，但是体形硕大，专门捕食个头小的猎物，它就是栎棘节腹泥蜂，也称大节腹泥蜂。在节腹泥蜂中，尤以它的个子最高、身材最壮。

9月下旬，膜翅目掘地虫开始建造住所，并准备把为幼虫捕猎的食物埋藏到住所里。它们选择住宅地点的标准十分苛刻，而且总会发生某种神秘莫测的变化，不同种类昆虫的住宅法则各不相同，而同一种类昆虫的住宅法则却是恒久不变的。杜福尔所说的节腹泥蜂，会选择在平坦的小路上建造住所，同时土质要是久经踩踏的、坚硬而紧实的，这样大雨倾盆时，住所也不会倒塌、变形或堵塞巷道。我们见到的节腹泥蜂则完全不同，它会把房子建在垂直的地方，只要在施工时进行一些小小的变动，就能规避大部分有可能损毁巷道的危险了；这样一来，它的住宅建筑用地就比较容易选择了，随便哪个地方都可以，无论是略带黏土的松软地，还是柔软易碎的沙石中，在这些地点建造，它的挖掘工作也简单很多。它们对建筑用地的唯一要求就是干燥，而且每天大部分时间都要有阳光照射。所以我们的膜翅目昆虫的家通常会安在道路旁的陡峭边坡处，或者经过雨水冲刷后的柔软沙地形成的沟壑的侧面。在卡班特拉附近有一个叫"凹路"的地方，那里有很多这样的地势，我就是在那儿观察到许许多多栎棘节腹泥蜂的，还搜集到了很多关于它们历史的事实资料。

对于栎棘节腹泥蜂来说，仅仅选择垂直的地方安家是不够的，因为面对深秋时节的一场场如期而至的秋雨，它们必须采取有效的措施防御雨水的侵袭。某些像房檐一样凸出的硬砂岩片在土里形成了一个拳头大小的洞，它们以此为避难所，在这样的遮雨条件下、在这样的洞底修筑巷道，同时，它的房屋也增添了一个前厅。虽然这些昆虫没有群居特性，但一小部分喜欢聚集在一起，通常是10只左右成一组，至少我在它们家中看到的情况是这样的。它们大多数的洞口都相距甚远，有的却紧挨在一起。

在天空晴朗、阳光灿烂的时候，去看看这些勤劳的矿工工作是再合适不过了。在洞穴里，它们有的用大颚耐心地挑拣着几粒砾石，然后把这些石子扔到洞外；有的用跗骨上锋利的耙刮着走廊的墙壁，然后倒退着把刮下来的一堆泥屑扫到洞外，泥屑就像溪水一样，从陡坡的侧面缓缓流淌下来。我就是凭借这些从巷道排出的细流才发现了节腹泥蜂的踪迹，进而找到它们的住所的。另外一些昆虫，也许太累了，也许是已经完成了某项艰苦的工作，此刻正在住所的天然遮雨棚下休息，有的擦拭着触角和翅膀；有的则在洞口纹丝不动，只探出黄黑相间的方形大脸；也有的正低声嗡嗡叫着，在胭脂虫栎附近的灌木丛上飞来飞去，引得那只一直守候在它们住所旁的雄蜂很快就尾随而来了，于是一段美好的姻缘即将开始。可是这个时候通常有另一只雄蜂来横加破坏，试图取代这个幸运者而成为雌蜂的伴侣。低沉的嗡嗡声变得越来越激烈了，两只雄蜂厮打翻滚起来，到处尘土飞扬，直到其中一只甘拜下风，战争才停止。雌蜂在不远处平心静气地观看着比赛，最后，它选择了胜利的一方，然后它们一起比翼双飞了，飞到遥远的灌木丛上去过清净的日子。雄蜂所做的就是这些，它比雌蜂的个子小一半，可数目几乎与雌蜂一样多。它们在雌蜂的住所旁徘徊，却从不加入到挖洞的艰苦劳动中来，至于为供应蜂房的粮食而捕猎这样的更为艰苦一些的工作，它就更敬而远之了。

没几天工夫，它们就挖好了一个巷道，尤其是那些前一年用过的巷道，稍微修理一下就可以继续使用。在我所知的范围内，其他种的节腹泥蜂没有一处可以当作遗产一代一代传下去的固定住所，它们是名副其实的流浪吉卜赛人，无论流浪到哪儿，只要土壤合适，就在那里定居。可是栎棘节腹泥蜂十分眷恋旧居，连那个曾被前人用于遮雨的砂岩片，都被它们充分利用了；它到祖先挖掘过的沙基上挖掘，对前人的建筑物进行一番改造更新，只是它们的巷道更隐秘了，向更深处延伸下去，使得我们如想查看就要费一些力气。巷道的直径很大，能放进一个大拇指，昆虫即使抱着它的猎物也能在里面活动自如。我们会在后面看到它捕捉猎物的过程。巷道先是水平走向，但是在一两分米深处时突然拐了一道弯，略为倾斜地一会儿朝这个方向，一会儿朝另一个方向延伸下去。整个巷道除了水平部分和拐弯处外，其余的似乎是依据挖掘土壤时的难易程度而变化的，从巷道最深处变化不定、蜿蜒曲折的走向就可以看出这一点。我检测的洞共有半米长，巷道尽头是蜂房，数量不多，每间蜂房里备有五六只鞘翅目昆虫作为食物。不过我们还是暂且把这些工程建造上的细节放下吧，先去看看那些更令我们振奋不已的事情。

这种节腹泥蜂总是捕猎身形硕大的小眼方喙象来喂养幼虫，这是一种象虫科昆虫。掠夺者用腿将沉重的猎物环抱，肚子贴着肚子，头靠着头，往洞口飞去。在洞口不远处，它会笨拙地落下来，步行过去。这时，在垂直的或者坡度很陡的平面上，膜翅目昆虫用大颚十分艰辛地拖着猎物，栽了一个又一个跟头，动不动就连同它的猎物一起滚到斜坡底下。但不知疲倦的母亲并没有灰心丧气，满身泥土的它终于带着一刻也不曾离手的战利品走进了洞里。栎棘节腹泥蜂抱着这么重的东西行走很艰难，飞起来却很容易。它高超的飞行能力太令人钦佩了，原本就很粗壮的它竟然能抱着一个跟它差不多大小但比它重很多的猎物飞行。我曾怀着强烈的好奇心，分别测量了栎棘节腹泥蜂和它的猎物的重量，结果我发现，前者一般重150毫克，而后者的平均重量为250毫克，几乎是它的两倍。

这些数字有力地证明了这个捕猎者是多么的强壮有力。所以当我在好奇心的驱使下靠它太近，使它受到惊吓时，它决定带着它那宝贵的战利品一起逃走。那时，它敏捷而利落地用腿抱起猎物飞走了，飞到了我看不见的高处，这令我赞叹不已。当然，它并不是每次都想逃。我曾用一根麦秸把它弄倒，在这个既不会伤害它又能让它松开猎物的妙计下，我迅速地把猎物抢了过来。节腹泥蜂遭到抢劫后，四处寻觅，钻进洞里，很快又出来，然后飞走重新捕猎。不到10分钟的工夫，这个灵敏的掠夺者又得到了一个战利品。对它们进行劫持后，我总是擅自把这些猎物据为己有。我曾用这样的办法连续抢劫同一只泥蜂8次，而它次次都会重新开始那徒劳无功的远征。它的矢志不移和顽强的毅力使我失去了耐心，于是，我把第9次抢到的猎物最终归还了它。

通过劫持以及私闯蜂巢行窃的方法，我得到了近100只象虫科昆虫。虽然读过杜福尔关于捕食吉丁的节腹泥蜂的习性这篇文章后，我已经完全明白了其中的缘由，但是当我看到亲手搜集来的材料时，还是十分惊讶。如果说吉丁的捕猎者可以捕捉同一类昆虫的任意一种，那么这种栎棘节腹泥蜂捕捉的则自始至终都是小眼方喙象。我在清点我的战利品时，只发现了一个例外，看到了一只交替方喙象，仅此一例。尽管它与小眼方喙象是同类，但我经常探查节腹泥蜂的住所，再也没有见过它。在接下来的研究中，我发现了白色甜菜象，这是第二个例外，也是最后一个例外。这种猎物味道更鲜美、更甜润，这是不是它只吃一种昆虫的原因呢？幼虫是否因为吃惯了这种野味的汁，认为它最合胃口，而从别的昆虫那儿根本品尝不到呢？我不这么认为。如果说杜福尔的节腹泥蜂能任意捕猎各种吉丁，是因为所有吉丁的味道都一样，那么象虫科昆虫也应当如此啊，它们的营养特征和味道也应当是相同的。所以，它们做出这种奇怪的选择，原因在于猎物体积的大小决定了它们需要付出多少辛劳。我们的这种节腹泥蜂在同类昆虫中堪称巨人，而它唯独喜欢捕食小眼方喙象，是因为这种象虫是我们这个地区内体形最大的昆虫，也许也是最常见的。但是假如这种节腹泥蜂捕不到自己最爱的猎物，不得不捕猎别的虫子时，就不在乎猎物的大小了，我们见到的两个例外就是证明。

另外，在昆虫界中，绝对不是只有这一种节腹泥蜂是靠捕食大吻管科昆虫——象虫来生存，许多别种节腹泥蜂也会根据各自的体形、力气和捕猎可能性等条件的不同，捕食相应的象虫科昆虫，因为这些虫的类别、种别、形状、大小，千差万别。

我们知道的沙地节腹泥蜂，就是用类似的食物来喂养它的幼虫。我就曾在它的洞穴里找到过直条根瘤象（豌豆叶象）、长腿根瘤象、细长短喙象、耳象。大耳节腹泥蜂的猎物有草莓耳象和带刺叶象。铁色节腹泥蜂的餐桌上经常有鼠灰色叶象、带刺叶象、直条根瘤象、槭虎象。槭虎象的身色有时呈现出那种美丽的金属蓝，但大多数时候是闪烁着金光的金铜色，它会把葡萄叶卷成雪茄状。有时，我会在一间蜂房里看到7只这样闪闪发光的昆虫，这时，地下洞穴里盛宴的奢华程度，简直可以与吉丁捕猎者埋藏起来的浑身披金戴银的吉丁相媲美了。其他种类的节腹泥蜂，特别是那些最弱小的种类，偏爱吃小一些的野味，并靠增加食物数量来弥补食物体积上的不足。比如四带节腹泥蜂，它在每个蜂房里堆放的圆腹梨象的数目多达30只；当然，如果它们有幸遇见直条根瘤象、鼠灰色叶象这类体形较大的象虫科昆虫，也不会熟视无睹。大唇节腹泥蜂也吃小体积的食物。此外，朱尔节腹泥蜂是我们这个地区最小的节腹泥蜂，它的主要猎物是最小的象甲科昆虫——圆腹梨象和谷仓豆象，这些野味与捕猎者娇小的身材十分相配。除了这些食物之外，最后我还要补充一点，某些节腹泥蜂，比如缀锦节腹泥蜂则有自己的一套饮食标准，它捕捉某些膜翅目昆虫来饲养它的幼虫，它的这种偏爱超出了我们讨论的范围，暂且不去谈论了。

可见，在这8种以鞘翅目昆虫为食的节腹泥蜂中，有7种捕食象虫，1种捕食吉丁。这些膜翅目昆虫为什么把捕猎的范围局限得这么小呢？它又为什么非以这种食物为食呢？从外形来看，吉丁和象虫毫无相似点，难道是因为它们在特征上存在着某些相似之处，使得它们都成为同类的膜翅目幼虫的食物吗？在其他的某些猎物之间，一定存在味道、营养质量上的差异，这个是幼虫能够做出明确判断的，但是它如此喜爱这种食物，除了有美食方面的原因外，一定还有更重要的原因。

对于节腹泥蜂能够使饲养幼虫的昆虫尸体长时间保持新鲜的问题，杜福尔已经给出了十分详细的解释，我就不再多说了。我的象虫有从地下挖来的，也有从捕猎者手下抢来的，尽管它们永远一动不动，但它们都保存得完整无损：色泽新鲜，膜和最小的关节都很柔软，内脏完好。最值得一提的是，即使用放大镜观察，也看不出它们的身上有任何一点伤痕。虽然它们早已没有了生命迹象，但是您看到这一切后，甚至会怀疑它是否真是一具死尸；于是，人们会以为这只昆虫随时有动起来的可能。一般来说，如果天气炎热，死后的昆虫只要存放几个小时就会被烤得干化，一碰就碎掉；而如果天气潮湿，它们就会腐烂发霉。我曾经在不采用任何预防措施的情况下，把节腹泥蜂捕捉到的吉丁和象甲放在玻璃管或者纸袋内一个多月。经过这么长的时间后，它们的内脏仍然像开始时那么新鲜，解剖时也很容易，跟活的动物没什么两样，这可真是太奇怪了。不，在这些事实面前，我们无法相信昆虫已经死去，无法相信它是靠防腐剂的作用才保持新鲜的；生命应该还在它们身上，这是潜伏着的消极的生命，是植物性的生命。就是因为有了这种生命，它才能在与化学力量的腐蚀性斗争中取胜，从而使机体不腐烂。它们还有生命，只不过不能活动而已，像被氯仿和乙醚麻醉了一样，于是出现了我们眼前的这个奇迹，这个奇迹是由神经系统的神秘特性引起的。

很明显，这种植物性生命虽然受到了惊扰而有所减弱，但仍旧在默默地发挥着作用。对此我是有事实证据的，我的那些象虫，虽然它们再也不会醒来，却没有真的死亡，而是处于一种沉睡状态。就在它沉睡的第一个星期内，仍有正常的间歇性排便，直到肠里排空了，排便才终止。尸体解剖证实了这一点。昆虫表现出来的微弱的生命迹象并不仅限于此；虽然它们几乎永远丧失了对外界刺激的反应能力，可我还能引起它的一点波澜。

我把几个刚从地里挖来的、纹丝不动的象虫放到一个装有木屑的小瓶里，木屑上浸了几滴苯。我惊奇地看到，一刻钟后，它们的腿动了动。有那么一瞬间我还以为可以使它们起死回生呢，可这是完全不可能的，它们的腿部动作只是即将消失的反应能力的回光返照，很快便停止了，此后再也无法激活。后来，我又进行了类似的实验，那些象虫从死亡几小时的到死亡了三四天的，全都出现了同样的情况。不过昆虫死亡的时间越久，激起它们动作的时间越长，而且这些动作总是从前部蔓延到后部。首先是触角慢慢摆动几下，然后是前跗骨抖动起来一齐摆动，接着是第二对跗骨，最后第三对跗骨也跟着动了起来。只要跗骨动起来，各个附属部分就都胡乱地动起来，直至最后全都停下来；而这种运动后的停止，通常发生得比较突然。如果昆虫死亡的时间已经很久，跗骨的摆动就不会传到较远的部位，腿就会一直不动。

昆虫死亡10天后，我使用了同样的方法来测试，却无法令象虫动起来；于是我使用电流来刺激它。电流的力量很大，不仅能引起肌肉收缩，还能使那些在苯蒸气下都无法激活的部位动起来。具体的操作过程是：把本生灯的一两个部件装在通了电的细针上，然后把一根针的针尖放在昆虫腹部的末端，另一根的针尖放在它的颈下，通上电以后，跗骨颤抖起来，所有的爪子都弯曲收缩到腹下；切断电流后，爪子又放松出来，伸直恢复原位。刚开始的几天，它们的伸张动作明显而有力，后来逐渐减弱，一段时间后就没有反应了。在第10天的时候，我仍能激起它做出明显的动作；第15天的时候，尽管昆虫的膜还很柔软，内脏还很新鲜，但是它对电流的刺激已经没有任何反应了。我在一些死亡的鞘翅目昆虫的身上也进行了同样的实验，比如被苯或二氧化硫窒息的琵琶甲、楔天牛、青杨黑天牛，然后比较它们的反应。在它们死后两小时的时候，我就已经无法激起它们的任何反应了，而那些被捕猎者害死的象虫，在半生半死的中间状态下度过好几天后，我却可以很容易地让它们动起来。

这些事实证明了昆虫彻底死亡的假设是不成立的，与昆虫尸体依靠某种防腐液的作用而保持新鲜的假设也背道而驰。这些事实只有一种解释：昆虫因为受到伤害而无法活动，它的反应能力在突然遭到麻痹后慢慢消失；与此同时，它的植物性功能相对比较顽强，所以消失得比较慢，因此它的内脏能保持得完好无损，以供幼虫在需要时享用。

此外，我们还应该注意一点，那就是这些象虫是怎么被谋杀的。很显然，节腹泥蜂的毒蜇针很关键。可是，象虫身披坚硬的甲胄，甲胄各处又拼合得十分紧密，毒蜇针究竟刺在了哪里，又是怎么刺进去的呢？在那些死于毒蜇针之下的昆虫身上，即使用放大镜也看不出任何谋杀的迹象。所以只能通过直接研究膜翅目昆虫的谋杀手段来解决，面对这个问题，杜福尔知难而退了，而我也曾觉得无能为力。可我还是进行了尝试，虽然摸索的过程很漫长，但庆幸的是我终于找到了答案。

节腹泥蜂从洞穴里飞出去捕猎时，没有固定的方向，时而飞向这边，时而飞向那边，然后不一定从哪个方向得到了猎物飞了回来。它们的搜寻范围并不明确，总是泛泛地进行；但是它们往返一次几乎不超过10分钟，所以范围应该不是很大，何况这中间还包括它发现猎物、向猎物进攻、杀死猎物的时间。于是我尽量在它的洞穴周围寻找，希望能找到某只正在捕猎的节腹泥蜂。可我花费了一个下午的时间却毫无收获，于是我觉得这种寻找是不可取的，因为即使有幸撞见几个狩猎者正分散在各处捕猎，可它们飞得那么快，眨眼间就会无影无踪，尤其是在种满葡萄树和橄榄树的地方就更不好观察、更无法追踪了。所以，我放弃了这个方法。

\newpage

\chapter{高明的杀手}

\begin{figure}[htbp]
    \centering
    \includegraphics[width=0.8\textwidth]{images/00315.jpeg}
\end{figure}

前面膜翅目昆虫已经向我们展示了它蜇针的蜇入点，帮我们解开了这个秘密。可是问题解决了吗？当然没有，还差得远呢。现在，让我们回过头来，暂时把昆虫刚刚告诉我们的事儿放到一边，开始想一下节腹泥蜂的问题吧。问题就是：节腹泥蜂是如何在地下的一个蜂房里储存足够的猎物，然后在上面产卵，使新生的幼虫吃这些食物的。

为幼虫供应食物的问题，乍一看似乎很简单，可是仔细一想，就会发现其实困难重重。比如，我们人类打猎是靠猎枪，捕到的猎物总是鲜血淋漓、伤痕累累的。膜翅目昆虫对猎物比人类要挑剔无数倍，它要求猎物完好无损，保持好的形状和色泽，薄膜没有破裂，更没有裂开的伤口和恐怖丑陋的死相。它的猎物也果然像活昆虫一样新鲜，蝶翅上精细的彩色鳞片一片不少，可那翅膀如果被我们的手指轻轻一碰，就会使上面的彩色花纹坏掉、脱落。想想看吧，即使是昆虫已经死了，已经成了死尸一具，要做到这一点也很难啊！用脚踩死一只昆虫，人人都能做到，可是要把昆虫干净利落地杀死又丝毫看不出任何伤痕，这可不是人人都能做到的。一只生命力顽强的小动物，即使头被砍下来，也会扑腾一阵子，可是假如在不弄坏它身体的前提下把它杀死，怕是人们就不知道如何是好了！也只有实验昆虫学家才会想到用麻痹的手段。但是即使采用最基本的方法，用苯或二氧化硫的蒸气来麻痹，也不一定成功。因为在中毒的情况下，昆虫挣扎的时间反而太长，挣扎得太久，它的蝶翅就会失去光彩。所以，我们必须采取更强烈的方法，比如让浸着氰化钾的纸带慢慢地散发出氰化氢，或者更好一些也可以使用一硫化碳的蒸气，这种化学物质对捕猎昆虫的人没有伤害。

可见，我们要想杀死一只昆虫，并达到节腹泥蜂在很短时间、很灵活的动作下就达到的效果，就要利用一系列的化学手段才能办到！

一具尸体！可这根本不是幼虫的日常食物，因为幼虫只吃鲜肉，如果它的野味哪怕有一点点臭味，它都会感到恶心，无法忍受。幼虫需要的是一点也没有腐烂的鲜肉，变味就是腐烂的开始。可是，不能把活蹦乱跳的猎物贮藏在蜂房里，就像我们的船上和列车上要为旅客提供新鲜食物一样，是不能把牲畜养在船里或车里的。把娇嫩的卵产在活的猎物中间会有什么后果？幼虫如此娇弱，轻轻一碰都有可能害死它，如果它整整几个星期都处于那些挥舞着带铁刺的长腿的鞘翅目昆虫中间会怎么样？膜翅目幼虫需要的是既像死亡时一样安静又像有生命时一样新鲜的内脏，这似乎是永远无法解决的一对矛盾。面对这样的食物难题，即使人类拥有最广博的知识也解决不了，实验昆虫学家也会承认束手无策；而节腹泥蜂的食橱告诉我们，这一切对它来说根本不算什么。

假设有一所解剖学家和生理学家科学院，设想在一次大会上，像弗卢朗、玛让迪、贝尔纳这样一些人正在讨论这个问题。为了做到食物长时间一动不动，同时又长时间不腐烂变质，人们头一个想法，最自然、最简单的想法，那就是食物罐头。人们会提出使用防腐液，就像著名的朗德学者在他的吉丁问题上设想的那样；人们会设想，膜翅目昆虫的毒液具有卓绝的杀菌防腐效力，但这些奇特的效力还有待证明。用不知为何物的防腐液来代替不知如何进行的鲜肉保存，这也许就是这个学者会议的结论，就像朗德博物学家所做的结论一样。

如果人们强调并明确指出，幼虫需要的不是罐头，因为罐头肉没有鲜肉那种可以颤动的特性，它们要的是一种尽管已了无生机却仍然活着的猎物，那么这个学者会议就会认定它是采用了麻痹的办法。"没错，就是这办法！必须把昆虫麻痹，必须使它既不能活动又没有死亡。"只有一个方法可以达到这个目的：在合适而精准的某个或某些部位上动手，损坏、折断昆虫的神经器官。

即使这个办法可以，那么如果让一个对深奥的解剖学知之甚多的人做，就未必能成功。能使昆虫麻痹却不会让它死掉的神经器官是如何排列的呢？首先，这神经器官在哪里？当然是从头部开始排布下来的，就像高级动物的脑子和脊髓一样。"这简直大错特错，"我们的同行们会这么说，"昆虫就像是一个翻转过来的动物，它用背部走路，也就是说，它的脊髓不是在背部，而是在腹部，顺着肺和肚子分列排布。所以神经器官在它的腹部，要麻痹昆虫只能在那里动手。"

这个困难解决了，又一个更严重的困难出现了。对于解剖学家来说，他手握解剖刀，想插到哪里就插到哪里，即使遇到障碍，他也能排除。可膜翅目昆虫没有选择的余地，它的猎物鞘翅目昆虫披挂着坚硬的甲胄，它的手术刀就是纤细而脆弱的蜇针，甲胄是完全能够抵挡这个武器的。因此，它的手术刀只有几个可供选择的部位，就是那些只覆盖着一层薄膜的关节，那些薄膜是毫无抵抗力的。但是，蜇针虽然能刺进猎物的肢体关节，却完全达不到目的，因为蜇到这些部位，最多就是局部麻醉，无法使它全身麻醉进而抑制它的全部运动器官。如果有可能的话，膜翅目昆虫需要的是一蜇致命，使对方立刻失去活动能力，不做长久的挣扎；它也不能连续蜇好几下，因为蜇的次数太多，就会危及猎物的性命。因为猎物的神经中枢是掌握运动的枢纽，神经就是从那儿分布到各个运动器官上去的，所以它一定要刺在神经中枢上。可是，这个神经中枢形成了很多的核或神经节，幼虫的神经节较多，成虫的较少，这些神经节分布于腹部的中轴线上，彼此距离不等，被神经髓质的双重饰带串联成了一个念珠串。在所有已发育完全的昆虫身上，胸部神经节也就是支配翅膀和腿部运动的神经节，共有3个。捕猎者需要刺的就是这些点，用某种方式摧毁这些点的动作，这样猎物也就无法存活了。

膜翅目昆虫的蜇针如此软弱，以至于它能蜇入的只有两个地方：一处是猎物的颈与前胸之间的关节；另一处是猎物的前胸和胸部其余部位之间的关节，也就是第一对腿和第二对腿之间的关节。颈关节不太合适，它离腿根附近的支配腿部活动的神经节太远。因此，它能刺入的部位就只剩一处了，并且只有这一处。有着高深学识的克劳德·贝尔纳等人在法兰西科学院阐述这个问题时，就是这么说的。膜翅目昆虫蜇针的刺入点就在腹部中线上的第一对腿和第二对腿之间。那么昆虫是得到怎样高明的指引而知道这个问题的呢？

在众多部位中选择一个最脆弱的部位刺入，只有熟知昆虫解剖学结构的生理学家才能做到。可是，只做到这一点还远远不够，膜翅目昆虫还要克服一个更大的困难，然而它竟然做到了，它那超群绝伦的本领真是令人震惊。我们说过，发育完全的昆虫的运动器官是由三处神经节支配的，它们通常是彼此分隔的，尽管也有聚在一处的情况，但十分罕见。不仅如此，这些中心节各自独立地支配着相应部位的运动，所以从立即麻痹的效果来说，如果一个中心点受损，只会导致它支配的相应的肢体部位瘫痪，而不会影响到其他神经节，更不会使受其他神经节支配的肢体瘫痪。如果用蜇针一个接一个地攻击这3个逐渐靠后的运动中枢，而且只通过第一对腿和第二对腿之间的那一个关节点，似乎是办不到的。因为蜇针很短，而且在这样的情况下也很难把针刺入。是的，某些鞘翅目昆虫胸部的3个神经节靠得很近，可还有些鞘翅目昆虫的最后两个神经节是完全联结、粘连、融合在一起的。我们可以看到，随着这些神经节越来越靠近、越来越集中，所引起的运动就越来越完善、越来越和谐，同时，也更易于受到攻击，这样的猎物正是节腹泥蜂需要的。这些鞘翅目昆虫的运动神经中枢挨得太近了，几乎挤到了一起，甚至粘成了一团，最终牢不可分，所以只要一针刺入就能使它立即瘫痪；即使需要多刺几下，要刺的神经节也都是聚集在一起的，甚至可以说，完全集中在蜇针的针尖下面了。

问题的关键是，哪些鞘翅目昆虫容易被刺瘫呢？克劳德·贝尔纳关于生命和器官的杰出理论只是泛泛而谈，已经无法解释这个问题，无法引导我们辨别这些昆虫。我向每一个有可能读过相关资料的生理学家请教，如果他不去查找资料，他能否说出哪些鞘翅目昆虫的神经节点是集中在一起的呢？即使查找资料，他能否立即找到需要的资料呢？现在我们涉猎的已经是专家们详尽研究的细节了，我们已经离开了大多数人走的大路，走上只有少数人熟悉的小路上来了。

相关资料，我在E.布朗夏尔先生《关于鞘翅目昆虫的神经系统》的杰作中找到了。我在该书中看到，符合神经器官集中这种特性的昆虫中，首先就是金龟子。但是它不是节腹泥蜂的猎物，这可能是因为大多数金龟子都体形较大，节腹泥蜂无法向它们展开进攻，也搬不动；加上许多金龟子都以粪便为生，而膜翅目昆虫十分爱洁净，是无论如何也不会到粪便里寻找金龟子的。另外，阎虫科昆虫也是属于运动神经中枢靠得比较近的，但是它们也十分肮脏，寄生在臭气熏天的死尸中，所以也应弃而不谈。至于棘胫小蠹，个子又太小了，所以最后只剩下吉丁和象虫是符合条件的。

在这个漆黑一片的疑团中，忽然现出了一缕光线，实在太令人惊喜了！在无数鞘翅目昆虫中，节腹泥蜂可以捕猎的似乎只有象虫和吉丁这两类，它们每一个条件都十分吻合。挑剔的狩猎者讨厌恶臭和污秽的环境，而它们正生活在远离恶臭和污秽的地方；它们种类繁多，形态千差万别，体态与捕猎者差不多大，捕猎者有自由选择的机会；与其他鞘翅目昆虫相比，它们支配腿和翅膀运动的神经中枢集中在一处，对膜翅目昆虫来说能够一刺即中。象虫胸部的三个神经节紧挨在一起，后两个神经节几乎融为一体了。在这相同的部位，吉丁的第二个和第三个神经节混成一团，又离第一个神经节很近。节腹泥蜂以鞘翅目昆虫为食物是早已得到证实的，同时，有8种节腹泥蜂是从不捕捉别的野味，只捕猎吉丁和象虫的！所以，各种节腹泥蜂的洞穴里堆放着的猎物，虽然表面看毫无相似之处，其实在内部结构上是有共同点的，那就是它们的神经器官都是集中在一处的。

即使拥有再广博的知识，也不能做出比这更高明的选择了，虽然这种选择的背后困难重重，但都被灵活巧妙地战胜了，以至于人们不由得思考起来：自己是否被一己之见蒙骗了？是不是先人的理论概念掩藏了事实的真相？是不是原本自己在异想天开却用生花的妙笔把幻想出来的东西写得跟真的一样？一个科学成果只有在经过各种方式的反复实验和证实后，才能真正确立。所以，让我们通过实验来证实节腹泥蜂生理学手术的真实性吧！如果能用人工的方法得到一只昆虫，使它像被膜翅目昆虫的蜇针刺过一样，一动不动并长时间保持着鲜活的状态；如果用节腹泥蜂捕捉到的鞘翅目昆虫，或者用别的神经节也十分集中的鞘翅目昆虫，都能制造出这种奇迹，而用神经节彼此独立分开的鞘翅目昆虫办不到，那么无论取证有多难，我们都会承认，膜翅目昆虫在其本能的无意识的支配下，确实显示出了卓越的科学本领。好吧，我们来看看实验结果是怎样的。

动手术的办法十分简单：用一根针，如果要更合适些就用一支锋利的金属笔尖，在笔尖上蘸一小滴腐蚀性液体，然后轻轻刺入昆虫的第一对腿与前胸的连接点上，把液体注入它前胸的运动神经中心。我使用的液体是氨水，别的液体只要具有同样强而有力的作用也可以。金属笔尖上的氨水好像一小滴墨水一样。我把笔刺进去以后，出现了不同的效果，而这完全是由昆虫的胸部神经节是彼此靠近还是彼此隔开决定的。对于胸部神经节彼此靠近的昆虫，我用的实验对象是金龟子科昆虫，即圣甲虫和宽颈金龟子；吉丁科昆虫棕色吉丁；还有象虫，尤其是那些我亲眼看到捕猎者追捕的方喙象。对于胸部神经节彼此隔开的昆虫，我选择的实验对象是步甲科昆虫，分别是小红夜蛾、步甲、强步甲、心步甲等；天牛科昆虫，包括楔天牛、沟胫天牛；还有杨树叶甲科昆虫，主要是琵琶甲、盗虻。

在金龟子、吉丁和象虫身上，效果是立竿见影的，致命的药滴刚碰到神经中枢，昆虫的任何行动都戛然而止了，并且没有抽搐一下。这种颇具毁灭性的后果产生得如此神速，甚至比节腹泥蜂的蜇刺都要快。没有什么比顽强有力的圣甲虫突然一动不动更令人惊异的了。蘸着氨水的有毒的金属笔尖在这方面与膜翅目昆虫蜇针的作用相似。被人工戳刺的金龟子、吉丁和象虫不只是一动不动了，还在三个星期、一个月，甚至两个月内都像被蜇针刺过一样，所有的关节都能弯曲，内脏也十分新鲜。在开始的几天，这些昆虫仍能像活着时一样正常排便，在电流的刺激下会动起来。总之，它们的表现跟被节腹泥蜂谋杀的鞘翅目昆虫的表现完全一样。也就是说，被节腹泥蜂蜇刺的猎物产生的状态，与用氨水故意破坏昆虫的胸部神经中枢造成的状态完全相同。可是，因为不能把昆虫保存得完好无损的原因归于氨水，所以必须抛弃任何注射防腐剂的想法，而应当承认，昆虫虽然一动不动，却并没有真正死去，它还有一线生机，在一段时间内，它的各个器官仍像活着时一样新鲜，然后才逐渐腐化直至最后腐烂。可在另外一些情况下，氨水只能麻痹昆虫的腿部，这时，液体的毒性并没有蔓延得很远，昆虫的触角还能轻微地动，甚至在注射了氨水一个多月后，只要轻轻碰碰它，昆虫还会迅速地把触角缩回去，这显然证明了在没有活力的躯体上，生命还没有消失。被节腹泥蜂蜇伤的象虫，也常常有触角抖动的情况。

注射氨水能使金龟子、象虫和吉丁立即终止运动，可是并不是都能使它们昏迷休克。如果刺入得太深，或者注入的氨水滴毒性太强，昆虫就会当即死掉，两三天后，就只剩下一具散发着臭气的尸体了。相反，如果刺得太浅，昆虫在昏迷了一段或长或短的时间后，会苏醒过来，还会恢复某些运动功能。就连那个掠夺者有时也会像人一样出错，因为我曾经亲眼看到过一只被膜翅目掘地虫蜇过的猎物又活了过来。我们在下面即将讲到黄翅飞蝗泥蜂的故事。它的洞里堆放着中了它的毒针的小蟋蟀。我取出了3只可怜的蟋蟀，它们的肌肉十分松弛，种种迹象表明它们已经死了，但其实都还活着。我把这些蟋蟀放在一个瓶子里，好好地保存起来，在大约3个星期内，它们都从未动过。后来，有两只发霉了，而第3只的某些部位，比如它的触角和嘴巴的某些部位则逐渐恢复了活力，尤其令人惊异的是，它的第一对腿又能运动了。既然灵巧的膜翅目昆虫也偶尔会出错，使它的猎物死而复生，我们又怎能要求笨拙的人类做的实验永远万无一失呢？

在第二类鞘翅目昆虫中，也就是胸部神经节彼此隔开的昆虫身上，氨水产生的效果则完全不同。步甲科昆虫似乎最不容易受到伤害。一只粗壮的圣甲虫被刺一针后，会立即无法动弹。可如果是个头比它小的步甲科昆虫受到这一刺，则顶多会引起一阵剧烈而无规则的抽搐。然后，它会慢慢平静下来，休息几小时后，又恢复正常的运动功能，就像从来没受过什么刺激似的。即使对同一只昆虫再做两次、三次或四次实验，结果还是一样，直到它伤得太重而真正死掉了，因为不久它就干瘪、腐烂了。

氨水在杨树叶甲科昆虫和天牛的身上产生的效果最明显。只在它们体内注射了一小滴腐蚀性液体，它们就立刻一动不动了，抽搐了几下后好像就彻底死了。但是像金龟子、象虫和吉丁那种持续很久的麻醉状态，在它们身上只是暂时的，到了第二天它们又能活动了，而且比原来更有活力；只有当氨水的分量足够大时，才会使它们彻底不能动弹，可是这时候昆虫已经真正死亡了，因为它的尸体很快就腐烂了。能使神经节彼此接近的鞘翅目昆虫持久麻醉的办法，在神经节彼此隔开的鞘翅目昆虫身上却毫无效果，对于后一类昆虫来说，麻醉只是暂时现象，到第二天就会完全消失。

论证已经够清楚的了。捕捉鞘翅目昆虫的节腹泥蜂在选择猎物时应用的理论，只有知识最渊博的生理学家和技术最高超的解剖学家才能领悟出来。不能说这只是一次偶然的巧合，因为这样的一致已经不是用一句偶然性能解释的。

\newpage

\chapter{黄翅飞蝗泥蜂}

\begin{figure}[htbp]
    \centering
    \includegraphics[width=0.8\textwidth]{images/00322.jpeg}
\end{figure}

鞘翅目昆虫身披坚硬的盔甲，掠夺者的蜇针只能从它的一个部位刺入。谋杀者对盔甲各处连接点的位置一清二楚，于是从这些地方刺入毒针，它的行刺对象主要是象虫和吉丁这类昆虫，因为它们的神经器官十分集中，刺一下就可以把3个运动神经中枢同时刺伤。但如果昆虫的表层是柔软而没有硬甲的，膜翅目昆虫在行凶时无论刺到它的哪个部位都没有关系。结果会怎样呢？膜翅目昆虫在蜇刺时是否有什么讲究呢？杀人犯在杀人时会直刺心脏，以减少受害者的反抗，减少麻烦，那么这个强盗是否也会像节腹泥蜂那样宁可刺伤运动神经节呢？假如是这样，并且假使这些昆虫的神经节彼此独立，各自发挥着作用，那么一个神经节被麻痹了以后，其他神经节没有被麻痹，又会发生什么情况呢？上述这些问题会在蟋蟀的捕猎者——黄翅飞蝗泥蜂的故事中得到解答。

接近7月末的时候，一直守护着卵的黄翅飞蝗泥蜂，从地下摇篮中飞了出来。罗兰蓟是一种十分普通却十分强壮的植物，整个8月，它经受着烈日的炙烤，依然昂首挺立；黄翅飞蝗泥蜂在它那带刺茎的枝头上飞来飞去，寻找蜜汁。但是对黄翅飞蝗泥蜂来说，这种无忧无虑的生活非常短暂，进入9月后，它就要开始进行辛苦的挖掘和狩猎工作了，它通常把家安在道路两侧的边坡上，空间不是很大。当然，那里必须具备两个条件：易于挖掘的沙土和充足的阳光。此外，它没有采取任何措施遮挡它的住所，来抵挡秋雨和冬霜。一块无遮无挡、饱受风雨侵袭的水平场地是它的最佳选择，但是这块地一定要有阳光普照。如果在它掘地工程进行的过程中突然下了一场暴雨，那它就惨了；因为到了第二天，还没建完的地道就会被沙土堵塞，凌乱不堪，最终它只得放弃。

通常情况下，黄翅飞蝗泥蜂是不会独自施工的，而是10只、20只甚至更多的伙伴们一起对选定的场所进行开发。您如果想了解这些勤劳矿工们忙碌的工作、灵敏的跳跃、迅疾的动作，就必须连续几天紧盯着这样的村落。

矿工们用它们那被林奈称为"犹如利刃"的前腿，像耙子一样快速地挖着土。即使是一只小狗也不会像它们那样兴致勃勃地耙地。与此同时，每个工人都欢快地唱着歌，歌声十分尖厉，时断时续，时而随着双翅和胸腔的振动而抑扬顿挫，就好像一群欢乐工作的伙伴们用歌声来互相鼓劲一样。工地上尘土飞扬，细碎的尘埃落在了它们轻轻抖动的羽翼上，大一些的沙砾则被它们滚到了离工地较远的地方。当遇到大块的不好耙的沙砾时，黄翅飞蝗泥蜂就使出一股猛劲儿，发出一声高亢的叫声，仿佛伐木工人在挥动斧头时喊出的"嗨哟"声。工人们腿颚并用，加倍努力，很快就挖好了一个能容身的小洞。接下来，它开始一会儿挖土，一会儿把挖出来的泥土扒到身后去。在这两项急促的交替运动中，飞蝗泥蜂不是一步步往前走，而是像被弹簧弹出去似的往前冲；它蹦蹦跳跳，腹部微微抽动，触角也一颤一颤的，全身都震得发响。现在，矿工完全进到了地下，离开了我们的视线，可还能听见它那不知疲倦的歌声，偶尔还能看见它把沙土推到洞口的后腿。飞蝗泥蜂有时也会停歇一会儿，或者飞到阳光下抖掉身上的尘土，以免关节上灰土太多影响灵活性，或者到四周巡视一番。

它中止工作的时间不会很长，因此虽然它时有停歇，也会在几个小时内就挖好洞穴。那时，飞蝗泥蜂会在洞口高唱凯歌，并对工程做最后的修整，刮刮不平的地方，弄掉几颗几乎只有它们才注意到的微小土粒。

在我看到的许多黄翅飞蝗泥蜂群中，有一种给我的印象最深刻。养路工人在挖路面一侧的小沟时，把挖出来的湿润的泥土堆在一条大路旁。这些土堆中有1.5米高的锥形土堆，湿土已被太阳晒干。飞蝗泥蜂十分钟情于这个地方，在这里建了一个村落，那是我见过的居民最多的村落。整个土堆，从下到上，洞穴密布，从表面看去就像一块大海绵。生活在这块大海绵上的居民，忙得热火朝天，它们忙忙碌碌、你来我往，会让人不禁联想到某个正在施工的工地。它们用触角把蟋蟀拖到这个锥形城市的斜坡上，存放到蜂巢的食品储存间里。尘土顺着挖掘的巷道流出，满面灰土的矿工不时出现在走廊口上，一会儿进一会儿出，忙个不停。有时您会看到一只忙里偷闲的黄翅飞蝗泥蜂爬上堆顶，似乎要从高处观赏它的杰作。多美妙的场景啊，我真想把这个大大的村落连同它的居民一起搬走！可这只是一个不现实的梦，我怎么可能把那又高又大的土堆连根拔起搬回家去呢！

还是回过头看看那在大自然的土地上工作的飞蝗泥蜂吧，这才是最常见的。洞一挖好，飞蝗泥蜂就开始捕猎了。现在，这种膜翅目昆虫外出捕猎了，我们趁机好好欣赏欣赏它的住所吧。我们说过，飞蝗泥蜂喜欢在平坦的地方群居。但它们的建筑所也不总是平坦如砥，有的地方凸出，生长着一簇草皮或者蒿属植物；有的地方则有皱褶，被植物的细根死死地板结起来。飞蝗泥蜂的家就安在这些皱褶的侧面上。洞穴的入口处先是一个水平的门厅，深两三法寸，这是通向隐藏所的通道，也是储存食物和孵育幼虫的地方。天气不好时，飞蝗泥蜂就藏在门厅里。夜间，这是它的藏身所；白天，这是它的休息室，从洞口探出它那表情丰富的面孔和肆无忌惮的大眼睛。走过门厅就是一个急转弯，坡度较缓，向下延伸至两三法寸处。弯道坡度的尽头是一个椭圆形的蜂房，直径较长，这条水平线就是最长的轴线。蜂房墙壁没有涂抹任何东西；尽管家徒四壁，却仍能看出它们在建筑时是十分认真的。这儿的沙土都被压得很坚实，地板、天花板、墙壁也都经过仔细修整，轻易不会坍塌，也不会因表面过于粗糙而伤害幼虫稚嫩的表皮。这个蜂房与过道相通的入口十分狭窄，勉强能容黄翅飞蝗泥蜂与猎物一起通过。

飞蝗泥蜂会在第一个蜂房内产下一个卵，同时储备下充足的食物，然后就封住入口。当然，它并不是要抛弃这个家，而是在第一个蜂房旁再挖第二个洞，然后同样产卵、备食，紧接着再挖第三个甚至第四个。

到了这时候，飞蝗泥蜂才把堆在洞口的残屑都搬回洞里，把洞外留下的痕迹全都销毁。一个洞穴一般有3个蜂房，也有2个蜂房的，但不常见，4个蜂房的就更不常见了。可是根据对飞蝗泥蜂的尸体解剖可知，它的卵有30个，这就需要10个蜂窝。然而，它们9月才开始筑巢，在月底就要结束，所以建造一个蜂窝和准备食物的时间最多才两三天。在如此短暂的时间内，勤劳的昆虫要挖好洞穴，捕猎12只蟋蟀，还要把猎物千辛万苦地运回来放进仓库，最后还要封住洞口，简直是分秒不停啊！况且，如果遇上刮风的天气或者阴雨连绵的日子，就无法捕猎，甚至什么都干不了。因此不难想象，黄翅飞蝗泥蜂的房子是不可能与栎棘节腹泥蜂的长巷道相比的，栎棘节腹泥蜂的长巷道无比牢固，似乎让人觉得它可以住上百年。栎棘节腹泥蜂的牢固住所是它的宝贵遗产，会代代相传，而它的后代们也会一年比一年挖得更深，所以当我想查看它们的住宅时，总是累出一身大汗，并且即便是我十分卖力，用上我全部的挖掘工具，也总挖不到尽头。黄翅飞蝗泥蜂的洞穴不是这样的，它没有可以继承的遗产，只能白手起家、事必躬亲，而且要迅速干完。它建住所就像搭帐篷一样匆忙，好像用完一天，第二天就要收起来似的。为了弥补这方面的不足，住在那覆盖着一层薄沙土的住所里的幼虫们，穿上了三四层防水外套，这可是它们的独创，连它们的母亲都创造不出来，这可比节腹泥蜂那轻如薄翼的茧安全多了。

一只狩猎的飞蝗泥蜂嗡嗡叫着回来了，它用大颚咬着一只蟋蟀的触角，在与住所一沟之隔的灌木丛上停了下来。那蟋蟀比它要肥胖笨重几倍，它咬着这重物累得筋疲力尽，歇息片刻后，又用腿夹住俘虏，奋力一跃，飞过沟壑，重重地降落在我眼前的这个飞蝗泥蜂村落中。剩下的路程它完全步行。虽然我就坐在一旁，可这只膜翅目昆虫毫不畏惧。它横跨在猎物身上，咬住猎物的触角，昂首挺胸，大踏步地向前走去。如果地面十分平坦，它拖运起来就很容易。但如果遇到盘根错节的草木，它突然被绊住，便会束手无策地惊呆在那里，那样子有趣极了。它往前走走，向后退退，绞尽脑汁，最终依靠翅膀的力量才巧妙地绕了过去，战胜了困难。这情景真是妙趣横生啊！它终于把蟋蟀拖到目的地了，它的触角已经到达蜂巢洞口。这时，飞蝗泥蜂放下猎物，急忙进到洞里。几秒钟后，它又回来了，探出头，欢快地尖叫一声，一把抓起脚下的蟋蟀的触角，迅速拖到了巢穴深处。

它为什么用这么麻烦的方法运蟋蟀呢？完全可以不用先进到住所里，然后再返回来把洞口的猎物运进去，可以像在平地上那样把蟋蟀一直拖进洞里嘛。地道是容得下它们一起通过的，或者也可以它走在前面，把蟋蟀拖在后面，可它为什么不这样做呢？我们目前观察到的各种膜翅目掠夺者，都是不讲究这些繁文缛节的，它们用大颚和中间的那两条腿把猎物抱在腹下，然后就径直地进到了洞穴深处。杜福尔观察到的节腹泥蜂是第一个把工作复杂化的，它先把吉丁放在洞口，自己倒退着进入地道，然后用颚咬住猎物拖到洞里。但是，在同样的情况下，节腹泥蜂采取的方法与蟋蟀捕捉者的方法截然不同。为什么在把猎物运进洞里之前，一定要先进去检查一遍呢？难道黄翅飞蝗泥蜂本性谨慎，每次把猎物送进洞里之前都要先检查一番，看看在它外出时是否有恬不知耻的寄生虫闯入？那么，这种寄生虫是什么呢？各种双翅目的小飞虫都掠夺成性，尤其是弥寄蝇，它们总是守候在膜翅目昆虫的洞口，只要有合适的机会就会把卵产在它的猎物身上；但是它们都不会偷偷闯入，也不敢擅自进入那漆黑的过道里，因为如果它们不幸撞见物主，就会为莽撞之举付出沉重的代价。毫不例外，黄翅飞蝗泥蜂也会受到弥寄蝇的抢掠，可是弥寄蝇决不会进到飞蝗泥蜂的巢洞里去使坏。况且它们不用进到洞里也有充足的时间把卵产在蟋蟀身上！因为如果它们行事足够谨慎，完全可以利用黄翅飞蝗泥蜂把猎物扔在洞口的机会，把繁衍后代的希望托付给蟋蟀。黄翅飞蝗泥蜂绝对不是因为这个原因才先进到洞里的，一定潜藏着某种更大的危险。

下面是观察到的唯一可以对这个问题做出一点说明的事实。在一群埋头苦干的黄翅飞蝗泥蜂中间，是不允许有别的膜翅目昆虫混入的。一天，我突然在蜂群中发现了一个异类，那是一只黑色步甲虫。这个夹杂在蜂群里的不速之客，从容自若、不慌不忙地搬运着沙粒、干草茎碎屑和其他细小材料，盖在了黄翅飞蝗泥蜂窝旁边的一个洞口上，这个洞口与蜂窝口大小差不多。黑色步甲虫干得很认真，任何人也不会想到这个搬运工的卵就埋在那洞口下。一只黄翅飞蝗泥蜂显得有些忐忑不安，它大概是这个洞穴的合法所有者，每当有别的昆虫闯入，它就会扑上去追赶。可它又突然像受了惊吓一样跑出来，黑色步甲虫跟在后面，仍然泰然自若地继续着自己的工作。我观察了这两只昆虫争夺的洞，在洞里发现了一个蜂房，里边有4只蟋蟀。我认为这些蟋蟀远远不是一只黑色步甲幼虫能消受的，因为黑色步甲的个头太小了，至少比黄翅飞蝗泥蜂小一半。我几乎对此深信不疑了。见黑色步甲虫名正言顺、聚精会神地封住洞口，一开始您可能会以为它是这个洞的主人呢，其实它是个地地道道的掠夺者。可是为什么黄翅飞蝗泥蜂个头比它高、力气比它大，却心甘情愿地让它抢了自己的巢穴而不加攻击，只象征性地前去赶一赶？甚至当这个目中无人的不速之客转身要走出洞穴时，前去驱赶的黄翅飞蝗泥蜂反而仓皇逃窜呢？难道昆虫与人一样，要想成功首先要大胆、大胆、再大胆吗？这个掠夺者确实是胆大包天，到现在还不急不慢地在黄翅飞蝗泥蜂面前转悠，而飞蝗泥蜂呢，站在那里愁容满面，不敢向强盗扑过去。

此外，这种被初步确认为寄生虫的黑色步甲虫，我在别的地方曾多次看到它用触角推着一只蟋蟀。那是它凭自己的本事得来的猎物吗？我本来是愿意相信的；可是它沿着路上的车辙漫无目的地走着，好像在寻找合适的洞穴，它那犹豫不决的表现让我满腹狐疑。

它真的曾经进行了艰苦的挖掘吗？可我从未见过它的挖掘成果。还有更可疑的事，我曾见到它把猎物扔掉，可能是没有找到合适的存放猎物的洞穴，它不知道该如何处理猎物吧。在我看来，这样浪费粮食的行为恰恰说明了这财产不是它的劳动所得，因此我怀疑那蟋蟀是它趁黄翅飞蝗泥蜂进洞里时偷来的。我对某种步甲也有同样的怀疑，它们的腹部跟白边飞蝗泥蜂一样也长着一条白带，并且也用蝗虫饲养幼虫。我没看见过它挖地道，却见过它拖着一只蝗虫，飞蝗泥蜂大概不会同意这战利品归它所有的。某些不同种类的昆虫享用着同一种食物，令我思考起这些战利品的合法性来。最后，由于我之前对它们的怀疑或多或少地损害了它的名誉，为了弥补我的过失，我要公开说，我曾亲眼看到跗猴步甲蜂堂堂正正地捉到了一只还没长翅膀的小蝗虫，我看到它挖了巢房，用自己的英勇战斗得到猎物。

因此，对于黄翅飞蝗泥蜂先下到洞底，然后才把猎物运进去的原因，我只能提出一些我的不太成熟的推理。除了把趁它不在时擅闯私宅的寄生虫驱赶出去外，它是否还有别的目的呢？我不得而知，谁能把本能的千百种表现形式都解释清楚呢？人类的智慧太贫乏了，是无法了解黄翅飞蝗泥蜂的智慧的！

不管怎么样，这种表现形式是被证明了永远不变的。因此，我要把那个令我激动万分的实验告诉你们。实验是这样的：当黄翅飞蝗泥蜂进去巡视住所时，我把它暂时丢弃在洞口的蟋蟀挪到了几法寸远的地方。当黄翅飞蝗泥蜂又飞回来时，发出了像往常一样的鸣叫，它东张西望，最后发现了远处的猎物，于是从洞里出来，又把蟋蟀抓回原来的位置上。可是当它放好后，又独自下去了。我采取了同样的方法，飞蝗泥蜂每次都同样懊丧，随之又把猎物放在洞口，然后独自走进洞里；次次如此，只要我还有继续实验下去的耐心，它的战术就一成不变。我一次又一次地在同一只飞蝗泥蜂身上进行了40次这样的实验，我的执着最终被飞蝗泥蜂的固执战胜了，它的战术自始至终没有一丝一毫的改变。

在同一个飞蝗泥蜂村，我对所有我感兴趣的飞蝗泥蜂进行了实验，而它们都无一例外地展示出了百折不挠的顽强毅力，这引起了我久久的思索。我认为，昆虫很显然受一种发自本能的禀性支配，环境是不能发生一丁点改变的；它的反应一成不变，它也许无法依靠自己的能力来获得应对的经验。但是接下来的观察证实了我的这种看法有些绝对。

第二年，我在一个适当的时机来到同一地点进行观察。在挖掘洞穴方面，后代继承了上一辈的住所，连建筑的方法也忠诚地继承了上一辈的传统。我也对它们进行了把蟋蟀放远的实验，实验结果都一样。这一年的飞蝗泥蜂与上一年的飞蝗泥蜂完全没有两样，都执着地进行着它们的动作，尽管那是白费力气的。当我在错误认识的泥潭中陷得越来越深时，突然遇到了一个幸运的转机，我在离第一个飞蝗泥蜂村落较远的地方发现了另一个飞蝗泥蜂群。于是，我又开始进行实验。刚开始实验的结果与我之前所得到的结果并无两样，可是经过了两三次的实验后，飞蝗泥蜂直接跨到蟋蟀身上，用大颚咬住它的触角，把它一下子拖到洞里去了。最傻的人是谁呢？是实验者，他的阴谋被诡计多端的膜翅目昆虫识破了。在别的洞穴里，我的计划也或早或晚被挫败了，它们把挪到了远处的猎物直接搬到家里，而不是固执地先把它拖到洞口放一会儿再运进去。这说明了什么？我今天观察的昆虫属于另一个血统，子孙后代总是在祖辈们选好的居所里生存，而这些村民比我上一年观察到的那些更灵活多变。习性是代代相传的，因为祖先的特性各有不同，所以有的部族灵活些，有的部族则愚钝些。黄翅飞蝗泥蜂跟人类一样，其才智会因出生地点的不同而有所不同。

第二天，我又到别的地方进行蟋蟀测试，结果每次都取得了成功。因为跟我第一次观察时的情况一样，我又遇到了一个头脑简单的部族，一群愚蠢至极的村民。

\newpage

\chapter{匕首三击}

\begin{figure}[htbp]
    \centering
    
\includegraphics[width=0.8\textwidth]{images/00329.jpeg}
\end{figure}

毫无疑问，黄翅飞蝗泥蜂在杀死蟋蟀时使出了它的看家本领，因此十分有必要看看它捕杀猎物的过程。在观察节腹泥蜂时，我进行了很多尝试，并且让我受益匪浅，所以我将在节腹泥蜂身上采用的卓有成效的方法，运用到了黄翅飞蝗泥蜂的身上。这方法就是在很短的时间内将它已经杀死的猎物取走，用另外一只活的猎物代替。我们在前面已经看到，黄翅飞蝗泥蜂在进入地洞之前会把猎物先扔在洞口，然后独自走到洞底去，这就是我们偷天换日的绝好机会。如果黄翅飞蝗泥蜂能勇敢无畏地爬到您的手指边，甚至爬到您的手上来抢那只活的替代猎物，就再完美不过了，因为我们可以在近距离内清楚地看到整个杀戮的全部细节。

要想找到活蟋蟀并不难，因为它们总是躲在石头下面晒太阳，所以几乎随便掀起一块石头就能看到它们一大堆挤在一起。它们是一些小蟋蟀，翅膀尚未长成，本领尚未练就，还不会像成年蟋蟀那样挖掘深深的隐避所来躲避黄翅飞蝗泥蜂。我可以尽情捕捉，要多少都可以，很快我就捉到了足够数量的蟋蟀。一切都准备好以后，我爬上观察所的高处，在黄翅飞蝗泥蜂村落中间的高地上安静地等待着。

一个猎手捕猎归来了，它把蟋蟀放在住所的入口处，独自走到了洞里。我趁机迅速拿走这只死蟋蟀，把我捉来的活蟋蟀放在洞口不远处。猎手从洞底返回了，它看了几秒钟就跑过去抓那只猎物。我瞪大眼睛，屏住呼吸，专心致志地观察着。我是无论如何也不会错过这场即将上演的悲剧的。蟋蟀惊恐万分，蹦跳着慌忙逃窜。飞蝗泥蜂朝它猛扑过去，双方扭成一团，顿时尘土飞扬。两个决斗者轮流占上风，一时难分高下。可是最终猎手胜利了，蟋蟀被打翻在地，足爪乱舞，双颚乱咬。

战争结束后，猎手立即着手处理战利品。它反向趴在对手的肚子上，用大颚咬着蟋蟀腹部末端的一块肉，用前足将蟋蟀那竭力挣扎的肥后腿遏制住，同时用中足勒住猎物抽动着的肋部，后足则像两根杠杆似的按在蟋蟀的脸上，蟋蟀脖子上的关节都大大地张开了。这时，为防止被蟋蟀的颚咬到，飞蝗泥蜂把腹部弯成一个90度角的凹面。我激动万分地看着，飞蝗泥蜂第一下刺在了猎物的脖子里，第二下刺在它胸部前两节的关节间，然后又刺向了腹部。说时迟，那时快，刹那间，一场凶杀就结束了。飞蝗泥蜂稍微整理了一下衣襟后，准备把牺牲品运到住所去，而奄奄一息的蟋蟀腿脚还在抖动着。

前面我对飞蝗泥蜂的捕猎过程只做了一个简单的介绍，现在我们就仔细领略一下它那无比绝伦的战术吧。节腹泥蜂的捕猎对象几乎没有进攻性武器，在争斗中处于被动地位，甚至连逃脱都很难，它们身上的坚甲是唯一可以防身的武器，可是凶杀者对坚甲的致命点了如指掌。然而，这里的情况却有着天壤之别！黄翅飞蝗泥蜂的捕猎对象长着凶狠的大颚，这大颚一旦咬住猎手，就能将它碎尸万段；同时，蟋蟀的两条腿粗壮有力，上面有两排尖锐锋利的锯齿，它们完全可以依靠双腿的弹跳逃之夭夭，或者袭击对手，把黄翅飞蝗泥蜂狠狠地踢翻在地。所以你们会看到，飞蝗泥蜂在用蜇针蜇刺之前，采取了十分小心谨慎的预防措施。被害者仰面朝天倒在地上，无法依靠后腿的力量逃逸了，而假如它处在正常的姿势下遭到攻击，就一定会逃走，像象虫在遭受节腹泥蜂攻击时那样。它那带锯齿的大腿被黄翅飞蝗泥蜂的前足死死地压住，丧失了攻击性；它的双颚被飞蝗泥蜂的后腿远远地支开并顶住，虽然张得很大，态势逼人，却无法对敌人构成任何威胁。但是对于黄翅飞蝗泥蜂来说，只做这些是远远不够的，这不能保证它一定能躲开猎物的进攻，它还需要死死地勒住猎物，使它一动不能动，以便蜇针能把毒汁准确地注入适当的部位；也许黄翅飞蝗泥蜂就是为了勒住蟋蟀，才咬住它腹部末端的肉。太神奇了，我们即使调动所有的想象力去制定一份进攻计划，也找不到比这更好的办法，就算是古代角斗场上的角斗士，在与对手搏击时采用的方法，也不一定比像这样经过深思熟虑的更精妙。

我在前面说过，蜇针在被害者身上刺了好几下，首先刺在脖子上，然后刺在前胸，最后刺在接近肚子根部的地方。这三刺刺得干净利落，正是本能赋予它的高超本领和万无一失的手段。让我们回忆一下前面对节腹泥蜂的研究得出的主要结论吧。黄翅飞蝗泥蜂幼虫赖以生存的食物，虽然一动不动，却并没有真正死亡。这些猎物只是全身或者局部被麻醉了，虽然动物性生命已在不同程度上被消灭了，但植物性生命——营养器官的生命还久未消失，所以猎物不会腐化，以至于很久以后，幼虫的食物还保持着新鲜。为了造成这种麻醉，膜翅目猎手使用的手段相当高明，甚至当今如此先进的科学才能向实验生理学家提供这样的建议，即用蜇针的毒液破坏指挥运动器官的神经中枢。另外，我们知道，在节肢动物的神经系统中，各个中枢或神经节在某种程度上是彼此独立的，所以某个神经节的损坏，只能引起，也可以说至少能引起它控制的相应部位的瘫痪；各个神经节彼此距离越远，情况就越明显。相反，如果各个神经节都聚集在一处，那么只要损坏共同的神经节，就会引起受神经分支支配的所有部位的瘫痪。吉丁和象虫就是这样，节腹泥蜂只要向它们的胸部神经中枢刺一下，就能使它们完全瘫痪。让我们来解剖一只蟋蟀看看吧！蟋蟀的三对脚是如何活动的呢？早在解剖学家得出结论之前，黄翅飞蝗泥蜂就已经知道了我们要探究的秘密：蟋蟀的三个神经中枢彼此相隔很远。由此可见，用蜇针连续刺三次是符合常理的。高高在上的科学啊，您自叹弗如了吧！

蟋蟀被黄翅飞蝗泥蜂刺伤后，与被节腹泥蜂蜇针蜇伤的象虫一样，虽然表面看是死亡了，但其实并没有真死。为证实节腹泥蜂的猎物方喙象还残留着生命，我曾采用了人工的实验方法，但在猎物外皮十分柔软的情况下，它内部残存着的弱小生命是能够明显地反映出来的，因此我也就不必再采取任何措施去证实了。在凶杀事件后，如果我们持续一个星期、半个月甚至更长的时间观察一只仰卧着的蟋蟀，就会看到它的腹部经过长久的间歇后，有很强的搏动。还经常会看到它的触须在微微抖动，触角和腹肌的运动也很明显，彼此岔开，然后突然并拢。如果把刺伤的蟋蟀存放在玻璃管里，可以完好地保鲜一个半月。黄翅飞蝗泥蜂的幼虫生长不足半个月就会把自己密封在茧里，所以它们在整个宴会过程中都有鲜肉可吃。

捕猎工作结束了。一个洞穴储备了三四只蟋蟀作为食物。蟋蟀堆放得井井有条，全部都是背朝下，头部位于洞穴的尽头处，脚在洞口。它们在每只蟋蟀身上产下了一只卵。最后要做的就是封住洞口了，它们把挖洞时堆在洞口的沙土堆快速地往后一扫，填充到通道里。飞蝗泥蜂总是用前腿扒开残土堆，把大个的沙砾挑拣出来，然后用大颚叼走，用以加固易碎的洞壁。如果它在触手可及的地方找不到合适的沙砾，就会到附近的其他地方去找，它挑选得十分认真，就像泥瓦匠在挑选建筑材料似的。植物的残枝断根、枯枝败叶在这个时候都发挥作用了。不一会儿，地洞的痕迹就从地面上消失了，如果不留心做个标记，怎么仔细地寻找也找不到这个地下住所。这个洞封好以后，它们开始挖下一个，挖好后放上食物，把它封起来。卵产完以后，飞蝗泥蜂又开始快乐无忧地四处游荡了，直至初冬乍冷时，它充实的一生才结束。

黄翅飞蝗泥蜂完成了它的使命，可我还想看看它的武器。它那制造毒汁的器官，由两根分成许多细枝的管子组成，两个管子都插到了一个梨形的共用贮汁库或者说贮壶里。

从贮壶里伸出来的一条纤细的管，深入到了蜇针的轴线中，把毒汁输送到蜇针的末梢。蜇针细极了，与黄翅飞蝗泥蜂高大的身材相比，特别是与它蜇到蟋蟀身上产生的效果相比，实在细得令人惊叹。与蜜蜂长着倒齿的蜇针不同，黄翅飞蝗泥蜂的针尖是光滑的。原因很简单：蜜蜂的蜇针主要在自己受到侮辱而施展报复时使用，而这对它本身也有致命危险；因为蜇针的倒齿在刺入对手的体内以后，很容易被伤口钩住而无法拔出，这样一来，蜜蜂的腹腔末端就会被扯得开膛破肚。假如黄翅飞蝗泥蜂在第一次捕猎时就被它的武器害死了，那它还要这样的武器做什么呢？它的蜇针主要是在给幼虫捕猎食物时才使用的，专门用以刺伤猎物。即使带锯齿的蜇针每次都能轻易拔出，我想大概也不会有哪只黄翅飞蝗泥蜂愿意在蜇针上长齿的。对它来说，蜇针不是耀武扬威的武器，在复仇的时候拔出这匕首，无疑是一大快事，可快意的代价是沉重的，睚眦必报的蜜蜂经常会为此付出生命的代价。黄翅飞蝗泥蜂的蜇针是工作器械，一个劳动工具，它关系着幼虫的存亡，因此在与猎物搏斗时，这工具必须便于使用，既能迅速地刺入猎物的体内，又能很痛快地抽出来，所以从这点上看，光滑的刀刃比有倒钩的刀刃更符合要求。

黄翅飞蝗泥蜂能够以迅雷不及掩耳之势将强壮的猎物打倒，我很想亲自感受一下被它蜇到究竟有多疼。好吧！我试了，我跟您说，这太不可思议了，它的针刺得一点儿也不疼，简直没法跟暴虐的蜜蜂和胡蜂的蜇针相比。我没用镊子，而是大胆地用手指把它抓住，因为我在研究中还会用到它。我可以很确定地说，各种节腹泥蜂、大头泥蜂，甚至看一眼都令人毛骨悚然的身体硕大的土蜂，以及我能观察到的所有膜翅目强盗，蜇得都不痛。但这不包括蜘蛛的捕猎者——蛛蜂，尽管它们蜇得也远不及蜜蜂蜇得痛。

最后要说明的是：我们知道，膜翅目昆虫的蜇针可以用于自卫，比如胡蜂，在遇到无法无天的入侵者时，它会义无反顾地猛扑上去，给对手的鲁莽行为以严厉的打击。相反，以蜇针为捕猎工具的膜翅目昆虫，脾气十分温和，就像它们意识到了自己贮壶里的毒汁对于子女有多么重要似的。这毒汁是保护种族的工具，是赖以生存的手段，所以它们只是在庄重严肃的狩猎场合，才小心翼翼地使用它，而不是用来炫耀自己敢于报复的勇气。当我置身于各种黄翅飞蝗泥蜂的部族之中，捣毁它们的家，抢掠它们的幼虫和食物时，我从没被蜇过。它只在被抓时，才不得不使用那件武器，而且假如不是我把手腕或者别的比手指更娇嫩的部位伸到它的蜇针下，它都未必会刺进皮里！

\begin{figure}[htbp]
    \centering
    \includegraphics[width=0.8\textwidth]{images/00330.jpeg}
\end{figure}

金铃蟋：直翅目蟋蟀科的一种，体形娇小敏捷，鸣声清亮。

\begin{figure}[htbp]
    \centering
    \includegraphics[width=0.8\textwidth]{images/00331.jpeg}
\end{figure}

纸巢黄蜂与黄翅飞蝗泥蜂同属膜翅目昆虫。其触角、翅、跗节呈橘黄色；体乌黑发亮，有黄条纹和成对的斑点。

\begin{figure}[htbp]
    \centering
    \includegraphics[width=0.8\textwidth]{images/00332.jpeg}
\end{figure}

虎头蜂：民间的俗称，属胡蜂类，并不是正式的学名，头大得像老虎，性情凶猛，身体长有虎斑纹，因此得名"虎头蜂"；又因为虎头蜂窝巢形状很大，像鸡笼一样，也被称为"鸡笼蜂"。

\newpage

\chapter{幼虫和蛹}

\begin{figure}[htbp]
    \centering
    
\includegraphics[width=0.8\textwidth]{images/00333.jpeg}
\end{figure}

黄翅飞蝗泥蜂的卵是白颜色的，圆柱形，微呈弯弧状，三四毫米宽。卵不是随便产在猎物身上的哪个部位都可以的，自始至终都有一个固定的绝佳产地，那就是横放在蟋蟀的胸膛上，稍稍靠边，在第一对脚和第二对脚之间。白边飞蝗泥蜂的卵和朗格多克飞蝗泥蜂的卵也产在类似的部位，前者产在蝗虫的胸膛上，后者则在距螽的胸膛上。因为我从没见到它们的产卵位置有什么变化，所以我断定，这个部位对于幼虫的安全一定十分重要。

卵产下三四天后就孵化了。一层极细极薄的膜裂开后，一只娇弱的小虫出现在我们眼前，它浑身像水晶一样明亮晶亮，后面稍稍有些鼓胀，前面像被勒住了，从后到前逐渐变细；它的身体两侧各有一条白色的小细带，主要是由支气管构成的。这个衰弱的小生命横躺着，就像卵刚产下时那样。它的头好像被放在了卵前端被固定的位置上，身体的其余部位则只是依靠在猎物身上，并没有与猎物融合在一起。由于它的身体是透明的，我们一眼就能看出它体内有快速的起伏运动，蠕动波一波接一波，很有规则地进行着，这些波从身体的中间开始，分别向前、向后逐渐蔓延开来。这是消化道的起伏运动，因为它正在猎物身上大口大口地吮吸着汁液。

让我们来仔细看看这个引人注目的场面吧。黄翅飞蝗泥蜂的蜂房里堆着三四只蟋蟀，这些猎物纹丝不动地仰卧着。在朗格多克飞蝗泥蜂的蜂房里只有一只猎物，那是一只身体富态、体形较大的距螽。如果让幼虫离开它那得以汲取生命源泉的部位，它就必死无疑；如果它不幸掉落下来，就会没命，因为它十分虚弱，一点力气也没有，根本无法走动，又怎么能回到它吮汁的位置呢？那只猎物只要轻轻一动，就可以将这个从它的内脏吮吸汁液的幼虫抖落下去，可是这庞然大物却任其摆布，连一下表示抗议的颤抖都没有。我当然知道，它被麻醉了，凶杀者的小针使它的大腿失去了效用；但它被蜇的时间不长，身上那些没被刺到的部位多少还有点感觉，也存留着一丝活动能力：腹部微微颤动，大颚一张一合，腹部的肌肉和触角左右摇晃。假如幼虫咬到这些尚有一丝感觉的部位，比如咬到大颚旁，或者咬在肉嫩汁鲜、最适合给虚弱的小虫先吃的肚皮上，会发生什么情况呢？如果蟋蟀、蝗虫、距螽被咬到致命之处，它们的皮肤至少会有一丝轻微的颤抖，而这轻微的颤抖就足以将那衰弱的幼虫抖下身去，幼虫也就一定活不成了，因为它就掉在那对儿可怕的大钳子下面啊！

但是对幼虫来说，猎物被黄翅飞蝗泥蜂蜇过的胸部是最安全的地方，可以躲避那些危险。实验者在新捕捉到的猎物的胸部，也只有在它的胸部，才可以用针尖随意刺探，到处戳洞，而使受刑者感觉不到一丝一毫的疼痛。所以卵永远只会产在这里，幼虫也总是从这儿开始吃它的猎物。这个部位使蟋蟀即使被咬到了也感觉不到疼，所以它始终一动不动。慢慢地，当幼虫啃咬的范围逐渐扩大，到达猎物的敏感部位时，它会发出本能的挣扎，但是已经来不及了，它被麻醉得太久了，而且敌人正日益壮大。这就是卵总是一成不变地产在被蜇针刺伤的胸膛上的原因，但不是在胸部正中间，因为对于幼虫来说那里的皮也许太厚了，所以它们把卵产在胸侧靠近腿根的地方，那儿的皮更细滑些。母蜂的抉择是多么合情合理又细致周到啊，在一片漆黑的地洞里，它能准确辨别猎物的身体部位，并选定了唯一适合的部位产下了它的卵。

我曾饲养过黄翅飞蝗泥蜂的幼虫，并把从蜂房里拿来的蟋蟀一只接一只地喂给它，就这样，这个婴儿在我的从不间断的密切关注下迅速长大了。就像我刚才所说的，幼虫在第一只蟋蟀，也就是它破卵而出的那只蟋蟀身上，向猎手蜇针第二次所刺的第一对腿和第二对腿之间的部位发起了进攻。不出几天，弱小的幼虫就在猎物的胸膛上挖出了一个有它半个身子深的洞。这时，我们就会看到，那被活活啃咬的蟋蟀徒劳地颤抖着它的触角和腹部肌肉，大颚白费力气地一会儿张开一会儿闭合，某只脚甚至会轻微地动一动；可是，它被它的敌人堂而皇之地掏空了内脏，而不能给它任何惩戒。这对于瘫痪了的蟋蟀来说，无异于一场恐怖的梦魇！

六七天以后，幼虫的第一份粮食吃光了；那只猎物只剩下了皮包骨，骨架的形状和构造几乎保持着原来的样子。这时，幼虫大约有12毫米长了，它从蟋蟀胸腔上的洞里钻了出来。在钻出来的过程中，它蜕了一次皮，而蜕下的皮一般就放在它钻出来的洞口上。蜕皮以后，歇息片刻，它就开始吃第二份粮食了。现在的幼虫已经强壮了，不再把蟋蟀那软弱无力的动作放在眼里，而且蟋蟀也越来越麻木，直至最后连反抗的意愿都没有了。所以幼虫可以在毫无防御措施的情况下对它发起进攻，并且通常从肉最嫩、汁最多的肚子开始。很快，它就开始吃第三只蟋蟀了，然后是第四只。在吃第四只的时候，它只用12个小时就吃光了。这四只猎物最后都只剩下咬不动的外皮了，外皮上都是被咬出的很多零零碎碎的洞，所有能吃的都被掏空了。如果这时再给幼虫拿来第五只蟋蟀，那么它根本不会予以理睬，甚至碰都不碰一下，这可不是想节食，而是因为要排泄了。注意到了吗？从开始直到现在，幼虫还没排过便，它的小肠内装着四只蟋蟀，已经快要胀裂了。

因此，即使是新粮食也不足以勾起它贪吃的欲望，它现在想给自己造一个丝质房间了。总之，它的这顿盛宴持续了10～12天。此时，幼虫已经有15～30毫米长，最宽的部分可达5～6毫米。幼虫的形状大多是后部略宽，越往前越逐渐细下去，大多数膜翅目幼虫都是这个形状。幼虫算上头部共有14节，它的头很小，大颚疲软无力，会让人们怀疑这大颚是如何起到刚刚那样的啃咬作用的。在这14个节段中，中间的那些节段有气门。它的外衣白底泛着淡黄色，上面布满白垩般的白点。

我们在前面已经知道了，幼虫在吃第二只蟋蟀时是从肚子开始的，因为那是猎物身上汁液最多、肉最嫩的部分。就像小孩吃面包片时总会先吃果酱，然后才不情愿地啃起面包一样，幼虫也先吃最可口的部位——猎物腹部的内脏，然后再耐着性子从角质外壳里掏肉吃，以便有空的时候慢慢消化。不过刚破卵而出的娇嫩小虫，在开始时对食物并没有这么贪婪。它是先吃面包，然后才吃果酱的。因为它没有别的选择，第一口只能咬在母亲产卵的猎物胸部。这个部位虽然多少有点硬，但是绝对安全，因为母亲的蜇针在猎物的胸部刺了三下，已经使那里完全丧失活力了。猎物的其他部位总是免不了偶尔有些痉挛性的抽动，虚弱的小虫很容易被甩下去，从而置身于一个随时可能丧命的危险境地。因为它掉落在了一堆猎物之中，这些猎物的后腿长着锯齿，时不时地会猛地一踢或一蹬，虽然踢蹬的次数越来越少，间隔的时间越来越长，但是它们的大颚还是会咬人的。所以母亲选择产卵地点时，完全是出于对幼虫安全的考虑，而并非根据它的食欲来决定的。对于这个问题，我还有一个疑问。幼虫的第一份粮食，也就是产有虫卵的那只蟋蟀，可能比其他的蟋蟀更容易给幼虫带来危险。首先，幼虫还十分柔弱；其次，猎物刚刚捕来不久，有充足的条件表现出它尚存的生命。所以，第一只猎物应当麻醉得更彻底些，黄翅飞蝗泥蜂刺它三下是有必要的。但是其他的猎物将随着时间的流逝越来越麻木，而与此同时，幼虫的身体越来越强壮，何必也刺三下呢？从幼虫吃第一份食物开始，猎物被麻醉的程度就越来越深，那么，对它们只刺一下、两下不就可以了吗？对于膜翅目昆虫来说，毒液是十分珍贵的，它们是不会随便浪费的；这相当于它们狩猎的子弹，应该节约使用。或者说，尽管我看到过它们对同一个猎物连刺三下，但在别的地方，我也看到过只刺两下的。是的，黄翅飞蝗泥蜂腹部的针尖颤抖着，似乎也在寻找合适的部位刺第三下，可这第三下我没见到，至于它到底刺没刺不得而知。所以我更愿意相信：它们只刺第一只蟋蟀三针，为了节约，其他的只刺两针。下面对捕捉毛虫的砂泥蜂的研究，将证实我的观点。

幼虫在吃完了最后一只蟋蟀后，就开始忙着织茧了，并且只需48小时就大功告成了。此后，这位工人就待在任何人也无法侵入的万无一失的隐蔽所内，经历着它生命中必须经历的那种深深的麻木状态，度过这种半睡半醒、半生半死的日子，经过10个月的脱胎换骨才破茧而出。它的茧的复杂程度十分罕见，除了外部有一层粗糙的网状物以外，还清晰地分成三个层次，好像三个套在一起的茧。现在，我们来欣赏一下这个丝质建筑物的结构吧。

建筑物的最外层是带网格的粗纱，像蜘蛛网一样，幼虫先把自己关在里面，像攀在吊床上一样，以便更舒适地织造真正意义上的网。这个用来充当脚手架的、匆忙之中编造的网，是残缺不全的，由随便抛出来的丝编成，并掺和着沙粒、土块和幼虫吃剩的残羹剩饭，还带着蟋蟀那鲜血淋漓的大腿、脚和头颅骨。再往里是封套，那应该算是第一层茧，由淡棕色毡状膜构成，十分细腻、柔滑，有着不规则的褶皱。几根随便抛出的丝线将脚手架和外套联结在一起。这外套好像圆柱形的钱包，四面是密封的，因为它要装的东西很少，这个外套又太大，所以表面出了褶皱。

紧挨着这一层的是一个"塑料匣子"，明显比包在它外面的钱包小，形状近似于圆柱，上端是圆形，幼虫的头就放在那儿；下端呈钝锥状。匣子是淡红棕色的，但下端锥体的颜色更深些。它看上去十分坚硬，可是轻轻一按就裂开了，只有锥极用手指按不破，里面可能装着某种硬物吧。打开这匣子可以看到，它是由两层紧紧贴在一起但彼此能够分开的东西组成的。外层跟上面的钱袋一样，是丝毡；内层，或者说整个茧的第三层，好像一种深紫棕色的发光涂料一样，容易破碎，手感柔软，质地好像也与茧的其他地方不同。用放大镜可以看到，它不像外套那样是丝毡的，而是一种特殊的清漆涂料，其来源很不同寻常，下面将会谈到。至于茧的锥极为什么那么坚硬，原因在于其紫黑色塞子。这种紫黑色塞子是由一种极易破碎的材料做成的，上面有许多黑点，闪闪烁烁的。其实，这是幼虫在整个蛹期内一次性排泄的干粪便团。就是这粪团才使得茧的锥极颜色比较深的。这个复杂的住所平均长度为27毫米，最宽处为9毫米。

现在回过头来看看茧内部涂抹着的紫色清漆吧。刚开始，我以为这清漆是由丝腺产生的，丝腺在编织完丝的双重匣子和脚手架后排出清漆。为了使自己更加确信这一点，我对一只已经结束纺织工作但还没有开始涂漆的幼虫进行了解剖。可是，在幼虫的丝腺里，我根本没有找到一丁点紫色液体的痕迹；这种颜色只存在于消化道里，消化道里满胀着苋红色的东西，下面我们还会在茧的粪便塞上见到这种颜色。除了这种苋红色外，其他都是白色或者微带黄色的。幼虫用它的粪便残余来涂刷它的茧，完全出乎我的意料。我坚信这粉刷浆是消化道的产物，同时我猜测它用嘴排出胃里的苋红色物体，然后用它做清漆涂料。可这只是我的猜测，无法确定，因为手脚笨拙的我错失了好几次证实这个问题的好机会。它只在最后一道工序结束后，才把消化的残余物揉成一团排出去。幼虫之所以做出这一令人反感的举动，非把粪便留在住所里不可，就是这个原因。

无论如何，清漆层的用途是毋庸置疑的；它完全不透水，能保护幼虫不受潮湿，而幼虫母亲为它挖的隐蔽所很不牢固，极易受潮。我们知道，幼虫是埋在敞露的沙土底下不足几法寸深的地方的。为了弄清这涂着清漆的茧的抗潮力有多大，我把茧在水中放了好几天，结果它的内部一点儿也没湿润。黄翅飞蝗泥蜂茧的这种多层次构造和巧妙布置，能够保护摇摇欲坠的巢里的幼虫。节腹泥蜂的茧则不同，是埋在半米深以下的、干燥的砂岩层隐蔽所下的。

它的茧的形状就像一个很长的梨，被切掉了较细的一端，只剩一个丝质的外套，十分纤薄、细腻，以至于透过外套可以看到里面的幼虫。在我进行的这么多次昆虫学观察中，我发现幼虫的本领与母亲的本领总是互补的。如果洞穴隐藏得够深，遮挡得够好，茧就用轻盈的材料制造；如果洞穴很浅，饱受风吹雨打，茧的结构就会非常结实。

9个月过去了，在此期间，茧内的工作都是秘密进行的。我不知道幼虫的变态是如何进行的，因此只能越过这个阶段等待成蛹的阶段，于是我从9月末一直等到第二年的7月初。这时，幼虫刚把褪了色的皮蜕掉；蛹是个过渡阶段，也可以说它是一个已完成变态但尚在襁褓中的昆虫，正在安静地等待着一个月后的苏醒。它的腿、触角、嘴和那没有长成的翅膀，就像液态的水晶一样光亮剔透，并有规则地摊在胸部和腹部下面。它身体的其余部分呈现出浊白色，就是白色中带着一些淡淡的黄色；腹部中间的四个节段，每边都有狭窄而圆钝的伸出部分。最后一节末端上，有类似于圆扇面的膨胀叠片，下面有两个并排着的锥形乳突，这一切组成了一个分布在腹部周围的附属器官。这就是那个柔弱的生物的特征，它为了蜕变成黄翅飞蝗泥蜂，必须穿着半黑半红的服装，然后再把紧紧包裹着它的薄皮蜕掉。

我曾经日复一日地等待着蛹的出现，观察着它颜色的变化，并进行实验，想验证太阳这个大自然万事万物的光源宝库、这个色彩斑斓的调色板，是否会影响它的变化。为此，我把一些蛹从茧里取出放到玻璃瓶里。然后，把一些瓶子放在昏暗的地方，让蛹处于自然条件下，以便将来进行对比；另一些瓶子则悬挂在一面白墙上，让它们天天接受太阳光的强烈照射。在这样截然不同的条件下，两组蛹的颜色变化是相同的；或者说，即使有某些轻微的不同，也是接受阳光照射的那些蛹的颜色变化较小。因此，与植物的情况完全不同，阳光的照射不会影响或加速昆虫颜色的变化。事实就是如此，在色彩最斑斓绚丽的昆虫世界中，比如吉丁和步甲，人们可能以为它们那美丽的颜色是从阳光那儿采来的，但事实上，它们是在漆黑的地下，或在那被虫蛀的百年老树的树干深处调制出来的。

第一个颜色变化出现在眼睛上，眼睛上最先有了一个带颜色的线条，然后角质的复眼相继从白色变成淡黄褐色，又变成深灰色，最后变成黑色。前额顶部的单眼的颜色也随之发生了改变，这时蛹身的其余部分还保持着原来的白色，丝毫没有改变。必须指出，在动物界中，眼睛这个最纤细的器官最早成熟是普遍现象。不久，中胸和后胸之间那条沟上出现了一道烟黑色，24小时后，整个中胸的背部都变成黑色了。与此同时，前胸的那一片颜色逐渐模糊了；后胸上部的中央部分出现了一个黑点，大颚被一层铁色覆盖。胸部两端的胸节颜色逐渐加深，直至蔓延到头部和臀部。只需要一天时间，头和胸部两端的胸节就会从烟黑色变成深黑色。这时，腹部的颜色变化越来越快：前部腹节的边沿染上了金黄色，而后部腹节镶了一道灰黑色的边儿。最后，触角和腿的颜色越来越深，变成了黑色，腹底完全成了橘红色，末端是黑色。此时，除了跗节和嘴是透明的棕红色以及还未长成的翅膀是灰黑色以外，它的一身华丽彩服就已搭配成功了。再过24小时，蛹就要彻底挣脱束缚了。

只要六七天，蛹的颜色就基本固定下来了，只有眼睛的变色是早在半个月以前就开始的。通过这些概括性的描述，就可以很容易掌握颜色变化规律了。我们看到，除了复眼和单眼像所有的高级动物那样提前完成变色外，其他部位的颜色都是从中央部位——胸部开始的，从那里开始，向四周扩散，先是蔓延至胸部的其他部位，然后是头部和腹部，最后推及至附属器官、触角和腿。跗节和嘴的变色更晚些，翅膀则是出了匣子之后才开始有颜色的。现在黄翅飞蝗泥蜂已经穿着打扮完毕，只等待脱掉茧壳了。那是一件精美绝伦的紧身薄膜，把它身体构造的每个细节都一览无余地凸现出来，成虫的形状和颜色几乎都没能遮盖住。在完成蜕变的最后一个动作之前，黄翅飞蝗泥蜂突然从昏睡中苏醒过来，猛烈地乱动着，似乎要唤回那麻木已久的肢体的生命活力。它的腹部一伸一缩，腿猛地伸开后又弯曲起来，接着又伸开，想用力地把各个关节都伸得直直的。昆虫用头和腹尖支撑身体，肚子朝上，使劲地抖动了很多次，把颈关节以及连接腹部与胸部的腿关节撑开。它的努力终于没有白费，这样艰难的体操大概进行了一刻钟后，紧身薄膜终于在四面八方的拉扯下，从脖子处、腿关节周围、靠近腹部肉茎的地方，总之在身体各个剧烈活动的部位波及的地方，都破裂开了。

\newpage

\chapter{高超的理论}

\begin{figure}[htbp]
    \centering
    
\includegraphics[width=0.8\textwidth]{images/00341.jpeg}
\end{figure}

飞蝗泥蜂的种类很多，可大多数品种我们国家都没有。据我所知，这种昆虫在法国只有三种，全都喜欢生活在阳光充沛的、长着橄榄树的炎热地区。它们就是黄翅飞蝗泥蜂，白边飞蝗泥蜂和朗格多克飞蝗泥蜂。可是，观察者会发现一件有趣的事，这三种掠夺者在选择食物时是严格遵守动物学规则的。它们给幼虫选择的食物都是直翅目昆虫，它们选择的三种食物分别是：蟋蟀、蝗虫、距螽。

这三种猎物的外表差异很大，要想把它们归于一类或者发现它们的共同之处，就必须具有经过动物学训练的敏锐眼力，或者至少有像飞蝗泥蜂那样的专业眼光。您不妨对蟋蟀和蝗虫进行一下比较：蟋蟀的头又圆又大，短小粗壮，全身乌黑，后侧大腿上装饰着一条红绶带；蝗虫呈淡灰色，身材修长苗条，头很小，呈锥形，长长的后腿轻轻一蹦就跳跃起来了，并能依靠它那折成扇子形的翅膀来保持这种飞腾。接下来，你们再把这两种昆虫跟距螽做一番比较吧。距螽的背上背着两件乐器，那是两个凹形的蚌壳状的铙钹，发出的声音十分刺耳。它笨重地拖着肥胖的肚子，它的肚子上嫩绿色和奶黄色的环节相间，末端长着一把长匕首。通过对这三者的比较，您一定也会有跟我一样的看法，那就是：各种飞蝗泥蜂选择的食物各不相同，但都属于同一动物学的类别范畴，这需要专业的眼光才能做到，而这眼光是令科学家而并不只是普通人都钦佩不已的。

这种偏爱（这偏爱的范围似乎是某个分类的立法者所确定的，比如某个拉特雷依）实在太怪异了，不知道生活在国外的另一些飞蝗泥蜂是否也捕捉同类猎物，如果能探究一下，将会是很有趣的。可是遗憾的是，这方面的资料十分稀少，甚至连很多同类昆虫的材料也相当匮乏。之所以会有这种遗憾，主要是因为人们普遍采用的方法都过于肤浅。人们抓住一只昆虫后，用一根长大头针把它钉在一个软木底的盒子上，然后把一个写着拉丁名字的标签系在它的腿上，有关这个昆虫的一切研究似乎就完成了。我不想用这种方法来了解昆虫史。人们告诉我，知道某种昆虫的触角有多少关节、翅膀有多少翅脉、腹部或胸部的某个区域有多少根毛，是毫无实际意义的。我只有在了解它的生活方式以及它的本能和习性后，才算真正认识了这种昆虫。

而您会看到，有些细节描述起来十分冗长而且不易理解，如果只用三言两语直接说出来，反而更清楚。比如您如果给我介绍朗格多克飞蝗泥蜂，一定会先描述它的翅膀上有多少翅脉、这些翅脉是如何分布的，然后描述它的肘肋条和回反肋条，接着描述它的外貌特征，介绍它这儿是黑色，那儿是铁色，翅膀末端呈烟棕色，还会说它这里有一块黑丝绒，那里有一块银白色的绒毛，另外的某个部位则是光滑平整的。描述得十分精准、细致，并且可以看得出描述者具有敏锐的观察力和高度的耐心，但是太长了，而且并不能把每个细节都说得很清楚。因此，即使是一个老练的人也免不了会在细节上啰唆起来，这是避免不了的。但是您只要在枯燥乏味的描述中加上"捕距螽"三个字，一切就都清楚了。我也就知道您所说的是哪种飞蝗泥蜂了，并且一定不会弄错，因为只有它是捕捉这种猎物的。要想给人打开这样一道启迪之门，需要做些什么呢？需要真正的观察，而不是把昆虫变成一个个昆虫标本。

这个话题暂且说到这儿，我们现在就来查阅一下关于外国飞蝗泥蜂捕捉对象方面的材料吧！当然，我们迄今为止掌握的相关材料十分匮乏。翻开拉普勒蒂埃·德·圣法古的《膜翅目昆虫史》，我看到阿尔及利亚省的黄翅飞蝗泥蜂和白边飞蝗泥蜂在这个地中海以外的地方仍保留着它们在法国的生活习性。无论是在生长着棕榈树的地方，还是在生长着橄榄树的地方，它们的捕捉对象都是直翅目昆虫。虽然茫茫大海将它们分隔两地，但是这些卡比利亚和柏柏尔的猎手们与它们的普罗旺斯同胞一样，都捕捉同一种昆虫。我在书上还看到了第四种飞蝗泥蜂——非洲飞蝗泥蜂，它们在奥朗日郊区捕捉的是蝗虫。我还记得曾在哪本书上读到过第五种飞蝗泥蜂，它们出没在里海附近的草原上，也以蝗虫为猎物。如此看来，地中海附近共有五种不同的飞蝗泥蜂，它们都捕捉直翅目昆虫来喂养自己的幼虫。

现在，让我们越过赤道走得更远一些，到地球另一半的毛里求斯和留尼汪群岛去看一看。在那儿，我们将看到一种膜翅目昆虫克罗翁。它与飞蝗泥蜂属于同一类，而且样子十分相似。它经常出没在船上和殖民地港口，专门捕食祸害粮食的害虫——卡凯拉克。这些卡凯拉克不是别的，就是蜚蠊属昆虫，就连我们这里的居民也经常被它们当中的某一种骚扰。这种昆虫臭气熏天，大概没人不认识吧！每到夜晚时分，它总靠自己那好像臭虫一样扁平的身体，从家具空隙和壁板缝里钻出来，到所有能让它饱餐一顿的食物那里去。这就是我们房屋里的蜚蠊，它的那副模样令人作呕，可对于克罗翁来说，它可是心爱之物。那么，飞蝗泥蜂的这位近亲为什么会选择卡凯拉克为猎物呢？原因很简单：卡凯拉克虽然样子长得像臭虫，其实却跟蟋蟀、距螽和蝗虫一样，也是一种直翅目昆虫。仅凭借我知道的这些来历不同的虫子也许就能得出结论：所有的飞蝗泥蜂都捕猎直翅目昆虫。即使我们不下这个具有普遍性的结论，也能看出大多数飞蝗泥蜂的幼虫吃的是什么。

它们做出如此奇特的选择，一定是有原因的。可原因是什么呢？它们严格地把日常食物的选择范围局限在同一目昆虫中，但它们是如何决定在这里以一身臭气的卡凯拉克为食，在别的地方以干燥但可口的蝗虫为食，而在另一些地方又以肥美的蟋蟀或距螽为食呢？

我承认我对这个问题一无所知，一点也不明白，所以只好交给别人解决了。不过，我们发现，在昆虫中，直翅目的情况与哺乳动物的反刍类很像，它们天生拥有强壮的肚子和沉稳的性格，在草场上怡然自得地吃着草，肚子长得肥肥大大；它们数目众多，无处不在，行动迟缓，容易捕捉；此外，它们大小适中，是捕猎者的主菜。那么，在反刍类昆虫中，像飞蝗泥蜂这样强劲的、专门捕捉肥壮猎物的掠夺者，是否只能找到像反刍类家畜中的牛、羊那样肉质肥美又性格平和的牺牲品呢？也许是这样的吧，但仅仅是也许而已。

对于另一个同样重要的问题，我有了一个确定的想法，而不再是猜测。以直翅目昆虫为食的消费者们的饮食制度是一成不变的吗？假如它们最喜欢的野味没有了，它们会吃另外一种吗？这个时候，朗格多克飞蝗泥蜂是否固执地认为，除了距螽之外别的东西都不好吃呢？是否摆在白边飞蝗泥蜂餐桌上的只能是蝗虫，而摆在黄翅飞蝗泥蜂餐桌上的只能是蟋蟀呢？或者，它们会根据时间、地点和具体情况的不同，而各自选择别的食物代替最爱的食物呢？如果情况确实是这样的，并且能亲眼看到，那将是一件极有意义的事儿。

因为这些事实将向我们揭示出本能的指引究竟是绝对的、亘古不变的，还是会适时改变的，从而表明狩猎者在选择食物时是否可以自由抉择。

但是，狩猎范围扩大这样的假设是不适用于飞蝗泥蜂的，因为虽然它们完全能在直翅目昆虫中找到许多形状完全不同的昆虫，可是我看到它们总是对固定的一种猎物情有独钟，每种飞蝗泥蜂都捕猎它们各自特定的猎物。不过我有幸发现了一个幼虫食物发生改变的事例，仅此一例，我很愿意把它记录在案。这些经过认真观察的事实，有朝一日会为那些想构建昆虫心理学的人所用，成为他们可靠的第一手资料。

这个故事的发生地是在罗讷河畔的一个防波堤上。防波堤的一侧是波涛汹涌的河水；另一侧是郁郁葱葱的柳条树、杨柳和芦苇，这些树木中间是一条沙石小路。一只黄翅飞蝗泥蜂拖着一个猎物，蹦蹦跳跳地走来了。我发现了什么？这猎物竟然不是蟋蟀，而是一只蝗虫！可这膜翅目昆虫就是我再熟悉不过的那个钟情于蟋蟀的黄翅飞蝗泥蜂啊！我几乎不敢相信自己的眼睛了。洞穴就在不远处，昆虫钻进去，把它的战利品堆放好。我决定坐下来等待它再次出征，如果有必要，就等它几个小时，顺便看看是否有别的捕猎者也像它这么奇怪。我就这么坐着，横挡在这条小路的中央。这时，突然出现了两个光着头的新兵，脸上还充满着稚气，都说兵营生活会把人变成木头人，看来他们就是活脱脱的例子。他们一边聊着天，一定是在谈论同乡的事；一边漫不经心地用刀削着一根柳树棍。我的内心忽然产生了一种恐惧感：在公共道路上做实验可不是件容易的事儿，当您等待了几年的事儿终于要发生时，突然一个过路人闯来了，打乱了一切，毁掉了这个将来很难再有的机会！我忧心如焚地站起来，给两个新兵让路，我整个人躲到了柳树后面，才把那个狭窄的过道空出来。假如我走上前去对他们说："朋友，别从那儿走！"那可就太冲动了，会把事情搞得更糟。因为如果我那样说，他们一定会以为沙下埋着什么陷阱，然后就会向我提问，可我是无论如何也解释不清的。况且如果我这么说了，说不定还会引得这两个闲来无事的人也想看个究竟，而在这样的研究中，他们可是非常碍事的。于是我站起来，什么也没说，一切听天由命吧。咳！我的运气太坏了，他们那沉重的军靴不偏不倚地踩在了飞蝗泥蜂的洞穴上。我不禁打了一个冷战，就像自己被铁靴踩着了一样难受。

两个新兵走了，我开始抢救那个已被践踏成了废墟的洞穴。我在洞里找到了飞蝗泥蜂，它被踩瘸了；洞里除了它和我刚刚看到它拖进去的蝗虫外，还有两只蝗虫；在那个原本用来堆放蟋蟀的地方，共有3只蝗虫。它为什么会有这种异乎寻常的改变呢？是不是因为洞穴附近没有蟋蟀，所以身陷困境的膜翅目昆虫只好用蝗虫来代替，就像谚语所说的"没有斑鸫，将就吃乌鸫"呢？我不大相信这一点，因为没有任何事实可以说明这里缺乏它们最爱的野味。未来某个比我幸运的人将会解决这个未知的新问题。不论黄翅飞蝗泥蜂是出于迫切的需要，还是我不能窥探到的动机，无论如何，这至少证明它偶尔会用另一种猎物——蝗虫来代替它最喜爱的蟋蟀，后者虽然与前者在外表上并无相似之处，但也是直翅目昆虫的一种。

在奥朗日郊区，观察者拉普勒蒂埃·德·圣法古也曾目睹了类似的以蝗虫为食的情况，并根据这一观察对同一类飞蝗泥蜂的习性做了简单的介绍。他是在无意中看到一只黄翅飞蝗泥蜂拖着一只蝗虫的。这是否也跟我在罗讷河边见到的一样，只是黄翅飞蝗泥蜂的一个偶然行为呢？这究竟是个例外还是普遍规律？是不是奥朗日乡下没有蟋蟀，所以膜翅目昆虫要用蝗虫来代替呢？这个问题很关键，我必须提出来，可是我找不到答案。

在这里，我有必要把我从拉科代尔的《昆虫学导论》看到的一段话引用过来，我在后面会提出对这段话的不同看法。这段话是说："达尔文曾专门写了一本书，证明人和动物行动是受着同样的智力原则支配的。有一天他在花园里散步，突然看到小径上有一只飞蝗泥蜂捕捉到了一只个头跟它相仿的蝇。达尔文看到它用大颚咬断了猎物的头和腹部，然后把剩下的连着翅膀的胸部带起来飞走了，但是一阵风吹来，蝇的翅膀动起来，阻碍了飞蝗泥蜂的飞行，使它在原地打起转。于是它被迫又停落在小径上，把蝇的一只翅膀咬断，解决了飞行的隐患后，带着余下的猎物飞走了。显而易见，这事实含有推理的成分。本能会支配这只飞蝗泥蜂咬断猎物的翅膀，然后把它带到窝里去，同类的某些昆虫也是这样做的；但是其中包含了一系列思想和这些思想产生的后果，而这一切，如果它不是受理智的支配，那么就根本无法解释。"

在这个故事中，把理智强加给昆虫实在是过于轻率了，暂且不说是不是真理了，最起码连一点真实性都不具备。这并不在于行为本身，而在于行为的动机。达尔文看到了他说的那件事，只是他把故事的主人公弄错了，把故事本身和故事的意义也弄错了。总之，他完全错了，我可以证明这一点。

恕我直言，这位英国老学者对他如此大肆赞扬的生物应该是十分熟知的。那么我们就用"飞蝗泥蜂"这个词来说吧，这是被严格的科学定义的名词。在他的这个假设中，这只英国飞蝗泥蜂（假设英国确实有飞蝗泥蜂的话）为什么如此怪异，如此背离常规，选择蝇为猎物，而它的同胞们的野味则是与之完全不同的直翅目昆虫呢？在我看来，即使勉强接受苍蝇是飞蝗泥蜂的猎物这个事实，后面的事情也是根本不可能发生的。膜翅目掠夺者提供给它们幼虫的，不是尸体，而是一只仅仅被麻醉了的处于麻木状态的猎物，这已经是毋庸置疑的了；那么，被飞蝗泥蜂切断了头、腹部和翅膀的猎物对飞蝗泥蜂来说到底意味着什么呢？意味着它带回去的猎物只不过是一块死肉，它的腐臭味会充满整个蜂房，这对幼虫没有任何意义，可是孵化很快就要开始了。所以达尔文观察到的，不是真正的飞蝗泥蜂。那么，他看到的是什么呢？

它捕捉到的猎物被称为蝇，而蝇是个十分笼统的概念，双翅目昆虫中的大多数都可以这么叫，所以在前几种昆虫里，我们不知它指的究竟是哪种。飞蝗泥蜂这个词也许也被用在了同样不确定的地方。在20世纪末，达尔文的著作充斥了整个时代，人们用这个字眼不仅指严格意义上的飞蝗泥蜂，也指方头泥蜂科昆虫，其中有些昆虫捕捉双翅目昆虫蝇作为幼虫的食物，它们是连英国博物学家都不认识的膜翅目昆虫的猎物。那么，达尔文说的飞蝗泥蜂是一种方头泥蜂吗？也不是，因为捕捉双翅目的昆虫猎手和捕捉其他野味的猎手一样，在卵的孵化和幼虫完全发育的半个月或三个星期中，要求它们的猎物保持着新鲜、一动不动的半活状态。所有这些小贪吃鬼要的都是新鲜的肉而不是腐烂发臭的肉。这是我发现的从未有过任何例外的一条规律。因此，飞蝗泥蜂这个词甚至不能只从它原来的意义上去理解。

所以，上述那本书中所说的并不是真正符合科学的事实，而是一个需要破解的谜。我们继续来猜测谜底吧。很多种类的方头泥蜂在体形、形状以及黑黄相间的外衣方面，都与胡蜂十分相似，因此如果不是对昆虫学上的细微差别十分在行的人，是根本分辨不出来的。在任何一个没有专门研究过此类问题的人看来，方头泥蜂就是胡蜂。这位英国观察者以一种居高临下的态度看待事物，对于这个即将证实他的理论观点是多么卓越超群的并使理智与昆虫相结合的事实，他觉得是微不足道的小事儿，没有必要去仔细考察。那么他会不会也犯了错误，犯了一个情有可原的但与事实完全相悖的错误：把胡蜂错当成方头泥蜂了呢？我几乎要坚信不疑了，以下就是我的理由。

虽然胡蜂不是永远，但至少经常以某种昆虫为食物来喂养它的幼虫；它并不是提前把猎物全部堆放在蜂房里，而是每天分很多次给幼虫们喂食；它用嘴一口一口地喂，就像母鸟喂雏鸟一样。喂食时，雌胡蜂先用大颚把昆虫捣碎，然后咀嚼成细酱；而这种颇受小家伙喜爱的肉酱的原材料就是双翅目昆虫，尤其是普通的蝇；如果有新鲜的肉，那就是意外收获，它一定会充分利用。谁没看到过胡蜂胆大包天地钻进我们的厨房，或者突然扑到肉店的砧板上，叼起一块合意的肉迅速飞走，把这美味的战利品喂给它的幼虫呢？当一缕阳光透过半闭的百叶窗洒在房间的地板上，苍蝇在阳光的照耀下甜美地睡着午觉或者掸着翅膀上的灰尘时，谁没见过胡蜂突然闯进来扑向双翅目昆虫，用大颚把它咬碎后带着这块野味逃走呢？这又是食肉类婴儿的一顿美食。

得到美食后，它们有时在掠夺地上，有时在路上，有时要带到窝里将其肢解。苍蝇的翅膀是毫无价值的，被切断扔掉了；腿里的汁液很少，有时也被抛弃了；剩下的是或者连在一起或者分割开来的尸体：头、胸、腹。胡蜂把它细细地咀嚼一遍，做成一盘美味的羹，供幼虫享用。我曾用一盘苍蝇酱来代替这种羹，喂我饲养的昆虫幼婴。我的实验对象是一个胡蜂窝，那是一个建造在一个灌木丛枝丫上的小蜂房，好像灰紫色玫瑰花萼的形状。我用一块大理石板当烹饪器具，把野味洗净，去掉翅膀、腿等坚硬部位，然后在大理石板上做起了苍蝇酱；最后，用一根细麦秸当匙，把美食放在麦秸上，挨个地给蜂房里每一只好像雏鸟一样的半张着大颚的婴儿喂食。我小时候喜欢饲养小雀儿，可那时也没有现在干得这么起劲儿、这么好。就这样，在我的耐心坚持下，一切都进展得很顺利，而在这聚精会神和细致周到的喂食过程中，我的耐心受到了严峻的考验。

\newpage

\chapter{朗格多克飞蝗泥蜂}

\begin{figure}[htbp]
    \centering
    
\includegraphics[width=0.8\textwidth]{images/00351.jpeg}
\end{figure}

化学家在认真地制订了研究计划后，就在恰当的时刻搅拌他的反应剂，同时点燃曲颈瓶下的火。他是时间、地点、环境的主人。他可以自由选择工作的时间，然后躲在与外界隔绝的实验室里，不会受到任何打扰；他随心所欲地制造出一切他想要的环境；他探寻着无机的自然的秘密，只要愿意，他可以在任何时候让它产生化学作用。

活的自然秘密则不是解剖学上的关于结构的秘密，而是鲜活的生命，尤其是本能的秘密，所以这对观察者来说，困难要更大、更微妙，人们不但无法自由支配时间，还要受到季节、月份、小时乃至于每分每秒的约束。机会一旦出现，就该毫不犹豫地抓住，因为这机会也许以后都不会再有了。加上这样的机会往往是在不经意间来到的，所以人们根本毫无准备。除非您立刻准备好一个小型实验器材，然后制订计划、设计战术、想出妙计，如果您的灵感来得相当快，您没有错失良机，那么您就太幸运了，因为这机会通常只出现在极力寻找它的人眼前。您必须耐心地日复一日地等待，有时蹲在烈日暴晒的沙坡上，有时在夹在陡坡之间的好似烘箱似的小路上等待，有时爬到砂岩的陡壁处，而那里人迹罕至，会令您不免胆战心惊。

假如您有办法把观察站设在一棵橄榄树下，它稀疏的叶子能够为您挡烈阳，那您就要感谢命运的恩赐了，是命运让您过上了如此美妙的日子，您得到的简直是一座伊甸园。您一定要仔细盯着，这个地方很棒，而且说不定哪一刻机会就来了。机会来了，是的，有点晚，但终究是来了。啊！如果现在能独自一人坐在宁静的书房专门研究自己的课题，没有过路人打扰，该多好啊！那些过路人见您聚精会神地盯着一个点，他却什么也没看到，就会停下来向您问东问西，把您当作拿着榛树魔棒探索宝藏的人；或者疑心更大者，把您当作图谋不轨的人，以为您正用咒语从地下寻找装满钱的旧罐子！如果您在他心目中有着基督徒的样子，那他会崇拜您，您看哪里，他也看哪里，并且面带笑容，那样子，简直是他正怀着一副仁慈之心端详着您这位全神贯注观察苍蝇的人。如果这个令人厌恶的参观者，尽管在内心中对您嗤之以鼻，却终于走开了，而没有破坏您的观察现场，没有把我遇到的那两个新兵的鞋跟带来的灾难扩大，您就算幸运了。

您忙碌着，可您又无法说出您忙的事儿，这种情形即使不会令过路人满腹狐疑，也会引起乡警的怀疑。乡警是乡村里的法律代表，是最不讲情面的人。

他已经监视您很久了，他经常看到您无缘无故地在这儿走走，到那儿逛逛，好像心事重重的样子；他还经常见您在地上寻找，小心翼翼地在一处凹陷的道路上拍打着沟壁，于是，他就怀疑您了。他认定您可能是吉卜赛人、流浪汉、可疑的闲逛者、偷庄稼的人，至少是个怪人。如果您带着植物标本箱，会被他误认为是偷猎者用来装白蚰的箱子，您根本无法改变他的看法，他认为您根本没把狩猎法和所有权放在眼里，并且认为您打算把附近所有的兔子都捕光。您可要小心啊，您就算口干舌燥，也千万不要伸手去摘身旁的葡萄，因为那个佩戴着乡警牌子的人可能就在那儿，他会觉得自己很幸运，终于能将您这个使他迷惑不解的人记录在案了。

我可以这样说，我从没干过这样的坏事。可是有一天，我正趴在沙土上全神贯注地观看一只飞蝗泥蜂操持家务时，突然听到身后传来一声呵斥："我以法律的名义命令您跟我走！"这是英国人的乡警，他一直在找我的把柄，可又找不到，他非常想解开心中的谜团，于是决意对我提出粗暴的警告。我只好进行解释，但这个可怜的家伙根本不信。"嘿！嘿！"他说，"您说您到这儿来烤着太阳就是为了看苍蝇飞，我是无论如何也不可能相信的。我一直盯着您呢，这您是知道的！如果再让我看到您这样，我可就真带您走了！"他走了。我始终觉得，他是因为看到了我的红绶带才走开的。在我为昆虫学或植物学而远征的途中，以及其他的一些同类小事上，我也认为是这条红绶带帮了我的忙。我觉得，即使事实并不一定是那样的也差不多，但是当我到万杜山采集植物标本时，那个带路人比村警好相处，那头驴子也比较温驯。

但是这条猩红色的小带子，并不能把昆虫学家在公路上做实验时遭受的全部磨难都消除掉。让我们举一个典型例子吧。有一天，天刚蒙蒙亮，我就躲藏在了一条小山谷的石头上，准备探访我的观察对象——朗格多克飞蝗泥蜂。三个前去采摘葡萄的女子从那里经过时，扫了一眼坐在那里沉思的人，并礼貌地问了一声好，对方也给予了礼貌的回答。太阳落山时，那几个收葡萄的女子又从那里经过，头上顶着的篮子里装满了葡萄；而那个沉思者仍坐在那块石头上，眼睛一眨不眨地盯着同一个地方。我纹丝不动地在这个荒无人烟的地方待了这么久，一定令她们十分奇怪。当她们从我面前经过时，我看到她们当中的一个人用手指指着额头，小声地对其他人说："一个不会害人的傻瓜，太可怜了！"然后三个人都在胸前画了一个十字。

她说"傻瓜"，是的，一个"傻瓜"。一个傻瓜，一个不会伤害别人但已经失去了理智的可怜人，所以这些女子都心怀怜悯地画了个十字。对她们来说，傻瓜是被主打上了印记的人。说什么呢！我心想，这简直是糟践人。您专心致志地在昆虫身上探寻着关于本能和理智的东西，可在这些善良的女子心目中，您就是一个失去了理智的人！这真是莫大的侮辱！算了，不去想了。在普罗旺斯语中，"可怜"这个词代表了最高程度的怜悯，这发自内心的"可怜"使我很快忘掉了"傻瓜"的称谓。

如果我刚才讲的这个小小的遭遇，没有破坏读者的兴致，没有让您觉得扫兴，那么我想邀请您到这三个采葡萄的女子经过的那个峡谷去看一看。朗格多克飞蝗泥蜂在造窝时，不是从某个地点成群结队地来，而是形单影只、稀稀拉拉地来，并且在长途跋涉的流浪过程中，随遇而安地来到某个地方就把家安在那里。它的同行黄翅飞蝗泥蜂总是寻找那些有同伴聚集的、热闹的劳动工地，它却更享受孤独，喜欢一个人居住的寂静。它的步态越来越稳重而谨慎，它的身体越来越强壮，衣着也越来越暗淡，它总是离群索居而不在乎别人在干什么，甚至对同伴也不屑一顾，是飞蝗泥蜂族中的真正的卓然自立者。前者善于群居，后者喜欢独居；从这个明显的区别中，就足以看出它们各自的特性了。

这也说明，要观察朗格多克飞蝗泥蜂难度更大些。对于这种飞蝗泥蜂，不会有经过长久谋划的实验，因为当第一次尝试失败后，就不要再指望在同样的情景下对第二只、第三只无休止地实验下去。假如您事先准备了观察器材，假如您准备了一块猎物，想用它代替飞蝗泥蜂的猎物，那么恐怕，或者可以肯定，捕猎者不会出现；而当它终于出现时，您的器材已经无法使用了。所以，一切都要在它出现时的一瞬间匆忙迅速地准备好，可我不完全具备这些能力。

我们应当相信，观察地点是对的，因为我已经连续几次在这些地方看到飞蝗泥蜂在阳光普照的葡萄叶上小憩。昆虫平躺在叶子上，美美地享受着阳光带给它的乐趣，偶尔还会嗡地叫上一声，一副乐不可支的样子。它舒服极了，不禁扭动起身子，用脚尖快速拍打身下的叶子，发出击鼓一样的声音，就像一阵狂风暴雨猛打着树叶似的。这种欢快的击鼓声，传到了几步路外的地方。接着它一动不动了，但是很快又发出了一阵乱舞和弹动，尽情地表达着它的快乐。我太了解这些阳光的狂热爱好者了。它们给幼虫建的窝刚挖了一半，就突然丢下工作，到附近的葡萄架上晒太阳，然后极不情愿地回到窝里，粗枝大叶地扫一扫，最后还是决定弃工而去了，因为它实在抵挡不了葡萄叶上的快乐对它的诱惑。

也许这个舒适的休息场所还是一个观察站，膜翅目昆虫正在那儿环顾四周，认真寻找着它的猎物。其实它想找的野味是吃葡萄的距螽，这些距螽分散在葡萄藤或者随便哪个荆棘丛上，而飞蝗泥蜂又独独喜欢那种肚子撑得鼓鼓的母距螽，因为这样的母距螽的肚子里塞满了虫卵，绝对够肥美。

我们不必多说那些关于反复奔波、徒劳探究和长久等待的事儿了，飞蝗泥蜂是怎样出现的，我们就怎样向读者介绍吧。看！飞蝗泥蜂出现在凹陷的道路上，道路两侧是高耸的陡坡。它徒步走来，挥动着翅膀，拖着沉重的捕获物。距螽的触角像一条细长的线，而对捕猎者来说，这正好可以用来当套车绳。飞蝗泥蜂昂着头，用大颚咬着一根从它腿间穿过的触角，它拖着的猎物仰面朝天。假如地面坑坑洼洼，不利于它这样运输，膜翅目昆虫就会抱起那庞然大物飞上一小段路程，其间只要遇到路面平坦的地方，它就步行。我从未见过它像那些灵活的、善于长途飞行的昆虫一样，双腿夹着猎物飞很远的距离。那些昆虫，比如泥蜂和节腹泥蜂，前者抱着双翅目昆虫，后者抱着象虫，在空中也许能飞上1千米，当然，它们的猎物与笨重的距螽相比要轻得多。因此，朗格多克飞蝗泥蜂在猎物沉重的事实面前，不得不在整个或者大部分路程中，采用那种既缓慢又困难重重的徒步运输的方式。

通常情况下，膜翅目掘地虫是先挖洞然后再供应粮食的，可是同样碍于猎物又大又重的原因，它的工作程序打乱了。掠夺者有足够的力气搬动猎物，膜翅目昆虫飞行运输的方法十分便利，可以随意选择住所位置，猎手也完全可以到很远的地方捕猎。它俘获了猎物后，就迅速飞回家，远近对它来说都没有太大影响。所以，它宁愿把它的诞生地——那个它的父辈们曾生活过的地方当作它的家。那个窝里有已经挖好的深深的巷道可以继承，那是几辈人辛勤工作的结晶；它把那些巷道稍加修整，作为连通新卧室的通道，就加固了卧室的安全性，比自己每年重新从地面开始挖掘的更牢固。节腹泥蜂和捕食蜜蜂的大头泥蜂就是这样的。如果父辈的老屋不够牢固，不能年复一年地抵挡风雨侵袭并且代代相传，膜翅目掘地虫就必须每年都依靠自己重新挖洞，那么，至少膜翅目昆虫认为它挖的洞比先人的住所更安全。所以它挖自己的巷道，把每条巷道作为通往蜂房群的走廊，减少了孵卵所需的劳动量。

这样的工作方式产生的尽管不是真正意义上的社会（因为它们没有出于共同目的的协调劳动），但至少是一些聚居的中心点。昆虫在那里能看到自己的同类、伙伴，这样，每个人都会更加干劲十足。确实如此，我们发现，同宗的小部落一起劳动时，与孤独的掘地虫单独劳动相比，积极性是不同的：整群的昆虫在一起，就像在万头攒动的工地上工作一样，一派热火朝天的景象；掘地虫独自工作时却孤独无聊，动作迟缓。昆虫跟人一样，是会互相带动和感染的，能够被榜样激励。

由此可以得出结论：对于掠夺者来说，如果猎物轻，它就可以做长距离的飞行运输，这样，膜翅目昆虫就可以随意选择建造洞穴的地点。它更喜欢到它的诞生地上，把那些老房子中的走廊修葺成通往若干个蜂房的过道。因为出生地点十分接近，所以造成了它们同类之间的聚居，它们毗邻而居，彼此成了劳动的动力。这种生活的出现首先是由运输容易决定的。让我们这样比喻吧，人类不也是这样的吗？在难以通行的山路边，人们建造的茅屋都是零零散散的；如果在便于行走的大路边，人们便容易聚居在一起进而形成一个人口繁多的城市；如果有铁路来消除距离，人们便集结在类似伦敦、巴黎这样庞大的蜂窝里。

朗格多克飞蝗泥蜂则截然不同。它的猎物是沉重的距螽，仅这一只距螽就与别的掠夺者飞行多次所得的食物总量不相上下了。节腹泥蜂和别的飞行速度快的掠夺者需要分期完成的工作，它只需一次就可以完成。它带着沉重的猎物无法做长途飞行，只能不辞辛劳地依靠脚步的力量把猎物慢慢拖回家。仅凭这一点，住所的位置就不得不由猎物的捕获地决定。也就是说，它是先有了猎物，然后才有了住房的。这样一来，它们就不可能有可供共同选择的聚会地点了，也不可能有同类相邻而居，更没有各个部落在劳动中竞相表现从而互相激励的情况了。朗格多克飞蝗泥蜂的家总是随处安置的，然后它就在那里寂寞独处，虽然始终都兢兢业业，可是一个人干活儿总是打不起精神。掘地虫先是找到猎物，然后发动进攻，接着把它麻醉，最后才考虑筑窝的事。它会在尽可能靠近捕获猎物的地方选择一处称心的地点，然后迅速挖好送给未来幼虫的卧室，准备立即迎接卵和食物的到来。这就是我的观察结果，下面我将摘要加以介绍。

我看到朗格多克飞蝗泥蜂总是独自在挖洞，有时待在古墙下被滚石砸落的灰沙堆里，有时待在一个由凸出的砂岩形成的隐蔽所里，而那恶狠狠的单眼蜥蜴正打算通过这个像前厅一样的隐蔽所侵入它的巢穴。这里阳光普照，像烤箱一样闷热。土地十分易于挖掘，因为它是由拱顶逐渐脱落的古老灰尘形成的。它用大颚当铲子挖土，用跗节当耙子清扫，房间很快就挖好了。于是，飞蝗泥蜂这种掘地虫飞起来了，可是飞得很慢，没有扇动它那强劲有力的翅膀，因为它并不打算做长途飞行。我们可以用眼睛追踪到它的身影，看到它落在了大约10米远的地方。有的时候，它如果决定步行远走，就会风风火火地离开洞穴朝某个地点走去，即便我们莽撞地跟上去，它也一点儿不会被我们影响。它或者步行或者飞行到想去的地方，从它那迟疑不决的步态和顾盼流离的表情中可以断定，它在寻觅着什么。它寻觅着，终于找到了，或者应该说是重新找到了。它重新找到的是一只已经处于半麻醉状态但跗节、触角、产卵管还在动的距螽。这一定是前不久被朗格多克飞蝗泥蜂刺了几下的那个猎物。因为带着重物到处寻找住所十分不方便，所以它对猎物动了麻醉手术后就离开了；有时候，它就把猎物扔在捕猎现场，或者放在某个显眼的草丛里，以便将来容易找到；它对自己的记忆力十分自信，坚信自己过一会儿还能再回到放置战利品的地方。接着它就开始四处寻找称心的挖洞点了。住所挖好后，它就立即去找猎物，没费什么力气就找到了；现在它正准备把猎物运到住所去。它跨在猎物身上，抓住猎物的一条触角，有时也会把两条触角同时抓住，然后依靠大颚和腰部的力量拖着它上路。

有时，这个搬运工会一口气跑完全程；但更多的时候，它会突然把重物丢在中途，独自飞快地跑回家。也许是它突然想起大门的入口不够宽，这个庞然大物进不去；也许是它忽然想起某些细节上的缺陷会影响食物的储存吧。果然如此，这位工人在修缮它的房子：扩大入口的宽度，铲平门口的道路，把拱顶加固一些。这一切只需要用跗节拍打几下就可以了。做完这些事以后，它才跑回去找丢在半路上的距螽，那距螽就仰面躺在离它几步远的地方。搬运又开始了。刚走了不远，飞蝗泥蜂那灵活的小脑袋似乎又想起了一件事：只是忙着检查大门了，却忘记去室内看一看。里面是否一切正常呢？于是，它又一次把距螽扔在半路上，独自跑回家。对室内进行了一番检查后，它还会顺便用跗节这把镘刀抹抹四壁，做最后的修补。这些细致的修整没有耽误膜翅目昆虫很多时间，它很快就又返回猎物这里，抓起它的触角前进。这一次会走完全程吗？我不敢保证。我曾见到两只比它更多疑的飞蝗泥蜂，也许它们对于建筑上的小事儿更健忘吧，为了消除疑虑，它们曾经把战利品扔在路上五六次，自己跑回洞里五六次，或者做一些小修补，或者只是到洞里看一眼。当然，也有的飞蝗泥蜂是带着猎物径直走到洞里的，一路上几乎连歇都不歇一下。我在这里还要补充一句，膜翅目昆虫在跑回去修补住所的时候，它会频频地回头张望那丢在半路上的距螽，看看是否有人去招惹它。它如此谨小慎微，不免令人想起了圣甲虫，谨慎的圣甲虫不时地从它那尚未挖完的大厅里出来，拍打拍打珍贵的粪球，每出来一次都把粪球推得离自己更近一些。

从上面叙述的事实中，我们可以很容易得到一个结论：任何一只朗格多克飞蝗泥蜂，无论是开始挖掘时，还是用跗节初步清扫尘土时，还是住所已经挖好时，它都会时而步行、时而飞行地进行一场短距离的远征，以便更好地拥有那只已被蜇刺到麻醉状态的猎物。由此我们可以信心十足地得出一个结论：膜翅目昆虫首先扮演的是猎手的角色，然后才是挖掘工的角色；它住所的建造地点完全是由捕猎的地点决定的。

\newpage

\chapter{本能的技能}

\begin{figure}[htbp]
    \centering
    \includegraphics[width=0.8\textwidth]{images/00357.jpeg}
\end{figure}

我毫不怀疑，朗格多克飞蝗泥蜂在麻醉它的猎物时，跟捕猎者在捕获蟋蟀时用的办法一样，即将蜇针刺入距螽的胸部好几次，以保证刺中它的胸部神经节。也许它早就掌握了刺伤猎物神经中枢的方法，而我也早就承认它的手术做得十分高明，动作既熟练又敏捷。所有的膜翅目强盗都对这种技法了然于胸，它们的毒针可不是毫无用处的摆设。不过我得承认，我还没看见过它的这种谋杀壮举，而这种遗憾完全是因为朗格多克飞蝗泥蜂的习性太孤僻了。

在黄翅飞蝗泥蜂窝的聚集地，它们总是先挖好窝，然后再放上食物，您只要在那儿等一会儿，就能看到带着猎物纷纷归来的捕猎者们。这时，您可以轻而易举地用一只活猎物代替它们运来的准备储存起来的猎物，并且只要您愿意，重复多少次这样的实验都可以。因为您无论在任何时候都能遇到您的观察对象，所以完全可以事先把一切都准备就绪。可是要观察朗格多克飞蝗泥蜂，情况就不同了，那些成功的条件都不具备。您如果事先准备好器材然后专门去寻找它，那几乎是徒劳无功的，因为习性孤独的昆虫总是一个个地消失在广袤的大地上；而且就算您有幸遇到它，也通常是在它无所事事的时候，所以您从它那儿还是什么也得不到。我们再说一遍，朗格多克飞蝗泥蜂几乎总是在我们不经意或毫无准备的情况下拖着它的距螽出现的。

当那个有利时机到来时，我们可以用替换猎物的方法从捕猎者那里探知它使用蜇针的秘密。这个唯一有利的时机来了，那就让我们迅速准备好一只活距螽当替代品吧！赶快，时间紧迫，如果再拖延几分钟，食物就要被送进窝里了，就要与天赐良机失之交臂了。这时还有必要描述我是如何埋怨自己运气不好，如何埋怨自己没有事先准备一只小小饵物吗？当我日思夜想的观察材料出现在自己眼前时，我却眼看着它溜走！我没有可以用来代替飞蝗泥蜂猎物的东西，无法从它那儿换取秘密！那么，请您想一想，您只有几分钟的时间，却要四处寻找飞蝗泥蜂猎物的替代品距螽，而我们要找到距螽至少需要三天的时间！这是根本办不到的事，可我却尝试了两次。啊！如果乡警看到我这时正在葡萄树下疯也似地跑着，那么对他来说简直是个千载难逢的好机会，他终于可以抓到一个偷农作物的人了，并且终于可以记录下他的口供了！我急得四处奔走，被树藤绊住了，可我才没空理会这些葡萄藤和葡萄串呢！我不惜一切代价只为得到一只距螽，并且要立即得到。我在如此匆忙进行的远征中，曾经得到过一只距螽。我当时喜出望外，可没想到，等待我的竟然是痛苦和失望。

如果我能及时赶回来，如果朗格多克飞蝗泥蜂还在忙着搬运猎物，那我就成功了！谢天谢地！一切都如我所想。膜翅目昆虫还在拖着它的猎物前进，离窝还有一段距离。我用镊子轻轻地从后面夹起它的猎物往后扯，遭到了猎手的抵抗，它乱舞着触须，不肯放弃。我用力一拉，使得猎手倒退了几步，可还是无济于事，飞蝗泥蜂还是不肯松口。还好，我随身带着一个小剪刀，这是我昆虫学小行囊中的必备物品。我用剪刀迅速一剪，剪断了距螽那像绳子一样的长触角。飞蝗泥蜂继续朝前走着，但突然停了下来，因为它惊奇地发现重物的重量变轻了很多。确实如此，它现在拖着的是已经被我巧妙剪断的触角，真正的重物——那个身宽体胖的猎物还在后面，并且已经被调了包，被那个活虫子替代了。膜翅目昆虫回转身，丢下光溜溜的触须，按原路返回。它来到替代品前，审视着这个猎物，半信半疑地把它翻过来，然后停下来，用唾液沾湿一条腿，擦了擦眼睛。它或许正在思索："天啊！我是人老眼花了还是睡着了？那东西根本不是我的。到底是谁、是什么东西骗了我？"不管怎样，飞蝗泥蜂并没有急着用大颚咬我的猎物，而是站在一旁，好像根本不打算捕捉它。为了刺激它，我用手指把昆虫捏到它跟前，我甚至试着让昆虫的触须触碰它的牙齿。我了解它是个勇敢的家伙，我知道它会毫不犹豫地从我的手指上把这个替代品夺走的。

怎么？飞蝗泥蜂对我献上的食物竟然不以为然，非但没有去咬我放在它眼前的东西，还不停地往后退去。我又一次把距螽放在地上，摆在它跟前。距螽这时已经一动不动了，对危险毫无意识。我们成功了？唉！没有，飞蝗泥蜂居然像个懦夫一样，继续后退，直到最后飞走了。从那以后，我就再也没有看到它。这个令我无比兴奋和激动的实验，就这样不了了之了。

后来，在我观察了更多的洞穴之后，我终于慢慢明白了失败的原因，也明白了为什么飞蝗泥蜂会断然拒绝我献上的猎物。我在供应食物时，无一例外地，总是捕捉雌距螽，因为它装着一肚子丰富美味的卵，这大概是幼虫最喜欢的食物了；而我在葡萄树下匆匆忙忙地寻觅时，捉到的是另一性别的。我放在飞蝗泥蜂眼前的是雄距螽，而在食物这么重要的问题上，膜翅目昆虫要比我敏锐得多，所以它拒绝了。"这就是给我幼虫的晚餐？把我当成什么了！"这位精明的美食家的眼光太敏锐了，它能辨别出雌性的肉比较细嫩，而雄性的肉相对比较粗糙！它的眼光多么犀利啊，雌性和雄性昆虫在形状、颜色上是一样的，可它竟能立即识别出来！雌性昆虫的肚尖上带着刀，那是产卵管，专门负责把卵产到地洞里；毋庸置疑，这是从外表上能将它与雄性区分开来的唯一特征。这个特征是瞒不过飞蝗泥蜂那敏锐的眼睛的；也正因为如此，当我实验中的那只膜翅目昆虫看到猎物时，揉了揉眼睛，百思不解：当初我抓到的分明是长着刀的，现在刀子没了。面对这种突如其来的变化，飞蝗泥蜂那个小脑袋里会想些什么呢？

现在让我们再来看看膜翅目昆虫吧。它建造好住所以后，就要把丢在不远处的那只已被它动了麻痹手术的猎物找回来。这时的距螽与被黄翅飞蝗泥蜂麻痹的蟋蟀的状态一样，这有力地说明了它也是胸部被蜇刺。不过猎物还能动，还有一定的活力，只是不能全身协调运动。昆虫因为无法站立起来而侧躺着或仰躺着。它快速地摆动着长长的触须和触角；它的大颚一张一合，啃咬的力量跟正常时一样大；腹部深深地一起一伏着；产卵管突然缩到了肚子下面，差不多贴到肚子上了；它的腿还在动，但只是懒洋洋地乱踢乱蹬着，腿部中间好像比其他部位麻痹得更严重。如果用针尖来刺激它，它会全身乱抖起来，拼命想站起来却不能。总之，除了不能站立以外，昆虫还是充满生命力的。因此它的麻醉完全是局部的，只是腿的某些部位不能正常运动了。关于这种并未完全丧失活动能力的原因，是猎物神经系统的某种特殊性决定的，还是膜翅目昆虫只蜇了它一下，而不是像蟋蟀的捕猎者那样刺遍了猎物胸部的每个神经节呢？我不能确定。

虽然猎物在颤动着、抽搐着、毫无章法地乱动着，但以它目前的情况来看，对食用它的幼虫是构不成任何危害的。我曾经从朗格多克飞蝗泥蜂的窝里取出了一只距螽，它十分有力地挣扎着，就像刚刚被局部麻醉似的；而刚刚孵出来几小时的小幼虫，还十分虚弱，却可以安全无忧地啃咬这只庞然大物，就像一个侏儒在啃噬一个巨人似的。这一切都得益于幼虫母亲对产卵点的选择恰到好处。我前面说过，黄翅飞蝗泥蜂把卵产在蟋蟀胸部靠边一点上，在第一对腿和第二对腿之间。白边飞蝗泥蜂蜇刺点的选择也大致如此；而朗格多克飞蝗泥蜂选择的产卵点则稍稍靠后一些，在一条后大腿的根部附近。这种一致性说明，这三种飞蝗泥蜂都具有令人敬佩的高超本领，能判断出卵产在哪里才安全。

现在，让我们看一看被囚禁在洞里的距螽吧。它仰面朝天地躺着，根本翻不了身，拼命地挣扎着、扑腾着，腿在空中乱踢着，可是房间太小了，这些腿无法用墙壁作为支撑点。猎物的抽搐对于小幼虫来说根本没有危险可言；幼虫所处的位置无论猎物的跗节、大颚、产卵管还是触须都碰不到，可以说是一个完全一动不动的、连皮肤都丝毫不会颤动的部位。只要距螽不能移动、不能翻身、不能站立，这里就是永远安全的，而这些条件是已经具备了的。

假如同时有好几个猎物，而麻醉的程度又不能更深一些，那么幼虫就有危险了。当然，对于幼虫来说，它第一个要进攻的昆虫是绝对安全的，因为它所处的位置是不会受到这只猎物攻击的；但它要格外小心旁边的其他猎物，这些猎物偶尔伸伸腿就有可能将腿上的刺戳到它的肚子里。黄翅飞蝗泥蜂把三四只蟋蟀堆在同一间蜂巢里大概就是出于这个原因吧，因为这样一来，它的那些猎物都几乎一动不能动了。至于朗格多克飞蝗泥蜂，它在每个洞里只放一只猎物，所以它允许距螽的身体大部分可以动，只是不允许它们移动位置或站立起来，这可以使它节省很多毒液，但是对于这一点我还无法证实。

如果说，处于半麻醉状态下的距螽对于放置在它身上几乎毫无行动能力的幼虫并非毫无危险，那么对于朗格多克飞蝗泥蜂这个将要把它搬运到住所的猎手来说，就不是没有危险了。首先，猎物几乎还有使用跗节的能力，它在被拖运的时候，会突然抓住路上的某个草茎，从而极大地阻碍搬运的顺利进行。庞大的重物已经把飞蝗泥蜂压得气喘吁吁了，在草茎多的地方就会更加疲惫不堪，结果因为猎物牢牢地抓住了什么东西只好无奈地放弃。但这个麻烦简直不足挂齿。距螽的大颚完全正常，咬起来跟平常一样有力。当掠夺者处于搬运的姿势时，面对猎物那可怕的大钳子的正是纤弱的身体。飞蝗泥蜂抓在了靠近猎物触须根部的位置上，而此时，肚子朝上的猎物正用大嘴巴对着飞蝗泥蜂的胸部或腹部。飞蝗泥蜂挺直它那长长的腿，昂首向前，我坚信它对身下那半张着的大颚一定会格外小心；但是它一旦稍有疏忽，一步失足，或因一点小小的意外，就会被那两把强劲有力的钳子碰到，到那时，这钳子是不会失去这个报复的好机会的。所以，即使不是在所有的情况下，但至少在遇到某些困难的情况下，应当消除这些钳子带来的危险隐患，消除猎物腿部的钩子可能给运输增加的阻力隐患。

飞蝗泥蜂怎么才能做到这一点呢？在这个问题上，即使是人类的专家也会无计可施，会因为实验毫无结果而茫然无措，甚至会觉得根本无法做到。向飞蝗泥蜂学习吧。它没有人教，也从没见别人做过，就完全掌握了手术的技巧。它知道或者说它好像知道神经生理学最深奥的秘密。它知道猎物的头颅下有一环好像高级动物的大脑那样的神经核。它知道嘴部的活动是受这个神经分布的源头支配的，还知道这是神经中枢，只有这里发出指令，肌肉才会活动。最后它还知道，通过破坏这类神经就能制止猎物的一切反抗，因为那昆虫已经不再有抵抗的愿望了。至于手术的方式，对于飞蝗泥蜂来说就更容易了，只要我们向它学习，完全可以试试它的方法。但在这里，飞蝗泥蜂使用的不再是蜇针，它做出了一个聪明的决定，用按压的方式代替毒刺。我们要向它的英明决定致敬，因为我们很快就会看到，在昆虫的智慧面前人类是多么的无知，而承认自己的无知又是多么的明智。我观看了一场精彩的表演，并当场用铅笔做了记录，我怕另做描绘反而无法把这位手术大师超群绝伦的技法表现出来，所以就把这篇笔记原原本本地抄了下来。

当飞蝗泥蜂感觉到猎物抓住草茎拼命抵抗时，停了下来，开始进行奇怪的手术，就像是要给它狠狠一击彻底结束它的痛苦似的。膜翅目昆虫跨在猎物身上，用力把猎物颈背处脖子的关节扳得很大，然后用大颚咬住它的脖子，并尽可能往前在它的头颅下搜索，在外部却没有留下任何伤痕，就这样，它咬住了猎物的脑神经节，不停地压迫。做完这个手术后，猎物就一动不动，不再有任何反抗了。而在此之前，它的腿虽然不能做出像走路时那样的协调运动，却还是会为了不被拉走而狠狠地拖拽住路边的什么东西。

很明显，飞蝗泥蜂用颚尖在猎物的头颅下搜寻并压迫它的大脑，但是不会损伤它纤细柔软的颈膜。没有伤口也没有流血，它只是在猎物体外压一压。当然，我把那个手术后一动不动的距螽收藏了起来，以便有空时查看它的手术结果。同样要说明的是，我立即用飞蝗泥蜂刚刚教我的办法在活距螽身上进行了反复实验。在此，我要把我的实验结果和膜翅目昆虫手术的结果做一下比较。

我用镊子夹着两只距螽，按压它们的脑神经节，它们果然很快就与飞蝗泥蜂的猎物一样了。只是当我用针尖刺激时，它们会发出刺耳的声音，腿还会懒懒地动一动。毋庸置疑，这种区别是因为我的手术对象的胸部神经节在此前没有受过伤害，而飞蝗泥蜂的距螽的胸部曾被刺过。除了这个重要的区别外，我还勉强算得上是一个合格的学生，能够很好地效仿我的老师——飞蝗泥蜂在生理学方面的本领。

我得承认，我为自己能做得几乎跟昆虫一样好而沾沾自喜。

一样好？我怎么能这么说？还是不要急着说这样的话吧，我要向飞蝗泥蜂学习的东西还有很多呢。事实上，我动过手术的那两只昆虫很快就死了，并且是真正意义上的死亡，四五天后，出现在我眼前的已经是两具臭气熏天的尸体了。飞蝗泥蜂的距螽呢？还用说吗？它甚至在手术10天后还保持着新鲜，并且完全符合幼虫对猎物的食用要求。不仅如此，飞蝗泥蜂在猎物头颅下的手术才结束几个小时，距螽的腿、触须、产卵管、大颚就能像什么事也没发生一样乱动起来。总而言之，昆虫又恢复了被飞蝗泥蜂咬住脑部之前的状态。这乱动一直保持着，只是会逐渐衰弱下去。飞蝗泥蜂只是暂时麻痹它的猎物而已，它完全可以在这段时间内，在猎物继续反抗前，把猎物拖到窝里；可我，自以为与它不相上下了，其实只不过是个既蠢笨又鲁莽的蹩脚的外科医生而已，因为我杀死了我的猎物。飞蝗泥蜂用它那无可复制的灵巧手法，熟练地压迫着猎物的头脑，使它麻醉了几个小时；而我，因为无知所以动作拙劣，也许我用镊子夹碎了这个纤细的器官，而这可是生命的原动力啊。如果我失败以后还有什么可以不必羞愧满面的理由，那就是，我坚信很少有人能像我这样与这些敏捷的生物比敏捷。

现在我来解释为什么飞蝗泥蜂不用蜇针伤害猎物的脑部神经节。因为在这个生命力的核心点上，只要注入一滴毒液就会使猎物全身静止，那么死亡也会随之而来。可是猎手并不是要猎物死亡；幼虫也根本不需要完全失去生命的猎物，不需要腐烂发臭的尸体；猎手需要的只是一种暂时麻木、暂时陷入昏沉状态的、在搬运途中不会反抗的猎物。它采取了一种只在实验生理学实验室里才惯用的脑部压迫法，并取得了成功。它像弗卢朗那样：把动物的脑袋剥露出来，然后施压，使动物立即丧失智力、意愿、敏感度和活动能力。当压迫停止后，一切又会恢复正常。距螽就是这样，当它从那巧妙压迫产生的麻醉感中慢慢苏醒过来时，又恢复了残余的生命活力。猎物的头部神经受到大颚的按压，但并不致命，然后慢慢恢复活力，从完完全全的昏沉状态中苏醒过来。我们必须承认，这真是一种可怕的科学！

在昆虫学研究中，真是世事难料：您拼命追求的，往往得不到；您已经忘记的，它却主动来敲门。为了探寻朗格多克飞蝗泥蜂猎杀距螽的过程，我多少次无功而返，又做了多少白费力气的事儿！20年过去了，这些写好的东西已经交到了出版社，这个月初的一天（1878年8月8日），我的儿子埃米尔突然冲进我的书房。"快！"他说，"快来！院子门前的梧桐树下有一只飞蝗泥蜂正拖着它的猎物！"埃米尔因为读过我写的东西所以知道这件事，他把我们夜间进行的准备工作当作一种趣事，尤其是他在我们的田野生活中曾看到过类似的事情。他看得一点没错。我迅速跑出去，看到了一只朗格多克飞蝗泥蜂正拖着一只麻木的距螽。它向附近的鸡窝走去，好像准备爬上鸡窝的墙壁，把窝筑在屋顶上的瓦片下；我曾在几年前看到过同一场景：同样是一只带着猎物的飞蝗泥蜂，爬上了同一个地方，把窝建在了一块接合得不是很好的瓦片弯曲处。眼前这只膜翅目昆虫，说不定就是我曾见过的那只艰难攀登的飞蝗泥蜂的后代呢。

\newpage

\chapter{本能的无知}

\begin{figure}[htbp]
    \centering
    
\includegraphics[width=0.8\textwidth]{images/00367.jpeg}
\end{figure}

飞蝗泥蜂已经给我们证明，它在某种无意识状态下，也就是在本能的支配下，行动是多么准确，技术是多么精湛；现在它将向我们证明，在遇到类似于稍微偏离正常道路的情况时，它又是多么的束手无策、才智枯竭，甚至与逻辑相悖。这就是本能的特征。这是一对有些怪异的矛盾：精妙绝伦的技能与深深的无知并存。因为本能，所以无论困难多大，它都可以做成一切事情。蜜蜂在建造那由三个菱形组成的六角形蜂房时，准确无误地解决了最大值与最小值等难题，如果人类遇到这样的难题，必须借助深奥的代数学才能解决。因为膜翅目昆虫的幼虫必须以猎物维生，所以即使是掌握了最精妙的解剖学和生理学的人，也不能与它那杰出的猎杀技术相提并论。如果行为在动物掌控的正常范围内，那么出于本能，它们可以解决一切难题；相反，如果行为超出了它们的正常实践范围，那么同样因为本能，即使是最简单的问题，对它们来说也是极其艰难的。昆虫那高度清醒的头脑，令我们惊奇、震撼，并拍案叫绝；可是不一会儿，面对最简单的却超出了它正常的实践范围的情况时，它又愚蠢得令我们难以理解。关于这一点，飞蝗泥蜂将为我们提供生动的事例。

我们仔细观察一下它是如何把距螽拖到窝里的吧！假如我们够幸运，也许会看到一个小场面，现在我就来描述一下这个场面。当这只膜翅目昆虫走进岩石下的隐蔽所时，看到了一只食肉类的修女螳螂正在一根草茎上歇息。这螳螂看似在虔诚地念经，其实却暗藏杀机。也许飞蝗泥蜂意识到了这埋伏在路上的强盗有多么危险，所以它放下猎物，勇敢地冲向螳螂，准备狠狠地揍它几下，把它赶走，或者吓吓它，让它不敢轻举妄动。可是那强盗一动不动，紧闭着前后臂那两把制造死亡的大锯。飞蝗泥蜂又从螳螂躺着的草茎旁走过。看看它头部的方向，就知道它十分警惕，它要用充满威力的目光震慑敌人，让它待在原地不动。这样的勇气是应该得到回报的：被放下的猎物没有遇到意外事故。

关于修女螳螂得多说几句。这种昆虫在普罗旺斯语中称为lou-prego Dieocu，意思是"祈祷上帝的昆虫"。它长着一对嫩绿色的长翅膀，好似大风帆一样，它总是高昂着头仰望天空，臂膀在胸前折叠交叉，这一切使它看上去就像一个正专心祈祷的修女。可事实上，它是一种嗜杀成性、凶残无比的昆虫。虽然它并不十分喜爱各种膜翅目掘地虫的工地，可还是会经常光顾。它在飞蝗泥蜂窝附近的荆棘丛上窥伺，等待着天赐良机，好半路劫杀某些到来者，这样它就可以得到双重收获，既得到猎手，又得到猎手的猎物。它的耐心是要经受时间的考验的。飞蝗泥蜂满腹狐疑，一直提防着，但是它还是逐渐放松了警惕，莫名其妙地犯起了糊涂。突然，螳螂像痉挛似的一抖，半张开翅膀，发出了一声巨响，把已经走近的飞蝗泥蜂吓了一跳，飞蝗泥蜂也因为害怕而迟疑起来。螳螂像弹簧似的，猛地把带着锯齿的前臂缩到同样长有锯齿的胳膊上，把飞蝗泥蜂夹在两个锯齿之间了，好像叼着饵物的狼被夹在了捕狼器的夹板中一样。这时，螳螂并没有松开它那凶猛的机器，而是一小口一小口地啃起了它的猎物。这种"祈上帝"昆虫就是这样凝神、祈祷、沉思的。

在我的印象中，我看到过修女螳螂屠杀的场景，我们不妨在此讲述一下。事情发生在泥蜂的工地上，这种泥蜂专门捕猎蜜蜂。因为这些膜翅目掘地虫的幼虫以蜜蜂为生，所以它们总趁蜜蜂忙着采集花粉和花蜜的时候，把它从花朵上抓来。如果泥蜂抓来的蜜蜂体内装满了蜜，那么它在储存起蜜蜂之前，会在路上或洞口先压迫蜜蜂的嗦囊，把甜美的糖浆挤出来，然后舔着这个不幸者的舌头饱吮一顿，糖浆源源不断地从垂死的蜜蜂嘴里流了出来。凶手挤压垂死者的肚子，把里面的蜂蜜挤空，当作自己的美味。如此折磨垂死者，真是太狠毒了。如果吃蜜蜂的泥蜂这么做有错，那我真该斥责它一番。血腥的美宴进行正酣的时候，我看到螳螂把膜翅目昆虫连同它的猎物一起抓住了，一个强盗遭到了另一个强盗的拦路抢劫。这个场面更恐怖，当螳螂抓住泥蜂，用它带着双重锯子的尖端戳穿泥蜂的肚子并开始啃噬时，泥蜂仍在吸吮着蜜蜂流出的蜜，尽管它正在遭受着死亡的痛苦，却还是不肯放弃到嘴的美味。我们快点把这个丑恶的场面遮盖上吧。

让我们继续看看飞蝗泥蜂的情况吧。在做更具体的描述前，我们应该先了解一下它的窝。它的窝建造在细沙里，或者更准确地说是在一个天然隐蔽所的尘土中。窝的过道很短，只有一两法寸长，中间没有拐弯，与唯一一间椭圆形的宽敞房间径直相通。总之，这是一个在匆忙之中挖成的洞穴，十分简陋，根本与精雕细刻的豪宅不沾边。

我曾经说过，飞蝗泥蜂的住所建得这么简陋，而且每个窝里只有一个房间、一个蜂房，是因为它是先抓到猎物才开始挖洞的，并且在挖洞的时候把猎物暂时丢在了狩猎场上。可是谁能预料到，猎手在第二次捕猎时会遭遇怎样的命运呢！所以洞穴必须建造在捕猎场所的附近。如果要储存第二只距螽，现在的住所就离得太远了，无法进行明天的工作。所以，它每抓到一只猎物，就要挖一个新洞，建造一个只有一间房间的新窝。这窝是随遇而安的，有时在这里，有时在那里。

介绍完这个情况以后，我们来做一些实验，看看当我们给飞蝗泥蜂创造一些新环境时，它会怎么做。

第一个实验。一只飞蝗泥蜂正拖着猎物，离窝还有几法寸远。我没有惊动它，而是用剪刀轻轻剪断距螽的触角，我们知道，飞蝗泥蜂是用触角当拖运的缰绳的。因为拖着的重担突然减轻，它奇怪地来到猎物身边，一把抓住触角的根部，也就是剪剩下的那一小节。太短了，几乎不到10毫米，但是没关系，对于飞蝗泥蜂来说，这已经够用了，于是，它咬住剩下的这一小段缰绳继续拖运。为了不伤害泥蜂，我又小心地剪起了那两段触角，这一次是贴着猎物的头顶盖剪的。昆虫在它熟悉的部位找不到可以抓的东西，就在旁边抓起猎物的一根长触须继续拖起来，并且没有对这种运输方式的改变感到奇怪。我允许它这么做。很快猎物被拖到了目的地，于是，膜翅目昆虫把猎物的头部放在洞口，独自走进窝里，在把食物储存起来之前，对蜂房的内部做了一番短暂的巡查。这不禁令人想起黄翅飞蝗泥蜂在同一种情况下采取的方法。我利用这个猎物被抛弃的短暂时间，迅速剪掉了它所有的触须，并把它放在了稍远一点，离窝大概有一步路的地方。飞蝗泥蜂上来了，它发现猎物在窝的门槛处，就径直跑过去。它上下左右地找遍了猎物的头部，却找不到任何可以抓住的东西。接下来，它在绝望之中做了一个无奈的尝试，把大颚张得大大的，试图咬住距螽的头；可是它的钳子张得不够大，根本夹不住这么大的东西，从又圆又滑的头颅上摔了下来。它又反复尝试了很多次，可终究没有任何结果。现在它承认是白费力气了。它走开了一点儿，似乎放弃努力了。它似乎已经垂头丧气了；它用后腿蹭蹭翅膀，把前跗节放到嘴里舔一舔，然后揉揉眼睛；在我看来这就表明膜翅目昆虫已经放弃了。

其实距螽的身上除了触角和触须外，还有别的可以抓住的部位。它有六条腿，还有产卵管，当然，这些器官都比较小，很难整个咬住当作拖运的支撑点。我相信，对于储存猎物这项工作来说，拉着触角先把头拖进去是最合适不过的；但是拉一条腿，尤其是前腿，同样也可以很容易地把猎物拖进去，因为洞口很宽，过道又短得几乎跟没有一样。可是，飞蝗泥蜂为什么不尝试一下抓住猎物六个跗节中的某一个，或者抓住它产卵管的端部，反而拼命做不切实际的、荒谬至极的事儿，用它那极短的大颚去咬猎物那巨大的脑袋呢？对这一点，它难道连想都没有想过吗？那么让我们想方设法提醒提醒它吧。

我把距螽的一条腿或者腹部的那把刀的末端放到飞蝗泥蜂的大颚下。飞蝗泥蜂却固执地拒绝；我一再引诱它，还是无济于事。它抓不住猎物的触角，又不知道去抓猎物的腿，飞蝗泥蜂进退两难了。这个猎手太奇怪了！也许是我的打扰，以及刚刚发生的这件不寻常的事影响了它的正常思维吧。那么，我们让飞蝗泥蜂单独和它的猎物在洞口待一会儿，让它在没人打扰的情况下安静思考解决的办法吧！于是，我丢下飞蝗泥蜂，继续走我的路。两小时后，我回到原处时，飞蝗泥蜂已经不在了，它的窝一直打开着，而距螽仍在我放的位置上躺着。由此，我们可以得出结论：膜翅目昆虫始终没有尝试过，它走了，把住所和猎物全都抛弃了，其实它只要抓住猎物的一条腿，这些就都还是它的东西。这种与弗卢朗不相上下的昆虫，刚刚还让我们对它的高超技能赞不绝口，因为它会压迫猎物的大脑使之昏迷，可当它面对习惯以外的情况时，尽管很简单，它也会愚蠢到令人无法想象的地步。它能精准地用蜇针刺中猎物前胸的神经节，用大颚压迫猎物的脑神经节；它会用带毒的蜇刺让神经的生命力永远消失，并通过压迫让猎物暂时昏迷，这一切它都掌控得恰到好处；可是它不知道，当抓不住一个部位时，可以抓住另一个部位。它根本不知道可以不抓触须而抓腿。它只知道抓触须或者抓猎物头部的别的丝状物，比如触角。如果没有这些绳子，它的种族就要遭殃了，因为这个小小的难题它根本无法解决。

第二个实验。膜翅目昆虫已经把食物储存在窝里，虫卵也产下了，此刻正忙着把窝封住。它倒退着用前跗节打扫起门前，把一堆堆尘土抛到住所门口。它的动作十分灵敏，从它肚子底下穿过的尘土，划出了一道抛物线，好像液体的网一样连续不断。飞蝗泥蜂不时用大颚挑选几粒沙子、几个小石子，然后插入土块中，用头顶住，用大颚挤压，把它们垒到一起。这道墙砌好以后，洞口的门很快就看不出来了。在它工作的过程中，我插了一手。我把飞蝗泥蜂拿开，用刀小心地清扫那条短过道，取走封门的材料，使蜂房与外部恢复畅通。但是我没有破坏它的建筑，而是用镊子把距螽从蜂房里取出来，当时距螽的头部位于窝的尽头处，产卵管位于门口。膜翅目昆虫像往常一样，把卵产在牺牲品的胸部，也就是一条后腿的根部位置，这表明膜翅目昆虫对它的窝做了最后的加工，以后再也不会回来了。

做完了这些事，并把取来的猎物小心地放到盒子里以后，我把位置让给了飞蝗泥蜂，而飞蝗泥蜂一直在一旁默默地看着它的住所被洗劫。它看见门打开了，就走了进去，在里面待了片刻后，又继续做起那项被我打断的工作，认真地封住蜂房门口，重新倒退着往门口扫地、运沙，始终认认真真地堆砌着，仿佛在做着一项重要的工作。洞口又一次封好了，昆虫抖抖身子，似乎对它的杰作十分满意，扫了一眼后就飞走了。

飞蝗泥蜂应该知道洞里的猎物已经被取走了，因为它刚才进去过一次，而且还足足待了一会儿。可是它在那所空无一物的住所里巡视了一周后，还是把蜂房重新封了起来，那股细致劲儿就像什么事也没发生过一样。它是准备把这个窝留到将来使用吗？还是想再带一只猎物回来，再在这里产卵吗？果真如此的话，它封窝就是为了防止某位不速之客趁它不在时闯入，那么它的这种做法就是一种十分谨慎的预防措施。防止别的掘地虫破坏自己的房子，这也许是保护房间的明智之举。对于某些掠夺成性的膜翅目昆虫来说，如果它们在工作中不得不暂停一段时间，那么它就确实会把门先封起来，防止别人进入。以蜜蜂为食的泥蜂的窝是一个竖井，我曾看到它在准备捕猎时，以及在太阳下山准备停工时，就是这样用一块平板小石头把蜂房的门封起来的。当然，那只是一个简单的工程，只用一块小石头盖在井口上。这样，昆虫回来时只要搬开那块小石头，就可以自由出入了，并且那是瞬间就可以办成的事。相反，我们刚才看到飞蝗泥蜂建造的是牢固的栅栏、坚实的砌体，整个过道里尘土和砾石一层夹着一层。这种建筑是会长久保留的，而不是暂时的防御工程，建筑者在建筑过程中的细致程度就是证明。况且从飞蝗泥蜂的行为方式来看，说它还会再回来利用这个住所，似乎很值得怀疑。我认为这已经是一个被证明了的不争的事实：飞蝗泥蜂将在别的地方捕捉猎物，也会在别的地方挖洞储存距螽。不过这毕竟只是推理，让我们看看真实的实验结果吧，实验比推理更有说服力。这件事发生后，我等了近一个星期，假如飞蝗泥蜂把窝顶细细地封起来是为了将来飞回来进行第二次产卵的话，那么我想看到这件事的发生。结果事实证明：它的窝一直完好地封存着，但是里面没有食物，没有虫卵，更没有幼虫。这个证明是绝对真实的，膜翅目昆虫再没有回来过。

被抢劫的飞蝗泥蜂进入窝里以后，看着这个空无一物的房间，十分镇定自若，就像根本没发现那个刚刚还占据着整个蜂房的庞大猎物已经被抢走了似的。它真没发现食物和卵已经被抢走了吗？它在杀猎物时，观察力那么敏锐，难道它的头脑会如此愚蠢，连蜂房里已经什么都没有了也没发现吗？我不敢说它愚蠢。因为它其实已经发现了。但是，这个窝已经被抢劫一空了，它也不打算再飞回来利用，那么它为什么还要继续封窝，并且那么一丝不苟呢？封门的工作是毫无意义并且荒谬至极的，但是，昆虫还是热情洋溢地完成了这项工作，就像这是决定幼虫未来的一件大事一样。昆虫的每一种行为都是受本能指引的、不可或缺的。因为它刚刚做过某件事，所以跟这件事相关的另一件事就必须做，从而补充前一件事或者为补充前一件事而铺路；同时这两个行为又是紧密相关的，所以做了第一件就一定要做第二件，即使在发生了某种特殊情况时，第二件事已经变得不合时宜了，甚至会损害自己的利益。现在，这个窝因为猎物和幼虫被抢走了而变得毫无用处。因为飞蝗泥蜂不会再飞回来了，所以它将永远也没有用处了，那么飞蝗泥蜂究竟为什么要把窝堵住呢？对于这种有悖逻辑的行为，只能看作对之前行为的补充，并且非做不可。在正常情况下，飞蝗泥蜂会先捕捉猎物，然后产卵，最后把窝封住。现在，捕猎这件事已经做过了，虽然猎物被我从蜂房里取走了，但是没关系，它已经做过了，卵也产完了，到了封窝的环节了。昆虫就是这样按部就班地做着它的事，内心没有一点儿判断力，并且一点儿也不认为它现在做的是无用功。

第三个实验。在正常条件下，能对一切都了然于胸；而一旦情况有变，又立刻变得一无所知，这就是昆虫向我们展示的既奇怪又矛盾的现象。下面我举的这个例子也是我从飞蝗泥蜂那儿得到的，可以证实这一说法。

白边飞蝗泥蜂捕猎的对象是中等大小的蝗虫，在它的住所附近，各个种类的蝗虫应有尽有，它可以不必挑选。因为蝗虫很多，所以捕猎不需要长途跋涉。当白边飞蝗泥蜂把那个像竖井一样的窝建好以后，只需在住所附近走动一小会儿，就能发现那些正在日光下觅食的蝗虫。它扑向蝗虫，制止蝗虫乱踢乱蹬的同时，把蜇针刺向蝗虫；这对于飞蝗泥蜂来说，只是转瞬之间的事。猎物扑腾几下它那或者呈胭脂红色或者呈天蓝色的翅膀，乱蹬了几下腿，就一动不动了。现在它要把猎物运到窝里去，并且要徒步运输。在进行这种艰苦作业时，白边飞蝗泥蜂采用了与它的两个同胞一样的方法，用大颚咬住猎物的一根触角，两腿抱着猎物，拖着往前走。在路上一旦遇到草丛，白边飞蝗泥蜂就从一根草茎跳到或飞到另一根草茎上去，并且始终牢牢咬住它的猎物。最后，当它来到离窝几步远的地方时，就像朗格多克飞蝗泥蜂那样干起来，只不过它总是不以为然的样子，所以不如后者那么庄重。白边飞蝗泥蜂把猎物扔在了路上，虽然住所并没有任何面临危险的迹象，它还是匆匆忙忙地奔向井口，并把头探进去好几次，有时甚至会往里走几步，然后返回来把蝗虫拖过来一点儿；接着又扔下猎物，第二次巡视起竖井来，就这样反反复复了很多次，而且每次都是那么匆忙。

\newpage

\chapter{登上万杜山}

\begin{figure}[htbp]
    \centering
    
\includegraphics[width=0.8\textwidth]{images/00375.jpeg}
\end{figure}

巍然耸立的万杜山是普罗旺斯的不毛之峰，四周受到各种大气因素的影响,它雄奇险峻，是阿尔卑斯和比利牛斯山边界处的最高峰。人们可以根据它来深入研究气候分布的植物种类。山底生长着郁郁葱葱的橄榄树和各种半木本植物，它们十分怕冷，比如百里香，它那芳香的气味需要南方地区的太阳的照耀；山顶上会有至少半年的时间都覆盖着皑皑白雪，个别地方长着那种来自于极地海滩的北方的花。沿着山坡往上走大半天光景，您就会接连看到生长在这里的主要植物种类，这些植物都是您往常要在同一经度上从南到北长途跋涉才能遇到的。刚动身时，您脚下踩着的是一簇簇百里香，它们散发着香脂般的气味，连绵不绝，像地毯一样铺满了山脉低处的圆山丘；过几个小时，您脚下就是生长着对生叶的虎耳草暗色的小垫了，这是植物学家在7月的斯匹茨卑尔根海边登岸后看到的第一种植物。在低处，您将在篱笆下采摘一朵像石榴树那样鲜红的花朵，石榴树是非洲气候的忠实拥护者；在高处，您将采撷到一种小小的毛茸茸的虞美人，它的花茎长在碎石渣下面，开着黄色的叶瓣较宽的花。这种虞美人除了长在万杜山顶的斜坡上以外，在格陵兰和北海海峡广阔的冰雪天地里也能见到。

这种对比鲜明的景象，令人每次到那儿都有一种新鲜感。虽然迄今为止我已经登过25次山，却远远没有满足。1865年8月我第23次登山。我们一行8人，3个人是为了植物学考察，5个人是要到山上走走，看看高处的风景。我们那5个对植物研究毫无兴致的同伴，后来没有一个人愿意陪我再去一趟，因为这场远征十分艰难，是看日出的乐趣不能弥补的。

可以打个形象的比喻，把万杜山比作一堆修路的碎石。如果把这堆石头垒到2 000米高，给它一个与高度成正比的底座，再在它那白色的石灰岩体表上植入一片漆黑的森林，就形成了整座山的印象。这个碎石堆，有零碎的碎石，有大块的石块，矗立在一马平川的平原上，没有斜坡也没有台阶。这种大平台把平原一层层地摞起，使它变得易于攀登。一走上碎石嶙峋的山路后，就开始登山了，这段路上连最好的路面也不如新铺的碎石路好走，再往后越来越难走，一直到海拔1 912米的山顶。鲜嫩的草地，奔腾的小溪，苔衣青石和百年老树的巨大树荫，这些可以为山岭增添无穷魅力的景色，在万杜峰上全都看不到。这里有的只是无穷无尽的剥落的岩层，脚踩上去就会将它踩碎，发出近乎金属碰撞的干硬的咔嚓声，石片一片接一片地滑下。在万杜山上，碎石流代替了瀑布，岩石窸窸窣窣的坍塌声代替了潺潺的流水声。

我们现在到达了万杜山山脚下的一个地方，名叫贝都安。我们已经与向导交涉好，并约定了出发时间，随身带哪些食物也商量好了。我们要想方设法睡上一觉，因为明天到了山上就会度过一个不眠之夜。可是睡觉对我来说真难，我总是睡不着，反而更加疲惫不堪。读者中如果有人想去登万杜山做植物学考察，请千万听我一句劝告，不要在星期天傍晚到达贝都安，因为这样你们才能躲开客栈里人来人往的嘈杂声，没完没了的高谈阔论声，弹子球房里连续不断的弹子球撞击声，觥筹交错的叮叮当当声，酒后小调的哼唱声，过往行人的夜曲声，隔壁酒吧里铜管乐器的喧闹声，以及其他的在这不需劳动的快乐日子里避免不了的各种折磨。否则，在接下来的一个星期中，你们能得到更好的休息吗？但愿能，可我不敢保证。反正我是一夜没合眼。那个专门为我们准备食品的生锈的烤肉架，在我的房外整整转动了一夜，吱吱嘎嘎地响个不停，而我跟那可恶的机器仅有一板之隔。天已经微微亮了，驴子在窗下大声嚎叫起来，看来是起床的时候了，跟没睡没什么区别。把粮食和行李装好以后，我们的向导吆喝了一声："驾！吁！"我们就出发了，此时是凌晨4点钟。万杜山的向导总管特里布勒牵着骡子和驴子走在人群的最前面。我的那些热衷植物学的伙伴们，在黎明的微光下，仔细探寻着路边的植物；其他人闲聊着；我跟在队尾，肩上挂着晴雨计，手上拿着笔记本和铅笔。

我的晴雨计本来是用于记录各种植物的生长纬度的，竟然成了大家跟朗姆酒葫芦接吻的借口。一有人发现某种特别的植物就会大喊："快，用一下晴雨计！"于是我们全都一股脑围住了酒葫芦，然后才把我肩上的物理仪器拿来。清晨的凉爽和山高路远，大家很喜欢这样来看晴雨计，以至于酒葫芦里烈酒的下降速度比晴雨计里水银柱的下降速度还快。考虑到不能只顾眼前，我不得不减少查看这托里切利管的次数。

气温越来越低，刚才的树种逐渐消失了，最先消失的是绿橄榄树和橡树，然后是葡萄树和扁桃树，接着是桑树、核桃树、白栎树。黄杨开始变得多起来，俯拾皆是。我们进入了一个植物种类十分单调的地区，这里是山毛榉生长地带的下限，作物停止生长，主要的植物是山地风轮菜。风轮菜纤细的叶子里充满了香精油的味道，又苦又涩，当地人们俗称它为"驴梨"。我们带的食物中有一些小乳酪上就洒着这种味道刺鼻的佐料。不止一个人想吃这些乳酪了，不止一个人因为饥饿而把目光投向了骡子背上存放粮食的鞍囊上。清晨的艰苦操练激发了我们的食欲，哪里是食欲，简直就是饥肠辘辘的饥饿感，贺拉斯将这种情况称为"胃的焦躁不安"。我教我的同事们怎样在下一个休息地到达之前应付胃的焦躁不安。我指给他们看生长在乱石中的一种叶子呈箭矢状的小酸味食物，并且以身试法吃了一大口。大家先是对我的建议嗤之以鼻，我置之不理，随便他们取笑。但是紧接着，他们争先恐后地采摘起这珍贵的酸模叶了。

我们嘴里咀嚼着酸酸的叶子，到达了山毛榉的生长地带，最先见到的是一些枝条藤蔓拖延在地上的灌木，稀稀疏疏地披散在山坡上；紧接着又出现了一些小矮树，一棵挨着一棵地挤在一起；最后看到的是一片枝干又粗又壮的山毛榉，结成了一片茂密幽深的树林，那儿的土壤是钙质的土块。这些山毛榉树冬天受积雪重压，一整年都饱受密斯托拉风的强吹猛刮，很多树枝都折断了，只剩下扭扭曲曲的奇怪形状，甚至有的已躺倒在地。穿过这片森林用去了一个多小时的时间，远远望去，这森林就像一条围在万杜山山腰上的黑色腰带一样。再往前走，山毛榉又呈现出稀稀疏疏分布的情况，原来这里是山毛榉生长地带的上限。只咀嚼酸模叶毕竟不是长久之计，当我们到达事先约定的休息地，大家都长舒了一口气，开始准备吃午饭。

我们的休息地是拉格拉斯泉，一股从地下冒出来的涓涓泉水注入到了这些首尾相连的山毛榉树搭成的长凹槽里；山里的牧人常常把羊群赶到这里饮水。泉水的温度是7℃，对于我们这些从平原如火炉般的酷暑中走出来的人来说，真是清凉得难以想象。地面好像被阿尔卑斯山的植物铺成了一条地毯，一泓泉水淌在地毯中央，长着百里香叶子的指甲草点缀着美丽的地毯，它那宽大而纤薄的苞就像银色的鳞片，散发着白色的光芒。从鞍囊里拿出食物，从稻草层中取出酒瓶。这边摆放着不易破碎的器皿以及涂着蒜汁的羊后腿和面包；淡味的小鸡放在另一侧，当空空如也的肚子填饱以后，可以吃这小鸡打打牙祭。在不远处可以算得上是上座的地方，摆着加了山地风轮菜为调料的万杜乳酪，也就是驴梨小乳酪；乳酪旁放着阿尔红香肠，肥肉条和整块的梨像大理石似的镶嵌在它那粉红色的肉泥中；摆在这边角落里的有绿橄榄，还在流着卤水、用油做佐料的橄榄；在那边的空地上，放着卡瓦翁的西瓜，有白瓤的，有橘黄色瓤的，每个人可以依据自己的喜好挑选食用；在另一处，摆着鲲鱼罐头，就着美味的小牛腿肉，可以猛喝鱼汤；最后，酒瓶子都放到了冰凉的水槽里镇上。我们忘记带什么了吗？是的，我们忘记带最重要的餐后点心玉葱了，拿东西直接蘸着盐就可以吃。刚开始，那两个巴黎人看到这么丰盛的一大桌菜肴，惊得目瞪口呆；不一会儿，他们就要开始赞不绝口了。好了，一切都准备就绪，开始吃吧！

于是，大家发疯似的大吃起来，这可能是我们一生中最难忘的一次进餐。刚开始吃的时候，大家都饥不择食，抓起一块块羊后腿、一片片面包就往嘴里塞，速度快得惊人。每个人都焦虑地看着眼前的食物，心想："照这样的吃法，今晚够吃了，可明天吃什么？"当然，大家都在心里嘀咕，没人说出来。起先，我们都一声不地狼吞虎咽着；这会儿已经不那么饿了，开始一边吃一边闲聊起来，为明天担忧的心情也渐渐缓解下来了。大家开始称赞起膳食总管来，认为他安排得十分合理，事先就预料到了我们会如此饥饿，所以把一切都准备得相当充足，完全有备无患。现在，大家开始像美食家那样评价食物了。一个伙伴用刀尖戳起一个个橄榄，赞不绝口；另一个则一边对鲲鱼罐头大加赞扬，一边在面包片上切开赭石色的小鱼；第三个伙伴正兴致勃勃地谈论起红香肠；最后大家异口同声称赞起了那还没有一个手掌大的驴梨乳酪。吃完后，大家点燃了烟斗和雪茄，躺在草地上，暖洋洋的太阳正好晒在肚皮上。

已经休息一个小时了，起来吧！时间紧迫，必须继续赶路了。向导独自带着行李沿着树林边往西边走去，那儿有一条可容牲畜通过的山路。海拔大约1 550米的地方是山毛榉生长的上限处，那里有一个羊棚或者称为大房子的地方，向导就到那里等我们。羊棚是用山石垒成的大房子，房顶盖着草，可供牲畜和人在这里过夜。我们继续爬山，来到山脊处，这样顺着山脊攀上山顶就能少费些力气。太阳落山后，我们将从山顶下到羊棚，向导肯定早已等候在那里了。这就是经向导提议并征得了我们同意的计划。

我们登上山脊了，这个斜坡比较舒缓，一直向南延伸下去，一眼望不到边，我们刚刚就是顺着这条斜坡爬过来的。向北坡望去，一片苍茫，气势雄伟，山坡时而是直上直下的陡壁，时而是令人毛骨悚然的陡峭阶梯。梯级虽然仅高1.5米，却十分陡峭；如果抛出一块石头，它会一蹦一蹦地一直滚到谷底，在半路上都不会停下来。山下的土鲁朗克河河床像一条明晃晃的带子，清晰可见。当我的同伴们撼动起一块岩石，把它推入深谷，听着它发出的一阵轰响时，我见到了我的老相识毛刺砂泥蜂。它躲藏在一块扁平的大石头下面。以前我在平原的道路路坡上看到这种砂泥蜂时，它总是孤零零的，可是在靠近万杜山山峰的地方，它们有上百只密密麻麻地挤在一起。

我正要研究这么多砂泥蜂聚集在一起的原因时，突然刮起了南部地区特有的风，今天清晨，这风已经让我们胆战心惊了，现在又忽然乌云翻滚，要下起雨来。我们还没来得及躲开，瓢泼大雨就倾盆而下了，雨水迷蒙了我们的眼睛，只能看到两步路远。不巧的是，我们当中的一个人，我最好的朋友德拉库尔独自去寻找山里的一种稀有植物岩生大戟，到现在还没回来。我们用双手围成话筒放在嘴边，一起放开嗓子大声叫喊着他，可是没有回声。我们的喊声被呼啸的风声和哗哗的倾盆暴雨声淹没了。既然迷路者听不到我们的喊声，我们就开始寻找他。在漆黑的云层遮蔽下，我们只要走出两步路，就谁也看不见谁了，而我们7个人中，只有我对这里最熟悉。为了不再有人丢失，我们大家手挽着手前进，我走在这串人群的最前面，就这样我们转了几分钟，真正地玩起了捉迷藏的游戏，可还是没找到他。德拉库尔了解万杜山的气候，可能在乌云压过来的时候，他借着最后一线亮光，迅速跑回羊棚去了。我们赶快也到羊棚去看看吧，我们浑身被雨淋透了；斜纹布裤已经贴在了腿上，就像人体的第二张皮似的。

这时，一个更大的难题出现了：我们来来回回、往往返返地寻找德拉库尔时，我已经像被蒙住双眼一样在原地打转，根本辨不清方向了，我甚至一点也不知道，哪边是右山脊。我问问这个，问问那个，大家的意见都不同，都不确定。我们中没人能断定哪边是北方，哪边是南方。我真的没有，从来没有像此时此刻这样，意识到辨别东西南北方位的重要性。我们周围是茫茫的灰云，我们的脚下，只能辨认出哪儿有斜坡，往哪儿延伸。但是应该走哪个斜坡呢？必须先选对路，然后才能往下走。否则，我们一旦失算往北坡走去，就会跌入万丈深渊，摔个粉身碎骨。那深渊，我们刚才看了一眼都胆战心惊，如果跌下去，一个都不会活命。那时候，在那几分钟里，我茫然无措地站在那里，陷入了深深的痛苦之中。

大部分人说："我们就待在原地，等雨停了再说。""这个主意实在糟糕透了！"另一部分人反驳道，我也是其中的一个反对者。雨会持续很久，像我们现在淋成这样，到了夜里一冷，都会被冻僵的。我敬爱的朋友贝尔纳·威尔罗，是为了跟我一起攀登万杜山而专程从巴黎植物园赶过来的，他十分镇定沉稳，坚信我能带领大家脱离险境。为了不引起别人的恐慌，我把他拉到一边，向他坦白了我内心的忧虑。我们进行了一场十分秘密的谈话，试图用思考的罗盘代替缺失的磁针。他问道："刚才乌云是从南面来的吗？""确实是从南面来的。""虽然风向难辨，但是雨是稍微由南向北偏斜的是吗？""是的，当时我还能辨别方向，我注意了这一点。我们可以由此辨别方向吗？我们从雨点落下的方向下山。""我想过这样做，可是我不确信。风力太弱了，根本无法判断风向。而且也许这是旋风，当乌云笼罩山顶时，通常会刮旋风。我们根本无法断定，风向是否从开始到现在一直没变过，现在的风不是从北边刮来的。""您的怀疑很有道理。那该怎么办呢？""是啊，怎么办，怎么办，问题就在这儿。啊，我有一个想法了：如果风向一直没变，那么我们应当主要是身体的左侧淋湿，而且只要我们没有迷失方向的话，雨是从这边打来的。如果风向有变，那我们浑身上下应该差不多一样湿。我们想想再决定好吗？""好！""可是如果我弄错了呢？""您不会弄错的。"

只消三言两语，大家都知道是怎么回事了。大家开始摸摸自己的衣服，不是摸外面的衣服，因为那不准确，而是摸最贴身的衣服。当我听到大家纷纷说左边湿得比右边严重时，我长舒了一口气。我们又手拉手连成一队，我走在最前面，威尔罗断后，力求不使一个人掉队。在出发之前，我再次对我的朋友说："哎，我们会不会有危险？""危险就危险，我们跟着您走。"于是我们不管三七二十一，一头扎进了前途未卜的探索中。

\newpage

\chapter{迁徙者}

\begin{figure}[htbp]
    \centering
    
\includegraphics[width=0.8\textwidth]{images/00382.jpeg}
\end{figure}

我曾说过，在大约1 800米海拔的万杜山顶峰，我有过一次昆虫学考察的好机会，如果这样的机会能经常出现，让人进行系统研究，一定会取得伟大的成就。但不幸的是，这样的考察我只进行过一次，此后就再无法做深入研究了，所以我对这次的考察结果还存在很多疑惑，这就交给未来的考察者们用事实来证实我的猜测吧。

毛刺砂泥蜂在平原上也是比较常见的，但一般都是孤零零地出现在山间的小路边或沙坡上，或者忙着挖掘竖井，或者忙着搬运沉重的毛虫。它总是喜欢独来独往，就像朗格多克飞蝗泥蜂一样，所以当我在即将到达万杜山顶的地方，发现同一块石板下聚集着这么多膜翅目昆虫时，感到十分惊讶。现在我看到的，不是我一直以来见到的形单影只的单独个体，而是一个数目庞大的群落。现在，我们来探讨一下它们这么聚居的原因吧。

对于膜翅目掘地虫来说，这种例外情况是十分罕见的。当春天来了，毛刺砂泥蜂就开始筑窝了。接近3月底时，如果天气开始变暖，最迟4月上旬，当蟋蟀已初具成虫形态、正在家门口痛苦地蜕掉幼虫皮时，当备受诗人喜爱的水仙第一次盛开、鹀鸟在高高的柳树枝头卖弄轻柔绵远的歌声时，毛刺砂泥蜂正忙着为它的幼虫挖掘住所，储备粮食。对于别的种类的砂泥蜂和绝大多数掠夺成性的膜翅目昆虫来说，只在秋天，也就是9、10月间才开始做这些工作。这种昆虫比一般的膜翅目昆虫提前6个月筑窝的情况，引起了我的思考。

我在想，这些在4月初就开始筑窝的毛刺砂泥蜂是不是当年我见到的昆虫；也就是说，这些春天的劳动者们是不是就是前3个月才完成蜕变并从茧中钻出来的那些昆虫呢？通常情况下，所有掘地虫变成成虫后，离开地下住宅以及为它们的幼虫准备粮食，都是在同一季节进行的。大多数擅长狩猎的膜翅目昆虫，都是在6、7月间从它们幼时居住的地下拱廊中出来，8、9、10月才开始进行矿工和猎手的工作。

这种法则对于毛刺砂泥蜂来说是不是也同样适用呢？它是否也是在同一个季节里完成变态并开始从事为幼虫储存食物的工作呢？这很值得怀疑，因为如果膜翅目昆虫在3月底就开始忙着筑窝，那它就应该在冬天，最迟在2月底完成蜕变并从茧中钻出来。可是气候的寒冷令我们无法接受这样的结论。当刺骨的密斯托拉风大约刮了半个月，大地冻得硬邦邦的时候，当伴随着凛冽的寒风开始飘扬起鹅毛大雪时，蛹期的艰难变态是不可能完成的，而成虫也不可能在这时候离开茧的隐蔽所。只有在夏天的明媚阳光下，土地既温暖又润湿，成虫才会舍得离开它的地下住宅。

假如我知道毛刺砂泥蜂第一次从地洞里出来的时间，那对我将极有帮助,可是十分遗憾，我无从得知。我长年累月记下的笔记中，没有这方面的说明（因为这类研究总是取决于那些难以预料的机遇，所以难免有模糊混乱的地方），而我今天意识到了这个问题的重要性，所以我要把我的材料综合起来，写下这些情况。我在笔记中找到，沙地砂泥蜂在6月5日孵化，银色砂泥蜂在6月20日孵化，可是没有任何关于毛刺砂泥蜂孵化期的介绍。这是被我疏忽的细节。这两种昆虫的孵化期符合一般规律，成虫是在天气炎热的时候出现的。通过类推的方法，我认为毛刺砂泥蜂也是同一时期破茧而出的。

可是，为什么我们看到的这种砂泥蜂在3月底或4月初的时候就开始建造住所呢？答案很简单，这些膜翅目昆虫不是当年破茧而出的，而是上一年。它们在一般情况下，也就是6、7月从茧里出来，经过了一个冬季，到了春天就开始立刻筑窝。总而言之，这些是已经经历了一个冬天的昆虫，实验证实了这一结论。

一年又一年，各种采蜜的膜翅目昆虫在阳光普照的竖直的土块或沙块上繁衍着后代，它们在壁上挖出一个又一个洞，由一条条走廊连接起来，好似迷宫一样，又好像一个巨大的海绵。在寒冷的冬季，人们只要耐心寻找，一定会看到毛刺砂泥蜂蜷缩着身子在住所的凹陷处惬意地晒着太阳，也许是孤孤单单的一只，也许是三三两两的一群，它们悠闲地等待着晴朗天气的到来。在寒冷的冬季，当鹀鸟和蟋蟀开始鸣叫时，这些优美的膜翅目昆虫已经为山间小路旁的草地增添了一股生机。只要我愿意，我是可以尽情享受这些小乐趣的。如果风平浪静，阳光再稍稍强烈些，这些怕冷的昆虫就会走出隐蔽所，欢快地来到洞口沐浴温暖的阳光；也有的会冒险来到外面，边擦拭着翅膀，边缓缓地走过海绵状的沙层表面。灰色的小蜥蜴，在住所的旧墙被暖阳晒得热乎乎的时候，也是如此。

但是在冬天，即使是在十分暖和的隐蔽所内，也根本找不到节腹泥蜂、飞蝗泥蜂、大头泥蜂以及一切需要为自己的幼虫捕猎肉食的膜翅目昆虫。因为它们经过秋天的一番劳作之后，全都死了，而在隆冬季节，它们的种族只剩下在地洞深处冬眠的幼虫了。因此，毛刺砂泥蜂是一种罕见的昆虫，只有它在炎热的季节孵化，然后在某个温暖的隐蔽所度过寒冬，这也是它们出现在初春时节的原因。

让我们根据这些资料，试着解释一下万杜山顶上的毛刺砂泥蜂为什么会群居吧。这么庞大的一群砂泥蜂躲藏在石块下究竟干什么呢？它们是想把这儿当作过冬的隐藏所，就这样蜷缩在石板下，等待适宜它们的季节到来吗？所有迹象都表明这是不可能的。因为动物并不是在炎热的8月冬眠的。缺乏从花朵里吮吸的蜜汁也不可能是理由。9月的阵雨很快就要来临，在伏天暂时停止生长的植物就要开始茁壮成长，装点着一个跟春天一样五彩缤纷的美丽田野。大多数膜翅目昆虫在这个季节都会过得十分欢快，毛刺砂泥蜂也不例外，不可能在这个时候睡眠。

另外，万杜山高大陡峭、崎岖不平，密斯托拉风呼啸着横扫而过，有时会把山毛榉和冷杉连根拔起；在山顶上，刺骨的北风会持续刮上6个月之久，把雪花吹得漫天飞扬；山峰上基本被寒云冷雾笼罩。怎么能想象这种喜爱温暖的昆虫会把这里当作过冬的隐藏所呢？这跟让它去北海海角的冰上过冬没什么区别。不，毛刺砂泥蜂不可能在这里过冬。我们见到的蜂群也许只是路过这里。每当有下雨的迹象（这些迹象我们看不到，可这些对大气变化十分敏感的毛刺砂泥蜂能提前感觉出来）时，蜂群就躲到石头下面等待雨停后再飞走。可是，它们来自哪里？要到哪儿去呢？

就是在8月，主要是9月这段时间里，迁徙的小鸟们从它们喜爱的地方，那是比我们这里更凉爽、树木更多、更安静的地方，也是它们产卵的地方，一个地方一个地方地往南飞，到达我们这个生长着橄榄树的酷热的地方。它们抵达的日子几乎是固定的，先后次序一点儿不变，就好像它们根据只有自己明白的日历制定了这个黄道吉日似的。它们在我们的平原上停留的那几天，是最富足的日子，因为这里有许多它们专门捕食的昆虫。这些鸟来到了我们的田地里，一片接一片地寻找那些被犁耙耕出来的暴露在地上的小虫，这是它们的美味；这样吃下去，没几天它们的臀部就长得肥肥的了，并且就像一个丰富的粮仓一样，里面装满了各种极富营养的储藏品，以供将来不时之需；最后，在储备下足够的旅途食物后，它们继续往南飞，前往那个没有冬季的、一整年都有昆虫的地方：西班牙和意大利南部、地中海上的岛屿、非洲。这是阿尔吉尔人进行狩猎、品尝野味的快乐时期。

长翅百灵最先抵达，我们这里称它们为"克雷乌"。8月初的时候，就会看到这种鸟在田地里寻觅狗尾草的种子，那是一种对农作物有害的禾本科植物。稍有惊动，它就立即飞走，发出难听的咕噜声，它的普罗旺斯名字就形象地模拟出了这种声音。在它之后不久来到的是石鹏，它在原先的苜蓿地里从容地捕猎着象虫、蝗虫、蚂蚁。紧接着，枝头的贵宾——阿尔吉尔的名鸟飞来了；到了9月，鹃鸟——也叫白尾雀,到来了，它是一种难得的美味，所有品尝过它的人都称赞不已。就连罗马著名诗人马提亚尔在他的诗中赞美过的那种颇受罗马饕餮之徒喜爱的燕雀，也比不上白尾雀鲜美的脂肪球，不过这些白尾雀吃得太多了，已经过于肥美了。这种鸟捕猎各种昆虫，我收集到的博物学资料中记录了它肚子里的东西，其中可以找到各种幼虫以及象虫、蝗虫、砂潜、龟甲、叶甲、蟋蟀、球螋、蚂蚁、蜘蛛、鼠妇、蜗牛、赤马陆和其他种类的各种各样的昆虫。为了帮助消化这些野味，它还要吃些葡萄、树莓以及猩红色的欧亚山茱萸的浆果。白尾雀飞起来时张开了尾巴，尾巴上洁白的羽毛，使它看上去就像只飘逸的蝴蝶一样美丽。它从一块土地飞到另一块土地，不停地吃着各种美味的食物，所以它才长得这么肥胖。

在增肥术方面，有一种唯一可以与它相媲美的鸟类，这就是与它同时迁徙的、同样爱吃昆虫的灌木丛中的雀形目小鸟——鹨鸟，书本上都用这种并不恰当的名字，而牧人则几乎都用"肥腿"来称呼它，意思是特别肥的鸟。仅从名称上看，就能知道它的基本特征了。没有任何一种鸟能长得这么肥。当到了一定的时候，这鸟浑身上下直至翅膀、脖子、脑袋下全都会长满板油，就像一小块一小块牛油似的。它太不幸了，专门爱吃象虫，所以浑身长满了脂肪。它费了很大的力气才从一株桑树飞到另一株桑树上，因为太胖了，它几乎要窒息了，停在浓密的树叶丛中喘着粗气。

10月，半灰半白的灰鹡鸰飞来了，它胸前长着黑色绒毛，脖子长长的，身体也又细又长。这种鸟，优雅地摇摆着尾巴，迈着小步子蹦蹦跳跳地跟在农夫套车的马脚下，在新开垦的田野里啄食害虫。云雀几乎在同一时期来到，起初是一小批先遣部队先到，好像侦察兵一样，然后大部队才会到来，占据了麦田和新开垦的土地。那里有许多狗尾草种子，这是它们最喜欢的食物。此时，在平原上，一棵棵草茎上挂满了洁白的白霜和露珠，在初升的阳光照射下，像镜子似地反射出耀眼的光芒。此时，猎人放飞的枭鸟，飞了一小段距离就俯冲下来，紧接着它转动着那双惊恐的眼睛，又猛地直冲云霄。云雀猛扑下来时，看到那闪闪发光的犁耙或者那巨大的飞鸟，十分惊异。云雀就在您面前十几步远的地方，垂着两爪，张开翅膀，就像圣灵的图像那样。是时候了，瞄准开枪吧！我祝愿读者们在这场狩猎中玩得畅快淋漓。

草地鹨是跟云雀一起到来的，俗名叫西西，这也是一个拟声词，形象地模拟出了它低沉的鸣叫声，它总是混在云雀群中与云雀一道飞来。这种鸟是最让枭鸟痴迷的了，它不断拍打着翅膀，在枭鸟周围飞翔。这些迁徙者们，我们是很难见到第二次的。因为它们大部分只是在我们这儿歇一歇脚：这里食物丰富，尤其是有许多昆虫，所以它们就在这里待上几个星期，吃得又强又壮，肥肥胖胖的，然后再继续往南飞。另外还有一些鸟会在我们的平原上选择一个地点过冬，因为这里很少下雪，甚至在三九天，地上也能找到许许多多的小种子。云雀就是其中的一种，它在麦地和新开垦的田地里觅食；草地鹨也是如此，它更愿意到苜蓿田和草地去觅食。

云雀几乎在整个法国都很常见，可它的住所不在沃克吕兹平原。生活在沃克吕兹平原的是凤头百灵，也叫"黏着的云雀"，筑路工人和养路工人对它最为熟悉。它寻找合适的产卵地点，根本不必到遥远的北方，在沃克吕兹平原的邻省德龙省就经常能见到这种鸟的窝。所以，在秋冬季节占领我们平原的云雀中，很可能有很多是从德龙省飞来的，而不是来自更遥远的地方。它们只要飞到邻省，就能找到没有积雪的平原和许许多多的小种子吃。

在我看来，万杜山顶附近的砂泥蜂群就是这种短距离的迁徙者。我已经证实了，这种膜翅目昆虫是蜕变为成虫以后躲在某个隐蔽所中度过严冬的，并且在4月初就开始建造住所了。它跟云雀一样，也要避寒。它不怕没有食物，因为它就算不吃东西也可以坚持到春暖花开的时节，但它十分怕冷，以至于必须采取措施防御严寒。所以它从冰天雪地的地方飞走；它像鸟类一样成群结队地迁徙，不远万里地来到古老的城墙上，或者有阳光照耀的南方沙滩上寻找住所。寒冬过后，这群昆虫又都飞回了原处，或者只有一部分飞回去。这就是万杜山上会有砂泥蜂群的原因。它们是一群迁徙的部落，来自于寒冷的德龙省，为了飞到生长着橄榄树的炎热的南方平原上，它们飞过了土鲁朗克深深的大峡谷，可是突然要下雨，就暂时在山顶处歇息一下。由此看来，毛刺砂泥蜂是为了躲过严冬才不得不进行迁徙的。当小鸟开始迁徙时，毛刺砂泥蜂也成群结队地开始搬迁了，它们也会离开寒冷之地，来到炎热之地。翻过几座山，越过几道岭，它们就到达了目的地。

关于昆虫在高地上奇怪地聚居的情况，我还曾收集到另外两个例子。10月的时候，我发现万杜山顶的小教堂居住着许许多多七星瓢虫，这种昆虫在民间俗称为慈悲虫。小教堂屋顶有多少块石板，这些昆虫就在石头上造了多少面墙壁，它们彼此紧密地挨在一起，从几步路以外看起来，那粗糙丑陋的建筑物就像个珊瑚球似的。聚集在这里的七星瓢虫，就像恒河里的沙粒一样多，我估计不出到底有多少只。这些吃蚜虫的昆虫，绝对不是被食物吸引到这海拔近2 000米的万杜山顶上来的。因为这里的植物十分贫乏，蚜虫是决不会冒险到这么高的地方来的。

还有一次是在6月，我在万杜山附近的海拔734米的圣阿曼高原上，又看到了类似这样的奇怪聚集，不过这次的数目要少得多。在高原的最高处，在悬崖的陡壁边上，竖立着一个以砌石为底座的十字架。正是在这底座的各面和底座的基石上，与在万杜山上相同的膜翅目昆虫——七星瓢虫成群地聚集在一起。这些虫子大多数一动不动；但是在阳光强烈的地方，新来的占位者和原先占位者之间在那临时的圣坛上，总是在不断地交换着位置，原先的占位者飞走，过一会儿再回来。

跟万杜山一样，圣阿曼高原上也没有任何线索能够告诉我，在这干涸的土地上既没有蚜虫，也没有别的吸引七星瓢虫的东西，它们为什么会如此超乎寻常地聚集在一起；也没有任何迹象能够告诉我，在高山的砌石工程上，众多昆虫的这种聚会究竟隐藏着什么秘密。这里还有别的昆虫迁徙的例子吗？有没有像燕子一样的，在出发之前大家先聚集在一起呢？这儿是不是聚集点，成群结队的七星瓢虫要从这儿出发到食物更丰富的地方呢？这种可能性非常大，但也十分奇怪，因为众所周知，七星瓢虫向来是不喜欢旅行的。当我们看到它猎杀蔷薇花上的绿色小虫和蚕豆上的黑色小虫时，我们肯定它是喜欢家居而不爱外出的，可是它们却依靠那短短的翅膀，成群结队地来到万杜山顶聚会，而且即使是雨燕，也只不过是在极有热情的时候才会飞到那儿去。它们为什么在这么高的地方聚会呢？为什么喜欢在这些堆砌起来的石头上居住呢？

\newpage

\chapter{砂泥蜂}

\begin{figure}[htbp]
    \centering
    \includegraphics[width=0.8\textwidth]{images/00389.jpeg}
\end{figure}

身体纤长，体态轻巧，腹部末端又细又窄，好像一根细线系在身上似的，身披一身黑衣，肚皮上有一块装饰性的红披巾，这就是砂泥蜂的大致样子。它们的形状和颜色与黄翅飞蝗泥蜂十分相近，可是习性有着天壤之别。飞蝗泥蜂捕捉直翅目昆虫，比如蝗虫、距螽、蟋蟀，而砂泥蜂喜欢毛虫。因为猎物的不同，它们在本能的捕杀手段上也十分不同。

砂泥蜂这个词听上去有些刺耳，所以我想挑挑刺儿。这个名字的意思是"沙之友"，不过这个术语太绝对化了，而且可能并不准确。沙子，这里所说的是干燥的、呈颗粒状的并且缓缓流动的纯沙，它真正的朋友是捕猎苍蝇的泥蜂，我在这里要介绍的毛虫捕捉者却根本不喜欢流动的纯沙；相反，它们见到这种流沙就会远远地躲开，因为这种沙只要稍一碰触，就会倒塌。它们在把食物和虫卵封在蜂房里之前，挖掘的竖井必须是畅通无阻的，所以挖掘竖井的地方应该比较坚硬、结实，以免刚挖完不久就被堵住了。它们需要的是一块容易挖掘的土质疏松的地方，那里的沙只需用黏土和石灰就可以牢牢粘住。山路边那长满了青草的斜坡，每天都有阳光的普照，这才是它们最喜欢的地方。在这些地方，一到春天，尤其是4月初的时候就有毛刺砂泥蜂了；而当9、10月来临时，沙地砂泥蜂、银色砂泥蜂和柔丝砂泥蜂也纷纷出现了。我在这里综合叙述一下这四种砂泥蜂的情况。

这四种砂泥蜂的地洞都是像井一样垂直的洞，里面的直径最多有一根粗粗的鹅毛管那么粗，深约0.5分米。底部就是蜂房，向来只有一间，比连通蜂房的竖井稍微大一点。总之，这个住宅真是再简单不过了，一次就可以轻而易举地挖成；幼虫就像黄翅飞蝗泥蜂的一样，有一个四层壳的茧，它在洞里就是靠这个茧壳抵御严冬的。砂泥蜂的挖掘工程是独自完成的，它安安稳稳、不急不忙地挖着，从来没有特别热情高涨的时候。它以前跗节为耙，以大颚充当挖掘工具。假如遇到难以拔出来的沙粒，我们就会听到从井底传来它的翅膀和身体振动发出的沙沙声，十分尖利，好像在大声喊叫似的。它每隔一小段时间就会来到地面，牙齿叼着一颗挖出来的沙砾，飞起来扔到稍远的地方去，以防影响工程。在它挖出来的沙砾中，有的沙子的形状和体积十分引人注意，至少砂泥蜂没有把它们扔到稍远的地方去，而是用脚搬运，放在井口的一边。原来这些是比较上乘的材料，是现成的砾石，将来可以用于封闭住所。

它在进行外部施工时，谨慎而细致，高高地翘着身子，腹部挂在一条长长的肉茎末端，每当要转身时，它都要整个掉个头，就像一条线一样，一端固定住了，另一端则十分准确地转动起来。如果它必须把它认为没必要的碎屑扔到远处，它就默默地一小块一小块地倒退着扔。砂泥蜂的头部总是最后从竖井里出来的，这可能是为了避免调转身体，以节省时间吧。这种动作十分呆板，就像自动装置似的，可腹部有长肉茎的沙地砂泥蜂和柔丝砂泥蜂却干得非常认真。如果腹部鼓胀得像梨子一样大，拴在一根带子的末端，那么调转起来就很艰难，一旦动作过于迅猛，纤细的肉茎就会折断，所以这些砂泥蜂的动作准确无误。如果它们需要飞，就倒退着出来，以免再进行一次调头转身。相反，毛刺砂泥蜂腹部的肉茎比较短，在挖地穴时像大部分掘地虫一样，动作洒脱而灵敏。它可以自由行动，因为没有肚子的障碍。

住宅挖好后，一到晚上，甚至一到太阳照不到的时候，砂泥蜂就要来到它在挖掘时堆积的小砾石堆上寻找一番，选一块满意的石子；如果在那儿找不到，就到附近去找，反正总能很快找到一块。这是一块平平的小石板，直径比井口略大一圈。它用大颚把石板搬过来，暂时盖在洞口上。这扇门十分坚实，能保护它的住所不受侵扰。第二天，当阳光普照着四周的斜坡，天气暖和，便于捕猎时，砂泥蜂是一定能找到它的地洞的。它咬着一条已被麻醉的毛虫的脖子，用腿把毛虫拖了回来；它掀开石板（这石板看上去跟周围的小石子别无二样，可只有它清楚区别在哪儿），把猎物放到井底，然后产下卵，接着把井口附近的残渣碎屑扫进竖井里，最后把住所永远地封了起来。

我好几次都看到，当太阳已经落山或者更晚些的时候，只能等到第二天再把食物储备到井底的情况下，沙地砂泥蜂和银色砂泥蜂就先把地洞封闭起来。膜翅目昆虫封住了洞口，我也只好第二天再来观察了。不过我事先在这里画了个记号，选好标线和基准点，插上几根树枝当标杆，以便在竖井填满后也能找到。因为如果我第二天来晚了，或者我让膜翅目昆虫有空利用白天时间，那么地穴总是已经储备好了食物，并且已经封闭起来了。

砂泥蜂的记忆力实在令人震惊。每当它工作到了很晚，不得不把余下的工作推到明天做时，它不在已经挖好的竖井里过夜，而是暂时抛弃这个新家，用一块小石板将井口盖住，然后才走开。它对这地方并不熟悉，也并不比别的地方更了解，因为砂泥蜂跟朗格多克飞蝗泥蜂一样，它随遇而安，飞到哪儿就把卵产在哪儿。它无意中飞到一个地方，如果它喜欢那儿的土壤，就会把洞挖在那里。现在砂泥蜂走开了，它去哪儿了？谁知道呢，也许在附近的花朵上休憩吧。在花朵上，在沉沉的暮色中，它会从花冠里吮吸一滴蜜汁；就像我们的矿工一样，在漆黑的巷道里累得疲惫不堪，晚上要喝一瓶酒来补充体力。砂泥蜂走开了，可能远一些也可能近一些，一站站地走到花窖里去。它度过了傍晚、夜间和清晨，可是它必须返回它的地洞里去完成未完的工程；头天夜里，它一朵花一朵花地尽情畅饮，清晨又来回走动开始捕猎，现在它必须回到地洞里。胡蜂能回到它的窝里，蜜蜂能回到它的蜂箱里，这一点不令我感到奇怪。它的窝和它的蜂箱是固定住宅，因为它们飞来飞去无数次了，对这条路已经熟悉了；可是砂泥蜂对这地方是很陌生的，却要在离开那么久以后再回到地洞里。它的竖井所在的地方，也许是它头一天才第一次到过的，可它今天必须再回到那里，它现在已经分辨不了方向了，何况还背负着沉重的猎物。可它对地形记得一清二楚，有时甚至精准得令我肃然起敬。砂泥蜂径直朝它的地洞奔去，就像这周围的小路它已经走了成千上万次似的。

有时，它会迟疑很久，寻找很多次。如果情况严重时，沉重的猎物会妨碍它快速寻找，它就把猎物放在高处的百里香上或一束草上，以便于过一会儿再来找时能够容易找到。放下猎物以后，砂泥蜂卸下了负担，继续积极地寻找。随着砂泥蜂走过的地方，我用铅笔画出它走的路线图。从所画的路线可以看到，那是一条毫无规则的线路，有弧线形，有锐角形，有凹曲的支线，也有辐射的支线，一环套一环地打着结、画着圈，彼此交叉，总之是个名副其实的迷宫。这路线如此复杂，说明迷途的砂泥蜂惶恐不安，不知所措。

终于，竖井找到了，石板也掀起来了，现在它必须回去寻找毛虫了；可是如果砂泥蜂在找窝的时候，来来往往走出了很远，再回到毛虫那儿就得费一番周折。虽然砂泥蜂把猎物放在了容易发现的地方，可它似乎预知到当它再回来寻找猎物时，会遇到很大的麻烦。所以一旦寻找住所的时间太久，它就会果断地停止寻找，义无反顾地回到毛虫那儿去，摸一摸，咬一咬，似乎证实一下这是不是它原来的猎物，是不是它当初放在这里的财产；然后又匆匆忙忙地奔到搜索地点，搜寻一会儿，紧接着它又一次放弃搜索，返回来查看查看猎物；然后再第三次去搜索。我认为它这样一再返回来，是为了让自己记住存放猎物的地点。

当它寻找地洞的路线十分复杂时，它就用这个办法；但是通常情况下，昆虫可以顺利地回到它前一天挖的竖井里。然而，竖井的地点并不是它之前就熟悉的，而是它随便确定的某个地方，是它随便走到哪里就挖在哪里的。它牢牢地记住了这个地点，并依靠记忆来指引行动。我下面就要介绍一下它的记忆力发挥了多么神奇的作用。对我而言，要想第二天再找到用小石片遮盖的竖井，我可不敢只靠记忆。我必须用笔记下来，画个草图，标明路线，竖立标杆，总之，我必须借助一整套详细的几何学知识。

用石板把地穴暂时封起来的办法，似乎只有沙地砂泥蜂和银色砂泥蜂会用，其他两种砂泥蜂则不会。至少我从没见到后两种砂泥蜂的住所盖过盖子。毛刺砂泥蜂则根本就不要这种暂时的封闭物，根据我的观察，毛刺砂泥蜂总是先捕捉猎物，然后在捕猎地点附近挖洞，以便及时把食物储存起来，完全没必要浪费力气盖个盖子。至于柔丝砂泥蜂不盖盖子的行为，我想可能另有原因。其他三种砂泥蜂在每个地穴里只放一只毛虫，而它则放五只，当然，体积都很小。就像我们忘掉关门一样，柔丝砂泥蜂或许也忘记把石板盖在竖井上了，因为它在很短的时间内就要下到井里至少五次。

其他三种砂泥蜂只给每只幼虫一只毛虫。不过，这些毛虫以庞大的体积弥补了数量上的匮乏。因为猎物肥肥胖胖，完全能够满足幼虫的食欲。我就曾经从沙地砂泥蜂的嘴里夺来过一只是猎手体重15倍的毛虫。15倍，只要我们想想猎手咬着这么沉重的猎物脖颈上的皮，克服地上的千难万苦，拖回去要花费的巨大力气，就知道这是个多么了不起的数字了。把别的膜翅目昆虫和它的猎物同时放到秤上称量会发现，掠夺者和它的猎物之间没有差距这么悬殊的。从地穴里挖出来的，以及从砂泥蜂爪下看到的食物颜色是千变万化的，这说明这三种掠夺者并没有特别偏爱的猎物，见到毛虫就捕捉，只要这毛虫的身材适中，不太大也不太小，而且属于夜蛾一类就行了。一身灰色、在浅地表啃咬植物根茎的毛虫是最常见的猎物。

在关于砂泥蜂的所有故事中，最主要的、最令我感兴趣的就是它捕捉猎物的方式，以及它是如何为了幼虫的安全而将猎物弄到无法伤人的状态的。它捕捉的猎物，确实与我们至今见到的吉丁、象虫、蝗虫、距螽等牺牲品在机体上有很大的差异。毛虫是由一系列相似的环或节段组成的。其中头三个环上有真正的脚，这些脚就是未来夜蛾的脚；其他的环上有膜状的脚，或者说假脚，这些脚只有毛虫才有，夜蛾则没有；另外一些环则没有脚。每个环都有神经核，也称神经节，是产生感觉和支配动作的中枢；这样，除去头颅里类似大脑的神经圈以外，神经分布系统有十二个彼此分隔开的不同的神经节。

这与象虫和吉丁的情况不同，象虫和吉丁的神经节是聚集在一处的，只要刺一下就可以造成全身的麻醉；与飞蝗泥蜂刺蟋蟀的情况也不同，飞蝗泥蜂是在蟋蟀的胸部神经节上刺了一下又一下从而使它一动不能动的。夜蛾的毛虫不是只有一个神经集中点，也不是只有三个神经中枢，而是有十二个由于环节相隔而彼此分离的神经节，这些神经节位于腹部，像一个一个念珠似的排列在身体的中轴线上。另外，在低等生物中，同一器官大量重复，因散乱而失去力量的情况已是普遍规律；这些神经核彼此之间是完全独立的，每个神经核支配着它所属的那个节段的功能，其相邻节段的功能紊乱对这一节段的影响是微乎其微的。所以，即便毛虫的一个环节失去了活动能力和感觉，其他的环节也仍然是完好无损的，仍可以自由活动，并保持灵敏的感觉。这些情况足以说明，这种膜翅目昆虫捕杀猎物的方法具有重要的研究价值。

这种方式虽然很有研究价值，观察起来却困难重重。砂泥蜂从来是独来独往的，所以它们零零散散地分布在广袤之地上，而且总是在十分偶然的情况下才能遇到一只。这一切使得我们不可能像研究朗格多克飞蝗泥蜂那样，对它也提前准备好实验。观察砂泥蜂必须要长久地等待时机，要有百折不挠的毅力，当您已经不抱有任何希望而希望突然来到时，要善于抓住良机。我一年又一年地等待着这样的机会；有一天机会突然出现了，我可以很容易地观察到它的一举一动，就连细节也看得清清楚楚的，这大概是对我长久以来苦苦等待的一个补偿吧。

在开始研究的时候，我曾经亲眼看到过两次砂泥蜂猎杀毛虫的过程，虽然它的动作十分迅捷，我却看到它把蜇针刺入了猎物的第五或第六节段上，而且只刺了一下就大获成功了。为了证实这一观察结果，我产生了一个想法：掠夺者正忙着把毛虫拖到地洞里去，既然我没有亲眼看到它蜇刺的整个过程，就应该把猎手的毛虫偷来，看看它到底被刺在了哪个节段上。可是我根本没有用放大镜观察，因为任何放大镜都发现不了牺牲品上的一点伤痕。我采取的办法是：我拿来一根细针，用针尖挨个节段地在这只一动不动的毛虫身上进行探测，这样，我可以根据昆虫对疼痛的反应程度来测试它的敏感度。结果当针尖刺进它的第五或第六节段，甚至整个戳穿时，毛虫还是毫无反应；而如果刺在这个节段的前面或后面，哪怕轻轻刺一秒钟，毛虫就开始乱扭身子，挣扎起来；刺入的节段离第五、六节越远，它挣扎得越厉害。特别是靠近后部末端的地方，只要稍一碰触，毛虫就开始扭动。由此可见，砂泥蜂蜇针只刺了一下，并且刺在了猎物的第五或第六节段上。

现在，为了防止毛虫逃走，让它一动不动，砂泥蜂是否会把蜇针刺到八个具有运动器官的每一个节段上呢？对于非常弱小的猎物，它会特别采取这种过分的预防措施吗？肯定不会，因为蜇针只刺一下就够了，但是蜇针必须刺到中心点上，毒汁产生的麻木感会在最短的时间内从那里逐渐扩散到有脚的节段内。所以为了进行这唯一一次注射，毫无疑问，应该选择在把两组运动环节分隔开的第五或第六节段。经过逻辑推理得出的蜇刺点，就是昆虫实际上采用的蜇刺点。

最后，我们要指出，砂泥蜂的卵仍是毫不例外地产在了猎物被完全麻醉的那个节段上。在这个部位，而且只有在这个部位上，小幼虫才可以任意啃咬猎物而不会引起猎物的扭动，从而避免了猎物对它造成的伤害。在这个部位，毒针的蜇刺不会使猎物产生任何反应，幼虫的啮咬也同样不会使它产生任何反应。猎物就这样一直安安静静地待在那里，最后，当幼虫长得身强力壮时，可以往猎物的其他部位进攻了，而这时它也有了自卫能力，同样不会被猎物伤害。

我在后来的多次观察研究中产生了怀疑，不是怀疑我的观察结果，而是怀疑这种观察结果普遍性的延伸。一些纤弱的量地虫和身材矮小的毛虫，只要被刺一下就丧失了伤害幼虫的能力，尤其是当掠夺者的蜇针刺在上面提到的第五或第六节段上时，这当然是极有可能的，而且不论是我直观的观察，还是用针尖来测试它的敏感性，这一点都已经得到了证实。可是沙地砂泥蜂，特别是毛刺砂泥蜂捕捉的猎物，十分庞大，我说过，这猎物的重量是掠夺者的15倍。蜇刺这种身材巨大的猎物的方法与蜇刺虚弱的量地虫的方法一样吗？为了制服这庞然大物，使它无法伤害幼虫，只刺一针可以吗？如果这可怕的灰色毛虫用它那强劲的臀部撞击蜂房的墙壁，会不会伤害虫卵或者小幼虫呢？我们无法想象，当这个刚刚孵化出来的小生命与那只巨龙面对面地待在狭小的地洞里时到底会发生什么事儿，不要忘了，这巨龙还可以自由地时而卷起或伸直它那弯弯曲曲的身子。

对毛虫进行了敏感性的探测实验以后，我对此更加怀疑了。当膜翅目昆虫的蜇针刺在猎物身上第五、六节段以外的其他环节时，柔丝砂泥蜂和银色砂泥蜂的小猎物就会拼命地挣扎个不停；而沙地砂泥蜂，尤其是毛刺砂泥蜂捕捉到的肥胖的毛虫，不论被刺到中间、前面还是后面，随便哪个地方，它都一动不动。它们的身体不会扭动，尾部不会突然上卷，蜇针只不过让它的皮肤轻微地颤动了一下而已，这说明毛虫还存留着一丝感觉。当用这种巨大的猎物喂养幼虫时，为了幼虫的安全考虑，必须使猎物的运动能力和感觉能力几乎全部丧失。在把猎物运到地洞里之前，砂泥蜂已使这团东西丧失活力了，但并没有使它真正死去。

我曾有幸亲眼看到了砂泥蜂给粗胖的毛虫动手术的过程。我跟我的一个朋友（唉，遗憾的是，他之后不久就去世了）一道从安格尔高原上下来，为圣甲虫设置陷阱，想考察考察它的智慧。这时，我们突然发现了一只毛刺砂泥蜂，它在一丛百里香下忙碌着。我们立即在这个劳作着的膜翅目昆虫附近趴下。昆虫并没有因为我们在场而惊恐不已，它来到我的袖子上歇息了一会儿，见它的两个邻居纹丝不动，认定不会伤害到自己，就又回到百里香下面去了。我们算是老朋友了，我知道这种大胆的亲密关系意味着什么，膜翅目昆虫正忙着做某件重要的事。让我们等着看吧，马上就会知道真相了。

砂泥蜂扒开百里香根茎处的土壤，拔出植物那纤细的侧茎，把头钻进了掀起来的小土块下面。它在百里香周围不停地忙碌着，一会儿跑到这儿，一会儿跑到那儿，探测着一切能使它进入灌木下面的缝隙。它不是在挖掘某个住宅，而是正在捕猎某种住在地下的猎物；看到它的动作，我们不禁想起了一条赶兔子出窝的狗。果然，一条肥肥胖胖的灰色毛虫不知道地面上发生了什么事，它心里忐忑不安，被砂泥蜂追得步步紧逼，决心离开地下来到地上。这下它惨了，猎手猛地扑了上来，抓住它后颈的皮，无论它怎么挣扎都无济于事。砂泥蜂跨在这庞然大物的背上，翘起腹部，就像一个对患者的解剖学结构了然于胸的外科大夫那样，有条不紊、不慌不忙地用手术刀在受害者的腹部，从第一节段到最后一个节段挨个刺了一遍。所有环节都不放过，无论这环节上有没有脚，它都要刺一下，而且按照从前到后的顺序，顺次刺下来。

这就是我看到的情景，而且是在我悠然自得、十分便利的情况下看到的，要想得到一个无可厚非的观察结果必须有这样的条件。膜翅目昆虫的动作十分精准，连科学家也会赞叹不绝；它知道人类几乎永远不可能知道的事；它对猎物的神经器官了如指掌，毛虫有多少个神经节，它就蜇刺多少下。我刚才说它知道和了解，我应该说，它的行为看上去就好像它知道并了解一样。它的行为是完全受到上天的启发的。昆虫根本不知道它在做什么，不过是顺从它本能的指引而已。但是这种至高无上的天启从何而来呢？这一切能用返祖性、行为选择、生存竞争这些理论来解释吗？我和我的朋友认为，这个问题自始至终都有那种难以言传的逻辑的最有说服力的启示，这逻辑以神启的法则操纵着这个世界，并指引着无意识的行为。我们被这真理之光深深地打动了，一种难以名状的泪水夺眶而出。

\begin{figure}[htbp]
    \centering
    
\includegraphics[width=0.8\textwidth]{images/00390.jpeg}
    \caption{柔丝砂泥蜂：体长18～24毫米，捕猎对象是量地虫，每年9、10月才能活跃。}
\end{figure}

\begin{figure}[htbp]
    \centering
    
\includegraphics[width=0.8\textwidth]{images/00391.jpeg}
    \caption{采蜜的蜂：蜂在采集花蜜时，通常选择盛开的花朵，因为这时的花朵中花蜜或分泌物含量较高。}
\end{figure}

\begin{figure}[htbp]
    \centering
    
\includegraphics[width=0.8\textwidth]{images/00392.jpeg}
    \caption{斜纹夜蛾：鳞翅目夜蛾科。成虫体长16～21毫米，翅展37～42毫米。}
\end{figure}

\begin{figure}[htbp]
    \centering
    
\includegraphics[width=0.8\textwidth]{images/00393.jpeg}
    \caption{夜蛾毛虫：夜蛾毛虫有十二个彼此分离的神经节，它们排列在身体的中轴线上，而砂泥蜂依然可以对其一击致命。}
\end{figure}

\begin{figure}[htbp]
    \centering
    
\includegraphics[width=0.8\textwidth]{images/00394.jpeg}
    \caption{桑枝尺蠖幼虫：桑枝尺蠖属尺蠖蛾科，其幼虫和夜蛾幼虫都称为量地虫。幼虫初孵化时为淡绿色，后变成灰褐色，静止时依枝斜立，好似一桑枝。}
\end{figure}

\begin{figure}[htbp]
    \centering
    \includegraphics[width=0.8\textwidth]{images/00395.jpeg}
    \caption{毛虫的神经系统：毛虫共有两根神经索，每段体节上都布满瘤状的神经节，神经节上布满纤细的神经末梢。}
\end{figure}

\newpage

\chapter{泥蜂}

\begin{figure}[htbp]
    \centering
    \includegraphics[width=0.8\textwidth]{images/00396.jpeg}
\end{figure}

伊萨尔森林\footnote{伊萨尔森林：距离阿维尼翁市约6千米。——译者注}是我最喜欢的观察地点，离阿维尼翁不远，位于罗讷河的右岸，与杜朗斯河的河口相对。千万不要曲解这个词的意思，很多人听到森林这个词会想到那铺着一层青苔地毯的土壤，或者想到长满密密麻麻的百年树龄的乔木林，阳光透过树叶的缝隙投下斑驳的阴影。然而，在炎热的夏季平原上，油橄榄树稀稀落落的，上面传来了蝉的聒噪声，那树影斑驳、凉气宜人的隐蔽所是找不到的。生长在伊萨尔森林里的，只有那一人高的绿色矮橡树，树木零零落落地分布着，树荫根本不足以抵挡酷暑的热气。在七八月份的三伏天里，我连续几个下午坐在矮林中的最佳观察点上，只能撑起一把大伞遮阳。然而就是这把大伞，竟然在我后来的观察中帮了我大忙，我会在合适的时候讲起这件事的。假如我在长途跋涉中忘记带上这把累赘，那么我能抵御太阳的唯一方法就是躺在某个沙丘的后面，而当太阳穴被太阳晒得青筋爆出时，我只得把头插进兔子洞里。在伊萨尔森林，我就是用这些办法乘凉的。

地面几乎是不毛之地，连木质的植物丛也没有。覆盖在地表的是流动性很大的干燥的细沙，在长有绿色橡树树根或树桩的地方，在风力的作用下，细沙堆成了一个个小沙丘。因为细沙十分易于流动，稍微有个下陷就会坍塌下来，所以沙的表面总是能自动恢复得匀称完整，沙丘的斜面也总是很平整。如果把手指插进沙里再拔出来，会立刻有细沙塌落下来，填平凹处，而沙面又恢复到了原来的样子，一点也没有凌乱的痕迹。但是在一定的深度上，如果最近一次的雨水刚下过不久，细沙还比较潮湿，就会粘连在一起，具有一定的黏度，可以在这个时候挖个小洞，洞壁和洞顶也不会塌陷。这里艳阳高照，阳光炽热，膜翅目昆虫的耙子一耙，沙坡就自然而然地塌下来，丰富的野味可当作幼虫的食物。这里人迹罕至，没人能破坏这里的宁静。在这块专属泥蜂的乐园里，真是万事俱备。让我们看看这种动作敏捷的昆虫的杰作吧。

如果读者也愿意陪我一起坐在伞下，或者像我一样钻进兔子洞里，那么他就会看到这样一幅7月末的场面：一只泥蜂不知从哪里突然冒了出来，它不进行任何巡查就随随便便地停在某个地方。在我看来，它停的位置跟沙地的其他地方没有任何不同。它的前腿长着一排排十分有力的纤毛，像扫把，像毛刷，又像钉耙似的，挖着它的地下住所。它用它靠后的四只脚支撑身体，最后两只脚稍稍叉开，前脚交替耙着或清扫着流动的沙土。它跗节的绞盘上就算长着弹簧，动作大概也不会比现在更精准而迅捷了。沙土从肚子下部抛到了身后，穿过后腿的拱孔，源源不断地喷射出来，像涓涓的溪流一样，在空中划出一道抛物线，落到了2米开外的地方。连续5～10分钟的时间里，它一刻不停地抛射着沙土，落下的沙土十分密集，这充分表明它的动作十分迅速敏捷。我不知道还有哪类昆虫的动作能像它一样灵活，而这种灵活一点也不妨碍它将动作做得轻盈优雅又收放自如。它一会儿在这边进进退退，一会儿又到那边进进退退，抛出的沙柱从未中断过。

它挖掘的地方土质比较松软。随着膜翅目昆虫的挖掘，周围的沙土不断地塌落下来把洞填满。塌落的沙中夹杂着碎木屑、烂叶根以及稍微大一些的石粒。泥蜂用大颚咬着这些大石粒，倒退着把它们扔到稍远的地方去；然后再回来把洞清扫干净，但是它每次都是轻轻一扫而过，并不会扫得很深，可能它根本不打算深入到地下去吧。它为什么只是单个在地面上劳作呢？如果只粗略地看上一眼就想知道答案，那是不可能的，但是我与这些可爱的膜翅目昆虫共度了几天的光景，在我把观察到的零散的材料综合起来之后，我觉得我仿佛知道了它这样劳作的原因。

这种膜翅目昆虫的窝肯定在地下几法寸深的地方；那建造在既新鲜又牢固的沙土中的小窝里，有一个卵，可能还有一只幼虫，泥蜂母亲每天用蝇喂养它。蝇是泥蜂幼虫的唯一食物。所以泥蜂母亲应该可以随时用脚抱着幼儿的粮食飞进窝里，就像猛禽抓着麋鹿、野猪这些美味后，随时钻进巢穴里喂养它的幼禽一样。但是如果猛禽的洞穴建造在某个难以进入的凸出的岩石上，那么它要返回时的唯一困难就是捕获的猎物太重或者不利于搬运；可是，泥蜂每次要进窝里都必须像矿工那样进行艰苦的挖掘，重新开辟巷道。因为随着它的进入，沙土会坍塌下来，巷道无疑就自动堵塞了。在这个地下住所里，幼虫居住的那个阔绰的蜂房的洞壁是唯一不会塌陷的，这蜂房建造的地点，就在它进食了半个月的美味后排泄出的废物中。泥蜂母亲要想到巷道尽头处的房间里，或者外出去打猎，必须经过狭窄的前厅，所以，前厅每次都会塌陷下来；至少得挖开干沙里的巷道，那部分已经因为一次次的进进出出而变得更加松软了。正是因为这些原因，膜翅目昆虫每次进洞或出洞时，都要在堆落的沙土中为自己挖一条通道。

相对来说，从洞里出来是比较容易的，哪怕是第一次破沙土而出也不难。刚开始的时候，沙子虽然有点儿坚硬，但是昆虫也能活动自如，它在将它掩埋的遮蔽物下是安全的，它可以从容不迫地运用它的跗节和大颚。可是，从外面进到洞里就完全不同了。泥蜂用前脚抱着猎物，猎物紧贴着它的肚皮，这使得它活动起来十分不便，也无法自如地使用它的工具； 更糟糕的是，那些卑鄙无耻的、堪称真正的盗贼的寄生虫，就潜藏在洞的四周，它们窥伺着泥蜂母亲，等它辛辛苦苦地带着猎物返回，并且眼看就要消失在巷道里的时候，突然飞过来把卵产到猎物身上。一旦强盗的阴谋得逞了，泥蜂的幼虫就会因为这些贪婪的共栖者而饿死。泥蜂似乎意识到了这个危险，所以它为了顺利地迅速返回窝里而采取了一些防御措施。堵塞在洞口的沙土，只要用头一顶，再用前脚快速一扫，就可以扒开。所以，泥蜂对留在住所四周的材料做了细致的筛选。当泥蜂母亲有闲暇，而且温度适宜，幼虫也有足够的食物无须它照料时，它就耙耙门前的土；因为它进入家门时危险重重，所以必须提前把一切可能阻挡进洞通道的碎木屑、大沙粒、烂树叶都统统剔除出去。我们刚才看到它勤勤恳恳、忙碌地工作，就是在进行这样的筛选。为了便于它进入地洞，它把前厅的材料都彻底检查了一遍，把一切妨碍走路的东西都剔出去，清除掉。昆虫的动作如此灵敏，工作起来如此快乐，它不正是用自己这种独特的方式表达着做母亲的幸福感、表达着它能照顾好这个存放着虫卵的家的自豪感吗？

因为这种膜翅目昆虫只注重保护家的外部，而不准备进到大厅里去，加上洞里的一切都井井有条，所以它没什么可着急的了。我们苦苦地等待着，却毫无收获，这时候的昆虫大概没什么可告诉我们的了。那么，我们来参观一下它的地洞吧。在泥蜂经常停留的地方，我们用刀慢慢地刮着沙丘，很快就发现洞口处的前厅被堵塞了，不过从被搅乱的沙土的特殊形状上，还是能够辨认出这里有一段通道的。这个过道有手指那么粗，顺应着沙土的性质和地势的起伏情况，呈现出或直或弯，或长或短的形状，总共有2～3分米长，通向这个凉爽的沙中的唯一一个房间。房间四周的墙壁上，没有涂抹砂浆来防止坍塌，也没有用泥浆来使粗糙的表面变得光滑。只要保证屋顶在幼虫的整个成长期内不会坍塌就可以了；当幼虫进入那个类似保险的坚硬的茧里以后，墙壁是否坍塌就无关紧要了。我们一会儿会看到它造茧的过程。蜂房的工程十分粗糙简陋，整个工程只不过是粗略地挖一挖而已，没有固定的形状，天花板也压得很低，整个空间才勉强能容下两三个核桃。

蜂房里只存放着一只猎物，个头儿很小，根本无法满足贪吃的幼虫的胃口。这个猎物就是叉叶绿蝇，它全身金绿，专门吃腐烂的肉。这种属于双翅目昆虫的食物，纹丝不动地躺在那儿，它是真正死亡了吗？还是仅仅被麻醉了？这个问题在后面就会知道，我们眼前的这个猎物，在它的肋部有一个稍稍有些弯曲的圆柱形白色虫卵，长约2毫米。这就是泥蜂的卵。通过泥蜂母亲的行为方式，我们可以预料它的住所里没有什么要紧事，事实也确实如此。它产下虫卵，并已经按照一定的比例，为将在24小时后诞生的虚弱的幼虫配备好了食物。在很长一段时间内，泥蜂母亲都不会再回到窝里，它只是守护在周围，另外也有可能会另挖一个洞，继续产卵。它产下一个又一个卵，每个卵都单独占据了一个蜂房。

带喙的泥蜂最初的食物只是一只小小的猎物，然而并不是只有它具备这样的特点，其他的各类泥蜂也都如此。我们随意打开一个刚刚产卵不久的泥蜂的窝，都会看到卵紧紧地粘在一只双翅目昆虫的肋部，一只就足够了。不仅如此，这些初生婴儿的猎物个个都十分矮小，可能这是泥蜂母亲专门为它的幼虫挑选的，因为这些小猎物通常比较鲜嫩，更适合虚弱的幼虫食用吧。但是，泥蜂母亲这样选择也可能是出于另一个目的，那就是要新鲜的食物。关于这个问题，我们将在后面进行考察。上桌的第一道菜大多都不丰盛，而且菜的种类还会因地洞附近是否有其他种类的猎物而发生变化。有时是叉叶绿绳，有时是厩蜇蝇或者某种尾蛆蝇，有时是披着黑毛绒衣服的柔弱的蜂虻；其中尤以肚子细长的斐洛福尼蝇最常见。

对于贪吃的幼虫来说，只有一只双翅目昆虫为食，确实太少了。不过，我们从这个具有普遍性的事实中，已经能够看出泥蜂最独特的特性了。在每一个蜂房里，膜翅目昆虫都会为幼虫储存维持生计所需的食物，其数目完全可以满足幼虫在整个孵化期内的需要。所以，它把卵产在一只猎物身上以后，就把住所封闭起来，从此以后再也不回来了。慢慢地，幼虫孵化出来了，默默地长大，它一眼就看到面前那份专门为它准备的食物。可是这个规律不适于泥蜂。蜂房在一开始就只有一只野味，自始至终都只有一只，体积很小，卵就产在猎物身上。产完卵以后，母亲就离开地洞了，而地洞会自动封闭，不过在泥蜂母亲离开之前，它会仔仔细细地把洞外的地面耙平整，以使任何人都看不出地下有个洞穴。

两三天过去了，新生的小幼虫吃着母亲事先为它准备的优质食物。在此期间，泥蜂母亲就守候在洞穴附近，它一会儿在瘦姬蜂流淌出汁液的头上舔一舔，一会儿在炎热的沙土上躺一会儿。毫无疑问，它保护着住所，担负起了住所警卫员的工作。有时，它会挑起洞口的沙子，然后飞得杳无踪影。或许它是忙着赶去另外一个地方挖掘蜂房，并以同样的方式为幼虫储备粮食吧。但是从它为幼虫准备的食物上可以看出，它无论飞走多久，都不曾忘记它的小幼虫，因为它捕获的猎物都是精心筛选的。母亲的本能让它清楚地知道婴儿的口粮会在什么时候吃完，什么时候需要添置新食物，于是它就会在那个时候准时地回到窝里。它能一下子就找到没有任何痕迹的洞口。这真是太令人钦佩了！这一次，它抱着一个稍微大一些的猎物来到地下了。可是刚把猎物放下，它就又离开了，到洞外等候着为幼虫供应第三次食物。幼虫狼吞虎咽地吃着，以至于第三次供应食物的时间很快就到来了。于是，泥蜂母亲又一次带着新食物来到了洞里。

幼虫的发育期大约有两周的时间，在这段时间内，它的食物都会在它需要的时候准时送到，幼虫长得越大，母亲送食物的时间间隔越短。半个月后，为了满足这个贪吃者的食欲，母亲就要忙个手忙脚乱了，而幼虫在蜕下去的爪和腹部角质的环节等残骸中，蠢笨地拖动着它的肚子。母亲不停地把猎物运回来，然后再不停地出去捕猎新的猎物。总之，泥蜂给幼虫的食物不是事先积存下来的，而是每天出去捕猎，再回来饲养它的幼虫的，就像鸟一样，母鸟喂养雏鸟的食物就是一口一口带回来的。有很多证据都证明，对于这些为幼虫捕捉猎物的膜翅目昆虫来说，这种饲养方法是很不同寻常的。我在前面说过，泥蜂把卵产在蜂房里，蜂房里只有一只小小的双翅目昆虫，并且从来都只有一只，绝无例外。要证实这一点并不难，不必等待什么特殊时机，下面就是一个有力的证据。

膜翅目昆虫会为幼虫事先准备好食物，现在就让我们来看看它们的窝吧。如果我们在昆虫拖着猎物进入地洞的时候前来观察，就会发现蜂房里有很多猎物。它们已经开始储备食物了，但这个时候它们连卵还没产下呢，更不会有幼虫了，因为它们只有把食物准备妥当以后，才会产卵。卵产下以后，它们就把蜂房封存起来，并且再不回来了。因此，只有当幼虫的洞穴不再需要母亲巡逻时，我们才有可能在堆积着的食物旁边找到幼虫。相反，假如我们在一只泥蜂带着猎物回到洞穴的时候，去查看它的住所，就会发现洞穴里有一堆食物残渣，一只或大或小的幼虫就躺在这些残渣中间。所以，泥蜂母亲现在带来的这份粮食是为了让幼虫把它的盛宴继续下去，这顿盛宴已经持续很久了，母亲还会不断提供新的猎物，让盛宴进行下去。如果我们在幼虫发育的最后阶段进行观察（我常常能够如愿以偿），就会在一大堆食物残渣上看到一只又肥又壮的幼虫，直到这时，它的母亲还会为它带来新鲜食物。母亲不断地为幼虫供应食物，直到幼虫胖得浑身滚圆，体内溢满了葡萄酒似的肉粥，躺在吃剩的猎物的残翅断爪上拒绝进食时，它才会彻底离开蜂房。

母亲每次捕猎回来都只带一只双翅目昆虫进到洞里。根据幼虫发育结束后留在蜂房里的食物残屑数量，也许能知道泥蜂母亲一共为幼虫提供了多少只猎物，进而知道这只膜翅目昆虫自产卵以来一共回到窝里多少次。但不幸的是，留下的这些残羹剩饭已经在幼虫饥饿的时候给反复咀嚼了，大部分都已经无法辨认了。不过如果我们这时打开一个尚在发育中的幼虫的蜂房，就可以清点出食物的数量。因为这些猎物有的还十分完整或近乎完整，更常见的情况是，食物被咬成了一段一段的，但保存完好，完全可以辨认出来。用这种方法统计出来的数字虽然不是十分精准，但也大得令人震惊。泥蜂母亲要怎样四处奔劳，才能为幼虫提供这么多的食物啊！下面就是我观察到的一份清单。

9月末，朱尔泥蜂的幼虫有成虫的1/3大了，我在这只幼虫的身边发现了下面这些食物：6只弥寄蝇，其中2只是完好无损的，剩下4只仅剩残肢断臂；4只彩色蚜蝇，其中2只是完好的，另外2只已经残破；3只黑服弥寄蝇，全都是完整的，其中1只是母亲刚刚运回来的，我正是通过它的这次送食才找到地洞的；2只粉蝇，1只完整，1只被啃咬过；1只蜂虻，被压得粉身碎骨；2只中带弥寄蝇，已沦为碎片；此外，还有2只花粉蝇，也已被压碎成片。一共有20只猎物。这张菜单确实够丰富的；可是此时的幼虫还不足成虫的1/3大，由此推断，在幼虫的整个发育期，它进食的昆虫可能会达到60只之多。

要对这个庞大的数字进行核实十分容易：我可以代替泥蜂母亲，为它的幼虫送去母性的关爱，给它提供食物，让它吃得酒足饭饱。我把蜂房转移到了一个铺着一层细沙的纸盒里，并把幼虫放到这张沙土床上。幼虫的表皮十分稚嫩，我搬动它的时候必须格外小心。我把它母亲事先为它备好的食物全部搬过来，放在幼虫周围，连一点儿碎渣都没漏掉。最后我用双手捧着纸盒，小心翼翼地往家里走去，在长达几千米的路途中，我尽量避免任何震荡，以免把蜂房弄乱，进而伤害我的幼虫。假如有人看到我疲倦地走在尼姆\footnote{尼姆：法国南部朗格多克－鲁西莱大区加尔省省会。——译者注}的公路上，小心翼翼地捧着我此番长途跋涉的唯一成果——一只肚子里塞满了苍蝇的肮脏的小虫，肯定会取笑我，认为我不正常。

我顺利地到家了，幼虫像什么事也没发生一样，继续静静地吃着它的双翅目昆虫。饲养到了第三天时，窝里的食物吃光了，幼虫把它那张尖嘴伸进猎物的残骸里搜寻着，可是没找到一丁点儿合胃口的东西；唯一抓住的那块东西十分坚硬，是角质的细片，于是，它气急败坏地扔掉了。现在，该由我为它提供食物了。我要提供给囚徒的食物就是可以信手拈来的双翅目昆虫。我把这些虫子用手指捏死，但没有捏碎。第一份粮食是3只尾蛆蝇和1只麻蝇。24小时内，它们全被吃光了。第二天，我给它提供了2只尾蛆蝇和4只苍蝇，这足够它当天食用了，但没有剩余。

我继续像这样每天给幼虫更丰盛些的食物。到了第九天，幼虫拒绝吃任何东西，开始结茧了。这八天盛宴吃的食物有62只虫，主要是尾蛆蝇和苍蝇，再加上原先在蜂房里找到的20只完整或残缺的猎物，总共为82只。

我在喂养幼虫的时候，大概不像它们的母亲那样注意卫生，分配食物时也不如它们的母亲更注重节制、审慎和节约；我把幼虫每天所需的食物一次性给出，完全由它们自由支配，这可能太浪费了。我必须承认，在某些情况下，在由它们的母亲照看的蜂房里，是不会发生这种事的。我的笔记如实地记载了下面的事实：在杜朗斯河冲积层的沙土中，我打开一个窝，大眼泥蜂刚带着一只麻蝇进去。在窝的尽头处，我找到了一只幼虫以及许多的昆虫残屑和几只完整的双翅目昆虫，其中也包括我亲眼看见泥蜂带进去的那一只。不过应该注意，占据这些猎物一半数量的麻蝇放在了蜂房的尽头，靠近幼虫的嘴边；而另外一半猎物还放在过道里，也就是蜂房的入口处，这样就不会被发育不完全的幼虫吃掉。因此，我认为当幼虫母亲捕猎的收获十分丰富时，就会把一些猎物暂时放在蜂房的门口，把那里当作一个临时仓库，当阴天下雨它无法正常工作时，就到那里取来一只。

它分配食物时十分节约，避免了无端的浪费，可我做不到这一点，也许我让幼虫吃得过于丰盛了。于是，我将捕猎数目减少到60只，并且只有中等个子的昆虫，比如苍蝇、麻蝇。如果昆虫的体积不是特别大的话，这大概就是幼虫母亲给幼虫提供的双翅目昆虫的数目。这个地区的所有泥蜂都是这样，只有铁爪泥蜂和带齿泥蜂例外，这两种泥蜂特别喜欢吃虻，数量从1只到12只不等，具体由双翅目昆虫的体积大小而定，因为各种虻的体积差别很悬殊。

下面分别描述一下我此次在六种泥蜂窝里观察到的双翅目昆虫的情况，以结束对食物性质的介绍。

1.橄榄树泥蜂。这种泥蜂我只在卡瓦伊翁见过一次，它的食物是叉叶绿蝇。以下五种泥蜂在阿维尼翁都十分常见。

2.大眼泥蜂。它通常把卵产在双翅目昆虫的身上，它为幼虫捕猎的食物有厩蜇蝇、红粉蝇、蚜蝇、圆形丽蝇、弥寄蝇、麻蝇、家蝇。其中尤以厩蜇蝇最常见，我曾在一个窝里就发现了五六十只。

3.跗节泥蜂。这种泥蜂也把卵产在双翅目昆虫身上，也捕猎卵蜂虻、蜂虻、厄奈斯尾蛆蝇、墓地尾蛆蝇、拟蜂蝇、蚜蝇。其中，蜂虻和卵蜂虻是最受它喜爱的食物。

4.朱尔泥蜂（9月、10月）。它的卵是产在粉蝇上的，主要以蚜蝇、红色弥寄蝇、中带弥寄蝇、黑服弥寄蝇、花粉蝇、红粉蝇、葡萄树蠹蛾、叉叶绿蝇、长足弥寄蝇、蜂虻为食。

5.铁爪泥蜂。这种泥蜂最爱吃虻，通常把卵产在蚜蝇、叉叶绿蝇身上，此后，它专门为幼虫捕捉各种肥硕的虻属的猎物。

6.带齿泥蜂。这种泥蜂也喜欢捕捉虻。我从未见它吃过别的猎物，也不知道它还把卵产在虻之外的什么双翅目昆虫身上。

食物的多样化说明泥蜂并不是专门吃某种昆虫，在捕猎时遇到哪种双翅目昆虫它就捕捉来。但是，它好像也有稍微偏爱一些的食物，比如第一种泥蜂特别爱吃叉叶绿蝇，第二种最喜欢厩蜇蝇，第三种和第五种则比较喜欢虻。

\begin{figure}[htbp]
    \centering
    \includegraphics[width=0.8\textwidth]{images/00397.jpeg}
    \caption{红腰泥蜂：膜翅目泥蜂科昆虫，主要捕食寄生昆虫。}
\end{figure}

\begin{figure}[htbp]
    \centering
    \includegraphics[width=0.8\textwidth]{images/00398.jpeg}
    \caption{叉叶绿蝇：又称金蝇，为丽蝇科的一属，分布广泛，通常在肉类或腐肉中繁殖。}
\end{figure}

\begin{figure}[htbp]
    \centering
    \includegraphics[width=0.8\textwidth]{images/00399.jpeg}
    \caption{姬蜂：膜翅目姬蜂科昆虫。分布广泛。形状、颜色、大小各异，许多种的腹部细长而弯曲，像黄蜂，触角较长，节较多。产卵器通常比身体长。}
\end{figure}

\begin{figure}[htbp]
    \centering
    \includegraphics[width=0.8\textwidth]{images/00400.jpeg}
    \caption{朱尔泥蜂：体长17～23毫米。大腮底部是黄色，触角为黑色。分布在欧洲中部、非洲北部等地。}
\end{figure}

\begin{figure}[htbp]
    \centering
    \includegraphics[width=0.8\textwidth]{images/00401.jpeg}
    \caption{铁爪泥蜂：膜翅目蜂类。体长17～24毫米。捕捉蝇或牛虻等双翅目昆虫作为幼虫的食物。高速飞行时善于在空中急停或急回转。}
\end{figure}

\begin{figure}[htbp]
    \centering
    \includegraphics[width=0.8\textwidth]{images/00402.jpeg}
    \caption{盗虻：双翅目盗虻科昆虫。体细长，脚长而有细毛。头顶在复眼间向下凹陷。触角3节，末端有一端刺。}
\end{figure}

\backmatter

\end{document}